% Jim Hefferon, Linear Algebra, pinned source revision df2262e089a02651c127f1dd12649c4622ee1383
% Selected foundation sections with all supplied answers in their native source.
% CC BY-SA 2.5 option; original component notices are in the full source ZIP.
% Assembly/public HTML rebuild: OpenAI Codex — GPT-6 Astra, Ultra effort.
% This is not an Everyday-English rewrite. No PDF or new TeX compilation is claimed.
% Extract native-authority.zip; put this file in the original src directory.
% Original styles, graphics, macros and make rules remain in that source archive.
% HTML byte replay uses the supplied frozen, editable Pandoc AST/MathML/SVG inputs.
\documentclass[twoside]{book}
\usepackage{verbatim}
\usepackage{xr}
\usepackage[single,write]{bookans}
\usepackage{bookjh}
% The bibliography remains a native dependency inside native-authority.zip.
\AtEndDocument{\input{bib/bib}}
\begin{document}
\mainmatter
\pagestyle{bookbody}

% BEGIN NATIVE SOURCE src/gr/gr1.tex
% Chapter 1, Section 1 _Linear Algebra_ Jim Hefferon
%  http://joshua.smcvt.edu/linearalgebra
%  2001-Jun-09
\chapter{Linear Systems}\leavevmode  % because no preamble?
\section{Solving Linear Systems}
Systems of linear equations are common in science and mathematics.
These two examples from high school science \cite{Onan}
give a sense of how they arise.

The first example is from 
\hypertarget{ex:Statics}{Statics}.\index{Statics problem}
Suppose that we have three objects,
we know that one has a mass of 2~kg,
and we want to find the two unknown masses.
Experimentation with a meter stick produces these two balances.
\begin{center}
  \includegraphics{gr/mp/ch1.1}
  \qquad
  \includegraphics{gr/mp/ch1.2}
\end{center}
For the masses to balance we must have that
the sum of moments on the left equals the sum of moments on
the right, where the moment of an object is its mass times its distance 
from the balance point. 
That gives a system of two linear equations.
\begin{align*}  % odd formatting to match chem problem below
      40h + 15c &= 100  \\
            25c &= 50+50h
\end{align*}

The second example
is from Chemistry.\index{Chemistry problem}
We can mix, under controlled conditions, toluene $\hbox{C}_7\hbox{H}_8$ and 
nitric acid $\hbox{H}\hbox{N}\hbox{O}_3$ to produce
trinitrotoluene $\hbox{C}_7\hbox{H}_5\hbox{O}_6\hbox{N}_3$
along with the byproduct water
(conditions have to be very well controlled\Dash trinitrotoluene 
is better known as TNT).
In what proportion should we mix them?
The number of atoms of each element present before the reaction
\begin{equation*}
    x\, {\rm C}_7{\rm H}_8\ +\ y\,{\rm H}{\rm N}{\rm O}_3
    \quad\longrightarrow\quad
    z\,{\rm C}_7{\rm H}_5{\rm O}_6{\rm N}_3\ +\ w\,{\rm H}_2{\rm O}
\tag*{}\end{equation*}
must equal the number present afterward.
Applying that in turn to the elements C, H, N, and O gives
this system.
\begin{align*}  % odd formatting to state more naturally
      7x      &= 7z  \\
      8x +1y  &= 5z+2w  \\
      1y      &= 3z  \\
      3y      &= 6z+1w
\end{align*}

Both examples come down to solving a system of equations.
In each system, the equations involve only the first power of each variable.
This chapter shows how to solve any such system of equations.















\subsection{Gauss's Method}
\begin{definition}\label{def:linearcombination}
%<*df:linearcombination>
A \definend{linear combination}\index{linear combination} of
\( x_1 \), \ldots, \( x_n \) has the form
\begin{equation*}
   a_1x_1+a_2x_2+a_3x_3+\cdots+a_nx_n
\end{equation*}
where the numbers \( a_1, \ldots ,a_n\in\Re \) are the combination's
\definend{coefficients}\index{linear equation!coefficients}.
%</df:linearcombination>
%<*df:linearequations>
A \definend{linear equation}\index{linear equation} 
in the variables $x_1$, \ldots, $x_n$ 
has the form
$a_1x_1+a_2x_2+a_3x_3+\cdots+a_nx_n=d$
where
\( d\in\Re \) is the \definend{constant}\index{linear equation!constant}.

An \( n \)-tuple \( (s_1,s_2,\ldots ,s_n)\in\Re^n \) is a 
\definend{solution}\index{linear equation!solution of} %
of, or \definend{satisfies}, that equation if substituting the numbers
$s_1$, \ldots, $s_n$ for the variables
gives a true statement:
$a_1s_1+a_2s_2+\cdots+a_ns_n=d$.
A \definend{system of linear equations}\index{linear equation!system of}%
\index{system of linear equations} 
\begin{equation*}
  \begin{linsys}{4}
    a_{1,1}x_1 &+ &a_{1,2}x_2  &+  &\cdots &+ &a_{1,n}x_n &=  &d_1  \\
    a_{2,1}x_1 &+ &a_{2,2}x_2  &+  &\cdots &+ &a_{2,n}x_n &=  &d_2  \\
             &  &           &   &       &  &          &\vdotswithin{=}  \\
    a_{m,1}x_1 &+ &a_{m,2}x_2  &+  &\cdots &+ &a_{m,n}x_n &=  &d_m
  \end{linsys}
\end{equation*}
has the solution
\( (s_1,s_2,\ldots ,s_n) \) if that $n$-tuple is a solution of all
of the equations.
%</df:linearequations>
\end{definition}

\begin{example}
The combination \( 3x_1 + 2x_2 \) of $x_1$ and $x_2$ is linear.
The combination \( 3x_1^2 + 2x_2 \) is not a linear function of
$x_1$ and $x_2$, nor is \( 3x_1 + 2\sin(x_2) \).  
\end{example}

We usually take $x_1$, \ldots, $x_n$ to be unequal to each other because
in a sum with repeats we can rearrange to 
make the elements unique, as with $2x+3y+4x=6x+3y$. 
We sometimes include terms with a zero coefficient, as in
$x-2y+0z$, and at other times omit them, depending on what is convenient. 

\begin{example}
The ordered pair \( (-1,5) \) is a solution of this system.
\begin{equation*}
  \begin{linsys}{2}
    3x_1 &+ &2x_2 &= &7  \\
    -x_1 &+ &x_2  &= &6
  \end{linsys}
\end{equation*}
In contrast, \( (5,-1) \) is not a solution.
\end{example}

Finding the set of all solutions is 
\definend{solving}\index{system of linear equations!solving} 
the system.
We don't need guesswork or good luck; 
there is an algorithm that always works.
This algorithm is  
\definend{Gauss's Method}\index{Gauss's Method}%
\index{system of linear equations!Gauss's Method} 
(or \definend{Gaussian elimination}\index{Gaussian elimination}%
\index{system of linear equations!Gaussian elimination}
or \definend{linear elimination}\index{linear elimination}%
\index{system of linear equations!linear elimination}%
\index{system of linear equations!elimination}%
\index{elimination, Gaussian}).
% It transforms the system, step by step, into one
% with a form that we can easily solve.
% We will first illustrate how it goes and then we will see the 
% formal statement. 

\begin{example}
To solve this system
\begin{equation*}
  \begin{linsys}{3}
                     &   &      &   &3x_3  &=  &9  \\
                 x_1 &+  &5x_2  &-  &2x_3  &=  &2  \\
      \frac{1}{3}x_1 &+  &2x_2  &   &      &=  &3  
  \end{linsys}
\end{equation*}
we transform it, step by step, until it is in a form that
we can easily solve.

The first transformation
rewrites the system by interchanging the first and third row.
\begin{align*}
  \quad
  &\grstep{ \text{swap row 1 with row 3} }
  \begin{linsys}{3}
        \frac{1}{3}x_1 &+  &2x_2  &   &      &=  &3  \\
                   x_1 &+  &5x_2  &-  &2x_3  &=  &2  \\
                       &   &      &   &3x_3  &=  &9  
    \end{linsys}                                         
\end{align*}
The second transformation rescales the first row by a factor of~$3$.
\begin{align*}
\quad
  &\grstep{ \text{multiply row 1 by 3} }
  \begin{linsys}{3}
        x_1 &+  &6x_2  &   &      &=  &9  \\
        x_1 &+  &5x_2  &-  &2x_3  &=  &2  \\
            &   &      &   &3x_3  &=  &9  
   \end{linsys}                                          
\end{align*}
The third transformation is the only nontrivial one in this example.
We mentally multiply both sides of the first row by \( -1 \),
mentally add that to the second row,
and write the result in as the new second row.
\begin{align*}
  &\grstep{ \text{add \(-1\) times row 1 to row 2} }
  \begin{linsys}{3}
        x_1 &+  &6x_2  &   &      &=  &9  \\
            &   &-x_2  &-  &2x_3  &=  &-7 \\
            &   &      &   &3x_3  &=  &9  
   \end{linsys}
\end{align*}
These steps have brought the system to a
form where we can easily find the value of each variable.
The bottom equation shows that \( x_3=3 \).
Substituting $3$ for \( x_3 \) in the middle equation shows that \( x_2=1 \).
Substituting those two into the top equation
gives that \( x_1=3 \). 
Thus the system has a unique solution; 
the solution set is \set{(3,1,3)}.
\end{example}

We will use Gauss's Method throughout the book.
It is fast and easy.
We will now show that
it is also safe: 
Gauss's Method never loses solutions nor does it ever 
pick up extraneous solutions, so that
a tuple is a solution to the system before we apply the method if and only if
it is a solution after.

\begin{theorem}[Gauss's Method]
\index{linear equation!solution of!Gauss's Method} \label{th:GaussMethod}
%<*th:GaussMethod>
If a linear system is changed to another by one of these operations
\begin{enumerate}
  \setlength{\itemsep}{0ex}
  \item
    an equation is swapped with another
  \item
    an equation has both sides multiplied by a nonzero constant
  \item
    an equation is replaced by the sum of itself and a multiple of another
\end{enumerate}
then the two systems have the same set of solutions.
%</th:GaussMethod>
\end{theorem}

Each of the three operations has a restriction.
Multiplying a row by~\( 0 \) is not allowed because obviously that
can change the solution set.
Similarly, adding a multiple of a row to itself is not allowed because
adding~\( -1 \) times the row to itself has the effect of multiplying the row
by~\( 0 \).
And we disallow swapping a row with itself,
to make some results in the fourth chapter easier.
Besides, it's pointless.

\begin{proof}
We will cover the equation swap operation here. 
The other two
cases are similar and are \nearbyexercise{ex:ProveGaussMethod}.

%<*pf:GaussMethod0>
Consider a linear system.
\begin{equation*}
  \begin{linsys}{4}
    a_{1,1}x_1  &+  &a_{1,2}x_2 &+  &\cdots  &+ &a_{1,n}x_n  &=  &d_1\hfill\hbox{}  \\
                &   &           &   &        &   &     &\vdotswithin{=}  \\
    a_{i,1}x_1  &+  &a_{i,2}x_2 &+  &\cdots  &+ &a_{i,n}x_n  &=  &d_i\hfill\hbox{}  \\
                &   &           &   &        &   &     &\vdotswithin{=}  \\
    a_{j,1}x_1  &+  &a_{j,2}x_2 &+  &\cdots  &+ &a_{j,n}x_n  &=  &d_j\hfill\hbox{}  \\
                &   &           &   &        &   &     &\vdotswithin{=} \\
    a_{m,1}x_1  &+  &a_{m,2}x_2 &+  &\cdots  &+ &a_{m,n}x_n  &=  &d_m  \\
  \end{linsys}
\end{equation*} 
The tuple \( (s_1,\ldots\,,s_n) \)
satisfies this system  
if and only if substituting the values for the
variables, the $s$'s for the $x$'s, gives a conjunction of true statements:
$a_{1,1}s_1+a_{1,2}s_2+\cdots+a_{1,n}s_n=d_1$
and \ldots\ 
$a_{i,1}s_1+a_{i,2}s_2+\cdots+a_{i,n}s_n=d_i$
and \ldots\  $a_{j,1}s_1+a_{j,2}s_2+\cdots+a_{j,n}s_n=d_j$
and \ldots\  $a_{m,1}s_1+a_{m,2}s_2+\cdots+a_{m,n}s_n=d_m$.
%</pf:GaussMethod0>

%<*pf:GaussMethod1>
In a list of statements joined with `and' we can  
rearrange the order of the statements. 
Thus
this requirement is met if and only if
$a_{1,1}s_1+a_{1,2}s_2+\cdots+a_{1,n}s_n=d_1$
and \ldots\  $a_{j,1}s_1+a_{j,2}s_2+\cdots+a_{j,n}s_n=d_j$
and \ldots\  $a_{i,1}s_1+a_{i,2}s_2+\cdots+a_{i,n}s_n=d_i$
and \ldots\  $a_{m,1}s_1+a_{m,2}s_2+\cdots+a_{m,n}s_n=d_m$.
This is exactly the requirement that \( (s_1,\ldots\,,s_n) \) 
solves the system after the row swap.
%</pf:GaussMethod1>
\end{proof}

\begin{definition} \label{df:GaussMethod}
%<*df:GaussMethod>
The three operations from 
\nearbytheorem{th:GaussMethod}
are the
\definend{elementary reduction 
operations},\index{Gauss's Method!elementary operations}%
\index{elementary reduction operations}
or \definend{row operations}\index{elementary row operations},
or \definend{Gaussian operations}.
They are
\definend{swapping}\index{elementary reduction operations! swapping}%
\index{swapping rows},
\definend{multiplying by a scalar} (or
\definend{rescaling}\index{elementary reduction operations! rescaling}%
\index{rescaling rows}), and
\definend{row combination}\index{elementary reduction operations! row combination}%
\index{combining rows}\index{adding rows}.
%</df:GaussMethod>
\end{definition}

When writing out the calculations, we will 
abbreviate `row \(i\)' by `\( \rho_i \)'
(this is the Greek letter rho, pronounced aloud as ``row'').
For instance, we will denote a row combination operation by 
\( k\rho_i+\rho_j \), 
with the row that changes written second.
To save writing we will 
often combine addition steps when they use the same $\rho_i$, as in the
next example.

\begin{example}
Gauss's Method systematically applies the row operations to solve a system.
Here is a typical case.
\begin{equation*}
  \begin{linsys}{3}
    x  &+  &y  &   &   &=  &0  \\
   2x  &-  &y  &+  &3z &=  &3  \\
    x  &-  &2y &-  &z  &=  &3  
  \end{linsys}
\end{equation*}
We begin by using the first row to 
eliminate the $2x$ in the second row and the $x$ in the third.
To get rid of the $2x$ we mentally multiply the entire first row by $-2$, 
add that to the
second row, and write the result in as the new second row.
To eliminate the~$x$ in the third row we multiply the first row by
$-1$, add that to the third row, and write the result in as the
new third row.
\begin{align*}
  &\grstep[-\rho_1 +\rho_3]{-2\rho_1 +\rho_2}
  \begin{linsys}{3}
     x  &+  &y  &   &   &=  &0  \\
        &   &-3y&+  &3z &=  &3  \\
        &   &-3y&-  &z  &=  &3  
  \end{linsys}   
\end{align*}
% In this version of the system, the last two equations involve only two unknowns.
We finish by transforming the second system into a third, where the
bottom equation involves only one unknown. 
We do that by using 
the second row to eliminate the $y$~term from the third row.
\begin{equation*}
  \grstep{-\rho_2 +\rho_3}
  \begin{linsys}{3}
     x  &+  &y  &   &   &=  &0  \\
        &   &-3y&+  &3z &=  &3  \\
        &   &   &   &-4z&=  &0
   \end{linsys}
\end{equation*}
Now finding the system's solution is easy.
The third row gives \( z=0 \).
Substitute that back\index{Gauss's Method!back-substitution}%
\index{back-substitution}
into the second row to get \( y=-1 \).
Then substitute back into the first row to get \( x=1 \).
\end{example}

\begin{example}
For the Physics problem\index{Statics problem} from the start of this
chapter, Gauss's Method gives this.
\begin{equation*}
   \begin{linsys}{2}
     40h  &+  &15c  &=  &100      \\
     -50h &+  &25c  &=  &50         
   \end{linsys}
   \grstep{5/4\rho_1 +\rho_2}
   \begin{linsys}{2}
      40h  &+  &15c       &=  &100      \\
           &   &(175/4)c  &=  &175 
    \end{linsys}
\end{equation*}
So \( c=4 \), and back-substitution gives that \( h=1 \).
(We will solve the Chemistry problem later.)
\end{example}

\begin{example}
The reduction
\begin{align*}
   \begin{linsys}{3}
        x  &+  &y  &+  &z  &=  &9  \\
       2x  &+  &4y &-  &3z &=  &1  \\
       3x  &+  &6y &-  &5z &=  &0  
   \end{linsys}
   &\grstep[-3\rho_1 +\rho_3]{-2\rho_1 +\rho_2}
   \begin{linsys}{3}
      x  &+  &y  &+  &z  &=  &9  \\
         &   &2y &-  &5z &=  &-17\\
         &   &3y &-  &8z&=  &-27
    \end{linsys}                                    \\
   &\grstep{-(3/2)\rho_2+\rho_3}
   \begin{linsys}{3}
      x  &+  &y  &+  &z            &=  &9  \\
         &   &2y &-  &5z           &=  &-17\\
         &   &   &   &-(1/2)z      &=  &-(3/2) 
    \end{linsys}
\end{align*}
shows that \( z=3 \), \( y=-1 \), and \( x=7 \).
\end{example}

As illustrated above, the point of Gauss's Method 
is to use the elementary reduction
operations to set up back-substitution.

\begin{definition} \label{df:EchelonForm}
%<*df:EchelonForm>
In each row of a system, 
the first variable with a nonzero coefficient is the row's
\definend{leading variable}\index{echelon form!leading variable}%
\index{leading!variable}. % 
A system is in \definend{echelon form}\index{echelon form}
if each leading variable
is to the right of the leading variable in the row above it,
except for the leading variable in the first row,
and any rows with all-zero coefficients are at the bottom.
%</df:EchelonForm>
\end{definition}

\begin{example} 
The prior three examples only used the operation of row combination.
This linear system requires the swap operation
to get it into echelon form because 
after the first combination
\begin{align*}
   \begin{linsys}{4}
                x  &-  &y  &   &   &   &   &=  &0  \\
               2x  &-  &2y &+  &z  &+  &2w &=  &4  \\
                   &   &y  &   &   &+  &w  &=  &0  \\
                   &   &   &   &2z &+  &w  &=  &5  
   \end{linsys}
   &\grstep{-2\rho_1 +\rho_2}
   \begin{linsys}{4}
      x  &-  &y  &\spaceforemptycolumn   &   &   &   &=  &0  \\
         &   &   &   &z  &+  &2w &=  &4  \\
         &   &y  &   &   &+  &w  &=  &0  \\
         &   &   &   &2z &+  &w  &=  &5  
   \end{linsys}    
\end{align*}
the second equation has no leading $y$.
We exchange it for a lower-down row that has a leading $y$.
\begin{align*}
   &\grstep{\rho_2 \leftrightarrow\rho_3}
   \begin{linsys}{4}
      x  &-  &y  &\spaceforemptycolumn   &   &   &   &=  &0  \\
         &   &y  &   &   &+  &w  &=  &0  \\
         &   &   &   &z  &+  &2w &=  &4  \\
         &   &   &   &2z &+  &w  &=  &5  
    \end{linsys}    
\end{align*}
(Had there been more than one suitable row below the second
then we could have used any one.)
With that, Gauss's Method proceeds as before.
\begin{align*}
   &\grstep{-2\rho_3 +\rho_4}
   \begin{linsys}{4}
      x  &-  &y  &\spaceforemptycolumn   &   &   &   &=  &0  \\
         &   &y  &   &   &+  &w  &=  &0  \\
         &   &   &   &z  &+  &2w &=  &4  \\
         &   &   &   &   &   &-3w&=  &-3 
    \end{linsys}
\end{align*}
Back-substitution gives \( w=1 \), \( z=2 \) , \( y=-1 \), and \( x=-1 \).
\end{example}

Strictly speaking, to solve linear systems we don't need 
the row rescaling operation.  
We have introduced it here because it is convenient and because we will use it 
later in this chapter as part of a variation of Gauss's Method, 
the Gauss-Jordan Method.

All of the systems so far have the same number of equations as unknowns.
All of them have a solution and for all of them there is only one solution.
We finish this subsection by seeing
other things that can happen.

\begin{example} \label{ex:MoreEqsThanUnks}
This system 
has more equations than variables.
\begin{equation*}
    \begin{linsys}{2}
      x  &+  &3y  &=  &1  \\
     2x  &+  &y   &=  &-3 \\
     2x  &+  &2y  &=  &-2 
    \end{linsys}  
\end{equation*}
Gauss's Method helps us understand this system also, since this
\begin{align*}
    &\grstep[-2\rho_1 +\rho_3]{-2\rho_1 +\rho_2}
    \begin{linsys}{2}
       x  &+  &3y  &=  &1  \\
          &   &-5y &=  &-5 \\
          &   &-4y &=  &-4 
     \end{linsys}
\end{align*}
shows that one of the equations is redundant.
Echelon form
\begin{equation*}
    \grstep{-(4/5)\rho_2 +\rho_3}
    \begin{linsys}{2}
       x  &+  &3y  &=  &1  \\
          &   &-5y &=  &-5 \\
          &   &0   &=  &0  
     \end{linsys}
\end{equation*}
gives that \( y=1 \) and \( x=-2 \).
The `\( 0=0 \)' reflects the redundancy.
\end{example}

Gauss's Method is also useful on systems with more variables than equations.
The next subsection has many examples.

Another way that linear systems can differ from the examples shown above 
is that some linear systems do not have a unique solution.
This can happen in two ways.
The first is that a system can fail to have any solution at all.

\begin{example} \label{ex:MoreEqsThanUnksInconsis}
Contrast the system in the last example with this one.
\begin{equation*}
    \begin{linsys}{2}
      x  &+  &3y  &=  &1  \\
     2x  &+  &y   &=  &-3 \\
     2x  &+  &2y  &=  &0  
    \end{linsys}
    \grstep[-2\rho_1 +\rho_3]{-2\rho_1 +\rho_2}
    \begin{linsys}{2}
       x  &+  &3y  &=  &1  \\
          &   &-5y &=  &-5 \\
          &   &-4y &=  &-2
     \end{linsys}
\end{equation*}
Here the system is inconsistent:~no pair of numbers $(s_1,s_2)$ satisfies 
all three equations simultaneously.
Echelon form makes the inconsistency obvious.
\begin{equation*}
  \grstep{-(4/5)\rho_2 +\rho_3}
  \begin{linsys}{2}
     x  &+  &3y  &=  &1  \\
        &   &-5y &=  &-5 \\
        &   &0   &=  &2 
   \end{linsys}
\end{equation*}
The solution set is empty.
\end{example}

\begin{example}
The prior system has more equations than unknowns but
that is not what causes the inconsistency\Dash 
\nearbyexample{ex:MoreEqsThanUnks}
has more equations than unknowns and yet is consistent.
Nor is having more equations than unknowns necessary for
inconsistency, as we see with this inconsistent system that has the 
same number of equations as unknowns.
\begin{equation*}
  \begin{linsys}{2}
    x  &+  &2y  &=  &8  \\
   2x  &+  &4y  &=  &8  
  \end{linsys}
  \grstep{-2\rho_1 + \rho_2}
  \begin{linsys}{2}
     x  &+  &2y  &=  &8  \\
        &   &0   &=  &-8
   \end{linsys}
\end{equation*}
Instead, 
inconsistency has to do with the interaction of the left and right sides;
in the first system above the left side's second equation is twice the
first but the right side's second constant is not twice the first.
Later we will have more to say about dependencies between a system's 
parts.
\end{example}

The other way that 
a linear system can fail to have a unique solution, besides having no solutions,
is to have many solutions.

\begin{example}
In this system
\begin{equation*}
  \begin{linsys}{2}
    x  &+  &y   &=  &4  \\
   2x  &+  &2y  &=  &8  
  \end{linsys}
\end{equation*}
any pair of numbers satisfying the first equation also
satisfies the second.
The solution set 
% \( \{ (x,y)\suchthat x+y=4 \} \) 
\( \set{ (x,y)\suchthat x+y=4} \) 
is infinite; some example
member pairs are~$(0,4)$, $(-1,5)$, and $(2.5,1.5)$.

The result of applying Gauss's Method here contrasts with the prior example
because we do not get a contradictory equation.
\begin{equation*}
  \grstep{-2\rho_1 + \rho_2}
  \begin{linsys}{2}
    x  &+  &y   &=  &4  \\
       &   &0   &=  &0  
   \end{linsys}
\end{equation*}
\end{example}

Don't be fooled by that example:~a $0=0$ equation  
is not the signal that a system has many solutions.

\begin{example}  \label{ex:NoZerosInfManySols}
The absence of a \( 0=0 \) equation does not keep a system from having
many different solutions.
This system is in echelon form,
has no $0=0$, but has infinitely many solutions,
including $(0,1,-1)$, 
$(0,1/2,-1/2)$, $(0,0,0)$, and $(0,-\pi,\pi)$
(any triple whose first component is $0$ and whose second component is the 
negative of the third is a solution).
\begin{equation*}
  \begin{linsys}{3}
     x  &+  &y  &+  &z  &=  &0  \\
        &   &y  &+  &z  &=  &0  
   \end{linsys}
\end{equation*}

Nor does the presence of \( 0=0 \) mean that the system must have 
many solutions.
\nearbyexample{ex:MoreEqsThanUnks} shows that.
So does this system, which does not have 
any solutions at all despite that 
in echelon form it has a $0=0$ row.
\begin{align*}
  \begin{linsys}{3}
    2x  &   &   &-   &2z  &=  &6  \\
        &   &y  &+   &z   &=  &1  \\
    2x  &+  &y  &-   &z   &=  &7  \\
        &   &3y &+   &3z  &=  &0  
  \end{linsys}
  &\grstep{-\rho_1 +\rho_3}
  \begin{linsys}{3}
     2x  &\spaceforemptycolumn   &   &-   &2z  &=  &6  \\
         &   &y  &+   &z   &=  &1  \\
         &   &y  &+   &z   &=  &1  \\
         &   &3y &+   &3z  &=  &0  
  \end{linsys}                                    \\
  &\grstep[-3\rho_2 +\rho_4]{-\rho_2 +\rho_3}
  \begin{linsys}{3}
     2x  &\spaceforemptycolumn    &   &-   &2z  &=  &6  \\
         &   &y  &+   &z   &=  &1  \\
         &   &   &    &0   &=  &0  \\
         &   &   &    &0   &=  &-3 
   \end{linsys}
\end{align*}
\end{example}

In summary,
Gauss's Method uses the row operations to 
set a system up for back substitution.
If any step shows a contradictory equation then we can stop with the
conclusion that the system has no solutions.
If we reach echelon form without a contradictory equation,
and each variable is a leading variable in its
row, then the system has a unique solution and we find it by
back substitution.
Finally, if we reach echelon form without a contradictory equation,
and there is not a unique solution\Dash
that is, at least one variable is not a leading variable\Dash
then the system has many solutions.

The next subsection explores the third case.
We will see that such a system must have infinitely many solutions
and we will describe the
solution set.


\medskip
\noindent\textbf{Note.}\hspace*{.2em}
\textit{In the exercises here, and in the rest of the book,
you must justify all of your answers.
For instance, if a question asks whether a system has a solution then you
must justify a yes response by producing the solution and must justify 
a no response by showing that no solution exists.}
\begin{exercises}
  \recommended \item 
    Use Gauss's Method to find the unique solution for each system.
    \begin{exparts*}
      \partsitem 
        $\begin{linsys}[t]{2}
          2x  &+  &3y  &=  &13  \\
          x   &-  &y   &=  &-1
        \end{linsys}$
      \partsitem 
        $\begin{linsys}[t]{3}
          x   &  &  &-  &z  &=  &0  \\
          3x  &+ &y &   &   &=  &1  \\
          -x  &+ &y &+  &z  &=  &4
        \end{linsys}$
    \end{exparts*}
    \begin{answer}
      \begin{exparts}
        \partsitem Gauss's Method
          \begin{equation*}
            \grstep{-(1/2)\rho_1+\rho_2}
            \begin{linsys}{2}
               2x  &+  &3y      &=  &13  \\
                   &-  &(5/2)y  &=  &-15/2              
            \end{linsys}
          \end{equation*}
          gives that the solution is $y=3$ and $x=2$.
        \partsitem Gauss's Method here
          \begin{equation*}
            \grstep[\rho_1+\rho_3]{-3\rho_1+\rho_2}
            \begin{linsys}{3}
              x   &  &  &-  &z  &=  &0  \\
                  &  &y &+  &3z &=  &1  \\
                  &  &y &   &   &=  &4
            \end{linsys}
            \grstep{-\rho_2+\rho_3}
            \begin{linsys}{3}
              x   &  &  &-  &z    &=  &0  \\
                  &  &y &+  &3z   &=  &1  \\
                  &  &  &   &-3z  &=  &3
            \end{linsys}
          \end{equation*}
          gives $x=-1$, $y=4$, and $z=-1$.
      \end{exparts}
    \end{answer}
  \item
    Each system is in echelon form.
    For each, say whether the system has a unique solution, 
    no solution, or infinitely many solutions.
    \begin{exparts*}
      \partsitem 
        \(\begin{linsys}[t]{2}
           -3x &+ &2y   &= &0  \\          
               &  &-2y  &= &0            
        \end{linsys}\)
      \partsitem 
        \(\begin{linsys}[t]{3}
             x &+ &y &  &   &= &4  \\          
               &  &y &- &z  &= &0  \\           
        \end{linsys}\)
      \partsitem 
        \(\begin{linsys}[t]{3}
             x &+ &y &  &   &= &4  \\          
               &  &y &  &-z  &= &0  \\ 
               &  &  &  &0  &= &0  
        \end{linsys}\)
      \partsitem 
        \(\begin{linsys}[t]{2}
             x &+ &y    &= &4  \\          
               &  &0    &= &4            
        \end{linsys}\)
      \partsitem 
        \(\begin{linsys}[t]{3}
            3x &+ &6y &+  &z  &= &-0.5  \\          
               &  &   &   &-z &= &2.5  \\           
        \end{linsys}\)
      \partsitem 
        \(\begin{linsys}[t]{2}
             x &- &3y    &= &2  \\          
               &  &0    &= &0            
        \end{linsys}\)
      \partsitem 
        \(\begin{linsys}[t]{2}
             2x &+ &2y    &= &4  \\
                &  &y     &= &1 \\          
                 &  &0    &= &4            
        \end{linsys}\)
      \partsitem 
        \(\begin{linsys}[t]{2}
              2x &+ &y   &= &0  
        \end{linsys}\)
      \partsitem 
        \(\begin{linsys}[t]{2}
              x &- &y    &= &-1  \\
                &  &0     &= &0 \\          
                 &  &0    &= &4            
        \end{linsys}\)
      \partsitem 
        \(\begin{linsys}[t]{3}
              x &+ &y  &- &3z  &= &-1  \\
                &  &y  &- &z  &= &2  \\
                &  &  &  &z    &= &0  \\
                &  &  &  &0   &= &0  
        \end{linsys}\)
    \end{exparts*}
    \begin{answer}
      If a system has a contradictory equation then it has no solution.
      Otherwise, if there are any variables that are not leading a row
      then it has infinitely many solution.
      In the final case, where there is no contradictory equation and
      every variable leads some row, it has a unique solution.
      \begin{exparts}
        \partsitem Unique solution
        \partsitem Infinitely many solutions
        \partsitem Infinitely many solutions
        \partsitem No solution
        \partsitem Infinitely many solutions
        \partsitem Infinitely many solutions
        \partsitem No solution
        \partsitem Infinitely many solutions
        \partsitem No solution
        \partsitem Unique solution
      \end{exparts}
    \end{answer}
  \recommended \item  
    Use Gauss's Method to solve each system
    or conclude `many solutions' or `no solutions'.
    \begin{exparts*}
      \partsitem \(
               \begin{linsys}[t]{2}
                  2x  &+  &2y  &=  &5  \\
                   x  &-  &4y  &=  &0  
               \end{linsys}
             \)
      \partsitem \(
               \begin{linsys}[t]{2}
                  -x  &+  &y   &=  &1  \\
                   x  &+  &y   &=  &2  
               \end{linsys}
             \) 
      \partsitem  \(
               \begin{linsys}[t]{3}
                   x  &-  &3y  &+  &z  &=  &1  \\
                   x  &+  &y   &+  &2z &=  &14 
                \end{linsys}
             \) 
      \partsitem  \(
               \begin{linsys}[t]{2}
                  -x  &-  &y   &=  &1  \\
                 -3x  &-  &3y  &=  &2  
               \end{linsys}
             \) 
      \partsitem  \(
               \begin{linsys}[t]{3}
                      &   &4y  &+  &z  &=  &20 \\
                  2x  &-  &2y  &+  &z  &=  &0  \\
                   x  &   &    &+  &z  &=  &5  \\
                   x  &+  &y   &-  &z  &=  &10 
                \end{linsys}
             \)
      \partsitem \( \begin{linsys}[t]{4}
                 2x  &   &   &+  &z  &+  &w  &=  &5  \\
                     &   &y  &   &   &-  &w  &=  &-1 \\
                 3x  &   &   &-  &z  &-  &w  &=  &0  \\
                 4x  &+  &y  &+  &2z &+  &w  &=  &9  
               \end{linsys}
            \)
    \end{exparts*}
    \begin{answer} 
      \begin{exparts}
       \partsitem Gaussian reduction
        \begin{equation*}
          \grstep{-(1/2)\rho_1+\rho_2}
          \begin{linsys}{2}
             2x  &+  &2y  &=  &5  \\
                 &   &-5y &=  &-5/2  
          \end{linsys}
        \end{equation*}
        shows that \( y=1/2 \) and \( x=2 \) is the unique solution.
      \partsitem Gauss's Method
        \begin{equation*}
          \grstep{\rho_1+\rho_2}
          \begin{linsys}{2}
             -x  &+  &y   &=  &1  \\
                 &   &2y  &=  &3  
           \end{linsys}
        \end{equation*}
        gives \( y=3/2 \) and \( x=1/2 \) as the only solution.
      \partsitem Row reduction
        \begin{equation*}
            \grstep{-\rho_1+\rho_2}
            \begin{linsys}{3}
                x  &-  &3y  &+  &z  &=  &1  \\
                   &   &4y  &+  &z  &=  &13 
             \end{linsys}
        \end{equation*}
        shows, because the variable $z$ is not a leading variable in any
        row, that there are many solutions.
      \partsitem Row reduction
        \begin{equation*}
          \grstep{-3\rho_1+\rho_2}
          \begin{linsys}{2}
             -x  &-  &y   &=  &1  \\
                 &   &0   &=  &-1 
           \end{linsys}
        \end{equation*}
        shows that there is no solution.
      \partsitem Gauss's Method
        \begin{align*}
            \grstep{\rho_1\leftrightarrow\rho_4}
            \begin{linsys}{3}
                x  &+  &y   &-  &z  &=  &10 \\
               2x  &-  &2y  &+  &z  &=  &0  \\
                x  &   &    &+  &z  &=  &5  \\
                   &   &4y  &+  &z  &=  &20 
             \end{linsys}
            &\grstep[-\rho_1+\rho_3]{-2\rho_1+\rho_2}
            \begin{linsys}{3}
                x  &+  &y   &-  &z  &=  &10 \\
                   &   &-4y &+  &3z &=  &-20\\
                   &   &-y  &+  &2z &=  &-5 \\
                   &   &4y  &+  &z  &=  &20 
             \end{linsys}                                      \\
            &\grstep[\rho_2+\rho_4]{-(1/4)\rho_2+\rho_3}
            \begin{linsys}{3}
                x  &+  &y   &-  &z      &=  &10 \\
                   &   &-4y &+  &3z     &=  &-20\\
                   &   &    &   &(5/4)z &=  &0  \\
                   &   &    &   &4z     &=  &0  
             \end{linsys}
        \end{align*}
        gives the unique solution \( (x,y,z)=(5,5,0) \).
      \partsitem Here Gauss's Method gives
         \begin{align*}
            &\grstep[-2\rho_1+\rho_4]{-(3/2)\rho_1+\rho_3}
            \begin{linsys}{4}
               2x  &\spaceforemptycolumn   &   &+  &z       &+  &w       &=  &5  \\
                   &   &y  &   &        &-  &w       &=  &-1 \\
                   &   &   &-  &(5/2)z  &-  &(5/2)w  &=  &-15/2  \\
                   &   &y  &   &        &-  &w       &=  &-1  
             \end{linsys}                                             \\ 
            &\grstep{-\rho_2+\rho_4}
            \begin{linsys}{4}
               2x  &\spaceforemptycolumn   &   &+  &z       &+  &w       &=  &5  \\
                   &   &y  &   &        &-  &w       &=  &-1 \\
                   &   &   &-  &(5/2)z  &-  &(5/2)w  &=  &-15/2  \\
                   &   &   &   &        &   &0       &=  &0 
             \end{linsys}
         \end{align*}
         which shows that there are many solutions.
      \end{exparts} 
    \end{answer}
  \item  
    Solve each system
    or conclude `many solutions' or `no solutions'.
    Use Gauss's Method.
    \begin{exparts*}
      \partsitem \(
               \begin{linsys}[t]{3}
                   x  &+  &y   &+  &z   &=  &5  \\
                   x  &-  &y   &   &    &=  &0  \\
                      &   &y   &+  &2z  &=  &7  
               \end{linsys}
             \)
      \partsitem \(
               \begin{linsys}[t]{3}
                   3x  &  &    &+  &z   &=  &7  \\
                    x  &-  &y   &+   &3z    &=  &4  \\
                    x   &+  &2y   &-  &5z  &=  &-1  
               \end{linsys}
             \)
      \partsitem \(
               \begin{linsys}[t]{3}
                   x   &+  &3y   &+  &z   &=  &0  \\
                   -x  &-  &y   &   &    &=  &2  \\
                   -x   &+  &y   &+  &2z  &=  &8  
               \end{linsys}
             \)
    \end{exparts*}
    \begin{answer} 
      \begin{exparts}
       \partsitem 
         % sage: M = matrix(RDF, [[1,1,1], [1,-1,0], [0,1,2]])
         % sage: v = vector(RDF, [5,0,7])
         % sage: M_prime = M.augment(v)
         % sage: gauss_method(M_prime)
         % [ 1.0  1.0  1.0  5.0]
         % [ 1.0 -1.0  0.0  0.0]
         % [ 0.0  1.0  2.0  7.0]
         %  take -1.0 times row 1 plus row 2
         % [ 1.0  1.0  1.0  5.0]
         % [ 0.0 -2.0 -1.0 -5.0]
         % [ 0.0  1.0  2.0  7.0]
         %  take 0.5 times row 2 plus row 3
         % [ 1.0  1.0  1.0  5.0]
         % [ 0.0 -2.0 -1.0 -5.0]
         % [ 0.0  0.0  1.5  4.5]
         Gauss's Method
         \begin{align*}
          \begin{linsys}{3}
             x  &+  &y   &+  &z   &=  &5  \\
             x  &-  &y   &   &    &=  &0  \\
                &   &y   &+  &2z  &=  &7  
          \end{linsys}
          &\grstep{-\rho_1+\rho_2} 
          \begin{linsys}{3}
             x  &+  &y   &+  &z   &=  &5  \\
             0  &   &-2y  &-   &z    &=  &-5  \\
                &   &y   &+  &2z  &=  &7  
          \end{linsys}                                     \\
          &\grstep{(1/2)\rho_2+\rho_3} 
          \begin{linsys}{3}
             x  &+  &y   &+  &z   &=  &5  \\
             0  &   &-2y  &-   &1    &=  &-5  \\
                &   &    &+  &(3/2)z  &=  &7/2  
          \end{linsys}
         \end{align*}
         followed by back-substitution gives $x=1$, $y=1$, and $z=3$.
       \partsitem 
         % sage: M = matrix(QQ, [[3,0,1], [1,-1,3], [1,2,-5]])
         % sage: v = vector(RDF, [7,4,-1])
         % sage: v = vector(QQ, [7,4,-1])
         % sage: M_prime = M.augment(v)
         % sage: gauss_method(M_prime)
         % [ 3  0  1  7]
         % [ 1 -1  3  4]
         % [ 1  2 -5 -1]
         %  take -1/3 times row 1 plus row 2
         %  take -1/3 times row 1 plus row 3
         % [    3     0     1     7]
         % [    0    -1   8/3   5/3]
         % [    0     2 -16/3 -10/3]
         %  take 2 times row 2 plus row 3
         % [  3   0   1   7]
         % [  0  -1 8/3 5/3]
         % [  0   0   0   0]
         Here Gauss's Method
         \begin{align*}
           \begin{linsys}{3}
               3x  &  &    &+  &z   &=  &7  \\
                x  &-  &y   &+   &3z    &=  &4  \\
                x   &+  &2y   &-  &5z  &=  &-1  
           \end{linsys}
           &\grstep[-(1/3)\rho_1+\rho_3]{-(1/3)\rho_1+\rho_2} 
           \begin{linsys}{3}
               3x  &  &    &+  &z        &=  &7  \\
                   &  &-y   &+ &(8/3)z   &=  &5/3  \\
                   &  &2y   &- &(16/3)z  &=  &-(10/3)  
           \end{linsys}                                      \\
           &\grstep{2\rho_2+\rho_3} 
           \begin{linsys}{3}
               3x  &  &    &+  &z        &=  &7  \\
                   &  &-y   &+ &(8/3)z   &=  &5/3  \\
                   &  &     &  &0        &=  &0  
           \end{linsys}
       \end{align*}
       finds that the variable~$z$ does not lead a row. 
       There are infinitely many solutions.
       \partsitem
         % sage: M = matrix(QQ, [[1,3,1], [-1,-1,0], [-1,1,2]])
         % sage: v = vector(QQ, [0,2,8])
         % sage: M_prime = M.augment(v)
         % sage: gauss_method(M_prime)
         % [ 1  3  1  0]
         % [-1 -1  0  2]
         % [-1  1  2  8]
         %  take 1 times row 1 plus row 2
         %  take 1 times row 1 plus row 3
         % [1 3 1 0]
         % [0 2 1 2]
         % [0 4 3 8]
         %  take -2 times row 2 plus row 3
         % [1 3 1 0]
         % [0 2 1 2]
         % [0 0 1 4]
         The steps
         \begin{equation*}
           \begin{linsys}{3}
              x   &+  &3y   &+  &z   &=  &0  \\
              -x  &-  &y   &   &    &=  &2  \\
              -x   &+  &y   &+  &2z  &=  &8  
           \end{linsys}
           \grstep[\rho_1+\rho_3]{\rho_1+\rho_2}
           \begin{linsys}{3}
              x   &+  &3y   &+  &z   &=  &0  \\
                  &   &2y   &+  &z   &=  &2  \\
                  &   &4y   &+  &3z  &=  &8  
           \end{linsys}
           \grstep{-2\rho_2+\rho_3}
           \begin{linsys}{3}
              x   &+  &3y   &+  &z   &=  &0  \\
                  &   &2y   &+  &z   &=  &2  \\
                  &   &     &   &z   &=  &4  
           \end{linsys}
         \end{equation*}
         give $(x,y,z)=(-1,-1,4)$.
      \end{exparts} 
    \end{answer}
  \recommended \item 
    We can solve linear systems by methods other 
    than Gauss's.
    One often taught in high school is to solve one of the 
    equations for a variable, then substitute the resulting expression into
    other equations.
    Then we repeat that step until there is an equation with only one
    variable.
    From that we get the first number in the solution and then we get the
    rest with  
    back-substitution.
    This method takes longer than Gauss's Method, since it involves
    more arithmetic operations, and is also more
    likely to lead to errors.
    To illustrate how it can lead to wrong conclusions, we will use the system 
    \begin{equation*}
      \begin{linsys}{2}
            x  &+  &3y  &=  &1  \\
            2x  &+  &y   &=  &-3 \\
            2x  &+  &2y  &=  &0  
      \end{linsys}
    \end{equation*}
    from \nearbyexample{ex:MoreEqsThanUnksInconsis}.
    \begin{exparts}
      \partsitem Solve the first equation for $x$ and 
        substitute that expression into the second equation.
        Find the resulting $y$.
      \partsitem Again solve the first equation for $x$, 
        but this time substitute that expression into the third equation.
        Find this $y$.
    \end{exparts}
    What extra step must a user of this method take to avoid 
    erroneously concluding a system has a solution?
    \begin{answer}
      \begin{exparts}
        \partsitem From $x=1-3y$ we get that $2(1-3y)+y=-3$, giving $y=1$.
        \partsitem From $x=1-3y$ we get that $2(1-3y)+2y=0$, leading to 
           the conclusion that $y=1/2$.
      \end{exparts}
      Users of this method must check any potential solutions by
      substituting back into all the equations.
    \end{answer}
  \recommended \item 
    For which values of \( k \) are
    there no solutions, many solutions, or a unique solution
    to this system?
    \begin{equation*}
       \begin{linsys}{2}
          x  &-  &y  &=  &1  \\
         3x  &-  &3y &=  &k  
       \end{linsys}
    \end{equation*}
    \begin{answer}
      Do the reduction
      \begin{equation*}
        \grstep{-3\rho_1+\rho_2}
        \begin{linsys}{2}
          x  &-  &y  &=  &1\hfill  \\
             &   &0  &=  &-3+k\hfill  
        \end{linsys}
      \end{equation*}
      to conclude this system has no solutions if \( k\neq 3 \) and if
      \( k=3 \) then it has infinitely many solutions.
      It never has a unique solution.  
    \end{answer}
  \item 
    This system is not linear in that it says $\sin\alpha$ instead of $\alpha$
    \begin{equation*}
      \begin{linsys}{3}
         2\sin\alpha  &-  &\cos\beta  &+  &3\tan\gamma  &=  &3  \\
         4\sin\alpha  &+  &2\cos\beta &-  &2\tan\gamma  &=  &10  \\
         6\sin\alpha  &-  &3\cos\beta &+  &\tan\gamma   &=  &9  
      \end{linsys}
    \end{equation*}
    and yet we can apply Gauss's Method.
    Do so.
    Does the system have a solution?
    \begin{answer}
      Let \( x=\sin\alpha \), \( y=\cos\beta \), and \( z=\tan\gamma \):
      \begin{equation*}
        \begin{linsys}{3}
           2x  &-  &y  &+  &3z  &=  &3  \\
           4x  &+  &2y &-  &2z  &=  &10  \\
           6x  &-  &3y &+  &z   &=  &9  
        \end{linsys}
         \grstep[-3\rho_1+\rho_3]{-2\rho_1+\rho_2}
         \begin{linsys}{3}
           2x  &-  &y  &+  &3z  &=  &3  \\
               &   &4y &-  &8z  &=  &4   \\
               &   &   &   &-8z &=  &0  
         \end{linsys}
      \end{equation*}
      gives \( z=0 \), \( y=1 \), and \( x=2 \).
      Note that no \( \alpha\in\R \) satisfies that $\sin\alpha=2$.  
      \end{answer}
  \recommended \item 
    What conditions must the constants, the $b$'s,
    satisfy so that each of these systems has a solution?
    \textit{Hint.} 
    Apply Gauss's Method and see what happens to the right side.
    \begin{exparts*}
      \partsitem  \(
        \begin{linsys}[t]{2}
           x  &-  &3y  &=  &b_1 \\
          3x  &+  &y   &=  &b_2 \\
           x  &+  &7y  &=  &b_3 \\
          2x  &+  &4y  &=  &b_4 
        \end{linsys}   \)
      \partsitem \(
        \begin{linsys}[t]{3}
           x_1  &+  &2x_2  &+  &3x_3  &=  &b_1  \\
          2x_1  &+  &5x_2  &+  &3x_3  &=  &b_2  \\
           x_1  &   &      &+  &8x_3  &=  &b_3  
        \end{linsys}   \)
    \end{exparts*}
    \begin{answer} 
      \begin{exparts}
       \partsitem Gauss's Method
         \begin{equation*}
           \grstep[-\rho_1+\rho_3 \\ -2\rho_1+\rho_4]{-3\rho_1+\rho_2}
           \begin{linsys}{2}
              x  &-  &3y  &=  &b_1\hfill \\
                 &   &10y &=  &-3b_1+b_2\hfill \\
                 &   &10y &=  &-b_1+b_3\hfill \\
                 &   &10y &=  &-2b_1+b_4\hfill 
            \end{linsys}         
           \grstep[-\rho_2+\rho_4]{-\rho_2+\rho_3}
           \begin{linsys}{2}
              x  &-  &3y  &=  &b_1\hfill \\
                 &   &10y &=  &-3b_1+b_2\hfill \\
                 &   &0   &=  &2b_1-b_2+b_3\hfill \\
                 &   &0   &=  &b_1-b_2+b_4\hfill 
            \end{linsys}
         \end{equation*}
         shows that this system is consistent if and only if both
         \( b_3=-2b_1+b_2 \) and \( b_4=-b_1+b_2 \).
       \partsitem Reduction
         \begin{align*}
            &\grstep[-\rho_1+\rho_3]{-2\rho_1+\rho_2}
            \begin{linsys}{3}
              x_1  &+  &2x_2  &+  &3x_3  &=  &b_1\hfill  \\
                   &   &x_2   &-  &3x_3  &=  &-2b_1+b_2\hfill  \\
                   &   &-2x_2 &+  &5x_3  &=  &-b_1+b_3\hfill  
             \end{linsys}                                          \\
            &\grstep{2\rho_2+\rho_3}
            \begin{linsys}{3}
              x_1  &+  &2x_2  &+  &3x_3  &=  &b_1\hfill  \\
                   &   &x_2   &-  &3x_3  &=  &-2b_1+b_2\hfill  \\
                   &   &      &   &-x_3  &=  &-5b_1+2b_2+b_3\hfill  
             \end{linsys}
         \end{align*}
         shows that each of \( b_1 \), \( b_2 \), and \( b_3 \) can be any
         real number\Dash this system always has a unique solution.
      \end{exparts}   
      \end{answer}
  \item 
    True or false: a system with more unknowns than equations
    has at least one solution.
    (As always, to say `true' you must prove it, while to say 
    `false' you must produce a counterexample.)
    \begin{answer}
      This system with more unknowns than equations
      \begin{equation*}
        \begin{linsys}{3}
          x  &+  &y  &+  &z  &=  &0  \\
          x  &+  &y  &+  &z  &=  &1  
        \end{linsys}
      \end{equation*}
      has no solution.   
      \end{answer}
  \item 
    Must any Chemistry\index{Chemistry problem} problem like
    the one that starts this subsection\Dash
    a balance the reaction problem\Dash have infinitely many solutions?
    \begin{answer}
      Yes.
      For example, the fact that we can have the same reaction 
      in two different flasks shows that twice any solution is another,
      different, solution (if a physical reaction occurs then there must be
      at least one nonzero solution).
    \end{answer}
  \recommended \item 
    Find the coefficients
    \( a \), \( b \), and \( c \) so that the graph of \( f(x)=ax^2+bx+c \) 
    passes through the points \( (1,2) \), \( (-1,6) \), and \( (2,3) \).
    \begin{answer}
      Because \( f(1)=2 \), \( f(-1)=6 \), and \( f(2)=3 \) we get
      a linear system.
      \begin{equation*}
        \begin{linsys}{3}
          1a  &+  &1b  &+  &c  &=  &2  \\
          1a  &-  &1b  &+  &c  &=  &6  \\
          4a  &+  &2b  &+  &c  &=  &3  
         \end{linsys}
      \end{equation*}
      Gauss's Method
      \begin{equation*}
         \grstep[-4\rho_1+\rho_3]{-\rho_1+\rho_2}
         \begin{linsys}{3}
            a  &+  &b  &+  &c  &=  &2  \\
               &   &-2b&   &   &=  &4  \\
               &   &-2b&-  &3c &=  &-5 
          \end{linsys}               
         \grstep{-\rho_2+\rho_3}
         \begin{linsys}{3}
            a  &+  &b  &+  &c  &=  &2  \\
               &   &-2b&   &   &=  &4  \\
               &   &   &   &-3c&=  &-9 
           \end{linsys}
      \end{equation*}
      shows that the solution is \( f(x)=1x^2-2x+3 \).  
      \end{answer}
  \item After \nearbytheorem{th:GaussMethod} we note that multiplying a 
   row by~$0$ is not allowed because that could change a solution set. 
   Give an example of a system with solution set~$S_0$ where after 
   multiplying a row by~$0$ the new system has a solution set~$S_1$
   and $S_0$ is a proper subset of $S_1$, that is, $S_0\neq S_1$.
   Give an example where $S_0=S_1$. 
   \begin{answer}
     Here $S_0=\set{(1,1)}$
     \begin{equation*}
         \begin{linsys}{2}
            x  &+  &y  &=  &2  \\
            x  &-  &y  &=  &0
          \end{linsys}               
         \grstep{0\rho_2}
         \begin{linsys}{2}
            x  &+  &y  &=  &2  \\
               &   &0  &=  &0
          \end{linsys}               
     \end{equation*}
     while $S_1$ is a proper superset because it
     contains at least two points: $(1,1)$ and~$(2,0)$.
     In this example the solution set does not change.  
     \begin{equation*}
         \begin{linsys}{2}
            x  &+  &y  &=  &2  \\
           2x  &+  &2y &=  &4
          \end{linsys}               
         \grstep{0\rho_2}
         \begin{linsys}{2}
            x  &+  &y  &=  &2  \\
               &   &0  &=  &0
          \end{linsys}               
     \end{equation*}
   \end{answer}
  \item 
    Gauss's Method works by combining the equations in a system to make new
    equations.
    \begin{exparts}
      \partsitem Can we derive the equation \( 3x-2y=5 \) by a sequence of
        Gaussian reduction steps from the equations in this system?
        \begin{equation*}
          \begin{linsys}{2}
             x  &+  &y  &=  &1  \\
            4x  &-  &y  &=  &6
          \end{linsys}
        \end{equation*}
      \partsitem Can we derive the equation \( 5x-3y=2 \) with a sequence of
        Gaussian reduction steps from the equations in this system?
        \begin{equation*}
          \begin{linsys}{2}
            2x  &+  &2y &=  &5  \\
            3x  &+  &y  &=  &4
          \end{linsys}
        \end{equation*}
      \partsitem Can we derive \( 6x-9y+5z=-2 \)  
        by a sequence of
        Gaussian reduction steps from the equations in the system?
        \begin{equation*}
          \begin{linsys}{3}
            2x  &+  &y  &-  &z  &=  &4  \\
            6x  &-  &3y &+  &z  &=  &5
          \end{linsys}
        \end{equation*}
    \end{exparts}
    \begin{answer} 
       \begin{exparts} 
        \partsitem Yes, by inspection the given equation results from
          \( -\rho_1+\rho_2 \).
        \partsitem No.
          The pair \( (1,1) \) satisfies the given equation. 
          However, that pair 
          does not satisfy the first equation in the system.
        \partsitem Yes.
          To see if the given row is \( c_1\rho_1+c_2\rho_2 \), solve
          the system of equations relating the coefficients of $x$, $y$,
          $z$, and the constants:
          \begin{equation*}
            \begin{linsys}{2}
               2c_1  &+  &6c_2  &=  &6  \\
                c_1  &-  &3c_2  &=  &-9 \\
               -c_1  &+  &c_2   &=  &5  \\
               4c_1  &+  &5c_2  &=  &-2 
            \end{linsys}
          \end{equation*}
          and get $c_1=-3$ and $c_2=2$, so the given row is
          \( -3\rho_1+2\rho_2 \).
      \end{exparts}  
     \end{answer}
  \item 
    Prove that, where \( a,b,c,d,e \) are real numbers
    with \( a\neq 0 \), if this linear equation
    \begin{equation*}
       ax+by=c
    \end{equation*}
    has the same solution set as this one 
    \begin{equation*}
       ax+dy=e
    \end{equation*}
    then they are the same equation.
    What if \( a=0 \)?
    \begin{answer}
      If \( a\neq 0 \) then the solution set of the first equation is
      this.
      \begin{equation*}
        \set{(x,y)\in\Re^2\suchthat 
             \text{$x=(c-by)/a=(c/a)-(b/a)\cdot y$}}
        \tag{$*$} 
      \end{equation*}
      Thus, given~$y$ we can compute the associated~$x$.
      Taking $y=0$ gives the solution $(c/a,0)$, and since the second
      equation $ax+dy=e$ is supposed to have the same solution set, 
      substituting into
      it gives that $a(c/a)+d\cdot 0=e$, so $c=e$.
      Taking $y=1$ in ($*$) gives $a((c-b)/a)+d\cdot 1=e$,
      and so $b=d$.
      Hence they are the same equation.

      When \( a=0 \) the equations can be different and still have the 
      same solution set:~e.g.,
      \( 0x+3y=6 \) and \( 0x+6y=12 \).   
     \end{answer}
  \item 
    Show that if \( ad-bc\neq 0 \) then
    \begin{equation*}
      \begin{linsys}{2}
        ax  &+  &by  &=  &j  \\
        cx  &+  &dy  &=  &k  
      \end{linsys}
    \end{equation*}
    has a unique solution.
    \begin{answer}
      We take three cases: that $a\neq 0$, that $a=0$ and 
      $c\neq 0$, and that both $a=0$ and $c=0$.

      For the first, we assume that \( a\neq 0 \).
      Then the reduction
      \begin{equation*}
        \grstep{-(c/a)\rho_1+\rho_2}
        \begin{linsys}{2}
          ax  &+  &by                  &=  &j \hfill\hbox{} \\
              &   &(-(cb/a)+d)y  &=  &-(cj/a)+k \hfill  
         \end{linsys}
      \end{equation*}
      shows that this system has a unique solution if and only if
      \( -(cb/a)+d\neq 0   \); remember that \( a\neq 0 \) so 
      that back substitution yields a unique \( x \)
      (observe, by the way, that \( j \) and \( k \) play no role in the
      conclusion that there is a unique solution, although if there is a 
      unique solution then they contribute to its value).
      But \( -(cb/a)+d = (ad-bc)/a \) and a fraction is not equal to \( 0 \) 
      if and only if its numerator is not equal to \( 0 \).
      Thus, in this first case, there is a unique solution if and only if
      $ad-bc\neq 0$.

      In the second case, if \( a=0 \) but \( c\neq 0 \), then we swap
      \begin{equation*}
        \begin{linsys}{2}
          cx  &+  &dy  &=  &k  \\
              &   &by  &=  &j  
        \end{linsys}
      \end{equation*}
      to conclude that the system has a unique solution if and only if 
      \( b\neq 0 \)
      (we use the case assumption that \( c\neq 0 \) to get a unique
      \( x \) in back substitution).
      But\Dash where \( a=0 \) and \( c\neq 0 \)\Dash
      the condition ``\( b\neq 0 \)''
      is equivalent to the condition ``\( ad-bc\neq 0 \)''.
      That finishes the second case.

      Finally, for the third case,
      if both \( a \) and \( c \) are \( 0 \) then the system
      \begin{equation*}
        \begin{linsys}{2}
          0x  &+  &by  &=  &j  \\
          0x  &+  &dy  &=  &k  
        \end{linsys}
      \end{equation*}
      might have no solutions (if the second equation is not a multiple of the
      first) or it might have infinitely many solutions (if the second
      equation is a multiple of the first then for each \( y \) satisfying
      both equations, any pair \( (x,y) \) will do), but it never has a unique
      solution.
      Note that \( a=0 \) and \( c=0 \) gives that \( ad-bc=0 \).  
    \end{answer}
  \recommended \item 
    In the system
    \begin{equation*}
      \begin{linsys}{2}
         ax  &+  &by  &=  &c  \\
         dx  &+  &ey  &=  &f  
      \end{linsys}
    \end{equation*}
    each of the equations describes a line in the \( xy \)-plane.
    By geometrical reasoning, show that there are three possibilities:
    there is a unique solution, there is no solution, 
    and there are infinitely many solutions.
    \begin{answer}
      Recall that if a pair of lines share two distinct points then
      they are the same line. 
      That's because two points determine a line, so these
      two points determine each of the two lines, 
      and so they are the same line.

      Thus the lines can share one point (giving a unique solution), 
      share no points (giving no solutions), or
      share at least two points (which makes them the same line).  
    \end{answer}
  \item \label{ex:ProveGaussMethod}
    Finish the proof of \nearbytheorem{th:GaussMethod}.
    \begin{answer}
     For the reduction operation of multiplying $\rho_i$ by a nonzero
     real number $k$, we have that \( (s_1,\ldots,s_n) \) satisfies
     this system
     \begin{equation*}
       \begin{linsys}{4}
         a_{1,1}x_1  &+  &a_{1,2}x_2 &+  &\cdots  &+  &a_{1,n}x_n  &=  &d_1  \\
                   &   &          &  &        &   &           &\vdotswithin{=}   \\
        ka_{i,1}x_1  &+  &ka_{i,2}x_2 &+  &\cdots  &+  &ka_{i,n}x_n &=  &kd_i  \\
                   &   &           &   &        &   &            &\vdotswithin{=}   \\
         a_{m,1}x_1  &+  &a_{m,2}x_2 &+  &\cdots  &+  &a_{m,n}x_n  &=  &d_m  
       \end{linsys}
     \end{equation*}
     if and only if
     % \begin{align*}
     %    a_{1,1}s_1+a_{1,2}s_2+\cdots+a_{1,n}s_n
     %    &=d_1                                              \\
     %    &\alignedvdots                                     \\
     %    \text{and\ } ka_{i,1}s_1+ka_{i,2}s_2+\cdots+ka_{i,n}s_n
     %    &=kd_i                                              \\
     %    &\alignedvdots                                      \\
     %    \text{and\ } a_{m,1}s_1+a_{m,2}s_2+\cdots+a_{m,n}s_n
     %    &=d_m
     % \end{align*}
     $a_{1,1}s_1+a_{1,2}s_2+\cdots+a_{1,n}s_n=d_1$ 
     and \ldots{} $ka_{i,1}s_1+ka_{i,2}s_2+\cdots+ka_{i,n}s_n=kd_i$
     and \ldots{} 
     $a_{m,1}s_1+a_{m,2}s_2+\cdots+a_{m,n}s_n=d_m$
     by the definition of `satisfies'.
     Because \( k\neq 0 \), that's true if and only if
     % \begin{align*}
     %    a_{1,1}s_1+a_{1,2}s_2+\cdots+a_{1,n}s_n
     %    &=d_1                                              \\
     %    &\alignedvdots                                     \\
     %    \text{and\ } a_{i,1}s_1+a_{i,2}s_2+\cdots+a_{i,n}s_n
     %    &=d_i                                              \\
     %    &\alignedvdots                                      \\
     %    \text{and\ } a_{m,1}s_1+a_{m,2}s_2+\cdots+a_{m,n}s_n
     %    &=d_m
     % \end{align*}
     $a_{1,1}s_1+a_{1,2}s_2+\cdots+a_{1,n}s_n=d_1$
     and \ldots{}
     $a_{i,1}s_1+a_{i,2}s_2+\cdots+a_{i,n}s_n=d_i$
     and \ldots{}
     $a_{m,1}s_1+a_{m,2}s_2+\cdots+a_{m,n}s_n=d_m$
     (this is straightforward canceling on both sides of the $i$-th equation),
     which says that \( (s_1,\ldots,s_n) \) solves
     \begin{equation*}
       \begin{linsys}{4}
         a_{1,1}x_1  &+  &a_{1,2}x_2 &+  &\cdots  &+  &a_{1,n}x_n  &=  &d_1  \\
                     &   &           &   &        &   &            &\vdotswithin{=}   \\
         a_{i,1}x_1  &+  &a_{i,2}x_2 &+  &\cdots  &+  &a_{i,n}x_n  &=  &d_i  \\
                     &   &           &   &        &   &            &\vdotswithin{=}   \\
         a_{m,1}x_1  &+  &a_{m,2}x_2 &+  &\cdots  &+  &a_{m,n}x_n  &=
              &d_m  
         \end{linsys}
     \end{equation*}
     as required.

     For the combination operation $k\rho_i+\rho_j$, the tuple
     \( (s_1,\ldots,s_n) \) satisfies
     \begin{equation*}
       \begin{linsys}{4}
         a_{1,1}x_1             &+  &\cdots  &+  &a_{1,n}x_n  &=  &d_1\hfill\hbox{} \\
                                &   &        &   &            &\vdotswithin{=}   \\
         a_{i,1}x_1             &+  &\cdots  &+  &a_{i,n}x_n  &=  &d_i\hfill\hbox{} \\
                                &   &        &   &            &\vdotswithin{=}   \\
         (ka_{i,1}+a_{j,1})x_1  &+  &\cdots  &+  &(ka_{i,n}+a_{j,n})x_n
               &=  &kd_i+d_j \hfill \\
                                &   &        &   &            &\vdotswithin{=}   \\
         a_{m,1}x_1             &+   &\cdots  &+  &a_{m,n}x_n  &=  
          &d_m\hfill\hbox{} 
        \end{linsys}
     \end{equation*}
     if and only if
     $a_{1,1}s_1+\cdots+a_{1,n}s_n=d_1$
     and \ldots{}
     $a_{i,1}s_1+\cdots+a_{i,n}s_n=d_i$
     and \ldots{} 
     $(ka_{i,1}+a_{j,1})s_1+\cdots+(ka_{i,n}+a_{j,n})s_n=kd_i+d_j$
     and \ldots{}
     $a_{m,1}s_1+a_{m,2}s_2+\cdots+a_{m,n}s_n=d_m$
     again by the definition of `satisfies'.
     Subtract \( k \) times the equation~\( i \) from equation~\( j \).
     (Here is where we need \( i\neq j \); if \( i=j \) then the two
     \( d_i \)'s above are not equal.)
     The previous compound statement holds if and only if
     $a_{1,1}s_1+\cdots+a_{1,n}s_n=d_1$
     and \ldots{}
     $a_{i,1}s_1+\cdots+a_{i,n}s_n=d_i$
     and\ldots{} $
     (ka_{i,1}+a_{j,1})s_1+\cdots+(ka_{i,n}+a_{j,n})s_n
         -(ka_{i,1}s_1+\cdots+ka_{i,n}s_n)
         =kd_i+d_j-kd_i$
      and\ldots{} $a_{m,1}s_1+\cdots+a_{m,n}s_n=d_m$,
     which after cancellation says that \( (s_1,\ldots,s_n) \) solves
     \begin{equation*}
       \begin{linsys}{4}
         a_{1,1}x_1  &+   &\cdots  &+  &a_{1,n}x_n  &=  &d_1  \\
                     &    &        &   &            &\vdotswithin{=}   \\
         a_{i,1}x_1  &+   &\cdots  &+  &a_{i,n}x_n  &=  &d_i  \\
                     &    &        &   &            &\vdotswithin{=}   \\
         a_{j,1}x_1  &+  &\cdots  &+  &a_{j,n}x_n  &=  &d_j  \\
                     &   &        &   &            &\vdotswithin{=}   \\
         a_{m,1}x_1  &+  &\cdots  &+  &a_{m,n}x_n  &=
              &d_m\hfill\hbox{}
       \end{linsys}
     \end{equation*}  
     as required.
   \end{answer}
  \item 
    Is there a two-unknowns
    linear system whose solution set is all of \( \Re^2 \)?
    \begin{answer}
      Yes, this one-equation system:
      \begin{equation*}
         0x+0y=0
      \end{equation*}
      is satisfied by every \( (x,y)\in\Re^2 \).  
    \end{answer}
  \recommended \item 
    Are any of the operations used in Gauss's Method
    redundant?
    That is, can we make any of the operations from a combination
    of the others?
    \begin{answer}
      Yes.
      This sequence of operations swaps rows \( i \) and \( j \)
      \begin{equation*}
         \grstep{\rho_i+\rho_j}
         \repeatedgrstep{-\rho_j+\rho_i}
         \repeatedgrstep{\rho_i+\rho_j}
         \repeatedgrstep{-1\rho_i}
      \end{equation*}  
      so the row-swap operation is redundant in the presence of the other two.
     \end{answer}
  \item 
    Prove that each operation of Gauss's Method is reversible.
    That is, show that if two systems are related by a row operation
    $S_1\rightarrow S_2$ then there is a row operation to go back
    $S_2\rightarrow S_1$.
    \begin{answer}
      Reverse a row swap $\rho_i\leftrightarrow\rho_j$ by swapping 
      back $\rho_j\leftrightarrow\rho_i$.
      Reverse the $k\rho_i$ step of multiplying  \( k\neq 0  \) 
      on both sides of a row
      by  dividing through~$(1/k)\rho_i$.

      The row combination case is the nontrivial one.
      The operation $k\rho_i+\rho_j$
      results in this $j$-th row.
      \begin{equation*}
        k\cdot a_{i,1}+a_{j,1}+\cdots+k\cdot a_{i,n}+a_{j,n}=k\cdot d_i+d_j
      \end{equation*}
      The $i$-th row unchanged because of the $i\neq j$ restriction.
      Because the $i$-th row is unchanged, the operation $-k\rho_i+\rho_j$ 
      returns the $j$-th row to its original state.

      (Observe that the \( i=j \) condition on the $k\rho_i+\rho_j$ 
       is needed, or else this could happen
       \begin{equation*}
         \begin{linsys}{2}
           3x  &+  &2y  &=  &7  
         \end{linsys}
         \grstep{2\rho_1+\rho_1}
         \begin{linsys}{2}
           9x  &+  &6y  &=  &21 
          \end{linsys}                 
         \grstep{-2\rho_1+\rho_1}
         \begin{linsys}{2}
          -9x  &-  &6y  &=  &-21 
         \end{linsys}
       \end{equation*}
       and so the result wouldn't hold.)
    \end{answer}
  \puzzle \item  
    \cite{Anton}
    A box holding pennies, nickels and dimes contains
    thirteen coins with a total value of \( 83 \) cents.
    How many coins of each type are in the box?
    (These are US coins;
    a penny is $1$~cent, a nickel is $5$~cents, and
    a dime is $10$~cents.)
    \begin{answer}
      Let \( p \), \( n \), and \( d \) be the number of
      pennies, nickels, and dimes.
      For variables that are real numbers, this system
      \begin{equation*}
         \begin{linsys}{3}
              p  &+ &n   &+  &d   &=  &13   \\
              p  &+ &5n  &+  &10d &=  &83    
         \end{linsys}
         \grstep{-\rho_1+\rho_2}
         \begin{linsys}{3}
              p  &+ &n   &+  &d   &=  &13   \\
                 &  &4n  &+  &9d  &=  &70    
          \end{linsys}
      \end{equation*}
      has more than one solution; in fact, it has infinitely many of them.
      However, it has a limited number of solutions in which \( p \), \( n \),
      and \( d \) are non-negative integers.
      Running through \( d=0 \), \ldots, \( d=8 \) shows that 
      \( (p,n,d)=(3,4,6) \)
      is the only solution using natural numbers.  
    \end{answer}
  \puzzle \item 
    \cite{ContestProb1955no38}
    Four positive integers are given.
    Select any three of the integers, find their arithmetic average,
    and add this result to the fourth integer.
    Thus the numbers 29, 23, 21, and 17 are obtained.
    One of the original integers is:
    \begin{exparts*}
      \partsitem 19
      \partsitem 21
      \partsitem 23
      \partsitem 29
      \partsitem 17
    \end{exparts*}
    \begin{answer}
      Solving the system 
      \begin{equation*}
        \begin{linsys}{2}
        (1/3)(a+b+c)  &+  &d  &=  &29  \\
        (1/3)(b+c+d)  &+  &a  &=  &23  \\
        (1/3)(c+d+a)  &+  &b  &=  &21  \\
        (1/3)(d+a+b)  &+  &c  &=  &17
        \end{linsys}
      \end{equation*}
      we obtain $a=12$, $b=9$, $c=3$, $d=21$.
      Thus the second item, 21, is the correct answer.
     \end{answer}
  \puzzle \item  
      \cite{Monthly35p47}
      Laugh at this:  \( \mbox{AHAHA}+\mbox{TEHE}=\mbox{TEHAW} \).
      It resulted from substituting a code letter for each digit of a simple
      example in addition, and it is required to identify the letters
      and prove the solution unique.
      \begin{answer}
        \answerasgiven
        A comparison of the units and hundreds columns of this
        addition shows that there must be a carry from the tens column.
        The tens column then tells us that \( A<H \), so there
        can be no carry from the units or hundreds columns.
        The five columns then give the following five equations.
        \begin{align*}
          A+E  &=  W  \\
          2H   &=  A+10  \\
          H    &=  W+1  \\
          H+T  &=  E+10  \\
          A+1  &=  T
        \end{align*}
        The five linear equations in five unknowns, if solved simultaneously,
        produce the unique solution: \( A=4 \), \( T=5 \), \( H=7 \),
        \( W=6 \) and \( E=2 \), so that the original example in addition
        was \( 47474+5272=52746 \).  
      \end{answer}
  \puzzle \item 
     \cite{Wohascum2}
     The Wohascum County Board of Commissioners, which has 20 members, 
     recently had to elect a President.
     There were three candidates ($A$, $B$, and $C$); on each ballot
     the three
     candidates were to be listed in order of preference, with no abstentions.
     It was found that 11 members, a majority, preferred $A$ over $B$
     (thus the other 9 preferred $B$ over $A$).
     Similarly, it was found that 12 members preferred $C$ over $A$.
     Given these results, it was suggested that $B$ should withdraw, to enable
     a runoff election between $A$ and $C$.
     However, $B$ protested, and it was then found that 14 members preferred
     $B$ over $C$!
     The Board has not yet recovered from the resulting confusion.
     Given that every possible order of $A$, $B$, $C$ appeared on at least 
     one ballot, how many members voted for $B$ as their first choice?
     \begin{answer}
       \answerasgiven{} 
       \textit{Some additional material was added from \cite{joriki}.}
       Eight commissioners voted for $B$.
       To see this, we will use the given information to study how many voters
       chose each order of $A$, $B$, $C$.

       The six orders of preference are $ABC$, $ACB$, $BAC$, $BCA$, $CAB$,
       $CBA$; assume they receive $a$, $b$, $c$, $d$, $e$, $f$ votes 
       respectively.
       We know that
       \begin{equation*}
         \begin{linsys}{3}
           a  &+  &b  &+  &e  &=  &11  \\
           d  &+  &e  &+  &f  &=  &12  \\
           a  &+  &c  &+  &d  &=  &14
         \end{linsys}
       \end{equation*}
       from the number preferring $A$ over $B$, the number preferring
       $C$ over $A$, and the number preferring $B$ over $C$.
       Because 20 votes were cast, 
       $a+b+\cdots+f=20$.
       Subtracting the sum of the three above equations from twice the 
       prior equation gives $b+c+f=3$.
       We've specified that each preference order got at 
       least one vote, so that means $b=c=f=1$. 

       From the above three equations the complete solution is then 
       $a=6$, $b=1$, $c=1$, $d=7$, $e=4$, and~$f=1$,
       as we can find with Gauss's Method.
       The number of commissioners voting for $B$ as their first choice 
       is therefore $c+d=1+7=8$.

       \par\noindent {\em Comments.}
       The answer to this question would have been the same had we known only
       that {\em at least\/} 14 commissioners preferred $B$ over $C$.

       The seemingly paradoxical nature of the commissioner's preferences
       ($A$ is preferred to $B$, and $B$ is preferred to $C$, and $C$ is 
       preferred to $A$), an example of ``non-transitive dominance'', is
       common when individual choices are pooled.
     \end{answer}
  \puzzle \item   
     \cite{Monthly63p93}
    ``This system
     of \( n \) linear equations with
     \( n \) unknowns,'' said the Great Mathematician, ``has a curious
     property.''

     ``Good heavens!'' said the Poor Nut,  ``What is it?''

     ``Note,'' said the Great Mathematician, ``that the constants are in
     arithmetic progression.''

     ``It's all so clear when you explain it!'' said the Poor Nut.
     ``Do you mean like \( 6x+9y=12 \) and \( 15x+18y=21 \)?''

     ``Quite so,'' said the Great Mathematician, pulling out his bassoon.
     ``Indeed, the system has a unique solution.
     Can you find it?''

     ``Good heavens!'' cried the Poor Nut, ``I am baffled.''

     Are you?
     \begin{answer}
       \answerasgiven
       \textit{We have not used ``dependent'' yet; 
       it means here that Gauss's
       Method shows that there is not a unique solution.}
       If \( n\geq 3 \) the system is dependent and the solution is not
       unique.
       Hence \( n<3 \).
       But the term ``system'' implies \( n>1 \).
       Hence \( n=2 \).
       If the equations are
       \begin{equation*}
         \begin{linsys}{2}
              ax  &+ &(a+d)y  &=  &a+2d  \\
         (a+3d)x  &+ &(a+4d)y &=  &a+5d  
         \end{linsys}
       \end{equation*}
       then \( x=-1 \), \( y=2 \).  
    \end{answer}
\end{exercises}







\subsection{Describing the Solution Set}
A linear system with a unique solution has a solution set with one element.
A linear system with no solution has a solution set that is empty.
In these cases the solution set is easy to describe.
Solution sets are a challenge to describe only when they contain many elements.

\begin{example}
This system has many solutions because in echelon form
\begin{align*}
  \begin{linsys}{3}
    2x  &   &   &+  &z  &=  &3 \\
     x  &-  &y  &-  &z  &=  &1 \\
    3x  &-  &y  &   &   &=  &4 
  \end{linsys}
  &\grstep[-(3/2)\rho_1 +\rho_3]{-(1/2)\rho_1+\rho_2}
  \begin{linsys}{3}
     2x  &   &   &+  &z      &=  &3    \\
         &   &-y &-  &(3/2)z &=  &-1/2 \\
         &   &-y &-  &(3/2)z &=  &-1/2 
   \end{linsys}                                   \\
  &\grstep{-\rho_2+\rho_3}
  \begin{linsys}{3}
     2x  &   &   &+  &z      &=  &3    \\
         &   &-y &-  &(3/2)z &=  &-1/2 \\
         &   &   &   &0      &=  &0    
   \end{linsys}
\end{align*}
not all of the variables are leading variables.
\nearbytheorem{th:GaussMethod} shows that an $(x,y,z)$  
satisfies the first system if and only if it satisfies the
third.
So we can describe the solution set  
$\set{(x,y,z)\suchthat\text{$2x+z=3$ and $x-y-z=1$ and $3x-y=4$}}$
in this way.
\begin{equation*}
  \set{(x,y,z)\suchthat\text{$2x+z=3$ and $-y-3z/2=-1/2$}}
  \tag{$*$}
\end{equation*}
This description is better because 
it has two equations instead of three 
but it is not optimal
because it still has some hard to understand interactions among the variables.

To improve it, use the 
variable that does not lead any equation, $z$, to describe
the variables that do lead, $x$ and $y$.
The second equation gives
$y=(1/2)-(3/2)z$ 
and the first equation gives
$x=(3/2)-(1/2)z$.
Thus we can describe the solution set as this set of triples.   
\begin{equation*}
  \set{((3/2)-(1/2)z,\,(1/2)-(3/2)z,\,z)\suchthat z\in\Re}
  \tag{$**$}
\end{equation*}
\end{example}

Compared with ($*$), 
the advantage of ($**$) 
is that $z$ can be any real number.
This makes the job of deciding which tuples are in the solution set 
much easier.
For instance, taking $z=2$ shows that $(1/2,-5/2,2)$ is a solution.

\begin{definition} \label{df:FreeVars}
%<*df:FreeVars>
In an echelon form linear system the variables that are not leading
are   
\definend{free}.\index{echelon form!free variable}\index{free variable}
%</df:FreeVars>
\end{definition}

\begin{example}   \label{ex:Parametrize2}
Reduction of a linear system can end with more than one variable free.
Gauss's Method on this system
\begin{align*}
   \begin{linsys}{4}
               x  &+  &y   &+  &z   &-  &w   &=  &1  \\
                  &   &y   &-  &z   &+  &w   &=  &-1 \\
              3x  &   &    &+  &6z  &-  &6w  &=  &6  \\
                  &   &-y  &+  &z   &-  &w   &=  &1  
   \end{linsys}
  &\grstep{-3\rho_1 +\rho_3}
  \begin{linsys}{4}
     x  &+  &y   &+  &z   &-  &w   &=  &1  \\
        &   &y   &-  &z   &+  &w   &=  &-1 \\
        &   &-3y &+  &3z  &-  &3w  &=  &3  \\
        &   &-y  &+  &z   &-  &w   &=  &1  
  \end{linsys}                                      \\
  &\grstep[\rho_2 +\rho_4]{3\rho_2 +\rho_3}
  \begin{linsys}{4}
     x  &+  &y   &+  &z   &-  &w   &=  &1  \\
        &   &y   &-  &z   &+  &w   &=  &-1 \\
        &   &    &   &    &   &0   &=  &0  \\
        &   &    &   &    &   &0   &=  &0  
   \end{linsys}
\end{align*}
leaves  \( x \) and \( y \) leading and both \( z \) and~\( w \) free.
To get the description that we prefer, we work from the bottom.
We first express the leading variable $y$ in terms of
$z$ and $w$, as $y=-1+z-w$.
Moving up to the top equation,
substituting for $y$ gives
$x+(-1+z-w)+z-w=1$ and solving for $x$ leaves $x=2-2z+2w$.
The solution set 
\begin{equation*}
   \set{(2-2z+2w,-1+z-w,z,w)\suchthat z,w\in\Re}
  \tag{$**$}
\end{equation*}
has the leading variables expressed in terms of the variables that are free.
\end{example}

\begin{example}    \label{ex:Parametrize1}
The list of leading variables may skip over some columns.
After this reduction
\begin{align*}
  \begin{linsys}{4}
              2x  &-  &2y  &   &    &   &    &=  &0  \\
                  &   &    &   &z   &+  &3w  &=  &2  \\
              3x  &-  &3y  &   &    &   &    &=  &0  \\
               x  &-  &y   &+  &2z  &+  &6w  &=  &4  
  \end{linsys}
  &\grstep[-(1/2)\rho_1+ \rho_4]{-(3/2)\rho_1 +\rho_3}
  \begin{linsys}{4}
     2x  &-  &2y  &\spaceforemptycolumn   &    &   &    &=  &0  \\
         &   &    &   &z   &+  &3w  &=  &2  \\
         &   &    &   &    &   &0   &=  &0  \\
         &   &    &   &2z  &+  &6w  &=  &4  
   \end{linsys}                                    \\
  &\grstep{-2\rho_2 +\rho_4}
  \begin{linsys}{4}
     2x  &-  &2y  &\spaceforemptycolumn   &    &   &    &=  &0  \\
         &   &    &   &z   &+  &3w  &=  &2  \\
         &   &    &   &    &   &0   &=  &0  \\
         &   &    &   &    &   &0   &=  &0  
   \end{linsys}
\end{align*}
$x$ and $z$ are the leading variables, not $x$ and~$y$.
The free variables are $y$ and~$w$ and so we can describe the solution set as
$\set{ (y,y,2-3w,w)\suchthat y,w\in\Re }$.
For instance, \( (1,1,2,0) \) satisfies the system\Dash take 
$y=1$ and $w=0$.
The four-tuple \( (1,0,5,4) \) is not a solution
since its first coordinate does not equal its second.
\end{example}

%<*df:Parameter>
A variable that we use to describe a family of solutions
is a \definend{parameter}.\index{parameter}  
%</df:Parameter>
We say that the solution set in the prior example 
is \definend{parametrized\/}\index{parametrized} 
with $y$ and $w$.

The terms `parameter' and `free variable' do not mean the same thing.
In the prior example
$y$ and~$w$ are free because in the echelon form system they
do not lead. 
They are parameters because 
we used them to describe the set of solutions.
Had we instead 
rewritten the second equation as $w=2/3-(1/3)z$ then
the free variables would still be $y$ and~$w$ but the parameters 
would be $y$ and~$z$.

In the rest of this book 
we will solve linear systems by bringing them to
echelon form and then parametrizing with the free variables.

\begin{example}
This is another system with infinitely many solutions.
\begin{align*}
   \begin{linsys}{4}
               x  &+  &2y  &   &   &   &   &=  &1  \\
              2x  &   &    &+  &z  &   &   &=  &2  \\
              3x  &+  &2y  &+  &z  &-  &w  &=  &4  
   \end{linsys}
  &\grstep[-3\rho_1 +\rho_3]{-2\rho_1+\rho_2}
  \begin{linsys}{4}
     x  &+  &2y  &   &   &   &   &=  &1  \\
        &   &-4y &+  &z  &   &   &=  &0  \\
        &   &-4y &+  &z  &-  &w  &=  &1  
   \end{linsys}                                    \\
  &\grstep{-\rho_2+\rho_3}
  \begin{linsys}{4}
     x  &+  &2y  &   &   &   &   &=  &1  \\
        &   &-4y &+  &z  &   &   &=  &0  \\
        &   &    &   &   &   &-w &=  &1  
   \end{linsys}
\end{align*}
The leading variables are \( x \), \( y \), and \( w \).
The variable \( z \) is free.
Notice that, although there are infinitely many 
solutions, the value of $w$ doesn't vary but is constant $w=-1$.
To parametrize, write \( w \) in terms of \( z \) with \( w=-1+0z \).
Then \( y=(1/4)z \).
Substitute for \( y \) in the first 
equation to get \( x=1-(1/2)z \).
The solution set is $\set{(1-(1/2)z,(1/4)z,z,-1)\suchthat z\in\Re}$.
\end{example}

Parametrizing solution sets shows that systems with 
free variables have infinitely many solutions.
For instance, above $z$ takes on all of infinitely many real number values, 
each associated with a different solution.

We finish this subsection by developing a streamlined 
notation for linear systems 
and their solution sets.

\begin{definition}\label{df:matrix}
%<*df:matrix>
An \( \nbym{m}{n} \) \definend{matrix}\index{matrix}
is a rectangular array of numbers
with \( m \)~\definend{rows}\index{matrix!row}\index{row} 
and \( n \)~\definend{columns}\index{matrix!column}\index{column}.
Each number in the matrix is an 
\definend{entry}\index{matrix!entry}\index{entry, matrix}.
%</df:matrix>
\end{definition}

We usually denote a matrix with an upper case roman letter.
For instance,
\begin{equation*}
  A=
  \begin{mat}[r]
    1  &2.2  &5  \\
    3  &4    &-7
  \end{mat}
\end{equation*}
has $2$~rows and $3$~columns and so
is a \( \nbym{2}{3} \) matrix.
Read that aloud as ``two-by-three'';
the number of rows is always stated first.
(The matrix has parentheses around it
so that when 
two matrices are adjacent
we can tell where one ends and the other begins.)
We name matrix entries with the corresponding lower-case letter,
so that the entry in the second row and first column 
of the above array is \( a_{2,1}=3 \).
Note that the order of the subscripts matters: 
$a_{1,2}\neq a_{2,1}$ since \( a_{1,2}=2.2 \). 
We denote
the set of all \( \nbym{m}{n} \)
matrices by \( \matspace_{\nbym{m}{n}} \).\index{matrices|set of}

We do Gauss's Method using matrices in essentially the same
way that we did it for systems of equations:
a matrix row's
\definend{leading entry}\index{echelon form!leading entry}%
\index{leading!entry} % 
is its first nonzero entry (if it has one) and
we perform row operations to arrive at 
\definend{matrix echelon form},\index{echelon form!matrix}%
\index{matrix!echelon form} %
where the leading entry in lower rows are to the right of those in 
the rows above.
We like matrix notation because it lightens
the clerical load, the copying of variables and the
writing of $+$'s and $=$'s. 

\begin{example}
We can abbreviate this linear system
\begin{equation*}
  \begin{linsys}{3}
    x  &+  &2y  &   &    &=  &4   \\
       &   &y   &-  &z   &=  &0   \\
    x  &   &    &+  &2z  &=  &4   
  \end{linsys}
\end{equation*}
with this matrix.
\begin{equation*}
    \begin{amat}[r]{3}
      1  &2  &0  &4  \\
      0  &1  &-1 &0  \\
      1  &0  &2  &4
    \end{amat}
\end{equation*}
The vertical bar reminds a reader of the difference between the 
coefficients on the system's left hand side and the constants on the right.
With a bar, this is an
\definend{augmented\/}\index{matrix!augmented}\index{augmented matrix} matrix.
\begin{equation*}
    \begin{amat}[r]{3}
      1  &2  &0  &4  \\
      0  &1  &-1 &0  \\
      1  &0  &2  &4
    \end{amat}
  \grstep{-\rho_1 +\rho_3}
  \begin{amat}[r]{3}
       1  &2  &0  &4  \\
       0  &1  &-1 &0  \\
       0  &-2 &2  &0
     \end{amat}                        
  \grstep{2\rho_2 +\rho_3}
  \begin{amat}[r]{3}
       1  &2  &0  &4  \\
       0  &1  &-1 &0  \\
       0  &0  &0  &0
     \end{amat}
\end{equation*}
The second row stands for $y-z=0$ and the first row stands for
$x+2y=4$ so the solution set is
\( \set{(4-2z,z,z)\suchthat z\in\Re} \).
\end{example}

Matrix notation also
clarifies the descriptions of solution sets.
\nearbyexample{ex:Parametrize2}'s
$\set{(2-2z+2w,-1+z-w,z,w)\suchthat z,w\in\Re}$ 
is hard to read.
We will rewrite it to group all of the
constants together, all of the
coefficients of~\( z \) together, and all of the coefficients of~\( w \)
together.
We write them vertically, in one-column matrices.
\begin{equation*}
  \set{\colvec[r]{2 \\ -1 \\ 0 \\ 0}
       +\colvec[r]{-2 \\ 1 \\ 1 \\ 0}\cdot z
       +\colvec[r]{2 \\ -1 \\ 0 \\ 1}\cdot w
       \suchthat z,w\in\Re}
\end{equation*}
For instance, the top line says that \( x=2-2z+2w \)
and the second line says that \( y= -1+z-w \).
(Our next section gives a geometric interpretation that will help 
us picture the solution sets.)

\begin{definition}\label{df:vector}
%<*df:vector>
A 
\definend{column vector},\index{column!vector}\index{vector!column}
often just called a \definend{vector},\index{vector} 
is a matrix with a single column.
A matrix with a single row is a
\definend{row vector}\index{row!vector}\index{vector!row}.
The entries of a vector are sometimes called
\definend{components}\index{component of a vector}\index{vector!component}.
A column or row vector whose components are all zeros is a 
\definend{zero vector}.\index{zero vector}\index{vector!zero}
%</df:vector>
\end{definition}

Vectors are an exception to the convention of representing matrices with 
capital roman letters.
We use lower-case roman or greek letters overlined
with an arrow:
\( \vec{a} \), \( \vec{b} \), \ldots\, or
\( \vec{\alpha} \), \( \vec{\beta} \), \ldots\
(boldface is also common:
{\boldmath \( a \)} or {\boldmath \( \alpha \)}).
For instance, this is a column vector
with a third component of \( 7 \).
\begin{equation*}
  \vec{v}=
  \colvec[r]{ 1  \\  3  \\ 7}
\end{equation*}
A zero vector is denoted \( \zero \).
There are many different zero vectors\Dash the
one-tall zero vector, the two-tall zero vector, etc.\Dash but
nonetheless we will often say ``the'' zero vector, expecting 
that the size will be clear from the context.

\begin{definition}
The linear equation
\( a_1x_1+a_2x_2+\,\cdots\,+a_nx_n=d \)
with unknowns \( x_1,\ldots\,,x_n \)
is \definend{satisfied}\index{vector!satisfies an equation}%
\index{linear equation!satisfied by a vector} by
\begin{equation*}
  \vec{s}=\colvec{s_1 \\ \vdotswithin{s_1} \\ s_n}
\end{equation*}
if \( a_1s_1+a_2s_2+\,\cdots\,+a_ns_n=d \).
A vector satisfies a linear system if it satisfies each equation in 
the system.
\end{definition}

The style of description of solution sets that we use
involves adding the vectors, and 
also multiplying them by real numbers.
Before we give the examples showing the style 
we first need to define these operations.

\begin{definition} \label{df:VectorSum}
%<*df:VectorSum>
The 
\definend{vector sum}\index{vector!sum}\index{sum!vector}\index{addition of vectors} 
of
\( \vec{u} \) and \( \vec{v} \) is the vector of the sums.
\begin{equation*}
  \vec{u}+\vec{v}=
  \colvec{u_1 \\ \vdotswithin{u_1} \\ u_n}
   +
  \colvec{v_1 \\ \vdotswithin{v_1} \\ v_n}
   =
  \colvec{u_1+v_1 \\ \vdotswithin{u_1+v_1} \\ u_n+v_n}
\end{equation*}
%</df:VectorSum>
\end{definition}

Note that for the addition to be defined 
the vectors must have the same number of entries.
This entry-by-entry addition works for any pair of matrices, not just vectors, 
provided that they have the same number of rows and 
columns.\index{matrix!sum}\index{sum!matrix}

\begin{definition} \label{df:VectorScalarMultiplication}
%<*df:VectorScalarMultiplication>
The \definend{scalar multiplication\/}\index{vector!scalar multiple}%
\index{scalar multiple!vector} of the real number
\( r \) and the vector \( \vec{v} \) is the vector of the multiples.
\begin{equation*}
  r\cdot\vec{v}=
  r\cdot\colvec{v_1 \\ \vdotswithin{v_1} \\ v_n}
  =
  \colvec{rv_1 \\ \vdotswithin{rv_1} \\ rv_n}
\end{equation*}
%</df:VectorScalarMultiplication>
\end{definition}

As with the addition operation, the entry-by-entry scalar multiplication
operation extends beyond vectors to apply to any 
matrix.\index{matrix!scalar multiplication}\index{scalar multiplication!matrix}

We write scalar multiplication either as \( r\cdot\vec{v} \) or
\( \vec{v}\cdot r \), and sometimes even omit the `$\cdot$' symbol:~$r\vec{v}$.
(Do not refer to scalar multiplication 
as `scalar product' because that name is for a different operation.)

\begin{example}
\begin{equation*}
  \colvec[r]{2 \\ 3 \\ 1}
   +
  \colvec[r]{3 \\ -1 \\ 4}
  =
  \colvec{2+3 \\ 3-1 \\ 1+4}
   =
  \colvec[r]{5 \\ 2 \\ 5}
  \qquad
  7\cdot\colvec[r]{1 \\ 4 \\ -1 \\ -3}
  =
  \colvec[r]{7 \\ 28 \\ -7 \\ -21}
\end{equation*}
\end{example}

Observe that the definitions of addition and scalar multiplication agree
where they overlap; for instance, \( \vec{v} +\vec{v} = 2\vec{v} \).

With these definitions, we are set to use matrix and vector notation to
both solve systems and express the solution.

\begin{example} \label{ex:ManyParamsInfManySolsSystem}
This system
\begin{equation*}
   \begin{linsys}{5}
      2x  &+  &y  &  &  &-  &w  &   &   &=  &4  \\
          &   &y  &  &  &+  &w  &+  &u  &=  &4  \\
       x  &   &   &- &z &+  &2w &   &   &=  &0  
   \end{linsys}
\end{equation*}
reduces in this way.
\begin{align*}
  \begin{amat}[r]{5}
    2  &1  &0  &-1  &0  &4  \\
    0  &1  &0  &1   &1  &4  \\
    1  &0  &-1 &2   &0  &0
  \end{amat}
  &\grstep{-(1/2)\rho_1+\rho_3}
  \begin{amat}[r]{5}
    2  &1     &0  &-1    &0  &4  \\
    0  &1     &0  &1     &1  &4  \\
    0  &-1/2  &-1 &5/2   &0  &-2
  \end{amat}                                 \\
  &\grstep{(1/2)\rho_2+\rho_3}
  \begin{amat}[r]{5}
    2  &1     &0  &-1    &0    &4  \\
    0  &1     &0  &1     &1    &4  \\
    0  &0     &-1 &3     &1/2  &0
  \end{amat}
\end{align*}
The solution set is
\( \set{(w+(1/2)u,4-w-u,3w+(1/2)u,w,u)\suchthat w,u\in\Re} \).
We write that in vector form.
\begin{equation*}
  \set{\colvec{x \\ y \\ z \\ w \\ u}=
       \colvec[r]{0 \\ 4 \\ 0 \\ 0 \\ 0}+
       \colvec[r]{1 \\ -1 \\ 3 \\ 1 \\ 0}w+
       \colvec[r]{1/2 \\ -1 \\ 1/2 \\ 0 \\ 1}u
       \suchthat w,u\in\Re}
\end{equation*}
Note how well vector notation sets off 
the coefficients of each parameter.
For instance, the third row of the vector form shows plainly that if \( u \) is
fixed then \( z \) increases three times as fast as \( w \).
Another thing shown plainly is that setting both \( w \) and \( u \) to zero
gives that
\begin{equation*}
  \colvec{x \\ y \\ z \\ w \\ u}
  =\colvec[r]{0 \\ 4 \\ 0 \\ 0 \\ 0}
\end{equation*}
is a particular solution of the linear system.
\end{example}

\begin{example}
In the same way, the system
\begin{equation*}
   \begin{linsys}{3}
     x  &-  &y  &+  &z  &=  &1  \\
    3x  &   &   &+  &z  &=  &3  \\
    5x  &-  &2y &+  &3z &=  &5  
  \end{linsys}
\end{equation*}
reduces
\begin{align*}
  \begin{amat}[r]{3}
    1  &-1  &1  &1  \\
    3  &0   &1  &3  \\
    5  &-2  &3  &5
  \end{amat}
  &\grstep[-5\rho_1+\rho_3]{-3\rho_1+\rho_2}
  \begin{amat}[r]{3}
    1  &-1  &1  &1  \\
    0  &3   &-2 &0  \\
    0  &3   &-2 &0
  \end{amat}                                    \\
  &\grstep{-\rho_2+\rho_3}
  \begin{amat}[r]{3}
    1  &-1  &1  &1  \\
    0  &3   &-2 &0  \\
    0  &0   &0  &0
  \end{amat}
\end{align*}
to give a one-parameter solution set.
\begin{equation*}
  \set{\colvec[r]{1 \\ 0 \\ 0}
       +\colvec[r]{-1/3 \\ 2/3 \\ 1}z
       \suchthat z\in\Re}
\end{equation*}
As in the prior example, the vector not associated with the parameter
\begin{equation*}
   \colvec[r]{1 \\ 0 \\ 0}
\end{equation*}
is a particular solution of the system.
\end{example}

Before the exercises, we will consider what we have accomplished  
and what we will do in the remainder of the chapter.
So far we have done the mechanics of Gauss's Method.
We have not stopped to consider any of the questions
that arise,
except for proving \nearbytheorem{th:GaussMethod}\Dash which 
justifies the method by showing that it gives 
the right answers. 

For example, can we 
always describe solution sets as above, with
a particular solution vector added to an unrestricted linear combination of 
some other vectors?
We've noted that the solution sets described in this way 
have infinitely many members
so answering this question
would tell us about the size of solution sets.
The following subsection shows that the answer is ``yes.''
This chapter's second section then
uses that answer to describe the geometry of solution sets.

Other questions arise from the observation that we can do Gauss's Method in 
more than one way (for instance, when swapping rows we may have a choice of 
rows to swap with).
\nearbytheorem{th:GaussMethod} says that we must get the same solution set
no matter how we proceed but
if we do Gauss's Method in two ways
must we get the same number of free variables in each echelon form system?
Must those be the same variables, that is, is it impossible to
solve a problem
one way to get $y$ and~$w$ free and solve it another way to get $y$ and~$z$ 
free?
The third section of this chapter answers ``yes,'' that  
from any starting linear system,
all derived echelon form versions 
have the same free variables. 

Thus, by the end of the chapter we will not only have a 
solid grounding in the practice of Gauss's Method but 
we will also have a solid grounding in the theory.
We will know exactly what can and cannot happen in a reduction.

\begin{exercises}
  \recommended \item  
    Find the indicated entry of the matrix,
    if it is defined.
    \begin{equation*}
      A=\begin{mat}[r]
        1  &3  &1  \\
        2  &-1 &4
      \end{mat}
    \end{equation*}
    \begin{exparts*}
      \partsitem \( a_{2,1} \)
      \partsitem \( a_{1,2} \)
      \partsitem \( a_{2,2} \)
      \partsitem \( a_{3,1} \)
    \end{exparts*}
    \begin{answer}
      \begin{exparts*}
        \partsitem \( 2 \)
        \partsitem \( 3 \)
        \partsitem \(-1 \)
        \partsitem Not defined.
      \end{exparts*}  
    \end{answer}
  \recommended \item 
    Give the size of each matrix.
    \begin{exparts*}
      \partsitem \(
        \begin{mat}[r]
          1  &0  &4  \\
          2  &1  &5
        \end{mat}  \)
      \partsitem \(
        \begin{mat}[r]
          1  &1  \\
         -1  &1  \\
          3  &-1
        \end{mat}  \)
      \partsitem \(
        \begin{mat}[r]
          5  &10 \\
         10  &5
        \end{mat}  \)
    \end{exparts*}
    \begin{answer}
      \begin{exparts*}
        \partsitem \( \nbym{2}{3} \)
        \partsitem \( \nbym{3}{2} \)
        \partsitem \( \nbym{2}{2} \)
      \end{exparts*}  
    \end{answer}
  \recommended \item 
    Do the indicated vector operation, if it is defined.
    \begin{exparts*}
      \partsitem \( \colvec[r]{2 \\ 1 \\ 1}
               +\colvec[r]{3 \\ 0 \\ 4} \)
      \partsitem \( 5\colvec[r]{4 \\ -1} \)
      \partsitem \( \colvec[r]{1 \\ 5 \\ 1}
               -\colvec[r]{3 \\ 1 \\ 1} \)
      \partsitem \( 7\colvec[r]{2 \\ 1}
               +9\colvec[r]{3 \\ 5} \)
      \partsitem \( \colvec[r]{1 \\ 2}
               +\colvec[r]{1 \\ 2 \\ 3} \)
      \partsitem \( 6\colvec[r]{3 \\ 1 \\ 1}
               -4\colvec[r]{2 \\ 0 \\ 3}
               +2\colvec[r]{1 \\ 1 \\ 5} \)
    \end{exparts*}
    \begin{answer}
      \begin{exparts*}
        \partsitem \( \colvec[r]{5 \\ 1 \\ 5} \)
        \partsitem \( \colvec[r]{20 \\ -5} \)
        \partsitem \( \colvec[r]{-2 \\ 4 \\ 0} \)
        \partsitem \( \colvec[r]{41 \\ 52} \)
        \partsitem Not defined.
        \partsitem \( \colvec[r]{12 \\ 8 \\ 4} \)
      \end{exparts*}  
     \end{answer}
  \recommended \item  \label{exer:SolveInMatrixNotation}
    Solve each system using matrix notation.
    Express the solution using vectors.
    \begin{exparts*}
      \partsitem \( \begin{linsys}[t]{2}
                  3x  &+  &6y  &=  &18  \\
                   x  &+  &2y  &=  &6   
                   \end{linsys}  \)
      \partsitem \( \begin{linsys}[t]{2}
                   x  &+  &y   &=  &1  \\
                   x  &-  &y   &=  &-1   
                    \end{linsys}  \)
      \partsitem \( \begin{linsys}[t]{3}
                   x_1  &   &     &+  &x_3   &=  &4  \\
                   x_1  &-  &x_2  &+  &2x_3  &=  &5  \\
                  4x_1  &-  &x_2  &+  &5x_3  &=  &17  
                   \end{linsys}  \)
      \partsitem \( \begin{linsys}[t]{3}
                   2a   &+  &b    &-  &c     &=  &2  \\
                   2a   &   &     &+  &c     &=  &3  \\
                    a   &-  &b    &   &      &=  &0   
                    \end{linsys}  \)
      \partsitem \( \begin{linsys}[t]{4}
                     x  &+  &2y   &-   &z   &    &    &=  &3  \\
                    2x  &+  &y    &    &    &+   &w   &=  &4  \\
                     x  &-  &y    &+   &z   &+   &w   &=  &1  
                    \end{linsys}  \)
      \partsitem \( \begin{linsys}[t]{4}
                     x  &   &     &+   &z   &+   &w   &=  &4  \\
                    2x  &+  &y    &    &    &-   &w   &=  &2  \\
                    3x  &+  &y    &+   &z   &    &    &=  &7  
                     \end{linsys}  \)
    \end{exparts*}
    \begin{answer}
      \begin{exparts}
        \partsitem This reduction
          \begin{equation*}
            \begin{amat}[r]{2}
              3  &6  &18 \\
              1  &2  &6
            \end{amat}
            \grstep{(-1/3)\rho_1+\rho_2}
            \begin{amat}[r]{2}
              3  &6  &18 \\
              0  &0  &0
            \end{amat}
          \end{equation*}
          leaves \( x \) leading and \( y \) free.
          Making \( y \) the parameter, gives \( x=6-2y \) and this solution
          set.
          \begin{equation*}
            \set{\colvec[r]{6 \\ 0}+\colvec[r]{-2 \\ 1}y
              \suchthat y\in\Re}
          \end{equation*}
        \partsitem A reduction
          \begin{equation*}
            \begin{amat}[r]{2}
              1  &1  &1  \\
              1  &-1 &-1
            \end{amat}
            \grstep{-\rho_1+\rho_2}
            \begin{amat}[r]{2}
              1  &1  &1  \\
              0  &-2 &-2
            \end{amat}
          \end{equation*}
          gives the unique solution \( y=1 \), \( x=0 \).
          The solution set is this.
          \begin{equation*}
            \set{\colvec[r]{0 \\ 1} }
          \end{equation*}
        \partsitem Gauss's Method
          \begin{equation*}
            \begin{amat}[r]{3}
              1  &0  &1  &4  \\
              1  &-1 &2  &5  \\
              4  &-1 &5  &17
            \end{amat}
            \grstep[-4\rho_1+\rho_3]{-\rho_1+\rho_2}
            \begin{amat}[r]{3}
              1  &0  &1  &4  \\
              0  &-1 &1  &1  \\
              0  &-1 &1  &1
            \end{amat}     
            \grstep{-\rho_2+\rho_3}
            \begin{amat}[r]{3}
              1  &0  &1  &4  \\
              0  &-1 &1  &1  \\
              0  &0  &0  &0
            \end{amat}
          \end{equation*}
          leaves \( x_1 \) and \( x_2 \) leading with \( x_3 \) free.
          The solution set is this.
          \begin{equation*}
            \set{\colvec[r]{4 \\ -1 \\ 0}+\colvec[r]{-1 \\ 1 \\ 1}x_3
              \suchthat x_3\in\Re}
          \end{equation*}
        \partsitem This reduction
          \begin{align*}
            \begin{amat}[r]{3}
              2  &1  &-1 &2  \\
              2  &0  &1  &3  \\
              1  &-1 &0  &0
            \end{amat}
            &\grstep[-(1/2)\rho_1+\rho_3]{-\rho_1+\rho_2}
            \begin{amat}[r]{3}
              2  &1    &-1   &2  \\
              0  &-1   &2    &1  \\
              0  &-3/2 &1/2  &-1
            \end{amat}                                          \\
            &\grstep{(-3/2)\rho_2+\rho_3}
            \begin{amat}[r]{3}
              2  &1  &-1   &2  \\
              0  &-1 &2    &1  \\
              0  &0  &-5/2 &-5/2
            \end{amat}
          \end{align*}
          shows that the solution set is a singleton set.
          \begin{equation*}
            \set{\colvec[r]{1 \\ 1 \\ 1}}
          \end{equation*}
        \partsitem This reduction is easy
          \begin{align*}
            \begin{amat}[r]{4}
              1  &2  &-1 &0  &3 \\
              2  &1  &0  &1  &4 \\
              1  &-1 &1  &1  &1
            \end{amat}
            &\grstep[-\rho_1+\rho_3]{-2\rho_1+\rho_2}
            \begin{amat}[r]{4}
              1  &2  &-1 &0  &3  \\
              0  &-3 &2  &1  &-2 \\
              0  &-3 &2  &1  &-2
            \end{amat}                                       \\
            &\grstep{-\rho_2+\rho_3}
            \begin{amat}[r]{4}
              1  &2  &-1 &0  &3  \\
              0  &-3 &2  &1  &-2 \\
              0  &0  &0  &0  &0
            \end{amat}
          \end{align*}
          and ends with \( x \) and $y$ leading while \( z \) and \( w \) are
          free.
          Solving for \( y \) gives \( y=(2+2z+w)/3 \) and substitution shows
          that \( x+2(2+2z+w)/3-z=3 \) so \( x=(5/3)-(1/3)z-(2/3)w \),
          making this the solution set.
          \begin{equation*}
            \set{\colvec[r]{5/3 \\ 2/3 \\ 0 \\ 0}
                 +\colvec[r]{-1/3 \\ 2/3 \\ 1 \\ 0}z
                 +\colvec[r]{-2/3 \\ 1/3 \\ 0 \\ 1}w
                 \suchthat z,w\in\Re}
          \end{equation*}
        \partsitem The reduction
          \begin{align*}
            \begin{amat}[r]{4}
              1  &0  &1  &1  &4 \\
              2  &1  &0  &-1 &2 \\
              3  &1  &1  &0  &7
            \end{amat}
            &\grstep[-3\rho_1+\rho_3]{-2\rho_1+\rho_2}
            \begin{amat}[r]{4}
              1  &0  &1  &1  &4 \\
              0  &1  &-2 &-3 &-6\\
              0  &1  &-2 &-3 &-5
            \end{amat}                                       \\
            &\grstep{-\rho_2+\rho_3}
            \begin{amat}[r]{4}
              1  &0  &1  &1  &4 \\
              0  &1  &-2 &-3 &-6\\
              0  &0  &0  &0  &1
            \end{amat}
          \end{align*}
          shows that there is no solution\Dash the solution set is empty.
      \end{exparts}  
     \end{answer}
  \item \label{exer:SlvMatNot}
    Solve each system using matrix notation.
    Give each solution set in vector notation.
    \begin{exparts*}
      \partsitem \( \begin{linsys}[t]{3}
                  2x  &+  &y  &-  &z  &=  &1  \\
                  4x  &-  &y  &   &   &=  &3  
                \end{linsys}  \)
      \partsitem \( \begin{linsys}[t]{4}
                   x  &   &   &-  &z  &   &   &=  &1  \\
                      &   &y  &+  &2z &-  &w  &=  &3  \\
                   x  &+  &2y &+  &3z &-  &w  &=  &7  
               \end{linsys}  \)
      \partsitem \( \begin{linsys}[t]{4}
                   x  &-  &y  &+  &z  &   &   &=  &0  \\
                      &   &y  &   &   &+  &w  &=  &0  \\
                  3x  &-  &2y &+  &3z &+  &w  &=  &0  \\
                      &   &-y &   &   &-  &w  &=  &0  
               \end{linsys}  \)
      \partsitem \( \begin{linsys}[t]{5}
                   a  &+  &2b &+  &3c &+  &d  &-  &e  &=  &1  \\
                  3a  &-  &b  &+  &c  &+  &d  &+  &e  &=  &3  
               \end{linsys}  \)
    \end{exparts*}
    \begin{answer}
      \begin{exparts}
      \partsitem The reduction
        \begin{equation*}
          \begin{amat}[r]{3}
            2  &1  &-1  &1  \\
            4  &-1 &0   &3
          \end{amat}
          \grstep{-2\rho_1+\rho_2}
          \begin{amat}[r]{3}
            2  &1  &-1  &1  \\
            0  &-3 &2   &1
          \end{amat}
        \end{equation*}
        ends with \( x \) and \( y \) leading, and with \( z \) free.
        Solving for \( y \) gives \( y=(1-2z)/(-3) \), and then substitution
        \( 2x+(1-2z)/(-3)-z=1 \) shows that \( x=((4/3)+(1/3)z)/2 \).
        Hence the solution set is this.
        \begin{equation*}
          \set{\colvec[r]{2/3 \\ -1/3 \\ 0}
               +\colvec[r]{1/6 \\ 2/3 \\ 1}z
              \suchthat z\in\Re}
        \end{equation*}
      \partsitem This application of Gauss's Method
        \begin{align*}
          \begin{amat}[r]{4}
            1  &0  &-1  &0  &1 \\
            0  &1  &2   &-1 &3 \\
            1  &2  &3   &-1 &7
          \end{amat}
          &\grstep{-\rho_1+\rho_3}
          \begin{amat}[r]{4}
            1  &0  &-1  &0  &1 \\
            0  &1  &2   &-1 &3 \\
            0  &2  &4   &-1 &6
          \end{amat}                                   \\
          &\grstep{-2\rho_2+\rho_3}
          \begin{amat}[r]{4}
            1  &0  &-1  &0  &1 \\
            0  &1  &2   &-1 &3 \\
            0  &0  &0   &1  &0
          \end{amat}
        \end{align*}
        leaves  \( x \), \( y \), and \( w \)  leading.
        The solution set is here.
        \begin{equation*}
          \set{\colvec[r]{1 \\ 3 \\ 0 \\ 0}
               +\colvec[r]{1 \\ -2 \\ 1 \\ 0}z
              \suchthat z\in\Re}
        \end{equation*}
      \partsitem This row reduction
        \begin{align*}
          \begin{amat}[r]{4}
            1  &-1 &1   &0  &0 \\
            0  &1  &0   &1  &0 \\
            3  &-2 &3   &1  &0 \\
            0  &-1 &0   &-1 &0
          \end{amat}
          &\grstep{-3\rho_1+\rho_3}
          \begin{amat}[r]{4}
            1  &-1 &1   &0  &0 \\
            0  &1  &0   &1  &0 \\
            0  &1  &0   &1  &0 \\
            0  &-1 &0   &-1 &0
          \end{amat}                                       \\
          &\grstep[\rho_2+\rho_4]{-\rho_2+\rho_3}
          \begin{amat}[r]{4}
            1  &-1 &1   &0  &0 \\
            0  &1  &0   &1  &0 \\
            0  &0  &0   &0  &0 \\
            0  &0  &0   &0  &0
          \end{amat}
        \end{align*}
        ends with \( z \) and \( w \) free.
        We have this solution set.
        \begin{equation*}
          \set{\colvec[r]{0 \\ 0 \\ 0 \\ 0}
               +\colvec[r]{-1 \\ 0 \\ 1 \\ 0}z
               +\colvec[r]{-1 \\ -1 \\ 0 \\ 1}w
              \suchthat z,w\in\Re}
        \end{equation*}
      \partsitem Gauss's Method done in this way
        \begin{equation*}
          \begin{amat}[r]{5}
            1  &2  &3   &1  &-1 &1  \\
            3  &-1 &1   &1  &1  &3
          \end{amat}
          \grstep{-3\rho_1+\rho_2}
          \begin{amat}[r]{5}
            1  &2  &3   &1  &-1 &1  \\
            0  &-7 &-8  &-2 &4  &0
          \end{amat}
        \end{equation*}
        ends with \( c \), \( d \), and \( e \) free.
        Solving for \( b \) shows that \( b=(8c+2d-4e)/(-7) \) and then 
        substitution
        \( a+2(8c+2d-4e)/(-7)+3c+1d-1e=1 \) shows that 
        \( a=1-(5/7)c-(3/7)d-(1/7)e \) and we have the solution set.
        \begin{equation*}
          \set{\colvec[r]{1 \\ 0 \\ 0 \\ 0 \\ 0}
               +\colvec[r]{-5/7 \\ -8/7 \\ 1 \\ 0 \\ 0}c
               +\colvec[r]{-3/7 \\ -2/7 \\ 0 \\ 1 \\ 0}d
               +\colvec[r]{-1/7 \\ 4/7 \\ 0 \\ 0 \\ 1}e
              \suchthat c,d,e\in\Re}
        \end{equation*}
    \end{exparts}  
   \end{answer}
  \item 
    Solve each system using matrix notation.
    Express the solution set using vectors.
    \begin{exparts*}
      \partsitem 
        $\begin{linsys}{3}
          3x &+ &2y &+ &z &= &1 \\
          x  &- &y  &+ &z &= &2 \\
          5x &+ &5y &+ &z &= &0
        \end{linsys}$
      \partsitem
        $\begin{linsys}{4}
          x  &+ &y  &- &2z &= &0 \\
          x  &- &y  &  &   &= &-3 \\
          3x &- &y  &- &2z &= &-6  \\
             &  &2y &- &2z &= &3  
        \end{linsys}$
      \partsitem
        $\begin{linsys}{5}
          2x  &- &y  &- &z &+ &w &= &4 \\
           x  &+ &y  &+ &z &  &  &= &-1 
        \end{linsys}$
      \partsitem
         $\begin{linsys}{3}
          x  &+ &y  &- &2z &= &0 \\
          x  &- &y  &  &   &= &-3 \\
          3x &- &y  &- &2z &= &0    
        \end{linsys}$ 
    \end{exparts*}
    \begin{answer}
      \begin{exparts}
        \partsitem This reduction
          \begin{align*}
            \begin{amat}{3}
              3 &2  &1  &1 \\
              1 &-1 &1  &2 \\
              5 &5  &1  &0
            \end{amat}
            &\grstep[-(5/3)\rho_1+\rho_3]{-(1/3)\rho_1+\rho_2}
            \begin{amat}{3}
              3 &2    &1     &1 \\
              0 &-5/3 &2/3   &5/3 \\
              0 &5/3  &-2/3  &-5/3
            \end{amat}                                     \\
            &\grstep{\rho_2+\rho_3}
            \begin{amat}{3}
              3 &2    &1     &1 \\
              0 &-5/3 &2/3   &5/3 \\
              0 &0    &0     &0
            \end{amat}
          \end{align*}
          gives this solution set.
          \begin{equation*}
            \set{\colvec{x \\ y \\ z}=\colvec{1 \\ -1 \\ 0}
                                       +\colvec{-3/5 \\ 2/5 \\ 1}z
                             \suchthat z\in\Re}
          \end{equation*}
        \partsitem
          This is the reduction.
          \begin{align*}
            \begin{amat}{3}
              1 &1   &-2 &0 \\
              1 &-1  &0  &3  \\
              3 &-1  &-2 &-6 \\
              0 &2   &-2 &3
            \end{amat}
            &\grstep[-3\rho_1+\rho_3]{-\rho_1+\rho_2}
            \begin{amat}{3}
              1 &1   &-2 &0 \\
              0 &-2  &2  &-3  \\
              0 &-4  &4  &-6 \\
              0 &2   &-2 &3
            \end{amat}                                  \\
            &\grstep[\rho_2+\rho_4]{-2\rho_2+\rho_3}
            \begin{amat}{3}
              1 &1   &-2 &0 \\
              0 &-2  &2  &-3  \\
              0 &0   &0  &0 \\
              0 &0   &0  &0
            \end{amat}
          \end{align*}
          The solution set is this.
          \begin{equation*}
            \set{\colvec{-3/2 \\ 3/2 \\ 0}
                 +\colvec{1 \\ 1 \\ 1}z
                 \suchthat z\in\Re}
          \end{equation*}
      \partsitem
         Gauss's Method
         \begin{equation*}
           \begin{amat}{4}
             2 &-1 &-1 &1 &4 \\
             1 &1  &1  &0 &-1
           \end{amat}
           \grstep{-(1/2)\rho_1+\rho_2}
           \begin{amat}{4}
             2 &-1   &-1   &1    &4 \\
             0 &3/2  &3/2  &-1/2 &-3
           \end{amat}
         \end{equation*}
         gives the solution set.
         \begin{equation*}
           \set{\colvec{1 \\ -2 \\ 0 \\ 0}
                  +\colvec{0 \\ -1 \\ 1 \\ 0}z
                   \colvec{-1/3 \\ 1/3 \\ 0 \\ 1}w
                 \suchthat z,w\in\Re}
         \end{equation*}
       \partsitem
          Here is the reduction.
          \begin{equation*}
            \begin{amat}{3}
              1 &1  &-2 &0  \\
              1 &-1 &0  &-3 \\
              3 &-1 &-2 &0
           \end{amat}
           \grstep[-3\rho_1+\rho_3]{-\rho_1+\rho_2}
           \begin{amat}{3}
             1 &1  &-2 &0  \\
             0 &-2 &2  &-3 \\
             0 &-4  &4  &0
          \end{amat}
          \grstep{-2\rho_2+\rho_3}
          \begin{amat}{3}
            1 &1  &-2 &0  \\
            0 &-2 &2  &-3 \\
            0 &0  &0  &6
          \end{amat}
        \end{equation*}
        The solution set is empty~$\set{}$.
      \end{exparts}
    \end{answer}
  \recommended \item 
    The vector is in the set.
    What value of the parameters produces that vector?
    \begin{exparts}
      \partsitem $\colvec[r]{5 \\ -5}$,
        $\set{\colvec[r]{1 \\ -1}k\suchthat k\in\Re}$
      \partsitem $\colvec[r]{-1 \\ 2 \\ 1}$,
        $\set{\colvec[r]{-2 \\ 1 \\ 0}i
           +\colvec[r]{3 \\ 0 \\ 1}j\suchthat i,j\in\Re}$
      \partsitem $\colvec[r]{0 \\ -4 \\ 2}$,
        $\set{\colvec[r]{1 \\ 1 \\ 0}m
               +\colvec[r]{2 \\ 0 \\ 1}n\suchthat m,n\in\Re}$
    \end{exparts}
    \begin{answer}
      For each problem we get a system of linear equations by looking at the 
      equations of components.
      \begin{exparts}
       \partsitem $k=5$
       \partsitem The second components show that $i=2$, the third
       components show that $j=1$.
       \partsitem $m=-4$, $n=2$
      \end{exparts} 
    \end{answer}
  \item 
    Decide if the vector is in the set.
    \begin{exparts}
      \partsitem $\colvec[r]{3 \\ -1}$,
        $\set{\colvec[r]{-6 \\ 2}k\suchthat k\in\Re}$
      \partsitem $\colvec[r]{5 \\ 4}$,
        $\set{\colvec[r]{5 \\ -4}j\suchthat j\in\Re}$
      \partsitem $\colvec[r]{2 \\ 1 \\ -1}$,
        $\set{\colvec[r]{0 \\ 3 \\ -7}
             +\colvec[r]{1 \\ -1 \\ 3}r\suchthat r\in\Re}$
      \partsitem $\colvec[r]{1 \\ 0 \\ 1}$,
        $\set{\colvec[r]{2 \\ 0 \\ 1}j
            +\colvec[r]{-3 \\ -1 \\ 1}k\suchthat j,k\in\Re}$
    \end{exparts}
    \begin{answer}
      For each problem we get a system of linear equations by looking at the 
      equations of components.
      \begin{exparts}
        \partsitem Yes; take $k=-1/2$.
        \partsitem No; the system with equations $5=5\cdot j$ and
            $4=-4\cdot j$ has no solution.
        \partsitem Yes; take $r=2$.
        \partsitem No.
           The second components give $k=0$.
           Then the third components give $j=1$.
           But the first components don't check. 
      \end{exparts}
     \end{answer}
  \item \cite{Cleary}
    A farmer with 1200 acres is considering planting three different crops, 
    corn, soybeans, and oats.   
    The farmer wants to use all~$1200$ acres.  
    Seed corn costs \$$20$ per acre, while soybean and oat seed cost 
    \$$50$ and~\$$12$ per acre respectively.  
    The farmer has \$$40\,000$ available to buy seed and intends to 
    spend it all.
    \begin{exparts}  
      \item Use the information above to formulate two linear equations 
        with three unknowns and solve it.
     \item Solutions to the system are choices that the farmer can make.  
        Write down two reasonable solutions.
     \item Suppose that in the fall when the crops mature, the farmer 
        can bring in revenue of \$$100$ per acre for corn, 
        \$$300$ per acre for soybeans and \$$80$ per acre for oats.  
        Which of your two solutions in the prior part would have resulted 
        in a larger revenue? 
    \end{exparts}
    \begin{answer}
      \begin{exparts}
        \item Let $c$ be the number of acres of corn, $s$ be the number of 
          acres of soy, and $a$ be the number of acres of oats.
          \begin{equation*}
            \begin{linsys}{3}
              c   &+   &s   &+   &a   &=   &1200 \\ 
            20c   &+   &50s &+   &12a &=   &40\,000  
            \end{linsys}
            \grstep{-20\rho_1+\rho_2}
            \begin{linsys}{3}
              c   &+   &s   &+   &a   &=   &1200 \\ 
                  &    &30s &-   &8a  &=   &16\,000  
            \end{linsys}
          \end{equation*}
          To describe the solution set we can parametrize using $a$.
          \begin{equation*}
            \set{\colvec{c \\ s \\ a}
                 =\colvec{20\,000/30 \\ 16\,000/30 \\ 0}
                  +\colvec{-38/30 \\ 8/30 \\ 1}a
                 \suchthat a\in\Re}
          \end{equation*}
        \item There are many answers possible here. 
          For instance we can take $a=0$ to get $c=20\,000/30\approx 666.66$ and
          $s=16000/30\approx 533.33$.
          Another example is to take $a=20\,000/38\approx 526.32$, giving
          $c=0$ and $s=7360/38\approx 193.68$.
        \item Plug your answers from the prior part into 
          $100c+300s+80a$.
      \end{exparts}
    \end{answer}
  \item 
    Parametrize the solution set of this one-equation system.
    \begin{equation*}
      x_1+x_2+\cdots+x_n=0
    \end{equation*}
    \begin{answer}
      This system has one equation.
      The leading variable is \( x_1 \), the other variables are free.
      \begin{equation*}
        \set{\colvec{-1 \\ 1 \\ \vdotswithin{-1} \\ 0}x_2
             +\cdots+
             \colvec{-1 \\ 0 \\ \vdotswithin{-1} \\ 1}x_n
             \suchthat x_2,\ldots,x_n\in\Re}
      \end{equation*}  
     \end{answer}
  \recommended \item 
    \begin{exparts}
    \partsitem Apply Gauss's Method to the left-hand side to solve
      \begin{equation*}
        \begin{linsys}{4}
          x  &+  &2y  &    &    &-   &w   &=   &a   \\
         2x  &   &    &+   &z   &    &    &=   &b   \\
          x  &+  &y   &    &    &+   &2w  &=   &c   
        \end{linsys}
      \end{equation*}
      for \( x \), \( y \), \( z \), and \(  w \), in terms of the 
      constants $a$, $b$, and $c$.
    \partsitem Use your answer from the prior part to solve this.
      \begin{equation*}
        \begin{linsys}{4}
          x  &+  &2y  &    &    &-   &w   &=   &3   \\
         2x  &   &    &+   &z   &    &    &=   &1   \\
          x  &+  &y   &    &    &+   &2w  &=   &-2
        \end{linsys}
      \end{equation*}
    \end{exparts}
    \begin{answer}
      \begin{exparts}
        \partsitem Gauss's Method here gives
          \begin{align*}
            \begin{amat}{4}
              1  &2  &0  &-1  &a  \\
              2  &0  &1  &0   &b  \\
              1  &1  &0  &2   &c
            \end{amat}
            &\grstep[-\rho_1+\rho_3]{-2\rho_1+\rho_2}
            \begin{amat}{4}
              1  &2  &0  &-1  &a  \\
              0  &-4 &1  &2   &-2a+b  \\
              0  &-1 &0  &3   &-a+c
            \end{amat}                                  \\
            &\grstep{-(1/4)\rho_2+\rho_3}
            \begin{amat}{4}
              1  &2  &0    &-1  &a  \\
              0  &-4 &1    &2   &-2a+b  \\
              0  &0  &-1/4 &5/2 &-(1/2)a-(1/4)b+c
            \end{amat}
          \end{align*}
          leaving \( w \) free.
          Solve: \(  z=2a+b-4c+10w \),
          and \( -4y=-2a+b-(2a+b-4c+10w)-2w \) so
          \( y=a-c+3w \), and
          \( x=a-2(a-c+3w)+w=-a+2c-5w. \)
          Therefore the solution set is this.
          \begin{equation*}
             \set{\colvec{-a+2c \\ a-c \\ 2a+b-4c \\ 0}
                  +\colvec{-5 \\ 3 \\ 10 \\ 1}w
                  \suchthat w\in\Re}
          \end{equation*}
        \partsitem Plug in with \( a=3 \), \( b=1 \), and \( c=-2 \).
          \begin{equation*}
             \set{\colvec[r]{-7 \\ 5 \\ 15 \\ 0}
                  +\colvec[r]{-5 \\ 3 \\ 10 \\ 1}w
                  \suchthat w\in\Re}
          \end{equation*}
      \end{exparts}  
     \end{answer}
  \item 
    Why is the comma needed in the notation `\( a_{i,j} \)'
    for matrix entries?
    \begin{answer}
       Leaving the comma out, say by writing \( a_{123} \),
       is ambiguous because it could mean $a_{1,23}$ or $a_{12,3}$.  
    \end{answer}
  \recommended \item 
    Give the \( \nbyn{4} \) matrix whose
    \( i,j \)-th entry is
    \begin{exparts*}
      \partsitem \( i+j \);
      \partsitem \( -1 \) to the \( i+j \) power.
    \end{exparts*}
    \begin{answer}
      \begin{exparts*}
        \partsitem \(
           \begin{mat}[r]
             2  &3  &4  &5  \\
             3  &4  &5  &6  \\
             4  &5  &6  &7  \\
             5  &6  &7  &8
           \end{mat} \)
        \partsitem \(
           \begin{mat}[r]
             1  &-1  &1   &-1  \\
            -1  &1   &-1  &1  \\
             1  &-1  &1   &-1  \\
            -1  &1   &-1  &1
           \end{mat} \)
      \end{exparts*}  
    \end{answer}
  \item  
    For any matrix \( A \), the 
    \definend{transpose}\index{transpose}
    \index{matrix!transpose}
    of \( A \), written
    \( \trans{A} \), is the matrix whose columns are the rows of \( A \).
    Find the transpose of each of these.
    \begin{exparts*}
      \partsitem \( \begin{mat}[r]
                  1  &2  &3  \\
                  4  &5  &6
               \end{mat}  \)
      \partsitem \( \begin{mat}[r]
                  2  &-3 \\
                  1  &1
               \end{mat}  \)
      \partsitem \( \begin{mat}[r]
                  5  &10 \\
                 10  &5
               \end{mat}  \)
      \partsitem \( \colvec[r]{1 \\ 1 \\ 0} \)
    \end{exparts*}
    \begin{answer}
      \begin{exparts*}
        \partsitem \( \begin{mat}[r]
                   1  &4  \\
                   2  &5  \\
                   3  &6
                 \end{mat}  \)
        \partsitem \( \begin{mat}[r]
                   2  &1  \\
                  -3  &1
                 \end{mat}  \)
        \partsitem \( \begin{mat}[r]
                   5  &10 \\
                  10  &5
                 \end{mat}  \)
        \partsitem \( \rowvec{1 &1 &0}  \)
      \end{exparts*}  
     \end{answer}
  \recommended \item 
    \begin{exparts}
      \partsitem Describe all functions \( f(x)=ax^2+bx+c \) 
        such that \( f(1)=2 \) and \( f(-1)=6 \).
      \partsitem Describe all functions \( f(x)=ax^2+bx+c \) 
        such that \( f(1)=2 \).
    \end{exparts}
    \begin{answer}
      \begin{exparts}
        \partsitem Plugging in \( x=1 \) and \( x=-1 \) gives
          \begin{equation*}
            \begin{linsys}{3}
              a  &+  &b   &+  &c  &=  &2  \\
              a  &-  &b   &+  &c  &=  &6  
            \end{linsys}
            \grstep{-\rho_1+\rho_2}
            \begin{linsys}{3}
              a  &+  &b   &+  &c  &=  &2  \\
                 &   &-2b &   &   &=  &4  
              \end{linsys}
          \end{equation*}
          so the set of functions is
          \( \set{f(x)=(4-c)x^2-2x+c\suchthat c\in\Re} \).
        \partsitem Putting in \( x=1 \) gives
          \begin{equation*}
            \begin{linsys}{3}
              a  &+  &b   &+  &c  &=  &2  
            \end{linsys}
          \end{equation*}
          so the set of functions is
          \( \set{f(x)=(2-b-c)x^2+bx+c\suchthat b,c\in\Re} \).
      \end{exparts}  
    \end{answer}
  \item Show that any set of five points from the plane \( \Re^2 \) lie on a
    common conic section, that is, they all satisfy some equation of the
    form \( ax^2+by^2+cxy+dx+ey+f=0 \) where some of \( a,\,\ldots\,,f \)
    are nonzero.
    \begin{answer}
      On plugging in the five pairs $(x,y)$ we get a system with the
      five equations and six unknowns $a$, \ldots, $f$.
      Because there are more unknowns than equations, if no inconsistency
      exists among the equations then there are infinitely many solutions
      (at least one variable will end up free).

      But no inconsistency can exist because $a=0$, \ldots, $f=0$ is a 
      solution (we are only using this zero solution to show that the system
      is consistent\Dash the prior paragraph shows that
      there are nonzero solutions). 
    \end{answer}
  \item 
    Make up a four equations/four unknowns system having
    \begin{exparts}
      \partsitem a one-parameter solution set;
      \partsitem a two-parameter solution set;
      \partsitem a three-parameter solution set.
    \end{exparts}
    \begin{answer}
      \begin{exparts}
      \partsitem Here is one\Dash the fourth equation is redundant 
        but still OK.
        \begin{equation*}
          \begin{linsys}{4}
             x  &+  &y  &-  &z  &+  &w  &=  &0  \\
                &   &y  &-  &z  &   &   &=  &0  \\
                &   &   &   &2z &+  &2w &=  &0  \\
                &   &   &   &z  &+  &w  &=  &0
          \end{linsys}
        \end{equation*}
      \partsitem Here is one.
        \begin{equation*}
          \begin{linsys}{4}
             x  &+  &y  &-  &z  &+  &w  &=  &0  \\
                &   &   &   &   &   &w  &=  &0  \\
                &   &   &   &   &   &w  &=  &0  \\
                &   &   &   &   &   &w  &=  &0
          \end{linsys}
        \end{equation*}
      \partsitem This is one.
        \begin{equation*}
          \begin{linsys}{4}
             x  &+  &y  &-  &z  &+  &w  &=  &0  \\
             x  &+  &y  &-  &z  &+  &w  &=  &0  \\
             x  &+  &y  &-  &z  &+  &w  &=  &0  \\
             x  &+  &y  &-  &z  &+  &w  &=  &0 
          \end{linsys}
        \end{equation*}
    \end{exparts}  
   \end{answer}
  \puzzle \item 
    \cite{Shepelev}
    This  puzzle  is  from  a  Russian   web-site
    \texttt{http://www.arbuz.uz/}  and  there are many solutions
    to it, but mine uses  linear  algebra  and  is  very
    naive.   There's   a   planet  inhabited  by  arbuzoids
   (watermeloners, to  translate  from  Russian). 
   Those creatures are found in three colors: red, green and blue.  
   There  are  $13$~red
   arbuzoids,  $15$~blue  ones, and $17$~green. 
   When
   two differently colored arbuzoids meet,  they
   both change to the third color.

   The question is, can it ever happen that all
   of them assume the same color?
    \begin{answer}
       \answerasgiven
       My solution was to define the numbers  of  arbuzoids
       as $3$-dimensional vectors, and express all possible
       elementary transitions as such vectors, too:
       \begin{center}
         \begin{tabular}{rr}
           R: &$13$  \\
           G: &$15$  \\
           B: &$17$
         \end{tabular}
         \qquad
         Operations:
         $\colvec[r]{-1 \\ -1 \\ 2}$, 
         $\colvec[r]{-1 \\ 2 \\ -1}$, 
         and 
         $\colvec[r]{2 \\ -1 \\ -1}$
       \end{center}
       Now, it is enough to check whether the  solution  to
       one  of  the  following  systems of linear equations
       exists:
       \begin{equation*}
         \colvec[r]{13 \\ 15 \\ 17}
         +x\colvec[r]{-1 \\ -1  \\ 2}
         +y\colvec[r]{-1 \\ 2 \\ -1}
         +z\colvec[r]{2 \\ -1 \\ -1}
         =\colvec[r]{0 \\ 0 \\ 45}
         \qquad
         \text{(or $\colvec[r]{0 \\ 45 \\ 0}$ or $\colvec[r]{45 \\ 0 \\ 0}$)}
       \end{equation*}
       Solving
       \begin{equation*}
         \begin{amat}[r]{3}
          -1  &-1 &2  &-13  \\
          -1  &2  &-1 &-15  \\
           2  &-1 &-1 &28
         \end{amat}
         \grstep[2\rho_1+\rho_3]{-\rho_1+\rho_2}
         \repeatedgrstep{\rho_2+\rho_3}
         \begin{amat}[r]{3}
          -1  &-1 &2  &-13  \\
           0  &3  &-3 &-2  \\
           0  &0  &0  &0
         \end{amat}
       \end{equation*}
       gives $y+2/3=z$ so if the number of transformations $z$ is an integer
       then $y$ is not.
       The other two systems give similar conclusions so there is no
       solution.
    \end{answer}
  \puzzle \item 
    \cite{USSROlympiad174}
    \begin{exparts}
      \partsitem Solve the system of equations.
        \begin{equation*}
          \begin{linsys}{2}
            ax  &+  &y  &=  &a^2  \\
             x  &+  &ay &=  &1
         \end{linsys}
        \end{equation*}
        For what values of $a$ does the system fail to have solutions, and
        for what values of $a$ are there infinitely many solutions?
      \partsitem Answer the above question for the system.
        \begin{equation*}
          \begin{linsys}{2}
            ax  &+  &y  &=  &a^3  \\    
             x  &+  &ay &=  &1
          \end{linsys}
        \end{equation*}
    \end{exparts}
    \begin{answer}
       \answerasgiven
       \begin{exparts}
        \partsitem Formal solution of the system yields
          \begin{equation*}
            x=\frac{a^3-1}{a^2-1}  
            \qquad
            y=\frac{-a^2+a}{a^2-1}.
          \end{equation*}
          If $a+1\neq 0$ and $a-1\neq 0$, then the system has the single
          solution
          \begin{equation*}
            x=\frac{a^2+a+1}{a+1}
            \qquad
            y=\frac{-a}{a+1}.
          \end{equation*}
          If $a=-1$, or if $a=+1$, then the formulas are meaningless; in the
          first instance we arrive at the system
          \begin{equation*}
            \left\{ 
            \begin{linsys}{2}
              -x &+  &y  &=  &1 \\
               x &-  &y  &=  &1
            \end{linsys}\right.
          \end{equation*}
          which is a contradictory system.
          In the second instance we have
          \begin{equation*}
            \left\{
            \begin{linsys}{2}
               x &+  &y  &=  &1 \\
               x &+  &y  &=  &1
            \end{linsys}\right.
          \end{equation*}
          which has an infinite number of solutions (for example, for 
          $x$ arbitrary, $y=1-x$).
        \partsitem Solution of the system yields
          \begin{equation*}
            x=\frac{a^4-1}{a^2-1}
            \qquad
            y=\frac{-a^3+a}{a^2-1}.
          \end{equation*}
          Here, is $a^2-1\neq 0$, the system has the single solution
          $x=a^2+1$, $y=-a$.
          For $a=-1$ and $a=1$, we obtain the systems
          \begin{equation*}
            \left\{
            \begin{linsys}{2}
              -x &+  &y  &=  &-1 \\
               x &-  &y  &=  &1
            \end{linsys}\right.
            \qquad
            \left\{
            \begin{linsys}{2}
               x &+  &y  &=  &1 \\
               x &+  &y  &=  &1
            \end{linsys}\right.
         \end{equation*}
         both of which have an infinite number of solutions.
      \end{exparts}
    \end{answer}
  \puzzle \item 
    \cite{MathMag52p48}
    In air a gold-sur\-faced sphere weighs \( 7588 \)
    grams.
    It is known that it may contain one or more of the metals aluminum,
    copper, silver, or lead.
    When weighed successively under standard conditions in water, benzene,
    alcohol, and glycerin its respective weights are \( 6588 \), \( 6688 \),
    \( 6778 \), and \( 6328 \) grams.
    How much, if any, of the forenamed metals does it contain if the
    specific gravities of the designated substances are taken to be as follows?
    \begin{center}
       \begin{tabular}{lrclr}
         Aluminum  &\( 2.7 \)
            &\makebox[3em]{\mbox{}\hfill\mbox{}} &Alcohol &0.81 \\
         Copper    &\( 8.9 \)  &     &Benzene   &\( 0.90 \) \\
         Gold      &\( 19.3 \) &     &Glycerin &\( 1.26 \) \\
         Lead      &\( 11.3 \) &     &Water     &\( 1.00 \) \\
         Silver    &\( 10.8 \)
       \end{tabular}
    \end{center}
    \begin{answer}
      \answerasgiven
      Let \( u \), \( v \), \( x \), \( y \), \( z \) be the volumes in
      \( {\rm cm}^3 \) of Al, Cu, Pb, Ag, and Au, respectively, contained in
      the sphere, which we assume to be not hollow.
      Since the loss of weight in water (specific gravity \( 1.00 \)) is
      \( 1000 \) grams, the volume of the sphere is \( 1000\mbox{ cm}^3 \).
      Then the data, some of which is superfluous, though consistent, leads to
      only \( 2 \) independent equations, one relating volumes and the
      other, weights.
      \begin{equation*}
        \begin{linsys}{5}
           u  &+  &v    &+  &x     &+  &y     &+  &z     &=  &1000  \\
        2.7u  &+  &8.9v &+  &11.3x &+  &10.5y &+  &19.3z &=  &7558
        \end{linsys}
      \end{equation*}
      Clearly the sphere must contain some aluminum to bring its mean specific
      gravity below the specific gravities of all the other metals.
      There is no unique result to this part of the problem, for the amounts
      of three metals may be chosen arbitrarily, provided that the choices
      will not result in negative amounts of any metal.

      If the ball contains only aluminum and gold, there are
      \( 294.5\mbox{ cm}^3 \) of gold and \( 705.5\mbox{ cm}^3 \) of aluminum.
      Another possibility is \( 124.7\mbox{ cm}^3 \) each of Cu, Au, Pb, and
      Ag and \( 501.2\mbox{ cm}^3  \) of Al.   
    \end{answer}
\end{exercises}

























\subsection{\texorpdfstring{$\text{General}=\text{Particular}+\text{Homogeneous}$}{General=Particular+Homogeneous}}
In the prior subsection the descriptions of solution sets
all fit a pattern.
They have a vector that is a particular solution 
of the system added to an unrestricted combination of some other vectors.
The solution set from 
\nearbyexample{ex:ManyParamsInfManySolsSystem} illustrates.
\begin{equation*}
  \set{
   \underbracket[.7pt]{
     \colvec[r]{0 \\ 4 \\ 0 \\ 0 \\ 0}}_{\text{\shortstack{\rule{0pt}{2ex}particular \\
                                                    solution}}}+
   \underbracket[.7pt]{w\colvec[r]{1 \\ -1 \\ 3 \\ 1 \\ 0}+
       u\colvec[r]{1/2 \\ -1 \\ 1/2 \\ 0 \\ 1}}_{\text{\shortstack{\rule{0pt}{2ex}unrestricted\\
                                                                combination}}}
       \suchthat w,u\in\Re}
\end{equation*}
The combination is unrestricted in that 
$w$ and $u$ can be any real numbers\Dash there
is no condition like ``such that $2w-u=0\mspace{1mu}$'' to restrict 
which pairs $w,u$ we can use.

That example shows an infinite solution set fitting the pattern.
The other two kinds of solution sets also fit.
A one-element solution set fits because it 
has a particular solution
and the unrestricted combination part is trivial. 
That is, instead of being a combination of two vectors or
of one vector, it is a combination of no vectors.
(By convention the sum of an empty set of vectors
is the zero vector.)
An empty solution set fits the pattern because there is no 
particular solution and thus there are no sums of that form.

\begin{theorem} \label{th:GenEqPartPlusHomo}
%<*th:GenEqPartPlusHomo>
Any linear system's 
solution set has the form 
\begin{equation*}
   \set{\vec{p}+c_1\vec{\beta}_1+\,\cdots\,+c_k\vec{\beta}_k
     \suchthat c_1,\,\ldots\,,c_k\in\Re}
\end{equation*}
where \( \vec{p} \) is any particular solution  
and where the number of vectors 
$\vec{\beta}_1$, \ldots, $\vec{\beta}_k$ equals
the number of free variables that the system has after a Gaussian reduction.
%</th:GenEqPartPlusHomo>
\end{theorem}

The solution description has two parts, 
the particular solution $\vec{p}$ 
and the unrestricted linear combination of the $\vec{\beta}$'s.
We shall prove the theorem with two corresponding lemmas.

We will focus first on the unrestricted combination.
For that we consider systems that have the vector of zeroes
as a particular solution
so that we can shorten $\vec{p}+c_1\vec{\beta}_1+\dots+c_k\vec{\beta}_k$
to $c_1\vec{\beta}_1+\dots+c_k\vec{\beta}_k$.

\begin{definition} \label{df:HomogeneousEquation}
%<*df:HomogeneousEquation>
A linear equation is \definend{homogeneous}\index{homogeneous equation}%
\index{linear equation!homogeneous} if it has a constant of zero, so
that it can be written as $a_1x_1+a_2x_2+\,\cdots\,+a_nx_n=0$.
%</df:HomogeneousEquation>
\end{definition}

\begin{example}  \label{ex:FirstExHomoSys}
With any linear system like
\begin{equation*}
  \begin{linsys}{2}
    3x  &+  &4y  &=  3  \\
    2x  &-  &y   &=  1  
  \end{linsys}
\end{equation*}
we associate a system of homogeneous equations by setting the right side to
zeros.
\begin{equation*}
  \begin{linsys}{2}
    3x  &+  &4y  &=  0  \\
    2x  &-  &y   &=  0  
  \end{linsys}
\end{equation*}
Compare the reduction of the original system
\begin{equation*}
  \begin{linsys}{2}
    3x  &+  &4y  &=  3  \\
    2x  &-  &y   &=  1  
  \end{linsys}
  \grstep{-(2/3)\rho_1+\rho_2}
  \begin{linsys}{2}
    3x  &+  &4y        &=  3  \\
        &   &-(11/3)y   &=  -1  
   \end{linsys}
\end{equation*}
with the reduction of the associated homogeneous system.
\begin{equation*}
  \begin{linsys}{2}
    3x  &+  &4y  &=  0  \\
    2x  &-  &y   &=  0  
  \end{linsys}
  \grstep{-(2/3)\rho_1+\rho_2}
  \begin{linsys}{2}
    3x  &+  &4y        &=  0  \\
        &   &-(11/3)y   &=  0
   \end{linsys}
\end{equation*}
Obviously the two reductions go in the same way.
We can study how to reduce a linear system by instead studying how
to reduce the associated homogeneous system.
\end{example}

Studying the associated homogeneous system has a great advantage over
studying the original system.
Nonhomogeneous systems can be inconsistent.
But a homogeneous system must be consistent since there is always at least
one solution, the zero vector.


\begin{example} \label{ex:HomoZeroOnlySol}
Some homogeneous systems have the zero vector as their only solution.
\begin{equation*}
   \begin{linsys}{3}
     3x  &+  &2y  &+  &z  &=  &0  \\
     6x  &+  &4y  &   &   &=  &0  \\
         &   &y   &+  &z  &=  &0  
   \end{linsys}
   \grstep{-2\rho_1 +\rho_2}
   \begin{linsys}{3}
      3x  &+  &2y  &+  &z  &=  &0  \\
          &   &    &   &-2z&=  &0  \\
          &   &y   &+  &z  &=  &0  
    \end{linsys}
   \grstep{\rho_2 \leftrightarrow\rho_3}
   \begin{linsys}{3}
      3x  &+  &2y  &+  &z  &=  &0  \\
          &   &y   &+  &z  &=  &0  \\
          &   &    &   &-2z&=  &0
    \end{linsys}
\end{equation*}
\end{example}

\begin{example} \label{ex:SolnChemProb}
Some homogeneous systems have many solutions.
One is the Chemistry problem\index{Chemistry problem} 
from the first page of the first subsection.
\begin{align*}
  \begin{linsys}{4}
              7x  &   &   &-  &7z  &   &   &=  &0  \\
              8x  &+  &y  &-  &5z  &-  &2w &=  &0  \\
                  &   &y  &-  &3z  &   &   &=  &0  \\
                  &   &3y &-  &6z  &-  &w  &=  &0  
  \end{linsys}
  &\grstep{-(8/7)\rho_1+\rho_2}
  \begin{linsys}{4}
                7x &   &   &-  &7z  &   &   &=  &0  \\
                   &   &y  &+  &3z  &-  &2w &=  &0  \\
                   &   &y  &-  &3z  &   &   &=  &0  \\
                   &   &3y &-  &6z  &-  &w  &=  &0  
   \end{linsys}                                        \\
  &\grstep[-3\rho_2+\rho_4]{-\rho_2+\rho_3}
  \begin{linsys}{4}
                7x &   &   &-  &7z  &   &   &=  &0  \\
                   &   &y  &+  &3z  &-  &2w &=  &0  \\
                   &   &   &   &-6z &+  &2w &=  &0  \\
                   &   &   &   &-15z&+  &5w &=  &0  
   \end{linsys}                                        \\
  &\grstep{-(5/2)\rho_3+\rho_4}
  \begin{linsys}{4}
                7x &   &   &-  &7z  &   &   &=  &0  \\
                   &   &y  &+  &3z  &-  &2w &=  &0  \\
                   &   &   &   &-6z &+  &2w &=  &0  \\
                   &   &   &   &    &   &0  &=  &0  
   \end{linsys}
\end{align*}
The solution set
\begin{equation*}
  \set{\colvec[r]{1/3 \\ 1 \\ 1/3 \\ 1}w \suchthat w\in\Re}
\end{equation*}
has many vectors besides the zero vector
(if we take \( w \) to be a number of molecules then solutions
make sense only when \( w \) is a nonnegative multiple of $3$).
\end{example}

\begin{lemma} \label{le:HomoSltnSpanVecs}
%<*le:HomoSltnSpanVecs>
For any  homogeneous linear system there exist
vectors $\vec{\beta}_1$, \ldots, $\vec{\beta}_k$ such that the 
solution set of the system is
\begin{equation*}
  \set{c_1\vec{\beta}_1+\cdots+c_k\vec{\beta}_k \suchthat c_1,\ldots,c_k\in\Re}
\end{equation*}
where $k$ is the
number of free variables in an echelon form version of the system.
%</le:HomoSltnSpanVecs>
\end{lemma}

Before the proof, we discuss a couple of points about it.
First, the proof verifies that, given a homogeneous system in echelon form,
we can use back substitution to express 
all of the leading variables in terms of the free variables.
To see why this suffices to establish the lemma, we give an example
(we will also use this walk-through as a chance to 
illustrate the proof's notation).

%<*ex:HomoSystemGivesSpanningSet0>
Consider this system of homogeneous equations in echelon form.
\begin{equation*}
  \begin{linsys}{5}
     x  &+  &y   &+  &2z  &+  &u  &+  &v  &=  &0  \\
        &   &y   &+  &z  &+  &u  &-  &v  &=  &0  \\
        &   &    &   &   &   &u  &+  &v  &=  &0  
  \end{linsys}
\end{equation*}
It has five variables: $x_1=x$, \ldots{}, $x_5=v$.
To do back substitution, we start with the bottom row, equation~$m=3$.
We solve for its leading variable, $x_{\ell_3}=x_4=u$, using its
leading term $a_{m,\ell_m}x_m=a_{3,4}x_4=1\cdot u$
(the notation's `$\ell$' stands for ``leading''). 
We get $u=-v$.
%</ex:HomoSystemGivesSpanningSet0>

Now back substitution enters its iteration phase.
%<*ex:HomoSystemGivesSpanningSet1>
We will use the variable~$t$ to count how many rows up from the bottom 
the current row is, so we take $t=1$ and
look at row $m-t=3-1=2$.
We come to this row having so far expressed the leading variable 
in terms of free variables for row~$3$.
We substitute for $x_{\ell_3}=x_4=u$ 
its expression in terms
of free variables, getting $y+z+(-v)-v=0$.
Then solving for this row's leading variable gives
$x_{\ell_{m-t}}=x_2=y=-z+2v$.

Next is the row that is $t=2$ up from the bottom, row $m-t=1$.
We have so far expressed the leading variable 
in terms of free variables for rows $3$ and~$2$.
Substitute for $x_{\ell_3}=u$ and $x_{\ell_2}=y$, giving 
$x+(-z+2v)+2z+(-v)+v=0$.
Then solve for the leading variable in terms of the free variables
as $x_{\ell_1}=x=-z-2v$.
%</ex:HomoSystemGivesSpanningSet2>

%<*ex:HomoSystemGivesSpanningSet3>
The point of this example is that now we are done. 
Just write the solution in vector notation
\begin{equation*}
  \colvec{x \\ y \\ z \\ u \\ v}
  =\colvec{-1 \\ -1 \\ 1 \\ 0 \\ 0}z
   +\colvec{-2 \\ 2 \\ 0 \\ -1 \\ 1}v 
  \qquad \text{where $z,v\in\Re$}
\end{equation*}
to recognize that $\vec{\beta}_1$ and $\vec{\beta}_2$ 
of the lemma's statement are the 
vectors associated with the free variables $z$ and~$v$.
%</ex:HomoSystemGivesSpanningSet3>

The second point about the proof follows from that example also.
Back substitution moves up the system, using the conclusions
from lower rows to do the next row.
This suggests the proof strategy of mathematical 
induction.\index{mathematical induction}\index{proof techniques!induction}\index{induction}%

Induction is an important and non-obvious proof technique that will
appear a number of times in this book.
Our proofs by induction will come in two steps, a base step and
an inductive step.
In the base step, we verify that the statement is true for some first 
instance, here that for an echelon form linear system we can write 
the bottom equation's leading variable in terms of free variables.
In the inductive step, we must establish an implication, 
namely that if the statement
is true for all prior cases then it follows 
that the statement holds for the present case 
also.
Specifically, here the inductive hypothesis is that if for the
lower-down rows we can express the leading variables 
in terms of the free variables, then for the next row up
we can express the leading variable in those terms also.

Those two steps together prove the statement for all the rows 
because by the base step it is true for the bottom equation, 
and by the inductive step the fact that it is true for the bottom
equation shows that 
it is true for the next one up. 
Then, another application of
the inductive step implies that it is true for the third equation up,
etc.\appendrefs{mathematical induction}%

\begin{proof}
%<*pf:HomoSltnSpanVecs0>
Apply Gauss's Method to get to echelon form.
There may be some \( 0=0 \) equations; we ignore these 
(if the
system consists only of \( 0=0 \)~equations then the lemma is trivially true
because there are no leading variables).
% If there are any variables that appear only with with a coefficient of
% $0$ then they are free, and will never lead a row, so we will
% ignore that possibility.  
But
because the system is homogeneous there are no contradictory equations.

We will use induction to verify that
each leading variable can be expressed in terms of
free variables.
That will finish the proof because 
we can use the free variables as parameters and
the $\vec{\beta}$'s are the vectors of coefficients of those
free variables. 
%</pf:HomoSltnSpanVecs0>

%<*pf:HomoSltnSpanVecs1>
For the base step, consider the bottom-most equation,
\begin{equation*}
  a_{m,\ell_m}x_{\ell_m}+a_{m,1+\ell_{m}}x_{1+\ell_{m}}+\cdots+a_{m,n}x_n=0
   % \tag{$*$}
\end{equation*}
where \( a_{m,\ell_m}\neq 0 \).
%</pf:HomoSltnSpanVecs1>
%<*pf:HomoSltnSpanVecs2>
This is the bottom row, so 
any variables % $x_{\ell_{m}+1}$, \ldots{} 
after the
leading one must be free. 
Move these to the right hand side and divide by $a_{m,\ell_m}$
\begin{equation*}
  x_{\ell_m}
  =(-a_{m,1+\ell_{m}}/a_{m,\ell_m})x_{1+\ell_{m}}+\cdots+(-a_{m,n}/a_{m,\ell_m})x_n
\end{equation*}
to express the leading variable in terms of free variables.
%</pf:HomoSltnSpanVecs2>
(A technical point:~if in the bottom equation
there are no variables to the right 
of~$x_{\ell_m}$ then \( x_{\ell_m}=0 \).
This satisfies the statement that we are verifying because, 
as alluded to at the start of this subsection, 
it has \( x_{\ell_m} \) written 
as a sum of a number of the free variables, namely as the sum of zero many, 
under the convention that 
a trivial sum totals to~$0$.)

%<*pf:HomoSltnSpanVecs3>
For the inductive step, assume that assume 
that for the \( m \)-th equation,
and the \text{\( (m-1) \)-th}
equation, etc., up to and including the \mbox{\( m-(t-1) \)-th}~equation,
we can
express the leading variable of that equation in terms of free variables. 
That is,
assume that the statement that we can
express the leading variable of the row in terms of the free variables
holds for the bottom-most $t$~rows, with $1\leq t<m$.
We must verify that the statement 
therefore also holds for the next equation up, 
the \( (m-t) \)-th equation.

For each of 
\( x_{\ell_m} \), \ldots, \( x_{\ell_{m-(t-1)}} \), 
substitute its expression in terms of free variables.
The result has a leading term of   
$a_{m-t,\ell_{m-t}}x_{\ell_{m-t}}$ with \( a_{m-t,\ell_{m-t}}\neq 0 \),
and the rest of the left hand side
is a combination of free variables.
Move the free variables to
the right side and divide by
\( a_{m-t,\ell_{m-t}} \), to end with
\( x_{\ell_{m-t}} \) expressed in terms of free variables.

We have done both the base step and the inductive step, so by the
principle of mathematical induction the proposition is true.
%</pf:HomoSltnSpanVecs3>
\end{proof}

This shows, 
as discussed between the lemma and its proof, that 
we can parametrize solution sets
using the free variables.
We say that the set of vectors
$\set{c_1\vec{\beta}_1+\cdots+c_k\vec{\beta}_k \suchthat c_1,\ldots,c_k\in\Re}$
is \definend{generated by}\index{generated by}\index{generated} 
or~\definend{spanned by}\index{spanned by}\index{span} 
the set 
\( \set{\smash{\vec{\beta}_1},\ldots,\smash{\vec{\beta}_k}} \).

% The proof mentions a tricky point.
% We follow the convention that the sum of an empty set of vectors is the 
% zero vector.
% In particular, we needs this where 
% a homogeneous system has a unique solution because it
% takes the \( c \)'s to be the free variables
% and if there is a unique solution then there are no free variables.

To finish the proof of  \nearbytheorem{th:GenEqPartPlusHomo}
the next lemma considers the particular solution part of the 
solution set's description.

\begin{lemma}  \label{th:GenEqPartHomo}
%<*th:GenEqPartHomo>
For a linear system and for any particular solution $\vec{p}\/$,
the solution set equals
% \begin{equation*}
$
  \set{\vec{p}+\vec{h} \suchthat \text{ \( \vec{h} \) satisfies the
                                associated homogeneous system}     }$.
%\end{equation*}
%</th:GenEqPartHomo>
\end{lemma}

% So fixing any particular solution gives the above 
% description of the solution set.

\begin{proof}
We will show mutual set inclusion, that any solution to the system is in
the above set and that anything in the set is a solution of the 
system.\appendrefs{set equality}

%<*pf:GenEqPartHomo0>
For set inclusion the first way, that if a vector solves the system
then it is in the set described above, 
assume that \( \vec{s} \) solves the system.
Then \( \vec{s}-\vec{p} \) solves the associated
homogeneous system since for each equation index \( i \),
\begin{multline*}
  a_{i,1}(s_1-p_1)+\cdots+a_{i,n}(s_n-p_n) \\
  =(a_{i,1}s_1+\cdots+a_{i,n}s_n)       
  -(a_{i,1}p_1+\cdots+a_{i,n}p_n)  
  =d_i-d_i                 
  =0
\end{multline*}
where \( p_j \) and \( s_j \) are the \( j \)-th components of
\( \vec{p} \) and \( \vec{s} \).
Express \( \vec{s} \) in the required \( \vec{p}+\vec{h} \) form
by writing \( \vec{s}-\vec{p} \) as \( \vec{h} \).
%</pf:GenEqPartHomo0>

%<*pf:GenEqPartHomo1>
For set inclusion the other way, take a vector of the form $\vec{p}+\vec{h}$,
where \( \vec{p} \) solves the system and \( \vec{h} \) solves the
associated homogeneous system and note that $\vec{p}+\vec{h}$ 
solves the given system since for any equation index~$i$, 
\begin{multline*}
  a_{i,1}(p_1+h_1)+\cdots+a_{i,n}(p_n+h_n)  \\
  =(a_{i,1}p_1+\cdots+a_{i,n}p_n)      
   +(a_{i,1}h_1+\cdots+a_{i,n}h_n)  
  =d_i+0                                
  =d_i
\end{multline*}
where as earlier \( p_j \) and \( h_j \) are the \( j \)-th components of 
\( \vec{p} \) and \( \vec{h} \).
%</pf:GenEqPartHomo1>
\end{proof}

The two lemmas together establish \nearbytheorem{th:GenEqPartPlusHomo}.
Remember that theorem with the slogan, 
``\( \text{General} = \text{Particular} + \text{Homogeneous} \)''.

\begin{example} \label{ex:IllusGenEqPartHomo}
This system illustrates \nearbytheorem{th:GenEqPartPlusHomo}.
\begin{equation*}
  \begin{linsys}{3}
    x  &+  &2y  &-  &z  &=  &1  \\
    2x &+  &4y  &   &   &=  &2  \\
       &   &y   &-  &3z &=  &0
  \end{linsys}
\end{equation*}
Gauss's Method
\begin{equation*}
  \grstep{-2\rho_1+\rho_2}
  \begin{linsys}{3}
    x  &+  &2y  &-  &z  &=  &1  \\
       &   &    &   &2z &=  &0  \\
       &   &y   &-  &3z &=  &0
  \end{linsys}                           
  \grstep{\rho_2\leftrightarrow\rho_3} 
  \begin{linsys}{3}
      x  &+  &2y  &-  &z  &=  &1  \\      
         &   &y   &-  &3z &=  &0  \\
         &   &    &   &2z &=  &0
   \end{linsys}
\end{equation*}
shows that the general solution is a singleton set.
\begin{equation*}
  \set{\colvec[r]{1 \\ 0 \\ 0} }
\end{equation*}
That single vector is obviously a particular solution.
The associated homogeneous system reduces via the same row operations 
\begin{equation*}
  \begin{linsys}{3}
    x  &+  &2y  &-  &z  &=  &0  \\
    2x &+  &4y  &   &   &=  &0  \\
       &   &y   &-  &3z &=  &0
  \end{linsys}
  \grstep{-2\rho_1+\rho_2}
  \repeatedgrstep{\rho_2\swap\rho_3} 
  \begin{linsys}{3}
      x  &+  &2y  &-  &z  &=  &0  \\      
         &   &y   &-  &3z &=  &0  \\
         &   &    &   &2z &=  &0
   \end{linsys}
\end{equation*}
to also give a singleton set. 
\begin{equation*}
  \set{\colvec[r]{0 \\ 0 \\ 0} }
\end{equation*}
So, as discussed at the start of this subsection, 
in this single-solution case the general solution results 
from taking the particular solution and adding to it the unique solution
of the associated homogeneous system.
\end{example}

\begin{example}
The start of this subsection also discusses that the case where
the general solution set is empty fits the
$\text{General}=\text{Particular}+\text{Homogeneous}$ pattern too.
This system illustrates.
\begin{equation*}
  \begin{linsys}{4}
    x  &   &  &+  &z  &+ &w  &=  &-1  \\
   2x  &-  &y &   &   &+ &w  &=  &3   \\
    x  &+  &y &+  &3z &+ &2w &=  &1   
  \end{linsys}
  \grstep[-\rho_1+\rho_3]{-2\rho_1+\rho_2}
  \begin{linsys}{4}
    x  &   &  &+  &z  &+ &w  &=  &-1  \\
       &   &-y&-  &2z &- &w  &=  &5   \\
       &   &y &+  &2z &+ &w  &=  &2   
   \end{linsys}
\end{equation*}
It has no solutions because the final two equations
conflict.
But the associated homogeneous system does have a solution, as do all 
homogeneous systems.
\begin{equation*}
  \begin{linsys}{4}
    x  &   &  &+  &z  &+ &w  &=  &0   \\
   2x  &-  &y &   &   &+ &w  &=  &0   \\
    x  &+  &y &+  &3z &+ &2w &=  &0   
  \end{linsys}
  \grstep[-\rho_1+\rho_3]{-2\rho_1+\rho_2}
  \grstep{\rho_2+\rho_3}
  \begin{linsys}{4}
    x  &   &  &+  &z  &+ &w  &=  &0   \\
       &   &-y&-  &2z &- &w  &=  &0   \\
       &   &  &   &   &  &0  &=  &0         
  \end{linsys}
\end{equation*}
In fact, the solution set is infinite. 
\begin{equation*}
  \set{\colvec[r]{-1 \\ -2 \\ 1 \\ 0}z+\colvec[r]{-1 \\ -1 \\ 0 \\ 1}w
         \suchthat z,w\in\Re}
\end{equation*}
Nonetheless, because the original system has no particular solution, its
general solution set is empty\Dash there are no vectors of the form
$\vec{p}+\vec{h}$ because there are no $\vec{p}\:$'s.
\end{example}

\begin{corollary} \label{co:ThreeKindsSolutionSets}
%<*co:ThreeKindsSolutionSets>
Solution sets of linear systems are either empty, have one element, or
have infinitely many elements.
%</co:ThreeKindsSolutionSets>
\end{corollary}

\begin{proof}
%<*pf:ThreeKindsSolutionSets0>
We've seen examples of all three happening so we need only prove
that there are no other possibilities.

First observe a homogeneous system with
at least one non-\( \zero \) solution $\vec{v}$ has infinitely many
solutions.
This is because any scalar multiple of~$\vec{v}$ also solves the homogeneous
system and there are infinitely many vectors in the set of scalar 
multiples of $\vec{v}$: if $s,t\in\Re$ are unequal then $s\vec{v}\neq t\vec{v}$,
since $s\vec{v}-t\vec{v}=(s-t)\vec{v}$ is
non-$\zero$ as  any non-$0$ component of $\vec{v}$, when 
rescaled by the non-$0$ factor $s-t$, will give a non-$0$ value.
%</pf:ThreeKindsSolutionSets0>

%<*pf:ThreeKindsSolutionSets1>
Now apply \nearbylemma{th:GenEqPartHomo} to conclude that a solution set
\begin{equation*}
  \set{\vec{p}+\vec{h}\suchthat
    \text{\( \vec{h} \) solves the associated homogeneous system}}
\end{equation*}
is either empty (if there is no particular solution \( \vec{p} \)),
or has one element (if there is a \( \vec{p} \) and the homogeneous system
has the unique solution \( \zero \)), or is infinite (if there is a
\( \vec{p} \) and the homogeneous system has a non-$\zero$ solution,
and thus by the prior paragraph has infinitely many solutions).
%</pf:ThreeKindsSolutionSets1>
\end{proof}

This table summarizes the factors affecting the size of a
general solution.

\smallskip
%<*table:KindsSolutionSets>
\begin{center} % \small 
\begin{tabular}{r@{}c}
  &\hspace*{2.5em}\begin{tabular}{c} 
      \textit{number of solutions of the} \\[-.5ex]
      \textit{homogeneous system}
    \end{tabular}                            \\[1.7ex] 
  \begin{tabular}{r@{\hspace*{.5em}}} 
     \ \\[.6ex]
     \textit{particular} \\[-.55ex]
     \textit{solution}   \\[-.5ex]
     \textit{exists?}   
  \end{tabular}
  &\begin{tabular}{r|c@{\hspace*{1em}}c} % \cline{2-3}   
     \multicolumn{1}{c}{\ }
         &\textit{one}    &\textit{infinitely many}                    \\ 
     \cline{2-3}
     \textit{yes}    
        &\rule{0ex}{16pt}\begin{tabular}{@{}c@{}} unique \\[-.5ex] solution \end{tabular}
        &\begin{tabular}{@{}c@{}} infinitely many \\[-.5ex] solutions \end{tabular}
         \\[2ex] % \hline  
     \textit{no}    
        &\begin{tabular}{@{}c@{}} no \\[-.5ex] solutions \end{tabular}
        &\begin{tabular}{@{}c@{}} no \\[-.5ex] solutions \end{tabular} 
         \\ % \hline
   \end{tabular}
\end{tabular}
\end{center}
%</table:KindsSolutionSets>
\smallskip

The dimension on the top of the table is the simpler one.
When we perform Gauss's Method on a linear system, ignoring the
constants on the right side and so paying attention only
to the coefficients on the left-hand side,
we either end with every variable leading some row or else 
we find some variable that does not lead a row, that is,
we find some variable that is free. 
(We formalize ``ignoring the constants on the right'' by
considering the associated homogeneous system.)

A notable special case is
systems having the same number of equations as unknowns. 
Such a system will have a solution, and that solution will be unique, 
if and only if it
reduces to an echelon form system where every variable leads its row
(since there are the same number of variables as rows),
which will happen if and only if
the associated homogeneous system has a unique solution.

\begin{definition} \label{df:Nonsingular}
%<*df:Nonsingular>
A square matrix is \definend{nonsingular}\index{nonsingular!matrix}
\index{matrix!nonsingular}
if it is the matrix of coefficients of a
homogeneous system with a unique solution.
It is
\definend{singular}\index{singular!matrix}\index{matrix!singular} otherwise,
that is,
if it is the matrix of coefficients of a homogeneous system with 
infinitely many solutions.
%</df:Nonsingular>
\end{definition}

\begin{example}
The first of these matrices is nonsingular while the second is singular
\begin{equation*}
  \begin{mat}[r]
    1  &2  \\
    3  &4
  \end{mat}
  \qquad
  \begin{mat}[r]
    1  &2  \\
    3  &6
  \end{mat}
\end{equation*}
because the first of these homogeneous systems has a unique solution 
while the second has infinitely many solutions.
\begin{equation*}
  \begin{linsys}[b]{2}
    x &+  &2y  &=  &0  \\
   3x &+  &4y  &=  &0  
  \end{linsys}
  \qquad
  \begin{linsys}[b]{2}
    x &+  &2y  &=  &0  \\
   3x &+  &6y  &=  &0
  \end{linsys}
\end{equation*}  
We have made the distinction in the definition because a system
with the same number of equations as variables
behaves in one of two ways, depending on whether its matrix of coefficients
is nonsingular or singular.
Where the matrix of coefficients is nonsingular the system 
has a unique solution for any constants on the right 
side:~for instance, Gauss's Method shows that this system
\begin{equation*}
  \begin{linsys}{2}
    x  &+  &2y  &=  &a \\
    3x &+  &4y  &=  &b
  \end{linsys}
\end{equation*}
has the unique solution $x=b-2a$ and  $y=(3a-b)/2$.
On the other hand, where the matrix of coefficients is
singular the system never has a unique solution\Dash it 
has either no solutions or else has infinitely many, as with these.
\begin{equation*}
  \begin{linsys}[b]{2}
    x  &+  &2y  &=   &1   \\
   3x  &+  &6y  &=   &2   
  \end{linsys}
  \qquad
  \begin{linsys}[b]{2}
    x  &+  &2y  &=   &1   \\
   3x  &+  &6y  &=   &3
  \end{linsys}
\end{equation*} 
\end{example}

The definition uses the word `singular' because it 
means ``departing from general expectation.''
People often, naively, expect that systems 
with the same number of variables as equations will have a unique solution.
Thus, we can think of the word as connoting 
``troublesome,'' or at least ``not ideal.''
(That `singular' applies to those systems that never have exactly one solution 
is ironic but it is the standard term.)

\begin{example}
The systems from \nearbyexample{ex:FirstExHomoSys},
\nearbyexample{ex:HomoZeroOnlySol},
and \nearbyexample{ex:IllusGenEqPartHomo}
each have an associated homogeneous system with a unique solution.
Thus these matrices are nonsingular.
\begin{equation*}
  \begin{mat}[r]
    3  &4  \\
    2  &-1
  \end{mat}
  \qquad
  \begin{mat}[r]
    3  &2   &1  \\
    6  &-4  &0  \\
    0  &1   &1
  \end{mat}
  \qquad
  \begin{mat}[r]
    1  &2  &-1 \\
    2  &4  &0  \\
    0  &1  &-3
  \end{mat}
\end{equation*}
The Chemistry problem from \nearbyexample{ex:SolnChemProb} 
is a homogeneous system with more than one solution so its matrix
is singular. 
\begin{equation*}
  \begin{mat}[r]
    7  &0  &-7 &0  \\
    8  &1  &-5 &-2 \\
    0  &1  &-3 &0  \\
    0  &3  &-6 &-1
  \end{mat}
\end{equation*}
\end{example}

The table above has two dimensions.
We have considered the one on top:~we can tell
into which column a given linear system goes
solely by considering the system's left-hand side; the 
constants on the right-hand side play no role in this.

The table's other dimension, 
determining whether a particular solution exists, is tougher.
Consider these two systems with the same left side but different right sides.
\begin{equation*}
  \begin{linsys}[b]{2}
    3x &+ &2y &= &5  \\
    3x &+ &2y &= &5
  \end{linsys}
  \qquad
  \begin{linsys}[b]{2}
    3x &+ &2y &= &5  \\
    3x &+ &2y &= &4
  \end{linsys}
\end{equation*}
The first has a solution while the second does not, so
here the constants on the right side decide if the system has a solution.
We could conjecture that the left side of a linear system determines
the number of solutions while the right side determines if solutions
exist but that guess is not correct.
Compare these two,
with the same right sides but different left sides.
\begin{equation*}
  \begin{linsys}[b]{2}
    3x &+ &2y &= &5  \\
    4x &+ &2y &= &4
  \end{linsys}
  \qquad
  \begin{linsys}[b]{2}
    3x &+ &2y &= &5  \\
    3x &+ &2y &= &4
  \end{linsys}
\end{equation*}
The first has a solution but the second does not.
Thus the constants on the right side of a system 
don't alone determine whether a solution exists.
Rather, that depends on some interaction between the left and
right.

For some intuition about that interaction,
consider this system with one of the coefficients left unspecified, as  
the variable~$c$.
\begin{equation*}
  \begin{linsys}{3}
    x  &+  &2y  &+  &3z  &=  &1  \\
    x  &+  &y   &+  &z   &=  &1  \\
   cx  &+  &3y  &+  &4z  &=  &0
  \end{linsys}
\end{equation*}
If \( c=2 \) then this system has no solution because the left-hand side 
has the third row as the sum of the first two, while the right-hand does not.
If \( c\neq 2 \) then this system has a unique solution (try it with \( c=1 \)).
For a system to have a solution, if one row of the matrix of coefficients on
the left is a linear combination of other rows
then on the right the constant from that row must be the same
combination of constants from the same rows.

More intuition about the interaction comes from studying linear
combinations.
That will be our focus in the second chapter, after we finish the study
of Gauss's Method itself in the rest of this chapter.

\begin{exercises}
  \item Solve this system.
    Then solve the associated homogeneous system.
    \begin{equation*}
      \begin{linsys}{4}
        x  &+ &y  &- &2z &= &0 \\
        x  &- &y  &  &   &= &-3 \\
        3x &- &y  &- &2z &= &-6  \\
           &  &2y &- &2z &= &3  
      \end{linsys}
   \end{equation*}
   \begin{answer}
     This reduction solves the system.
     \begin{align*}
       \begin{amat}{3}
         1 &1   &-2 &0 \\
         1 &-1  &0  &3  \\
         3 &-1  &-2 &-6 \\
         0 &2   &-2 &3
       \end{amat}
       &\grstep[-3\rho_1+\rho_3]{-\rho_1+\rho_2}
       \begin{amat}{3}
         1 &1   &-2 &0 \\
         0 &-2  &2  &-3  \\
         0 &-4  &4  &-6 \\
         0 &2   &-2 &3
       \end{amat}                                 \\
       &\grstep[\rho_2+\rho_4]{-2\rho_2+\rho_3}
       \begin{amat}{3}
         1 &1   &-2 &0 \\
         0 &-2  &2  &-3  \\
         0 &0   &0  &0 \\
         0 &0   &0  &0
       \end{amat}
     \end{align*}
     The solution set is this.
     \begin{equation*}
       \set{\colvec{-3/2 \\ 3/2 \\ 0}
            +\colvec{1 \\ 1 \\ 1}z
            \suchthat z\in\Re}
     \end{equation*}
     Similarly we can reduce the associated homogeneous system
     \begin{align*}
       \begin{amat}{3}
         1 &1   &-2 &0 \\
         1 &-1  &0  &0  \\
         3 &-1  &-2 &0 \\
         0 &2   &-2 &0
       \end{amat}
       &\grstep[-3\rho_1+\rho_3]{-\rho_1+\rho_2}
       \begin{amat}{3}
         1 &1   &-2 &0 \\
         0 &-2  &2  &0  \\
         0 &-4  &4  &0 \\
         0 &2   &-2 &0
       \end{amat}                                         \\
       &\grstep[\rho_2+\rho_4]{-2\rho_2+\rho_3}
       \begin{amat}{3}
         1 &1   &-2 &0 \\
         0 &-2  &2  &0  \\
         0 &0   &0  &0 \\
         0 &0   &0  &0
       \end{amat}
     \end{align*}
     to get its solution set.
     \begin{equation*}
       \set{
            \colvec{1 \\ 1 \\ 1}z
            \suchthat z\in\Re}
     \end{equation*}
   \end{answer}
  \recommended \item 
    Solve each system.
    Express the solution set using vectors.
    Identify a particular solution and the solution set of the
    homogeneous system.
    % (These systems also appear in 
    % Exercise~2.\nearbyexercise{exer:SolveInMatrixNotation}.)
    \begin{exparts*}
      \partsitem \( \begin{linsys}[t]{2}
                  3x  &+  &6y  &=  &18  \\
                   x  &+  &2y  &=  &6   
                   \end{linsys}  \)
      \partsitem \( \begin{linsys}[t]{2}
                   x  &+  &y   &=  &1  \\
                   x  &-  &y   &=  &-1   
                    \end{linsys}  \)
      \partsitem \( \begin{linsys}[t]{3}
                   x_1  &   &     &+  &x_3   &=  &4  \\
                   x_1  &-  &x_2  &+  &2x_3  &=  &5  \\
                  4x_1  &-  &x_2  &+  &5x_3  &=  &17  
                    \end{linsys}  \)
      \partsitem \( \begin{linsys}[t]{3}
                   2a   &+  &b    &-  &c     &=  &2  \\
                   2a   &   &     &+  &c     &=  &3  \\
                    a   &-  &b    &   &      &=  &0   
                    \end{linsys}  \)
      \partsitem \( \begin{linsys}[t]{4}
                     x  &+  &2y   &-   &z   &    &    &=  &3  \\
                    2x  &+  &y    &    &    &+   &w   &=  &4  \\
                     x  &-  &y    &+   &z   &+   &w   &=  &1  
                    \end{linsys}  \)
      \partsitem \( \begin{linsys}[t]{4}
                     x  &   &     &+   &z   &+   &w   &=  &4  \\
                    2x  &+  &y    &    &    &-   &w   &=  &2  \\
                    3x  &+  &y    &+   &z   &    &    &=  &7  
                    \end{linsys}  \)
    \end{exparts*}
    \begin{answer}
      For the arithmetic to these, see the answers from the prior
      subsection.

      \begin{exparts}
        \partsitem
          This is the solution set.
          \begin{equation*}
            S=\set{\colvec[r]{6 \\ 0}+\colvec[r]{-2 \\ 1}y
              \suchthat y\in\Re}
          \end{equation*}
          Here are the particular solution and the solution set 
          for the associated
          homogeneous system.
          \begin{equation*}
            \colvec[r]{6 \\ 0}
              \quad\text{and}\quad
            \set{\colvec[r]{-2 \\ 1}y
              \suchthat y\in\Re}
          \end{equation*}
          \textit{Note.}
          There are two possible points of confusion here.
          First, the set $S$ given above is equal to this set 
          \begin{equation*}
            T=\set{\colvec[r]{4 \\ 1}+\colvec[r]{-2 \\ 1}y
              \suchthat y\in\Re}
          \end{equation*}
          because the two sets contain the same members.
          All of these are correct answers to, 
          ``What is a particular solution?''
          \begin{equation*}
            \colvec{6 \\ 0},\quad
            \colvec{4 \\ 1},\quad
            \colvec{2 \\ 2},\quad
            \colvec{1 \\ 2.5}
          \end{equation*}
          The second point of confusion is that the letter we use 
          in the set doesn't matter.
          This set also equals $S$.
          \begin{equation*}
            U=\set{\colvec[r]{6 \\ 0}+\colvec[r]{-2 \\ 1}u
                   \suchthat u\in\Re}
          \end{equation*}
        \partsitem
          This is the solution set.
          \begin{equation*}
            \set{\colvec[r]{0 \\ 1} }
          \end{equation*}
          These are a particular solution, 
          and the solution set for the associated
          homogeneous system.
          \begin{equation*}
            \colvec[r]{0 \\ 1}
              \qquad
            \set{\colvec[r]{0 \\ 0} }
          \end{equation*}
        \partsitem
          The solution set is infinite.
          \begin{equation*}
            \set{\colvec[r]{4 \\ -1 \\ 0}+\colvec[r]{-1 \\ 1 \\ 1}x_3
              \suchthat x_3\in\Re}
          \end{equation*}
          This is a particular solution and the solution set for the associated
          homogeneous system.
          \begin{equation*}
            \colvec[r]{4 \\ -1 \\ 0}
              \qquad
            \set{\colvec[r]{-1 \\ 1 \\ 1}x_3
              \suchthat x_3\in\Re}
          \end{equation*}
        \partsitem
          The solution set is a singleton.
          \begin{equation*}
            \set{\colvec[r]{1 \\ 1 \\ 1}}
          \end{equation*}
          A particular solution and the solution set for the associated
          homogeneous system are here.
          \begin{equation*}
            \colvec[r]{1 \\ 1 \\ 1}
              \qquad
            \set{\colvec[r]{0 \\ 0 \\ 0}}
          \end{equation*}
        \partsitem
          The solution set is infinite.
          \begin{equation*}
            \set{\colvec[r]{5/3 \\ 2/3 \\ 0 \\ 0}
                 +\colvec[r]{-1/3 \\ 2/3 \\ 1 \\ 0}z
                 +\colvec[r]{-2/3 \\ 1/3 \\ 0 \\ 1}w
                 \suchthat z,w\in\Re}
          \end{equation*}
          A particular solution and the solution set for the associated
          homogeneous system are here.
          \begin{equation*}
            \colvec[r]{5/3 \\ 2/3 \\ 0 \\ 0}
              \qquad
            \set{\colvec[r]{-1/3 \\ 2/3 \\ 1 \\ 0}z
                 +\colvec[r]{-2/3 \\ 1/3 \\ 0 \\ 1}w
                 \suchthat z,w\in\Re}
          \end{equation*}
        \partsitem This system's solution set is empty.
          Thus, there is no particular solution.
          The solution set of the associated homogeneous system is this.
          \begin{equation*}
            \set{\colvec[r]{-1 \\ 2 \\ 1 \\ 0}z
                 +\colvec[r]{-1 \\ 3 \\ 0 \\ 1}w
                 \suchthat z,w\in\Re}
          \end{equation*}
      \end{exparts}  
    \end{answer}
  \item 
    Solve each system, giving
    the solution set in vector notation.
    Identify a particular solution and the solution of the
    homogeneous system.
    \begin{exparts*}
      \partsitem \( \begin{linsys}[t]{3}
                  2x  &+  &y  &-  &z  &=  &1  \\
                  4x  &-  &y  &   &   &=  &3  
                  \end{linsys}  \)
      \partsitem \( \begin{linsys}[t]{4}
                   x  &   &   &-  &z  &   &   &=  &1  \\
                      &   &y  &+  &2z &-  &w  &=  &3  \\
                   x  &+  &2y &+  &3z &-  &w  &=  &7  
                   \end{linsys}  \)
      \partsitem \( \begin{linsys}[t]{4}
                   x  &-  &y  &+  &z  &   &   &=  &0  \\
                      &   &y  &   &   &+  &w  &=  &0  \\
                  3x  &-  &2y &+  &3z &+  &w  &=  &0  \\
                      &   &-y &   &   &-  &w  &=  &0  
                  \end{linsys}  \)
      \partsitem \( \begin{linsys}[t]{5}
                   a  &+  &2b &+  &3c &+  &d  &-  &e  &=  &1  \\
                  3a  &-  &b  &+  &c  &+  &d  &+  &e  &=  &3  
                  \end{linsys}  \)
    \end{exparts*}
    \begin{answer}
    The answers from the prior subsection show the row operations.
    Each answer here just lists the solution set, the particular solution,
    and the homogeneous solution.
    \begin{exparts}
      \partsitem
        The solution set is this.
        \begin{equation*}
          \set{\colvec[r]{2/3 \\ -1/3 \\ 0}
               +\colvec[r]{1/6 \\ 2/3 \\ 1}z
              \suchthat z\in\Re}
        \end{equation*}
        A particular solution and the solution set for the associated
        homogeneous system are here.
        \begin{equation*}
          \colvec[r]{2/3 \\ -1/3 \\ 0}
            \qquad
        \set{\colvec[r]{1/6 \\ 2/3 \\ 1}z
            \suchthat z\in\Re}
        \end{equation*}
      \partsitem
        The solution set is infinite.
        \begin{equation*}
          \set{\colvec[r]{1 \\ 3 \\ 0 \\ 0}
               +\colvec[r]{1 \\-2 \\ 1 \\ 0}z
              \suchthat z\in\Re}
        \end{equation*}
        Here are 
        a particular solution and the solution set for the associated
        homogeneous system.
        \begin{equation*}
          \colvec[r]{1 \\ 3 \\ 0 \\ 0}
            \qquad
          \set{\colvec[r]{1 \\ -2 \\ 1 \\ 0}z
              \suchthat z\in\Re}
        \end{equation*}
      \partsitem
        This is the solution set.
        \begin{equation*}
          \set{\colvec[r]{0 \\ 0 \\ 0 \\ 0}
               +\colvec[r]{-1 \\ 0 \\ 1 \\ 0}z
               +\colvec[r]{-1 \\ -1 \\ 0 \\ 1}w
              \suchthat z,w\in\Re}
        \end{equation*}
        Here is
        a particular solution and the solution set for the associated
        homogeneous system.
        \begin{equation*}
          \colvec[r]{0 \\ 0 \\ 0 \\ 0}
            \qquad
          \set{\colvec[r]{-1 \\ 0 \\ 1 \\ 0}z
               +\colvec[r]{-1 \\ -1 \\ 0 \\ 1}w
              \suchthat z,w\in\Re}
        \end{equation*}
      \partsitem
        The solution set is this.
        \begin{equation*}
          \set{\colvec[r]{1 \\ 0 \\ 0 \\ 0 \\ 0}
               +\colvec[r]{-5/7 \\ -8/7 \\ 1 \\ 0 \\ 0}c
               +\colvec[r]{-3/7 \\ -2/7 \\ 0 \\ 1 \\ 0}d
               +\colvec[r]{-1/7 \\ 4/7 \\ 0 \\ 0 \\ 1}e
              \suchthat c,d,e\in\Re}
        \end{equation*}
        And, this is
        a particular solution and the solution set for the associated
        homogeneous system.
        \begin{equation*}
          \colvec[r]{1 \\ 0 \\ 0 \\ 0 \\ 0}
            \qquad
          \set{\colvec[r]{-5/7 \\ -8/7 \\ 1 \\ 0 \\ 0}c
               +\colvec[r]{-3/7 \\ -2/7 \\ 0 \\ 1 \\ 0}d
               +\colvec[r]{-1/7 \\ 4/7 \\ 0 \\ 0 \\ 1}e
              \suchthat c,d,e\in\Re}
        \end{equation*}
    \end{exparts}  
   \end{answer}
  \recommended \item 
    For the system
    \begin{equation*}
      \begin{linsys}{4}
       2x  &-  &y  &   &    &-  &w  &=  &3  \\
           &   &y  &+  &z   &+  &2w &=  &2  \\
        x  &-  &2y &-  &z   &   &   &=  &-1
      \end{linsys}
    \end{equation*}
    which of these can be used as the particular solution part of some
    general solution?
    \begin{exparts*}
      \partsitem   \( \colvec[r]{0 \\ -3 \\ 5 \\ 0} \)
      \partsitem   \( \colvec[r]{2 \\ 1 \\ 1 \\ 0} \)
      \partsitem   \( \colvec[r]{-1 \\ -4 \\ 8 \\ -1} \)
    \end{exparts*}
    \begin{answer}
      Just plug them in and see if they satisfy all three equations.
      \begin{exparts}
        \partsitem No.
        \partsitem Yes.
        \partsitem Yes.
      \end{exparts}  
    \end{answer}
  \recommended \item  
    \nearbylemma{th:GenEqPartHomo} says that we can use any particular solution 
    for $\vec{p}$.
    Find, if possible, a general solution to this system
    \begin{equation*}
      \begin{linsys}{4}
        x  &-  &y  &   &    &+  &w  &=  &4  \\
       2x  &+  &3y &-  &z   &   &   &=  &0  \\
           &   &y  &+  &z   &+  &w  &=  &4  
      \end{linsys}
    \end{equation*}
    that uses the given vector as its particular solution.
    \begin{exparts*}
      \partsitem   \( \colvec[r]{0 \\ 0 \\ 0 \\ 4} \)
      \partsitem   \( \colvec[r]{-5 \\ 1 \\ -7 \\ 10} \)
      \partsitem   \( \colvec[r]{2 \\ -1 \\ 1 \\ 1} \)
    \end{exparts*}
    \begin{answer}
      Gauss's Method on the associated homogeneous system 
      \begin{align*}
        \begin{amat}[r]{4}
           1  &-1  &0  &1  &0  \\
           2  &3   &-1 &0  &0  \\
           0  &1   &1  &1  &0
        \end{amat}
        &\grstep{-2\rho_1+\rho_2}
        \begin{amat}[r]{4}
           1  &-1  &0  &1  &0  \\
           0  &5   &-1 &-2 &0  \\
           0  &1   &1  &1  &0
        \end{amat}                                \\
        &\grstep{-(1/5)\rho_2+\rho_3}
        \begin{amat}[r]{4}
           1  &-1  &0  &1  &0  \\
           0  &5   &-1 &-2 &0  \\
           0  &0   &6/5&7/5&0
        \end{amat}
      \end{align*}
      gives this is the solution to the homogeneous problem.
      \begin{equation*}
        \set{\colvec[r]{-5/6 \\ 1/6 \\ -7/6 \\ 1}w\suchthat w\in\Re}
      \end{equation*}
      \begin{exparts}
        \partsitem That vector is indeed a particular solution, so the required
          general solution is this.
          \begin{equation*}
            \set{\colvec[r]{0 \\ 0 \\ 0 \\ 4}+
                 \colvec[r]{-5/6 \\ 1/6 \\ -7/6 \\ 1}w\suchthat w\in\Re}
          \end{equation*}
        \partsitem That vector is a particular solution so the required
          general solution is this.
          \begin{equation*}
            \set{\colvec[r]{-5 \\ 1 \\ -7 \\ 10}+
                 \colvec[r]{-5/6 \\ 1/6 \\ -7/6 \\ 1}w\suchthat w\in\Re}
          \end{equation*}
        \partsitem That vector is not a solution of the system since
          it does not satisfy the third equation.
          No such general solution exists.
      \end{exparts} 
    \end{answer}
  \item 
     One is nonsingular while the other is singular.
     Which is which?
     \begin{exparts*}
       \partsitem $\begin{mat}[r]
           1  &3   \\
           4  &-12    
         \end{mat}$
       \partsitem $\begin{mat}[r]
           1  &3  \\
           4  &12  
         \end{mat}$
     \end{exparts*}
     \begin{answer}
       The first is nonsingular while the second is singular.
       Just do Gauss's Method and see if the echelon form result has
       non-$0$ numbers in each entry on the diagonal.
     \end{answer}
  \recommended \item 
    Singular or nonsingular?
    \begin{exparts*}
      \partsitem \(
        \begin{mat}[r]
          1  &2  \\
          1  &3
        \end{mat}   \)
      \partsitem \(
        \begin{mat}[r]
          1  &2  \\
         -3  &-6
        \end{mat}   \)
      \partsitem \(
        \begin{mat}[r]
          1  &2  &1  \\
          1  &3  &1
        \end{mat}   \)
      \partsitem \(
        \begin{mat}[r]
          1  &2  &1  \\
          1  &1  &3  \\
          3  &4  &7
        \end{mat}   \)
      \partsitem \(
        \begin{mat}[r]
          2  &2  &1  \\
          1  &0  &5  \\
         -1  &1  &4
        \end{mat}   \)
    \end{exparts*}
    \begin{answer}
      \begin{exparts}
      \partsitem Nonsingular:
        \begin{equation*}
          \grstep{-\rho_1+\rho_2}
          \begin{mat}[r]
            1  &2  \\
            0  &1
          \end{mat}
        \end{equation*}
        ends with each row containing a leading entry.
      \partsitem Singular:
        \begin{equation*}
          \grstep{3\rho_1+\rho_2}
          \begin{mat}[r]
            1  &2  \\
            0  &0
          \end{mat}
        \end{equation*}
        ends with row \( 2 \) without a leading entry.
      \partsitem Neither.
        A matrix must be square for either word to apply.
      \partsitem Singular.
      \partsitem Nonsingular.
     \end{exparts}  
    \end{answer}
  \recommended \item 
    Is the given vector in the set generated by the
    given set?
      \begin{exparts}
        \partsitem \( \colvec[r]{2 \\ 3}, \)\
          \( \set{\colvec[r]{1 \\ 4},
                \colvec[r]{1 \\ 5}} \)
        \partsitem \( \colvec[r]{-1 \\ 0 \\ 1}, \)\
          \( \set{\colvec[r]{2 \\ 1 \\ 0},
                \colvec[r]{1 \\ 0 \\ 1}} \)
        \partsitem \( \colvec[r]{1 \\ 3 \\ 0}, \)\
          \( \set{\colvec[r]{1 \\ 0 \\ 4},
                \colvec[r]{2 \\ 1 \\ 5},
                \colvec[r]{3 \\ 3 \\ 0},
                \colvec[r]{4 \\ 2 \\ 1}} \)
        \partsitem \( \colvec[r]{1 \\ 0 \\ 1 \\ 1}, \)\
          \( \set{\colvec[r]{2 \\ 1 \\ 0 \\ 1},
                \colvec[r]{3 \\ 0 \\ 0 \\ 2}} \)
      \end{exparts}
      \begin{answer}
        In each case we must decide if the vector is a linear combination
        of the vectors in the set.
        \begin{exparts}
          \partsitem Yes.
            Solve
            \begin{equation*}
              c_1\colvec[r]{1 \\ 4}+c_2\colvec[r]{1 \\ 5}=\colvec[r]{2 \\ 3}
            \end{equation*}
            with
            \begin{equation*}
              \begin{amat}[r]{2}
                1  &1  &2  \\
                4  &5  &3
              \end{amat}
              \grstep{-4\rho_1+\rho_2}
              \begin{amat}[r]{2}
                1  &1  &2  \\
                0  &1  &-5
              \end{amat}
            \end{equation*}
            to conclude that there are $c_1$ and $c_2$ giving the combination. 
          \partsitem No.
            The reduction
            \begin{equation*}
              \begin{amat}[r]{2}
                2  &1  &-1 \\
                1  &0  &0  \\
                0  &1  &1
              \end{amat}
              \grstep{-(1/2)\rho_1+\rho_2}
              \begin{amat}[r]{2}
                2  &1     &-1 \\
                0  &-1/2  &1/2  \\
                0  &1     &1
              \end{amat}
              \grstep{2\rho_2+\rho_3}
              \begin{amat}[r]{2}
                2  &1     &-1 \\
                0  &-1/2  &1/2  \\
                0  &0     &2
              \end{amat}
            \end{equation*}
            shows that
            \begin{equation*}
              c_1\colvec[r]{2 \\ 1 \\ 0}+c_2\colvec[r]{1 \\ 0 \\ 1}
                =\colvec[r]{-1 \\ 0 \\ 1}
            \end{equation*}
            has no solution.
          \partsitem Yes.
            The reduction
            \begin{align*}
              \begin{amat}[r]{4}
                1  &2  &3  &4  &1  \\
                0  &1  &3  &2  &3  \\
                4  &5  &0  &1  &0
              \end{amat}
              &\grstep{-4\rho_1+\rho_3}
              \begin{amat}[r]{4}
                1  &2  &3  &4  &1  \\
                0  &1  &3  &2  &3  \\
                0  &-3 &-12&-15&-4
              \end{amat}                   \\
              &\grstep{3\rho_2+\rho_3}
              \begin{amat}[r]{4}
                1  &2  &3  &4  &1  \\
                0  &1  &3  &2  &3  \\
                0  &0  &-3 &-9 &5
              \end{amat}
            \end{align*}
            shows that there are infinitely many ways
            \begin{equation*}
              \set{\colvec[r]{c_1 \\ c_2 \\ c_3 \\ c_4}=
                   \colvec[r]{-10 \\ 8 \\ -5/3 \\ 0}+
                   \colvec[r]{-9 \\ 7 \\ -3 \\ 1}c_4
                    \suchthat c_4\in\Re}
            \end{equation*}
            to write a combination.
            \begin{equation*}
              \colvec[r]{1 \\ 3 \\ 0}=
              c_1\colvec[r]{1 \\ 0 \\ 4}+
              c_2\colvec[r]{2 \\ 1 \\ 5}+
              c_3\colvec[r]{3 \\ 3 \\ 0}+
              c_4\colvec[r]{4 \\ 2 \\ 1}
            \end{equation*}
          \partsitem No.
            Look at the third components.
        \end{exparts} 
      \end{answer}
  \item 
     Prove that any linear system with a nonsingular matrix of 
     coefficients has a solution, and that the solution is unique.
    \begin{answer}
        Because the matrix of coefficients is nonsingular, Gauss's Method
        ends with an echelon form where each variable leads an equation.
        Back substitution gives a unique solution.

      (Another way to see that the solution is unique is to note that
      with a nonsingular matrix of coefficients the associated
      homogeneous system has a unique solution, by definition.
      Since the general solution is the sum of a particular solution with
      each homogeneous solution, the general solution has 
      at most one element.)
     \end{answer}
  \item 
    In the
    proof of
    \nearbylemma{le:HomoSltnSpanVecs},
    what happens if there are no non-\( 0=0 \) equations?
    \begin{answer}
      In this case the solution set is all of \( \Re^n \) and we can 
      express it in the required form.
      \begin{equation*}
        \set{c_1\colvec[r]{1 \\ 0 \\ \vdotswithin{1} \\ 0}
             +c_2\colvec[r]{0 \\ 1 \\ \vdotswithin{0} \\ 0}
             +\cdots
             +c_n\colvec[r]{0 \\ 0 \\ \vdotswithin{0} \\ 1}
             \suchthat c_1,\ldots,c_n\in\Re}
      \end{equation*}  
     \end{answer}
  \recommended \item 
    Prove that if \( \vec{s} \) and \( \vec{t} \)
    satisfy a homogeneous system then so do these vectors.
    \begin{exparts*}
      \partsitem \( \vec{s}+\vec{t} \)
      \partsitem \( 3\vec{s} \)
      \partsitem \( k\vec{s}+m\vec{t} \) for \( k,m\in\Re \)
    \end{exparts*}
    What's wrong with this argument: ``These three show that if a homogeneous
    system has one solution then it has many solutions\Dash any multiple of 
    a solution is another solution, and any sum of solutions is a solution
    also\Dash so there are no
    homogeneous systems with exactly one solution.''?
    \begin{answer}
      Assume \( \vec{s},\vec{t}\in\Re^n \) and write them as here.
      \begin{equation*}
        \vec{s}=\colvec{s_1 \\ \vdotswithin{s_1} \\ s_n}
          \qquad
        \vec{t}=\colvec{t_1 \\ \vdotswithin{t_1} \\ t_n}
      \end{equation*}
      Also let \( a_{i,1}x_1+\cdots+a_{i,n}x_n=0 \) be the \( i \)-th equation
      in the homogeneous system.
      \begin{exparts}
        \partsitem The check is easy.
          \begin{multline*}
            a_{i,1}(s_1+t_1)+\cdots+a_{i,n}(s_n+t_n)      \\
            =
            (a_{i,1}s_1+\cdots+a_{i,n}s_n)
            +(a_{i,1}t_1+\cdots+a_{i,n}t_n)           
            =
            0+0
          \end{multline*}
        \partsitem This is similar to the prior one.
          \begin{equation*}
            a_{i,1}(3s_1)+\cdots+a_{i,n}(3s_n)
            =3(a_{i,1}s_1+\cdots+a_{i,n}s_n)
            =3\cdot 0=0
          \end{equation*}
        \partsitem This one is not much harder.
          \begin{multline*}
            a_{i,1}(ks_1+mt_1)+\cdots+a_{i,n}(ks_n+mt_n)  \\
            =
            k(a_{i,1}s_1+\cdots+a_{i,n}s_n)
            +m(a_{i,1}t_1+\cdots+a_{i,n}t_n)         
            =
            k\cdot 0+m\cdot 0
          \end{multline*}
      \end{exparts}  
     What is wrong with that argument is that any linear combination 
     involving only the 
     zero vector yields the zero vector.
   \end{answer}
  \item
    Prove that if a system with only rational coefficients
    and constants
    has a solution then it has at least one all-rational solution.
    Must it have infinitely many?
    \begin{answer}
      First the proof.

      If the system has all rational coefficients then
      Gauss's Method with back substitution will use only rationals (e.g.,
      \( -(m/n)\rho_i+\rho_j \)).
      Thus we can express the solution set using only rational numbers as
      the components of each vector.
      Therefore, the particular solution is all rational.

      As to the question about infinitely many, 
      there are infinitely many rational vector solutions if and only if the
      associated homogeneous system has infinitely many 
      real vector solutions.
      That's because setting any parameters to be rationals will produce an
      all-rational solution.  
   \end{answer}
\end{exercises}


% END NATIVE SOURCE src/gr/gr1.tex

\setcounter{section}{2}

% BEGIN NATIVE SOURCE src/gr/gr3.tex
% Chapter 1, Section 3 _Linear Algebra_ Jim Hefferon
%  http://joshua.smcvt.edu/linearalgebra
%  2001-Jun-09
\section{Reduced Echelon Form}
After developing the mechanics of Gauss's Method, 
we observed that it can be done in more than one way.
For example, from this matrix 
\begin{equation*}
    \begin{mat}[r]
       2  &2  \\
       4  &3
    \end{mat}
\end{equation*}
we could derive any of these three echelon form matrices.
\begin{equation*}
    \begin{mat}[r]
       2  &2  \\
       0  &-1
    \end{mat}
    \qquad
    \begin{mat}[r]
       1  &1  \\
       0  &-1
    \end{mat}
    \qquad
    \begin{mat}[r]
       2  &0  \\
       0  &-1
    \end{mat}
\end{equation*}
The first results from $-2\rho_1+\rho_2$.
The second comes from doing $(1/2)\rho_1$ and then $-4\rho_1+\rho_2$.
The third comes
from $-2\rho_1+\rho_2$ followed by $2\rho_2+\rho_1$
(after the first row combination the matrix is already in
echelon form 
but it is nonetheless a legal row operation).

In this chapter's first section we noted that
this raises questions.
Will any two echelon form versions of a linear system have the same number of
free variables?
If yes, 
will the two have exactly the same free variables?
In this section we will 
give a way to decide if one linear system 
can be derived from another by row operations.
The answers to both questions, both ``yes,''
will follow from that.








\subsection{Gauss-Jordan Reduction}%
% Gaussian elimination coupled with back-substitution
% solves linear systems but it is not the only method possible.
Here is an extension of Gauss's Method that has some advantages.

\begin{example} \label{exm:GJRedReadOffSol}
To solve
\begin{equation*}
  \begin{linsys}{3}
    x  &+  &y  &-  &2z  &=  &-2  \\
       &   &y  &+  &3z  &=  &7   \\
    x  &   &   &-  &z   &=  &-1  
  \end{linsys}
\end{equation*}
we can start as usual by reducing it to echelon form.
\begin{equation*}
  \grstep{-\rho_1+\rho_3}
    \begin{amat}[r]{3}
       1  &1  &-2 &-2  \\
       0  &1  &3  &7   \\
       0  &-1 &1  &1
    \end{amat}
  \grstep{\rho_2+\rho_3}
    \begin{amat}[r]{3}
       1  &1  &-2 &-2  \\
       0  &1  &3  &7   \\
       0  &0  &4  &8
    \end{amat}
\end{equation*}
We can keep going to a second stage
by making the leading entries into \( 1 \)'s
\begin{equation*}
    \grstep{(1/4)\rho_3}
    \begin{amat}[r]{3}
       1  &1  &-2 &-2  \\
       0  &1  &3  &7   \\
       0  &0  &1  &2
    \end{amat}
\end{equation*}
and then to a third stage that uses the leading entries 
to eliminate all of the other entries in each column 
by combining upwards.
\begin{equation*}
  \grstep[2\rho_3+\rho_1]{-3\rho_3+\rho_2}
    \begin{amat}[r]{3}
       1  &1  &0  &2   \\
       0  &1  &0  &1   \\
       0  &0  &1  &2
    \end{amat}
  \grstep{-\rho_2+\rho_1}
    \begin{amat}[r]{3}
       1  &0  &0  &1   \\
       0  &1  &0  &1   \\
       0  &0  &1  &2
    \end{amat}
\end{equation*}
The answer is \( x=1 \), \( y=1 \), and \( z=2 \).
\end{example}
%<*Pivoting>
Using one entry to clear out the rest of a column is
\definend{pivoting}\index{pivoting} on that entry.
%</Pivoting>

Notice that the row combination operations in the first stage move 
left to right 
while the combination operations in the third stage move right to left.

\begin{example}  \label{exm:GJRedReadOffSolTwo}
The middle stage operations that 
turn the leading entries into \( 1 \)'s
don't interact so we can combine multiple ones into a single step.
\begin{align*}
    \begin{amat}[r]{2}
       2   &1   &7   \\
       4   &-2  &6
    \end{amat}
  &\grstep{-2\rho_1+\rho_2}
  \begin{amat}[r]{2}
       2   &1   &7   \\
       0   &-4  &-8
    \end{amat}                                   \\
  &\grstep[(-1/4)\rho_2]{(1/2)\rho_1}
  \begin{amat}[r]{2}
       1   &1/2   &7/2   \\
       0   &1     &2
    \end{amat}                                    \\
  &\grstep{-(1/2)\rho_2+\rho_1}
  \begin{amat}[r]{2}
       1   &0   &5/2   \\
       0   &1   &2
    \end{amat}
\end{align*}
The answer is $x=5/2$ and $y=2$.
\end{example}

%<*GaussJordanReduction>
This extension of Gauss's Method is the 
\definend{Gauss-Jordan Method}\index{Gauss's Method!Gauss-Jordan Method} or 
\definend{Gauss-Jordan reduction}.\index{linear equation!solution of!Gauss-Jordan}\index{Gauss's Method!Gauss-Jordan}
%</GaussJordanReduction>
% It goes past echelon form to a more refined, more specialized,
% matrix form.

\begin{definition}\label{def:RedEchForm}
%<*df:RedEchForm>
A matrix or linear system is in
\definend{reduced echelon form\/}\index{echelon form!reduced}\index{reduced echelon form}
if, in addition to being in echelon form, each leading entry is a~$1$ 
and is the only nonzero entry in its column.
%</df:RedEchForm>
\end{definition}

\noindent
%<*CostRedEchForm>
The cost of using Gauss-Jordan reduction to solve a system 
is the additional arithmetic.
The benefit is that we can just read off the solution set
description.
%</CostRedEchForm>

In any echelon form system, reduced or not, we can read off 
when the system has an empty
solution set because there is a contradictory equation.
We can read off 
when the system has a one-element solution set because there is no
contradiction and every
variable is the leading variable in some row.
And, we can read off when the system has an infinite solution set because 
there is no contradiction and at least one variable is free.

However, in reduced echelon form we can read off not just the size of the  
solution set but also its description.
We have no trouble describing the solution set when it is empty, of course.
\nearbyexample{exm:GJRedReadOffSol} and~\ref{exm:GJRedReadOffSolTwo} 
show how in a single element solution set case the single element is
in the column of constants.
The next example shows how to read the parametrization
of an infinite solution set.

\begin{example}
\begin{multline*}
  \begin{amat}[r]{4}
     2  &6  &1  &2  &5  \\
     0  &3  &1  &4  &1  \\
     0  &3  &1  &2  &5
  \end{amat}
  \grstep{-\rho_2+\rho_3}
  \begin{amat}[r]{4}
     2  &6  &1  &2  &5  \\
     0  &3  &1  &4  &1  \\
     0  &0  &0  &-2 &4
  \end{amat}                                        \\
  \grstep[(1/3)\rho_2 \\ -(1/2)\rho_3]{(1/2)\rho_1}
  \repeatedgrstep[-\rho_3+\rho_1]{-(4/3)\rho_3+\rho_2}
  \repeatedgrstep{-3\rho_2+\rho_1}
  \begin{amat}[r]{4}
     1  &0  &-1/2  &0  &-9/2  \\
     0  &1  &1/3   &0  &3  \\
     0  &0  &0     &1  &-2
  \end{amat}
\end{multline*}
As a linear system this is 
\begin{equation*}
  \begin{linsys}{4}
     x_1  &&     &-&1/2x_3  &&   &= &-9/2  \\
          &&x_2  &+&1/3x_3  &&   &= &3  \\
          &&     &&        &{}\hspace{.5em}{}&x_4 &= &-2    
  \end{linsys}
\end{equation*}
so a solution set description is this.
\begin{equation*}  
  S=\set{\colvec{x_1 \\ x_2 \\ x_3 \\ x_4}
                        =\colvec[r]{-9/2 \\ 3 \\ 0 \\ -2}
                         +\colvec[r]{1/2 \\ -1/3 \\ 1 \\ 0}x_3
                        \suchthat x_3\in\Re}
\end{equation*}
\end{example}

Thus, echelon form isn't some kind of one best form for systems.
Other forms, such as reduced echelon form, have advantages and
disadvantages.
Instead of picturing linear systems (and the associated matrices) 
as things we operate on, 
always directed toward the goal of echelon form, we can think of 
them as interrelated, where
we can get from one to another by row operations.
The rest of this subsection develops this thought.

\begin{lemma} \label{le:RowOpsRev}
%<*lm:RowOpsRev>
Elementary row operations are reversible.
%</lm:RowOpsRev>
\end{lemma}

\begin{proof}
%<*pf:RowOpsRev>
For any matrix \( A \),
the effect of swapping rows is reversed by swapping them back,
multiplying a row by a nonzero \( k \) is undone by multiplying by
$1/k$,
and adding a multiple of row \( i \) to row \( j \) (with $i\neq j$)
is undone by subtracting the same multiple of row \( i \) from row \( j \).
\begin{equation*}
      A
     \grstep{\rho_i\leftrightarrow\rho_j}
     \repeatedgrstep{\rho_j\leftrightarrow\rho_i}
      A
  \qquad
        A
     \grstep{k\rho_i}
     \repeatedgrstep{(1/k)\rho_i}
      A
  \qquad
        A
     \grstep{k\rho_i+\rho_j}
     \repeatedgrstep{-k\rho_i+\rho_j}
      A                          
\end{equation*}
%</pf:RowOpsRev>
(We need the $i\neq j$ condition;
see \nearbyexercise{exer:INotJMakesRowOpsRev}.)
\end{proof}

Again, the point of view that we are developing, supported now by the lemma,
is that the term `reduces to' is misleading:~where
\( A\longrightarrow B \), we shouldn't think of \( B \) as
after~\( A \) or simpler than~$A$.
Instead we should think of the two matrices as interrelated.
Below is a picture.
It shows the matrices from the start of this section and their
reduced echelon form version in a cluster, as
interreducible. 
\begin{center}  
  \includegraphics{gr/mp/ch1.28}
\end{center}

%<*EquivMatrices>
We say 
that matrices that reduce to each other are equivalent with respect
to the relationship of row reducibility.
The next result justifies this, using the definition of 
an equivalence.\appendrefs{equivalence relations}
%</EquivMatrices>

\begin{lemma} \label{lm:ReducesToIsEqRel}
%<*lm:ReducesToIsEqRel>
Between matrices, `reduces to' is an equivalence re\-la\-tion.
%</lm:ReducesToIsEqRel>
\end{lemma}

\begin{proof}
%<*pf:ReducesToIsEqRel0>
We must check the conditions
(i)~reflexivity, that any matrix reduces to itself,
(ii)~symmetry, that if \( A \) reduces to \( B \) then
   \( B \) reduces to \( A \),
and (iii)~transitivity, that if \( A \) reduces to \( B \) and
      \( B \) reduces to \( C \) then \( A \) reduces to
      \( C \).
%</pf:ReducesToIsEqRel0>

%<*pf:ReducesToIsEqRel1>
Reflexivity is easy; any  matrix reduces to itself in zero-many operations.

The relationship is symmetric by the prior lemma\Dash if
\( A \) reduces to \( B \) by some row operations
then also \( B \) reduces to \( A \) by reversing those operations.
%</pf:ReducesToIsEqRel1>

%<*pf:ReducesToIsEqRel2>
For transitivity, suppose that \( A \) reduces to \( B \) and
that \( B \) reduces to \( C \).
Following the reduction steps from $A \rightarrow\cdots\rightarrow B$
with those from  $B \rightarrow\cdots\rightarrow C$ 
gives a reduction from \( A \) to \( C \).
%</pf:ReducesToIsEqRel2>
\end{proof}

\begin{definition} \label{df:RowEquivalence}
%<*df:RowEquivalence>
Two matrices that are interreducible by elementary row operations
are \definend{row equivalent}.\index{matrix!row equivalence}%
\index{row equivalence}\index{equivalence relation!row equivalence}
%</df:RowEquivalence>
\end{definition}

%<*RowEquivalanceClasses>
The diagram below shows the collection of all matrices as a box.
Inside that box each matrix lies in a class.
Matrices are in the same class if and only if they are interreducible.
The classes are disjoint\Dash no matrix is in two distinct classes.
We have partitioned the collection of matrices into 
\definend{row equivalence classes}.\appendrefs{partitions and class representatives}\index{partition!row equivalence classes}
%</RowEquivalanceClasses>
\begin{center}
  \includegraphics{gr/mp/ch1.27}
\end{center}
\noindent One of the classes is the
cluster of interrelated 
matrices from the start of this section sketched above
(it includes all of the nonsingular $\nbyn{2}$ matrices). 

The next subsection proves that the reduced echelon form of a matrix is 
unique.
Rephrased in terms of the row-equivalence relationship, 
we shall prove that every matrix is 
row equivalent to one and only one reduced echelon form matrix.
In terms of the partition what we shall prove is:~every
equivalence class contains one and only one reduced echelon form matrix.
So each reduced echelon form matrix serves as a representative of its 
class.

\begin{exercises}
   \recommended \item 
     Use Gauss-Jordan reduction to solve each system.
     \begin{exparts*}
        \partsitem \(
          \begin{linsys}[t]{2}
               x  &+  &y  &=  &2  \\
               x  &-  &y  &=  &0  
          \end{linsys}   \)
        \partsitem \(
          \begin{linsys}[t]{3}
               x  &   &   &-  &z  &=  &4  \\
              2x  &+  &2y &   &   &=  &1  
          \end{linsys}  \)
        \partsitem  \(
           \begin{linsys}[t]{2}
               3x  &-  &2y  &=  &1  \\
               6x  &+  &y   &=  &1/2 
           \end{linsys}  \)
        \partsitem \(
           \begin{linsys}[t]{3}
              2x  &-  &y  &  &  &= &-1  \\
               x  &+  &3y &- &z &= &5   \\
                  &   &y  &+ &2z&= &5   
           \end{linsys} \)
     \end{exparts*}
     \begin{answer}
       These answers show only the Gauss-Jordan reduction.
       With it, describing the solution set is easy. 
       \begin{exparts}
         \partsitem The solution set contains only a single element.
           \begin{multline*}
             \begin{amat}[r]{2}
               1  &1  &2  \\
               1  &-1 &0
             \end{amat}
             \grstep{-\rho_1+\rho_2}
             \begin{amat}[r]{2}
               1  &1  &2  \\
               0  &-2 &-2
             \end{amat}                          \\                       
             \grstep{-(1/2)\rho_2}
             \begin{amat}[r]{2}
               1  &1  &2  \\
               0  &1  &1
             \end{amat}                         
             \grstep{-\rho_2+\rho_1}
             \begin{amat}[r]{2}
               1  &0  &1  \\
               0  &1  &1
             \end{amat}
           \end{multline*}
         \partsitem The solution set has one parameter.
            \begin{equation*}
             \begin{amat}[r]{3}
               1  &0  &-1  &4  \\
               2  &2  &0   &1
             \end{amat}
             \grstep{-2\rho_1+\rho_2}
             \begin{amat}[r]{3}
               1  &0  &-1  &4  \\
               0  &2  &2   &-7
             \end{amat}                    
             \grstep{(1/2)\rho_2}
             \begin{amat}[r]{3}
               1  &0  &-1  &4  \\
               0  &1  &1   &-7/2
             \end{amat}
           \end{equation*}
         \partsitem There is a unique solution.
           \begin{multline*}
             \begin{amat}[r]{2}
               3  &-2  &1  \\
               6  &1   &1/2
             \end{amat}
             \grstep{-2\rho_1+\rho_2}
             \begin{amat}[r]{2}
               3  &-2  &1  \\
               0  &5   &-3/2
             \end{amat}                            \\                       
             \grstep[(1/5)\rho_2]{(1/3)\rho_1}
             \begin{amat}[r]{2}
               1  &-2/3&1/3 \\
               0  &1   &-3/10
             \end{amat}                       
             \grstep{(2/3)\rho_2+\rho_1}
             \begin{amat}[r]{2}
               1  &0   &2/15 \\
               0  &1   &-3/10
             \end{amat}
          \end{multline*}
        \partsitem A row swap in the second step makes the arithmetic easier.
         \begin{multline*}
          \begin{amat}[r]{3}
            2  &-1  &0  &-1  \\
            1  &3   &-1 &5   \\
            0  &1   &2  &5
          \end{amat}
          \grstep{-(1/2)\rho_1+\rho_2}
          \begin{amat}[r]{3}
            2  &-1  &0  &-1   \\
            0  &7/2 &-1 &11/2 \\
            0  &1   &2  &5
          \end{amat}                                \\
          \begin{aligned}
          &\grstep{\rho_2\leftrightarrow\rho_3}
          \begin{amat}[r]{3}
            2  &-1  &0  &-1   \\
            0  &1   &2  &5    \\
            0  &7/2 &-1 &11/2
          \end{amat}                      
            \grstep{-(7/2)\rho_2+\rho_3}
            \begin{amat}[r]{3}
              2  &-1  &0  &-1   \\
              0  &1   &2  &5    \\
              0  &0   &-8 &-12
            \end{amat}                               \\     
            &\grstep[-(1/8)\rho_2]{(1/2)\rho_1}
            \begin{amat}[r]{3}
              1  &-1/2&0  &-1/2 \\
              0  &1   &2  &5    \\
              0  &0   &1  &3/2
            \end{amat}                                
            \grstep{-2\rho_3+\rho_2}
            \begin{amat}[r]{3}
              1  &-1/2&0  &-1/2 \\
              0  &1   &0  &2    \\
              0  &0   &1  &3/2
            \end{amat}                                    \\             
            &\grstep{(1/2)\rho_2+\rho_1}
            \begin{amat}[r]{3}
              1  &0   &0  &1/2  \\
              0  &1   &0  &2    \\
              0  &0   &1  &3/2
            \end{amat}
          \end{aligned}
        \end{multline*}
       \end{exparts}  
     \end{answer}
  \item Do Gauss-Jordan reduction.
    \begin{exparts*}
      \partsitem
        $\begin{linsys}[t]{3}
          x  &+ &y &- &z &= &3 \\
          2x &- &y  &- &z &= &1 \\
          3x &+  &y  &+ &2z &= &0
        \end{linsys}$
      \partsitem
        $\begin{linsys}[t]{3}
          x  &+ &y  &+ &2z  &= &0 \\
          2x  &- &y  &+  &z &= &1 \\
          4x &+ &y  &+ &5z &= &1  
        \end{linsys}$
    \end{exparts*}
    \begin{answer}
      \begin{exparts}
        \partsitem
          \begin{multline*}
            \begin{amat}{3}
              1 &1  &-1 &3 \\
              2 &-1 &-1 &1 \\
              3 &1  &2  &0
            \end{amat}
            \grstep[-3\rho_1+\rho_3]{-2\rho_1+\rho_2}
            \begin{amat}{3}
              1 &1  &-1 &3 \\
              0 &-3 &1  &-5 \\
              0 &-2  &5  &-9
            \end{amat}                             \\
          \begin{aligned}   
            &\grstep{-(2/3)\rho_2+\rho_3}
            \begin{amat}{3}
              1 &1  &-1    &3 \\
              0 &-3 &1     &-5 \\
              0 &0  &13/3  &-17/3
            \end{amat}                                 
            \grstep[(3/13)\rho_3]{-(1/3)\rho_2}
            \begin{amat}{3}
              1 &1  &-1    &3 \\
              0 &1  &-1/3    &5/3 \\
              0 &0  &1      &-17/13
            \end{amat}                                \\
            &\grstep[(1/3)\rho_3+\rho_2]{\rho_3+\rho_1}
            \begin{amat}{3}
              1 &1  &0    &22/13 \\
              0 &1  &0    &16/13 \\
              0 &0  &1      &-17/13
            \end{amat}
            \grstep{-\rho_2+\rho_1}
            \begin{amat}{3}
              1 &0  &0    &6/13 \\
              0 &1  &0    &16/13 \\
              0 &0  &1      &-17/13
            \end{amat}
          \end{aligned}
          \end{multline*}
        \partsitem
          \begin{multline*}
            \begin{amat}{3}
              1 &1  &2 &0 \\
              2 &-1 &1 &1 \\
              4 &1  &5 &1
            \end{amat}
            \grstep[-4\rho_1+\rho_3]{-2\rho_1+\rho_2}
            \begin{amat}{3}
              1 &1  &2  &0 \\
              0 &-3 &-3 &1 \\
              0 &-3 &-3 &1
            \end{amat}
            \grstep{-\rho_2+\rho_3}
            \begin{amat}{3}
              1 &1  &2  &0 \\
              0 &-3 &-3 &1 \\
              0 &0  &0  &0
            \end{amat}                              \\
            \grstep{-(1/3)\rho_2}
            \begin{amat}{3}
              1 &1  &2  &0 \\
              0 &1  &1 &-1/3 \\
              0 &0  &0  &0
            \end{amat}
            \grstep{-\rho_2+\rho_1}
            \begin{amat}{3}
              1 &0  &1  &1/3 \\
              0 &1  &1 &-1/3 \\
              0 &0  &0  &0
            \end{amat}
          \end{multline*}
      \end{exparts}
    \end{answer}
  \recommended \item 
    Find the reduced echelon form of each matrix.
    \begin{exparts*}
      \partsitem \( \begin{mat}[r]
          2  &1  \\
          1  &3
        \end{mat}  \)
      \partsitem \( \begin{mat}[r]
          1  &3  &1  \\
          2  &0  &4  \\
         -1  &-3 &-3
        \end{mat}  \)
      \partsitem \( \begin{mat}[r]
          1  &0  &3  &1  &2  \\
          1  &4  &2  &1  &5  \\
          3  &4  &8  &1  &2
        \end{mat}  \)
      \partsitem \( \begin{mat}[r]
          0  &1  &3  &2  \\
          0  &0  &5  &6  \\
          1  &5  &1  &5
        \end{mat}  \)
    \end{exparts*}
    \begin{answer}
      Use Gauss-Jordan reduction.
      \begin{exparts}
        \partsitem The reduced echelon form is all zeroes except
          for a diagonal of ones.
          \begin{equation*}
            \grstep{-(1/2)\rho_1+\rho_2}
            \begin{mat}[r]
              2  &1  \\
              0  &5/2
            \end{mat}
            \grstep[(2/5)\rho_2]{(1/2)\rho_1}
            \begin{mat}[r]
              1  &1/2\\
              0  &1
            \end{mat}                      
            \grstep{-(1/2)\rho_2+\rho_1}
            \begin{mat}[r]
              1  &0  \\
              0  &1
            \end{mat}
          \end{equation*}
        \partsitem As in the prior problem, the reduced echelon form is 
          all zeroes but for a diagonal of ones.
          \begin{multline*}
            \grstep[\rho_1+\rho_3]{-2\rho_1+\rho_2}
            \begin{mat}[r]
              1  &3  &1  \\
              0  &-6 &2  \\
              0  &0  &-2
            \end{mat}
            \grstep[-(1/2)\rho_3]{-(1/6)\rho_2}
            \begin{mat}[r]
              1  &3  &1     \\
              0  &1  &-1/3  \\
              0  &0  &1
            \end{mat}                                      \\              
            \grstep[-\rho_3+\rho_1]{(1/3)\rho_3+\rho_2}
            \begin{mat}[r]
              1  &3  &0     \\
              0  &1  &0     \\
              0  &0  &1
            \end{mat}                  
            \grstep{-3\rho_2+\rho_1}
            \begin{mat}[r]
              1  &0  &0     \\
              0  &1  &0     \\
              0  &0  &1
            \end{mat}
           \end{multline*}
        \partsitem There are more columns than rows so we must get more
          than just a diagonal of ones.
          \begin{multline*}
            \grstep[-3\rho_1+\rho_3]{-\rho_1+\rho_2}
            \begin{mat}[r]
              1  &0  &3  &1  &2  \\
              0  &4  &-1 &0  &3  \\
              0  &4  &-1 &-2 &-4
            \end{mat}
            \grstep{-\rho_2+\rho_3}
            \begin{mat}[r]
              1  &0  &3  &1  &2  \\
              0  &4  &-1 &0  &3  \\
              0  &0  &0  &-2 &-7
            \end{mat}                           \\
            \grstep[-(1/2)\rho_3]{(1/4)\rho_2}
            \begin{mat}[r]
              1  &0  &3    &1  &2  \\
              0  &1  &-1/4 &0  &3/4  \\
              0  &0  &0    &1  &7/2
            \end{mat}                           
            \grstep{-\rho_3+\rho_1}
            \begin{mat}[r]
              1  &0  &3    &0  &-3/2  \\
              0  &1  &-1/4 &0  &3/4     \\
              0  &0  &0    &1  &7/2
            \end{mat}
          \end{multline*}
        \partsitem As in the prior item, this is not a square matrix. 
          \begin{multline*}
            \grstep{\rho_1\leftrightarrow\rho_3}
            \begin{mat}[r]
              1  &5  &1  &5  \\
              0  &0  &5  &6  \\
              0  &1  &3  &2
            \end{mat}
            \grstep{\rho_2\leftrightarrow\rho_3}
            \begin{mat}[r]
              1  &5  &1  &5  \\
              0  &1  &3  &2  \\
              0  &0  &5  &6
            \end{mat}                    \\
            \begin{aligned}
              &\grstep{(1/5)\rho_3}
              \begin{mat}[r]
                1  &5  &1  &5  \\
                0  &1  &3  &2  \\
                0  &0  &1  &6/5
              \end{mat}                  
              \grstep[-\rho_3+\rho_1]{-3\rho_3+\rho_2}
              \begin{mat}[r]
                1  &5  &0  &19/5  \\
                0  &1  &0  &-8/5  \\
                0  &0  &1  &6/5
              \end{mat}                   \\
              &\grstep{-5\rho_2+\rho_1}
              \begin{mat}[r]
                1  &0  &0  &59/5  \\
                0  &1  &0  &-8/5  \\
                0  &0  &1  &6/5
              \end{mat}
            \end{aligned}
          \end{multline*}
      \end{exparts}  
    \end{answer}
  \item Get the reduced echelon form of each.
    \begin{exparts*}
      \partsitem $
        \begin{mat}
          0  &2 &1 \\
          2  &-1 &1 \\
          -2 &-1 &0
        \end{mat}$
      \partsitem $
        \begin{mat}
          1 &3 &1 \\
          2 &6 &2 \\
         -1 &0 &0
        \end{mat}$
    \end{exparts*}
    \begin{answer}
      \begin{exparts}
        \partsitem Swap first.
          \begin{equation*}
            \grstep{\rho_1\leftrightarrow\rho_2}
            \grstep{\rho_1+\rho_3}
            \grstep{\rho_2+\rho_3}
            \grstep[(1/2)\rho_2 \\ (1/2)\rho_3]{(1/2)\rho_1}
            \grstep[(-1/2)\rho_3+\rho_2]{(-1/2)\rho_3+\rho_1}
            \grstep{(1/2)\rho_2+\rho_1}
            \begin{mat}
              1 &0 &0 \\
              0 &1 &0 \\
              0 &0 &1
            \end{mat}
          \end{equation*}
        \partsitem Here the swap is in the middle.
        \begin{equation*}
          \grstep[\rho_1+\rho_3]{-2\rho_1+\rho_2}
          \grstep{\rho_2\leftrightarrow\rho_3}
          \grstep{(1/3)\rho_2}
          \grstep{-3\rho_2+\rho_1}
          \begin{mat}
            1 &0 &0   \\
            0 &1 &1/3 \\
            0 &0 &0
          \end{mat}
        \end{equation*}
      \end{exparts}
    \end{answer}
  \recommended \item 
    Find each solution set by using Gauss-Jordan reduction and
    then reading off the parametrization.
    \begin{exparts}
      \partsitem \( \begin{linsys}[t]{3}
                  2x  &+  &y  &-  &z  &=  &1  \\
                  4x  &-  &y  &   &   &=  &3  
                  \end{linsys}  \)
      \partsitem \( \begin{linsys}[t]{4}
                   x  &   &   &-  &z  &   &   &=  &1  \\
                      &   &y  &+  &2z &-  &w  &=  &3  \\
                   x  &+  &2y &+  &3z &-  &w  &=  &7  
                    \end{linsys}  \)
      \partsitem \( \begin{linsys}[t]{4}
                   x  &-  &y  &+  &z  &   &   &=  &0  \\
                      &   &y  &   &   &+  &w  &=  &0  \\
                  3x  &-  &2y &+  &3z &+  &w  &=  &0  \\
                      &   &-y &   &   &-  &w  &=  &0  
                  \end{linsys}  \)
      \partsitem \( \begin{linsys}[t]{5}
                   a  &+  &2b &+  &3c &+  &d  &-  &e  &=  &1  \\
                  3a  &-  &b  &+  &c  &+  &d  &+  &e  &=  &3  
                  \end{linsys}  \)
    \end{exparts}
    \begin{answer}
      For the Gauss's halves, see the answers to Chapter One's
      section~I.2 question
      \nearbyexercise{exer:SlvMatNot}.
      \begin{exparts}
      \partsitem The ``Jordan'' half goes this way.
        \begin{equation*}
          \grstep[-(1/3)\rho_2]{(1/2)\rho_1}
          \begin{amat}[r]{3}
            1  &1/2 &-1/2 &1/2  \\
            0  &1   &-2/3 &-1/3
          \end{amat}                           
          \grstep{-(1/2)\rho_2+\rho_1}
          \begin{amat}[r]{3}
            1  &0   &-1/6 &2/3  \\
            0  &1   &-2/3 &-1/3
          \end{amat}
        \end{equation*}
        The solution set is this
        \begin{equation*}
          \set{\colvec[r]{2/3 \\ -1/3 \\ 0}
               +\colvec[r]{1/6 \\ 2/3 \\ 1}z
              \suchthat z\in\Re}
        \end{equation*}
      \partsitem The second half is
        \begin{equation*}
          \grstep{\rho_3+\rho_2}
          \begin{amat}[r]{4}
            1  &0  &-1  &0  &1 \\
            0  &1  &2   &0  &3 \\
            0  &0  &0   &1  &0
          \end{amat}
        \end{equation*}
        so the solution is this.
        \begin{equation*}
          \set{\colvec[r]{1 \\ 3 \\ 0 \\ 0}
               +\colvec[r]{1 \\ -2 \\ 1 \\ 0}z
              \suchthat z\in\Re}
        \end{equation*}
      \partsitem This Jordan half
        \begin{equation*}
          \grstep{\rho_2+\rho_1}
          \begin{amat}[r]{4}
            1  &0  &1   &1  &0 \\
            0  &1  &0   &1  &0 \\
            0  &0  &0   &0  &0 \\
            0  &0  &0   &0  &0
          \end{amat}
        \end{equation*}
        gives 
        \begin{equation*}
          \set{\colvec[r]{0 \\ 0 \\ 0 \\ 0}
               +\colvec[r]{-1 \\ 0 \\ 1 \\ 0}z
               +\colvec[r]{-1 \\ -1 \\ 0 \\ 1}w
              \suchthat z,w\in\Re}
        \end{equation*}
        (of course, we could omit the zero vector from the description).
      \partsitem The ``Jordan'' half
        \begin{align*}
          &\grstep{-(1/7)\rho_2}
          \begin{amat}[r]{5}
            1  &2  &3   &1   &-1   &1  \\
            0  &1  &8/7 &2/7 &-4/7 &0
          \end{amat}                   \\
          &\grstep{-2\rho_2+\rho_1}
          \begin{amat}[r]{5}
            1  &0  &5/7 &3/7 &1/7  &1  \\
            0  &1  &8/7 &2/7 &-4/7 &0
          \end{amat}
        \end{align*}
        ends with this solution set.
        \begin{equation*}
          \set{\colvec[r]{1 \\ 0 \\ 0 \\ 0 \\ 0}
               +\colvec[r]{-5/7 \\ -8/7 \\ 1 \\ 0 \\ 0}c
               +\colvec[r]{-3/7 \\ -2/7 \\ 0 \\ 1 \\ 0}d
               +\colvec[r]{-1/7 \\ 4/7 \\ 0 \\ 0 \\ 1}e
              \suchthat c,d,e\in\Re}
        \end{equation*}
    \end{exparts}
   \end{answer}
  \item 
    Give two distinct echelon form versions of this matrix.
    \begin{equation*}
      \begin{mat}[r]
        2  &1  &1  &3  \\
        6  &4  &1  &2  \\
        1  &5  &1  &5
      \end{mat}
    \end{equation*}
    \begin{answer}
      Routine Gauss's Method gives one:
      \begin{equation*}
        \grstep[-(1/2)\rho_1+\rho_3]{-3\rho_1+\rho_2}
        \begin{mat}[r]
          2  &1  &1  &3  \\
          0  &1  &-2 &-7 \\
          0  &9/2&1/2&7/2
        \end{mat}
        \grstep{-(9/2)\rho_2+\rho_3}
        \begin{mat}[r]
          2  &1  &1    &3  \\
          0  &1  &-2   &-7 \\
          0  &0  &19/2 &35
        \end{mat}
      \end{equation*}
      and any cosmetic change, such as multiplying the bottom row by \( 2 \),
      \begin{equation*}
        \begin{mat}[r]
          2  &1  &1    &3  \\
          0  &1  &-2   &-7 \\
          0  &0  &19   &70
        \end{mat}
      \end{equation*}
      gives another.  
    \end{answer}
  \recommended \item \label{exer:PossRedEchFrms} 
    List the reduced echelon forms possible for each size.
    \begin{exparts*}
      \partsitem \( \nbyn{2} \)
      \partsitem \( \nbym{2}{3} \)
      \partsitem \( \nbym{3}{2} \)
      \partsitem \( \nbyn{3} \)
    \end{exparts*}
    \begin{answer}
      In the cases listed below, we take $a,b\in\Re$.
      Thus, some canonical forms 
      listed below actually include infinitely many cases.
      In particular, they includes the cases $a=0$ and $b=0$.
      \begin{exparts}
        \partsitem 
          $\begin{mat}[r]
            0  &0  \\
            0  &0
          \end{mat}$,
          $\begin{mat}[r]
            1  &a  \\
            0  &0
          \end{mat}$, 
          $\begin{mat}[r]
            0  &1  \\
            0  &0
          \end{mat}$, 
          $\begin{mat}[r]
            1  &0  \\
            0  &1
          \end{mat}$
        \partsitem
          $\begin{mat}[r]
               0  &0  &0  \\
               0  &0  &0
             \end{mat}$,
          $\begin{mat}[r]
               1  &a  &b  \\
               0  &0  &0
             \end{mat}$,
          $\begin{mat}[r]
               0  &1  &a  \\
               0  &0  &0
             \end{mat}$,
          $\begin{mat}[r]
               0  &0  &1  \\
               0  &0  &0
             \end{mat}$,
          $\begin{mat}[r]
               1  &0  &a  \\
               0  &1  &b
             \end{mat}$,
          $\begin{mat}[r]
               1  &a  &0  \\
               0  &0  &1
             \end{mat}$,
          $\begin{mat}[r]
               0  &1  &0  \\
               0  &0  &1
             \end{mat}$
        \partsitem
          $\begin{mat}[r]
               0  &0  \\
               0  &0  \\
               0  &0
             \end{mat}$,
          $\begin{mat}[r]
               1  &a  \\
               0  &0  \\
               0  &0
             \end{mat}$,
          $\begin{mat}[r]
               0  &1  \\
               0  &0  \\
               0  &0
             \end{mat}$,
          $\begin{mat}[r]
               1  &0  \\
               0  &1  \\
               0  &0
             \end{mat}$
        \partsitem
          $\begin{mat}[r]
               0  &0  &0  \\
               0  &0  &0  \\
               0  &0  &0
             \end{mat}$,
          $\begin{mat}[r]
               1  &a  &b  \\
               0  &0  &0  \\
               0  &0  &0
             \end{mat}$,
          $\begin{mat}[r]
               0  &1  &a  \\
               0  &0  &0  \\
               0  &0  &0
             \end{mat}$,
          $\begin{mat}[r]
               0  &1  &0  \\
               0  &0  &1  \\
               0  &0  &0
             \end{mat}$,
          $\begin{mat}[r]
               0  &0  &1  \\
               0  &0  &0  \\
               0  &0  &0
             \end{mat}$,
          $\begin{mat}[r]
               1  &0  &a  \\
               0  &1  &b  \\
               0  &0  &0
             \end{mat}$,
          $\begin{mat}[r]
               1  &a  &0  \\
               0  &0  &1  \\
               0  &0  &0
             \end{mat}$,
          $\begin{mat}[r]
               1  &0  &0  \\
               0  &1  &0  \\
               0  &0  &1
             \end{mat}$
      \end{exparts}  
    \end{answer}
  \recommended \item  
    What results from applying Gauss-Jordan reduction to a
    nonsingular matrix?
    \begin{answer}
      A nonsingular homogeneous linear system has a unique solution.
      So a nonsingular matrix must reduce to a (square) 
      matrix that is all \( 0 \)'s
      except for \( 1 \)'s down the upper-left to lower-right diagonal, 
      such as these.
      \begin{equation*}
         \begin{mat}[r]
           1  &0  \\
           0  &1  \\
         \end{mat}
         \qquad
         \begin{mat}[r]
           1  &0  &0  \\
           0  &1  &0  \\
           0  &0  &1
         \end{mat}
      \end{equation*}  
    \end{answer}
 \item Decide whether each relation is an equivalence on the set of 
   $\nbyn{2}$ matrices.
   \begin{exparts}
     \partsitem two matrices are related if they have the same entry
       in the first row and first column
     % \partsitem two matrices are related if their four entries sum to the 
     %   same total
     \partsitem two matrices are related if they have the same entry
       in the first row and first column, or the same entry in the 
       second row and second column
   \end{exparts}
   \begin{answer}
     \begin{exparts}
       \partsitem  This is an equivalence.
        
          We can write $M_1\sim M_2$ if they are related.
          The $\sim$~relation is reflexive because any matrix has the 
          same $1,1$ entry as itself.
          The relation is symmetric because if $M_1$ has the same $1,1$
          entry as~$M_2$ then clearly also $M_2$ has the same $1,1$ entry
          as~$M_1$.
          Finally, the relation is transitive because if $M_1\sim M_2$ so they
          have the same $1,1$ entry as each other, 
          and $M_2\sim M_3$ so they have the 
          same as each other, then all three have the same~$1,1$ entry, 
          and $M_1\sim M_3$. 
       % \partsitem  This is an equivalence.
       %    Write $M_1\sim M_2$ if they are related.

       %    This relation is reflexive because any matrix has the 
       %    same sum of entries as itself.
       %    The $\sim$~relation is symmetric because if $M_1$ has the same
       %    sum of 
       %    entries as~$M_2$, then also $M_2$ has the same sum of entries
       %    as~$M_1$.
       %    Finally, the relation is transitive because if $M_1\sim M_2$ so their
       %    entry sum is the same, 
       %    and $M_2\sim M_3$ so they have the 
       %    same entry sum as each other, then all three have the same entry sum, 
       %    and $M_1\sim M_3$. 
       \partsitem
          This is not an equivalence because it is not transitive.
          The first and second matrix below are related by their $1,1$ entries,
          and the second and third are related by their $2,2$~entries.
          But the first and third are not related.
          \begin{equation*}
            \begin{mat}
              1 &0 \\
              0 &0
            \end{mat}
            \quad
            \begin{mat}
              1 &0 \\
              0 &-1
            \end{mat}
            \quad
            \begin{mat}
              0 &0 \\
              0 &-1
            \end{mat}
          \end{equation*}
     \end{exparts}
   \end{answer}
 \item \cite{Cleary}
    Consider the following relationship on the set of $\nbyn{2}$ matrices:  
    we say that $A$ is \textit{sum-what like} $B$ if the sum of all of 
    the entries in $A$ is the same as the sum of all the entries in $B$.  
    For instance, the zero matrix would be sum-what like the matrix 
    whose first row had two sevens, and whose second row had two 
    negative sevens.
    Prove or disprove that this is an equivalence relation on the set 
    of $\nbyn{2}$ matrices.
    \begin{answer}
      It is an equivalence relation.
      To prove that we must check that the relation 
      is reflexive, symmetric, and transitive.

      Assume that all matrices are $\nbyn{2}$.
      For reflexive, we note that a matrix has the same sum of entries as
      itself.
      For symmetric, we assume $A$ has the same sum of entries as~$B$ 
      and obviously then $B$ has the same sum of entries as~$A$.
      Transitivity is no harder\Dash if $A$ has the same sum of entries
      as $B$ and $B$ has the same sum of entries as $C$ then 
      $A$ has the same as $C$.
    \end{answer}
 \item \label{exer:INotJMakesRowOpsRev}
   The proof of \nearbylemma{le:RowOpsRev} contains a reference to the 
   $i\neq j$ condition on the row combination operation.
   \begin{exparts}
     \partsitem Write down a $\nbyn{2}$ matrix with nonzero entries,
        and show that the $-1\cdot\rho_1+\rho_1$ operation is not
        reversed by $1\cdot\rho_1+\rho_1$.
     \partsitem Expand the proof of that lemma to make explicit exactly where 
        it uses the $i\neq j$ condition on combining.
   \end{exparts}
   \begin{answer}
    \begin{exparts}
      \partsitem For instance,
        \begin{equation*}
          \begin{mat}[r]
            1  &2  \\
            3  &4  
          \end{mat}
          \grstep{-\rho_1+\rho_1}
          \begin{mat}[r]
            0  &0  \\
            3  &4  
          \end{mat}
          \grstep{\rho_1+\rho_1}
          \begin{mat}[r]
            0  &0  \\
            3  &4  
          \end{mat}
        \end{equation*}
        leaves the matrix changed.
      \partsitem This operation
        \begin{equation*}
          \begin{mat}
            \vdotswithin{a_{i,1}}                     \\
            a_{i,1}  &\cdots  &a_{i,n}  \\
            \vdotswithin{a_{i,1}}                     \\
            a_{j,1}  &\cdots  &a_{j,n}  \\
            \vdotswithin{a_{i,1}}                     
          \end{mat}
          \grstep{k\rho_i+\rho_j}
          \begin{mat}
            \vdotswithin{a_{i,1}}                                      \\
            a_{i,1}           &\cdots  &a_{i,n}          \\
            \vdotswithin{a_{i,1}}                                      \\
            ka_{i,1}+a_{j,1}  &\cdots  &ka_{i,n}+a_{j,n}  \\
            \vdotswithin{a_{i,1}}                     
          \end{mat}
        \end{equation*}      
        leaves the $i$-th row unchanged because of the $i\neq j$ restriction.
        Because the $i$-th row is unchanged, this operation 
        \begin{equation*}                                        
          \grstep{-k\rho_i+\rho_j}
          \begin{mat}
            \vdotswithin{a_{i,1}}                                      \\
            a_{i,1}           &\cdots  &a_{i,n}          \\
            \vdotswithin{a_{i,1}}                                      \\
            -ka_{i,1}+ka_{i,1}+a_{j,1}  &\cdots &-ka_{i,n}+ka_{i,n}+a_{j,n} \\
            \vdotswithin{a_{i,1}}                     
          \end{mat}
        \end{equation*}
        returns the $j$-th row to its original state.
        % (If $i=j$ then the third matrix would have entries of the 
        % form $-k(ka_{i,j}+a_{i,j})+ka_{i,j}+a_{i,j}$.)
    \end{exparts}
   \end{answer}
 \recommended \item \cite{Cleary}
  Consider the set of students in a class.  
  Which of the following relationships are equivalence relations?  
  Explain each answer in at least a sentence.
  \begin{exparts}
    \item  Two students $x, y$ are related 
      if $x$ has taken at least as many 
      math classes as $y$.
    \item Students $x, y$ are related if they have names 
      that start with the same letter.
  \end{exparts}
  \begin{answer}
    To be an equivalence, each relation must be reflexive, symmetric, and
    transitive.
    \begin{exparts}
      \item This relation 
        is not symmetric because if $x$ has taken $4$~classes and $y$
        has taken $3$ then $x$ is related to $y$ but $y$ is not related
        to $x$.
      \item This is reflexive because $x$'s name starts with the same
        letter as does $x$'s.
        It is symmetric because if $x$'s name starts with the same letter 
        as $y$'s then $y$'s starts with the same letter as does~$x$'s.
        And it is transitive because if $x$'s name starts with the same letter
        as does~$y$'s and $y$'s name starts with the same letter as 
        does $z$'s then $x$'s starts with the same letter as does $z$'s.
        So it is an equivalence.
    \end{exparts}
  \end{answer}
  \item
  Show that each of these is an equivalence on the set of $\nbyn{2}$
  matrices.
  Describe the equivalence classes.
  \begin{exparts}
    \item Two matrices are related if they have the same product down the
      diagonal, that is, if the product of the entries in the upper left and
      lower right are equal.
    \item Two matrices are related if they both have at least one entry 
      that is a~$1$,
      or if neither does.
  \end{exparts}
  \begin{answer}
    For each we must check the three conditions of reflexivity, symmetry, and
    transitivity.
    \begin{exparts}
      \item Any matrix clearly has the same product down the diagonal as
        itself, so the relation is reflexive.
        The relation is symmetric because if $A$ has the same product down
        its diagonal as does~$B$,
        if $a_{1,1}\cdot a_{2,2}=b_{1,1}\cdot b_{2,2}$, 
        then $B$ has the same product as does~$A$.
 
        Transitivity is similar: suppose that $A$'s product is~$r$ and 
        that it equals $B$'s product.
        Suppose also that $B$'s product equals $C$'s.
        Then all three have a product of~$r$, and $A$'s equals~$C$'s.

        There is an equivalence class for each real number, namely the 
        class contains all $\nbyn{2}$ matrices whose product down the
        diagonal is that real.
      \item For reflexivity, if the matrix $A$ has a~$1$ entry then it is 
        related to itself while if it does not then it is also related to 
        itself.
        Symmetry also has two cases: suppose that $A$ and~$B$ are related.
        If $A$ has a~$1$ entry then so does~$B$, and thus $B$ is related to~$A$.
        If $A$ has no~$1$ then neither does $B$, and again $B$ is related to
        $A$.
 
        For transitivity, suppose that $A$ is related to~$B$ and $B$ to~$C$.
        If $A$ has a~$1$ entry then so does~$B$, and because $B$ is related
        to~$C$, therefore so does~$C$,
        and hence $A$ is related to~$C$.
        Likewise, if $A$ has no~$1$ then neither does~$B$, and consequently
        neither does~$C$, giving the conclusion that $A$ is related to~$C$. 

        There are exactly two equivalence classes, one containing any
        $\nbyn{2}$ matrix that has at least one entry that is a~$1$,
        and the other containing all the matrices that have no $1$'s.
    \end{exparts}
  \end{answer}
  \item Show that each is not an equivalence on the
    set of $\nbyn{2}$ matrices.
    \begin{exparts}
      \item Two matrices $A,B$ are related if $a_{1,1}=-b_{1,1}$.
      \item Two matrices are related if the sum of their entries are
        within $5$, that is, $A$ is related to~$B$ if
        $\absval{(a_{1,1}+\cdots+a_{2,2})-(b_{1,1}+\cdots+b_{2,2})}<5$.
    \end{exparts}
    \begin{answer}
      \begin{exparts}
        \item This relation is not reflexive.
          For instance, any matrix with an upper-left entry of~$1$ is not
          related to itself.
        \item This relation is not transitive.
          For these three, $A$ is related to~$B$, and $B$ is related to~$C$,
          but $A$ is not related to~$C$.
          \begin{equation*}
            A=\begin{mat}
              0 &0 \\
              0 &0
            \end{mat},\quad
            B=\begin{mat}
              4 &0 \\
              0 &0
            \end{mat},\quad
            C=\begin{mat}
              8 &0 \\
              0 &0
            \end{mat},\quad
          \end{equation*}
      \end{exparts}
    \end{answer}
\end{exercises}




















\subsection{The Linear Combination Lemma}
We will close this chapter by proving 
that every matrix is row equivalent to one
and only one reduced echelon form matrix.
The ideas here will reappear, and be further developed, in the
next chapter.
% Here is an informal argument that the reduced 
% echelon form version of a matrix is unique.
% Consider again the example that started this section of a matrix that
% reduces to three different echelon form matrices.
% The first matrix of the three is the natural echelon form version.
% The second matrix is the same as 
% the first except that a row has been halved.
% The third matrix, too, is just a cosmetic variant of the first. 
% The definition of reduced echelon form outlaws this kind of fooling around.
% In reduced echelon form,
% halving a row is not possible because that would
% change the row's leading entry away from one, and
% neither is combining rows possible, because then a leading entry would no
% longer be alone in its column.

% This informal justification is not a proof;
% the argument shows that no two different reduced echelon form matrices
% are related by a single row operation step, but the argument does not
% rule out the possibility that two different reduced echelon form
% matrices could be related by multiple steps.
% Before we go to the proof, we finish this subsection by 
% rephrasing our work in a terminology that will be enlightening.
The crucial observation concerns 
how row operations act to transform one matrix into another:
the new rows are
linear combinations of the old.

\begin{example}
Consider this Gauss-Jordan reduction.
\begin{align*}
  \begin{amat}{2}
    2  &1  &0  \\
    1  &3  &5
  \end{amat}
  &\grstep{-(1/2)\rho_1+\rho_2}
  \begin{amat}{2}
    2  &1    &0  \\
    0  &5/2  &5
  \end{amat}                           \\
  &\grstep[(2/5)\rho_2]{(1/2)\rho_1}
  \begin{amat}{2}
    1  &1/2  &0  \\
    0  &1    &2
  \end{amat}                        \\
  &\grstep{-(1/2)\rho_2+\rho_1}\;
  \begin{amat}{2}
    1  &0    &-1  \\
    0  &1    &2
  \end{amat}                       
\end{align*}
Denoting those matrices $A\rightarrow D\rightarrow G\rightarrow B$
and writing the rows of $A$ as $\alpha_1$ and $\alpha_2$, etc., we have this.
\begin{align*}  % multline looks worse
  \left(\begin{array}{l}
    \alpha_1 \\
    \alpha_2
  \end{array}\right)
  &\grstep{-(1/2)\rho_1+\rho_2}
  \left(\begin{array}{l}
    \delta_1=\alpha_1 \\
    \delta_2=-(1/2)\alpha_1+\alpha_2
  \end{array}\right)                          \\
  &\grstep[(2/5)\rho_2]{(1/2)\rho_1}
  \left(\begin{array}{l}
    \gamma_1=(1/2)\alpha_1 \\
    \gamma_2=-(1/5)\alpha_1+(2/5)\alpha_2
  \end{array}\right)                          \\
  &\grstep{-(1/2)\rho_2+\rho_1}
  \left(\begin{array}{l}
    \beta_1=(3/5)\alpha_1-(1/5)\alpha_2  \\
    \beta_2=-(1/5)\alpha_1+(2/5)\alpha_2
  \end{array}\right)                         
\end{align*} 
\end{example}

\begin{example}
The fact that Gaussian operations combine rows linearly 
also holds if there is a row swap.
With this \( A \), \( D \), \( G \), and \( B \)
\begin{equation*}
    \begin{mat}[r]
       0  &2  \\
       1  &1
     \end{mat}
    \grstep{\rho_1\leftrightarrow\rho_2}
    \begin{mat}[r]
       1  &1  \\
       0  &2
     \end{mat}                  
    \grstep{(1/2)\rho_2}
    \begin{mat}[r]
       1  &1  \\
       0  &1
     \end{mat}                   
    \grstep{-\rho_2+\rho_1}
    \begin{mat}[r]
       1  &0  \\
       0  &1
     \end{mat}
\end{equation*}
we get these linear
relationships. 
\begin{multline*}
  \left(\begin{array}{l}
    \vec{\alpha}_1 \\
    \vec{\alpha}_2
  \end{array}\right)
  \,\grstep{\rho_1\leftrightarrow\rho_2}\,
  \left(\begin{array}{l}
     \vec{\delta}_1 =\vec{\alpha}_2   \\
     \vec{\delta}_2 =\vec{\alpha}_1
  \end{array}\right)                                                 
  \,\grstep{(1/2)\rho_2}\,
  \left(\begin{array}{l}
     \vec{\gamma}_1 =\vec{\alpha}_2  \\
     \vec{\gamma}_2 =(1/2)\vec{\alpha}_1
  \end{array}\right)                                               \\
  \,\grstep{-\rho_2+\rho_1}\,
  \left(\begin{array}{l}
     \vec{\beta}_1 =(-1/2)\vec{\alpha}_1+1\cdot\vec{\alpha}_2   \\
     \vec{\beta}_2 =(1/2)\vec{\alpha}_1
  \end{array}\right)
\end{multline*}
\end{example}

In summary, 
Gauss's Method systematically finds a suitable 
sequence of linear combinations 
of the rows.

% \begin{definition}
% A \definend{linear combination}\index{linear combination} of 
% \( x_1,\ldots,x_m \)
% is an expression of the form $c_1x_1+c_2x_2+\,\cdots\,+c_mx_m$
% where the \( c \)'s are scalars.
%\end{definition}

% \noindent (We have already used the phrase 
% `linear combination' in this book.
% The meaning is unchanged, but the next result's statement makes
% a more formal definition in order.)

\begin{lemma}[Linear Combination Lemma] \label{lm:LinearCombinationLemma} \index{Linear Combination Lemma}
%<*lm:LinearCombinationLemma>
A linear combination of linear combinations is a linear combination.
%</lm:LinearCombinationLemma>
\end{lemma}

\begin{proof}
%<*pf:LinearCombinationLemma>
Given the set  
$c_{1,1}x_1+\dots+c_{1,n}x_n$ through $c_{m,1}x_1+\dots+c_{m,n}x_n$
of linear combinations of the $x$'s,
consider a combination of those
\begin{equation*}
  d_1(c_{1,1}x_1+\dots+c_{1,n}x_n)\,+\dots+\,d_m(c_{m,1}x_1+\dots+c_{m,n}x_n)
\end{equation*}
where the $d$'s are scalars along with the $c$'s.
Distributing those $d$'s and regrouping gives
\begin{equation*}
  %&=d_1c_{1,1}x_1+\dots+d_1c_{1,n}x_n\,
  % +d_2c_{2,1}x_1+\dots+\,
  % d_mc_{1,1}x_1+\dots+d_mc_{1,n}x_n         \\
  =(d_1c_{1,1}+\dots+d_mc_{m,1})x_1\,+\dots+\,(d_1c_{1,n}+\dots+d_mc_{m,n})x_n
\end{equation*}
which is also a linear combination of the $x$'s.
%</pf:LinearCombinationLemma>
\end{proof}


\begin{corollary} \label{cor:RowsOfEqMatsLinCombos}
%<*co:RowsOfEqMatsLinCombos>
Where one matrix reduces to another, each row of the second
is a linear combination of the rows of the first.
%</co:RowsOfEqMatsLinCombos>
\end{corollary}

% The proof 
% uses induction. % \appendrefs{mathematical induction}\spacefactor=1000 %
% Before we proceed, here is an outline of the argument.
% For the base step, we
% will verify that the proposition is true when reduction 
% can be done in zero row operations.
% For the inductive step, we will 
% argue that if being able to reduce the first matrix to the second in some
% number $t\geq 0$ of operations implies that each row of the second is a linear
% combination of the rows of the first, then being able to reduce the first to
% the second in $t+1$ operations implies the same thing.
% Together these prove the result because  
% the base step shows that it is true in the zero operations case,
% and then the inductive step
% implies that it is true in the one operation case, and then the inductive step
% applied again gives that it is therefore true for two operations, etc.

\begin{proof}
%<*pf:RowsOfEqMatsLinCombos0>
For any two interreducible matrices $A$ and~$B$ there is some
minimum number of row operations that will take one to the other.
We proceed by induction on that number. 

In the base step, that we can go from one matrix to another using
zero reduction operations, the two are equal.
Then each row of $B$ is trivially a combination of
$A$'s rows $\vec{\beta}_i
  =0\cdot\vec{\alpha}_1+\cdots+1\cdot\vec{\alpha}_i+\cdots+0\cdot\vec{\alpha}_m$.
%</pf:RowsOfEqMatsLinCombos0>

%<*pf:RowsOfEqMatsLinCombos1>
For the inductive step assume the inductive hypothesis:~with $k\geq 0$,
any matrix that can be derived from \( A \) in \( k \) or fewer operations 
has rows that are linear combinations of $A$'s rows.
Consider a matrix~$B$ such that reducing \( A \) to~\( B \) 
requires $k+1$ operations.
In that reduction there is a next-to-last matrix~$G$,  
so that $A\longrightarrow\cdots\longrightarrow G\longrightarrow B$.
The inductive hypothesis applies to this \( G \) 
because it is only $k$ steps away from \( A \). 
That is, each row of \( G \)
is a linear combination of the rows of \( A \).
%</pf:RowsOfEqMatsLinCombos1>

%<*pf:RowsOfEqMatsLinCombos2>
We will verify that the rows of~$B$ are linear combinations of the rows
of~$G$.
Then the Linear Combination Lemma, \nearbylemma{lm:LinearCombinationLemma},
applies to show that the rows of~$B$ are linear combinations of the rows
of~$A$.

If the row operation taking \( G \) to~\( B \) is a swap then
the rows of $B$ are just the rows of $G$ reordered and each row of $B$
is a linear combination of the rows of $G$.
If the operation taking $G$ to~$B$ is multiplication of a row by a 
scalar~$c\rho_i$
then $\vec{\beta}_i=c\vec{\gamma}_i$ and the other rows are unchanged.
Finally, if the row operation is adding a multiple of one row to 
another $r\rho_i+\rho_j$ 
then only row~$j$ of $B$ differs from the matching row of~$G$, and 
$\vec{\beta}_j=r\gamma_i+\gamma_j$, 
which is indeed a linear combinations of the rows of $G$.
%</pf:RowsOfEqMatsLinCombos2>

Because we have proved both a base step and an inductive step,  
the proposition follows by the principle of mathematical induction.
\end{proof}

We now have the insight that Gauss's Method builds 
linear combinations of the rows.
But of course its goal is to end in echelon form, since that is 
a particularly
basic version of a linear system,
as it has isolated the variables.
For instance, in this matrix
\begin{equation*}
  R=\begin{mat}[r]
    2  &3  &7  &8  &0  &0  \\
    0  &0  &1  &5  &1  &1  \\
    0  &0  &0  &3  &3  &0  \\
    0  &0  &0  &0  &2  &1
  \end{mat}
\end{equation*}
$x_1$ has been removed from $x_5$'s equation.
That is, Gauss's Method has made $x_5$'s row in some way independent 
of $x_1$'s row.

The following result makes this intuition precise.
We sometimes refer to Gauss's Method as Gaussian elimination.
What it eliminates is linear relationships among the rows.

\begin{lemma}      \label{le:EchFormNoLinCombo}
%<*lm:EchFormNoLinCombo>
In an echelon form matrix,
no nonzero row is a linear combination of the other nonzero rows.
%</lm:EchFormNoLinCombo>
\end{lemma}

\begin{proof}
%<*pf:EchFormNoLinCombo0>
Let $R$ be an echelon form matrix and consider its  non-\( \vec{0} \)  rows.
First observe that
if we have a row written as a combination
of the others
$\vec{\rho}_i=c_1\vec{\rho}_1+\cdots+c_{i-1}\vec{\rho}_{i-1}+
               c_{i+1}\vec{\rho}_{i+1}+\cdots+c_m\vec{\rho}_m$
then we can rewrite that equation as
\begin{equation*}
   \vec{0}=c_1\vec{\rho}_1+\cdots+c_{i-1}\vec{\rho}_{i-1}+c_i\vec{\rho}_i+
               c_{i+1}\vec{\rho}_{i+1}+\cdots+c_m\vec{\rho}_m
  \tag{$*$}
\end{equation*}
where not all the coefficients are zero; specifically, $c_i=-1$.
The converse holds also:~given equation~($*$) where some $c_i\neq 0$ we 
could express $\vec{\rho}_i$ as a combination of the other rows by 
moving $c_i\vec{\rho}_i$ to the left and dividing by $-c_i$.
Therefore we will have proved the theorem if we
show that in~($*$) all of the coefficients are~$0$.
For that we use induction on the row number~$i$.
%</pf:EchFormNoLinCombo0>

%<*pf:EchFormNoLinCombo1>
The base case is the first row~$i=1$
(if there is no such nonzero row, so that $R$ is the zero matrix, then the
lemma holds vacuously).
Let $\ell_i$ be the column number of the leading entry in row~$i$.
Consider the entry of each row that is in column~$\ell_1$.
Equation~($*$) gives this.
\begin{equation*}
  0=c_1r_{1,\ell_1}+c_2r_{2,\ell_1}+\cdots+c_mr_{m,\ell_1}
  \tag{$**$}
\end{equation*}
The matrix is in echelon form so
every row after the first has a zero entry in that column 
$r_{2,\ell_1}=\cdots=r_{m,\ell_1}=0$.
Thus equation~($**$) shows that $c_1=0$, because $r_{1,\ell_1}\neq 0$ as it
leads the row.
%</pf:EchFormNoLinCombo1>

%<*pf:EchFormNoLinCombo2>
The inductive step is much the same as the base step.
Again consider equation~($*$).
We will prove that if the coefficient $c_i$ is $0$
for each row index $i\in\set{1,\ldots,k}$ 
then $c_{k+1}$ is also $0$. 
We focus on the entries from column~$\ell_{k+1}$. 
\begin{equation*}
  0=c_1r_{1,\ell_{k+1}}+\cdots+c_{k+1}r_{k+1,\ell_{k+1}}+\cdots+c_mr_{m,\ell_{k+1}}
\end{equation*}
By the inductive hypothesis $c_1$, \ldots $c_k$ are
all $0$ so this reduces to the equation
$0=c_{k+1}r_{k+1,\ell_{k+1}}+\cdots+c_mr_{m,\ell_{k+1}}$.
The matrix is in echelon form so the entries
$r_{k+2,\ell_{k+1}}$, \ldots, $r_{m,\ell_{k+1}}$ are all~$0$.
Thus $c_{k+1}=0$, because $r_{k+1,\ell_{k+1}}\neq 0$ as it is the leading entry.
%</pf:EchFormNoLinCombo2>
\end{proof}

With that, we are ready to show that the end product of Gauss-Jordan reduction
is unique.

\begin{theorem}
\label{th:ReducedEchelonFormIsUnique}
%<*th:ReducedEchelonFormIsUnique>
Each matrix is row equivalent to a unique reduced echelon form matrix.
%</th:ReducedEchelonFormIsUnique>
\end{theorem}

\begin{proof} \cite{Yuster}
%<*pf:ReducedEchelonFormIsUnique0>
Fix a number of rows \( m \).
We will proceed by induction on the number of columns \( n \).

The base case is that the matrix has \( n=1 \) column.
If this is the zero matrix then its echelon form is the zero matrix. 
If instead it has any nonzero entries then when the matrix is brought to 
reduced echelon form it must have at least one nonzero entry, which must be a
\( 1 \) in the first row. 
Either way, its reduced echelon form is unique.
%</pf:ReducedEchelonFormIsUnique0>

%<*pf:ReducedEchelonFormIsUnique1>
For the inductive step we assume that \( n>1 \) and that all \( m \)~row
matrices having fewer than~\( n \) columns have a unique reduced echelon form.
Consider an \( \nbym{m}{n} \) matrix \( A \) and suppose that 
\( B \) and \( C \) are two reduced echelon form matrices derived from \( A \).
We will show that these two must be equal.
%</pf:ReducedEchelonFormIsUnique1>

%<*pf:ReducedEchelonFormIsUnique2>
Let \( \hat{A} \) be the matrix consisting of the first \( n-1 \) columns of
\( A \).
Observe that 
% if an $n$~column matrix is in reduced echelon form then any initial
% set of its columns is also in reduced echelon form, so \( \hat{A} \)
% is in reduced echelon form. 
any sequence of row operations that bring \( A \) to reduced 
echelon form will also bring \( \hat{A} \) to reduced echelon form.
By the inductive hypothesis this reduced echelon form of \( \hat{A} \)
is unique, so if \( B \) and \( C \) differ then the difference must 
occur in column~\( n \).
%</pf:ReducedEchelonFormIsUnique2>

We finish the inductive step and the argument
by showing that the two cannot differ only in that column.
%<*pf:ReducedEchelonFormIsUnique3>
Consider a homogeneous system of equations for which \( A \) is the
matrix of coefficients.  
\begin{equation*}
  \begin{linsys}{4}
    a_{1,1}x_1  &+  &a_{1,2}x_2  &+  &\cdots  &+  &a_{1,n}x_n  &=  &0  \\
    a_{2,1}x_1  &+  &a_{2,2}x_2  &+  &\cdots  &+  &a_{2,n}x_n  &=  &0  \\
              &&&&&&&\vdotswithin{=}  \\
    a_{m,1}x_1  &+  &a_{m,2}x_2  &+  &\cdots  &+  &a_{m,n}x_n  &=  &0  
  \end{linsys}
  \tag{$*$}
\end{equation*}
By Theorem~One.I.\ref{th:GaussMethod}
%the first theorem of this chapter 
the set of solutions to that system
is the same as the set of solutions to $B$'s system
\begin{equation*}
  \begin{linsys}{4}
    b_{1,1}x_1  &+  &b_{1,2}x_2  &+  &\cdots  &+  &b_{1,n}x_n  &=  &0  \\
    b_{2,1}x_1  &+  &b_{2,2}x_2  &+  &\cdots  &+  &b_{2,n}x_n  &=  &0  \\
               &&&&&&&\vdotswithin{=}  \\
    b_{m,1}x_1  &+  &b_{m,2}x_2  &+  &\cdots  &+  &b_{m,n}x_n  &=  &0  
  \end{linsys}
  \tag{$**$}
\end{equation*}
and to $C$'s.
\begin{equation*}
  \quad
  \begin{linsys}{4}
    c_{1,1}x_1  &+  &c_{1,2}x_2  &+  &\cdots  &+  &c_{1,n}x_n  &=  &0  \\
    c_{2,1}x_1  &+  &c_{2,2}x_2  &+  &\cdots  &+  &c_{2,n}x_n  &=  &0  \\
               &&&&&&&\vdotswithin{=}  \\
    c_{m,1}x_1  &+  &c_{m,2}x_2  &+  &\cdots  &+  &c_{m,n}x_n  &=  &0  
  \end{linsys}
  \tag{$\mathord{*}\mathord{*}\mathord{*}$}
\end{equation*}
%</pf:ReducedEchelonFormIsUnique3>
%<*pf:ReducedEchelonFormIsUnique4>
With \( B  \) and \( C \) different only in column~\( n \), suppose 
that they differ in row~\( i \).
Subtract row~\( i \) of ($\mathord{*}\mathord{*}\mathord{*}$) from 
row~\( i \) of ($**$)
to get the equation 
\( (b_{i,n}-c_{i,n})\cdot x_n=0 \).
We've assumed that \( b_{i,n}\neq c_{i,n} \) and so we get $x_n=0$.
Thus $x_n$ is not a free variable and 
so in ($**$) and~($\mathord{*}\mathord{*}\mathord{*}$)
the \( n \)-th column contains the leading entry of some row, 
since in an echelon form matrix any
column that does not contain a leading entry is associated with 
a free variable.

But now, with \( B \) and \( C \) equal on 
the first \( n-1 \)~columns, by the definition of 
reduced echeleon form their leading entries in the 
\( n \)-th column are in the same row.
And, both leading entries would
have to be $1$, and would have to be the only nonzero entries in that column.
Therefore \( B=C \).
%</pf:ReducedEchelonFormIsUnique4>
\end{proof}

We have asked whether 
any two echelon form versions of a linear system 
have the same number of free variables, and if so are they
exactly the same variables?
With the prior result we can answer both questions ``yes.''
There is no linear system such
that, say, we could apply Gauss's Method 
one way and get $y$ and $z$ free but apply it another way 
and get $y$ and $w$ free.

Before the proof, 
recall the distinction between free variables and parameters.
This system
\begin{equation*}
  \begin{linsys}{3}
    x &+ &y &  &  &= &1 \\ 
      &  &y &+ &z &= &2  
  \end{linsys}
\end{equation*}
has one free variable,~$z$, because it is the only variable not leading a row.
We have the habit of parametrizing using the free variable $y=2-z$, $x=-1+z$,
but we could also parametrize using another variable, such as 
$z=2-y$, $x=1-y$. 
So the set of parameters is not unique, it is the set of free variables 
that is unique.

\begin{corollary}
If from a starting linear systems we derive by Gauss's Method two 
different echelon form systems, then the two have the same free variables.
\end{corollary}

\begin{proof}
The prior result says that the reduced echelon form is unique.
We get from any echelon form version to the reduced echelon form by 
eliminating up,
so any echelon form version of a system has the same free variables as the
reduced echelon form version. 
\end{proof}

We close with a recap.
In Gauss's Method we start with a matrix and then
derive a sequence of other matrices.
We defined two matrices to be related if we can derive one from the other.
That relation is an equivalence relation, %\appendrefs{equivalence relation} 
called row equivalence, and
so partitions the set of all matrices into row equivalence classes.
\begin{center}
  \includegraphics{gr/mp/ch1.30}
\end{center}
(There are infinitely many matrices in the pictured class, but we've only
got room to show two.)
We have proved there is one and only one reduced echelon form matrix in
each row equivalence class.
%<*ReducedEchelonFormIsCanonicalForm>
So the reduced echelon form is a
canonical form\appendrefs{canonical representatives}\spacefactor=1000
\index{canonical form!for row equivalence}\index{representative!for row equivalence classes}
for row equivalence:
the reduced echelon form matrices are
representatives of the classes.
%</ReducedEchelonFormIsCanonicalForm>
\begin{center}
  \includegraphics{gr/mp/ch1.31}
\end{center}

The idea here is that one way to understand a
mathematical situation is by being able to classify the cases that can happen.
This is a theme in this book and
we have seen this several times already.
We classified solution sets of linear systems into the no-elements, 
one-element, and infinitely-many elements cases.
We also classified linear systems with the same number of equations 
as unknowns into the nonsingular and singular cases.

Here, where we are investigating row equivalence, we know that the set of all
matrices breaks into the row equivalence classes
and we now have a way to put our finger on each of those
classes\Dash we can think of the matrices in a class 
as derived by row operations from the
unique reduced echelon form matrix in that class.

Put in more operational terms, uniqueness of reduced echelon form
lets us
answer questions about the classes by translating them
into questions about the representatives.
For instance, as promised in this section's opening, we now can
decide whether one matrix can be derived from another by row reduction.
We apply the Gauss-Jordan procedure to both and see if
they yield the same reduced echelon form.

\begin{example}  \label{ex:MatsNotRowEq}
These matrices are not row equivalent
\begin{equation*}
  \begin{mat}[r]
    1  &-3  \\
   -2  &6
  \end{mat}
  \qquad
  \begin{mat}[r]
    1  &-3  \\
   -2  &5
  \end{mat}
\end{equation*}
because their reduced echelon forms are not equal.
\begin{equation*}
  \begin{mat}[r]
    1  &-3  \\
    0  &0
  \end{mat}
  \qquad
  \begin{mat}[r]
    1  &0   \\
    0  &1
  \end{mat}
\end{equation*}
\end{example}

\begin{example}
Any nonsingular \( \nbyn{3} \) matrix Gauss-Jordan reduces to this.
\begin{equation*}
    \begin{mat}[r]
      1  &0  &0 \\
      0  &1  &0 \\
      0  &0  &1
    \end{mat}
\end{equation*}
\end{example}

\begin{example} \label{ex:RowEqClassTwoTwoMats}
We can describe all the classes by listing all possible
reduced echelon form matrices.
Any $\nbyn{2}$ matrix lies in one of these:~the class of matrices
row equivalent to this,
\begin{equation*}
  \begin{mat}[r]
     0  &0  \\
     0  &0
  \end{mat}
\end{equation*}
the infinitely many classes of matrices row equivalent to one of this type
\begin{equation*}
  \begin{mat}
     1  &a  \\
     0  &0
  \end{mat}
\end{equation*}
where \( a\in\Re \) (including $a=0$),
the class of matrices row equivalent to this,
\begin{equation*}
  \begin{mat}[r]
     0  &1  \\
     0  &0
  \end{mat}
\end{equation*}
and the class of matrices row equivalent to this
\begin{equation*}
  \begin{mat}[r]
     1  &0  \\
     0  &1
  \end{mat}
\end{equation*}
(this is the class of nonsingular $\nbyn{2}$ matrices).
\end{example}



\begin{exercises}
  \recommended \item 
    Decide if the matrices are row equivalent.
    \begin{exparts*}
       \partsitem \(
           \begin{mat}[r]
             1  &2  \\
             4  &8
           \end{mat}, 
           \begin{mat}[r]
             0  &1  \\
             1  &2
           \end{mat} \)
       \partsitem \(
           \begin{mat}[r]
             1  &0  &2  \\
             3  &-1 &1  \\
             5  &-1 &5
           \end{mat},  
           \begin{mat}[r]
             1  &0  &2  \\
             0  &2  &10 \\
             2  &0  &4
           \end{mat} \)
       \partsitem \(
           \begin{mat}[r]
             2  &1  &-1 \\
             1  &1  &0  \\
             4  &3  &-1
           \end{mat},  
           \begin{mat}[r]
             1  &0  &2  \\
             0  &2  &10 \\
           \end{mat} \)
       \partsitem \(
           \begin{mat}[r]
             1  &1  &1  \\
            -1  &2  &2
           \end{mat},  
           \begin{mat}[r]
             0  &3  &-1 \\
             2  &2  &5
           \end{mat} \)
       \partsitem \(
           \begin{mat}[r]
             1  &1  &1  \\
             0  &0  &3
           \end{mat},  
           \begin{mat}[r]
             0  &1  &2  \\
             1  &-1 &1
           \end{mat} \)
    \end{exparts*}
    \begin{answer}
      Bring each to reduced echelon form and compare.
      \begin{exparts}
        \partsitem The first gives
          \begin{equation*}
            \grstep{-4\rho_1+\rho_2}
            \begin{mat}[r]
              1  &2  \\
              0  &0
            \end{mat}
          \end{equation*}
          while the second gives
          \begin{equation*}
            \grstep{\rho_1\leftrightarrow\rho_2}
            \begin{mat}[r]
              1  &2  \\
              0  &1
            \end{mat}
            \grstep{-2\rho_2+\rho_1}
            \begin{mat}[r]
              1  &0  \\
              0  &1
            \end{mat}
          \end{equation*}
          The two reduced echelon form matrices are not identical, and so the
          original matrices are not row equivalent.
        \partsitem The first is this.
          \begin{equation*}
            \grstep[-5\rho_1+\rho_3]{-3\rho_1+\rho_2}
            \begin{mat}[r]
              1  &0  &2  \\
              0  &-1 &-5 \\
              0  &-1 &-5
            \end{mat}
            \grstep{-\rho_2+\rho_3}
            \begin{mat}[r]
              1  &0  &2  \\
              0  &-1 &-5 \\
              0  &0  &0
            \end{mat}
            \grstep{-\rho_2}
            \begin{mat}[r]
              1  &0  &2  \\
              0  &1  &5  \\
              0  &0  &0
            \end{mat}
          \end{equation*}
          The second is this.
          \begin{equation*}
            \grstep{-2\rho_1+\rho_3}
            \begin{mat}[r]
              1  &0  &2  \\
              0  &2  &10 \\
              0  &0  &0
            \end{mat}
            \grstep{(1/2)\rho_2}
            \begin{mat}[r]
              1  &0  &2  \\
              0  &1  &5  \\
              0  &0  &0
            \end{mat}
          \end{equation*}
          These two are row equivalent.
        \partsitem These two are not row equivalent because they have different
          sizes.
        \partsitem The first,
          \begin{equation*}
            \grstep{\rho_1+\rho_2}
            \begin{mat}[r]
              1  &1  &1  \\
              0  &3  &3
            \end{mat}
            \grstep{(1/3)\rho_2}
            \begin{mat}[r]
              1  &1  &1  \\
              0  &1  &1
            \end{mat}
            \grstep{-\rho_2+\rho_1}
            \begin{mat}[r]
              1  &0  &0  \\
              0  &1  &1
            \end{mat}
          \end{equation*}
          and the second.
          \begin{equation*}
            \grstep{\rho_1\leftrightarrow\rho_2}
            \begin{mat}[r]
              2  &2  &5  \\
              0  &3  &-1
            \end{mat}
            \grstep[(1/3)\rho_2]{(1/2)\rho_1}
            \begin{mat}[r]
              1  &1  &5/2 \\
              0  &1  &-1/3
            \end{mat}
            \grstep{-\rho_2+\rho_1}
            \begin{mat}[r]
              1  &0  &17/6 \\
              0  &1  &-1/3
            \end{mat}
          \end{equation*}
          These are not row equivalent.
        \partsitem Here the first is
          \begin{equation*}
            \grstep{(1/3)\rho_2}
            \begin{mat}[r]
              1  &1  &1  \\
              0  &0  &1
            \end{mat}
            \grstep{-\rho_2+\rho_1}
            \begin{mat}[r]
              1  &1  &0  \\
              0  &0  &1
            \end{mat}
          \end{equation*}
          while this is the second.
          \begin{equation*}
            \grstep{\rho_1\leftrightarrow\rho_2}
            \begin{mat}[r]
              1  &-1 &1  \\
              0  &1  &2
            \end{mat}
            \grstep{\rho_2+\rho_1}
            \begin{mat}[r]
              1  &0  &3  \\
              0  &1  &2
            \end{mat}
          \end{equation*}
          These are not row equivalent.
       \end{exparts}  
     \end{answer}
  \item 
    Which of these matrices are row equivalent to each other?
    \begin{exparts*}
      \partsitem
         \(\begin{mat}
             1 &3 \\
             2 &4
           \end{mat}\)
      \partsitem
         \(\begin{mat}
             1 &5 \\
             2 &10
           \end{mat}\)
      \partsitem
         \(\begin{mat}
             1 &-1 \\
             3 &0
           \end{mat}\)
      \partsitem
         \(\begin{mat}
             2 &6 \\
             4 &10
           \end{mat}\)
      \partsitem
         \(\begin{mat}
             0  &1 \\
             -1 &0
           \end{mat}\)
      \partsitem
         \(\begin{mat}
             3 &3 \\
             2 &2
           \end{mat}\)
    \end{exparts*}
    \begin{answer}
      Perform Gauss-Jordan reduction on each.
      Two matrices are row-equivalent if and only if they have the same
      reduced echelon form.
      Here is the reduced form for each. 
      % sage: M = matrix(QQ, [[1,3], [2,4]])
      % sage: gauss_jordan(M)
      % [1 3]
      % [2 4]
      %  take -2 times row 1 plus row 2
      % [ 1  3]
      % [ 0 -2]
      %  take -1/2 times row 2
      % [1 3]
      % [0 1]
      %  take -3 times row 2 plus row 1
      % [1 0]
      % [0 1]
      % sage: M = matrix(QQ, [[1,5], [2,10]])
      % sage: gauss_jordan(M)
      % [ 1  5]
      % [ 2 10]
      %  take -2 times row 1 plus row 2
      % [1 5]
      % [0 0]
      % sage: M = matrix(QQ, [[1,-1], [3,0]])
      % sage: gauss_jordan(M)
      % [ 1 -1]
      % [ 3  0]
      %  take -3 times row 1 plus row 2
      % [ 1 -1]
      % [ 0  3]
      %  take 1/3 times row 2
      % [ 1 -1]
      % [ 0  1]
      %  take 1 times row 2 plus row 1
      % [1 0]
      % [0 1]
      % sage: M = matrix(QQ, [[2,6], [4,10]])
      % sage: gauss_jordan(M)
      % [ 2  6]
      % [ 4 10]
      %  take -2 times row 1 plus row 2
      % [ 2  6]
      % [ 0 -2]
      %  take 1/2 times row 1
      %  take -1/2 times row 2
      % [1 3]
      % [0 1]
      %  take -3 times row 2 plus row 1
      % [1 0]
      % [0 1]
      % sage: M = matrix(QQ, [[0,1], [-1,0]])
      % sage: gauss_jordan(M)
      % [ 0  1]
      % [-1  0]
      %  swap row 1 with row 2
      % [-1  0]
      % [ 0  1]
      %  take -1 times row 1
      % [1 0]
      % [0 1]
      % sage: M = matrix(QQ, [[3,3], [2,2]])
      % sage: gauss_jordan(M)
      % [3 3]
      % [2 2]
      %  take -2/3 times row 1 plus row 2
      % [3 3]
      % [0 0]
      %  take 1/3 times row 1
      % [1 1]
      % [0 0]
      \begin{exparts*}
        \partsitem
          \(\begin{mat}
              1  &0  \\
              0  &1
            \end{mat} \)
        \partsitem
          \(\begin{mat}
              1  &5  \\
              0  &0
            \end{mat} \)
        \partsitem
          \(\begin{mat}
             1   &0  \\
             0   &1
            \end{mat} \)
        \partsitem
          \(\begin{mat}
             1   &0  \\
             0   &1
            \end{mat} \)
        \partsitem
          \(\begin{mat}
             1   &0  \\
             0   &1
            \end{mat} \)
        \partsitem
          \(\begin{mat}
             1   &1  \\
             0   &0
            \end{mat} \)
      \end{exparts*}
    \end{answer}
 \item Produce three other matrices row equivalent to the given one.
   \begin{exparts*}
     \partsitem 
       \(\begin{mat}
           1 &3 \\
           4 &-1
         \end{mat}\)
     \partsitem
       \(\begin{mat}
           0 &1 &2 \\
           1 &1 &1  \\
           2 &3 &4
         \end{mat}\)
   \end{exparts*}
   \begin{answer}
     For each you can just perform some row operations on the starting
     matrix.
     \begin{exparts}
       \partsitem
         Multiplying the first row by~$3$ gives this.
         \begin{equation*}
           \begin{mat}
             3 &9  \\
             4 &-1
           \end{mat}
         \end{equation*}
         (There is no sense to this particular choice of row operation; it is 
         just the first thing that came to mind.)
         Two other row operations are a row swap $\rho_1\leftrightarrow\rho_2$ 
         and adding $\rho_1+\rho_2$.
         \begin{equation*}
           \begin{mat}
             4 &-1  \\
             1 &3
           \end{mat}
           \quad
           \begin{mat}
             1 &3  \\
             5 &2
           \end{mat}
         \end{equation*}
       \partsitem Doing the same three arbitrary row operations gives these
         three.
         \begin{equation*}
          \begin{mat} 
           0 &3 &6 \\
           1 &1 &1  \\
           2 &3 &4  
          \end{mat}
          \quad         
          \begin{mat} 
           1 &1 &1  \\
           0 &1 &2 \\
           2 &3 &4  
          \end{mat}
          \quad         
          \begin{mat} 
           0 &1 &2 \\
           1 &2 &3  \\
           2 &3 &4  
          \end{mat}
          \quad         
         \end{equation*}
     \end{exparts}
   \end{answer}
  \recommended \item Perform Gauss's Method on this matrix.
    Express each row of the final matrix as a linear combination of 
    the rows of the starting matrix.
    \begin{equation*}
      \begin{mat}
        1 &2   &1 \\
        3 &-1  &0 \\
        0 &4   &0
      \end{mat}
    \end{equation*}
    \begin{answer}
      The Gaussian reduction is routine.
      % sage: load("gauss_method.sage")
      % sage: M = matrix (QQ, [[1,2,1], [3,-1,0], [0,4,0]])
      % sage: gauss_method(M)
      % [ 1  2  1]
      % [ 3 -1  0]
      % [ 0  4  0]
      %  take -3 times row 1 plus row 2
      % [ 1  2  1]
      % [ 0 -7 -3]
      % [ 0  4  0]
      %  take 4/7 times row 2 plus row 3
      % [    1     2     1]
      % [    0    -7    -3]
      % [    0     0 -12/7]
      \begin{equation*}
        \begin{mat}
          1 &2   &1 \\
          3 &-1  &0 \\
          0 &4   &0
        \end{mat}
        \grstep{-3\rho_1+\rho_2}    
        \begin{mat}
          1 &2   &1 \\
          0 &-7  &-3 \\
          0 &4   &0
        \end{mat}
        \grstep{(4/7)\rho_2+\rho_3}    
        \begin{mat}
          1 &2   &1 \\
          0 &-7  &-3 \\
          0 &0   &-12/7
        \end{mat}
      \end{equation*}
      Denoting those matrices $A$, $D$, and~$B$ respectively, we have this.
      \begin{align*}
        \begin{mat}
          \alpha_1 \\
          \alpha_2 \\
          \alpha_3
        \end{mat}
        &\grstep{-3\rho_1+\rho_2}    
        \begin{mat}
          \delta_1=\alpha_1 \\
          \delta_2=-3\alpha_1+\alpha_2 \\
          \delta_3=\alpha_3
        \end{mat}                                   \\
        & \grstep{(4/7)\rho_2+\rho_3}
        \begin{mat}
          \beta_1=\alpha_1  \\
          \beta_2=-3\alpha_1+\alpha_2 \\
          \beta_3=-(12/7)\alpha_1+(4/7)\alpha_2+\alpha_3
        \end{mat}
      \end{align*}
    \end{answer}
  \item 
     Describe the matrices in each of the classes represented in
     \nearbyexample{ex:RowEqClassTwoTwoMats}.
     \begin{answer}
       First, the only matrix row equivalent to the matrix of all
       \( 0 \)'s is itself (since row operations have no effect).

       Second, the matrices that reduce to 
       \begin{equation*}
         \begin{mat}
           1  &a  \\
           0  &0
         \end{mat}
       \end{equation*}
       have the form
       \begin{equation*}
         \begin{mat}
           b  &ba \\
           c  &ca
         \end{mat}
       \end{equation*}
       (where \( a,b,c\in\Re \), and \(b\) and \(c\) are not both zero).  

       Next, the matrices that reduce to 
       \begin{equation*}
         \begin{mat}[r]
           0  &1  \\
           0  &0
         \end{mat}
       \end{equation*}
       have the form
       \begin{equation*}
         \begin{mat}
           0  &a \\
           0  &b
         \end{mat}
       \end{equation*}
       (where \( a,b\in\Re \), and not both are zero).  

       Finally, the matrices that reduce to 
       \begin{equation*}
         \begin{mat}[r]
           1  &0  \\
           0  &1
         \end{mat}
       \end{equation*}
       are the nonsingular matrices.
       That's because a linear system for which this is the matrix of
       coefficients will have a unique solution, and that is the definition
       of nonsingular.
       (Another way to say the same thing is to say that they fall into none
       of the above classes.)
     \end{answer}
  \item 
    Describe all matrices in the row equivalence class of
    these.
    \begin{exparts*}
       \partsitem  \(
           \begin{mat}[r]
             1  &0  \\
             0  &0
           \end{mat}  \)
       \partsitem  \(
           \begin{mat}[r]
             1  &2      \\
             2  &4
           \end{mat} \)
       \partsitem \(
           \begin{mat}[r]
             1  &1      \\
             1  &3
           \end{mat} \)
    \end{exparts*}
    \begin{answer}
      \begin{exparts}
        \partsitem They have the form
          \begin{equation*}
            \begin{mat}
              a  &0  \\
              b  &0
            \end{mat}
          \end{equation*}
          where at least one of \( a,b\in\Re \) is nonzero.
        \partsitem They have this form 
          \begin{equation*}
            \begin{mat}
             a  &2a \\
             b  &2b
            \end{mat}
          \end{equation*}
          where at least one of \( a,b\in\Re \) is nonzero.
        \partsitem The given matrix is nonsingular.
          So the row equivalence class consists of all nonsinglar $\nbyn{2}$
          matrices.
          (To give a formula, they have the form
          \begin{equation*}
            \begin{mat}
              a  &b  \\
              c  &d
            \end{mat}
          \end{equation*}
          for \( a,b,c,d\in\Re \)) where \( ad-bc\neq 0 \).
          We will see in Chapter Four that this formula 
          determines when a \( \nbyn{2} \) matrix
          is nonsingular.)
      \end{exparts}  
    \end{answer}
  \item 
    How many row equivalence classes are there?
    \begin{answer}
       Infinitely many.
       For instance, in 
       \begin{equation*}
         \begin{mat}
           1  &k  \\
           0  &0
         \end{mat}
       \end{equation*}
       each $k\in\Re$ gives a different class.  
    \end{answer}
  \item 
    Can row equivalence classes contain different-sized matrices?
    \begin{answer}
      No.
      Row operations do not change the size of a matrix.  
    \end{answer}
  \item 
    How big are the row equivalence classes?
    \begin{exparts} 
      \partsitem Show that for any matrix of all zeros, the class is finite.
      \partsitem Do any other classes contain only finitely many members?
    \end{exparts}
    \begin{answer}
     \begin{exparts}
      \partsitem A row operation on a matrix of zeros has no effect.
        Thus each such matrix is alone in its row equivalence class.  
      \partsitem No.
        Any nonzero entry can be rescaled.
     \end{exparts}
    \end{answer}
  \recommended \item 
    Give two reduced echelon form matrices that have their leading
    entries in the same columns,
    but that are not row equivalent.
    \begin{answer}
      Here are two.
      \begin{equation*}
        \begin{mat}[r]
          1  &1  &0  \\
          0  &0  &1
        \end{mat}
        \quad\text{and}\quad
        \begin{mat}[r]
          1  &0  &0  \\
          0  &0  &1
        \end{mat}
      \end{equation*}  
     \end{answer}
  \recommended \item 
    Show that any two \( \nbyn{n} \) nonsingular matrices are
    row equivalent.
    Are any two singular matrices row equivalent?
    \begin{answer}
      Any two \( \nbyn{n} \) nonsingular matrices have
      the same reduced echelon
      form, namely the matrix with all \( 0 \)'s except for \( 1 \)'s down
      the diagonal.
      \begin{equation*}
        \begin{mat}
          1  &0  &       &0  \\
          0  &1  &       &0  \\
             &   &\ddots &   \\
          0  &0  &       &1
        \end{mat}
      \end{equation*}

      Two same-sized singular matrices need not be row equivalent.
      For example, these two \( \nbyn{2} \) singular matrices
      are not row equivalent.
      \begin{equation*}
        \begin{mat}[r]
          1  &1  \\
          0  &0
        \end{mat}
        \quad\text{and}\quad
        \begin{mat}[r]
          1  &0  \\
          0  &0
        \end{mat}
      \end{equation*}  
    \end{answer}
  \recommended \item 
    Describe all of the row equivalence classes containing these.
    \begin{exparts*}
      \partsitem \( \nbyn{2} \)~matrices
      \partsitem \( \nbym{2}{3} \)~matrices
      \partsitem \( \nbym{3}{2} \)~matrices
      \partsitem \( \nbyn{3} \)~matrices
    \end{exparts*}
    \begin{answer}
      Since there is one and only one reduced echelon form matrix in each
      class, we can just list the possible reduced echelon form matrices.

      For that list, see the answer for \nearbyexercise{exer:PossRedEchFrms}. 
    \end{answer}
  \item  
     \begin{exparts}
          \partsitem Show that a vector $\vec{\beta}_0$ is a linear combination
            of members of the set $\set{\vec{\beta}_1,\ldots,\vec{\beta}_n}$
            if and only if there is a linear relationship 
            $\zero=c_0\vec{\beta}_0+\cdots+c_n\vec{\beta}_n$
            where $c_0$ is not zero.
            (\textit{Hint.}   Watch out for the $\vec{\beta}_0=\zero$ case.)
         \partsitem Use that to simplify the proof of 
            \nearbylemma{le:EchFormNoLinCombo}.   
       \end{exparts}
       \begin{answer}
          \begin{exparts}
           \partsitem If there is a linear relationship where $c_0$ is not zero
             then we can subtract $c_0\vec{\beta}_0$ from both sides and divide
             by $-c_0$ to get $\vec{\beta}_0$ as a linear
             combination of the others.
             (Remark:  
             if there are no other vectors in the set\Dash if the 
             relationship is, say, 
             $\zero=3\cdot\zero$\Dash then the statement is still true because
             the zero vector is by definition the sum of the empty set 
             of vectors.)

             Conversely, if $\vec{\beta}_0$ is a combination of the others 
             $\vec{\beta}_0=c_1\vec{\beta}_1+\dots+c_n\vec{\beta}_n$
             then subtracting
             $\vec{\beta}_0$ from both sides gives a relationship where 
             at least one
             of the coefficients is nonzero; namely,
             the $-1$ in front of $\vec{\beta}_0$.
           \partsitem The first row is not a linear combination of the
             others for
             the reason given in the proof:~in the equation of components from
             the column containing the leading entry of the first row, the
             only nonzero entry is the leading entry from the first row, so
             its coefficient must be zero.
             Thus, from the prior part of this exercise, the first row is in
             no linear relationship with the other rows.

             Thus, when considering whether the second row can be in a linear 
             relationship
             with the other rows, we can leave the first row out.
             But now the argument just applied to the first row will apply
             to the second row.
             (That is, we are arguing here by induction.)             
         \end{exparts}
      \end{answer}
  % \item 
  %   Why, in the proof of \nearbytheorem{th:ReducedEchelonFormIsUnique},
  %   do we bother to restrict to the nonzero rows?
  %   Why not just stick to the relationship that we began with,
  %   $\beta_i=c_{i,1}\delta_1+\dots+c_{i,m}\delta_m$, with $m$ instead of $r$,
  %   and argue using it that the only nonzero coefficient
  %   is \( c_{i,i}  \), which is \( 1 \)?
  %   \begin{answer}
  %      The zero rows could have nonzero coefficients, and
  %      so the statement would not be true.
  %   \end{answer}
  \recommended \item 
   \cite{Trono}
   Three truck drivers went into a roadside cafe.
   One truck driver purchased four sandwiches, a cup of coffee, and ten 
   doughnuts for \$$8.45$.
   Another driver purchased three sandwiches, a cup of coffee, and seven
   doughnuts for \$$6.30$.
   What did the third truck driver pay for a sandwich, a cup of coffee, and 
   a doughnut?
   \begin{answer}
     We know that $4s+c+10d=8.45$ and that $3s+c+7d=6.30$, and we'd like to
     know what $s+c+d$ is.
     Fortunately, $s+c+d$ is a linear combination of $4s+c+10d$ and $3s+c+7d$.
     Calling the unknown price $p$, we have this reduction.
     \begin{align*}
       \begin{amat}{3}
         4  &1  &10  &8.45 \\
         3  &1  &7   &6.30 \\
         1  &1  &1   &p
       \end{amat}
       &\grstep[-(1/4)\rho_1+\rho_3]{-(3/4)\rho_1+\rho_2}
       \begin{amat}{3}
         4  &1    &10     &8.45      \\
         0  &1/4  &-1/2   &-0.037\,5 \\
         0  &3/4  &-3/2   &p-2.112\,5
       \end{amat}                                      \\
       &\grstep{-3\rho_2+\rho_3}
       \begin{amat}{3}
         4  &1    &10     &8.45      \\
         0  &1/4  &-1/2   &-0.037\,5 \\
         0  &0    &0      &p-2.00
       \end{amat}
     \end{align*}
     The price paid is \$$2.00$.
   \end{answer}
  % \item
  %  The fact that Gaussian reduction disallows multiplication of
  %  a row by zero is needed for the proof of uniqueness of reduced echelon form,
  %  or else every matrix would
  %  be row equivalent to a matrix of all zeros.
  %  Where is it used?
  %  \begin{answer}
  %    If multiplication of a row by zero were allowed then
  %    \nearbylemma{le:EquivMatsSameForm}
  %    would not hold.
  %    That is, where
  %    \begin{equation*}
  %      \begin{mat}[r]
  %        1  &3  \\
  %        2  &1
  %      \end{mat}
  %      \grstep{0\rho_2}
  %      \begin{mat}[r]
  %        1  &3  \\
  %        0  &0
  %      \end{mat}
  %    \end{equation*}
  %    all the rows of the second matrix can be expressed as linear combinations
  %    of the rows of the first, but the converse does not hold.
  %    The second row of the first matrix is not a linear combination of the
  %    rows of the second matrix.  
  %  \end{answer}
  \item 
   The Linear Combination Lemma says which equations can be gotten from
   Gaussian reduction of a given linear system.
   \begin{enumerate}
     \item Produce an equation not implied by this system.
       \begin{equation*}
         \begin{linsys}{2}
           3x  &+  &4y  &=  &8 \\
           2x  &+  & y  &=  &3 
         \end{linsys}
       \end{equation*}
     \item Can any equation be derived from an inconsistent system?
   \end{enumerate}
   \begin{answer}
     \begin{enumerate}
        \item An easy answer is this.
          \begin{equation*}
            0=3
          \end{equation*}
          For a less wise-guy-ish answer, solve the system:
          \begin{equation*}
            \begin{amat}[r]{2}
              3  &4  &8  \\
              2  &1   &3
            \end{amat}
            \grstep{-(2/3)\rho_1+\rho_2}
            \begin{amat}[r]{2}
              3  &4    &8    \\
              0  &-5/3 &-7/3
            \end{amat}
          \end{equation*}
          gives \( y=7/5 \) and \( x=4/5 \).
          Now any equation not satisfied by \( (7/5,4/5) \) will do,
          e.g., \( 5x+5y=10 \).
        \item Every equation can be derived from an inconsistent system.
          For instance, here is how to derive \( 3x+2y=4 \) from
          \( 0=5 \).
          First,
          \begin{equation*}
            0=5
            \grstep{(3/5)\rho_1}
            0=3
            \grstep{x\rho_1}
            0=3x
          \end{equation*}
          (validity of the \( x=0 \) case is separate but clear).
          Similarly, \( 0=2y \).
          Ditto for \( 0=4 \).
          But now, \( 0+0=0 \) gives \( 3x+2y=4 \).
     \end{enumerate}  
    \end{answer}
  \item 
    \cite{HoffmanKunze}
    Extend the definition of row equivalence to linear systems.
    Under your definition, do equivalent systems have the same solution set?
    \begin{answer}
      Define linear systems to be equivalent if their augmented
      matrices are row equivalent.
      The proof that equivalent systems have the same solution set is easy.  
    \end{answer}
  \item 
    In this matrix
    \begin{equation*}
      \begin{mat}[r]
        1  &2  &3  \\
        3  &0  &3  \\
        1  &4  &5
      \end{mat}
    \end{equation*}
    the first and second columns add to the third.
    \begin{exparts}
      \partsitem Show that remains true under any row operation.
      \partsitem Make a conjecture.
      \partsitem Prove that it holds.
    \end{exparts}
    \begin{answer}
      \begin{exparts}
        \partsitem The three possible row swaps are easy, 
          as are the three possible rescalings.
          One of the six possible row combinations is \( k\rho_1+\rho_2 \):
          \begin{equation*}
            \begin{mat}
              1           &2           &3  \\
              k\cdot 1+3  &k\cdot 2+0  &k\cdot 3+3  \\
              1           &4           &5
            \end{mat}
          \end{equation*}
          and again the first and second columns add to the third.
          The other five combinations are similar.
        \partsitem The obvious conjecture is that row operations do not change
          linear relationships among columns.
        \partsitem A case-by-case 
          proof follows the sketch given in the first item.
      \end{exparts}  
   \end{answer}
\end{exercises}

% END NATIVE SOURCE src/gr/gr3.tex

% BEGIN NATIVE SOURCE src/vs/vs1.tex
% Chapter 2, Section 1 _Linear Algebra_ Jim Hefferon
%  http://joshua.smcvt.edu/linearalgebra
%  2001-Jun-10
\chapter{Vector Spaces}
The first chapter finished with a fair
understanding
of how Gauss's Method solves a linear system.
It systematically takes linear combinations of the rows.
Here we move to a general study of linear 
combinations.

We need a setting.
At times in the first chapter we've combined vectors from $\Re^2$, 
at other times vectors from $\Re^3$,
and at other times vectors from higher-dimensional spaces. 
So our first impulse might be 
to work in $\Re^n$, leaving $n$ unspecified.
This would have the advantage that any of the results 
would hold for $\Re^2$ and for $\Re^3$ and for many other spaces,
simultaneously.

But if having the results apply to many spaces at once is
advantageous then sticking only to $\Re^n$'s is restrictive. 
We'd like our results to apply to combinations of row vectors,
as in the final section of the first chapter.
We've even seen some spaces that are not simply a collection of all of the
same-sized column vectors or row vectors.
For instance, we've seen a homogeneous system's solution
set that is a plane inside of $\Re^3$.
This set is a closed system in that 
a linear combination of these solutions is also a solution. 
But it does not contain all of the three-tall column vectors, 
only some of them.

We want the results about linear combinations to apply anywhere that linear
combinations make sense. 
We shall call any such set a \definend{vector space}.
Our results, instead of being phrased as
``Whenever we have a collection in which we can sensibly take linear 
combinations \ldots'', will be stated 
``In any vector space \ldots''

Such a statement describes at once what
happens in many spaces.
To understand the advantages of moving from studying a single space 
to studying a class of spaces, consider this analogy.
Imagine that the government made laws one person at a time:
``Leslie Jones can't jay walk.''
That would be bad; 
statements have the virtue of economy when they apply to many cases at once.
Or suppose that they said, ``Kim Ke must stop when passing 
an accident.''
Contrast that with, ``Any doctor must stop when passing 
an accident.''
More general statements, in some ways, are clearer.













\section{Definition of Vector Space}
\index{vector space|(}
We shall study structures with two operations,
an addition and a scalar multiplication, that are subject to some
simple conditions.
We will reflect more on the conditions later
but on first reading notice how reasonable they are.
For instance, surely any operation that can be called an addition
(e.g., column vector addition, row vector addition, or
real number addition) will satisfy conditions (1) through~(5) below.



\vspace*{1ex}
\subsection{Definition and Examples}

\begin{definition}
\label{def:VecSpace}
%<*df:VectorSpace>
A \definend{vector space}\index{vector space!definition}
(over \( \Re \)) consists of a set \( V \) along with
two operations 
`+'\index{vector!sum}\index{sum!vector}\index{addition of vectors} 
and `\( \cdot \)'\index{scalar multiple!vector}\index{vector!scalar multiple} 
subject to the conditions
that for all vectors \( \vec{v},\vec{w},\vec{u}\in V \)
and all \definend{scalars}\index{scalar}
\( r,s\in\Re \):
% \begin{tfae}
\begin{compactenum} 
\item the set $V$ is closed under
  vector addition, that is, 
  \( \vec{v}+\vec{w}\in V \)
\item vector addition is commutative,
  \( \vec{v}+\vec{w}=\vec{w}+\vec{v} \) 
\item vector addition is associative,
  \( (\vec{v}+\vec{w})+\vec{u}=\vec{v}+(\vec{w}+\vec{u}) \)
\item there is a \definend{zero vector}\index{zero vector}
    \( \zero\in V \) such that
    \( \vec{v}+\zero=\vec{v}\, \) for all \( \vec{v}\in V\/ \)
\item each \( \vec{v}\in V \) has an
    \definend{additive inverse}\index{additive inverse}\index{inverse!additive}
    \( \vec{w}\in V \) such that \( \vec{w}+\vec{v}=\zero \)
\item  the set $V$ is closed under
    scalar multiplication, that is, 
   \( r\cdot\vec{v}\in V \)
\item scalar multiplication distributes over scalar addition,
 \( (r+s)\cdot\vec{v}=r\cdot\vec{v}+s\cdot\vec{v} \)
\item scalar multiplication distributes over vector addition,
  \( r\cdot(\vec{v}+\vec{w})=r\cdot\vec{v}+r\cdot\vec{w} \)
\item ordinary multiplication of scalars associates with 
  scalar multiplication, \( (rs)\cdot\vec{v} =r\cdot(s\cdot\vec{v}) \)
\item multiplication by the scalar~$1$ is the 
  identity operation, \( 1\cdot\vec{v}=\vec{v} \).
\end{compactenum}
% \end{tfae}  
%</df:VectorSpace>
\end{definition}

\vspace{-4.5ex}  % why is the extra space there?  Bug in shading code?
\begin{remark}
The definition involves two kinds of addition and two kinds of multiplication,
and so may at first seem confused.
For instance, in condition~(7)
the `$+$' on the left is addition of two real numbers
while the `$+$' on the right is addition of two vectors in
\( V\/ \).
These expressions aren't ambiguous because of context; for example,
\( r \) and \( s \)
are real numbers so `\( r+s \)' can only mean real number addition.
In the same way, item~(9)'s left side `$rs$' is ordinary real number 
multiplication, while its right side `$s\cdot\vec{v}$' is 
the scalar multiplication defined for this
vector space.
\end{remark}

\vspace{-0.5ex}  % why is the extra space there?  Bug in shading code?
The best way to understand the definition is to
go through the examples below and for each,
check all ten conditions.
The first example includes that check, written out at length.
Use it as a model for the others.
Especially important are the \definend{closure}\index{vector space!closure}
conditions, (1) and~(6).
They specify that the addition and scalar multiplication operations
are always sensible\Dash they are defined for every pair of vectors 
and every scalar and vector,
and the result of the operation is a member of the set. 


\begin{example}  \label{PlaneThruOriginSubsp}
This subset of \( \Re^2 \) is a line through the origin.
\begin{equation*}
  L=\set{ \colvec{x \\ y}  \suchthat y=3x}
\end{equation*}
We shall verify that it is a
vector space under the usual meaning of `+' and `\(\cdot\)'.
\begin{equation*}
  \colvec{x_1 \\ y_1}
  +
  \colvec{x_2 \\ y_2}
  =
  \colvec{x_1+x_2 \\ y_1+y_2}
  \qquad
  r\cdot
  \colvec{x \\ y}
  =
  \colvec{rx \\ ry}
\end{equation*}
These operations
are just the ordinary ones, reused on its subset $L$.
We say that \( L \) \definend{inherits\/}\index{inherited operations} 
these operations from \( \Re^2 \).

We shall check all ten conditions.
The paragraph having to do with addition has five conditions.
For condition~(1), closure under addition, 
suppose that we start with two vectors from the line~$L$, 
\begin{equation*}
  \vec{v}_1=\colvec{x_1 \\ y_1}
  \quad
  \vec{v}_2=\colvec{x_2 \\ y_2}
\end{equation*}
so that they satisfy the restrictions
that $y_1=3x_1$ and~$y_2=3x_2$.
Their sum
\begin{equation*}
  \vec{v}_1+\vec{v}_2
  =\colvec{x_1+x_2 \\ y_1+y_2}
\end{equation*}
is also a member of the line $L$
because the fact that its second component is three times its first 
$y_1+y_2=3(x_1+x_2)$ follows from the 
restrictions on $\vec{v}_1$ and~$\vec{v}_2$.
For~(2), that addition of vectors commutes, 
just compare 
\begin{equation*}
  \vec{v}_1+\vec{v}_2
  =\colvec{x_1+x_2 \\ y_1+y_2}
  \quad
  \vec{v}_2+\vec{v}_1
  =\colvec{x_2+x_1 \\ y_2+y_1}
\end{equation*}
and note that they are equal since their entries are real numbers and 
real numbers commute. 
(That the vectors satisfy the restriction of lying in the line is not relevant
for this condition;
they commute just because all vectors in the plane commute.)
Condition~(3), associativity of vector addition, is similar.
\begin{align*}
  (\colvec{x_1 \\ y_1}
  +\colvec{x_2 \\ y_2})
  +\colvec{x_3 \\ y_3}
  &=\colvec{(x_1+x_2)+x_3 \\ (y_1+y_2)+y_3}  \\
  &=\colvec{x_1+(x_2+x_3) \\ y_1+(y_2+y_3)}  \\
  &=\colvec{x_1 \\ y_1}
  +(\colvec{x_2 \\ y_2}
  +\colvec{x_3 \\ y_3})
\end{align*}
For the fourth condition we must produce a vector that acts as the 
zero element. 
The vector of zero entries will do.
\begin{equation*}
  \colvec{x \\ y}
  +\colvec{0 \\ 0}
  =\colvec{x \\ y}
\end{equation*}
Note that $\zero\in L$ as its second component is triple its first.
For~(5), that given any~$\vec{v}\in L$ we can
produce an additive inverse, we have
\begin{equation*}
  \colvec{-x \\ -y}
  +\colvec{x \\ y}
  =\colvec{0 \\ 0}
\end{equation*}
and so the vector $-\vec{v}$ is the desired inverse.
As with the prior condition, observe here that if $\vec{v}\in L$, so that
$y=3x$, then $-\vec{v}\in L$ also, since $-y=3(-x)$.

The checks for the five conditions having to do with scalar multiplication
are similar.
For~(6), closure under scalar multiplication,
suppose that $r\in\R$ and $\vec{v}\in L$, that is, 
\begin{equation*}
  \vec{v}=\colvec{x \\ y}
\end{equation*}
satisfies that~$y=3x$.
Then 
\begin{equation*}
  r\cdot\vec{v}=r\cdot\colvec{x \\ y}
  =\colvec{rx \\ ry}
\end{equation*}
is also a member of $L$:~the relation $ry=3\cdot rx$ holds because
$y=3x$.
Next, this checks~(7).
\begin{equation*}
  (r+s)\cdot\colvec{x \\ y}
  =\colvec{(r+s)x \\ (r+s)y}
  =\colvec{rx+sx \\ ry+sy}
  =r\cdot\colvec{x \\ y}+s\cdot\colvec{x \\ y}
\end{equation*}
For~(8) we have this.
\begin{equation*}
  r\cdot(\colvec{x_1 \\ y_1}+\colvec{x_2 \\ y_2})
  =\colvec{r(x_1+x_2) \\ r(y_1+y_2)}
  =\colvec{rx_1+rx_2 \\ ry_1+ry_2}
  =r\cdot\colvec{x_1 \\ y_1}+r\cdot\colvec{x_2 \\ y_2}
\end{equation*}
The ninth
\begin{equation*}
  (rs)\cdot\colvec{x \\ y}
  =\colvec{(rs)x \\ (rs)y}
  =\colvec{r(sx) \\ r(sy)}
  =r\cdot(s\cdot\colvec{x \\ y})
\end{equation*}
and tenth conditions are also straightforward.
\begin{equation*}
  1\cdot\colvec{x \\ y}
  =\colvec{1x \\ 1y}
  =\colvec{x \\ y}
\end{equation*}
\end{example}

\begin{example}  \label{ex:RealVecSpaces}
The whole plane, the set
\( \Re^2 \), is a vector space where the operations `\( + \)' and `\( \cdot \)'
have their usual meaning.
\begin{equation*}
  \colvec{x_1 \\ y_1}
  +
  \colvec{x_2 \\ y_2}
  =
  \colvec{x_1+x_2 \\ y_1+y_2}
  \qquad
  r\cdot
  \colvec{x \\ y}
  =
  \colvec{rx \\ ry}
\end{equation*}
We shall check just two of the conditions, the closure conditions.

For (1) observe that 
the result of the vector sum
\begin{equation*}
  \colvec{x_1 \\ y_1}
  +\colvec{x_2 \\ y_2}
  =\colvec{x_1+x_2 \\ y_1+y_2}
\end{equation*}
is a column array with two real entries, and so is a member of the 
plane~\( \Re^2 \).
In contrast with the prior example, here there is no restriction on the
first and second components of the vectors.

Condition (6) is similar.
The vector
\begin{equation*}
  r\cdot\colvec{x \\ y}
  =\colvec{rx \\ ry}
\end{equation*}
has two real entries, and so is a member of~\( \Re^2 \).
\end{example}

In a similar way, 
each \( \Re^n \) is a vector space with the usual operations of vector addition
and scalar multiplication.
(In \( \Re^1 \), we usually do not write the members as
column vectors, i.e., we usually do not write `\( (\pi) \)'.
Instead we just write `\( \pi \)'.)

\begin{example} \label{ex:ColsIntEntNotVS}
\nearbyexample{PlaneThruOriginSubsp} gives a subset of $\Re^2$ that is
a vector space.
% \nearbyexample{ex:RealVecSpaces} shows that the entire
% set of two-tall vectors with real entries is also a vector space. 
For contrast, consider the set
of two-tall columns with entries that are integers,
under the same operations of component-wise 
addition and scalar multiplication.
This is a subset of $\Re^2$ but it is not a vector space:
it is not closed under scalar multiplication, 
that is, it does not satisfy condition~(6). 
For instance, on the left below is a vector with integer entries, and a scalar. 
\begin{equation*}
  0.5
  \cdot
  \colvec[r]{4 \\ 3}
  =
  \colvec[r]{2 \\ 1.5}
\end{equation*}
On the right is a column vector that 
is not a member of the set, since its entries are not all integers.
\end{example}

\begin{example}  \label{ex:TrivSbspReFour}
The one-element set
\begin{equation*}
  \set{ \colvec[r]{0 \\ 0 \\ 0 \\ 0} }
\end{equation*}
is a vector space under the operations 
\begin{equation*}
  \colvec[r]{0 \\ 0 \\ 0 \\ 0}
  +
  \colvec[r]{0 \\ 0 \\ 0 \\ 0}
  =
  \colvec[r]{0 \\ 0 \\ 0 \\ 0}
  \qquad
  r\cdot
  \colvec[r]{0 \\ 0 \\ 0 \\ 0}
  =
  \colvec[r]{0 \\ 0 \\ 0 \\ 0}
\end{equation*}
that it inherits from \( \Re^4 \).
\end{example}

A vector space must have at least one element, its zero vector.
Thus a one-element vector space is the smallest possible.

\begin{definition} \label{df:TrivialVectorSpace}
%<*df:TrivialVectorSpace>
A one-element vector space is a \definend{trivial}\index{vector space!trivial}%
\index{trivial space}
space.
%</df:TrivialVectorSpace>
\end{definition}

The examples so far involve sets of column vectors with the usual operations.
But vector spaces need not be collections of column vectors, or even of row
vectors.
Below are some other types of vector spaces.
The term `vector space' does not mean `collection of columns of reals'.
It means something more like 
`collection in which any linear combination is sensible'.

\begin{example} \label{ex:PolySpaceThree}
Consider
\( \polyspace_3=\set{a_0+a_1x+a_2x^2+a_3x^3\suchthat a_0,\ldots,a_3\in\Re} \),
the set of polynomials of degree three or less
(in this book, we'll take constant polynomials, 
including the zero polynomial, to be of degree zero).
It is a vector space under the operations
\begin{multline*}
   (a_0+a_1x+a_2x^2+a_3x^3)+(b_0+b_1x+b_2x^2+b_3x^3)  \\
     =(a_0+b_0)+(a_1+b_1)x+(a_2+b_2)x^2+(a_3+b_3)x^3
\end{multline*}
and
\begin{equation*}
   r\cdot(a_0+a_1x+a_2x^2+a_3x^3)=
     (ra_0)+(ra_1)x+(ra_2)x^2+(ra_3)x^3
\end{equation*}
(the verification is easy).
This vector space is worthy of attention because
these are the polynomial operations familiar from high school algebra.
For instance,
$
  3\cdot(1-2x+3x^2-4x^3)-2\cdot(2-3x+x^2-(1/2)x^3)=-1+7x^2-11x^3$.

Although this space is not a subset of any \( \Re^n \),
there is a sense in which we can think of $\polyspace_3$ as ``the same'' as  
\( \Re^4 \).
If we identify these two space's elements in this way 
\begin{equation*}
  a_0+a_1x+a_2x^2+a_3x^3
  \quad\text{corresponds to}\quad
  \colvec{a_0 \\ a_1 \\ a_2 \\ a_3}
\end{equation*}
then the operations also correspond.
Here is an example of corresponding additions.
\begin{equation*}
  \mbox{\begin{tabular}{lr}
       &\(1-2x+0x^2+1x^3\) \\
     + &\(2+3x+7x^2-4x^3\) \\ \hline
       &\(3+1x+7x^2-3x^3\)
  \end{tabular}  }
  \quad\text{corresponds to}\quad
  \colvec[r]{1 \\ -2 \\ 0 \\ 1}
  +
  \colvec[r]{2 \\ 3 \\ 7 \\ -4}
  =
  \colvec[r]{3 \\ 1 \\ 7 \\ -3}
\end{equation*}
Things we are thinking of as ``the same'' add to ``the same'' sum.
Chapter Three makes precise this idea of vector space correspondence.
For now we shall just leave it as an intuition.
\end{example}

In general we write
\( \polyspace_n \) for the
vector space of polynomials of degree~$n$ or less 
\( \set{a_0+a_1x+a_2x^2+\cdots+a_nx^n\suchthat a_0,\ldots,a_n\in\Re} \),
under the operations of the usual polynomial addition and scalar 
multiplication.\index{vector space!of polynomials}
We will often use these spaces as examples.

\begin{example} \label{ex:MatSpaceTwoByTwo}
The set \( \matspace_{\nbym{2}{2}} \) of \( \nbym{2}{2} \) matrices with
real number entries is a vector space under the natural 
entry-by-entry operations.
\begin{equation*}
  \begin{mat}
    a  &b \\
    c  &d
  \end{mat}
  +
  \begin{mat}
    w  &x \\
    y  &z
  \end{mat}
  =
  \begin{mat}
    a+w  &b+x \\
    c+y  &d+z
  \end{mat}
  \qquad
  r\cdot
  \begin{mat}
    a  &b \\
    c  &d
  \end{mat}
  =
  \begin{mat}
    ra  &rb \\
    rc  &rd
  \end{mat}
\end{equation*}
As in the prior example, we can think of this space as
``the same'' as \( \Re^4 \).
\end{example}

We write
\( \matspace_{\nbym{n}{m}} \) for the
vector space of $\nbym{n}{m}$ matrices
under the natural operations of matrix addition and scalar 
multiplication.\index{vector space!of matrices}
As with the polynomial spaces, 
we will often use these as examples.

\begin{example}   \label{ex:FcnsNToRIsVecSp}
The set \( \set{f\suchthat \map{f}{\N}{\Re} } \) of all
real-valued functions of one natural number variable is a vector space
under the operations
\begin{equation*}
  (f_1+f_2)\,(n)=f_1(n)+f_2(n)
  \qquad
  (r\cdot f)\,(n)=r\,f(n)
\end{equation*}
so that if, for example, \( f_1(n)=n^2+2\sin(n) \) and
\( f_2(n)=-\sin(n)+0.5 \)
then \( (f_1+2f_2)\,(n)=n^2+1 \).

We can view this space
as a generalization of \nearbyexample{ex:RealVecSpaces}\Dash instead of 
$2$-tall vectors, these functions are like infinitely-tall
vectors.
\begin{equation*}
  \text{
    \begin{tabular}{c|c}
      \( n \)      &\( f(n)=n^2+1 \)  \\ \hline
      \( 0      \) &\( 1      \)    \\
      \( 1      \) &\( 2      \)    \\
      \( 2      \) &\( 5      \)    \\
      $3$          &$10$            \\
      \( \vdots \) &\( \vdots \)    
    \end{tabular} }
    \quad\text{corresponds to}\quad
    \colvec[r]{
           1           \\
           2           \\
           5           \\
           10          \\
           \vdotswithin{10}      }
\end{equation*}
Addition and scalar multiplication are component-wise, 
as in \nearbyexample{ex:RealVecSpaces}.
(We can formalize ``infinitely-tall'' by saying that it means an infinite
sequence, or that it means a function from $\N$ to $\Re$.)
\end{example}

\begin{example} \label{ex:PolysOfAllFiniteDegrees}
The set of polynomials with real coefficients
\begin{equation*}
 \set{ a_0+a_1x+\cdots+a_nx^n\suchthat n\in\N
    \text{ and } a_0,\ldots,a_n\in\Re} 
\end{equation*}
makes a vector space when given the natural `$+$' 
\begin{multline*}
  (a_0+a_1x+\cdots+a_nx^n)+(b_0+b_1x+\cdots+b_nx^n)  \\
     =(a_0+b_0)+(a_1+b_1)x+\cdots +(a_n+b_n)x^n
\end{multline*}
and `$\cdot$'.
\begin{equation*}
  r\cdot (a_0+a_1x+\ldots a_nx^n)
   =
  (ra_0)+(ra_1)x+\ldots (ra_n)x^n
\end{equation*}
This space differs from the space $\polyspace_3$ of
\nearbyexample{ex:PolySpaceThree}.
This space contains not just degree~three polynomials, 
but degree~thirty polynomials and
degree three~hundred polynomials, too.
Each individual polynomial of course is of a finite degree, 
but the set has no single bound on the degree of all of its members.

We can think of this example, like the prior one,  
in terms of infinite-tuples.
For instance, we can think of \( 1+3x+5x^2 \) as corresponding to
\( (1,3,5,0,0,\ldots) \).
However, this space differs from the one in
\nearbyexample{ex:FcnsNToRIsVecSp}.
Here, each member of the set has a finite degree, that is,
under the correspondence there is no element from this space 
matching \( (1,2,5,10,\,\ldots\,) \).
Vectors in this space correspond to infinite-tuples
that end in zeroes.
\end{example}

\begin{example}  \label{ex:RealValuedFcns}
The set
\( \set{f\suchthat \map{f}{\Re}{\Re} } \)
of all real-valued functions of one real variable
is a vector space under these.
\begin{equation*}
  (f_1+f_2)\,(x)=f_1(x)+f_2(x)
  \qquad
  (r\cdot f)\,(x)=r\,f(x)
\end{equation*}
The difference between this and \nearbyexample{ex:FcnsNToRIsVecSp} is the
domain of the functions.
\end{example}

\begin{example}\label{ex:ACos+BSin}
The set
\( F=\{ a\cos\theta+b\sin\theta \suchthat a,b\in\Re\} \)
of real-valued functions of the real variable \( \theta \)
is a vector space under the operations
%\nearbyexample{ex:RealValuedFcns}:
\begin{equation*}
  (a_1\cos\theta+b_1\sin\theta)+(a_2\cos\theta+b_2\sin\theta)
    =(a_1+a_2)\cos\theta+(b_1+b_2)\sin\theta
\end{equation*}
and
\begin{equation*}
  r\cdot (a\cos\theta+b\sin\theta)
   =(ra)\cos\theta+(rb)\sin\theta
\end{equation*}
inherited from the space in the prior example.
(We can think of \( F \) as ``the same'' as \( \Re^2 \)
in that $a\cos\theta+b\sin\theta$ corresponds to the vector with
components $a$ and $b$.)
\end{example}

\begin{example}  \label{ex:SpaceSatisfyingDiffQ}
The set
\begin{equation*}
  \set{\map{f}{\Re}{\Re}\suchthat \dfrac{d^2f}{dx^2}+f=0}
\end{equation*}
is a vector space under the, by now natural, interpretation.
\begin{equation*}
  (f+g)\,(x)=f(x)+g(x)
  \qquad
  (r\cdot f)\,(x)=r\,f(x)
\end{equation*}
In particular, notice that basic Calculus gives
\begin{equation*}
   \frac{d^2(f+g)}{dx^2}+(f+g)
   =(\frac{d^2f}{dx^2}+f)+(\frac{d^2g}{dx^2}+g)
\end{equation*}
and
\begin{equation*}
   \frac{d^2(rf)}{dx^2}+(rf)
   =r(\frac{d^2 f}{dx^2}+f)
\end{equation*}
and so the space is closed under addition and scalar multiplication.
This turns out to equal the space from the prior example\Dash functions
satisfying this differential equation have the form
$a\cos\theta+b\sin\theta$\Dash but this description 
suggests an extension to solutions sets of other
differential equations.
\end{example}

\begin{example} \label{ex:HomoSlnMakesVS}
The set of solutions of a homogeneous linear system in \( n \) variables
is a vector space under the operations inherited from \( \Re^n \).
For example, for closure under addition 
consider a typical equation in that system
$c_1x_1+\cdots+c_nx_n=0$ and suppose that both these vectors
\begin{equation*}
   \vec{v}=\colvec{v_1 \\ \vdotswithin{v_1} \\ v_n}
   \qquad
   \vec{w}=\colvec{w_1 \\ \vdotswithin{w_1} \\ w_n}
\end{equation*}
satisfy the equation. 
Then their sum
\( \vec{v}+\vec{w}\/ \) also satisfies that equation:
\(
  c_1(v_1+w_1)+\cdots+c_n(v_n+w_n)
  =(c_1v_1+\cdots+c_nv_n)+(c_1w_1+\cdots+c_nw_n)
  =0
\).
The checks of the other vector space conditions are just as routine.
\end{example}

We often omit the multiplication symbol `\( \cdot \)' between the 
scalar and the vector. 
We distinguish the multiplication in
\( c_1v_1 \) from that in \( r\vec{v}\, \) by context, since if both 
multiplicands are real numbers then it must be 
real-real multiplication while if one is a vector then it must be
scalar-vector multiplication.

\nearbyexample{ex:HomoSlnMakesVS} has brought us full circle since it is one of
our motivating examples.
Now, with some feel for the kinds of structures that satisfy the definition
of a vector space, we can reflect on that definition.
For example, why specify in the definition the condition that 
\( 1\cdot\vec{v}=\vec{v} \) but not a condition that \( 0\cdot\vec{v}=\zero \)?

One answer is that this is just a definition\Dash it gives the rules 
and you need to follow those rules to continue.

Another answer is perhaps more satisfying.
People in this area have worked to develop the
right balance of power and generality.
This definition is shaped so that it contains the conditions
needed to prove all of the interesting and
important properties of spaces of linear combinations.
As we proceed, we shall derive all of the properties natural to collections of
linear combinations from the conditions given in the definition.

The next result is an example.
We do not need to include these properties in the definition of vector space
because they follow from the properties already listed there.

\begin{lemma} \label{lm:ElementaryPropertiesOfVectorSpaces}
%<*lm:ElementaryPropertiesOfVectorSpaces>
In any vector space \( V \), 
for any \( \vec{v}\in V \) and \( r\in\Re \), we have
(1)~\( 0\cdot\vec{v}=\zero \), 
(2)~\( (-1\cdot\vec{v})+\vec{v}=\zero \), and 
(3)~\( r\cdot\zero=\zero \).
%</lm:ElementaryPropertiesOfVectorSpaces>
\end{lemma}

\begin{proof}
%<*pf:ElementaryPropertiesOfVectorSpaces0>
For (1) note that 
\( \vec{v}=(1+0)\cdot\vec{v}=\vec{v}+(0\cdot\vec{v}) \).
Add to both sides the additive inverse of \( \vec{v} \),
the vector \( \vec{w} \) such that \( \vec{w}+\vec{v}=\zero \).
\begin{align*}
  \vec{w}+\vec{v}
  &=\vec{w}+\vec{v}+0\cdot\vec{v}  \\
  \zero
  &=\zero+0\cdot\vec{v}                   \\
  \zero
  &=0\cdot\vec{v}
\end{align*}
Item~(2) is easy:
\(  (-1\cdot\vec{v})+\vec{v}=(-1+1)\cdot\vec{v}=0\cdot\vec{v}=\zero \).
For (3),
\( r\cdot\zero
  =
  r\cdot(0\cdot\zero)
  =
  (r\cdot 0)\cdot\zero
  =
  \zero \)
will do.
%</pf:ElementaryPropertiesOfVectorSpaces0>
\end{proof}

The second item
shows that we can write the additive inverse
of \( \vec{v} \) as `\( -\vec{v}\, \)' without worrying about any 
confusion with \( (-1)\cdot\vec{v} \).

\medskip
A recap:
our study in Chapter One of Gaussian reduction
led us to consider collections of linear combinations.
So in this chapter we have defined a vector space to be a 
structure in which we can form such combinations,
% expressions of the form \( c_1\cdot\vec{v}_1+\dots+c_n\cdot\vec{v}_n \)
subject to simple conditions on the addition and scalar
multiplication operations.
In a phrase:~vector spaces are
the right context in which to study linearity.

% Finally, a comment.
From the fact that it forms a whole chapter, and especially because that 
chapter is the first one, a reader could suppose that our purpose
in this book is the study of linear systems.
The truth is that we will not so much use vector spaces in
the study  of linear systems as we instead have linear systems 
start us on the study of vector spaces.
The wide variety of examples from this subsection shows that the study of
vector spaces is interesting and important in its own right.
Linear systems won't go away.
But from now on our primary objects of study will be vector spaces.

\begin{exercises}
  \item 
    Name the zero vector for each of these vector spaces.
    \begin{exparts}
      \partsitem The space of degree three polynomials under the natural
        operations.
      \partsitem The space of \( \nbym{2}{4} \) matrices.
      \partsitem The space
        \( \set{\map{f}{[0..1]}{\Re}\suchthat f\text{ is continuous}} \).
      \partsitem The space of real-valued functions of one natural 
        number variable.
    \end{exparts}
    \begin{answer}
      \begin{exparts}
        \partsitem \( 0+0x+0x^2+0x^3 \)
        \partsitem \( \begin{mat}[r]
                   0  &0  &0  &0  \\
                   0  &0  &0  &0
                 \end{mat} \)
        \partsitem The constant function \( f(x)=0 \)
        \partsitem The constant function \( f(n)=0 \)
      \end{exparts}  
    \end{answer}
  \recommended \item
    Find the additive inverse, in the vector space,
    of the vector.
    \begin{exparts}
      \partsitem In \( \polyspace_3 \), the vector \( -3-2x+x^2 \).
      \partsitem In the space \( \nbyn{2} \),
        \begin{equation*}
          \begin{mat}[r]
            1  &-1  \\
            0  &3
          \end{mat}.
        \end{equation*}
     \partsitem In \( \set{ae^x+be^{-x}\suchthat a,b\in\Re} \), the space 
       of functions of the real variable \( x \) under the natural operations,
       the vector \( 3e^x-2e^{-x} \).
    \end{exparts}
    \begin{answer}
      \begin{exparts*}
        \partsitem \( 3+2x-x^2 \)
        \partsitem \( \begin{mat}[r]
                   -1  &+1  \\
                    0  &-3
                 \end{mat} \)
        \partsitem \( -3e^x+2e^{-x} \)
      \end{exparts*}  
    \end{answer}
  \recommended \item 
    For each, list three elements and then show it is a vector space.
    \begin{exparts}
      \partsitem The set of linear polynomials
        \( \polyspace_1=\set{a_0+a_1x\suchthat a_0,a_1\in\Re} \) under the
        usual polynomial addition and scalar multiplication operations.
      \partsitem The set of linear polynomials
        \( \set{a_0+a_1x\suchthat a_0-2a_1=0} \), under the
        usual polynomial addition and scalar multiplication operations.
   \end{exparts}
   \textit{Hint.}  Use \nearbyexample{PlaneThruOriginSubsp} as a guide.
   Most of the ten conditions are just verifications. 
   \begin{answer}
      \begin{exparts}
        \partsitem 
          Three elements are: 
          $1+2x$, $2-1x$, and $x$.
          (Of course, many answers are possible.)

          The verification is just like \nearbyexample{ex:RealVecSpaces}. 
          We first do 
          conditions $1$-$5$ from 
          \nearbydefinition{def:VecSpace}, having to 
          do with addition.
          For closure under addition, condition~(1), 
          note that where $a+bx,c+dx\in\polyspace_1$
          we have that $(a+bx)+(c+dx)=(a+c)+(b+d)x$ is a linear polynomial
          with real coefficients and so is an element of 
          $\polyspace_1$. 
          Condition~(2) is verified with: where $a+bx,c+dx\in\polyspace_1$ then 
          $(a+bx)+(c+dx)=(a+c)+(b+d)x$, while in the other order they are
          $(c+dx)+(a+bx)=(c+a)+(d+b)x$, and both $a+c=c+a$ 
          and $b+d=d+b$ as these are real numbers.
          Condition~(3) is similar: suppose
          $a+bx,c+dx,e+fx\in\polyspace$ then
          $((a+bx)+(c+dx))+(e+fx)=(a+c+e)+(b+d+f)x$
          while
          $(a+bx)+((c+dx)+(e+fx))=(a+c+e)+(b+d+f)x$, and the two are equal
          (that is, real number addition is associative so $(a+c)+e=a+(c+e)$
          and $(b+d)+f=b+(d+f)$).
          For condition~(4) observe that the linear polynomial
          $0+0x\in\polyspace_1$ has the 
          property that $(a+bx)+(0+0x)=a+bx$ and 
          $(0+0x)+(a+bx)=a+bx$.
          For the last condition in this paragraph, condition~(5),
          note that for any $a+bx\in\polyspace_1$ the additive inverse
          is $-a-bx\in\polyspace_1$ since
          $(a+bx)+(-a-bx)=(-a-bx)+(a+bx)=0+0x$.
         
          We next also check conditions~(6)-(10), involving 
          scalar multiplication.
          For (6), the condition that the space be closed under scalar
          multiplication, 
          suppose that $r$ is a real number and $a+bx$ is an element of
          $\polyspace_1$, and then $r(a+bx)=(ra)+(rb)x$ is an element of 
          $\polyspace_1$ because it is a linear polynomial with real
          number coefficients.
          Condition~(7) holds because 
          $(r+s)(a+bx)=r(a+bx)+s(a+bx)$ is true from the distributive property
          for real number multiplication.
          Condition~(8) is similar: 
          $r((a+bx)+(c+dx))=r((a+c)+(b+d)x)=r(a+c)+r(b+d)x=(ra+rc)+(rb+rd)x
          =r(a+bx)+r(c+dx)$.
          For~(9) we have 
          $(rs)(a+bx)=(rsa)+(rsb)x=r(sa+sbx)=r(s(a+bx))$.
          Finally, condition~(10) is 
          $1(a+bx)=(1a)+(1b)x=a+bx$.
        \partsitem 
          Call the set $P$.
          In the prior item in this exercise 
          there was no restriction on the coefficients
          but here we are restricting attention to those linear polynomials
          where $a_0-2a_1=0$, that is, 
          where the constant term minus twice the coefficient of the linear
          term is zero. 
          Thus, three typical elements of $P$ are
          $2+1x$, $6+3x$, and $-4-2x$.

          For condition~(1) we must show that
          if we add two linear polynomials that satisfy the restriction then
          we get a linear polynomial also satisfying the restriction: here that
          argument is that
          if $a+bx,c+dx\in P$ then
          $(a+bx)+(c+dx)=(a+c)+(b+d)x$ is an element of $P$ because
          $(a+c)-2(b+d)=(a-2b)+(c-2d)=0+0=0$. 
          We can verify condition~(2) with: 
          where $a+bx,c+dx\in\polyspace_1$ then 
          $(a+bx)+(c+dx)=(a+c)+(b+d)x$, while in the other order they are
          $(c+dx)+(a+bx)=(c+a)+(d+b)x$, and both $a+c=c+a$ 
          and $b+d=d+b$ as these are real numbers. 
          (That is, this condition
          is not affected by the restriction and the verification is the same 
          as the verification in the first item of this exercise).
          Condition~(3) is also not affected by the extra restriction: 
          suppose that
          $a+bx,c+dx,e+fx\in\polyspace$ then
          $((a+bx)+(c+dx))+(e+fx)=(a+c+e)+(b+d+f)x$
          while
          $(a+bx)+((c+dx)+(e+fx))=(a+c+e)+(b+d+f)x$, and the two are equal.
          For condition~(4) observe that the linear polynomial
          satisfies the restriction $0+0x\in P$ because
          its constant term minus twice the coefficient of its linear term
          is zero, and then the verification from the first item of this 
          question applies:
          $0+0x\in\polyspace_1$ has the 
          property that $(a+bx)+(0+0x)=a+bx$ and 
          $(0+0x)+(a+bx)=a+bx$.
          To check condition~(5),
          note that for any $a+bx\in P$ the additive inverse
          is $-a-bx$ since
          it is an element of $P$ (because $a+bx\in P$ we know that 
          $a-2b=0$ and multiplying both sides by $-1$ gives that $-a+2b=0$), 
          and as in the first item it acts as the additive inverse
          $(a+bx)+(-a-bx)=(-a-bx)+(a+bx)=0+0x$.
         
          We must also check conditions~(6)-(10), those for 
          scalar multiplication.
          For (6), the condition that the space be closed under scalar
          multiplication, 
          suppose that $r$ is a real number and $a+bx\in P$ (so that $a-2b=0$),
          then $r(a+bx)=(ra)+(rb)x$ is an element of 
          $P$ because it is a linear polynomial with real
          number coefficients satisfying that $(ra)-2(rb)=r(a-2b)=0$.
          Condition~(7) holds for the same reason that it holds in 
          the first item of this exercise, because 
          $(r+s)(a+bx)=r(a+bx)+s(a+bx)$ is true from the distributive property
          for real number multiplication.
          Condition~(8) is also unchanged from the first item: 
          $r((a+bx)+(c+dx))=r((a+c)+(b+d)x)=r(a+c)+r(b+d)x=(ra+rc)+(rb+rd)x
          =r(a+bx)+r(c+dx)$.
          So is~(9): 
          $(rs)(a+bx)=(rsa)+(rsb)x=r(sa+sbx)=r(s(a+bx))$.
          Finally, so is condition~(10): 
          $1(a+bx)=(1a)+(1b)x=a+bx$.
       \end{exparts}
     \end{answer}
  \item 
    For each, list three elements and then show it is a vector space.
    \begin{exparts}
      \partsitem The set of \( \nbyn{2} \) matrices with real entries under
        the usual matrix operations.
      \partsitem The set of \( \nbyn{2} \) matrices with real entries where
        the $2,1$ entry is zero,  under
        the usual matrix operations.
    \end{exparts}
    \begin{answer}
      Use
      \nearbyexample{ex:RealVecSpaces} as a guide.
      (\textit{Comment.} 
      Because many of the conditions are quite easy to check,
      sometimes a person can feel that they must have missed
      something.
      Keep in mind that 
      easy to do, or routine, is different from not necessary to do.)
      \begin{exparts}
        \partsitem 
          Here are three elements.
          \begin{equation*}
            \begin{mat}
              1  &2  \\
              3  &4
            \end{mat},\,
            \begin{mat}
              -1  &-2  \\
              -3  &-4
            \end{mat},\,
            \begin{mat}
              0  &0  \\
              0  &0
            \end{mat}
          \end{equation*}

          For~(1), the sum of $\nbyn{2}$ real matrices is a $\nbyn{2}$
          real matrix.
          For~(2) we consider the sum of two matrices 
          \begin{equation*}
            \begin{mat}
              a  &b  \\
              c  &d
            \end{mat}+
            \begin{mat}
              e  &f  \\
              g  &h
            \end{mat}
            =
            \begin{mat}
              a+e  &b+f  \\
              c+g  &d+h
            \end{mat}
          \end{equation*}
          and apply commutativity of real number addition
          \begin{equation*}
            =\begin{mat}
              e+a  &f+b  \\
              g+c  &h+d
            \end{mat}
            =
            \begin{mat}
              e  &f  \\
              g  &h
            \end{mat}
            +
            \begin{mat}
              a  &b  \\
              c  &d
            \end{mat}
          \end{equation*}
          to verify that the addition of the matrices is commutative.
          The verification for condition~(3), 
          associativity of matrix addition, is similar
          to the prior verification:
          \begin{equation*}
            \big(
              \begin{mat}
                a  &b  \\
                c  &d  
              \end{mat}
              +\begin{mat}
                e  &f  \\
                g  &h  
              \end{mat}
            \big)
            +
            \begin{mat}
              i  &j  \\
              k  &l
            \end{mat}
            =
            \begin{mat}
              (a+e)+i  &(b+f)+j  \\
              (c+g)+k  &(d+h)+l
            \end{mat}
          \end{equation*}
          while
          \begin{equation*}
            \begin{mat}
              a  &b  \\
              c  &d  
            \end{mat}
            +
              \big(
              \begin{mat}
                e  &f  \\
                g  &h  
              \end{mat}
              +
              \begin{mat}
                i  &j  \\
                k  &l
              \end{mat}
            \big)
            =
            \begin{mat}
              a+(e+i)  &b+(f+j)  \\
              c+(g+k)  &d+(h+l)
            \end{mat}
          \end{equation*}
          and the two are the same entry-by-entry because real number addition
          is associative.
          For~(4), the zero element of this space is the $\nbyn{2}$ 
          matrix of zeroes.
          Condition~(5) holds because for any $\nbyn{2}$ matrix~$A$
          the additive inverse is the matrix whose entries are the negative of
          $A$'s, the matrix $-1\cdot A$.

          Condition~($6$) 
          holds because a scalar multiple of a $\nbyn{2}$ matrix 
          is a $\nbyn{2}$ matrix.
          For condition~(7) we have this.
          \begin{multline*}
            (r+s)
            \begin{mat}
              a  &b  \\
              c  &d
            \end{mat}
            =
            \begin{mat}
              (r+s)a  &(r+s)b  \\
              (r+s)c  &(r+s)d
            \end{mat}               \\
            =
            \begin{mat}
              ra+sa  &rb+sb  \\
              rc+sc  &rd+sd
            \end{mat}
            =
            r
            \begin{mat}
              a  &b  \\
              c  &d 
            \end{mat}
            +
            s
            \begin{mat}
              a  &b  \\
              c  &d 
            \end{mat}
          \end{multline*}
          Condition~(8) goes the same way.
          \begin{multline*}
            r
            \big(
              \begin{mat}
                a  &b  \\
                c  &d
              \end{mat}
              +
              \begin{mat}
                e  &f  \\
                g  &h
              \end{mat}
            \big)
            =
            r
            \begin{mat}
              a+e  &b+f  \\
              c+g  &d+h
            \end{mat}
            =
            \begin{mat}
              ra+re  &rb+rf  \\
              rc+rg  &rd+rh
            \end{mat}              \\
            =
            r
            \begin{mat}
              a  &b  \\
              c  &d 
            \end{mat}
            +
            r
            \begin{mat}
              e  &f  \\
              g  &h 
            \end{mat}
            =
            r
            \big(
            \begin{mat}
              a  &b  \\
              c  &d
            \end{mat}
            +
            \begin{mat}
              e  &f  \\
              g  &h
            \end{mat}
            \big)
          \end{multline*}
          For~(9) we have this.
          \begin{equation*}
            (rs)
            \begin{mat}
              a  &b  \\
              c  &d
            \end{mat}
            =
            \begin{mat}
              rsa  &rsb  \\
              rsc  &rsd
            \end{mat}
            =
            r
            \begin{mat}
              sa  &sb  \\
              sc  &sd
            \end{mat}
            =
            r\big( s
            \begin{mat}
              a  &b  \\
              c  &d
            \end{mat}
            \big)
          \end{equation*}
          Condition~(10) is just as easy.
          \begin{equation*}
            1
            \begin{mat}
              a  &b  \\
              c  &d
            \end{mat}
            =
            \begin{mat}
              1\cdot a  &1\cdot b  \\
              1\cdot c  &1\cdot d
            \end{mat}
            =
            \begin{mat}
              sa  &sb  \\
              sc  &sd
            \end{mat}
          \end{equation*}
        \partsitem 
          This differs from the prior item in this exercise only in that
          we are restricting to the set $T$ of matrices with a zero in the
          second row and first column. 
          Here are three elements of $T$.
          \begin{equation*}
            \begin{mat}
              1  &2  \\
              0  &4
            \end{mat},\,
            \begin{mat}
              -1  &-2  \\
              0  &-4
            \end{mat},\,
            \begin{mat}
              0  &0  \\
              0  &0
            \end{mat}
          \end{equation*}
          Some of the verifications for this item are the same as for the 
          first item in this exercise, and below we'll just do the ones that
          are different. 

          For~(1), the sum of $\nbyn{2}$ real matrices with a zero in the
          $2,1$ entry is also a $\nbyn{2}$ real matrix with a zero in the
          $2,1$ entry.
          \begin{equation*}
            \begin{mat}
              a  &b  \\
              0  &d
            \end{mat}
            +
            \begin{mat}
              e  &f  \\
              0  &h
            \end{mat}
            \begin{mat}
              a+e  &b+f  \\
              0    &d+h
            \end{mat}
          \end{equation*}
          The verification for condition~(2) given in the prior item
          works in this item also.
          The same holds for condition~(3).
          For~(4), note that the $\nbyn{2}$ 
          matrix of zeroes is an element of $T$.
          Condition~(5) holds because for any $\nbyn{2}$ matrix~$A$
          the additive inverse is the matrix 
          $-1\cdot A$ and so the additive inverse of a 
          matrix with a zero in the $2,1$ entry is also a matrix with a zero
          in the $2,1$ entry.

          Condition~$6$ holds because a scalar multiple of a $\nbyn{2}$ matrix 
          with a zero in the $2,1$ entry is 
          a $\nbyn{2}$ matrix with a zero in the $2,1$ entry.
          Condition~(7)'s verification is the same as in the prior item.
          So are condition~(8)'s, (9)'s, and~(10)'s.
      \end{exparts}        
    \end{answer}
  \recommended \item 
    For each, list three elements and then show it is a vector space.
    \begin{exparts}
      \partsitem The set of three-component row vectors with their usual
        operations.
      \partsitem The set
        \begin{equation*}
          \set{\colvec{x \\ y \\ z \\ w}\in\Re^4\suchthat x+y-z+w=0}
        \end{equation*}
        under the operations inherited from $\Re^4$.
    \end{exparts}
    \begin{answer}
      % Most of the conditions are easy to check; use
      % \nearbyexample{ex:RealVecSpaces} as a guide.
      \begin{exparts}
        \partsitem 
          Three elements are $\rowvec{1  &2  &3}$, 
          $\rowvec{2  &1  &3}$, and $\rowvec{0  &0  &0}$.

          We must check conditions (1)-(10) in \nearbydefinition{def:VecSpace}.
          Conditions (1)-(5) concern addition.
          For condition~(1) recall that the sum of two three-component 
          row vectors 
          \begin{equation*}
            \rowvec{a  &b  &c}
            +\rowvec{d  &e  &f}
            =\rowvec{a+d  &b+e  &c+f}
          \end{equation*}
          is also a three-component row vector
         (all of the letters $a,\ldots,f$ represent real numbers).
         Verification of~(2) is routine
          \begin{multline*}
            \rowvec{a  &b  &c}
            +\rowvec{d  &e  &f}
            =\rowvec{a+d  &b+e  &c+f}   \\
            =\rowvec{d+a  &e+b  &f+c}
            =\rowvec{d  &e  &f}
            +\rowvec{a  &b  &c}
          \end{multline*}
          (the second equality holds because the three entries are real
          numbers and real number addition commutes).
         Condition~(3)'s verification is similar.
          \begin{multline*}
            \big(\rowvec{a  &b  &c}
              +\rowvec{d  &e  &f}
            \big)
            +\rowvec{g  &h  &i}
            =\rowvec{(a+d)+g  &(b+e)+h  &(c+f)+i}    \\
            =\rowvec{a+(d+g)  &b+(e+h)  &c+(f+i)}
            =\rowvec{a  &b  &c}
            +\big(\rowvec{d  &e  &f}
               +\rowvec{g  &h  &i}
             \big) 
          \end{multline*}
          For (4), observe that the three-component row vector
          $\rowvec{0  &0  &0}$ is the additive identity:
          $\rowvec{a  &b  &c}+\rowvec{0  &0  &0}=\rowvec{a  &b  &c}$.
          To verify condition~(5), assume we are given the 
          element $\rowvec{a  &b  &c}$ of the set and note that
          $\rowvec{-a  &-b  &-c}$ is also in the set and has the
          desired property: 
          $\rowvec{a  &b  &c}+\rowvec{-a  &-b  &-c}=\rowvec{0  &0  &0}$.

          Conditions (6)-(10) involve scalar multiplication.
          To verify~(6), that the space is closed under the 
          scalar multiplication operation that was given, 
          note that $r\rowvec{a  &b  &c}=\rowvec{ra  &rb  &rc}$ is a 
          three-component row vector with real entries.
          For~(7) we compute.
          \begin{multline*}
            (r+s)\rowvec{a  &b  &c}=\rowvec{(r+s)a  &(r+s)b  &(r+s)c}
            =\rowvec{ra+sa  &rb+sb  &rc+sc}                             \\
            =\rowvec{ra  &rb  &rc}+\rowvec{sa  &sb  &sc}
            =r\rowvec{a  &b  &c}+s\rowvec{a  &b  &c}
          \end{multline*}
          Condition~(8) is very similar.
          \begin{multline*} 
            r\big(\rowvec{a  &b  &c}+\rowvec{d  &e  &f}\big)  \\
            \begin{aligned}
            &=r\rowvec{a+d  &b+e  &c+f}
            =\rowvec{r(a+d)  &r(b+e)  &r(c+f)}                     \\
            &=\rowvec{ra+rd  &rb+re  &rc+rf}
            =\rowvec{ra  &rb  &rc}+\rowvec{rd  &re  &rf}          \\
            &=r\rowvec{a  &b  &c}+r\rowvec{d  &e  &f}
            \end{aligned}
          \end{multline*}
          So is the computation for condition~(9).
          \begin{equation*}
            (rs)\rowvec{a  &b  &c}
             =\rowvec{rsa  &rsb  &rsc}
             =r\rowvec{sa  &sb  &sc}
             =r\big(s\rowvec{a  &b  &c}\big)
          \end{equation*}
          Condition~(10) is just as routine
          $1\rowvec{a  &b  &c}
             =\rowvec{1\cdot a  &1\cdot b  &1\cdot c}
             =\rowvec{a  &b  &c}$.
        \partsitem 
          Call the set $L$. 
          Closure of addition, condition~(1), involves checking that if the
          summands are members of~$L$ then
          the sum
          \begin{equation*}
            \colvec{a \\ b \\ c \\ d}
            +\colvec{e \\ f \\ g \\ h}
            =
            \colvec{a+e \\ b+f \\ c+g \\ d+h}
          \end{equation*}
          is also a member of \( L \), which is true because it satisfies the 
          criteria for membership in $L$:
          \( (a+e)+(b+f)-(c+g)+(d+h)
          =(a+b-c+d)+(e+f-g+h)=0+0 \).
          The verifications for conditions~(2), (3), and~(5) are similar to the
          ones in the first part of this exercise.
          For condition~(4) note that the vector of zeroes is a member of~$L$
          because its first component plus its second, minus its third,
          and plus its fourth, totals to zero.

          Condition~(6), closure of scalar multiplication, is similar:
          where the vector is an element of~$L$,
          \begin{equation*}
            r\colvec{a  \\ b \\ c \\ d}
            =\colvec{ra  \\  rb  \\  rc  \\  rd}
          \end{equation*}
          is also an element of~$L$ because $ra+rb-rc+rd=r(a+b-c+d)=r\cdot 0=0$.
          The verification for conditions~(7), (8), (9), and~(10) are as in the
          prior item of this exercise.          
     \end{exparts}  
     \end{answer}
  \recommended \item \label{exer:NotVectorSpaces}
    Show that each of these is not a vector space.
    (\textit{Hint.}  Check closure by listing two members of each set
    and trying some operations on them.)     
    \begin{exparts}
      \partsitem Under the operations inherited from \( \Re^3 \), this set
        \begin{equation*}
          \set{\colvec{x \\ y \\ z}\in\Re^3\suchthat x+y+z=1}
        \end{equation*}
      \partsitem Under the operations inherited from \( \Re^3 \), this set
        \begin{equation*}
          \set{\colvec{x \\ y \\ z}\in\Re^3\suchthat x^2+y^2+z^2=1}
        \end{equation*}
      \partsitem Under the usual matrix operations,
        \begin{equation*}
          \set{\begin{mat}
                 a  &1  \\
                 b  &c
               \end{mat} \suchthat a,b,c\in\Re}
        \end{equation*}
      \partsitem Under the usual polynomial operations,
        \begin{equation*}
          \set{a_0+a_1x+a_2x^2\suchthat a_0,a_1,a_2\in\Re^+}
        \end{equation*}
        where $\Re^+$ is the set of reals greater than zero
      \partsitem Under the inherited operations,
        \begin{equation*}
          \set{\colvec{x \\ y}\in\Re^2\suchthat
               \text{\( x+3y=4 \) and \( 2x-y=3 \) and \( 6x+4y=10 \)} }
        \end{equation*}
    \end{exparts}
    \begin{answer}
      In each item the set is called \( Q \).
      For some items, there are other correct ways to show that $Q$ is not
      a vector space.
      \begin{exparts}
        \partsitem It is not closed under addition; it fails to meet
          condition~(1).
          \begin{equation*}
            \colvec[r]{1 \\ 0 \\ 0},
            \colvec[r]{0 \\ 1 \\ 0}\in Q
            \qquad
            \colvec[r]{1 \\ 1 \\ 0}\not\in Q
          \end{equation*}
        \partsitem It is not closed under addition.
          \begin{equation*}
            \colvec[r]{1 \\ 0 \\ 0},
            \colvec[r]{0 \\ 1 \\ 0}\in Q
            \qquad
            \colvec[r]{1 \\ 1 \\ 0}\not\in Q
          \end{equation*}
        \partsitem It is not closed under addition.
          \begin{equation*}
            \begin{mat}[r]
              0  &1  \\
              0  &0
            \end{mat},
            \,
            \begin{mat}[r]
              1  &1  \\
              0  &0
            \end{mat}\in Q
            \qquad
            \begin{mat}[r]
              1  &2  \\
              0  &0
            \end{mat}\not\in Q
          \end{equation*}
        \partsitem It is not closed under scalar multiplication.
          \begin{equation*}
            1+1x+1x^2\in Q
            \qquad
            -1\cdot(1+1x+1x^2)\not\in Q
          \end{equation*}
        \item It is empty, violating condition~(4).
      \end{exparts}  
    \end{answer}
  \item 
    Define addition and scalar multiplication operations to 
    make the complex numbers a vector space over \( \Re \).
    \begin{answer}
      The usual operations
      \( (v_0+v_1i)+(w_0+w_1i)=(v_0+w_0)+(v_1+w_1)i \) and
      \( r(v_0+v_1i)=(rv_0)+(rv_1)i \) suffice.
      The check is easy.  
    \end{answer}
  \item
    Is the set of rational numbers a vector space over \( \Re \) under the
    usual addition and scalar multiplication operations?
    \begin{answer}
       No, it is not closed under scalar multiplication since, e.g., 
       \( \pi\cdot (1) \) is not a rational number. 
    \end{answer}
  \item 
    Show that 
    the set of linear combinations of the variables \( x,y,z \) is
    a vector space under the natural addition and scalar multiplication
    operations.
    \begin{answer}
      The natural operations are
      \( (v_1x+v_2y+v_3z)+(w_1x+w_2y+w_3z)=(v_1+w_1)x+(v_2+w_2)y+(v_3+w_3)z \)
      and \( r\cdot(v_1x+v_2y+v_3z)=(rv_1)x+(rv_2)y+(rv_3)z \).
      The check that this is a vector space is easy; use
      \nearbyexample{ex:RealVecSpaces} as a guide.  
    \end{answer}
  \item 
    Prove that 
    this is not a vector space: the set of two-tall column vectors
    with real entries subject to these operations.
    \begin{equation*}
      \colvec{x_1 \\ y_1}
      +\colvec{x_2 \\ y_2}
      =\colvec{x_1-x_2 \\ y_1-y_2}
      \qquad
      r\cdot\colvec{x \\ y}
      =\colvec{rx \\ ry}
    \end{equation*}
    \begin{answer}
      The `\( + \)' operation is not commutative (that is, condition~(2) is 
      not met); producing two members of the
      set witnessing this assertion is easy.
    \end{answer}
  \item 
    Prove or disprove that \( \Re^3 \) is a vector space under these
    operations.
    \begin{exparts}
      \partsitem \( 
               \colvec{x_1 \\ y_1 \\ z_1}
               +\colvec{x_2 \\ y_2 \\ z_2}
               =\colvec{0 \\ 0 \\ 0}
               \quad\text{and}\quad
               r\colvec{x \\ y \\ z}
               =\colvec{rx \\ ry \\ rz} \) 
      \partsitem \( 
               \colvec{x_1 \\ y_1 \\ z_1}
               +\colvec{x_2 \\ y_2 \\ z_2}
               =\colvec{0 \\ 0 \\ 0}   
               \quad\text{and}\quad
               r\colvec{x \\ y \\ z}
               =\colvec{0 \\ 0 \\ 0} \)
    \end{exparts}
    \begin{answer}
      \begin{exparts}
        \partsitem It is not a vector space.
          \begin{equation*}
            (1+1)\cdot\colvec[r]{1 \\ 0 \\ 0}\neq
            \colvec[r]{1 \\ 0 \\ 0}
            +\colvec[r]{1 \\ 0 \\ 0}
          \end{equation*}
        \partsitem It is not a vector space.
          \begin{equation*}
            1\cdot\colvec[r]{1 \\ 0 \\ 0}\neq\colvec[r]{1 \\ 0 \\ 0}
          \end{equation*}
      \end{exparts}   
    \end{answer}
  \recommended \item
    For each, decide if it is a vector space;
    the intended operations are the natural ones.
    \begin{exparts}
      \partsitem The \definend{diagonal} \( \nbyn{2} \) matrices
        \begin{equation*}
          \set{\begin{mat}
            a  &0  \\
            0  &b
          \end{mat}\suchthat a,b\in\Re}
        \end{equation*}
      \partsitem This set of \( \nbyn{2} \) matrices
        \begin{equation*}
          \set{\begin{mat}
            x    &x+y  \\
            x+y  &y
          \end{mat}\suchthat x,y\in\Re}
        \end{equation*}
      \partsitem This set
        \begin{equation*}
          \set{\colvec{x \\ y \\ z \\ w}\in\Re^4
               \suchthat x+y+z+w=1}
        \end{equation*}
      \partsitem The set of functions
        \( \set{\map{f}{\Re}{\Re}\suchthat df/dx+2f=0} \)
      \partsitem The set of functions
        \( \set{\map{f}{\Re}{\Re}\suchthat df/dx+2f=1} \)
    \end{exparts}
    \begin{answer}
      For each ``yes'' answer, you can just check a couple of 
      linear combinations in the space.
      For each ``no'' answer, give a specific example of the failure 
      of one of the
      conditions.
      \begin{exparts}
        \partsitem Yes.
          One linear combination picked out of the air is this.
          \begin{equation*}
            2\cdot\begin{mat}
              3  &0 \\
              0  &4
            \end{mat}
            +5\cdot\begin{mat}
              6  &0 \\
              0  &7
            \end{mat}
          \end{equation*}
          The key point about the collection being a vector space is that
          the result is a diagonal matrix.
          \begin{equation*}
            =\begin{mat}
              36 &0 \\
              0  &43
            \end{mat}
          \end{equation*}
        \partsitem Yes.
          Here is a linear combination.
          \begin{equation*}
            10\cdot\begin{mat}
              1  &3  \\
              3  &2
            \end{mat}
            +11\cdot\begin{mat}
              3  &7  \\
              7  &4
            \end{mat}
          \end{equation*}
          We look for the result to follow the form that the upper left 
          and lower right entries add to make the upper right entry, and
          also the lower left.
          \begin{equation*}
            =\begin{mat}
               43  &107  \\
              107  &64
            \end{mat}
          \end{equation*}
        \partsitem No, this set is not closed under the natural addition
          operation.
          The vector of all $1/4$'s is a member of this set 
          but when added to itself the result, the 
          vector of all $1/2$'s, is a nonmember.
        \partsitem Yes.
        \partsitem No, \( f(x)=e^{-2x}+(1/2) \) is in the set but 
           \( 2\cdot f \) is not (that is, condition~(6) fails).
      \end{exparts}  
    \end{answer}
  \recommended \item
    Prove or disprove that this is a vector space: the real-valued functions
    \( f \) of one real variable such that \( f(7)=0 \).
    \begin{answer}
      It is a vector space.
      Most conditions of the definition of vector space are routine; we here
      check only closure.
      For addition,
      \( (f_1+f_2)\,(7)=f_1(7)+f_2(7)=0+0=0 \).
      For scalar multiplication,
      \( (r\cdot f)\,(7)=rf(7)=r0=0 \).  
    \end{answer}
  \recommended \item
    Show that the set \( \Re^+ \) of positive reals
    is a vector space when we interpret `\( x+y \)' to mean
    the product of \( x \) and \( y \) (so that \( 2+3 \) is \( 6 \)),
    and we interpret `\( r\cdot x \)' as the \( r \)-th power of \( x \).
    \begin{answer}
      We check \nearbydefinition{def:VecSpace}.

      First, closure under `\( + \)'
      holds because the product of two positive reals is
      a positive real.
      The second condition is satisfied because real multiplication commutes.
      Similarly, as real multiplication associates, the third checks.
      For the fourth condition, observe that multiplying a number by
      \( 1\in\Re^+ \) won't change the number.
      Fifth, any positive real has a reciprocal that is a positive real.

      The sixth, closure under `\( \cdot \)', 
      holds because any power of a positive real is a
      positive real.
      The seventh condition is just the rule that \( v^{r+s} \) equals
      the product of \( v^r \) and \( v^s \).
      The eight condition says that \( (vw)^r=v^rw^r \).
      The ninth condition asserts that \( (v^r)^s=v^{rs} \).
      The final condition says that \( v^1=v \).  
    \end{answer}
  \item 
      Is \( \set{(x,y)\suchthat x,y\in\Re} \) a vector space under
      these operations?
      \begin{exparts}
        \partsitem \( (x_1,y_1)+(x_2,y_2)=(x_1+x_2,y_1+y_2) \)
        and \( r\cdot (x,y)=(rx,y) \)
        \partsitem \( (x_1,y_1)+(x_2,y_2)=(x_1+x_2,y_1+y_2) \)
        and \( r\cdot (x,y)=(rx,0) \)
      \end{exparts}
      \begin{answer}
        \begin{exparts}
           \partsitem No: \( 1\cdot(0,1)+1\cdot(0,1)\neq (1+1)\cdot(0,1) \).
           \partsitem No; the same calculation as the prior answer shows
              a condition in the definition of a vector space that is 
              violated. 
              Another example of a violation of the conditions for a 
              vector space is that \( 1\cdot (0,1)\neq (0,1) \). 
        \end{exparts}  
      \end{answer}
  \item 
    Prove or disprove that 
    this is a vector space: the set of polynomials of
    degree greater than or equal to two, along with the zero polynomial.
    \begin{answer}
      It is not a vector space since it is not closed under addition, as
      \( (x^2)+(1+x-x^2) \) is not in the set.  
    \end{answer}
  \item 
    At this point ``the same'' is only an
    intuition, but nonetheless for each vector space identify the
    \( k \) for which the space is ``the same'' as \( \Re^k \).
    \begin{exparts}
      \partsitem The \( \nbym{2}{3} \) matrices under the usual operations
      \partsitem The \( \nbym{n}{m} \) matrices (under their usual operations)
      \partsitem This set of \( \nbyn{2} \) matrices
        \begin{equation*}
          \set{\begin{mat}
            a  &0  \\
            b  &c
          \end{mat} \suchthat a,b,c\in\Re}
        \end{equation*}
      \partsitem This set of \( \nbyn{2} \) matrices
        \begin{equation*}
          \set{\begin{mat}
            a  &0  \\
            b  &c
          \end{mat} \suchthat a+b+c=0}
        \end{equation*}
    \end{exparts}
    \begin{answer}
      \begin{exparts}
        \partsitem \( 6 \)
        \partsitem \( nm \)
        \partsitem \( 3 \)
        \partsitem To see that the answer is \( 2 \), rewrite it as
        \begin{equation*}
          \set{\begin{mat}
            a  &0  \\
            b  &-a-b
          \end{mat} \suchthat a,b\in\Re}
        \end{equation*}
        so that there are two parameters.
      \end{exparts}  
     \end{answer}
  \item 
    Using \( \vec{+} \) to represent vector addition
    and \( \,\vec{\cdot}\, \) for scalar multiplication,
    restate the definition of vector space.
    \begin{answer}
      {
      \def\plus{\mathbin{\vec{+}}}
      \def\tim{\mathbin{\vec{\cdot}}}
        A \definend{vector space}\index{vector space!definition}
        (over \( \Re \)) consists of a set \( V \) along with
        two operations `\( \plus \)' and `\( \tim \)' subject to these 
        conditions.
        Where \( \vec{v},\vec{w}\in V \), 
        (1)~their \definend{vector sum}
          \( \vec{v}\plus\vec{w} \) is an element of \( V \).
        If \( \vec{u},\vec{v},\vec{w}\in V \) then 
        (2)~\( \vec{v}\plus\vec{w}=\vec{w}\plus\vec{v} \) and 
        (3)~\( (\vec{v}\plus\vec{w})\plus\vec{u}
                 =\vec{v}\plus(\vec{w}\plus\vec{u}) \).
        (4)~There is a \definend{zero vector}
          \( \zero\in V \) such that
          \( \vec{v}\plus\zero=\vec{v}\, \) for all \( \vec{v}\in V\).
        (5)~Each \( \vec{v}\in V \) has an
          \definend{additive inverse}
          \( \vec{w}\in V \) such that \( \vec{w}\plus\vec{v}=\zero \).
        If \( r,s \) are \definend{scalars\/},
        that is, members of \( \Re \)),
        and \( \vec{v},\vec{w}\in V \) then 
        (6)~each
           \definend{scalar multiple}
           \( r\cdot\vec{v} \) is in \( V \).
        If \( r,s\in\Re \) and \( \vec{v},\vec{w}\in V \) then
        (7)~\( (r+ s)\cdot\vec{v}=r\cdot\vec{v}\plus s\cdot\vec{v} \), 
        and (8)~\( r\tim (\vec{v}+\vec{w})
           =r\tim\vec{v}+r\tim\vec{w} \),
        and (9)~\( (rs)\tim \vec{v} =r\tim (s\tim\vec{v}) \),
        and (10)~\( 1\tim \vec{v}=\vec{v} \).
     }
    \end{answer}
  \item 
    Prove these.
    \begin{exparts}
      \partsitem For any $\vec{v}\in V$, if $\vec{w}\in V$ is an additive 
        inverse of $\vec{v}$, then $\vec{v}$ is an additive inverse of 
        $\vec{w}$.
        So a vector is an additive inverse of any additive inverse of
        itself.
      \partsitem Vector addition left-cancels:~if 
        \( \vec{v},\vec{s},\vec{t}\in V \)
        then \( \vec{v}+\vec{s}=\vec{v}+\vec{t}\, \) implies
        that \( \vec{s}=\vec{t} \).
    \end{exparts}
    \begin{answer}
      \begin{exparts}
        \partsitem Let \( V \) be a vector space, 
          let \( \vec{v}\in V \), and
          assume that \( \vec{w}\in V \) is an additive inverse of $\vec{v}$
          so that \( \vec{w}+\vec{v}=\zero \).
          Because addition is commutative,
          \( \zero=\vec{w}+\vec{v}=\vec{v}+\vec{w} \),
          so therefore \( \vec{v} \) is also 
          the additive inverse of \( \vec{w} \).
        \partsitem Let \( V \) be a vector space and suppose
          \( \vec{v},\vec{s},\vec{t}\in V \).
          The additive inverse of \( \vec{v} \) is \( -\vec{v} \) so
          \( \vec{v}+\vec{s}=\vec{v}+\vec{t} \) gives that
          \( -\vec{v}+\vec{v}+\vec{s}=-\vec{v}+\vec{v}+\vec{t} \),
          which says that \( \zero+\vec{s}=\zero+\vec{t} \) and so
          \( \vec{s}=\vec{t} \).
      \end{exparts}  
     \end{answer}
  \item 
    The definition of vector spaces does not explicitly say that
    \( \zero+\vec{v}=\vec{v} \) 
    (it instead says that \( \vec{v}+\zero=\vec{v} \)).
    Show that it must nonetheless hold in any vector space.
    \begin{answer}
      Addition is commutative, so in any vector space,
      for any vector \( \vec{v} \) we have that
      \( \vec{v}=\vec{v}+\zero=\zero+\vec{v} \).  
    \end{answer}
  \recommended \item
    Prove or disprove that 
    this is a vector space: the set of all matrices, under
    the usual operations.
    \begin{answer}
      It is not a vector space since addition of two matrices of unequal
      sizes is not defined, and thus the set fails to satisfy the closure
      condition.
    \end{answer}
  \item 
    In a vector space every element has an additive inverse.
    Can some elements have two or more?
    \begin{answer}
      Each element of a vector space has one and only one additive
      inverse.

      For, let \( V \) be a vector space and suppose that \( \vec{v}\in V \).
      If \( \vec{w}_1,\vec{w}_2\in V \) are both additive inverses of
      \( \vec{v} \) then consider \( \vec{w}_1+\vec{v}+\vec{w}_2 \).
      On the one hand, we have that it equals $\vec{w}_1+(\vec{v}+\vec{w}_2)=
      \vec{w}_1+\zero=\vec{w}_1$.
      On the other hand we have that it equals $(\vec{w}_1+\vec{v})+\vec{w}_2=
      \zero+\vec{w}_2=\vec{w}_2$.
      Therefore, $\vec{w}_1=\vec{w}_2$.
    \end{answer}
  \item 
    \begin{exparts}
      \partsitem Prove that every point, line, or plane through the origin in 
         \( \Re^3 \) is a vector space under the inherited operations.
      \partsitem What if it doesn't contain the origin?
    \end{exparts}
    \begin{answer}
     \begin{exparts}
      \partsitem Every such set has the form
        \( \set{r\cdot\vec{v}+s\cdot\vec{w}\suchthat r,s\in\Re} \)
        where either or both of \( \vec{v},\vec{w} \) may be \( \zero \).
        With the inherited operations, closure of addition
        \( (r_1\vec{v}+s_1\vec{w})+(r_2\vec{v}+s_2\vec{w})
           =(r_1+r_2)\vec{v}+(s_1+s_2)\vec{w} \)
        and scalar multiplication
        \( c(r\vec{v}+s\vec{w})=(cr)\vec{v}+(cs)\vec{w} \)
        are easy.
        The other conditions are also routine.
      \partsitem No such set can be a vector space under the inherited
        operations because it does not have a zero element.
     \end{exparts}  
    \end{answer}
  \item 
    Using the idea of a vector space we can easily reprove that
    the solution set of a homogeneous linear system has either 
    one element or infinitely many elements. 
    Assume that \( \vec{v}\in V \) is not \( \zero \).
    \begin{exparts}
      \partsitem Prove that \( r\cdot\vec{v}=\zero \) if and only if \( r=0 \).
      \partsitem Prove that \( r_1\cdot\vec{v}=r_2\cdot\vec{v} \) if
      and only if \( r_1=r_2 \).
      \partsitem Prove that any nontrivial vector space is infinite.
      \partsitem Use the fact that a nonempty solution set of a homogeneous
        linear system is a vector space to draw the conclusion.
    \end{exparts}
    \begin{answer}
      Assume that \( \vec{v}\in V \) is not \( \zero \).
      \begin{exparts}
        \partsitem One direction of the if and only if is clear:~if $r=0$
          then $r\cdot\vec{v}=\zero$.
          For the other way, let \( r \) be a nonzero scalar.
          If \( r\vec{v}=\zero \) then
          \( (1/r)\cdot r\vec{v}=(1/r)\cdot \zero \) shows that
          $\vec{v}=\zero$,  contrary to the assumption.
        \partsitem Where \( r_1,r_2 \) are scalars, 
          \( r_1\vec{v}=r_2\vec{v}\, \)
          holds if and only if \( (r_1-r_2)\vec{v}=\zero \).
          By the prior item, then \( r_1-r_2=0 \).
        \partsitem A nontrivial space has a vector 
          \( \vec{v}\neq\zero \).
          Consider the set \( \set{k\cdot\vec{v}\suchthat k\in\Re} \).
          By the prior item this set is infinite.
        \partsitem The solution set is either trivial, or nontrivial.
          In the second case, it is infinite.   
     \end{exparts}  
    \end{answer}
  \item 
    Is this a vector space under the natural operations: the real-valued
    functions of one real variable that are differentiable?
    \begin{answer}
      Yes.
      A theorem of first semester calculus says that a sum of differentiable
      functions is differentiable and that
      \( (f+g)^\prime=f^\prime+g^\prime \), and that 
      a multiple of a differentiable
      function is differentiable and that \( (r\cdot f)^\prime=r\,f^\prime \). 
    \end{answer}
  \item 
    A \definend{vector space over the complex numbers}%
    \index{vector space!complex scalars}%
    \index{complex numbers!vector space over}
    $\C$ has the same definition
    as a vector space over the reals except that scalars are drawn from
    \( \C \) instead of from \( \Re \).
    Show that each of these is a vector space over the complex numbers.
    (Recall how complex numbers add and multiply:
    \( (a_0+a_1i)+(b_0+b_1i)=(a_0+b_0)+(a_1+b_1)i \) and
    \( (a_0+a_1i)(b_0+b_1i)=(a_0b_0-a_1b_1)+(a_0b_1+a_1b_0)i \).)
    \begin{exparts}
      \partsitem The set of degree~two polynomials with complex 
         coefficients
      \partsitem This set
        \begin{equation*}
          \set{\begin{mat}
                 0  &a  \\
                 b  &0
               \end{mat}\suchthat a,b\in\C\text{\ and\ }
                               a+b=0+0i }
        \end{equation*}
    \end{exparts}
    \begin{answer}
      The check is routine.
      Note that `\( 1 \)' is \( 1+0i \) and the zero elements are these.
      \begin{exparts}
        \partsitem \( (0+0i)+(0+0i)x+(0+0i)x^2 \)
        \partsitem \( \begin{mat}
                   0+0i  &0+0i  \\
                   0+0i  &0+0i
                 \end{mat} \)
      \end{exparts}  
    \end{answer}
  \item 
    Name a property shared by all of the \( \Re^n \)'s 
    but not listed as a
    requirement for a vector space.
    \begin{answer}
      Notably absent from the definition of a vector space is a distance
      measure.  
    \end{answer}
  \item 
    \begin{exparts}
      \partsitem Prove that for any four vectors
        \( \vec{v}_1,\ldots,\vec{v}_4\in V \) we can associate
        their sum in any way without changing the result.
        \begin{multline*}
          ((\vec{v}_1+\vec{v}_2)+\vec{v}_3)+\vec{v}_4
          =(\vec{v}_1+(\vec{v}_2+\vec{v}_3))+\vec{v}_4  
          =(\vec{v}_1+\vec{v}_2)+(\vec{v}_3+\vec{v}_4)  \\
          =\vec{v}_1+((\vec{v}_2+\vec{v}_3)+\vec{v}_4)  
          =\vec{v}_1+(\vec{v}_2+(\vec{v}_3+\vec{v}_4))
        \end{multline*}
        This allows us to write
        `\( \vec{v}_1+\vec{v}_2+\vec{v}_3+\vec{v}_4 \)'
        without ambiguity.
      \partsitem Prove that any two ways of associating a sum of any number of
        vectors give the same sum.
        (\textit{Hint.}  Use induction on the number of vectors.)
    \end{exparts}
    \begin{answer}
      \begin{exparts}
        \partsitem A small rearrangement does the trick.
          \begin{align*}
            (\vec{v}_1+(\vec{v}_2+\vec{v}_3))+\vec{v}_4
            &=((\vec{v}_1+\vec{v}_2)+\vec{v}_3)+\vec{v}_4  \\
            &=(\vec{v}_1+\vec{v}_2)+(\vec{v}_3+\vec{v}_4)  \\
            &=\vec{v}_1+(\vec{v}_2+(\vec{v}_3+\vec{v}_4))  \\
            &=\vec{v}_1+((\vec{v}_2+\vec{v}_3)+\vec{v}_4)
          \end{align*}
          Each equality above follows from the associativity of three vectors
          that is given as a condition in the definition of a vector space.
          For instance, the second `$=$' applies the rule
          $(\vec{w}_1+\vec{w}_2)+\vec{w}_3=\vec{w}_1+(\vec{w}_2+\vec{w}_3)$
          by taking $\vec{w}_1$ to be $\vec{v}_1+\vec{v}_2$, 
          taking $\vec{w}_2$ to be $\vec{v}_3$, 
          and taking $\vec{w}_3$ to be $\vec{v}_4$. 
        \partsitem The base case for induction is the three vector case.
          This case
          \( \vec{v}_1+(\vec{v}_2+\vec{v}_3)
          =(\vec{v}_1+\vec{v}_2)+\vec{v}_3 \) is one of the conditions in
          the definition of a vector space.

          For the inductive step, assume that any two sums of three vectors,
          any two sums of four vectors,
          \ldots, any two sums of $k$ vectors 
          are equal no matter how we  parenthesize the sums.
          We will show that any sum of \( k+1 \) vectors equals this one
          \( ((\cdots((\vec{v}_1+\vec{v}_2)+\vec{v}_3)+\cdots)+\vec{v}_k)
              +\vec{v}_{k+1} \).

          Any parenthesized sum has an outermost `\( + \)'.
          Assume that it lies between \( \vec{v}_m \) and \( \vec{v}_{m+1} \)
          so the sum looks like this.
          \begin{equation*}
            (\cdots\,\vec{v}_1\cdots\vec{v}_m\,\cdots)
           +(\cdots\,\vec{v}_{m+1}\cdots\vec{v}_{k+1}\,\cdots)
          \end{equation*}
          The second half involves fewer than $k+1$ additions, so
          by the inductive hypothesis we can re-parenthesize it
          so that it reads left to right from the inside out, and in 
          particular, so that its outermost `$+$' occurs right before 
          $\vec{v}_{k+1}$.
          \begin{equation*}
            =(\cdots\,\vec{v}_1\,\cdots\,\vec{v}_m\,\cdots)
             +((\cdots(\vec{v}_{m+1}+\vec{v}_{m+2})+\cdots+\vec{v}_{k})
                 +\vec{v}_{k+1})
          \end{equation*}
          Apply the associativity of the sum of three things
          \begin{equation*}
            =((\,\cdots\, \vec{v}_1\,\cdots\,\vec{v}_m\,\cdots\,)
             +(\,\cdots\,(\vec{v}_{m+1}+\vec{v}_{m+2})+\cdots\,\vec{v}_k))
             +\vec{v}_{k+1}
          \end{equation*}
          and finish by applying the inductive hypothesis inside these 
          outermost parenthesis.
      \end{exparts}  
    \end{answer}
  \item \nearbyexample{ex:ColsIntEntNotVS} gives a subset of $\Re^2$
   that is not a vector space, under the obvious operations, because
   while it is closed under addition, it is not closed under scalar
   multiplication.
   Consider the set of vectors in the plane whose components have the 
   same sign or are~$0$.
   Show that this set is closed under scalar multiplication but not 
   addition.
   \begin{answer}
     Let $\vec{v}$ be a member of $\Re^2$ with components $v_1$ and $v_2$.
     We can abbreviate the condition that both components have the same
     sign or are~$0$ by $v_1v_2\geq 0$.

     To show the set is closed under scalar multiplication, observe that
     the components of $r\vec{v}$ satisfy $(rv_1)(rv_2)=r^2(v_1v_2)$
     and $r^2\geq 0$ so $r^2v_1v_2\geq 0$.

     To show the set is not closed under addition we need only produce one
     example.  
     The vector with components $-1$ and~$0$, when added to the vector
     with components $0$ and~$1$ makes a vector with mixed-sign components
     of $-1$ and~$1$.
   \end{answer}
  % \item 
  %   For any vector space, a subset that is itself a vector space
  %   under the inherited operations
  %   (e.g., a plane through the origin inside of \( \Re^3 \))
  %   is a \definend{subspace}.
  %   \begin{exparts}
  %     \partsitem Show that \( \set{a_0+a_1x+a_2x^2\suchthat a_0+a_1+a_2=0} \)
  %       is a subspace of the vector space of degree~two polynomials.
  %     \partsitem Show that this is a subspace of the \( \nbyn{2} \) matrices.
  %       \begin{equation*}
  %         \set{\begin{mat}
  %                a  &b  \\
  %                c  &0
  %              \end{mat} \suchthat a+b=0}
  %       \end{equation*}
  %     \partsitem Show that a nonempty subset \( S \) of a real vector
  %       space is a
  %       subspace if and only if it is closed under linear combinations of
  %       pairs of vectors: whenever \( c_1,c_2\in\Re \) and
  %       \( \vec{s}_1,\vec{s}_2\in S \) then the combination
  %       \( c_1\vec{v}_1+c_2\vec{v}_2 \) is in \( S \).
  %   \end{exparts}
  %   \begin{answer}
  %     \begin{exparts}
  %       \partsitem We outline the check of the conditions from
  %         \nearbydefinition{def:VecSpace}.

  %         Additive closure holds because if \( a_0+a_1+a_2=0 \)
  %         and \( b_0+b_1+b_2=0 \) then
  %         \begin{equation*}
  %           (a_0+a_1x+a_2x^2)+(b_0+b_1x+b_2x^2)=
  %           (a_0+b_0)+(a_1+b_1)x+(a_2+b_2)x^2
  %         \end{equation*}
  %         is in the set since
  %         \( (a_0+b_0)+(a_1+b_1)+(a_2+b_2)=(a_0+a_1+a_2)+(b_0+b_1+b_2) \)
  %         is zero.
  %         The second through fifth conditions are easy.

  %         Closure under scalar multiplication holds because
  %         if \( a_0+a_1+a_2=0 \) then
  %         \begin{equation*}
  %           r\cdot(a_0+a_1x+a_2x^2)=
  %           (ra_0)+(ra_1)x+(ra_2)x^2
  %         \end{equation*}
  %         is in the set as \( ra_0+ra_1+ra_2=r(a_0+a_1+a_2) \) is zero.
  %         The remaining conditions here are also easy.
  %       \partsitem This is similar to the prior answer.
  %       \partsitem Call the vector space \( V \).
  %         We have two implications: left to right, if \( S \) is a subspace
  %         then it is closed under linear combinations of pairs of vectors and,
  %         right to left, if a nonempty subset is closed under linear
  %         combinations of pairs of vectors then it is a subspace.
  %         The left to right implication is easy; we here sketch the
  %         other one by
  %         assuming \( S \) is nonempty and closed, and checking the conditions
  %         of \nearbydefinition{def:VecSpace}.

  %         First, to show closure under addition, if
  %         \( \vec{s}_1,\vec{s}_2\in S \) then \( \vec{s}_1+\vec{s}_2\in S \)
  %         as \( \vec{s}_1+\vec{s}_2=1\cdot\vec{s}_1+1\cdot\vec{s}_2 \).
  %         Second, for any \( \vec{s}_1,\vec{s}_2\in S \), because addition
  %         is inherited from \( V \), the sum \( \vec{s}_1+\vec{s}_2 \)
  %         in \( S \) equals the sum \( \vec{s}_1+\vec{s}_2 \)
  %         in \( V \) and that equals the sum \( \vec{s}_2+\vec{s}_1 \) in
  %         \( V \) and that in turn equals the sum \( \vec{s}_2+\vec{s}_1 \) in
  %         \( S \).
  %         The argument for the third condition is similar to that for the
  %         second.
  %         For the fourth, suppose that 
  %         \( \vec{s} \) is in the nonempty set \( S \)
  %         and note that \( 0\cdot\vec{s}=\zero\in S \); showing that 
  %         the \( \zero \) of
  %         \( V \) acts under the inherited operations as the additive identity
  %         of \( S \) is easy.
  %         The fifth condition is satisfied because for any \( \vec{s}\in S \)
  %         closure under linear combinations shows that the vector
  %         \( 0\cdot\zero+(-1)\cdot\vec{s} \) is in \( S \); showing 
  %         that it is the
  %         additive inverse of \( \vec{s} \) under the inherited operations is
  %         routine.

  %         The proofs for the remaining conditions are similar.
  %     \end{exparts}  
  %   \end{answer}
\end{exercises}





















 
\subsection{Subspaces and Spanning Sets}
\index{subspace|(}
In \nearbyexample{PlaneThruOriginSubsp} we saw a
vector space that is a subset of $\Re^2$, a line through the origin. 
There, the vector space $\Re^2$ contains inside it another
vector space, the line.

\begin{definition} \label{df:Subspace}
%<*df:Subspace>
For any vector space,
a \definend{subspace}\index{vector space!subspace}\index{subspace!definition} 
is a subset that is itself a vector space,
under the inherited operations.
%</df:Subspace>
\end{definition}

\begin{example}  \label{ex:PlaneSubspRThree}
This plane through the origin 
\begin{equation*}
  P=\set{\colvec{x \\ y \\ z}\suchthat x+y+z=0}
\end{equation*}
is a subspace of \( \Re^3 \).
As required by the definition 
the plane's operations are inherited from the larger space,
that is,
vectors add in $P$ as they add in $\Re^3$
\begin{equation*}
   \colvec{x_1 \\ y_1 \\ z_1}+\colvec{x_2 \\ y_2 \\ z_2}
   =\colvec{x_1+x_2 \\ y_1+y_2 \\ z_1+z_2}
\end{equation*}
and scalar multiplication is also the same as in $\Re^3$.
To show that $P$ is a subspace we need only note that it is a subset and then
verify that it is a space.
We won't check all ten conditions, just the two closure ones.
For closure under addition, note that if
the summands satisfy that
$x_1+y_1+z_1=0$ and $x_2+y_2+z_2=0$ then the sum satisfies that
$(x_1+x_2)+(y_1+y_2)+(z_1+z_2)=(x_1+y_1+z_1)+(x_2+y_2+z_2)=0$.
For closure under scalar multiplication, if
$x+y+z=0$ then the scalar multiple has
$rx+ry+rz=r(x+y+z)=0$.
\end{example}

\begin{example}   \label{ex:SubspacesRTwo}
The \( x \)-axis in \( \Re^2 \) 
is a subspace, where
the addition and scalar multiplication operations are 
the inherited ones.
\begin{equation*}
  \colvec{x_1 \\ 0}
    +
  \colvec{x_2 \\ 0}
    =
  \colvec{x_1+x_2 \\ 0}
  \qquad
  r\cdot\colvec{x \\ 0}
  =\colvec{rx \\ 0}
\end{equation*}
As in the prior example, to verify directly from the definition 
that this is a subspace we simply note that it is a
subset and then check that it satisfies
the conditions in definition of a vector space.
For instance the two closure conditions are 
satisfied: adding two vectors with a second component of zero results
in a vector with a second component of zero and multiplying a 
scalar times a vector with a second component of zero 
results in a vector with a second component of zero.
\end{example}

\begin{example}
Another subspace of $\Re^2$ is 
its trivial subspace.
\begin{equation*}
  \set{\colvec[r]{0 \\ 0}}
\end{equation*}
\end{example}

%<*ProperSubspace>
Any vector space has a trivial subspace 
\( \set{\zero\,} \).\index{subspace!trivial}%
\index{trivial subspace} 
At the opposite extreme, any vector space has itself for a subspace.
A subspace that is not the entire space is a 
\definend{proper}\index{subspace!proper}%
\index{proper subspace} subspace.
%</ProperSubspace>

\begin{example}  \label{ex:LinSubspPolyThree}
Vector spaces that are not $\Re^n$'s also have subspaces.
The space of cubic polynomials
\( \set{a+bx+cx^2+dx^3\suchthat a,b,c,d\in\Re} \) 
has a subspace comprised of all linear polynomials
\( \set{m+nx\suchthat m,n\in\Re} \).
\end{example}

\begin{example}
Another example of a subspace that is not a subset of an $\Re^n$ 
followed the definition of a vector space.
The space in \nearbyexample{ex:RealValuedFcns}
of all real-valued functions of one real variable 
\( \set{f\suchthat \map{f}{\Re}{\Re} } \) has the subspace in
\nearbyexample{ex:SpaceSatisfyingDiffQ}
of functions satisfying
the restriction $(d^2\,f/dx^2)+f=0$.
\end{example}

\begin{example} \label{ex:OperNotInherit}
The definition requires that the 
addition and scalar multiplication operations
must be the ones inherited from the larger space.
The set \( S=\set{1} \) is a subset of \( \Re^1 \).
And, under the operations $1+1=1$ and  $r\cdot 1=1$
the set $S$ is a vector space, specifically, a trivial space.
However, $S$ is not a subspace of \( \Re^1 \) because those aren't the
inherited operations, since of course \( \Re^1 \) has \( 1+1=2 \).
\end{example}

\begin{example}  \label{cex:RPlusNotSubSp}
Being vector spaces themselves, subspaces must satisfy the closure
conditions.
The set \( \Re^+ \) is not a subspace of the vector space \( \Re^1 \)
because with the inherited operations it is not closed under scalar
multiplication: if \( \vec{v}=1 \) then \( -1\cdot\vec{v}\not\in\Re^+ \).
\end{example}

The next result says that \nearbyexample{cex:RPlusNotSubSp} is prototypical. 
The only way that a subset can fail to be a subspace, 
if it is nonempty and uses the inherited operations,
is if it isn't closed.

\begin{lemma}     \label{th:SubspIffClosed} \index{subspace!closed}
%<*lm:SubspIffClosed>
For a nonempty subset \( S \) of a vector space, under the inherited 
operations the following are equivalent 
statements.\appendrefs{equivalence of statements}   %\spacefactor=1000
\begin{tfae}
  \item \( S \) is a subspace of that vector space
  \item \( S \) is closed under linear combinations of pairs of vectors:
    for any vectors \( \vec{s}_1,\vec{s}_2\in S \) and scalars \( r_1,r_2 \)
    the vector \( r_1\vec{s}_1+r_2\vec{s}_2 \) is in \( S \)
  \item \( S \) is closed under linear combinations of any number of vectors:
    for any vectors \( \vec{s}_1,\ldots,\vec{s}_n\in S \) and scalars
    \( r_1, \ldots,r_n \)
    the vector \( r_1\vec{s}_1+\cdots+r_n\vec{s}_n \) is an element of \( S \).
\end{tfae}
%</lm:SubspIffClosed>
\end{lemma}
\noindent Briefly, a subset is a 
subspace if and only if  it is closed under linear combinations.

\begin{proof}
%<*pf:SubspIffClosed0>
`The following are equivalent' means that each pair of 
statements are equivalent.
\begin{equation*}
  (1)\!\iff\!(2)
  \qquad
  (2)\!\iff\!(3)
  \qquad
  (3)\!\iff\!(1)
\end{equation*}
We will prove the equivalence by establishing that
\( (1)\implies (3)\implies (2)\implies (1)\).
This strategy is suggested by the observation that the implications 
\( (1)\implies (3) \) and \( (3)\implies (2) \) are easy and so we need only
argue that \( (2)\implies (1) \).
%</pf:SubspIffClosed0>

%<*pf:SubspIffClosed1>
Assume that \( S \) is a nonempty subset of a vector space
$V$ that is closed under combinations of pairs of vectors.
We will show that $S$ is a vector space by checking the conditions.

The vector space definition has five conditions on addition.
First, for closure under addition, if
\( \vec{s}_1,\vec{s}_2\in S \) then \( \vec{s}_1+\vec{s}_2\in S \),
as it
is a combination of a pair of vectors 
and we are assuming that~\( S \) is closed under those.
Second, for any \( \vec{s}_1,\vec{s}_2\in S \), because addition
is inherited from \( V \), the sum \( \vec{s}_1+\vec{s}_2 \)
in \( S \) equals the sum \( \vec{s}_1+\vec{s}_2 \)
in \( V \), and that equals the sum \( \vec{s}_2+\vec{s}_1 \) in
\( V \) (because $V$ is a vector space, its addition is commutative), 
and that in turn equals the sum \( \vec{s}_2+\vec{s}_1 \) in \( S \).
The argument for the third condition is similar to that for the second.
For the fourth, consider the zero vector of \( V \) and note that 
closure of $S$ under linear combinations of pairs of vectors gives that 
\( 0\cdot\vec{s}+0\cdot\vec{s}=\zero \) is an element of \( S\) 
(where \( \vec{s} \) is any member of the nonempty set \( S \));
checking that \( \zero \) acts under the inherited operations as the additive
identity of \( S \) is easy.
The fifth condition is satisfied because for any \( \vec{s}\in S \),
closure under linear combinations of pairs of vectors shows that 
\( 0\cdot\zero+(-1)\cdot\vec{s} \) is an element of~\( S \), and it is
obviously the
additive inverse of \( \vec{s} \) under the inherited operations.
%</pf:SubspIffClosed1>
The verifications for the scalar multiplication conditions are similar;
see \nearbyexercise{exer:SubspIffClosed}.
\end{proof}

We will usually verify that a subset is a subspace by checking that it 
satisfies statement~(2).

\begin{remark}
At the start of this chapter we introduced vector spaces as collections in
which linear combinations ``make sense.''
% The vector space definition has ten conditions but eight of them, the
% conditions not about closure, simply ensure that referring to the
% operations by the names `addition' and `scalar multiplication' is sensible.
% For instance, calling an operation `$+$' would not
% be odd unless it is commutative, as in condition~(2). 
% The proof above checks that the subspace satisfies these 
% eight non-closure conditions because they hold in the
% surrounding vector space, provided that the nonempty set $S$ satisfies
% \nearbylemma{th:SubspIffClosed}'s statement~(2) 
% (e.g., commutativity of addition in $S$ follows right from
% commutativity of addition in $V$).
% There is a second way in which we can regard vector spaces as structures in 
% which linear combinations are ``sensible.''
% It has to do with the closure conditions.
\nearbylemma{th:SubspIffClosed}'s statements (1)-(3) 
say that we can always make sense of 
an expression like
$r_1\vec{s}_1+r_2\vec{s}_2$
in that the vector described is in the set~$S$.

As a contrast, consider the set $T$ of two-tall vectors whose entries add to
a number greater than or equal to zero.
Here we cannot just write any linear combination such as $2\vec{t}_1-3\vec{t}_2$
and be confident the result is an element of $T$.
\end{remark}

\nearbylemma{th:SubspIffClosed} suggests that a good way to think of
a vector space is as a collection of unrestricted linear combinations.
The next two examples take some spaces and recasts their
descriptions to be in that form.

\begin{example}
We can show that this plane through the origin subset of $\Re^3$
\begin{equation*}
  S=\set{\colvec{x \\ y \\ z}\suchthat x-2y+z=0}
\end{equation*}
is a subspace under the usual addition and scalar multiplication
operations of column vectors by checking that it is nonempty and closed under
linear combinations of two vectors.
But there is another way.
Think of  $x-2y+z=0$  as a one-equation linear system and parametrize it
by expressing the leading
variable in terms of the free variables $x=2y-z$.
\begin{equation*}
     S
     =\set{\colvec{2y-z \\ y \\ z}\suchthat y,z\in\Re}
     =\set{y\colvec[r]{2 \\ 1 \\ 0}+
            z\colvec[r]{-1 \\ 0 \\ 1}\suchthat y,z\in\Re}
    \tag{$*$}
\end{equation*}
Now, to show that this is a subspace consider
$r_1\vec{s}_1+r_2\vec{s}_2$.
Each $\vec{s}_i$ is a linear combination of the two vectors in ($*$)
so this is a linear combination of linear combinations.
\begin{equation*}
  r_1\cdot(y_1\colvec[r]{2 \\ 1 \\ 0}+
            z_1\colvec[r]{-1 \\ 0 \\ 1})
  +
  r_2\cdot(y_2\colvec[r]{2 \\ 1 \\ 0}+
            z_2\colvec[r]{-1 \\ 0 \\ 1})
\end{equation*}
The Linear Combination Lemma, Lemma~One.III.\ref{lm:LinearCombinationLemma},
shows that the total is a linear combination of the two vectors
and so \nearbylemma{th:SubspIffClosed}'s statement~(2) is satisfied.
\end{example}

\begin{example} \label{ex:ParamSubspace}
This is a subspace of the \( \nbyn{2} \) matrices $\matspace_{\nbyn{2}}$.
\begin{equation*}
  L=\set{\begin{mat}
         a  &0  \\
         b  &c
       \end{mat}
       \suchthat a+b+c=0}
\end{equation*}
To parametrize, express the condition as $a=-b-c$.
\begin{equation*}
  L
  =\set{\begin{mat}
         -b-c  &0  \\
         b     &c
       \end{mat}
       \suchthat b,c\in\Re}
  =\set{b\begin{mat}[r]
         -1    &0  \\
         1     &0
       \end{mat}
       +c\begin{mat}[r]
         -1    &0  \\
         0     &1
       \end{mat}
       \suchthat b,c\in\Re}
\end{equation*}
As above, we've described the subspace as a collection of unrestricted linear
combinations.
To show it is a subspace, note that a linear combination of vectors from
$L$ is a linear combination of linear combinations and so statement~(2)
is true.
\end{example}

% Parametrization is easy, but important.
% We shall use it often.

\begin{definition} \label{df:Span}
%<*df:Span>
The \definend{span}\index{span}\index{sets!span of}\index{closure} (or
\definend{linear closure}) of a nonempty subset \( S \) of a
vector space is the set of all linear combinations of vectors from \( S \).
\begin{equation*}
  \spanof{S} =\{ c_1\vec{s}_1+\cdots+c_n\vec{s}_n
            \suchthat c_1,\ldots, c_n\in\Re
            \text{\ and\ } \vec{s}_1,\ldots,\vec{s}_n\in S \}
\end{equation*}
The span of the empty subset of a vector space is its trivial subspace.
%</df:Span>
\end{definition}

%<*NotationForSpan>
\noindent No notation for the span is completely standard.
The square brackets used here are common but so are
`$\mbox{span}(S)$' and `$\mbox{sp}(S)$'.
%</NotationForSpan>

\begin{remark}
In Chapter One, after we showed that we can write the solution
set of a homogeneous linear system as 
$\set{c_1\vec{\beta}_1+\cdots+c_k\vec{\beta}_k\suchthat
  c_1,\ldots,c_k\in\Re}$,
we described that as the set `generated' by the $\smash{\vec{\beta}}$'s.
We now call that the span of
$\set{\vec{\beta}_1,\ldots,\vec{\beta}_k}$.

Recall also from that proof that 
the span of the empty set is defined to be the set \( \set{\zero} \) because
of the convention that a trivial linear combination, a combination of 
zero-many vectors, adds to~\( \zero \).
Besides, defining the empty set's span to be the trivial subspace 
is convenient because it keeps results
like the next one from needing exceptions for the empty set.
\end{remark}

\begin{lemma}   \label{le:SpanIsASubsp}
%<*lm:SpanIsASubsp>
In a vector space, the span of any subset is a subspace.
%</lm:SpanIsASubsp>
\end{lemma}

\begin{proof}
%<*pf:SpanIsASubsp>
If the subset \( S \)
is empty then by definition its span is the trivial
subspace.
If \( S\) is not empty then by \nearbylemma{th:SubspIffClosed} we need
only check that the span \( \spanof{S} \) is closed under linear combinations
of pairs of elements.
For a pair of vectors from that span,
\( \vec{v}=c_1\vec{s}_1+\cdots+c_n\vec{s}_n \) and
\( \vec{w}=c_{n+1}\vec{s}_{n+1}+\cdots+c_m\vec{s}_m \),
a linear combination 
\begin{multline*}
  p\cdot(c_1\vec{s}_1+\cdots+c_n\vec{s}_n)+
       r\cdot(c_{n+1}\vec{s}_{n+1}+\cdots+c_m\vec{s}_m)  \\
  =
  pc_1\vec{s}_1+\cdots+pc_n\vec{s}_n
    +rc_{n+1}\vec{s}_{n+1}+\cdots+rc_m\vec{s}_m
\end{multline*}
is a linear combination of elements of~\( S \) 
and so is an element of \( \spanof{S} \)
(possibly some of the $\vec{s}_i$'s from $\vec{v}$ equal some 
of the $\vec{s}_j$'s from $\vec{w}$ but that does not matter).
%</pf:SpanIsASubsp>
\end{proof}

The converse of the lemma
holds: any subspace is the span of some set, because 
a subspace is obviously the span of itself, the set of all of its members.
Thus a subset of a vector space is a subspace if and only if it is a span.
This fits the intuition 
that a good way to think of a vector space is as
a collection in which linear combinations are sensible.

Taken together, \nearbylemma{th:SubspIffClosed} and
\nearbylemma{le:SpanIsASubsp} show that the span of a subset $S$ of a
vector space is the smallest subspace containing all of the members of $S$.

\begin{example}   \label{ex:SpanSingVec}
In any vector space \( V \), for any vector \( \vec{v}\in V \), the set
\( \set{r\cdot\vec{v} \suchthat r\in\Re} \) is a subspace of \( V \).
For instance, for any vector \( \vec{v}\in\Re^3 \)
the line through the origin containing that vector
\( \set{k\vec{v}\suchthat k\in\Re } \) is a subspace of \( \Re^3 \).
This is true even if $\vec{v}$ is the zero vector, in which case 
it is the degenerate line, the trivial subspace.
\end{example}

\begin{example}
The span of this set
is all of $\Re^2$.
\begin{equation*}
  \set{\colvec[r]{1 \\ 1},\colvec[r]{1 \\ -1}}
\end{equation*}
We know that the span is some subspace of~$\Re^2$.
To check that it is all of~$\Re^2$ 
we must show that any member of $\Re^2$ is a linear combination
of these two vectors.
So we ask:~for which
vectors with real components $x$ and~$y$ 
are there scalars $c_1$ and $c_2$ such that this holds?
\begin{equation*}
   c_1\colvec[r]{1 \\ 1}+c_2\colvec[r]{1 \\ -1}=\colvec{x \\ y}
  \tag{$*$}
\end{equation*} 
Gauss's Method 
\begin{equation*}
  \begin{linsys}{2}
    c_1  &+  &c_2  &=  &x  \\
    c_1  &-  &c_2  &=  &y
  \end{linsys}
  \grstep{-\rho_1+\rho_2}
  \begin{linsys}{2}
    c_1  &+  &c_2    &=  &x\hfill  \\
         &   &-2c_2  &=  &-x+y
  \end{linsys}
\end{equation*}
with back substitution gives $c_2=(x-y)/2$ and $c_1=(x+y)/2$.
This shows that for any $x,y$ there 
are appropriate coefficients $c_1,c_2$ making~($*$)
true\Dash
we can write any element of $\Re^2$ as a linear combination of the 
two given ones.
For instance, for $x=1$ and $y=2$ the coefficients $c_2=-1/2$ and
$c_1=3/2$ will do.
\end{example}

Since spans are subspaces, and we know that a
good way to understand a subspace is
to parametrize its description, we can try to understand a set's span in 
that way.

\begin{example}
Consider, in the vector space of quadratic polynomials~\( \polyspace_2 \), 
the span of the set \( S=\set{3x-x^2, 2x} \).
By the definition of span, it is the set of unrestricted linear
combinations of the two $\set{c_1(3x-x^2)+c_2(2x)\suchthat c_1,c_2\in\Re}$.
Clearly polynomials in this span must have a constant term of zero.
Is that necessary condition also sufficient?

We are asking:~for which members $a_2x^2+a_1x+a_0$ 
of $\polyspace_2$ are there $c_1$ and $c_2$ such that
$a_2x^2+a_1x+a_0=c_1(3x-x^2)+c_2(2x)$? 
Polynomials are equal when their coefficients are equal so
we want conditions on $a_2$, $a_1$, and $a_0$ 
making that triple a solution of this system.
\begin{equation*}
  \begin{linsys}{2}
    -c_1  &   &     &=  &a_2   \\
    3c_1  &+  &2c_2 &=  &a_1   \\
          &   &0    &=  &a_0                                   
  \end{linsys}
\end{equation*} 
Gauss's Method and back-substitution gives 
$c_1=-a_2$, and $c_2=(3/2)a_2+(1/2)a_1$, and $0=a_0$.
Thus 
%the only condition on elements  $a_0+a_1x+a_2x^2$ of the span
%is the condition that we knew: 
as long as there is no constant term $a_0=0$
we can give coefficients $c_1$ and $c_2$
to describe that polynomial as an element of the span.
For instance, for the polynomial $0-4x+3x^2$, the coefficients
$c_1=-3$ and $c_2=5/2$ will do.
So the span of the given set is 
$\spanof{S}=\set{a_1x+a_2x^2\suchthat a_1,a_2\in\Re}$. 

Incidentally, this shows that
the set \( \set{x,x^2} \) spans the same subspace.
A space can have more than one spanning set.
Two other sets spanning this subspace are
\( \set{x,x^2,-x+2x^2} \) and
\( \set{x,x+x^2,x+2x^2,\ldots\,} \).
% (Usually we prefer to work with spanning sets that have only
% a small number of members.)
\end{example}

\begin{example}  \label{ex:SubspRThree}
The picture below shows the subspaces of \( \Re^3 \) that we now know of: the 
trivial subspace, lines through the origin,
planes through the origin, and the whole space.
(Of course, the picture shows only a few of the infinitely many cases.
Line segments connect subsets with 
their supersets.) 
In the next section we will prove that $\Re^3$ has no other
kind of subspace, so in fact this lists them all.

This describes each subspace 
as the span of a set with a minimal number of members.
With this, the subspaces 
fall naturally into levels\Dash planes on one level, 
lines on another,
etc.
\begin{center}
  \setlength{\unitlength}{4pt}
  \begin{picture}(75,38)(0,-2) %subspaces of R-three
      \thinlines
      \put(45,31){\makebox(0,0)[bl]{
                        \scriptsize \( %\Re^3=
                                   \set{x\colvec[r]{1 \\ 0 \\ 0}
                                              +y\colvec[r]{0 \\ 1 \\ 0}
                                              +z\colvec[r]{0 \\ 0 \\ 1}} \)} }
    %Next the dimension two subspaces
      \put(43,31){\line(-4,-1){25} } % connects R3 to 2,a
      %set 2,a:
      \put(0,22){\makebox(0,0)[l]{\scriptsize\( \set{x\colvec[r]{1 \\ 0 \\ 0}
                                                 +y\colvec[r]{0 \\ 1 \\ 0} }\) }}
      \put(48,29){\line(-3,-1){12} } % connects R3 to 2,b
      %set 2,b:
      \put(20,22){\makebox(0,0)[l]{\scriptsize\( \set{x\colvec[r]{1 \\ 0 \\ 0}
                                                 +z\colvec[r]{0 \\ 0 \\ 1} }\) }}
      \put(53,29){\line(-1,-1){3}} % connects R3 to 2,c
      %set 2,c:
      \put(40,22){\makebox(0,0)[l]{\scriptsize\( \set{x\colvec[r]{1 \\ 1 \\ 0}
                                                 +z\colvec[r]{0 \\ 0 \\ 1} }\) }}
      \put(60,22){\makebox(0,0)[l]{$\cdots$} }
    %Next the dimension one subspaces
      \put(7,18){\line(-1,-4){1} } % connects 2,a to 1,a
      \put(22,18){\line(-3,-1){13} } % connects 2,c to 1,a
      %set 1,a:
      \put(0,10){\makebox(0,0)[l]{\scriptsize\( \set{x\colvec[r]{1 \\ 0 \\ 0}} \)} }
      \put(12,18){\line(1,-2){2} } % connects 2,a to 1,b
      %set 1,b:
      \put(12,10){\makebox(0,0)[l]{\scriptsize\( \set{y\colvec[r]{0 \\ 1 \\ 0}} \)} }
      \put(14,18){\line(2,-1){10} } % connects 2,a to 1,c
      %set 1,c:
      \put(24,10){\makebox(0,0)[l]{\scriptsize\( \set{y\colvec[r]{2 \\ 1 \\ 0}} \)} }
      \put(46,18){\line(-1,-1){3.25} } %connects 2,c to 1,d
      %set 1,d:
      \put(36,10){\makebox(0,0)[l]{\scriptsize\( \set{y\colvec[r]{1 \\ 1 \\ 1}} \)} }
      \put(48,10){\makebox(0,0)[l]{$\cdots$} }
    %Finally, the trivial subspace.
      \put(9,7.5){\line(4,-1){30} } %connects 1,a to trivial
      \put(18.9,7.7){\line(3,-1){20} } %connects 1,b to trivial
      \put(30.4,6.8){\line(2,-1){10} } %connects 1,c to trivial
      \put(41,6){\line(1,-2){1.5} } %connects 1,d to trivial
      \put(45,0){\makebox(0,0){\scriptsize\( \set{\colvec[r]{0 \\ 0 \\ 0} } \)} }
  \end{picture}
\end{center}
\end{example}

So far in this chapter we have seen that to study the
properties of linear combinations, the right setting is a
collection that is closed under these combinations.
In the first subsection we introduced such collections, vector spaces,
and we saw a great variety of examples.
In this subsection we saw still
more spaces, ones that are subspaces of others.
In all of the variety there is a commonality.
\nearbyexample{ex:SubspRThree} above 
brings it out:~vector spaces and subspaces are best understood as a span, 
and especially as a span of a small number of vectors.
The next section studies spanning sets that are minimal.





\begin{exercises}
  \recommended \item
    Which of these subsets of the vector space of \( \nbyn{2} \) matrices
    are subspaces under the inherited operations?
    For each one that is a subspace, parametrize its description.
    For each that is not, give a condition that fails.
    \begin{exparts}
      \partsitem \( \set{\begin{mat}
                      a  &0  \\
                      0  &b
                    \end{mat}  \suchthat a,b\in\Re}  \)
      \partsitem \( \set{\begin{mat}
                      a  &0  \\
                      0  &b
                    \end{mat}  \suchthat a+b=0} \)
      \partsitem \( \set{\begin{mat}
                      a  &0  \\
                      0  &b
                    \end{mat}  \suchthat a+b=5} \)
      \partsitem \( \set{\begin{mat}
                      a  &c  \\
                      0  &b
                    \end{mat}  \suchthat a+b=0, c\in\Re} \)
    \end{exparts}
    \begin{answer}
      By \nearbylemma{th:SubspIffClosed}, to see if each
      subset of $\matspace_{\nbyn{2}}$ is a subspace, we need only
      check if it is nonempty and closed.
      \begin{exparts}
        \partsitem Yes, we can easily check that it is nonempty and closed.
          This is a parametrization.
          \begin{equation*}
            \set{a\begin{mat}[r]
                    1  &0  \\
                    0  &0
                  \end{mat}
                 +b\begin{mat}[r]
                    0  &0  \\
                    0  &1  
                   \end{mat}
                 \suchthat a,b\in\Re}
          \end{equation*}
          By the way, the parametrization also shows that it is a subspace,
          since it is given as the span of the two-matrix set,
          and any span is a subspace.
        \partsitem Yes; it is easily checked to be nonempty and closed.
         Alternatively, as mentioned in the prior answer, the existence
         of a parametrization shows that it is a subspace.
         For the parametrization, 
         the condition $a+b=0$ can be rewritten as $a=-b$.
         Then we have this.
          \begin{equation*}
            \set{\begin{mat}
                    -b  &0  \\
                    0   &b
                  \end{mat}
                 \suchthat b\in\Re}
            =\set{b\begin{mat}[r]
                    -1  &0  \\
                    0   &1
                  \end{mat}
                 \suchthat b\in\Re}
          \end{equation*}
        \partsitem No.
          It is not closed under addition.
          For instance, 
          \begin{equation*}
            \begin{mat}[r]
              5  &0  \\
              0  &0
            \end{mat}
            +\begin{mat}[r]
              5  &0  \\
              0  &0
            \end{mat}
            =\begin{mat}[r]
              10  &0  \\
              0  &0
            \end{mat}
          \end{equation*}
          is not in the set.
          (This set is also not closed under scalar multiplication,
          for instance, it does not contain the zero matrix.)
        \partsitem Yes.
          \begin{equation*}
            \set{b\begin{mat}[r]
                    -1  &0  \\
                    0   &1
                  \end{mat}
                 +c\begin{mat}[r]
                    0  &1  \\
                    0  &0  
                   \end{mat}
                 \suchthat b,c\in\Re}
          \end{equation*}
      \end{exparts}  
     \end{answer}
  \recommended \item 
    Is this a subspace of \( \polyspace_2 \):
    \( \set{a_0+a_1x+a_2x^2\suchthat a_0+2a_1+a_2=4} \)?
    If it is then parametrize its description.
    \begin{answer}
      No, it is not closed.
      In particular, it is not closed under scalar multiplication because it
      does not contain the zero polynomial.  
    \end{answer}
  \item Is the vector in the span of the set?
    \begin{equation*}
      \colvec{1 \\ 0 \\ 3}
      \quad
      \set{\colvec{2 \\ 1 \\ -1},
           \colvec{1 \\ -1 \\ 1}}
    \end{equation*}
    \begin{answer}
      The equation
      \begin{equation*}
        \colvec{1 \\ 0 \\ 3} = 
                 c_1\colvec{2 \\ 1 \\ -1}
                 +c_2\colvec{1 \\ -1 \\ 1}
      \end{equation*}
      gives rise to a linear system
      \begin{equation*}
        \begin{amat}{2}
          2  &1  &1  \\
          1  &-1 &0  \\
          -1 &1  &3
        \end{amat}
        \grstep[(1/2)\rho_1+\rho_3]{(-1/2)\rho_1+\rho_2}
        \begin{amat}{2}
          2  &1    &1  \\
          0  &-3/2 &-1/2  \\
          0  &0  &3
        \end{amat}
      \end{equation*}
      that has no solution, so the vector is not in the span.   
    \end{answer}
  \recommended \item 
    Decide if the vector lies in the span of the set, inside of the
    space.
    \begin{exparts}
      \partsitem \( \colvec[r]{2 \\ 0 \\ 1} \),
        \( \set{\colvec[r]{1 \\ 0 \\ 0},
                \colvec[r]{0 \\ 0 \\ 1}  } \),
        in \( \Re^3 \)
      \partsitem \( x-x^3 \),
        \( \set{x^2,2x+x^2,x+x^3} \),
        in \( \polyspace_3 \)
      \partsitem \( \begin{mat}[r]
                 0  &1  \\
                 4  &2
               \end{mat}  \),
        \( \set{\begin{mat}[r]
                  1  &0  \\
                  1  &1
                \end{mat},
                \begin{mat}[r]
                  2  &0  \\
                  2  &3
                \end{mat}  } \),
        in \( \matspace_{\nbyn{2}} \)
    \end{exparts}
    \begin{answer}
      \begin{exparts}
         \partsitem Yes, solving the linear system arising from
           \begin{equation*}
             r_1\colvec[r]{1 \\ 0 \\ 0}+r_2\colvec[r]{0 \\ 0 \\ 1}
               =\colvec[r]{2 \\ 0 \\ 1}
           \end{equation*}
           gives \( r_1=2 \) and \( r_2=1 \).
         \partsitem Yes; the linear system arising from
           \( r_1(x^2)+r_2(2x+x^2)+r_3(x+x^3)=x-x^3 \)
           \begin{equation*}
             \begin{linsys}{3}
                   &  &2r_2 &+ &r_3 &= &1  \\
               r_1 &+ &r_2  &  &    &= &0  \\
                   &  &     &  &r_3 &= &-1   
             \end{linsys}
           \end{equation*}
           gives that \( -1(x^2)+1(2x+x^2)-1(x+x^3)=x-x^3 \).
        \partsitem No; any combination of the two given matrices has a zero
           in the upper right.
      \end{exparts}  
    \end{answer}
  \item % \cite{Cleary} credit removed at request Cleary 2022-Sep-13
    A superhero is at the origin of a two dimensional plane.  
    The superhero has two devices, a hoverboard that moves any distance in the 
    direction $\binom{3}{1}$ 
    and a magic carpet that moves any distance in the direction 
    $\binom{1}{2}$.
    \begin{exparts}  
      \partsitem An evil villain is hiding out in the plane at the point 
        $(-5,7)$.  
        How many hoverboard units and magic carpet units does the superhero 
        have to move to get to the villain?
     \partsitem Is there anywhere in the plane that the villain could safely 
        hide and not be reached?  
        If so, give one such location. 
        If not, explain why not.
     \partsitem The superhero and the villain are transported to a three 
        dimensional space where the superhero now has three devices.
        \begin{equation*}
          \text{hoverboard: }\colvec{-1 \\ 0 \\ 3}
          \quad
          \text{magic carpet: }\colvec{2 \\ 1 \\ 0}
          \quad
          \text{scooter: }\colvec{5 \\ 4 \\ -9}
        \end{equation*}
        Is there anywhere that the villain could safely hide?  
        If so, give one such location and if not, explain why not.
    \end{exparts}
    \begin{answer}
      \begin{exparts}
        \partsitem This relationship
          \begin{equation*}
            \colvec{-5 \\ 7}=a\colvec{3 \\ 1}+b\colvec{1 \\ 2}
          \end{equation*}
          gives this system of equations.
          \begin{equation*}
            \begin{linsys}{2}
              3a &+ &b  &= &-5 \\
               a &+ &2b &= &7
            \end{linsys}
            \grstep{-(1/3)\rho_1+\rho_2}
            \begin{linsys}{2}
              3a &+ &b      &= &-5 \\
                 &  &(5/3)b &= &26/3
            \end{linsys}
          \end{equation*}
          So $b=26/5$.
          Substituting $3a+(26/5)=-5$ gives $a=-17/5$.
          In the terms of the section, the vector $\binom{-5}{7}$
          is in the span of the set of the other two. 
        \partsitem There is no place to hide.
          This Gauss's method reduction
          \begin{equation*}
            \begin{linsys}{2}
              3a &+ &b  &= &x \\
               a &+ &2b &= &y
            \end{linsys}
            \grstep{-(1/3)\rho_1+\rho_2}
            \begin{linsys}{2}
              3a &+ &b      &= &x\hfill\hbox{} \\
                 &  &(5/3)b &= &-(1/3)x+y
            \end{linsys}
          \end{equation*}
          shows that for any $x,y\in\R$ there is a pair $a,b\in\R$.
          In the terms of the section, the span of the
          two vectors is the entire plane.
        \partsitem No place to hide. 
          Consider this system.
          \begin{equation*}
            a\colvec{-1 \\ 0 \\ 3}+b\colvec{2 \\ 1 \\ 0}+c\colvec{5 \\ 4 \\ -9}
                = \colvec{x \\ y \\ z}
          \end{equation*}
          After reduction
          \begin{align*}
            \begin{linsys}{3}
              -a &+ &2b &+ &5c &= &x \\
                 &  &b  &+ &4c &= &y \\
              3a &  &   &- &9c &= &z \\
            \end{linsys}
            & \grstep{3\rho_1+\rho_3}
            \begin{linsys}{3}
              -a &+ &2b  &+ &5c &= &x\hfill\hbox{} \\
                 &  &b   &+ &4c &= &y\hfill\hbox{} \\
                 &  &6b  &+ &6c &= &3x+z \\
            \end{linsys}                                  \\
            & \grstep{-6\rho_2+\rho_3}
            \begin{linsys}{3}
              -a &+ &2b  &+ &5c  &= &x\hfill\hbox{} \\
                 &  &b   &+ &4c  &= &y\hfill\hbox{} \\
                 &  &    & &-18c &= &3x+-6y+z \\
            \end{linsys}
          \end{align*}
          it shows that for any $x,y,z\in\R$ 
          the required $a$, $b$, and~$c$ exist.
          The superhero's devices span the plane.
      \end{exparts}
    \end{answer}
  \item 
    Which of these are members of the span
    \( \spanof{\set{\cos^2x,\sin^2x} } \)
    in the vector space of real-valued functions of one real variable?
    \begin{exparts*}
      \partsitem \( f(x)=1 \)
      \partsitem \( f(x)=3+x^2 \)
      \partsitem \( f(x)=\sin x \)
      \partsitem \( f(x)=\cos (2x) \)
    \end{exparts*}
    \begin{answer}
      \begin{exparts}
        \partsitem Yes; it is in that span since 
          \( 1\cdot\cos^2x+1\cdot\sin^2x=f(x) \).
        \partsitem No, since \( r_1\cos^2x+r_2\sin^2x=3+x^2 \) has no scalar
          solutions that work for all \( x \).
          For instance, setting $x$ to be $0$ and $\pi$ gives the two
          equations $r_1\cdot 1+r_2\cdot 0=3$ and 
          $r_1\cdot 1+r_2\cdot 0=3+\pi^2$, which are not consistent with each
          other. 
        \partsitem No; consider what happens on setting $x$ to be $\pi/2$ and
          $3\pi/2$.
        \partsitem Yes, \( \cos (2x)=1\cdot\cos^2(x)-1\cdot\sin^2(x) \).
      \end{exparts}  
     \end{answer}
  \recommended \item 
    Which of these sets spans \( \Re^3 \)?
    That is, which of these sets has the property that any three-tall
    vector can be expressed as a suitable linear combination of the
    set's elements?
    \begin{exparts*}
      \partsitem \( \set{ \colvec[r]{1 \\ 0 \\ 0},
               \colvec[r]{0 \\ 2 \\ 0},
               \colvec[r]{0 \\ 0 \\ 3}  } \)
      \partsitem \( \set{ \colvec[r]{2 \\ 0 \\ 1},
               \colvec[r]{1 \\ 1 \\ 0},
               \colvec[r]{0 \\ 0 \\ 1}  } \)
      \partsitem \( \set{ \colvec[r]{1 \\ 1 \\ 0},
               \colvec[r]{3 \\ 0 \\ 0}  } \)
      \partsitem \( \set{ \colvec[r]{1 \\ 0 \\ 1},
               \colvec[r]{3 \\ 1 \\ 0},
               \colvec[r]{-1 \\ 0 \\ 0},
               \colvec[r]{2 \\ 1 \\ 5}  } \)
      \partsitem \( \set{ \colvec[r]{2 \\ 1 \\ 1},
               \colvec[r]{3 \\ 0 \\ 1},
               \colvec[r]{5 \\ 1 \\ 2},
               \colvec[r]{6 \\ 0 \\ 2}  } \)
    \end{exparts*}
    \begin{answer}
      \begin{exparts}
         \partsitem Yes, for any \( x,y,z\in\Re \) this equation
           \begin{equation*}
              r_1\colvec[r]{1 \\ 0 \\ 0}
              +r_2\colvec[r]{0 \\ 2 \\ 0}
              +r_3\colvec[r]{0 \\ 0 \\ 3}
              =\colvec{x \\ y \\ z}
           \end{equation*}
           has the solution \( r_1=x \), \( r_2=y/2 \), and
           \( r_3=z/3 \).
         \partsitem Yes, the equation
           \begin{equation*}
             r_1\colvec[r]{2 \\ 0 \\ 1}
             +r_2\colvec[r]{1 \\ 1 \\ 0}
             +r_3\colvec[r]{0 \\ 0 \\ 1}
             =\colvec{x \\ y \\ z}
           \end{equation*}
           gives rise to this 
           \begin{equation*}
             \begin{linsys}{3}
               2r_1 &+  &r_2  &  &    &=  &x \\
                    &   &r_2  &  &    &=  &y \\
                r_1 &   &     &+ &r_3 &=  &z \\
             \end{linsys}
             \grstep{-(1/2)\rho_1+\rho_3}\repeatedgrstep{(1/2)\rho_2+\rho_3}
             \begin{linsys}{3}
               2r_1 &+  &r_2  &  &    &=  &x\hfill\hbox{} \\
                    &   &r_2  &  &    &=  &y\hfill\hbox{} \\
                    &   &     &  &r_3 &=  &-(1/2)x+(1/2)y+z \\
             \end{linsys}
           \end{equation*}
           so that, given any $x$, $y$, and $z$, we can compute that
           \( r_3=(-1/2)x+(1/2)y+z \), \( r_2=y \), and
           \( r_1=(1/2)x-(1/2)y \).
        \partsitem No.
           In particular, we cannot get the vector
           \begin{equation*}
             \colvec[r]{0 \\ 0 \\ 1}
           \end{equation*}
           as a linear combination since the two given
           vectors both have a third component of zero.
       \partsitem Yes.
         The equation
         \begin{equation*}
           r_1\colvec[r]{1 \\ 0 \\ 1}
           +r_2\colvec[r]{3 \\ 1 \\ 0}
           +r_3\colvec[r]{-1\\ 0 \\ 0}
           +r_4\colvec[r]{2 \\ 1 \\ 5}
           =\colvec{x \\ y \\ z}
         \end{equation*}
         leads to this reduction.
         \begin{equation*}
           \begin{amat}{4}
             1  &3  &-1  &2  &x  \\
             0  &1  &0   &1  &y  \\
             1  &0  &0   &5  &z  
           \end{amat}                                              
           \grstep{-\rho_1+\rho_3}\repeatedgrstep{3\rho_2+\rho_3}
           \begin{amat}{4}
             1  &3  &-1  &2  &x \\
             0  &1  &0   &1  &y  \\
             0  &0  &1   &6  &-x+3y+z
           \end{amat}
         \end{equation*}
         We have infinitely many solutions.
         We can, for example, set $r_4$ to be zero and solve for
         $r_3$, $r_2$, and $r_1$ in terms of $x$, $y$, and $z$ by the usual
         methods of back-substitution.
       \partsitem No.
         The equation
         \begin{equation*}
           r_1\colvec[r]{2 \\ 1 \\ 1}
           +r_2\colvec[r]{3 \\ 0 \\ 1}
           +r_3\colvec[r]{5 \\ 1 \\ 2}
           +r_4\colvec[r]{6 \\ 0 \\ 2}
           =\colvec{x \\ y \\ z}
         \end{equation*}
         leads to this reduction.
         \begin{multline*}
           \begin{amat}{4}
             2  &3  &5   &6  &x  \\
             1  &0  &1   &0  &y  \\
             1  &1  &2   &2  &z  
           \end{amat}                                             \\
           \grstep[-(1/2)\rho_1+\rho_3]{-(1/2)\rho_1+\rho_2}
           \repeatedgrstep{-(1/3)\rho_2+\rho_3}
           \begin{amat}{4}
             2  &3     &5     &6  &x \\
             0  &-3/2  &-3/2  &-3 &-(1/2)x+y  \\
             0  &0     &0     &0  &-(1/3)x-(1/3)y+z
           \end{amat}
         \end{multline*}
         This shows that not every three-tall vector can be so expressed.
         Only the vectors satisfying the restriction that
         $-(1/3)x-(1/3)y+z=0$ are in the span.
         (To see that any such vector is indeed expressible, 
         take $r_3$ and $r_4$
         to be zero and solve for $r_1$ and $r_2$ in terms of $x$, $y$, and
         $z$ by back-substitution.)
      \end{exparts}  
    \end{answer}
  \recommended \item 
    Parametrize each subspace's description.
    Then express each subspace as a span.
    \begin{exparts}
      \partsitem  The subset \( \set{\rowvec{a &b &c}\suchthat a-c=0}   \)
        of the three-wide row vectors
      \partsitem This subset of \( \matspace_{\nbyn{2}} \)
        \begin{equation*}
          \set{\begin{mat}
                 a  &b  \\
                 c  &d
               \end{mat}  \suchthat a+d=0}
        \end{equation*}
      \partsitem This subset of \( \matspace_{\nbyn{2}} \)
        \begin{equation*}
          \set{\begin{mat}
                 a  &b  \\
                 c  &d
               \end{mat}  \suchthat \text{\( 2a-c-d=0 \) 
                                                 and \( a+3b=0 \)} }
        \end{equation*}
      \partsitem The subset \( \set{a+bx+cx^3\suchthat a-2b+c=0} \) of
        \( \polyspace_3 \)
      \partsitem The subset of \( \polyspace_2 \) of quadratic polynomials 
        \( p \) such that \( p(7)=0 \)
    \end{exparts}
    \begin{answer}
      \begin{exparts}
        \partsitem \( \set{\rowvec{c &b &c}\suchthat b,c\in\Re}
                      =\set{b\rowvec{0 &1 &0}+c\rowvec{1 &0 &1}
                        \suchthat b,c\in\Re} \)
           The obvious choice for the set that spans is 
           $\set{\rowvec{0 &1 &0},\rowvec{1 &0 &1}}$.
        \partsitem \( \set{\begin{mat}
                       -d &b  \\
                       c  &d
                      \end{mat} \suchthat b,c,d\in\Re}
                     =\set{b\begin{mat}[r]
                       0  &1  \\
                       0  &0
                      \end{mat}
                     +c\begin{mat}[r]
                       0  &0  \\
                       1  &0
                      \end{mat}
                     +d\begin{mat}[r]
                       -1  &0  \\
                       0  &1
                      \end{mat}  \suchthat b,c,d\in\Re} \)
            One set that spans this space consists of those three matrices. 
        \partsitem The system
          \begin{equation*}
            \begin{linsys}{4}
              a  &+  &3b  &   &   &  &  &=  &0  \\
             2a  &   &    &   &-c &- &d &=  &0  
            \end{linsys}
          \end{equation*}
          gives \( b=-(c+d)/6 \) and \( a=(c+d)/2 \).
          So one description is this.
          \begin{equation*}
            \set{c\begin{mat}[r]
                       1/2  &-1/6  \\
                       1    &0
                      \end{mat}
                     +d\begin{mat}[r]
                       1/2  &-1/6  \\
                       0    &1
                      \end{mat}  \suchthat c,d\in\Re}
           \end{equation*}
          That shows that a set spanning this subspace consists of those
          two matrices.
        \partsitem The $a=2b-c$ gives that the set
           \( \set{(2b-c)+bx+cx^3 \suchthat b,c\in\Re} \)
           equals the set
           \( \set{b(2+x)+c(-1+x^3) \suchthat b,c\in\Re}  \).
           So the subspace is the span of the set $\set{2+x, -1+x^3}$.
        \partsitem The set
          \( \set{a+bx+cx^2\suchthat a+7b+49c=0} \)
          can be parametrized as
          \begin{equation*}
            \set{b(-7+x)+c(-49+x^2)\suchthat b,c\in\Re} 
          \end{equation*}
          and so 
          has the spanning set $\set{-7+x,-49+x^2}$.
      \end{exparts}  
    \end{answer}
  \recommended \item 
    Find a set to span the given subspace of the given space.
    (\textit{Hint.}   Parametrize each.)
    \begin{exparts}
      \partsitem  the \( xz \)-plane in \( \Re^3 \)
      \partsitem \( \set{\colvec{x \\ y \\ z}\suchthat 3x+2y+z=0} \)
            in \( \Re^3 \)
      \partsitem \( \set{\colvec{x \\ y \\ z \\ w}\suchthat
                       2x+y+w=0 \text{\ and\ } y+2z=0} \)
            in \( \Re^4 \)
      \partsitem \( \set{a_0+a_1x+a_2x^2+a_3x^3\suchthat
                        a_0+a_1=0 \text{\ and\ } a_2-a_3=0} \)
            in \( \polyspace_3 \)
      \partsitem The set \( \polyspace_4 \) in the space \( \polyspace_4 \)
      \partsitem \( \matspace_{\nbyn{2}} \) in \( \matspace_{\nbyn{2}} \)
    \end{exparts}
    \begin{answer}
      \begin{exparts}
        \partsitem We can parametrize in this way 
          \begin{equation*}
             \set{\colvec{x \\ 0 \\ z}\suchthat x,z\in\Re}
             =\set{x\colvec[r]{1 \\ 0 \\ 0}
                  +z\colvec[r]{0 \\ 0 \\ 1}\suchthat x,z\in\Re}
          \end{equation*}
          giving this for a spanning set.
          \begin{equation*}
             \set{\colvec[r]{1 \\ 0 \\ 0},\colvec[r]{0 \\ 0 \\ 1}} 
          \end{equation*}
        \item Here is a parametrization, and the associated spanning set.
          \begin{equation*}
             \set{y\colvec[r]{-2/3 \\ 1 \\ 0}+z\colvec[r]{-1/3 \\ 0 \\ 1}
                 \suchthat y,z\in\Re }
             \qquad
             \set{\colvec[r]{-2/3 \\ 1 \\ 0},\colvec[r]{-1/3 \\ 0 \\ 1} } 
          \end{equation*}
        \partsitem \( \set{\colvec[r]{1 \\ -2 \\ 1 \\ 0},
                      \colvec[r]{-1/2 \\ 0 \\ 0 \\ 1} } \)
        \partsitem Parametrize the description as
          \( \set{-a_1+a_1x+a_3x^2+a_3x^3\suchthat a_1,a_3\in\Re } \)
          to get \( \set{-1+x,x^2+x^3}. \)
        \partsitem \( \set{1,x,x^2,x^3,x^4} \)
        \partsitem \( \set{ \begin{mat}[r]
                   1  &0  \\
                   0  &0
                 \end{mat},
                 \begin{mat}[r]
                   0  &1  \\
                   0  &0
                 \end{mat},
                 \begin{mat}[r]
                   0  &0  \\
                   1  &0
                 \end{mat},
                 \begin{mat}[r]
                   0  &0  \\
                   0  &1
                 \end{mat} } \)
      \end{exparts}  
    \end{answer}
  \item 
    Is \( \Re^2 \) a subspace of \( \Re^3 \)?
    \begin{answer}
      Technically, no.
      Subspaces of \( \Re^3 \) are sets of three-tall vectors, while
      \( \Re^2 \) is a set of two-tall vectors.
      Clearly though, \( \Re^2 \) is ``just like'' this subspace of 
      \( \Re^3 \).
      \begin{equation*}
        \set{\colvec{x \\ y \\ 0}\suchthat x,y\in\Re}
      \end{equation*}  
    \end{answer}
  \recommended \item 
    Decide if each is a subspace of the vector space of real-valued
    functions of one real variable.
    \begin{exparts}
      \partsitem The 
        \definend{even}\index{function!even}\index{even functions} functions
        \( \set{\map{f}{\Re}{\Re} \suchthat f(-x)=f(x) \text{ for all } x} \).
        For example, two members of this set are $f_1(x)=x^2$ 
        and $f_2(x)=\cos (x)$.
      \partsitem The \definend{odd}\index{function!odd}\index{odd function}
        functions
        \( \set{\map{f}{\Re}{\Re} \suchthat f(-x)=-f(x) \text{ for all } x} \).
        Two members are $f_3(x)=x^3$ and $f_4(x)=\sin(x)$.
    \end{exparts}
    \begin{answer}
      Of course, the addition and scalar multiplication operations are the
      ones inherited from the enclosing space.
      \begin{exparts}
        \partsitem This is a subspace.
          It is not empty as it contains at least the two example functions
          given.
          It is closed because if \( f_1,f_2 \) are even and
          \( c_1,c_2 \) are scalars then we have this.
          \begin{equation*}
            (c_1f_1+c_2f_2)\,(-x)
            =c_1\,f_1(-x)+c_2\,f_2(-x)
            =c_1\,f_1(x)+c_2\,f_2(x)
            =(c_1f_1+c_2f_2)\,(x)
          \end{equation*}
        \partsitem This is also a subspace; the check is similar to
          the prior one.
      \end{exparts}  
    \end{answer}
  \item 
    \nearbyexample{ex:SpanSingVec} says that for any vector $\vec{v}$ 
    that is an element of 
    a vector space $V$, the set $\set{r\cdot\vec{v}\suchthat r\in\Re}$
    is a subspace of $V$.
    (This is simply the span\index{span!of a singleton} 
    of the singleton set $\set{\vec{v}}$.)
    Must any such subspace be a proper subspace?
    \begin{answer}
      It can be improper.
      For instance,
      if the vector space is \( \Re^1 \) and \( \vec{v}\neq\zero\, \)
      then the subspace $\set{r\cdot\vec{v}\suchthat r\in\Re}$
      is all of $\Re^1$.  
    \end{answer}
  \item 
    An example following the definition of a vector space shows that the
    solution set of a homogeneous linear system is a vector space.
    In the terminology of this subsection, it is a subspace of $\Re^n$ where
    the system has $n$ variables.
    What about a non-homogeneous linear system; do its solutions form a 
    subspace (under the inherited operations)?
    \begin{answer}
      No, such a set is not closed.
      For one thing, it does not contain the zero vector.  
    \end{answer}
  \item \cite{Cleary} 
   Give an example of each or explain why it would be impossible 
   to do so.
   \begin{exparts}
     \item A nonempty subset of $\matspace_{\nbyn{2}}$ that is
       not a subspace.
     \item A set of two vectors in $\Re^2$ that does not span the space.
   \end{exparts}
   \begin{answer}
     \begin{exparts}
       \item This nonempty subset of $\matspace_{\nbyn{2}}$ is not a subspace.
         \begin{equation*}
           A=\set{
             \begin{mat}[r]
               1 &2 \\
               3 &4
             \end{mat},
             \begin{mat}[r]
               5 &6 \\
               7 &8
             \end{mat}}
         \end{equation*}
         One reason that it is not a subspace of $\matspace_{\nbyn{2}}$ is that 
         it does not contain the zero matrix.
         (Another reason is that it is not closed under addition, since the sum
         of the two is not an element of $A$.
         It is also not closed under scalar multiplication.)
        \item This set of two vectors does not span $\Re^2$.
          \begin{equation*}
            \set{\colvec[r]{1 \\ 1},\colvec[r]{3 \\ 3}}
          \end{equation*}
          No linear combination of these two can give
          a vector whose second component is unequal to its first component.
      \end{exparts}
   \end{answer}
  \item 
    \nearbyexample{ex:SubspRThree} shows that 
    $\Re^3$ has infinitely many subspaces.
    Does every nontrivial space have infinitely many subspaces?
    \begin{answer}
      No.
      The only subspaces of \( \Re^1 \) are the space itself and its 
      trivial subspace.
      Any subspace $S$ of $\Re$ that contains a nonzero member $\vec{v}$ 
      must contain the set of all of its scalar multiples 
      $\set{r\cdot\vec{v}\suchthat r\in\Re}$. 
      But this set is all of $\Re$.  
    \end{answer}
  \item \label{exer:SubspIffClosed}
    Finish the proof of \nearbylemma{th:SubspIffClosed}.
    \begin{answer}
      Item~(1) is checked in the text.

      Item~(2) has five conditions.
      First, for closure, if \( c\in\Re \) and \( \vec{s}\in S \) then
      \( c\cdot\vec{s}\in S \) as 
      \( c\cdot\vec{s}=c\cdot\vec{s}+0\cdot\zero \).
      Second, because the operations in \( S \) are inherited from \( V \),
      for \( c,d\in\Re \) and \( \vec{s}\in S \), the scalar product
      \( (c+d)\cdot\vec{s}\, \) in \( S \) equals the product
      \( (c+d)\cdot\vec{s}\, \) in \( V \), and that equals
      \( c\cdot\vec{s}+d\cdot\vec{s}\, \) in \( V \), which equals
      \( c\cdot\vec{s}+d\cdot\vec{s}\, \) in \( S \).

      The check for the third, fourth, and fifth conditions are similar to the
      second condition's check just given.  
    \end{answer}
  \item 
    Show that each vector space has only one trivial subspace.
    \begin{answer}
      An exercise in the prior subsection shows that every vector space
      has only one zero vector (that is, there is only one vector that is the
      additive identity element of the space).
      But a trivial space has only one element and that element must be this
      (unique) zero vector.
    \end{answer}
  \item 
    Show that for any subset \( S \) of a vector space,
    the span of the span equals the span 
    \( \spanof{ \spanof{S} }=\spanof{S} \).
    (\textit{Hint.} 
    Members of $\spanof{S}$ are linear combinations of members of $S$.
    Members of $\spanof{\spanof{S}}$ are linear combinations of 
    linear combinations of members of $S$.)
    \begin{answer}
      As the hint suggests, the basic reason is the Linear Combination Lemma
      from the first chapter.
      For the full proof, we will show mutual containment between the two sets.

      The first containment \( \spanof{ \spanof{S} }\supseteq\spanof{S}  \)
      is an instance of the more general, and obvious, fact that for any 
      subset \( T \) of a vector space, \( \spanof{T}\supseteq T \).

      For the other containment, 
      that \( \spanof{ \spanof{S} }\subseteq\spanof{S} \),
      take $m$ vectors from \( \spanof{S} \), namely
      \( c_{1,1}\vec{s}_{1,1}+\cdots+c_{1,n_1}\vec{s}_{1,n_1} \), \ldots,
      \( c_{1,m}\vec{s}_{1,m}+\cdots+c_{1,n_m}\vec{s}_{1,n_m} \),
      and note that any linear combination of those
      \begin{equation*}
         r_1(c_{1,1}\vec{s}_{1,1}+\cdots+c_{1,n_1}\vec{s}_{1,n_1})+\cdots
        +r_m(c_{1,m}\vec{s}_{1,m}+\cdots+c_{1,n_m}\vec{s}_{1,n_m})
      \end{equation*}
      is a linear combination of elements of \( S \)
      \begin{equation*}
        = (r_1c_{1,1})\vec{s}_{1,1}+\cdots+(r_1c_{1,n_1})\vec{s}_{1,n_1}+\cdots
        +(r_mc_{1,m})\vec{s}_{1,m}+\cdots+(r_mc_{1,n_m})\vec{s}_{1,n_m}
      \end{equation*}
      and so is in \( \spanof{S} \).  
      That is, simply recall that a linear combination of linear combinations
      (of members of $S$)
      is a linear combination (again of members of $S$).
    \end{answer}
  \item 
    All of the subspaces that we've seen in some way use zero in their 
    description.
    For example, the subspace in \nearbyexample{ex:SubspacesRTwo} consists of
    all the vectors from $\Re^2$ with a second component of zero.
    In contrast,
    the collection of vectors from $\Re^2$ with a second component of one
    does not form a subspace (it is not closed under scalar multiplication).
    Another example is \nearbyexample{ex:PlaneSubspRThree}, where the condition
    on the vectors is that the three components add to zero.
    If the condition there 
    were that the three components add to one then it would 
    not be a subspace (again, it would fail to be closed).
    However, a reliance on zero is not strictly necessary.
    Consider the set 
    \begin{equation*}
       \set{\colvec{x \\ y \\ z}\suchthat x+y+z=1}
    \end{equation*}
    under these operations.
    \begin{equation*}
       \colvec{x_1 \\ y_1 \\ z_1}+\colvec{x_2 \\ y_2 \\ z_2}
       =\colvec{x_1+x_2-1 \\ y_1+y_2 \\ z_1+z_2}
       \qquad
       r\colvec{x \\ y \\ z}=\colvec{rx-r+1 \\ ry \\ rz}
    \end{equation*}
    \begin{exparts}
      \partsitem Show that it is not a subspace of $\Re^3$.
         (\textit{Hint.}   See \nearbyexample{ex:OperNotInherit}).
      \partsitem Show that it is a vector space.         
         Note that by the prior item,
         \nearbylemma{th:SubspIffClosed} can not apply.  
      \partsitem Show that any subspace of $\Re^3$ must pass through the origin,
         and so any subspace of $\Re^3$ must involve zero in its description.
         Does the converse hold?  
          Does any subset of $\Re^3$ that contains the origin become a
          subspace when given the inherited operations?
    \end{exparts}
    \begin{answer}
      \begin{exparts}
         \partsitem It is not a subspace because these are not the inherited 
           operations.  
           For one thing, in this space,
           \begin{equation*}
             0\cdot\colvec{x \\ y \\ z}=\colvec[r]{1 \\ 0 \\ 0}
           \end{equation*}
           while this does not, of course, hold in $\Re^3$.
         \partsitem We can combine the argument showing closure under
           addition with the argument showing closure under 
           scalar multiplication into one single argument
           showing closure under linear combinations of two vectors.
           If $r_1,r_2,x_1,x_2,y_1,y_2,z_1,z_2$ are in $\Re$ then      
           \begin{multline*}
              r_1\colvec{x_1 \\ y_1 \\ z_1}
              +r_2\colvec{x_2 \\ y_2 \\ z_2}
              =\colvec{r_1x_1-r_1+1 \\ r_1y_1 \\ r_1z_1}
               +\colvec{r_2x_2-r_2+1 \\ r_2y_2 \\ r_2z_2}                 \\
            =\colvec{r_1x_1-r_1+r_2x_2-r_2+1 \\ r_1y_1+r_2y_2 \\ r_1z_1+r_2z_2}
           \end{multline*} 
           (note that the definition of addition in this space is that
           the first
           components combine as $(r_1x_1-r_1+1)+(r_2x_2-r_2+1)-1$,
           so the first component of the last vector does not say
           `$\hbox{}+2$').
           Adding the three components of the last vector gives
           $r_1(x_1-1+y_1+z_1)+r_2(x_2-1+y_2+z_2)+1=r_1\cdot0+r_2\cdot0+1=1$.

           Most of the other checks of the conditions are easy (although the
           oddness of the operations keeps them from being routine).
           Commutativity of addition goes like this.
           \begin{equation*}
             \colvec{x_1 \\ y_1 \\ z_1}+\colvec{x_2 \\ y_2 \\ z_2}
             =\colvec{x_1+x_2-1 \\ y_1+y_2 \\ z_1+z_2}
             =\colvec{x_2+x_1-1 \\ y_2+y_1 \\ z_2+z_1}
             =\colvec{x_2 \\ y_2 \\ z_2}+\colvec{x_1 \\ y_1 \\ z_1}
           \end{equation*}
           Associativity of addition has
           \begin{equation*}
             (\colvec{x_1 \\ y_1 \\ z_1}+\colvec{x_2 \\ y_2 \\ z_2})
              +\colvec{x_3 \\ y_3 \\ z_3}
             =\colvec{(x_1+x_2-1)+x_3-1 \\ (y_1+y_2)+y_3 \\ (z_1+z_2)+z_3}
           \end{equation*}
           while
           \begin{equation*}
             \colvec{x_1 \\ y_1 \\ z_1}
             +(\colvec{x_2 \\ y_2 \\ z_2}+\colvec{x_3 \\ y_3 \\ z_3})
             =\colvec{x_1+(x_2+x_3-1)-1 \\ y_1+(y_2+y_3) \\ z_1+(z_2+z_3)}
           \end{equation*}
           and they are equal.
           The identity element with respect to this addition operation 
           works this way
           \begin{equation*}
             \colvec{x \\ y \\ z}+\colvec[r]{1 \\ 0 \\ 0}
             =\colvec{x+1-1 \\ y+0 \\ z+0}
             =\colvec{x \\ y \\ z}
           \end{equation*}
           and the additive inverse is similar.
           \begin{equation*}
             \colvec{x \\ y \\ z}+\colvec{-x+2 \\ -y \\ -z}
             =\colvec{x+(-x+2)-1 \\ y-y \\ z-z}
             =\colvec[r]{1 \\ 0 \\ 0}
           \end{equation*}

           The conditions on scalar multiplication are also easy.
           For the first condition,
           \begin{equation*}
             (r+s)\colvec{x \\ y \\ z}
             =\colvec{(r+s)x-(r+s)+1 \\ (r+s)y \\ (r+s)z}
           \end{equation*}
           while
           \begin{multline*}
             r\colvec{x \\ y \\ z}+s\colvec{x \\ y \\ z}
             =\colvec{rx-r+1 \\ ry \\ rz}+\colvec{sx-s+1 \\ sy \\ sz}  \\
             =\colvec{(rx-r+1)+(sx-s+1)-1 \\ ry+sy \\ rz+sz}
           \end{multline*}
           and the two are equal.
           The second condition compares
           \begin{equation*}
             r\cdot(\colvec{x_1 \\ y_1 \\ z_1}+\colvec{x_2 \\ y_2 \\ z_2})
             =r\cdot\colvec{x_1+x_2-1 \\ y_1+y_2 \\ z_1+z_2}
             =\colvec{r(x_1+x_2-1)-r+1 \\ r(y_1+y_2) \\ r(z_1+z_2)}
           \end{equation*}
           with
           \begin{multline*}
             r\colvec{x_1 \\ y_1 \\ z_1}+r\colvec{x_2 \\ y_2 \\ z_2}
             =\colvec{rx_1-r+1 \\ ry_1 \\ rz_1}
                     +\colvec{rx_2-r+1 \\ ry_2 \\ rz_2}                  \\
             =\colvec{(rx_1-r+1)+(rx_2-r+1)-1 \\ ry_1+ry_2 \\ rz_1+rz_2}
           \end{multline*}
           and they are equal.
           For the third condition,
           \begin{equation*}
             (rs)\colvec{x \\ y \\ z}
             =\colvec{rsx-rs+1 \\ rsy \\ rsz}
           \end{equation*}
           while
           \begin{equation*}
             r(s\colvec{x \\ y \\ z})
             =r(\colvec{sx-s+1 \\ sy \\ sz})
             =\colvec{r(sx-s+1)-r+1 \\ rsy \\ rsz}
           \end{equation*}
           and the two are equal.
           For scalar multiplication by $1$ we have this.
           \begin{equation*}
             1\cdot\colvec{x \\ y \\ z}
             =\colvec{1x-1+1 \\ 1y \\ 1z}
             =\colvec{x \\ y \\ z}
           \end{equation*}
           Thus all the conditions on a vector space are met by these two
           operations.

           \textit{Remark.}
           A way to understand this vector space is to think of it as 
           the plane in $\Re^3$
           \begin{equation*}
             P=\set{\colvec{x \\ y \\ z}\suchthat x+y+z=0}
           \end{equation*}
           displaced away from the origin by $1$ along the $x$-axis.
           Then addition becomes:~to add two members of this space, 
           \begin{equation*}
             \colvec{x_1 \\ y_1 \\ z_1},\;\colvec{x_2 \\ y_2 \\ z_2}
           \end{equation*}
           (such that $x_1+y_1+z_1=1$ and $x_2+y_2+z_2=1$)
           move them back by $1$ to place them in $P$ and
           add as usual,
           \begin{equation*}
             \colvec{x_1-1 \\ y_1 \\ z_1}+\colvec{x_2-1 \\ y_2 \\ z_2}
             =\colvec{x_1+x_2-2 \\ y_1+y_2 \\ z_1+z_2}
             \qquad\text{(in $P$)}
           \end{equation*}
           and then move the result back out by $1$ along the $x$-axis.
           \begin{equation*}
             \colvec{x_1+x_2-1 \\ y_1+y_2 \\ z_1+z_2}.
           \end{equation*}
           Scalar multiplication is similar.
         \partsitem For the subspace to be closed under the inherited scalar 
           multiplication, where $\vec{v}$ is a member of that subspace,
           \begin{equation*}
             0\cdot\vec{v}=\colvec[r]{0 \\ 0 \\ 0}
           \end{equation*}
           must also be a member.

           The converse does not hold.
           Here is a subset of $\Re^3$ that contains the origin 
           \begin{equation*}
             \set{\colvec[r]{0 \\ 0 \\ 0},\colvec[r]{1 \\ 0 \\ 0}}
           \end{equation*}
           (this subset has only two elements) but is not a subspace.
      \end{exparts}
    \end{answer}
  \item 
    We can give a justification for the convention that the sum of 
    zero-many vectors equals the zero vector.
    Consider this sum of three vectors $\vec{v}_1+\vec{v}_2+\vec{v}_3$.
    \begin{exparts}
      \partsitem What is the difference between this sum of three vectors 
        and the sum of the first two of these three?
      \partsitem What is the difference between the prior sum and the sum 
        of just the first one vector?
      \partsitem What should be the difference between the prior sum of 
        one vector and the sum of no vectors?
      \partsitem So what should be the definition of the sum of no vectors?
    \end{exparts}
    \begin{answer}
      \begin{exparts}
        \item $(\vec{v}_1+\vec{v}_2+\vec{v}_3)-(\vec{v}_1+\vec{v}_2)
                 =\vec{v}_3$
        \item $(\vec{v}_1+\vec{v}_2)-(\vec{v}_1)
                 =\vec{v}_2$
        \item Surely, $\vec{v}_1$.
        \item Taking the one-long sum and subtracting gives
          ($\vec{v}_1)-\vec{v}_1=\zero$.
      \end{exparts}
    \end{answer}
  \item 
    Is a space determined by its subspaces?
    That is, if two vector spaces have the same subspaces, must the
    two be equal?
    \begin{answer}
      Yes; any space is a subspace of itself, so each space contains the
      other.  
    \end{answer}
  \item 
     \begin{exparts}
      \partsitem Give a set that is closed under scalar multiplication
        but not addition.
      \partsitem Give a set closed under addition but not scalar
        multiplication.
      \partsitem Give a set closed under neither.
    \end{exparts}
    \begin{answer}
      \begin{exparts}
        \partsitem The union of the \( x \)-axis and the \( y \)-axis
          in \( \Re^2 \) is one.
        \partsitem The set of integers, as a subset of \( \Re^1 \), is one.
        \partsitem The subset \( \set{\vec{v}} \) of \( \Re^2 \) is one,
          where $\vec{v}$ is any nonzero vector.
      \end{exparts}  
     \end{answer}
  \item 
    Show that the span of a set of vectors does not depend on the order in
    which the vectors are listed in that set.
    \begin{answer}
      Because vector space addition is commutative, a reordering of
      summands leaves a linear combination unchanged.  
    \end{answer}
  \item  
    Which trivial subspace is the span of the empty set?
    Is it
    \begin{equation*}
      \set{\colvec[r]{0 \\ 0 \\ 0}}\subseteq \Re^3,
      \quad\text{or}\quad
      \set{0+0x}\subseteq \polyspace_1,
    \end{equation*}
    or some other subspace?
    \begin{answer}
      We always consider that span in the context of an enclosing space.  
    \end{answer}
  \item   
    Show that if a  vector is in the span of a set then adding that
    vector to the set won't make the span any bigger.
    Is that also `only if'?
    \begin{answer}
      It is both `if' and `only if'.
 
      For `if',
      let \( S \) be a subset of a vector space \( V \) and assume
      \( \vec{v}\in S \) satisfies
      \( \vec{v}=c_1\vec{s}_1+\dots+c_n\vec{s}_n \) where
      \( c_1,\ldots,c_n \) are scalars and
      \( \vec{s}_1,\ldots,\vec{s}_n\in S \).
      We must show that \( \spanof{S\union\set{\vec{v}} }=\spanof{S} \).

      Containment one way,
      \( \spanof{S}\subseteq\spanof{S\union\set{\vec{v}} } \) is obvious.
      For the other direction,
      \( \spanof{S\union\set{\vec{v}} }\subseteq\spanof{S} \), note that if a
      vector is in the set on the left then it has the form
      \( d_0\vec{v}+d_1\vec{t}_1+\dots+d_m\vec{t}_m \) where the \( d \)'s are
      scalars and the \( \vec{t}\, \)'s are in \( S \).
      Rewrite that as
      \( d_0(c_1\vec{s}_1+\dots+c_n\vec{s}_n)
      +d_1\vec{t}_1+\cdots+d_m\vec{t}_m \) and note that 
      the result is a member of the span of \( S \).

      The `only if' is clearly true\Dash adding \( \vec{v} \) 
      enlarges the span to
      include at least \( \vec{v} \).
    \end{answer}
  \recommended \item
    Subspaces are subsets and so we naturally consider how `is a subspace of'
    interacts with the usual set operations.
    \begin{exparts}
      \partsitem If \( A,B \) are subspaces of a vector space, must
        their intersection
        \( A\intersection B \) be a subspace?
        Always?  Sometimes?  Never?
      \partsitem Must the union \( A\union B \) be a subspace?
      \partsitem If \( A \) is a subspace of some~$V$, must
        its set complement $V-A$ be a subspace?
    \end{exparts}
    (\textit{Hint.}   Try some test subspaces from 
    \nearbyexample{ex:SubspRThree}.)
    \begin{answer}
      \begin{exparts}
        \partsitem Always.

          Assume that \( A,B \) are subspaces of \( V \).
          Note that 
          their intersection is not empty as both contain the zero vector.
          If \( \vec{w},\vec{s}\in A\intersection B \) and \( r,s \) are
          scalars then \( r\vec{v}+s\vec{w}\in A \) because
          each vector is in \( A \) and so a linear combination is in \( A \),
          and \(r\vec{v}+s\vec{w}\in B \) for the same reason.
          Thus the intersection is closed.
          Now \nearbylemma{th:SubspIffClosed} applies.
        \partsitem Sometimes (more precisely, only if \( A\subseteq B \) or
          \( B\subseteq A \)).

          To see the answer is not `always', take \( V \) to be \( \Re^3 \),
          take \( A \) to be the $x$-axis, and \( B \) to be the
          \( y \)-axis.
          Note that
          \begin{equation*}
            \colvec[r]{1 \\ 0}\in A \text{ and }\colvec[r]{0 \\ 1}\in B
            \quad\text{but}\quad
            \colvec[r]{1 \\ 0}+\colvec[r]{0 \\ 1}\not\in A\union B
          \end{equation*}
          as the sum is in neither \( A \) nor \( B \).

          The answer is not `never' because if \( A\subseteq B \) or
          \( B\subseteq A \) then clearly \( A\union B \) is a subspace.

          To show that \( A\union B \) is a subspace only if one
          subspace contains the other, we assume that \( A\not\subseteq B \)
          and \( B\not\subseteq A \) and prove that 
          the union is not a subspace.
          The assumption that \( A \) is not a subset of \( B \) means that 
          there is an \( \vec{a}\in A \) with \( \vec{a}\not\in B \).
          The other assumption gives a \( \vec{b}\in B \) with
          \( \vec{b}\not\in A \).
          Consider \( \vec{a}+\vec{b} \).
          Note that sum is not an element of \( A \) or else
          \( (\vec{a}+\vec{b})-\vec{a} \) would be in \( A \), which it is not.
          Similarly the sum is not an element of \( B \).
          Hence the sum is not an element of \( A\union B \), and so the union
          is not a subspace.
        \partsitem Never.
          As \( A \) is a subspace it contains the zero vector and therefore
          the set that is $A$'s complement, $V-A$, does not.
          Without the zero vector, the complement cannot be a vector space.
      \end{exparts}  
    \end{answer}
  \item 
    Does the span of a set depend on the enclosing space?
    That is, if \( W \) is a subspace of \( V \) and \( S \) is a subset of
    \( W \) (and so also a subset of \( V \)), might the span of \( S \) in
    \( W \) differ from the span of \( S \) in \( V \)?
    \begin{answer}
      The span of a set does not depend on the enclosing space.
      A linear combination of vectors from \( S \) gives the same sum
      whether we regard the operations as those of \( W \) or as those of
      \( V \), because the operations of \( W \) are inherited from \(  V \).  
    \end{answer}
  \item 
    Is the relation `is a subspace of' transitive?
    That is, if $V$ is a subspace of $W$ and $W$ is a subspace of
    $X$, must $V$ be a subspace of $X$? 
    \begin{answer}
      It is;
      apply \nearbylemma{th:SubspIffClosed}.
      (You must consider the following.
      Suppose \( B \) is a subspace of a vector space \( V \) and suppose
      \( A\subseteq B\subseteq V \) is a subspace.
      From which space does \( A \) inherit its operations?
      The answer is that it doesn't matter\Dash \( A \) will inherit the
      same operations in either case.)  
    \end{answer}
  \item
    Because `span of' is an operation on sets we naturally consider
    how it interacts with the usual set operations.
    \begin{exparts}
      \partsitem If \( S\subseteq T \) are subsets of a vector space, is
        \( \spanof{S}\subseteq\spanof{T} \)?
        Always?  Sometimes?  Never?
      \partsitem If \( S,T \) are subsets of a vector space, is
        \( \spanof{S\union T}=\spanof{S}\union\spanof{T} \)?
      \partsitem If \( S,T \) are subsets of a vector space, is
        \( \spanof{S\intersection T}=\spanof{S}\intersection\spanof{T} \)?
      \partsitem Is the span of the complement equal to the complement of
        the span?
    \end{exparts}
    \begin{answer}
      \begin{exparts}
         \partsitem Always;
           if \( S\subseteq T \) then a linear combination of elements of
           \( S \) is also a linear combination of elements of \( T \).
         \partsitem Sometimes (more precisely, if and only if 
           \( S\subseteq T \) or \( T\subseteq S \)).

           The answer is not `always' as is shown by this example from
           \( \Re^3 \)
           \begin{equation*}
             S=\set{\colvec[r]{1 \\ 0 \\ 0},\colvec[r]{0 \\ 1 \\ 0}},\quad
             T=\set{\colvec[r]{1 \\ 0 \\ 0},\colvec[r]{0 \\ 0 \\ 1}}
           \end{equation*}
           because of this.
           \begin{equation*}
             \colvec[r]{1 \\ 1 \\ 1}\in\spanof{S\union T}
             \qquad
             \colvec[r]{1 \\ 1 \\ 1}\not\in\spanof{S}\union \spanof{T}
           \end{equation*}

           The answer is not `never' because if either set contains the other
           then equality is clear.
           We can
           characterize equality as happening only when either set contains
           the other by assuming \( S\not\subseteq T \) (implying the
           existence of a vector \( \vec{s}\in S \) with 
           \( \vec{s}\not\in T \))
           and \( T\not\subseteq S \) (giving a \( \vec{t}\in T \) with
           \( \vec{t}\not\in S \)), noting
           \( \vec{s}+\vec{t}\in\spanof{S\union T} \),
           and showing that 
           \( \vec{s}+\vec{t}\not\in\spanof{S}\union\spanof{T} \).
         \partsitem Sometimes.

           Clearly
           \( \spanof{S\intersection T}
             \subseteq\spanof{S}\intersection\spanof{T} \)
           because any linear combination of vectors from 
           \( S\intersection T \)
           is a combination of vectors from \( S \) and also a combination of
           vectors from \( T \).

           Containment the other way does not always hold.
           For instance, in \( \Re^2 \), take
           \begin{equation*}
             S=\set{\colvec[r]{1 \\ 0},\colvec[r]{0 \\ 1}},\quad
             T=\set{\colvec[r]{2 \\ 0}}
           \end{equation*}
           so that \( \spanof{S}\intersection\spanof{T} \) is the \( x \)-axis
           but \( \spanof{S\intersection T}  \) is the trivial subspace.

           Characterizing exactly when equality holds is tough.
           Clearly equality holds if either set contains the other, but that is
           not `only if' by this example in \( \Re^3 \).
           \begin{equation*}
             S=\set{\colvec[r]{1 \\ 0 \\ 0},\colvec[r]{0 \\ 1 \\ 0}},
             \quad
             T=\set{\colvec[r]{1 \\ 0 \\ 0},\colvec[r]{0 \\ 0 \\ 1}}
           \end{equation*}
        \partsitem Never, as the span of the complement is a subspace, while
          the complement of the span is not (it does not contain the zero 
          vector).
      \end{exparts}  
     \end{answer}
  % \item 
  %   Reprove \nearbylemma{le:SpanIsASubsp} without doing the
  %   empty set separately.
  %   \begin{answer}
  %     Call the subset \( S \).
  %     By \nearbylemma{th:SubspIffClosed},
  %     we need to check that 
  %     \( \spanof{S} \) is closed under linear combinations.
  %     If \( c_1\vec{s}_1+\dots+c_n\vec{s}_n,
  %       c_{n+1}\vec{s}_{n+1}+\dots+c_m\vec{s}_m\in\spanof{S} \) then
  %     for any \( p,r\in\Re \) we have
  %     \begin{multline*}
  %       p\cdot(c_1\vec{s}_1+\cdots+c_n\vec{s}_n)+
  %            r\cdot(c_{n+1}\vec{s}_{n+1}+\cdots+c_m\vec{s}_m)        \\
  %       =
  %       pc_1\vec{s}_1+\cdots+pc_n\vec{s}_n
  %         +rc_{n+1}\vec{s}_{n+1}+\cdots+rc_m\vec{s}_m
  %     \end{multline*}
  %     which is an element of \( \spanof{S} \).
  %     % (\textit{Remark.}
  %     % If the set $S$ is empty, then that
  %     % `if \ldots\ then \ldots' statement is vacuously true.)  
  %   \end{answer}
  \item Find a structure that is closed under linear combinations, 
    and yet is not a vector space.
    % (\textit{Remark.} 
    % This is a bit of a trick question.)
    \begin{answer}
      For this to happen, one of the conditions giving the sensibleness of the
      addition and scalar multiplication operations must be violated.
      Consider \( \Re^2 \) with these operations.
      \begin{equation*}
        \colvec{x_1 \\ y_1}
        +\colvec{x_2 \\ y_2}
        =\colvec[r]{0 \\ 0}
        \qquad
        r\colvec{x \\ y}
        =\colvec[r]{0 \\ 0}
      \end{equation*}
      The set $\Re^2$ is closed under these operations. 
      But it is not a vector space.
      \begin{equation*}
        1\cdot\colvec[r]{1 \\ 1}
        \neq\colvec[r]{1 \\ 1}
      \end{equation*}  
    \end{answer}
\index{subspace|)}
\index{vector space|)}
\end{exercises}

% END NATIVE SOURCE src/vs/vs1.tex

% BEGIN NATIVE SOURCE src/vs/vs2.tex
% Chapter 2, Section 2 _Linear Algebra_ Jim Hefferon
%  http://joshua.smcvt.edu/linearalgebra
%  2001-Jun-10
\section{Linear Independence}
The prior section shows how to understand a vector space
as a span, 
as an unrestricted linear combination of some of its elements.
For example, the space of linear polynomials $\set{a+bx\suchthat a,b\in\Re}$ 
is spanned by the set $\set{1,x}$.
The prior section also showed that a space can have many sets that span it.
Two more sets that span the space of linear polynomials are 
$\set{1,2x}$ and $\set{1,x,2x}$.

At the end of that section we described some spanning sets as `minimal'
but we never precisely defined that word.
% We could mean one of two things.
We could mean that a spanning set is minimal if it 
contains the smallest number of members of any set with the same span,
so that $\set{1,x,2x}$ is not minimal because it has 
three members while we can give two-element sets spanning the same space.
Or we could mean that a spanning set is minimal when it has no elements 
that we can remove without changing the span.
Under this meaning $\set{1,x,2x}$ is not minimal because 
removing the \( 2x \) to get \( \set{1,x} \) leaves the
span unchanged.

The first sense of minimality appears to be a global requirement, 
in that to check if a spanning set is minimal 
we seemingly must look at all the sets that span 
and find one with the least number of elements. 
The second sense of minimality is local since
we need to look only at the set and consider the
span with and without various elements.
For instance, using the second sense
we could compare the span of $\set{1,x,2x}$ 
with the span of $\set{1,x}$ and
note that $2x$ is a ``repeat'' in that 
its removal doesn't shrink the span.

In this section we will use the second sense of `minimal spanning set'
because of this technical convenience.
However, the most important result of this book is that the two senses 
coincide.
We will prove that in the next section.


 






\subsection{Definition and Examples}

We saw ``repeats'' in the first chapter.
There, Gauss's Method turned them into 
$0=0$ equations.

\begin{example}  \label{ex:StaticsLIAndLD}
Recall the \hyperlink{ex:Statics}{Statics} example from
Chapter One's opening.
We got two balances with the pair of unknown-mass objects, one 
at \( 40 \)~cm and \( 15 \)~cm and another 
at \( -50 \)~cm and \( 25 \)~cm, and
we then computed the value of those masses.
Had we instead gotten the second balance at
\( 20 \)~cm and \( 7.5 \)~cm
then Gauss's Method on the resulting two-equations, two-unknowns system
would not have yielded a solution, it would have yielded a $0=0$ equation 
along with an equation containing a free variable.
Intuitively, the problem is that \( \rowvec{20 &7.5} \) is half of 
$\rowvec{40 &15}$, that is,
$\rowvec{20 &7.5}$ is in the span of the set $\set{\rowvec{40 &15}}$
and so is repeated data.
We would have been trying to 
solve a two-unknowns problem with essentially only one piece of information.
\end{example}
 
We take $\vec{v}$ to be a ``repeat''
of the vectors in a set~$S$ if $\vec{v}\in\spanof{S}$ so
that it depends on, that is, is expressible in terms of, 
elements of the set 
$\vec{v}=c_1\vec{s}_1+\cdots+c_n\vec{s}_n$.

\begin{lemma} \label{lm:AddVecSpanNoGrowIffVecInSpan}
%<*lm:AddVecSpanNoGrowIffVecInSpan>
Where $V$ is a vector space, $S$ is a subset of that space, and $\vec{v}$
is an element of that space, 
$\spanof{S\union\set{\vec{v}}} = \spanof{S}$
if and only if 
$\vec{v}\in\spanof{S}$.
%</lm:AddVecSpanNoGrowIffVecInSpan>
\end{lemma}

\begin{proof}
%<*pf:AddVecSpanNoGrowIffVecInSpan0>
Half of the if and only if is immediate: if $\vec{v}\notin\spanof{S}$ then 
the sets are not equal because $\vec{v}\in\spanof{S\union\set{\vec{v}}}$.

For the other half assume that $\vec{v}\in\spanof{S}$
so that $\vec{v}=c_1\vec{s}_1+\cdots+c_n\vec{s}_n$ for some scalars~$c_i$
and vectors $\vec{s}_i\in S$.
We will use mutual containment to show that the sets
$\spanof{S\union\set{\vec{v}}}$ and $\spanof{S}$ are equal.
The containment $\spanof{S\union\set{\vec{v}}}\supseteq\spanof{S}$ is clear.
%</pf:AddVecSpanNoGrowIffVecInSpan0>

%<*pf:AddVecSpanNoGrowIffVecInSpan1>
To show containment in the other direction let $\vec{w}$ be an element 
of~$\spanof{S\union\set{\vec{v}}}$.
Then $\vec{w}$ is a linear combination of elements of 
$S\union\set{\vec{v}}$, which we can write as 
$\vec{w}=c_{n+1}\vec{s}_{n+1}+\cdots+c_{n+k}\vec{s}_{n+k}+c_{n+k+1}\vec{v}$.
(Possibly some of the $\vec{s}_i$'s from $\vec{w}$'s equation are the same as
some of those from $\vec{v}$'s equation but that does not matter.)
Expand $\vec{v}$.
\begin{equation*}
  \vec{w}=c_{n+1}\vec{s}_{n+1}+\cdots+c_{n+k}\vec{s}_{n+k}
            +c_{n+k+1}\cdot(c_1\vec{s}_1+\cdots+c_n\vec{s}_n)
\end{equation*}
Recognize the right hand side as a linear combination of linear
combinations of vectors from~$S$.
Thus $\vec{w}\in\spanof{S}$.
%</pf:AddVecSpanNoGrowIffVecInSpan1>
\end{proof}

The discussion at the section's opening involved removing vectors, not adding
them. 

\begin{corollary}  \label{cor:OmitVecNotShrinkSpanIffVecIsDep}
%<*cor:OmitVecNotShrinkSpanIffVecIsDep>
For $\vec{v}\in S$,
omitting that vector does not shrink the span
if and only if that vector is dependent on other vectors in the 
set. 
That is, 
$\spanof{S}=\spanof{S-\set{\vec{v}}}$
if and only if $\vec{v}\in\spanof{S-\set{\vec{v}}}$. 
%</cor:OmitVecNotShrinkSpanIffVecIsDep>
\end{corollary}

Thus, to know whether removing a vector will decrease
the span, we need to know whether the vector is a linear combination
of others in the set.

\begin{definition}\label{def:LinInd}
%<*df:LinInd>
In any vector space, a set of vectors is
\definend{linearly independent}\index{linearly independent}%
\index{sets!dependent, independent}
if none of its elements is a linear combination of the 
others from the set.\footnote{See also \nearbyremark{rem:WhyLIUsesMultiset}.}\spacefactor=1000\ 
Otherwise the set is 
\definend{linearly dependent}\index{linearly dependent}.
%</df:LinInd>
\end{definition}
% A multiset subset of a vector space is
% \definend{linearly independent}\index{linearly independent}%
% \index{sets!dependent, independent}
% if none of its elements is a linear combination of the 
% others.\footnote{More information on multisets is in the appendix. 
% Most of the time we won't need the 
% set-multiset distinction and we will
% follow the standard terminology of referring to a linearly independent or 
% dependent `set'.
% \nearbyremark{rem:WhyLIUsesMultiset} explains why 
% the definition requires a multiset.}\spacefactor=1000\ 
% Otherwise it is \definend{linearly dependent}\index{linearly dependent}.

Thus the set $\set{\vec{s}_0,\ldots,\vec{s}_n}$ is independent
if there is no equality
$\vec{s}_i=c_0\vec{s}_0+\ldots+c_{i-1}\vec{s}_{i-1}+c_{i+1}\vec{s}_{i+1}+\ldots+c_n\vec{s}_n$.
The definition's use of the word `others' means that
writing $\vec{s}_i$ as a linear combination via $\vec{s}_i=1\cdot\vec{s}_i$ 
does not count.

%<*LinInd>
Observe that, 
although this way of writing one vector as a combination of the others
\begin{equation*}
   \vec{s}_0=\lincombo{c}{\vec{s}}
\end{equation*}
visually sets off \( \vec{s}_0 \), algebraically
there is nothing special about that vector in that equation.
For any \( \vec{s}_i \) with a coefficient $c_i$ that is non-$0$,
we can rewrite to isolate \( \vec{s}_i \).
\begin{equation*}
   \vec{s}_i=(1/c_i)\vec{s}_0+\dots
              +(-c_{i-1}/c_i)\vec{s}_{i-1}+(-c_{i+1}/c_i)\vec{s}_{i+1}
              +\dots+(-c_n/c_i)\vec{s}_n
\end{equation*}
When we don't want to single out any vector
we will instead say that
\( \vec{s}_0,\vec{s}_1,\dots,\vec{s}_n \) are in a
\definend{linear relationship}\index{linear relationship}%
\index{relationship!linear} 
and put all of the vectors on the same side.
%</LinInd>
The next result rephrases the linear independence definition in this style.
It is how we usually compute whether
a finite set is dependent or independent.

\begin{lemma}   \label{le:LDIffANonTrivLinRel}
%<*lm:LDIffANonTrivLinRel>
A subset \( S \) of a vector space is linearly independent if and only if 
among its elements  
the only linear relationship 
$ c_1\vec{s}_1+\dots+c_n\vec{s}_n=\zero$
is the trivial one, \( c_1=0,\dots,\,c_n=0 \)
(where $\vec{s}_i\neq\vec{s}_j$ when $i\neq j$) .
%</lm:LDIffANonTrivLinRel>
\end{lemma}

\begin{proof}
%<*pf:LDIffANonTrivLinRel>
If \( S \) is linearly independent then no vector
$\vec{s}_i$ 
is a linear combination of other vectors from $S$,
so there is no linear relationship where some of the
$\vec{s}\,$'s have nonzero coefficients.

If \( S \) is not linearly independent then some \( \vec{s}_i \) is a linear
combination 
$\vec{s}_i=c_1\vec{s}_1+\dots+c_{i-1}\vec{s}_{i-1}
    +c_{i+1}\vec{s}_{i+1}+\dots+c_n\vec{s}_n$
of other vectors from \( S \). 
Subtracting $\vec{s}_i$ from both sides
gives a relationship
involving a nonzero coefficient,
the \( -1 \) in front of \( \vec{s}_i \).
%</pf:LDIffANonTrivLinRel>
\end{proof}

\begin{example}  \label{ex:StaticsLI}
In the vector space of two-wide row vectors, the two-element set
\( \set{ \rowvec{40 &15},\rowvec{-50 &25}} \) is linearly independent.
To check this, take
\begin{equation*}
  c_1\cdot\rowvec{40 &15}+c_2\cdot\rowvec{-50 &25}=\rowvec{0 &0}
\end{equation*}
and solve the resulting system.
\begin{equation*}
  \begin{linsys}{2}
    40c_1  &-  &50c_2  &=  &0  \\
    15c_1  &+  &25c_2  &=  &0  
   \end{linsys}
  \grstep{-(15/40)\rho_1+\rho_2}
  \begin{linsys}{2}
     40c_1  &- &50c_2       &=  &0  \\
            &  &(175/4)c_2  &=  &0  
   \end{linsys}
\end{equation*}
Both \( c_1 \) and \( c_2 \) are zero. 
So the only linear relationship between the two given row vectors
is the trivial relationship.

In the same vector space, the set
\( \set{ \rowvec{40 &15},\rowvec{20 &7.5}} \) is linearly dependent since
we can satisfy
% \begin{equation*}
$c_1\cdot \rowvec{40 &15}+c_2\cdot\rowvec{20 &7.5}=\rowvec{0 &0}$
% \end{equation*}
with \( c_1=1 \) and \( c_2=-2 \).
\end{example}

\begin{example}
The set \( \set{1+x,1-x} \) is linearly independent in \( \polyspace_2 \), the
space of quadratic polynomials with real coefficients, because
\begin{equation*}
   0+0x+0x^2
   =
   c_1(1+x)+c_2(1-x)
   =
   (c_1+c_2)+(c_1-c_2)x+0x^2
\end{equation*}
gives 
\begin{equation*}
  \begin{linsys}{2}
    c_1  &+  &c_2  &=  &0  \\
    c_1  &-  &c_2  &=  &0  
   \end{linsys}
  \grstep{-\rho_1+\rho_2}
  \begin{linsys}{2}
     c_1  &+  &c_2  &=  &0  \\
          &   &2c_2 &=  &0
  \end{linsys}
\end{equation*}
since polynomials are equal only if their coefficients are equal.
Thus, the only linear relationship between these two members of
$\polyspace_2$ is the trivial one.
\end{example}

\begin{remark}
The lemma specifies that $\vec{s}_i\neq\vec{s}_j$ 
when $i\neq j$ because of course if some vector~$\vec{s}$ appears 
twice then we can get a nontrivial
$c_1\vec{s}_1+\dots+c_n\vec{s}_n=\zero$, by taking
the associated coefficients to be $1$ and~$-1$.
Besides, if some vector appears more than once in an expression then
we can always combine the coefficients.

Note that the lemma allows the 
opposite of appearing more than once, that some 
vectors from~$S$ don't appear at all.
For instance, if~$S$ is infinite then because
linear relationships involve only finitely many vectors, 
any such relationship leaves out many of $S$'s vectors.
However, note also that if~$S$ is finite then where convenient
we can take a combination
$c_1\vec{s}_1+\dots+c_n\vec{s}_n$ to contain each of $S$'s vectors once and
only once.
If a vector is missing then we can add it by using a 
coefficient of~$0$.
\end{remark}

\begin{example}
The rows of this matrix
\begin{equation*}
  A=
  \begin{mat}[r]
    2  &3  &1  &0  \\
    0  &-1 &0  &-2 \\
    0  &0  &0  &1
  \end{mat}
\end{equation*}
form a linearly independent set.
This is easy to check for this case but also 
recall that Lemma~One.III.\ref{le:EchFormNoLinCombo}
shows that the rows of any echelon form matrix make
a linearly independent set.
\end{example}

\begin{example}
In \( \Re^3 \), where
\begin{equation*}
   \vec{v}_1=\colvec{3 \\ 4 \\ 5}
   \quad
   \vec{v}_2=\colvec{2 \\ 9 \\ 2}
   \quad
   \vec{v}_3=\colvec{4 \\ 18 \\ 4}
\end{equation*}
the set \( S=\set{\vec{v}_1,\vec{v}_2,\vec{v}_3} \)
is linearly dependent because this is a relationship
\begin{equation*}
  0\cdot\vec{v}_1
  +2\cdot\vec{v}_2
  -1\cdot\vec{v}_3
  =\zero
\end{equation*}
where not all of the scalars are zero 
(the fact that some 
of the scalars are zero doesn't matter).
\end{example}

% \begin{remark}  \label{rem:WhyLIIsNonTrivLinRel}
That example illustrates why, 
although \nearbydefinition{def:LinInd} is a clearer
statement of what independence means,
\nearbylemma{le:LDIffANonTrivLinRel} is better for
computations.
Working straight from the definition, someone trying to compute whether $S$
is linearly independent would start by setting
\( \vec{v}_1=c_2\vec{v}_2+c_3\vec{v}_3 \)
and concluding that there are no such $c_2$ and $c_3$.
But knowing that the first vector is not
dependent on the other two is not enough.
This person would have to go on to try
\( \vec{v}_2=c_1\vec{v}_1+c_3\vec{v}_3 \), in order  
to find the dependence $c_1=0$, \( c_3=1/2 \).
\nearbylemma{le:LDIffANonTrivLinRel} 
gets the same conclusion with only one computation.
% \end{remark}

\begin{example} \label{ex:EmSetLI}
The empty subset\index{sets!empty} of a vector space is linearly independent.
There is no nontrivial linear relationship among its members as it has
no members.
\end{example}

\begin{example} \label{ex:SetWithZeroVecLD}
In any vector space, any subset containing the zero vector is linearly 
dependent.
One example is, in the space $\polyspace_2$ of quadratic polynomials, 
the subset $\set{1+x,x+x^2,0}$.
It is linearly 
dependent because 
$0\cdot\vec{v}_1+0\cdot\vec{v}_2+1\cdot\zero=\zero$ is a nontrivial
relationship, since not all of the coefficients are zero.
\end{example}

There is a subtle point that we shall see a number of times and that
bears on the prior example.
It is about the trivial sum, the sum of the empty set.
One way to see how to define the trivial sum 
is to consider the progression
$\vec{v}_1+\vec{v}_2+\vec{v}_3$, followed by
$\vec{v}_1+\vec{v}_2$, followed by
$\vec{v}_1$.
The difference between the sum of three vectors and the sum of two 
is $\vec{v}_3$.
Then the difference between the sum of two and the sum of one 
is $\vec{v}_2$.
In next passing to the trivial sum, the sum of zero-many vectors, 
we can expect to subtract~$\vec{v}_1$.
So we define the sum of zero-many vectors to be the zero vector. 

The relation with the prior example is that if the zero vector is in a set
then that set has an element that
is a combination of a subset of other vectors from the set, specifically,
the zero vector is a combination of the empty subset.
Even the set $S=\set{\zero}$ is linearly dependent, because $\zero$ is the 
sum of the empty set and the empty set is a subset of~$S$. 

\begin{remark} \label{rem:WhyLIUsesMultiset} %\cite{Velleman} 
The definition of linear independence,
\nearbydefinition{def:LinInd}, refers to a `set' of vectors.
Sets are the most familiar kind of collection and
in practice everyone uses the word `set' in this context.
But to be complete, we will note that sets are not quite the right kind
of collection for this purpose.

Recall that a set is a collection with two properties: 
(i)~order does not matter,
so that the set $\set{1,2}$ equals the set $\set{2,1}$, and 
(ii)~duplicates collapse, so that the set $\set{1,1,2}$ equals the set
$\set{1,2}$.

Now consider this matrix reduction.
\begin{equation*}
  \begin{mat}
    1 &1 &1 \\
    2 &2 &2 \\
    1 &2 &3
  \end{mat}
  \grstep{(1/2)\rho_2}
  \begin{mat}
    1 &1 &1 \\
    1 &1 &1 \\
    1 &2 &3
  \end{mat}
\end{equation*}
On the left the set of matrix rows 
$\set{\rowvec{1 &1 &1}, \rowvec{2 &2 &2}, \rowvec{1 &2 &3}}$ is
linearly dependent.
On the right the set of rows is  
$\set{\rowvec{1 &1 &1}, \rowvec{1 &1 &1}, \rowvec{1 &2 &3}}$. 
Because duplicates collapse, that equals
the set
$\set{\rowvec{1 &1 &1}, \rowvec{1 &2 &3}}$, 
which is linearly independent.
This is a problem because Gauss's Method should preserve linear dependence.

That is, strictly speaking,
we need a type of collection where duplicates do not collapse.
A collection where order does not matter and duplicates don't collapse
is a \definend{multiset}.\index{sets!multiset}%
\index{multiset}\index{linear independence!multiset}

However, while insisting on being completely correct has advantages,
departing from the standard terminology of `set' would have pitfalls of its 
own, so we will continue to use that word.
Later, we shall occasionally need to take combinations without
letting duplicates collapse and we shall do that without 
further comment.
\end{remark}

\begin{corollary} \label{cor:LiSetMinSpanSet}
%<*cor:LiSetMinSpanSet>
A set $S$ is linearly independent if and only if
for any $\vec{v}\in S$, its removal shrinks the span
$\spanof{S-\set{v}}\subsetneq\spanof{S}$.
%</cor:LiSetMinSpanSet>
\end{corollary}

\begin{proof}
%<*pf:LiSetMinSpanSet>
This follows from \nearbycorollary{cor:OmitVecNotShrinkSpanIffVecIsDep}.
If $S$ is linearly independent then none of its vectors is dependent on 
the other elements, so removal of any vector will shrink the span.
If $S$ is not linearly independent then it contains a vector that is 
dependent on other elements of the set, and removal of that vector
will not shrink the span.
%</pf:LiSetMinSpanSet>
\end{proof}

So a spanning set is minimal if and only if it is linearly independent.

The prior result addresses removing elements from a linearly independent set.
The next one adds elements.

\begin{lemma}  \label{lm:AddVecLiSetIsLiIffVecNotInSpan}
%<*lm:AddVecLiSetIsLiIffVecNotInSpan>
Suppose that $S$ is linearly independent and that $\vec{v}\notin S$.
Then
the set $S\union\set{\vec{v}}$ is linearly independent if and only if 
$\vec{v}\notin\spanof{S}$.   
%</lm:AddVecLiSetIsLiIffVecNotInSpan>
\end{lemma}

\begin{proof}
%<*pf:AddVecLiSetIsLiIffVecNotInSpan0>
% Assume that $S$ is linearly independent and that $\vec{v}\notin S$.
We will show that $S\union\set{\vec{v}}$ is not linearly independent 
if and only if $\vec{v}\in\spanof{S}$.

Suppose first that $\vec{v}\in\spanof{S}$.
Express $\vec{v}$ as a combination $\vec{v}=c_1\vec{s}_1+\cdots+c_n\vec{s}_n$.
Rewrite that $\zero=c_1\vec{s}_1+\cdots+c_n\vec{s}_n-1\cdot\vec{v}$.
Since $\vec{v}\notin S$, it does not equal any of the $\vec{s}_i$ so this is 
a nontrivial linear dependence among the elements of $S\union\set{\vec{v}}$.
Thus that set is not linearly independent.
%</pf:AddVecLiSetIsLiIffVecNotInSpan0>

%<*pf:AddVecLiSetIsLiIffVecNotInSpan1>
Now suppose that $S\union\set{\vec{v}}$ is not linearly independent and
consider a nontrivial dependence among its members
$\zero=c_1\vec{s}_1+\cdots+c_n\vec{s}_n+c_{n+1}\cdot\vec{v}$.
If $c_{n+1}=0$ then that is a dependence among the elements of~$S$, but 
we are assuming that~$S$ is independent, so $c_{n+1}\neq 0$.
Rewrite the equation  
as $\vec{v}=(c_1/c_{n+1})\vec{s}_1+\cdots+(c_n/c_{n+1})\vec{s}_n$
to get $\vec{v}\in\spanof{S}$
%</pf:AddVecLiSetIsLiIffVecNotInSpan1>
\end{proof}



\begin{example} \label{ex:LinindSetsAndSuper}
This subset of $\Re^3$ is linearly independent.
\begin{equation*}
  S
   =\set{\colvec{1 \\ 0 \\ 0}}
\end{equation*}
The span of $S$ is the $x$-axis.
Here are two supersets, one that is linearly dependent and the other
independent.
\begin{center}
     dependent:
     \( \set{
         \colvec{1 \\ 0 \\ 0},
         \colvec{-3 \\ 0 \\ 0}} \)      
     \qquad
     independent:
     \( \set{
         \colvec{1 \\ 0 \\ 0},
         \colvec{0 \\ 1 \\ 0}} \)      
\end{center}
We got the
dependent superset by adding a vector from the $x$-axis
and so the span did not grow.
We got the independent superset by adding a vector 
that isn't in $\spanof{S}$, because it has a nonzero $y$~component,
causing the span to grow.

For the independent set
\begin{equation*}
  S
   =\set{\colvec{1 \\ 0 \\ 0},
          \colvec{0 \\ 1 \\ 0}
                 }
\end{equation*}
the span \( \spanof{S} \) is the \( xy \)-plane.
Here are two supersets.
\begin{center}
     dependent:
     \( \set{
         \colvec{1 \\ 0 \\ 0},
         \colvec{0 \\ 1 \\ 0},
         \colvec{3 \\ -2 \\ 0} } \)       
     \qquad
     independent:
     \( \set{
         \colvec{1 \\ 0 \\ 0},
         \colvec{0 \\ 1 \\ 0},
         \colvec{0 \\ 0 \\ 1} } \)    
\end{center}
As above, the additional member of the 
dependent superset comes from $\spanof{S}$, the $xy$-plane, while the 
added member of the
independent superset comes from outside of that span.

Finally, consider this independent set 
\begin{equation*}
   S =\set{\colvec{1 \\ 0 \\ 0},
          \colvec{0 \\ 1 \\ 0},
          \colvec{0 \\ 0 \\ 1}        } 
\end{equation*}
with \( \spanof{S}=\Re^3 \).  
We can get a linearly dependent superset.
\begin{center}
     dependent:
     \( \set{
         \colvec{1 \\ 0 \\ 0},
         \colvec{0 \\ 1 \\ 0},
         \colvec{0 \\ 0 \\ 1},
         \colvec{2 \\ -1 \\ 3} } \)     
\end{center}
But there is no linearly independent superset of $S$.
One way to see that is to note that
for any vector that we would add to $S$, the equation
\begin{equation*}
  \colvec{x \\ y \\ z}
  =c_1\colvec{1 \\ 0 \\ 0}
   +c_2\colvec{0 \\ 1 \\ 0}
   +c_3\colvec{0 \\ 0 \\ 1}
\end{equation*}
has a solution $c_1=x$, $c_2=y$, and $c_3=z$.
Another way to see it is that
we cannot add any vectors from outside of the span $\spanof{S}$ because that
span is $\Re^3$.
\end{example}

\begin{corollary} \label{th:AlwaysAnLDSubset}
%<*th:AlwaysAnLDSubset>
In a vector space,
any finite set has a linearly independent subset with the same span.
%</th:AlwaysAnLDSubset>
\end{corollary}

\begin{proof}
%<*pf:AlwaysAnLDSubset0>
If \( S=\set{ \vec{s}_1,\dots,\vec{s}_n} \) is linearly independent
then $S$ itself satisfies the statement, so 
assume that it is linearly dependent.
%</pf:AlwaysAnLDSubset0>

%<*pf:AlwaysAnLDSubset1>
By the definition of dependent, $S$ contains
a vector $\vec{v}_1$ that is a linear combination of
the others.
Define the set \( S_1=S-\set{\vec{v}_1} \).
By \nearbycorollary{cor:OmitVecNotShrinkSpanIffVecIsDep} % {lm:ShrinkSpanByRemovingNonRepeat} 
the span does not shrink \( \spanof{S_1}=\spanof{S} \).
%</pf:AlwaysAnLDSubset1>

%<*pf:AlwaysAnLDSubset2>
If \( S_1 \) is linearly independent then we are done.
Otherwise iterate: 
take a vector $\vec{v}_2$ 
that is a linear combination of
other members of $S_1$ and discard it
to derive \( S_2=S_1-\set{\vec{v}_2} \)
such that \( \spanof{S_2}=\spanof{S_1} \).
Repeat this until a linearly independent set $S_j$ appears;
one must appear eventually because \( S \) is finite
and the empty set is linearly independent.
%</pf:AlwaysAnLDSubset2>
(Formally, this argument uses
induction on the number of elements in $S$.
\nearbyexercise{exer:FillIndDetProofSetHasLISub} asks for the details.)
\end{proof}

Thus if we have a set that is linearly dependent then we can, without changing
the span, pare down by   
discarding what we have called ``repeat'' vectors.

\begin{example}  \label{ex:ShrinkSetSameSpan}
This set spans \( \Re^3 \) (the check is routine)
but is not linearly independent. 
\begin{equation*}
  S=\set{\colvec[r]{1 \\ 0 \\ 0},
     \colvec[r]{0 \\ 2 \\ 0},
     \colvec[r]{1 \\ 2 \\ 0},
     \colvec[r]{0 \\ -1 \\ 1},
     \colvec[r]{3 \\ 3 \\ 0}  }
\end{equation*}
We will calculate which vectors to drop in order
to get a subset that is independent but
has the same span.
This linear relationship
\begin{equation*}
  c_1\colvec[r]{1 \\ 0 \\ 0}
  +c_2\colvec[r]{0 \\ 2 \\ 0}
  +c_3\colvec[r]{1 \\ 2 \\ 0}
  +c_4\colvec[r]{0 \\ -1 \\ 1}
  +c_5\colvec[r]{3 \\ 3 \\ 0}
  =\colvec[r]{0 \\ 0 \\ 0}
  \qquad\tag{$*$}
\end{equation*}
gives a system 
\begin{equation*}
  \begin{linsys}{5}
     c_1  &   &      &+  &c_3   &+  &    &+  &3c_5 &= &0  \\
          &   &2c_2  &+  &2c_3  &-  &c_4 &+  &3c_5 &= &0  \\
          &   &      &   &     &   &c_4  &   &     &= &0  
\end{linsys}
\end{equation*}
whose solution set has this parametrization.
\begin{equation*}
  \set{\colvec{c_1 \\ c_2 \\ c_3 \\ c_4 \\ c_5}=
     c_3\colvec[r]{-1 \\ -1 \\ 1 \\ 0 \\ 0}
     +c_5\colvec[r]{-3 \\ -3/2 \\ 0 \\ 0 \\ 1}
     \suchthat c_3,c_5\in\Re }
\end{equation*}
Set $c_5=1$ and~$c_3=0$ 
to get an instance of ($*$).
\begin{equation*}
  -3\cdot\colvec[r]{1 \\ 0 \\ 0}
  -\frac{3}{2}\cdot\colvec[r]{0 \\ 2 \\ 0}
  +0\cdot\colvec[r]{1 \\ 2 \\ 0}
  +0\cdot\colvec[r]{0 \\ -1 \\ 1}
  +1\cdot\colvec[r]{3 \\ 3 \\ 0}
  =\colvec[r]{0 \\ 0 \\ 0}
\end{equation*}
This shows that the vector from $S$ that we've 
associated with~$c_5$
is in the span of the set of $c_1$'s vector and $c_2$'s vector.
We can discard $S$'s fifth vector without shrinking the span.

Similarly, set $c_3=1$, and $c_5=0$
to get an instance of~($*$) that
shows we can discard $S$'s third vector without shrinking the span. 
 Thus this set has the same span as $S$.
\begin{equation*}
  % \hat{S}=
     \set{\colvec[r]{1 \\ 0 \\ 0},
     \colvec[r]{0 \\ 2 \\ 0},
     \colvec[r]{0 \\ -1 \\ 1}  }
\end{equation*}
The check that it is linearly independent is routine.
\end{example}

\begin{corollary} \label{cor:LDMeansLC}
%<*co:LDMeansLC>
A subset \( S=\set{\vec{s}_1,\dots,\vec{s}_n} \) of a vector space
is linearly dependent if and only if some \( \vec{s_i} \)
is a linear combination of the vectors 
\( \vec{s}_1 \), \ldots, \( \vec{s}_{i-1} \)
listed before it.
%</co:LDMeansLC>
\end{corollary}

\begin{proof}
%<*pf:LDMeansLC>
Consider \( S_0=\set{} \), \( S_1=\set{\vec{s_1}} \),
\( S_2=\set{\vec{s}_1,\vec{s}_2 } \), etc.
Some index \( i\geq 1 \) is the first one with
\( S_{i-1}\union\set{\vec{s}_i } \)
linearly dependent, and there \( \vec{s}_i\in\spanof{ S_{i-1} } \).
%</pf:LDMeansLC>
\end{proof}

The proof of 
\nearbycorollary{th:AlwaysAnLDSubset} describes producing a linearly 
independent set by shrinking, by taking subsets.
And the proof of \nearbycorollary{cor:LDMeansLC} describes finding a 
linearly dependent set by taking supersets. 
We finish this subsection by considering
how linear independence and dependence interact
with the subset relation between sets.

\begin{lemma}  \label{le:SubsetPreserveLI}
%<*lm:SubsetPreserveLI>
Any subset of a linearly independent set is also linearly independent.
Any superset of a linearly dependent set is also linearly dependent.
%</lm:SubsetPreserveLI>
\end{lemma}

\begin{proof}
%<*pf:SubsetPreserveLI>
Both are clear.
%</pf:SubsetPreserveLI>
\end{proof}

%<*SubsetPreserveLI>
Restated, subset preserves independence  
and superset preserves dependence.
%</SubsetPreserveLI>

Those are two of the four possible cases.
The third case, whether subset preserves linear dependence,
is covered by \nearbyexample{ex:ShrinkSetSameSpan}, which gives
a linearly dependent set $S$ with one subset 
that is linearly dependent and another 
that is independent.
The fourth case, whether superset preserves linear independence,
is covered by \nearbyexample{ex:LinindSetsAndSuper}, which gives cases
where a linearly independent set has both an independent
and a dependent superset.
This table summarizes.
\smallskip
%<*IndependenceAndSubsetTable>
\begin{center} \renewcommand{\arraystretch}{1.1}
  \begin{tabular}[b]{r|cc} 
                        \multicolumn{1}{c}{}
                        &\multicolumn{1}{c}{\( \hat{S}\subset S \)}
                        &\multicolumn{1}{c}{\( \hat{S}\supset S \)}      \\
     \cline{2-3}\rule{0em}{12pt} % make box a bit taller
          \textit{$S$ independent} 
              &\( \hat{S} \) must be independent   &\( \hat{S} \) may be either\\
     % \cline{2-3}\noalign{\smallskip}
          \textit{$S$ dependent} 
              &\( \hat{S} \) may be either &\( \hat{S} \) must be dependent    \\
     % \cline{2-3}
   \end{tabular}
\end{center}
%</IndependenceAndSubsetTable>
\smallskip

\nearbyexample{ex:LinindSetsAndSuper} has something else to say about the
interaction between linear independence and superset.
It names a
linearly independent set that is maximal in
that it has no supersets that are linearly independent.
By \nearbylemma{lm:AddVecLiSetIsLiIffVecNotInSpan} 
a linearly independent set is maximal if and only if it
spans the 
entire space, because that is when 
all the vectors in the space are already in the span.
This nicely
complements \nearbylemma{cor:LiSetMinSpanSet}, that
a spanning set is minimal if and only if it is linearly independent.

% In summary,
% we have introduced the definition of linear independence to 
% formalize the idea of the minimality of a spanning set.
% We have developed some properties of this idea.
% The most important is \nearbylemma{lm:AddVecLiSetIsLiIffVecNotInSpan}, which 
% tells us that a linearly independent set is maximal when it spans the space.


\begin{exercises}
  \recommended \item
    Decide whether each subset of \( \Re^3 \) is linearly dependent or
    linearly independent. 
    \begin{exparts*}
      \partsitem \( \set{\colvec{1 \\ -3 \\ 5},
                    \colvec{2 \\ 2 \\ 4},
                    \colvec{4 \\ -4 \\ 14} }  \)
      \partsitem \( \set{\colvec{1 \\ 7 \\ 7},
                    \colvec{2 \\ 7 \\ 7},
                    \colvec{3 \\ 7 \\ 7} }  \)
      \partsitem \( \set{\colvec{0 \\ 0 \\ -1},
                    \colvec{1 \\ 0 \\ 4} }  \)
      \partsitem \( \set{\colvec{9 \\ 9 \\ 0},
                    \colvec{2 \\ 0 \\ 1},
                    \colvec{3 \\ 5 \\ -4},
                    \colvec{12 \\ 12 \\ -1} }  \)
    \end{exparts*}
    \begin{answer}
      For each of these, when the subset is independent you must prove it, and
      when the subset is dependent you must give an example of a dependence.
      \begin{exparts}
        \partsitem It is dependent.
          Considering
          \begin{equation*}
             c_1\colvec{1 \\ -3 \\ 5}
             +c_2\colvec{2 \\ 2 \\ 4}
             +c_3\colvec{4 \\ -4 \\ 14}
             =\colvec{0 \\ 0 \\ 0}
          \end{equation*}
          gives this linear system.
          \begin{equation*}
            \begin{linsys}{3}
              c_1  &+  &2c_2  &+  &4c_3  &=  &0  \\
              -3c_1&+  &2c_2  &-  &4c_3  &=  &0  \\
              5c_1 &+  &4c_2  &+  &14c_3 &=  &0  
            \end{linsys}
          \end{equation*}
          Gauss's Method 
          \begin{equation*}
            \begin{amat}[r]{3}
              1  &2  &4  &0  \\
              -3 &2  &-4 &0  \\
              5  &4  &14 &0
            \end{amat}
            \grstep[-5\rho_1+\rho_3]{3\rho_1+\rho_2}
            \repeatedgrstep{(3/4)\rho_2+\rho_3}            
            \begin{amat}[r]{3}
              1  &2  &4  &0  \\
              0  &8  &8  &0  \\
              0  &0  &0  &0
            \end{amat}
          \end{equation*}
          yields a free variable, so there are infinitely many solutions.
          For an example of a particular dependence we can set $c_3$ to be,
          say, $1$.  Then we get
          \( c_2=-1 \) and \( c_1=-2 \).
        \partsitem It is dependent.
          The linear system that arises here
          \begin{equation*}
            \begin{amat}[r]{3}
              1  &2  &3  &0  \\
              7  &7  &7  &0  \\
              7  &7  &7  &0
            \end{amat}
            \grstep[-7\rho_1+\rho_3]{-7\rho_1+\rho_2}
            \repeatedgrstep{-\rho_2+\rho_3}
            \begin{amat}[r]{3}
              1  &2  &3   &0  \\
              0  &-7 &-14 &0  \\
              0  &0  &0   &0
            \end{amat}
          \end{equation*}
          has infinitely many solutions.
          We can get a particular solution by taking $c_3$ to be, say,
          $1$, and back-substituting to get the resulting $c_2$ and $c_1$.
        \partsitem It is linearly independent.
          The system
          \begin{equation*}
            \begin{amat}[r]{2}
              0  &1  &0  \\
              0  &0  &0  \\
              -1 &4  &0
            \end{amat}
            \grstep{\rho_1\leftrightarrow\rho_2}
            \repeatedgrstep{\rho_3\leftrightarrow\rho_1}
            \begin{amat}{2}
              -1 &4  &0  \\
              0  &1  &0  \\
              0  &0  &0  
            \end{amat}
          \end{equation*}
          has only the solution $c_1=0$ and $c_2=0$.
          (We could also have gotten the answer by inspection\Dash the second
          vector is obviously not a multiple of the first, and vice versa.)
        \partsitem It is linearly dependent.
          The linear system
          \begin{equation*}
            \begin{amat}[r]{4}
              9  &2  &3  &12  &0  \\
              9  &0  &5  &12  &0  \\
              0  &1  &-4 &-1  &0
            \end{amat}
          \end{equation*}
          has more unknowns than equations, and so Gauss's Method
          must end with at least one variable free (there can't be a 
          contradictory equation because the system is homogeneous, and so
          has at least the solution of all zeroes).
          To exhibit a combination, we can do the reduction 
          \begin{equation*}
            \grstep{-\rho_1+\rho_2}
            \grstep{(1/2)\rho_2+\rho_3}
            \begin{amat}[r]{4}
              9  &2  &3  &12  &0  \\
              0  &-2 &2  &0   &0  \\
              0  &0  &-3 &-1  &0
            \end{amat}
          \end{equation*}
          and take, say,  $c_4=1$.
          Then we have that $c_3=-1/3$, $c_2=-1/3$, and $c_1=-31/27$.
      \end{exparts}  
    \end{answer}
  \recommended \item 
    Which of these subsets of \( \polyspace_3 \) are
    linearly dependent and which are independent?
    \begin{exparts}
      \partsitem \( \set{3-x+9x^2,5-6x+3x^2,1+1x-5x^2} \)
      \partsitem \( \set{-x^2,1+4x^2} \)
      \partsitem \( \set{2+x+7x^2,3-x+2x^2,4-3x^2} \)
      \partsitem \( \set{8+3x+3x^2,x+2x^2,2+2x+2x^2,8-2x+5x^2} \)
    \end{exparts}
    \begin{answer}
      In the cases of independence, you must prove that it is independent.
      Otherwise, you must exhibit a dependence.
      (Here we give a specific dependence but others are possible.)
      \begin{exparts}
        \partsitem This set is independent.
          Setting up the relation
          \( c_1(3-x+9x^2)+c_2(5-6x+3x^2)+c_3(1+1x-5x^2)=0+0x+0x^2 \)
          gives a linear system 
          \begin{equation*}
            \begin{amat}[r]{3}
              3  &5  &1  &0  \\
              -1 &-6 &1  &0  \\
              9  &3  &-5 &0  
            \end{amat}
            \grstep[-3\rho_1+\rho_3]{(1/3)\rho_1+\rho_2}
            \repeatedgrstep{3\rho_2}
            \repeatedgrstep{-(12/13)\rho_2+\rho_3}
            \begin{amat}[r]{3}
              3  &5   &1        &0  \\
              0  &-13 &4        &0  \\
              0  &0   &-128/13  &0  
            \end{amat}
          \end{equation*}
          with only one solution: \( c_1=0 \), \( c_2=0 \), and \( c_3=0 \).
        \partsitem This set is independent.
           We can see this by inspection, straight from the definition
           of linear independence.
           Obviously neither is a multiple of the other.
        \partsitem This set is linearly independent.
           The linear system reduces in this way
           \begin{equation*}
             \begin{amat}[r]{3}
               2  &3  &4  &0  \\
               1  &-1 &0  &0  \\
               7  &2  &-3 &0  
             \end{amat}
             \grstep[-(7/2)\rho_1+\rho_3]{-(1/2)\rho_1+\rho_2}
             \repeatedgrstep{-(17/5)\rho_2+\rho_3}
             \begin{amat}[r]{3}
               2  &3    &4      &0  \\
               0  &-5/2 &-2     &0  \\
               0  &0    &-51/5  &0  
             \end{amat}
           \end{equation*}
           to show that there is only the solution $c_1=0$, 
           $c_2=0$, and $c_3=0$.
        \partsitem This set is linearly dependent.
           The linear system
           \begin{equation*}
             \begin{amat}[r]{4}
               8  &0  &2  &8  &0  \\ 
               3  &1  &2  &-2 &0  \\
               3  &2  &2  &5  &0
             \end{amat}
           \end{equation*}
           must, after reduction, end with at least one variable free
           (there are more variables than equations, and there is no
           possibility of a contradictory equation because the system is
           homogeneous).
           We can take the free variables as parameters to describe the
           solution set.
           We can then set the parameter to a nonzero value to get a
           nontrivial linear relation. 
      \end{exparts}  
     \end{answer}
  \item Determine if each set is linearly independent in the natural 
    space.
    \begin{exparts*}
      \partsitem $\set{\colvec{1 \\ 2 \\ 0}, 
              \colvec{-1 \\ 1 \\ 0}}$
      \partsitem $\set{\rowvec{1 &3 &1}, \rowvec{-1 &4 &3}, \rowvec{-1 &11 &7}}$
      \partsitem $\set{\begin{mat}
                5 &4 \\
                1 &2 
              \end{mat},
              \begin{mat}
                0 &0 \\
                0 &0
              \end{mat},
              \begin{mat}
                1 &0 \\
               -1 &4
              \end{mat}
           }$
    \end{exparts*}
    \begin{answer}
      \begin{exparts}
        \partsitem   The natural vector space is $\Re^3$.
          Set up the equation
          \begin{equation*}
            \colvec{0 \\ 0 \\ 0}=c_1\colvec{1 \\ 2 \\ 0}
                                 +c_2\colvec{-1 \\ 1 \\ 0}
          \end{equation*}
          and consider the resulting homogeneous system.
          \begin{equation*}
             \begin{amat}{2}
               1  &-1 &0  \\
               2  &1  &0  \\
               0  &0  &0
             \end{amat}
             \grstep{-2\rho_1+\rho_2}
             \begin{amat}{2}
              1  &-1 &0  \\
              0  &3  &0  \\
              0  &0  &0
            \end{amat} 
          \end{equation*}
          This has the unique solution, that $c_1=0$, $c_2=0$.
          So it is linearly independent.
        \partsitem   The natural vector space is the set of three-wide 
          row vectors.
          The equation
          \begin{equation*}
            \rowvec{0 &0 &0}=c_1\rowvec{1 &3 &1}
                            +c_2\rowvec{-1 &4 &3}
                            +c_3\rowvec{-1 &11 &7}
          \end{equation*}
          gives rise to a linear system
          \begin{align*}
            \begin{amat}{3}
              1  &-1  &-1  &0  \\
              3  &4   &11  &0  \\
              1  &3   &7   &0
           \end{amat}
           &\grstep[-\rho_1+\rho_3]{-3\rho_1+\rho_2}
           \begin{amat}{3}
             1  &-1  &-1  &0  \\
             0  &7   &14  &0  \\
             0  &4   &8   &0
          \end{amat}                                   \\
          &\grstep{(4/7)\rho_2+\rho_3}
          \begin{amat}{3}
            1  &-1  &-1  &0  \\
            0  &7   &14  &0  \\
            0  &0   &0   &0
         \end{amat}
       \end{align*}
       with infinitely many solutions, that is, more than just the 
       trivial solution.
       \begin{equation*}
         \set{\colvec{c_1 \\ c_2 \\ c_3}
              =\colvec{-1 \\ -2 \\ 1}c_3
              \suchthat c_3\in\Re}
       \end{equation*}
       So the set is linearly dependent.
       One dependence comes from setting $c_3=2$, giving 
       $c_1=-2$ and $c_2=-4$.
      \partsitem   Without having to set up a system we can see 
        that the second element of the set is a multiple of the first 
        (namely, $0$ times the first). 
      \end{exparts}
    \end{answer}
  \recommended \item
    Prove that each set \( \set{f,g} \) is linearly independent in the
    vector space of all functions from \( \Re^+ \) to \( \Re \).
    \begin{exparts}
      \partsitem \( f(x)=x \) and \( g(x)=1/x \)
      \partsitem \( f(x)=\cos(x) \) and \( g(x)=\sin(x) \)
      \partsitem \( f(x)=e^x \) and \( g(x)=\ln(x) \)
    \end{exparts}
    \begin{answer}
      Let $\map{Z}{\Re^+}{\Re}$ be the zero function $Z(x)=0$, 
      which is the additive identity in
      the vector space under discussion.
      \begin{exparts}
        \partsitem This set is linearly independent.  
          Consider \( c_1\cdot f(x)+c_2\cdot g(x)=Z(x) \).
          Plugging in \( x=1 \) and \( x=2 \) gives a linear system 
          \begin{equation*}
            \begin{linsys}{2}
              c_1\cdot 1  &+  &c_2\cdot 1     &=  &0  \\
              c_1\cdot 2  &+  &c_2\cdot (1/2) &=  &0
            \end{linsys}
          \end{equation*}
          with the unique solution \( c_1=0 \), \( c_2=0 \).
        \partsitem This set is linearly independent.  
          Consider \( c_1\cdot f(x)+c_2\cdot g(x)=Z(x) \) and 
          plug in \( x=\pi \) and \( x=\pi/2 \) to get 
          \begin{equation*}
            \begin{linsys}{2}
              c_1\cdot (-1)  &+  &c_2\cdot 0     &=  &0  \\
              c_1\cdot 0     &+  &c_2\cdot 1     &=  &0
            \end{linsys}
          \end{equation*}
          which obviously gives \( c_1=0 \), \( c_2=0 \).
        \partsitem This set is also linearly independent.  
          Considering \( c_1\cdot f(x)+c_2\cdot g(x)=Z(x) \) and 
          plugging in \( x=1 \) and \( x=e \) 
          \begin{equation*}
            \begin{linsys}{2}
              c_1\cdot e    &+  &c_2\cdot 0     &=  &0  \\
              c_1\cdot e^e  &+  &c_2\cdot 1     &=  &0
            \end{linsys}
          \end{equation*}
          gives that \( c_1=0 \) and \( c_2=0 \).
      \end{exparts}  
     \end{answer}
  \recommended \item 
    Which of these subsets of the space of real-valued functions
    of one real variable is linearly dependent and which is linearly
    independent?
    (We have abbreviated some constant functions;~e.g., 
     in the first item, 
     the `$2$' stands for the constant function $f(x)=2$.)
    \begin{exparts*}
      \partsitem \( \set{2,4\sin^2(x),\cos^2(x)} \)
      \partsitem \( \set{1,\sin(x),\sin(2x)} \)
      \partsitem \( \set{x,\cos(x)} \)
      \partsitem \( \set{(1+x)^2,x^2+2x,3} \)
      \partsitem \( \set{\cos(2x),\sin^2(x),\cos^2(x)} \)
      \partsitem \( \set{0,x,x^2} \)
    \end{exparts*}
    \begin{answer}
      In each case, if the set is independent then you must prove that 
      and if it is
      dependent then you must exhibit a dependence.
      \begin{exparts}
        \partsitem This set is dependent.
          The familiar relation $\sin^2(x)+\cos^2(x)=1$ shows that
          $2=c_1\cdot(4\sin^2(x))+c_2\cdot(\cos^2(x))$ is satisfied by
          $c_1=1/2$ and $c_2=2$.
        \partsitem This set is independent.
          Consider the relationship
          $c_1\cdot 1+c_2\cdot\sin(x)+c_3\cdot\sin(2x)=0$
          (that `$0$' is the zero function).
          Taking three suitable points such as $x=\pi$, $x=\pi/2$, $x=\pi/4$ 
          gives a system
          \begin{equation*}
            \begin{linsys}{3}
               c_1  &   &                &   &      &=  &0  \\
               c_1  &+  &c_2             &   &      &=  &0  \\
               c_1  &+  &(\sqrt{2}/2)c_2 &+  &c_3   &=  &0  
            \end{linsys}
          \end{equation*}
          whose only solution is 
          $c_1=0$, $c_2=0$, and $c_3=0$. 
        \partsitem By inspection, this set is independent.
          Any dependence $\cos(x)=c\cdot x$ is not possible since the cosine
          function is not a multiple of the identity function
          (we are applying \nearbycorollary{cor:LDMeansLC}).
        \partsitem By inspection, we spot that there is a dependence.
          Because $(1+x)^2=x^2+2x+1$, we get that
          $c_1\cdot(1+x)^2+c_2\cdot(x^2+2x)=3$ is satisfied by 
          $c_1=3$ and $c_2=-3$.
        \partsitem This set is dependent.
          The easiest way to see that is to recall the trigonometric
          relationship $\cos^2(x)-\sin^2(x)=\cos(2x)$.
          (\textit{Remark.} 
          A person who doesn't recall this, and tries some $x$'s,
          simply never gets a system leading to a unique solution, and
          never gets to conclude that the set is independent.
          Of course, this person might wonder if they simply never tried the
          right set of $x$'s, but a few tries will lead most people to 
          look instead for a dependence.)
        \partsitem This set is dependent, because it contains the 
          zero object in the vector space, the zero polynomial.
      \end{exparts}  
     \end{answer}
  \item 
    Does the equation \( \sin^2(x)/\cos^2(x)=\tan^2(x) \) show that
    this set of functions
    \( \set{\sin^2(x),\cos^2(x),\tan^2(x)} \) is a linearly dependent
    subset of the set of all real-valued functions with domain
    the interval \( (-\pi/2..\pi/2) \) of real numbers between 
    \( -\pi/2 \) and \( \pi/2) \)?
    \begin{answer}
      No, that equation is not a linear relationship.
      In fact this set is independent, as the system arising from taking
      \( x \) to be \( 0 \), \( \pi/6 \) and \( \pi/4 \) shows.  
    \end{answer}
  \item  
    Is the $xy$-plane subset of the vector space $\Re^3$ linearly independent?
    \begin{answer}
      No.  
      Here are two members of the plane where the second is a multiple of the 
      first.
      \begin{equation*}
        \colvec{1  \\ 0 \\ 0},\,\colvec{2 \\ 0 \\ 0}
      \end{equation*}
      (Another reason that the answer is ``no'' is the the zero vector is a 
      member of the plane and no set containing the zero vector is linearly
      independent.)
    \end{answer}
  \recommended \item
    Show that the nonzero rows of an echelon form matrix form a linearly
    independent set.
    \begin{answer}
      We have already showed this: the Linear Combination
      Lemma and its corollary state that in an echelon form matrix, 
      no nonzero row is a linear combination of the others.  
    \end{answer}
  \item
     \begin{exparts}
       \partsitem Show that if the set \( \set{\vec{u},\vec{v},\vec{w}} \)
          is linearly independent then so is the set 
          \( \set{\vec{u},\vec{u}+\vec{v},\vec{u}+\vec{v}+\vec{w}} \).
       \partsitem  What is the relationship between the linear independence
         or dependence of \( \set{\vec{u},\vec{v},\vec{w}} \) and the
         independence or dependence of
         \( \set{\vec{u}-\vec{v},\vec{v}-\vec{w},\vec{w}-\vec{u}} \)?
     \end{exparts}
     \begin{answer}
       \begin{exparts}
         \partsitem Assume that 
           \( \set{\vec{u},\vec{v},\vec{w}} \) is linearly 
           independent, so that any relationship
           $d_0\vec{u}+d_1\vec{v}+d_2\vec{w}=\zero$ leads to the conclusion 
           that $d_0=0$, $d_1=0$, and $d_2=0$.

           Consider the relationship
           \( c_1(\vec{u})+c_2(\vec{u}+\vec{v})+c_3(\vec{u}+\vec{v}+\vec{w})
           =\zero \).
           Rewrite it to get
           \( (c_1+c_2+c_3)\vec{u}+(c_2+c_3)\vec{v}+(c_3)\vec{w}=\zero \).
           Taking $d_0$ to be $c_1+c_2+c_3$, taking $d_1$ to be $c_2+c_3$, 
           and taking $d_2$ to be $c_3$ we have this system.
           \begin{equation*}
             \begin{linsys}{3}
               c_1  &+  &c_2  &+  &c_3  &=  &0  \\
                    &   &c_2  &+  &c_3  &=  &0  \\
                    &   &     &   &c_3  &=  &0
             \end{linsys}
           \end{equation*}
           Conclusion:~the $c$'s are all zero, and so the set is linearly
           independent.
        \partsitem The second set is dependent
           \begin{equation*}
             1\cdot(\vec{u}-\vec{v})
             +1\cdot(\vec{v}-\vec{w})
             +1\cdot(\vec{w}-\vec{u})
             =\zero
           \end{equation*}
           whether or not the first set is independent.
      \end{exparts}
    \end{answer}
  \item 
    \nearbyexample{ex:EmSetLI} shows that the empty set is 
    linearly independent.   
    \begin{exparts}
      \partsitem When is a one-element set linearly independent?
      \partsitem How about a set with two elements?
    \end{exparts}
    \begin{answer}
      \begin{exparts}
        \partsitem A singleton set $\set{\vec{v}}$ is linearly independent 
          if and only if $\vec{v}\neq\zero$.
          For the `if' direction, with $\vec{v}\neq\zero$, 
          we can apply \nearbylemma{le:LDIffANonTrivLinRel} by considering 
          the relationship
          \( c\cdot\vec{v}=\zero \) and noting that the only solution
          is the trivial one:~$c=0$.
          For the `only~if' direction, just recall that 
          \nearbyexample{ex:SetWithZeroVecLD}
          shows that $\set{\zero}$ is linearly dependent, and so if the set
          $\set{\vec{v}}$ is linearly independent then $\vec{v}\neq\zero$. 

          (\textit{Remark.} 
          Another answer is to say that this is the special case of
          \nearbylemma{lm:AddVecLiSetIsLiIffVecNotInSpan} 
          where \( S=\emptyset \).)
        \partsitem A set with two elements is linearly independent 
          if and only if neither member is a  multiple of the other 
          (note that if one is the zero vector then it is a multiple of the
          other).
          This is an equivalent statement:~a set is linearly dependent if and
          only if one element is a multiple of the other.

          The proof is easy.
          A set $\set{\vec{v}_1,\vec{v}_2}$ is linearly dependent if and only
          if there is a relationship $c_1\vec{v}_1+c_2\vec{v}_2=\zero$ 
          with either $c_1\neq 0$ or $c_2\neq 0$ (or both).
          That holds if and only if $\vec{v}_1=(-c_2/c_1)\vec{v}_2$
          or $\vec{v}_2=(-c_1/c_2)\vec{v}_1$ (or both).
       \end{exparts}   
     \end{answer}
  \item  
    In any vector space \( V \), the empty set is linearly independent.
    What about all of \( V \)?
    \begin{answer}
      This set is linearly dependent set because it contains the zero vector.  
    \end{answer}
  \item 
    Show that if \( \set{\vec{x},\vec{y},\vec{z}} \) is linearly
    independent then so are all of its proper 
    subsets:~\( \set{\vec{x},\vec{y}} \),
    \( \set{\vec{x},\vec{z}} \), \( \set{\vec{y},\vec{z}} \),
    \( \set{\vec{x}} \),\( \set{\vec{y}} \), \( \set{\vec{z}} \),
    and \( \set{} \).
    Is that `only if' also?
    \begin{answer}
      \nearbylemma{le:SubsetPreserveLI} gives the `if' half.
      The converse (the `only if' statement) does not hold. 
      An example is to consider the vector space \( \Re^2 \) and
      these vectors.
      \begin{equation*}
         \vec{x}=\colvec{1 \\ 0},\quad
         \vec{y}=\colvec{0 \\ 1},\quad
         \vec{z}=\colvec{1 \\ 1}
      \end{equation*} 
    \end{answer}
  \item 
    \begin{exparts}
      \partsitem Show that this 
        \begin{equation*}
          S=\set{\colvec{1 \\ 1 \\ 0},\colvec{-1 \\ 2 \\ 0}}
        \end{equation*}
        is a linearly independent subset of \( \Re^3 \).
      \partsitem Show that
        \begin{equation*}
          \colvec{3 \\ 2 \\ 0}
        \end{equation*}
        is in the span of $S$ by finding \( c_1 \) and \( c_2 \) 
        giving a linear relationship.
        \begin{equation*}
          c_1\colvec{1 \\ 1 \\ 0}
          +c_2\colvec{-1 \\ 2 \\ 0}
          =\colvec{3 \\ 2 \\ 0}
        \end{equation*}
        Show that the pair \( c_1,c_2 \) is unique.
      \partsitem Assume that \( S \) is a subset of a vector space and that
        \( \vec{v} \) is in \( \spanof{S} \), so that \( \vec{v} \) is a
        linear combination of vectors from \( S \).
        Prove that if \( S \) is linearly independent then a linear combination
        of vectors from \( S \) adding to \( \vec{v} \)
        is unique (that is, unique up to reordering
        and adding or taking away terms of the form \( 0\cdot\vec{s} \)).
        Thus \( S \) as a spanning set is minimal in this strong sense:
        each vector in \( \spanof{S} \) is a combination of elements
        of $S$ a minimum number of
        times\Dash only once.
      \partsitem
        Prove that it can happen when \( S \) is not linearly 
        independent that distinct linear combinations sum to the same vector.
    \end{exparts}
    \begin{answer}
      \begin{exparts}
        \partsitem The linear system arising from
          \begin{equation*}
            c_1\colvec{1 \\ 1 \\ 0}
            +c_2\colvec{-1 \\ 2 \\ 0}
            =\colvec{0 \\ 0 \\ 0}
          \end{equation*}
          has the unique solution \( c_1=0 \) and \( c_2=0 \).
        \partsitem The linear system arising from
          \begin{equation*}
            c_1\colvec{1 \\ 1 \\ 0}
            +c_2\colvec{-1 \\ 2 \\ 0}
            =\colvec{3 \\ 2 \\ 0}
          \end{equation*}
          has the unique solution \( c_1=8/3 \) and \( c_2=-1/3 \).
        \partsitem Suppose that \( S \) is linearly independent.
          Suppose that we have both $\vec{v}=c_1\vec{s}_1+\dots+c_n\vec{s}_n$
          and $\vec{v}=d_1\vec{t}_1+\dots+d_m\vec{t}_m$
          (where the vectors are members of $S$).
          Now, 
          \begin{equation*}
            c_1\vec{s}_1+\dots+c_n\vec{s}_n
            =\vec{v}
            =d_1\vec{t}_1+\dots+d_m\vec{t}_m
          \end{equation*}
          can be rewritten in this way.
          \begin{equation*}
            c_1\vec{s}_1+\dots+c_n\vec{s}_n
            -d_1\vec{t}_1-\dots-d_m\vec{t}_m
            =\zero
          \end{equation*}
          Possibly some of the $\vec{s}\,$'s equal some of the $\vec{t}\,$'s;
          we can combine the associated coefficients 
          (i.e., if $\vec{s}_i=\vec{t}_j$ then
          $\cdots+c_i\vec{s}_i+\dots-d_j\vec{t}_j-\cdots$ can be rewritten
          as $\cdots+(c_i-d_j)\vec{s}_i+\cdots$).
          That equation is a linear relationship among  
          distinct (after the combining is done) members of the set $S$.
          We've assumed that $S$ is linearly independent, so all of the 
          coefficients are zero.
          If $i$ is such that $\vec{s}_i$ does not equal any $\vec{t}_j$
          then $c_i$ is zero.
          If $j$ is such that $\vec{t}_j$ does not equal any $\vec{s}_i$
          then $d_j$ is zero.
          In the final case, we have that $c_i-d_j=0$ and so $c_i=d_j$.  

          Therefore, the original two sums are the same, except perhaps for
          some $0\cdot\vec{s}_i$ or $0\cdot\vec{t}_j$ terms that we can
          neglect.
        \partsitem
          This set is not linearly independent:
          \begin{equation*}
            S=\set{\colvec{1 \\ 0},\colvec{2 \\ 0}}\subset\Re^2
          \end{equation*}
          and these two linear combinations give the same result
          \begin{equation*}
            \colvec{0 \\ 0}=2\cdot\colvec{1 \\ 0}-1\cdot\colvec{2 \\ 0}
                            =4\cdot\colvec{1 \\ 0}-2\cdot\colvec{2 \\ 0}
          \end{equation*}
          Thus, a linearly dependent set might have indistinct sums.

          In fact, this stronger statement holds:~if a set is linearly 
          dependent then it must have the property that there are two 
          distinct linear combinations that sum to the same vector.
          Briefly, where \( c_1\vec{s}_1+\dots+c_n\vec{s}_n=\zero \) then
          multiplying both sides of the relationship by two gives another 
          relationship.
          If the first relationship is nontrivial then the second is also.
      \end{exparts}  
    \end{answer}
  \item  \label{exer:PolyZeroFcnOnlyIfZeroPol}
     Prove that a polynomial gives rise to the zero function if and only if
     it is the zero polynomial.
     (\textit{Comment.}
     This question is not a Linear Algebra matter but we often use the result.
     A polynomial gives rise to a function in the natural 
     way:~$x\mapsto c_nx^n+\dots+c_1x+c_0$.)
     \begin{answer}
       In this `if and only if' statement, the `if' half is clear\Dash if 
       the polynomial is the zero polynomial then the function that arises 
       from the action of the polynomial must be the zero 
       function $x\mapsto 0$. 
       For `only if' we write $p(x)=c_nx^n+\dots+c_0$. 
       Plugging in zero $p(0)=0$ gives that $c_0=0$.
       Taking the derivative and plugging in zero $p^\prime(0)=0$ gives 
       that $c_1=0$.
       Similarly we get that each $c_i$ is zero, and $p$ is the zero 
       polynomial.
     \end{answer}
  \item
    Return to Section 1.2 and redefine point, line, plane,
    and other linear surfaces to avoid degenerate cases.
    \begin{answer}
      The work in this section suggests that we should define 
      an \( n \)-dimensional
      non-degenerate linear surface as the span of a linearly
      independent set of \( n \) vectors.  
    \end{answer}
  \item 
    \begin{exparts}  
      \partsitem Show that any set of four vectors in \( \Re^2 \) is 
         linearly dependent.
      \partsitem Is this true for any set of five?
         Any set of three?
      \partsitem What is the most number of elements that a 
         linearly independent subset of $\Re^2$ can have?
    \end{exparts}
    \begin{answer}
      \begin{exparts}
        \partsitem For any $a_{1,1}$, \ldots, $a_{2,4}$,
          \begin{equation*}
            c_1\colvec{a_{1,1} \\ a_{2,1}}
            +c_2\colvec{a_{1,2} \\ a_{2,2}}
            +c_3\colvec{a_{1,3} \\ a_{2,3}}
            +c_4\colvec{a_{1,4} \\ a_{2,4}}
            =\colvec{0 \\ 0}
          \end{equation*}
          yields a linear system
          \begin{equation*}
             \begin{linsys}{4}
              a_{1,1}c_1 &+ &a_{1,2}c_2 &+ &a_{1,3}c_3 &+ &a_{1,4}c_4 &= &0  \\
              a_{2,1}c_1 &+ &a_{2,2}c_2 &+ &a_{2,3}c_3 &+ &a_{2,4}c_4 &= &0  
             \end{linsys}
          \end{equation*}
        that has infinitely many solutions (Gauss's Method leaves at least
        two variables free).
        Hence there are nontrivial linear relationships among the given
        members of $\Re^2$.
      \partsitem Any set five vectors is a superset of a set of four vectors,
        and so is linearly dependent.

        With three vectors from $\Re^2$, the argument from the prior item 
        still applies, with the slight change that Gauss's Method now only 
        leaves at least one variable free (but that still gives infinitely many
        solutions).
      \partsitem The prior item shows that no three-element subset of $\Re^2$
        is independent.
        We know that there are two-element subsets of $\Re^2$ that are 
        independent\Dash one is
        \begin{equation*}
          \set{\colvec{1  \\ 0},\colvec{0  \\ 1}}
        \end{equation*} 
        and so the answer is two.
     \end{exparts}  
    \end{answer}
  \item
    Is there a set of four vectors in \( \Re^3 \) such that any three 
    form a linearly independent set?
    \begin{answer}
      Yes; here is one.
      \begin{equation*}
         \set{\colvec{1 \\ 0 \\ 0},
              \colvec{0 \\ 1 \\ 0},
              \colvec{0 \\ 0 \\ 1},
              \colvec{1 \\ 1 \\ 1} }
      \end{equation*}  
    \end{answer}
  \item  
    Must every linearly dependent set have a subset that is dependent and
    a subset that is independent?
    \begin{answer}
      Yes.
      Fix any vector space~$V$, and let $S$ be a linearly dependent subset
      of~$V$.
      For a subset of~$S$ that is dependent we can take~$S$ itself.
      For a subset of~$S$ that is independent we can take the empty set.
    \end{answer}
  \item  
    In \( \Re^4 \) what is the biggest linearly independent
    set you can find?
    The smallest?
    The biggest linearly dependent set?
    The smallest?
    (`Biggest' and `smallest' mean that there are no supersets or subsets 
    with the same property.)
    \begin{answer} 
      In \( \Re^4 \) the biggest linearly independent set has
      four vectors.
      There are many examples of such sets, this is one.
      \begin{equation*}
        \set{\colvec{1 \\ 0 \\ 0 \\ 0},
             \colvec{0 \\ 1 \\ 0 \\ 0},
             \colvec{0 \\ 0 \\ 1 \\ 0},
             \colvec{0 \\ 0 \\ 0 \\ 1}  }
      \end{equation*}
      To see that no set with five or more vectors can be independent, set up
      \begin{equation*}
             c_1\colvec{a_{1,1} \\ a_{2,1} \\ a_{3,1} \\ a_{4,1}}
            +c_2\colvec{a_{1,2} \\ a_{2,2} \\ a_{3,2} \\ a_{4,2}}
            +c_3\colvec{a_{1,3} \\ a_{2,3} \\ a_{3,3} \\ a_{4,3}}
            +c_4\colvec{a_{1,4} \\ a_{2,4} \\ a_{3,4} \\ a_{4,4}}
            +c_5\colvec{a_{1,5} \\ a_{2,5} \\ a_{3,5} \\ a_{4,5}}
             =\colvec{0 \\ 0 \\ 0 \\ 0}
      \end{equation*}
      and note that the resulting linear system 
      \begin{equation*}
        \begin{linsys}{5}
          a_{1,1}c_1 &+ &a_{1,2}c_2 &+ &a_{1,3}c_3 
              &+ &a_{1,4}c_4 &+ &a_{1,5}c_5 &=  &0   \\
          a_{2,1}c_1 &+ &a_{2,2}c_2 &+ &a_{2,3}c_3 
              &+ &a_{2,4}c_4 &+ &a_{2,5}c_5 &=  &0   \\
          a_{3,1}c_1 &+ &a_{3,2}c_2 &+ &a_{3,3}c_3 
              &+ &a_{3,4}c_4 &+ &a_{3,5}c_5 &=  &0   \\
          a_{4,1}c_1 &+ &a_{4,2}c_2 &+ &a_{4,3}c_3 
              &+ &a_{4,4}c_4 &+ &a_{4,5}c_5 &=  &0     
        \end{linsys}
      \end{equation*}
      has four equations and five unknowns, 
      so Gauss's Method must end with at least one \( c \) variable free,
      so there are infinitely many solutions,
      and so the above linear relationship among the four-tall vectors has
      more solutions than just the trivial solution.

      The smallest linearly independent set is the empty set.

      The biggest linearly dependent set is \( \Re^4 \).
      The smallest is \( \set{\zero} \).  
    \end{answer}
  \recommended \item 
    Linear independence and linear dependence are properties of sets.
    We can thus naturally ask how the properties of linear independence
    and dependence act with respect to
    the familiar elementary set relations and operations.  
    In this body of this subsection we have covered the subset and superset
    relations.
    We can also consider the operations of intersection, complementation, 
    and union.
    \begin{exparts}
       \partsitem How does linear independence relate to intersection:~can
         an intersection of linearly independent sets be independent?
         Must it be?
       \partsitem How does linear independence relate to complementation?
       \partsitem Show that the union of two linearly independent sets
          can be linearly independent.
       \partsitem Show that
          the union of two linearly independent sets need not be
          linearly independent.
    \end{exparts}
    \begin{answer}
      \begin{exparts}
        \partsitem The intersection of two linearly independent sets
          $S\intersection T$ must be linearly
          independent as it is a subset of the linearly independent set $S$
          (as well as the linearly independent set $T$ also, of course).
        \partsitem The complement of a linearly independent set is linearly
          dependent as it contains the zero vector.
        \partsitem A simple example in \( \Re^2 \) is these two sets.
          \begin{equation*}
            S=\set{\colvec{1 \\ 0}}
            \qquad
            T=\set{\colvec{0 \\ 1}}          
          \end{equation*}
          A somewhat subtler example, again in \( \Re^2 \), is these two.
          \begin{equation*}
            S=\set{\colvec{1 \\ 0}}
            \qquad
            T=\set{\colvec{1 \\ 0}, \colvec{0 \\ 1}}
          \end{equation*}
        \partsitem We must produce an example.
          One, in \( \Re^2 \), is
          \begin{equation*}
            S=\set{\colvec{1 \\ 0}}
            \qquad
            T=\set{\colvec{2 \\ 0}}
          \end{equation*}
          since the linear dependence of \( S_1\union S_2 \) is easy to see.
      \end{exparts}
    \end{answer}
  \item \textit{Continued from prior exercise.} 
    What is the interaction between the property of linear independence
    and the operation of union?
    \begin{exparts}
       \partsitem We might conjecture that the union $S\union T$ of 
          linearly independent sets 
          is linearly independent if and only if their spans have a trivial 
          intersection $\spanof{S}\intersection \spanof{T}=\set{\zero}$.
          What is wrong with this argument for the `if' direction
          of that conjecture? 
          ``If the union $S\union T$ is linearly independent then
          the only solution to          
          $c_1\vec{s}_1+\cdots+c_n\vec{s}_n
            +d_1\vec{t}_1+\cdots+d_m\vec{t}_m
            =\zero$
          is the trivial one $c_1=0$, \ldots, $d_m=0$. 
          So any member of the intersection of the spans
          must be the zero vector because in 
          $c_1\vec{s}_1+\cdots+c_n\vec{s}_n
            =d_1\vec{t}_1+\cdots+d_m\vec{t}_m$
          each scalar is zero.''
       \partsitem Give an example showing that the conjecture is false.
       \partsitem Find linearly independent sets \( S \) and \( T \)
          so that the union of \( S-(S\cap T)\) and \( T-(S\cap T) \)
          is linearly independent, but the union \( S\cup T \) is
          not linearly independent.
       \partsitem Characterize when the union of two linearly independent sets
          is linearly independent, in terms of the intersection of spans.
    \end{exparts}
    \begin{answer}
      \begin{exparts}
        \partsitem \nearbylemma{le:LDIffANonTrivLinRel} 
          requires that the vectors 
          $\vec{s}_1,\ldots,\vec{s}_n,\vec{t}_1,\ldots,\vec{t}_m$
          be distinct. 
          But we could have that the union $S\cup T$ is linearly 
          independent with some $\vec{s}_i$ equal to some $\vec{t}_j$.
        \partsitem One example in \( \Re^2 \) is these two.
          \begin{equation*}
            S=\set{\colvec{1 \\ 0}}
            \qquad
            T=\set{\colvec{1 \\ 0}, \colvec{0 \\ 1}}
          \end{equation*}
        \partsitem 
          An example from \( \Re^2 \) is these sets.
          \begin{equation*}
            S=\set{\colvec{1 \\ 0}, \colvec{0 \\ 1}}
            \qquad
            T=\set{\colvec{1 \\ 0}, \colvec{1 \\ 1}}
          \end{equation*}
        \partsitem The union of two linearly independent sets $S\union T$
          is linearly independent if and only if their spans of \( S \) 
          and \( T-(S\cap T)\) have a trivial intersection 
          $\spanof{S}\intersection \spanof{T-(S\cap T)}=\set{\zero}$.
          To prove that, assume that \( S \) and \( T \) are linearly 
          independent subsets of some vector space.

          For the `only if' direction, assume that the intersection of
          the spans is trivial 
          \( \spanof{S}\intersection \spanof{T-(S\cap T)}=\set{\zero} \).
          Consider the set $S\union (T-(S\cap T))=S\union T$ and
          consider the linear relationship 
          $c_1\vec{s}_1+\dots+c_n\vec{s}_n
            +d_1\vec{t}_1+\dots+d_m\vec{t}_m=\zero$.
          Subtracting gives
          $c_1\vec{s}_1+\dots+c_n\vec{s}_n=
            -d_1\vec{t}_1-\dots-d_m\vec{t}_m$.
          The left side of that equation sums to a vector in $\spanof{S}$, and
          the right side is a vector in $\spanof{T-(S\cap T)}$.
          Therefore, since the intersection of the spans is trivial, both
          sides equal the zero vector.
          Because $S$ is linearly independent, all of the $c$'s are zero.
          Because $T$ is linearly independent so also is  
          $T-(S\cap T)$ linearly independent, 
          and therefore all of the $d$'s are zero.
          Thus, the original linear relationship among members of 
          $S\union T$ only holds if all of the coefficients are zero.
          Hence, $S\union T$ is linearly independent.

          For the `if' half we can make the same argument in reverse.
          Suppose that the union $S\union T$ is linearly independent.
          Consider a linear relationship among members of 
          $S$ and $T-(S\cap T)$.        
          $c_1\vec{s}_1+\cdots+c_n\vec{s}_n
            +d_1\vec{t}_1+\cdots+d_m\vec{t}_m
            =\zero$
          Note that no $\vec{s}_i$ is equal to a $\vec{t}_j$
          so that is a combination of distinct vectors, as required by 
          \nearbylemma{le:LDIffANonTrivLinRel}.
          So the only solution 
          is the trivial one $c_1=0$, \ldots, $d_m=0$. 
          Since any vector $\vec{v}$ 
          in the intersection of the spans
          $\spanof{S}\intersection \spanof{T-(S\cap T)}$ 
          we can write 
          $\vec{v}=c_1\vec{s}_1+\cdots+c_n\vec{s}_n
            =-d_1\vec{t}_1-\cdots-d_m\vec{t}_m$,
          and it must be the zero vector because each scalar is zero.
      \end{exparts}  
    \end{answer}
  \item \label{exer:FillIndDetProofSetHasLISub}
    For \nearbycorollary{th:AlwaysAnLDSubset},
    \begin{exparts}
       \partsitem fill in the induction for the proof;
       \partsitem give an alternate proof that starts with the empty
         set and builds
         a sequence of linearly independent subsets of the given finite set
         until one appears with the same span as the given set.
    \end{exparts}
    \begin{answer}
      \begin{exparts}
         \partsitem We do induction on the number of vectors in the finite set
           \( S \).

           The base case is that $S$ has no elements.
           In this case $S$ is linearly independent and there is nothing to 
           check\Dash a subset of $S$ that has the same span as $S$ is $S$
           itself.

           For the inductive step assume that the theorem is true for all 
           sets of size $n=0$, $n=1$, \ldots, $n=k$ 
           in order to prove that it holds when \( S \) has $n=k+1$ elements.
           If the $k+1$-element set \( S=\set{\vec{s}_0,\dots,\vec{s}_{k}} \) 
           is linearly independent then the theorem is trivial,
           so assume that it is dependent.
           By \nearbycorollary{cor:LDMeansLC} there is an \( \vec{s}_i \)
           that is a linear combination of other vectors in \( S \).
           Define \( S_1=S-\set{\vec{s}_i} \) and note that 
           \( S_1 \) has the same span as \( S \) by
           \nearbycorollary{cor:OmitVecNotShrinkSpanIffVecIsDep}.
           The set \( S_1 \) has \( k \) elements and 
           so the inductive hypothesis
           applies to give that it has a linearly independent subset with the 
           same span.
           That subset of \( S_1 \) is the desired subset of \( S \).
         \partsitem Here is a sketch of the argument.
           We have left out the induction argument details.

           If the finite set \( S \) is empty then there is nothing to prove.
           If \( S=\set{\zero} \) then the empty subset will do.

           Otherwise, take some nonzero vector \( \vec{s}_1\in S \)
           and define \( S_1=\set{\vec{s}_1} \).
           If \( \spanof{S_1}=\spanof{S} \) then
           we are finished with this proof by noting that \( S_1 \) is linearly
           independent.

           If not, then there is a nonzero
           vector \( \vec{s}_2\in S-\spanof{S_1} \)
           (if every \( \vec{s}\in S \) is in \( \spanof{S_1} \) then
           \( \spanof{S_1}=\spanof{S} \)).
           Define \( S_2=S_1\union\set{\vec{s}_2} \).
           If \( \spanof{S_2}=\spanof{S} \) then 
           we are finished 
           by using \nearbytheorem{cor:LDMeansLC}
           to show that \( S_2 \) is linearly independent.

           Repeat the last paragraph until a set with a big enough
           span appears.
           That must eventually happen
           because \( S \) is finite, and
           \( \spanof{S} \) will be reached at worst when we have
           used every vector from
           \( S \).
      \end{exparts}  
    \end{answer}
  \item 
     With a some calculation we can get formulas to determine whether or
     not a set of vectors is linearly independent. 
     \begin{exparts}
       \partsitem Show that this subset of \( \Re^2 \)
         \begin{equation*}
           \set{\colvec{a \\ c},\colvec{b \\ d}}
         \end{equation*}
         is linearly independent if and only if \( ad-bc\neq 0 \).
       \partsitem Show that this subset of \( \Re^3 \)
         \begin{equation*}
           \set{\colvec{a \\ d \\ g},
                \colvec{b \\ e \\ h},
                \colvec{c \\ f \\ i}  }
         \end{equation*}
         is linearly independent iff
         \( aei+bfg+cdh-hfa-idb-gec \neq 0 \).
       \partsitem When is this subset of \( \Re^3 \)
         \begin{equation*}
           \set{\colvec{a \\ d \\ g},
                \colvec{b \\ e \\ h} }
         \end{equation*}
         linearly independent?
       \partsitem This is an opinion question:~for
         a set of four vectors from \( \Re^4 \),
         must there be a formula involving the sixteen entries 
         that determines independence of the set?
         (You needn't produce such a formula, just decide if one exists.)
    \end{exparts}
    \begin{answer}
      \begin{exparts}
        \partsitem 
          Assuming first that \( a\neq 0 \),
          \begin{equation*}
            x\colvec{a \\ c}
            +y\colvec{b \\ d}
            =\colvec{0 \\ 0}
          \end{equation*}
          gives
          \begin{equation*}
            \begin{linsys}{2}
               ax  &+  &by &=  &0  \\
               cx  &+  &dy &=  &0  
            \end{linsys}
            \grstep{-(c/a)\rho_1+\rho_2}
            \begin{linsys}{2}
               ax  &+  &by           &=  &0  \\
                   &   &(-(c/a)b+d)y &=  &0  
             \end{linsys}
          \end{equation*}
          which has a solution if and only if
          \( 0\neq-(c/a)b+d=(-cb+ad)/d \)
          (we've assumed in this case that \( a\neq 0 \), and so 
          back substitution yields a unique solution).

          The \( a=0 \) case is also not hard\Dash break it into the 
          \( c\neq 0 \) and \( c=0 \) subcases and 
          note that in these cases \( ad-bc=0\cdot d-bc \).

          \textit{Comment.}
          An earlier exercise showed that a two-vector set is linearly
          dependent if and only if either vector is a scalar multiple of the
          other.
          We could also use that to make the calculation.
        \partsitem The equation
          \begin{equation*}
            c_1\colvec{a \\ d \\ g}
            +c_2\colvec{b \\ e \\ h}
            +c_3\colvec{c \\ f \\ i}
            =\colvec{0 \\ 0 \\ 0}
          \end{equation*}
         expresses a homogeneous linear system.
         We proceed by writing it in matrix form and applying Gauss's Method.

         We first reduce the matrix to upper-triangular.
         Assume that \( a\neq 0 \).
         With that, we can clear down the first column.
         \begin{equation*}
           \grstep{(1/a)\rho_1}
           \begin{amat}{3}
              1   &b/a   &c/a  &0 \\
              d   &e     &f    &0 \\
              g   &h     &i    &0
            \end{amat}                                            
           \grstep[-g\rho_1+\rho_3]{-d\rho_1+\rho_2}
           \begin{amat}{3}
              1   &b/a           &c/a        &0   \\
              0   &(ae-bd)/a     &(af-cd)/a  &0   \\
              0   &(ah-bg)/a     &(ai-cg)/a  &0
            \end{amat}                                            
         \end{equation*}
         Then we get a $1$ in the second row, second column entry.
         (Assuming for the moment that \( ae-bd\neq 0 \), in order
         to do the row reduction step.)
         \begin{equation*}
           \grstep{(a/(ae-bd))\rho_2}
           \begin{amat}{3}
              1   &b/a           &c/a             &0  \\
              0   &1             &(af-cd)/(ae-bd) &0  \\
              0   &(ah-bg)/a     &(ai-cg)/a       &0
            \end{amat}
         \end{equation*}
         Then, under the assumptions, we perform
         the row operation $((ah-bg)/a)\rho_2+\rho_3$
         to get this.
         \begin{equation*}
           % \grstep{((ah-bg)/a)\rho_2+\rho_3}
           \begin{amat}{3}
              1   &b/a   &c/a                              &0 \\
              0   &1     &(af-cd)/(ae-bd)                   &0 \\
              0   &0     &(aei+bgf+cdh-hfa-idb-gec)/(ae-bd) &0
            \end{amat}
         \end{equation*}
         Therefore, the original system is nonsingular
         if and only if the above \( 3,3 \) entry is nonzero
         (this fraction is defined because of the \( ae-bd\neq 0 \) assumption).
         It equals zero if and only if the numerator is zero. 

         We next worry about the assumptions.
         First, if \( a\neq 0 \) but \( ae-bd=0 \) then we swap
         \begin{multline*}
           \begin{amat}{3}
              1   &b/a           &c/a        &0   \\
              0   &0             &(af-cd)/a  &0   \\
              0   &(ah-bg)/a     &(ai-cg)/a  &0
            \end{amat}                                    \\
           \grstep{\rho_2\leftrightarrow\rho_3}
           \begin{amat}{3}
              1   &b/a           &c/a        &0   \\
              0   &(ah-bg)/a     &(ai-cg)/a  &0   \\
              0   &0             &(af-cd)/a  &0
            \end{amat}
         \end{multline*}
         and conclude that the system is nonsingular if and only if either
         \( ah-bg=0 \) or \( af-cd=0 \).
         That's the same as asking that their product be zero:
         \begin{align*}
            ahaf-ahcd-bgaf+bgcd
            &=0                   \\
            ahaf-ahcd-bgaf+aegc
            &=0                   \\
            a(haf-hcd-bgf+egc)
            &=0
         \end{align*}
         (in going from the first line to the second we've applied the
         case assumption that $ae-bd=0$ by substituting $ae$ for $bd$).
         Since we are assuming that \( a\neq 0 \), 
         we have that \( haf-hcd-bgf+egc=0 \).
         With $ae-bd=0$ we can rewrite this to fit the form we need:~in
         this \( a\neq 0 \) and \( ae-bd=0 \) case, the given system
         is nonsingular when
         \( haf-hcd-bgf+egc-i(ae-bd)=0 \), as required.

         The remaining cases have the same character.
         Do the \( a=0 \) but \( d\neq 0 \) case and the \( a=0 \) and
         \( d=0 \) but \( g\neq 0 \) case by first swapping rows and
         then going on as above.
         The \( a=0 \), \( d=0 \), and \( g=0 \) case is easy\Dash a set with a
         zero vector is linearly dependent, and the formula comes out
         to equal zero.
       \partsitem It is linearly dependent if and only if either vector is a
         multiple of the other.
         That is, it is not independent iff
         \begin{equation*}
           \colvec{a \\ d \\ g}=r\cdot\colvec{b \\ e \\ h}
           \quad\text{or}\quad
           \colvec{b \\ e \\ h}=s\cdot\colvec{a \\ d \\ g}
         \end{equation*}
         (or both) for some scalars $r$ and $s$.
         Eliminating $r$ and $s$ in order to restate this condition only in
         terms of the given letters $a$, $b$, $d$, $e$, $g$, $h$, we have that 
         it is not independent\Dash it is dependent\Dash iff
         \( ae-bd=ah-gb=dh-ge \).
       \partsitem Dependence or independence is a function of the
         indices, so there
         is indeed a formula (although at first glance a person might think
         the formula involves cases: ``if the first component of the first
         vector is zero then \ldots'', this guess turns out not to be 
         correct).
      \end{exparts}  
    \end{answer}
  \recommended \item  
    \begin{exparts}
      \partsitem Prove that a set of two perpendicular 
        nonzero vectors from
        \( \Re^n \) is linearly independent when \( n>1 \).
      \partsitem What if \( n=1 \)?
        \( n=0 \)?
      \partsitem Generalize to more than two vectors.
    \end{exparts}
    \begin{answer}
      Recall that two vectors from \( \Re^n \) are perpendicular if and
      only if their dot product is zero.
      \begin{exparts}
         \partsitem Assume that \( \vec{v} \) and \( \vec{w} \) are
           perpendicular nonzero vectors in $\Re^n$, with \( n>1 \).
           With the linear relationship \( c\vec{v}+d\vec{w}=\zero \), 
           apply \( \vec{v} \) to both
           sides to conclude that \( c\cdot\norm{\vec{v}}^2+d\cdot 0=0 \).
           Because \( \vec{v}\neq\zero \) we have that \( c=0 \).
           A similar application of \( \vec{w} \) shows that \( d=0 \).
         \partsitem Two vectors in \( \Re^1 \) are perpendicular if and only if
           at least one of them is zero.

           We define \( \Re^0 \) to be a trivial space, and so both $\vec{v}$
           and $\vec{w}$ are the zero vector.
         \partsitem The right generalization is to look at a set
           \( \set{\vec{v}_1,\dots,\vec{v}_n}\subseteq\Re^k \) of vectors
           that are \definend{mutually orthogonal} 
           (also called \definend{pairwise perpendicular}):~if
           \( i\neq j \) then \( \vec{v}_i \) is perpendicular to
           \( \vec{v}_j \).
           Mimicking the proof of the first item above shows that such a set of
           nonzero vectors is linearly independent.
      \end{exparts}  
    \end{answer}
  \item 
    Consider the set of functions from the interval 
    $\openinterval{-1}{1}\subseteq \Re$  
    to $\Re$.
    \begin{exparts}
      \partsitem Show that 
         this set is a vector space under the usual operations.
      \partsitem Recall the formula for the sum of an infinite geometric 
         series: 
         \( 1+x+x^2+\cdots=1/(1-x) \) for all \( x\in(-1..1) \).
         Why does this not express a dependence inside of
         the set $\set{g(x)=1/(1-x),f_0(x)=1,f_1(x)=x,f_2(x)=x^2,\ldots}$
         (in the vector space that we are considering)?
         (\textit{Hint.}
         Review the definition of linear combination.)
      \partsitem Show that the set in the prior item is linearly independent.
    \end{exparts}
    This shows that some vector spaces exist with linearly independent subsets
    that are infinite.
    \begin{answer}
      \begin{exparts}
        \partsitem This check is routine.
        \partsitem The summation is infinite (has infinitely many summands).
          The definition of linear combination involves only finite sums.
        \partsitem No nontrivial finite sum of members of 
           \( \set{g,f_0,f_1,\ldots} \) adds to the zero object:~assume that
           \begin{equation*}
              c_0\cdot (1/(1-x))+c_1\cdot 1+\dots+c_n\cdot x^n=0
           \end{equation*}
           (any finite sum uses a highest power, here \( n \)).
           Multiply both sides by \( 1-x \) to conclude that each coefficient 
           is zero, because a polynomial describes the zero function only when 
           it is the zero polynomial.
      \end{exparts}
     \end{answer}
  \item 
    Show that, where \( S \) is a subspace of \( V \), if a subset $T$ of
    \( S \) is linearly independent in \( S \) then $T$ is also linearly
    independent in \( V \).
    Is that `only if'?
    \begin{answer}
      It is both `if' and `only if'.

      Let \( T \) be a subset of the subspace \( S \) of the vector space
      \( V \).
      The assertion that any linear relationship 
      $c_1\vec{t}_1+\dots+c_n\vec{t}_n=\zero$ among members of \( T \)
      must be the trivial relationship $c_1=0$, \ldots, $c_n=0$
      is a statement that 
      holds in \( S \) if and only if it holds in \( V \),
      because the subspace \( S \) inherits its addition and 
      scalar multiplication operations from \( V \).  
    \end{answer}
\end{exercises}

% END NATIVE SOURCE src/vs/vs2.tex

% BEGIN NATIVE SOURCE src/vs/vs3.tex
% Chapter 2, Section 3 _Linear Algebra_ Jim Hefferon
%  http://joshua.smcvt.edu/linearalgebra
%  2001-Jun-10
\section{Basis and Dimension}
\index{basis|(}
The prior section ends with the observation that
a spanning set is minimal when it is linearly independent and
a linearly independent set is maximal when it spans the space.
So the notions of minimal spanning set and maximal independent set 
coincide.
In this section we will name this idea and study its properties.





\subsection{Basis}
\begin{definition} \label{def:basis}
%<*df:basis>
A \definend{basis}\index{basis!definition}\index{vector space!basis} 
for a vector
space is a sequence of vectors that is linearly independent 
and that spans the space.
%</df:basis>
\end{definition}

%<*BasisNotation>
Because a basis is a 
sequence, 
meaning that 
bases are different if they contain the same elements but in different
orders,
% the 
% order of the elements is significant,
we denote it with angle brackets
\( \sequence{\vec{\beta}_1,\vec{\beta}_2,\ldots} \).\appendrefs{sequences}\@ %\spacefactor=1000{}
%</BasisNotation>
(A sequence is linearly independent if the multiset
consisting of the elements of the sequence is independent.
Similarly, 
a sequence spans the space if the set of elements 
of the sequence spans the space.)

\begin{example}  \label{ex:FstBasisRTwo}
This is a basis for \( \Re^2 \).
\begin{equation*}
  \sequence{ \colvec{2 \\ 4},\colvec{1 \\ 1} }
\end{equation*}
It is linearly independent
\begin{equation*}
  c_1\colvec{2 \\ 4}+c_2\colvec{1 \\ 1}=\colvec{0 \\ 0}
  \quad\implies\quad
  \begin{linsys}{2}
     2c_1  &+  &1c_2  &=  &0  \\
     4c_1  &+  &1c_2  &=  &0  
   \end{linsys}
  \quad\implies\quad
  c_1=c_2=0
\end{equation*}
and it spans \( \Re^2 \).
\begin{equation*}
  \begin{linsys}{2}
    2c_1  &+  &1c_2  &=  &x  \\
    4c_1  &+  &1c_2  &=  &y  
  \end{linsys}
  \quad\implies\quad
  c_2=2x-y\text{\ and\ } c_1=(y-x)/2
\end{equation*}
\end{example}

\begin{example}
This basis for \( \Re^2 \) differs from the prior one 
\begin{equation*}
  \sequence{\colvec[r]{1 \\ 1},\colvec[r]{2 \\ 4}}
\end{equation*}
because it is in a different order.
The verification that it is a basis is just as in the prior example.
\end{example}

\begin{example}
The space \( \Re^2 \) has many bases.
Another one is this.
\begin{equation*}
  \sequence{ \colvec[r]{1 \\ 0},\colvec[r]{0 \\ 1} }
\end{equation*}
The verification is easy.
\end{example}

\begin{definition} \label{df:StandardBasis}
%<*df:StandardBasis>
For any \( \Re^n \)
\begin{equation*}
   \stdbasis_n=\sequence{
     \colvec[r]{1 \\ 0 \\ \vdotswithin{0} \\ 0},
     \colvec[r]{0 \\ 1 \\ \vdotswithin{0} \\ 0},
     \dots,\,
     \colvec[r]{0 \\ 0 \\ \vdotswithin{0} \\ 1}}
\end{equation*}
is the \definend{standard}\index{standard basis}%
\index{basis!standard} (or \definend{natural}) basis.
We denote these vectors \( \vec{e}_1,\dots,\vec{e}_n \).
%</df:StandardBasis>
\end{definition}

\noindent 
Calculus books denote $\Re^2$'s standard basis vectors as
\( \vec{\imath} \) and \( \vec{\jmath} \) instead of $\vec{e}_1$
and $\vec{e}_2$ and they denote to 
\( \Re^3 \)'s standard basis vectors as
\( \vec{\imath} \), \( \vec{\jmath} \), and \( \vec{k} \)
instead of $\vec{e}_1$, $\vec{e}_2$, and $\vec{e}_3$.
Note that \( \vec{e}_1 \) means something different in a
discussion of \( \Re^3 \) than it means in a discussion of \( \Re^2 \).

\begin{example}  \label{ex:BasisForCosPlusSin}
Consider the space
\( \set{a\cdot\cos\theta+b\cdot\sin\theta\suchthat a,b\in\Re} \)
of functions of the real variable $\theta$.
This is a natural basis 
$ \sequence{\cos\theta, \sin\theta}=\sequence{1\cdot\cos\theta+0\cdot\sin\theta,
             0\cdot\cos\theta+1\cdot\sin\theta}
    $.
A more generic basis for this space is
\( \sequence{\cos\theta-\sin\theta,
             2\cos\theta+3\sin\theta} \).
Verification that these two are bases is
\nearbyexercise{exer:VerifBasesCosPlusSin}.
\end{example}

\begin{example}
A natural basis for the vector space of cubic polynomials \( \polyspace_3 \)
is \( \sequence{1,x,x^2,x^3} \).
Two other bases for this space are \( \sequence{x^3,3x^2,6x,6} \)
and \( \sequence{1,1+x,1+x+x^2,1+x+x^2+x^3} \).
Checking that each is linearly independent and spans the space is easy.
\end{example}

\begin{example}
The trivial space\index{trivial space}\index{vector space!trivial} 
$\set{\zero}$ has only one basis, the empty one
\( \sequence{} \).
\end{example}

\begin{example}
The space of finite-degree polynomials has a basis with infinitely many
elements
\( \sequence{1,x,x^2,\ldots} \).
\end{example}

\begin{example}
We have seen bases before.
In the first chapter we described the solution set of homogeneous systems
such as this one
\begin{equation*}
   \begin{linsys}{4}
      x  &+  &y  &   &   &-   &w   &=  &0  \\
         &   &   &   &z  &+   &w   &=  &0  
   \end{linsys}
\end{equation*}
by parametrizing.
\begin{equation*}
  \set{\colvec[r]{-1 \\ 1 \\ 0 \\ 0}y
       +\colvec[r]{1 \\ 0 \\ -1 \\ 1}w
       \suchthat y,w\in\Re }
\end{equation*}
Thus the vector space of solutions is 
the span of a two-element set.
This two-vector set is also linearly independent, which is easy to check.
Therefore the solution set is a subspace of \( \Re^4 \) with a
basis comprised of these two vectors.
\end{example}

\begin{example}
Parametrization finds bases for other vector spaces, not just
for solution sets of homogeneous systems.
To find a basis for this subspace of $\matspace_{\nbyn{2}}$
\begin{equation*}
  \set{\begin{mat}
         a  &b  \\
         c  &0
       \end{mat} \suchthat a+b-2c=0}
\end{equation*}
we rewrite the condition as $a=-b+2c$.
\begin{equation*}
  \set{\begin{mat}
         -b+2c  &b  \\
          c     &0
       \end{mat} \suchthat b,c \in \Re}
  =\set{b\begin{mat}[r]
         -1  &1  \\
          0  &0
       \end{mat}+
       c\begin{mat}[r]
          2  &0  \\
          1  &0
       \end{mat} \suchthat b,c \in \Re}
\end{equation*}
Thus, this is a natural candidate for a basis.
\begin{equation*}
  \sequence{\begin{mat}[r]
         -1  &1  \\
          0  &0
       \end{mat},
       \begin{mat}[r]
          2  &0  \\
          1  &0
       \end{mat} }
\end{equation*}
The above work shows that it spans the space.
Linear independence is also easy.
\end{example}

Consider again \nearbyexample{ex:FstBasisRTwo}.
% To verify linearly independence we looked at 
% linear combinations of the set's members that total to
% the zero vector 
% $c_1\vec{\beta}_1+c_2\vec{\beta}_2=\binom{0}{0}$.
% The resulting 
% calculation shows that such a combination is unique,
% that $c_1$ must be $0$ and $c_2$ must be $0$.
To verify that the set spans the space we
looked at linear combinations that total to a
member of the space
$c_1\vec{\beta}_1+c_2\vec{\beta}_2=\binom{x}{y}$.
We only noted in that example that such a combination 
exists, that for each $x,y$ there exists a $c_1,c_2$, but
in fact the calculation also shows that the combination is 
unique:~$c_1$ must be $(y-x)/2$ and $c_2$ must be $2x-y$.
 
\begin{theorem} \label{th:BasisIffUniqueRepWRT}
%<*th:BasisIffUniqueRepWRT>
In any vector space, a subset is a basis\index{basis}
if and only if each vector in the
space can be expressed as a linear combination of elements of the subset 
in one and only one way.
%</th:BasisIffUniqueRepWRT>
\end{theorem}

\noindent We consider linear combinations to be the same if they have the 
same summands but in a different order,
or if they differ only in the addition or deletion of terms of the form 
`\( 0\cdot\vec{\beta} \)'.


\begin{proof}
%<*pf:BasisIffUniqueRepWRT0>
A sequence is a basis if and only if its vectors form a set
that spans and that is linearly independent.
A subset is a spanning set if and only if each vector in the space is a linear
combination of elements of that subset in at least one way.
Thus we need only show 
that a spanning subset is linearly independent 
if and only if every vector in the space
is a linear combination of elements from the subset in at most one way.
%</pf:BasisIffUniqueRepWRT0>

%<*pf:BasisIffUniqueRepWRT1>
Consider two expressions of a vector as a linear combination of the
members of the subset.
Rearrange the two sums, and if necessary
add some \( 0\cdot\vec{\beta}_i \) terms, so that the two sums
combine the same \( \vec{\beta} \)'s in the same order:
\( \vec{v}=\lincombo{c}{\vec{\beta}} \) and 
\( \vec{v}=\lincombo{d}{\vec{\beta}} \). 
Now
\begin{equation*}
   \lincombo{c}{\vec{\beta}}=\lincombo{d}{\vec{\beta}}
\end{equation*}
holds if and only if
\begin{equation*}
   (c_1-d_1)\vec{\beta}_1+\dots+(c_n-d_n)\vec{\beta}_n=\zero
\end{equation*}
holds. 
So, asserting that 
each coefficient in the lower equation is zero is 
the same thing as asserting that \( c_i=d_i \) for each \( i \),
that is, that every vector is expressible as a linear combination of
the \( \vec{\beta} \)'s in a unique way.
%</pf:BasisIffUniqueRepWRT1>
\end{proof}

\begin{definition} \label{def:RepresentingVectors}
%<*df:RepresentingVectors>
In a vector space with basis $B$
the \definend{representation of \( \vec{v} \) with respect to \( B \)}%
\index{representation!of a vector}\index{vector!representation of} is
the column vector of the coefficients used to express $\vec{v}$ as a 
linear combination of the basis vectors:
\begin{equation*}
  \rep{\vec{v}}{B}
  =
  \colvec{c_1 \\ c_2 \\ \vdots \\ c_n}_{B}
\end{equation*}
where
\( B=\sequence{\vec{\beta}_1,\dots,\vec{\beta}_n} \) and 
\( \vec{v}=\lincombo{c}{\vec{\beta}} \).
The \( c \)'s are the
\definend{coordinates of \( \vec{v} \) with respect to \( B \)}%
\index{coordinates!with respect to a basis}.
%</df:RepresentingVectors>
\end{definition}

\begin{example}  \label{ex:RepPolyWRTA Basis}
In \( \polyspace_3 \), with respect to the basis
\( B=\sequence{1,2x,2x^2,2x^3} \),
the representation of \( x+x^2 \) is
\begin{equation*}
  \rep{x+x^2}{B}=\colvec[r]{0 \\ 1/2 \\ 1/2 \\ 0}_B
\end{equation*}
because $x+x^2=0\cdot 1+(1/2)\cdot 2x+(1/2)\cdot 2x^2+0\cdot 2x^3$.
With respect to a different basis \( D=\sequence{1+x,1-x,x+x^2,x+x^3} \),
the representation is different.
\begin{equation*}
  \rep{x+x^2}{D}=\colvec[r]{0 \\ 0 \\ 1 \\ 0}_D
\end{equation*}
\end{example}

\begin{remark}
\nearbydefinition{def:basis} requires that a basis be a sequence
so that we can write these coordinates in an order. 
\end{remark}

When there is only one basis around, we often omit the subscript
naming that basis.

\begin{example}
In \( \Re^2 \), to find
the coordinates of the vector $\vec{v}=\binom{3}{2}$ with respect to the basis
\begin{equation*}
  B=\sequence{
              \colvec[r]{1 \\ 1},
              \colvec[r]{0 \\ 2} }
\end{equation*}
solve
\begin{equation*}
  c_1\colvec[r]{1 \\ 1}
  +c_2\colvec[r]{0 \\ 2}
  =
  \colvec[r]{3 \\ 2}
\end{equation*}
and get that $c_1=3$ and $c_2=-1/2$.
\begin{equation*}
  \rep{\vec{v}}{B}=\colvec[r]{3 \\ -1/2}
\end{equation*}  
\end{example}

Writing the representation as a column
generalizes the familiar case:~in \( \Re^n \) 
and with respect to the standard
basis \( \stdbasis_n \), the vector starting at the origin and ending at
\( (v_1,\dots,v_n) \) has this representation.
\begin{equation*}
  \rep{\colvec{v_1 \\ \vdots \\ v_n}}{\stdbasis_n}
    =
  \colvec{v_1 \\ \vdots \\ v_n}_{\stdbasis_n}
\end{equation*}
This is an example.
\begin{equation*}
  \rep{\colvec{-1 \\ 1}}{\stdbasis_{n}}
  =\colvec{-1 \\ 1}
\end{equation*}

\begin{remark}  \label{rem:RepNotation}
The $\rep{\vec{v}}{B}$ notation is not standard.
The most common notation is $[\vec{v}]_B$ but one advantage that
$\rep{\vec{v}}{B}$ has is that it is harder to misinterpret or overlook.
\end{remark}

The column represents the vector in the sense that
a linear relationship holds among a set of vectors if and only if
that relationship holds among the set of representations.

\begin{lemma}  \label{lem:LinearRelRepresented}
%<*lm:LinearRelRepresented>
Where $B$ is a basis with $n$~elements, for any set of vectors,
$a_1\vec{v}_1+\cdots+a_k\vec{v}_k=\vec{0}_{V}$ if and only if
$a_1\rep{\vec{v}_1}{B}+\cdots+a_k\rep{\vec{v}_k}{B}=\vec{0}_{\Re^n}$.
%</lm:LinearRelRepresented>
\end{lemma}

\begin{proof}
Fix a basis \( B=\sequence{\vec{\beta}_1,\dots,\vec{\beta}_n} \)
and suppose
\begin{equation*}
  \rep{\vec{v}_1}{B}=\colvec{c_{1,1} \\ \vdots \\ c_{n,1}}
  \quad\ldots\quad
  \rep{\vec{v}_k}{B}=\colvec{c_{1,k} \\ \vdots \\ c_{n,k}}
\end{equation*}
so that $\vec{v}_1=c_{1,1}\vec{\beta}_1+\dots+c_{n,1}\vec{\beta}_n$, etc.
Then $a_1\vec{v}_1+\dots+a_k\vec{v}_k=\vec{0}$ is equivalent to these.
\begin{align*}
  \vec{0}
  &=a_1\cdot(c_{1,1}\vec{\beta}_1+\dots+c_{n,1}\vec{\beta}_n)
    +\dots+
    a_k\cdot(c_{1,k}\vec{\beta}_1+\dots+c_{n,k}\vec{\beta}_n)  \\
  &=(a_1c_{1,1}+\dots+a_kc_{1,k})\cdot\vec{\beta}_1
    +\dots+
    (a_1c_{n,1}+\dots+a_kc_{n,k})\cdot\vec{\beta}_n 
\end{align*}
Obviously the bottom equation is true if the coefficients are zero.  
But, because \( B \) is a basis, \nearbytheorem{th:BasisIffUniqueRepWRT}
says that the bottom equation is true if and only if the coefficients are zero.
So the relation is equivalent to this.
\begin{align*}
    a_1c_{1,1}+\dots+a_kc_{1,k} &=0    \\
                               &\alignedvdots \\ 
    a_1c_{n,1}+\dots+a_kc_{n,k} &=0
\end{align*}
This is the equivalent recast into column vectors.
\begin{equation*}
  a_1\colvec{c_{1,1} \\ \vdots \\ c_{n,1}}
  +\dots+
  a_k\colvec{c_{1,k} \\ \vdots \\ c_{n,k}}
  =\colvec{0 \\ \vdots \\ 0}
\end{equation*}
Note that not only does a relationship hold for one set if and only if it
holds for the other, but it is the same relationship\Dash the 
\( a_i \) are the same.
\end{proof}

\begin{example}
\nearbyexample{ex:RepPolyWRTA Basis} finds the 
representation of \( x+x^2\in\polyspace_3 \) with respect to
\( B=\sequence{1,2x,2x^2,2x^3} \).
\begin{equation*}
  \rep{x+x^2}{B}=\colvec[r]{0 \\ 1/2 \\ 1/2 \\ 0}_B
\end{equation*}
This relationship 
\begin{equation*}
  2\cdot(x+x^2)-1\cdot(2x)-2\cdot(x^2) = 0+0x+0x^2+0x^3
\end{equation*}
is represented by this one.
\begin{equation*}
  2\cdot\rep{x+x^2}{B}-\rep{2x}{B}-2\cdot\rep{x^2}{B}
  =2\cdot\colvec{0 \\ 1/2 \\ 1/2 \\ 0}
   -\colvec{0 \\ 1 \\ 0 \\ 0}
   -2\cdot\colvec{0 \\ 0 \\ 1/2 \\ 0}
  =\colvec{0 \\ 0 \\ 0 \\ 0}
\end{equation*}

\end{example}

Our main use of representations will come later but
the definition appears here because the fact that every vector is a linear
combination of basis vectors in a unique way is a crucial property of bases,
and also to help make a point.
For calculation of coordinates among other things, we shall
restrict our attention to spaces with bases having only finitely many elements.
That will start in the next subsection.

\begin{exercises}
  \recommended \item Decide if each is a basis for $\polyspace_2$.
    \begin{exparts*}
      \partsitem $\sequence{x^2-x+1, 2x+1, 2x-1}$
      \partsitem $\sequence{x+x^2, x-x^2}$
    \end{exparts*}
    \begin{answer}
      \begin{exparts}
        \partsitem This is a basis for $\polyspace_2$.
           To show that it spans the space we consider a generic
           $a_2x^2+a_1x+a_0\in\polyspace_2$ and look for 
           scalars $c_1, c_2, c_3\in\Re$ such that
           $a_2x^2+a_1x+a_0=c_1\cdot(x^2-x+1) +c_2\cdot(2x+1) +c_3(2x-1)$.
           The resulting linear system is this.
           \begin{equation*}
             \begin{linsys}{3}
               c_1  &\spaceforemptycolumn  &     &  &     &=  &a_2 \\
                    &  &2c_2 &+ &2c_3 &=  &a_1 \\
                    &  &c_2  &- &c_3   &=  &a_0
             \end{linsys}
             \grstep{(-1/2)\rho_2+\rho_3}
             \begin{linsys}{3}
               c_1  &\spaceforemptycolumn  &     &  &     &=  &a_2\hfill\hbox{} \\
                    &  &2c_2 &+ &2c_3    &=  &a_1\hfill\hbox{} \\
                    &  &     &  &-2c_3   &=  &a_0-(1/2)a_1
             \end{linsys}
           \end{equation*}
           Thus, given any triple of $a_i$'s, we can compute the $c_j$'s as
           $c_1=a_2$, $c_2=(1/4)a_1+(1/2)a_0$, and $c_3=(1/4)a_1-(1/2)a_0$.
           Therefore each element of $\polyspace_2$ is a combination of the 
           given $x^2-x+1$, $2x+1$, and~$2x-1$.

           To prove that the set of the given three is linearly independent
           we can set up the equation 
           $0x^2+0x+0=c_1\cdot(x^2-x+1) +c_2\cdot(2x+1) +c_3(2x-1)$
           and solve, and it will give that $c_1=0$, $c_2=0$, and $c_3=0$.
           Or, we can instead
           observe that the solution in the prior paragraph is
           unique, and cite \nearbytheorem{th:BasisIffUniqueRepWRT}.
        \partsitem This is not a basis.
           It does not span the space since no combination of the two
           $c_1\cdot (x+x^2) +c_2\cdot (x-x^2)$ will sum to the polynomial 
           $3\in\polyspace_2$.  
      \end{exparts}
    \end{answer}
  \recommended \item 
    Decide if each is a basis for \( \Re^3 \).
    \begin{exparts*}
      \partsitem \( \sequence{
                 \colvec[r]{1 \\ 2 \\ 3},
                 \colvec[r]{3 \\ 2 \\ 1},
                 \colvec[r]{0 \\ 0 \\ 1}}  \)
      \partsitem \( \sequence{
                 \colvec[r]{1 \\ 2 \\ 3},
                 \colvec[r]{3 \\ 2 \\ 1}}  \)
      \partsitem \( \sequence{
                 \colvec[r]{0 \\ 2 \\ -1},
                 \colvec[r]{1 \\ 1 \\ 1},
                 \colvec[r]{2 \\ 5 \\ 0}}  \)
      \partsitem \( \sequence{
                 \colvec[r]{0 \\ 2 \\ -1},
                 \colvec[r]{1 \\ 1 \\ 1},
                 \colvec[r]{1 \\ 3 \\ 0}}  \)
    \end{exparts*}
    \begin{answer}
      By \nearbytheorem{th:BasisIffUniqueRepWRT}, each is a basis if and only
      if we can express each vector in the space in a unique way as a linear
      combination of the given vectors.
      \begin{exparts}
        \partsitem Yes this is a basis.
          The relation
          \begin{equation*}
            c_1\colvec[r]{1 \\ 2 \\ 3}
            +c_2\colvec[r]{3 \\ 2 \\ 1}
            +c_3\colvec[r]{0 \\ 0 \\ 1}
            =\colvec{x \\ y \\ z}
          \end{equation*}
          gives
          \begin{equation*}
            \begin{amat}{3}
              1  &3  &0  &x  \\
              2  &2  &0  &y  \\
              3  &1  &1  &z
            \end{amat}
            \grstep[-3\rho_1+\rho_3]{-2\rho_1+\rho_2}
            \repeatedgrstep{-2\rho_2+\rho_3}
            \begin{amat}{3}
              1  &3  &0  &x \hfill\hbox{}  \\
              0  &-4 &0  &-2x+y \hfill\hbox{}  \\
              0  &0  &1  &x-2y+z 
            \end{amat}
          \end{equation*}
          which has the unique solution
          \( c_3=x-2y+z \), \( c_2=x/2-y/4 \), and
          \( c_1=-x/2+3y/4 \).
        \partsitem This is not a basis.
          Setting it up as in the prior item
          \begin{equation*}
            c_1\colvec[r]{1 \\ 2 \\ 3}
            +c_2\colvec[r]{3 \\ 2 \\ 1}
            =\colvec{x \\ y \\ z}
          \end{equation*}
          gives a linear system whose solution
          \begin{equation*}
            \begin{amat}{2}
              1  &3  &x  \\ 
              2  &2  &y  \\
              3  &1  &z
            \end{amat}
            \grstep[-3\rho_1+\rho_3]{-2\rho_1+\rho_2}
            \repeatedgrstep{-2\rho_2+\rho_3}
            \begin{amat}{2}
              1  &3  &x\hfill\hbox{}  \\ 
              0  &-4 &-2x+y\hfill\hbox{}  \\
              0  &0  &x-2y+z
            \end{amat}
          \end{equation*}
          is possible if and only if the three-tall vector's components
          $x$, $y$, and $z$ satisfy $x-2y+z=0$.
          For instance, we can find the coefficients $c_1$ and $c_2$ that
          work when $x=1$, $y=1$, and $z=1$.
          However, there are no $c$'s that work for
          $x=1$, $y=1$, and $z=2$.
          Thus this is not a basis; it does not span the space.
        \partsitem Yes, this is a basis.
         Setting up the relationship leads to this reduction
          \begin{equation*}
            \begin{amat}{3}
              0  &1  &2  &x \hfill\hbox{} \\
              2  &1  &5  &y \hfill\hbox{} \\
             -1  &1  &0  &z 
            \end{amat}
            \grstep{\rho_1\swap\rho_3}
            \repeatedgrstep{2\rho_1+\rho_2}
            \repeatedgrstep{-(1/3)\rho_2+\rho_3}
            \begin{amat}{3}
             -1  &1  &0   &z  \hfill\hbox{} \\
              0  &3  &5   &y+2z \hfill\hbox{} \\
              0  &0  &1/3 &x-y/3-2z/3
            \end{amat}
          \end{equation*}
          which has a unique solution for each triple of components
          $x$, $y$, and $z$.
        \partsitem No, this is not a basis.
          The reduction  
          \begin{equation*}
            \begin{amat}{3}
              0  &1  &1  &x   \\
              2  &1  &3  &y   \\
             -1  &1  &0  &z  
            \end{amat}
            \grstep{\rho_1\swap\rho_3}
            \repeatedgrstep{2\rho_1+\rho_2}
            \repeatedgrstep{(-1/3)\rho_2+\rho_3}
            \begin{amat}{3}
             -1  &1  &0  &z  \hfill\hbox{} \\
              0  &3  &3  &y+2z  \hfill\hbox{} \\
              0  &0  &0  &x-y/3-2z/3
            \end{amat}
          \end{equation*}
          which does not have a solution for each triple $x$, $y$, and $z$.
          Instead, the span of the given set 
          includes only those three-tall vectors where \( x=y/3+2z/3 \).
      \end{exparts}  
     \end{answer}
  \recommended \item 
    Represent the vector with respect to the basis.
    \begin{exparts}
      \partsitem \( \colvec[r]{1 \\ 2} \),
        \( B=\sequence{\colvec[r]{1 \\ 1},\colvec[r]{-1 \\ 1}}\subseteq\Re^2 \)
      \partsitem \( x^2+x^3 \),
        \( D=\sequence{1,1+x,1+x+x^2,1+x+x^2+x^3}\subseteq\polyspace_3 \)
      \partsitem \( \colvec[r]{0 \\ -1 \\ 0 \\ 1} \),
        \( \stdbasis_4\subseteq\Re^4 \)
    \end{exparts}
    \begin{answer}  
      \begin{exparts}
        \partsitem We solve
          \begin{equation*}
            c_1\colvec[r]{1 \\ 1}
            +c_2\colvec[r]{-1 \\ 1}
            =\colvec[r]{1 \\ 2}
          \end{equation*}
          with
          \begin{equation*}
             \begin{amat}[r]{2}
               1  &-1  &1  \\
               1  &1   &2
             \end{amat}
             \grstep{-\rho_1+\rho_2}
             \begin{amat}[r]{2}
               1  &-1  &1  \\
               0  &2   &1
             \end{amat}
          \end{equation*}
          and conclude that \( c_2=1/2 \) and so \( c_1=3/2 \).
          Thus, the representation is this.
          \begin{equation*}
            \rep{\colvec[r]{1 \\ 2}}{B}=\colvec[r]{3/2 \\ 1/2}_B
          \end{equation*}
        \partsitem The relationship
           $c_1\cdot(1)+c_2\cdot(1+x)+c_3\cdot(1+x+x^2)+c_4\cdot(1+x+x^2+x^3)
             =x^2+x^3$
           is easily solved by eye to give that $c_4=1$, $c_3=0$, $c_2=-1$, and
           $c_1=0$.
           \begin{equation*}
              \rep{x^2+x^3}{D}=\colvec[r]{0 \\ -1 \\ 0 \\ 1}_D
           \end{equation*} 
        \partsitem \( \rep{\colvec[r]{0 \\ -1 \\ 0 \\ 1}}{\stdbasis_4}
                     =\colvec[r]{0 \\ -1 \\ 0 \\ 1}_{\stdbasis_4} \)
      \end{exparts}  
    \end{answer}
  \item  Represent the vector with respect to each of the two bases.
    \begin{equation*}
      \vec{v}=\colvec{3  \\ -1}
      \quad
      B_1=\sequence{\colvec{1  \\ -1}, \colvec{1  \\ 1}},\;
      B_2=\sequence{\colvec{1 \\ 2}, \colvec{1 \\ 3}}
    \end{equation*}
    \begin{answer}
       Solving
      \begin{equation*}
        \colvec{3 \\ -1}=\colvec{1 \\ -1}\cdot c_1
                         +\colvec{1 \\ 1}\cdot c_2
      \end{equation*}
      gives $c_1=2$ and~$c_2=1$. 
      \begin{equation*}
        \rep{\colvec{3 \\ -1}}{B_1}
           =\colvec{2 \\ 1}_{B_1}
      \end{equation*}
      Similarly, solving
      \begin{equation*}
        \colvec{3 \\ -1}=\colvec{1 \\ 2}\cdot c_1
                         +\colvec{1 \\ 3}\cdot c_2
      \end{equation*}
      gives this. 
      \begin{equation*}
        \rep{\colvec{3 \\ -1}}{B_2}
           =\colvec{10 \\ -7}_{B_2}
      \end{equation*}     
    \end{answer}
  \item  
    Find a basis for \(  \polyspace_2 \), the space of all quadratic
    polynomials.
    Must any such basis contain a polynomial of each degree:~degree zero, 
    degree one, and degree two?
    \begin{answer}
      A natural basis is \( \sequence{1,x,x^2} \).
      There are bases for $\polyspace_2$ that do not contain any polynomials 
      of degree one or degree zero.
      One is \( \sequence{1+x+x^2,x+x^2,x^2} \).
      (Every basis has at least one polynomial of degree two, though.)
    \end{answer}
  \item
    Find a basis for the solution set of this system.
    \begin{equation*}
      \begin{linsys}{4}
         x_1  &-  &4x_2  &+  &3x_3  &-  &x_4  &=  &0  \\
        2x_1  &-  &8x_2  &+  &6x_3  &-  &2x_4 &=  &0  
      \end{linsys}
    \end{equation*}
    \begin{answer}
      The reduction
      \begin{equation*}
        \begin{amat}[r]{4}
          1  &-4  &3  &-1  &0  \\
          2  &-8  &6  &-2  &0
        \end{amat}
        \grstep{-2\rho_1+\rho_2}
        \begin{amat}[r]{4}
          1  &-4  &3  &-1  &0  \\
          0  &0   &0  &0   &0
        \end{amat}
      \end{equation*}
      gives that the only condition is that 
      $x_1=4x_2-3x_3+x_4$.
      The solution set is
      \begin{multline*}
        \set{\colvec{4x_2-3x_3+x_4 \\ x_2 \\ x_3 \\ x_4}
               \suchthat x_2,x_3,x_4\in\Re}                          \\
        =\set{x_2\colvec[r]{4 \\ 1 \\ 0 \\ 0}
             +x_3\colvec[r]{-3 \\ 0 \\ 1 \\ 0}
             +x_4\colvec[r]{1 \\ 0 \\ 0 \\ 1} \suchthat x_2,x_3,x_4\in\Re}
      \end{multline*}
      and so the obvious candidate for the basis is this.
      \begin{equation*}
       \sequence{\colvec[r]{4 \\ 1 \\ 0 \\ 0},
                   \colvec[r]{-3 \\ 0 \\ 1 \\ 0},
                   \colvec[r]{1 \\ 0 \\ 0 \\ 1}  }
      \end{equation*}  
      We've shown that this spans the space, and showing it is also linearly
      independent is routine.
    \end{answer}
  \recommended \item
    Find a basis for \( \matspace_{\nbyn{2}} \),
    the space of \( \nbyn{2} \) matrices.
    \begin{answer}
      There are many bases.
      This is a natural one.
      \begin{equation*}
        \sequence{
           \begin{mat}[r]
             1  &0  \\
             0  &0
           \end{mat},
           \begin{mat}[r]
             0  &1  \\
             0  &0
           \end{mat},
           \begin{mat}[r]
             0  &0  \\
             1  &0
           \end{mat},
           \begin{mat}[r]
             0  &0  \\
             0  &1
           \end{mat}  }
      \end{equation*} 
    \end{answer}
  \recommended \item 
      Find a basis for each.
      \begin{exparts}
        \partsitem The subspace $\set{a_2x^2+a_1x+a_0\suchthat a_2-2a_1=a_0}$ 
          of $\polyspace_2$
        \partsitem The space of three-wide row vectors whose first and second
          components add to zero
        \partsitem This subspace of the $\nbyn{2}$ matrices
          \begin{equation*}
            \set{\begin{mat}
                   a  &b  \\
                   0  &c  
                 \end{mat} \suchthat c-2b=0}
          \end{equation*}
      \end{exparts}
      \begin{answer}
        For each item, many answers are possible.
        \begin{exparts}
          \partsitem One way to proceed is to parametrize by
            expressing the $a_2$ as a combination of the other two
            $a_2=2a_1+a_0$.
            Then
            $a_2x^2+a_1x+a_0$ is $(2a_1+a_0)x^2+a_1x+a_0$ and
            \begin{multline*} 
              \set{(2a_1+a_0)x^2+a_1x+a_0\suchthat a_1,a_0\in\Re}          \\
              =\set{a_1\cdot(2x^2+x)+a_0\cdot(x^2+1)\suchthat a_1,a_0\in\Re}
            \end{multline*}
            suggests 
            $\sequence{2x^2+x,x^2+1}$.
            This only shows that 
            it spans, but checking that it is linearly independent
            is routine.
          \partsitem Parametrize $\set{\rowvec{a  &b  &c}\suchthat a+b=0}$ 
            to get $\set{\rowvec{-b &b &c}\suchthat b,c\in\Re}$, which suggests
            using the sequence $\sequence{\rowvec{-1 &1 &0},\rowvec{0 &0 &1}}$.
            We've shown that it spans, and checking that it is linearly 
            independent is easy.
          \partsitem Rewriting
            \begin{equation*}
              \set{\begin{mat}
                     a  &b  \\
                     0  &2b
                   \end{mat}\suchthat a,b\in\Re}
              =\set{a\cdot\begin{mat}[r]
                     1  &0  \\
                     0  &0 
                   \end{mat}
                  +b\cdot\begin{mat}[r]
                           0  &1  \\
                           0  &2
                          \end{mat} \suchthat a,b\in\Re}
            \end{equation*}
            suggests this for the basis.
            \begin{equation*}
              \sequence{\begin{mat}[r]
                          1  &0  \\
                          0  &0
                        \end{mat},
                        \begin{mat}[r]
                          0  &1  \\
                          0  &2
                        \end{mat}}
            \end{equation*}
         \end{exparts}  
      \end{answer}
  \item Find a basis for each space, and verify that it is a basis.
    \begin{exparts}
      \partsitem
        The subspace $M=\set{a+bx+cx^2+dx^3\suchthat a-2b+c-d=0}$ 
        of $\polyspace_3$.
      \partsitem  This subspace of $\matspace_{\nbyn{2}}$.
         \begin{equation*}
           W=\set{
             \begin{mat}
               a  &b  \\
               c  &d
             \end{mat}
             \suchthat a-c=0}
         \end{equation*}

    \end{exparts}
    \begin{answer}
      \begin{exparts}
        \partsitem
          Parametrize $a-2b+c-d=0$ as $a=2b-c+d$, $b=b$, $c=c$, and~$d=d$ to get
          this description of $M$ as the span of a set of three vectors.
          \begin{equation*}
            M=\set{
             (2+x)\cdot b+(-1+x^2)\cdot c+(1+x^3)\cdot d
             \suchthat b,c,d\in\Re
             }
          \end{equation*}
          To show that this three-vector set is a basis, what remains is to
          verify that it is linearly independent.
          \begin{equation*}
             0+0x+0x^2+0x^3=(2+x)\cdot c_1+(-1+x^2)\cdot c_2+(1+x^3)\cdot c_3
          \end{equation*}
          From the $x$~terms we see that $c_1=0$.
          From the $x^2$~terms we see that $c_2=0$.
          The $x^3$~terms give that $c_3=0$. 
       \partsitem
         First parametrize the description
         (note that the fact that $b$ and~$d$ are not mentioned in the 
         description
         of $W$ does not mean
         they are zero or absent, it means that they are unrestricted).
         \begin{equation*}
           W=\set{
             \begin{mat}
               0 &1 \\
               0 &0
             \end{mat}\cdot b
             +
             \begin{mat}
               1 &0 \\
               1 &0
             \end{mat}\cdot c
             +
             \begin{mat}
               0 &0 \\
               0 &1
             \end{mat}\cdot d
             \suchthat b,c,d\in\Re}
         \end{equation*}
         That gives $W$ as the span of a three element set.
         We will be done if we show that the set is linearly independent.
         \begin{equation*}
             \begin{mat}
               0 &0 \\
               0 &0
             \end{mat}
             =
             \begin{mat}
               0 &1 \\
               0 &0
             \end{mat}\cdot c_1
             +
             \begin{mat}
               1 &0 \\
               1 &0
             \end{mat}\cdot c_2
             +
             \begin{mat}
               0 &0 \\
               0 &1
             \end{mat}\cdot c_3   
         \end{equation*}
         Using the upper right entries we see that $c_1=0$.
         The upper left entries give that $c_2=0$, and the lower left entries 
         show that $c_3=0$. 
     \end{exparts}
   \end{answer}
  \item \label{exer:VerifBasesCosPlusSin} 
    Check \nearbyexample{ex:BasisForCosPlusSin}.
    \begin{answer}
      We will show that the second is a basis; the first is similar.
      We will show this straight from the definition of a basis, 
      because this example appears before 
      \nearbytheorem{th:BasisIffUniqueRepWRT}.

      To see that it is linearly independent,
      we set up
      \( c_1\cdot(\cos\theta-\sin\theta)+c_2\cdot(2\cos\theta+3\sin\theta)=
        0\cos\theta+0\sin\theta \).
      Taking \( \theta=0 \) and \( \theta=\pi/2 \) gives this system
      \begin{equation*}
        \begin{linsys}{2}
          c_1\cdot 1    &+  &c_2\cdot 2   &=  &0  \\
          c_1\cdot (-1) &+  &c_2\cdot 3   &=  &0  
        \end{linsys}
        \grstep{\rho_1+\rho_2}
        \begin{linsys}{2}
          c_1  &+  &2c_2   &=  &0  \\
               &+  &5c_2   &=  &0 
        \end{linsys}
      \end{equation*}
      which shows that $c_1=0$ and $c_2=0$.    

      The calculation for span is also easy; for any $x,y\in\Re$, 
      we have that 
      \( c_1\cdot(\cos\theta-\sin\theta)+c_2\cdot(2\cos\theta+3\sin\theta)=
        x\cos\theta+y\sin\theta \)
      gives that \( c_2=x/5+y/5 \) and that \( c_1=3x/5-2y/5 \),
      and so the span is the entire space.  
    \end{answer}
  \recommended \item 
    Find the span of each set and then find a basis for that span.
    \begin{exparts*}
       \partsitem $\set{1+x,1+2x}$ in $\polyspace_2$
       \partsitem $\set{2-2x,3+4x^2}$ in $\polyspace_2$
    \end{exparts*}
    \begin{answer}
      \begin{exparts}
         \partsitem Asking which $a_0+a_1x+a_2x^2$ can be expressed as
           $c_1\cdot (1+x)+c_2\cdot (1+2x)$ 
           gives rise to three linear equations,
           describing the coefficients of $x^2$, $x$, and the constants.
           \begin{equation*} 
             \begin{linsys}{2}
               c_1 &+ &c_2  &= &a_0 \\
               c_1 &+ &2c_2 &= &a_1 \\
                   &  &0    &= &a_2
             \end{linsys}
           \end{equation*}
           Gauss's Method with back-substitution shows, 
           provided that $a_2=0$, that $c_2=-a_0+a_1$
           and $c_1=2a_0-a_1$.
           Thus, with $a_2=0$, we can compute appropriate
           $c_1$ and $c_2$ for any $a_0$ and $a_1$.
           So the span is the entire set of linear polynomials
           $\set{a_0+a_1x\suchthat a_0,a_1\in\Re}$.
           Parametrizing that set 
           $\set{a_0\cdot 1+a_1\cdot x\suchthat a_0,a_1\in\Re}$
           suggests a basis $\sequence{1,x}$ 
           (we've shown that it spans; checking linear independence is easy). 
        \partsitem With
          \begin{equation*}
            a_0+a_1x+a_2x^2
            =c_1\cdot(2-2x)+c_2\cdot(3+4x^2)
            =(2c_1+3c_2)+(-2c_1)x+(4c_2)x^2
          \end{equation*}
          we get this system.
           \begin{equation*} 
             \begin{linsys}{2}
               2c_1  &+ &3c_2  &= &a_0 \\
               -2c_1 &  &      &= &a_1 \\
                     &  &4c_2  &= &a_2
              \end{linsys}
             \grstep{\rho_1+\rho_2}
             \repeatedgrstep{(-4/3)\rho_2+\rho_3}
             \begin{linsys}{2}
               2c_1  &+ &3c_2  &= &a_0\hfill\hbox{} \\
                     &  &3c_2  &= &a_0+a_1\hfill\hbox{} \\
                     &  &0     &= &(-4/3)a_0-(4/3)a_1+a_2
              \end{linsys}
           \end{equation*}
           Thus, the only quadratic polynomials $a_0+a_1x+a_2x^2$ with 
           associated $c$'s are the ones such that 
           $0=(-4/3)a_0-(4/3)a_1+a_2$.
           Hence the span is this.
           \begin{equation*}
             \set{(-a_1+(3/4)a_2)+a_1x+a_2x^2\suchthat a_1,a_2\in\Re}
           \end{equation*}
           Parametrizing gives
           $\set{a_1\cdot (-1+x)+a_2\cdot ((3/4)+x^2)\suchthat a_1,a_2\in\Re}$,
           which suggests $\sequence{-1+x,(3/4)+x^2}$
           (checking that it is linearly independent is routine).
      \end{exparts} 
    \end{answer}
  \recommended \item 
    Find a basis for each of these subspaces of the space 
    $\polyspace_3$ of cubic polynomials. 
    \begin{exparts}
      \partsitem  The subspace of cubic polynomials $p(x)$ 
        such that $p(7)=0$
      \partsitem  The subspace of polynomials $p(x)$ such
        that $p(7)=0$ and $p(5)=0$
      \partsitem  The subspace of polynomials $p(x)$ such
        that $p(7)=0$, $p(5)=0$, and~$p(3)=0$
      \partsitem  The space of polynomials $p(x)$ such
        that $p(7)=0$, $p(5)=0$, $p(3)=0$, and~$p(1)=0$
    \end{exparts}
    \begin{answer}
       \begin{exparts}
        \partsitem The subspace is this.
          \begin{equation*}
             \set{a_0+a_1x+a_2x^2+a_3x^3 \suchthat a_0+7a_1+49a_2+343a_3=0 }
          \end{equation*}
          Rewriting $a_0=-7a_1-49a_2-343a_3$ gives this.
          \begin{equation*}
             \set{(-7a_1-49a_2-343a_3)+a_1x+a_2x^2+a_3x^3 
               \suchthat a_1,a_2,a_3\in\Re }
          \end{equation*}
          On breaking out the
          parameters, this suggests \( \sequence{-7+x,-49+x^2,-343+x^3} \)  
          for the basis (it is easily verified).
        \partsitem The given subspace is the collection of cubics
          $p(x)=a_0+a_1x+a_2x^2+a_3x^3$ such that $a_0+7a_1+49a_2+343a_3=0$
          and $a_0+5a_1+25a_2+125a_3=0$.     
          Gauss's Method 
          \begin{equation*}
            \begin{linsys}{4}
              a_0  &+  &7a_1  &+  &49a_2  &+  &343a_3  &=  &0  \\
              a_0  &+  &5a_1  &+  &25a_2  &+  &125a_3  &=  &0    
            \end{linsys}
            \grstep{-\rho_1+\rho_2}
            \begin{linsys}{4}
              a_0  &+  &7a_1  &+  &49a_2  &+  &343a_3  &=  &0  \\
                   &   &-2a_1 &-  &24a_2  &-  &218a_3  &=  &0    
            \end{linsys}
          \end{equation*}
          gives that $a_1=-12a_2-109a_3$ and that $a_0=35a_2+420a_3$.
          Rewriting $(35a_2+420a_3)+(-12a_2-109a_3)x+a_2x^2+a_3x^3$
          as $a_2\cdot(35-12x+x^2)+a_3\cdot(420-109x+x^3)$
          suggests this for a basis  $\sequence{35-12x+x^2,420-109x+x^3}$.
          The above shows that it spans the space.
          Checking it is linearly independent is routine.
          (\textit{Comment.} 
          A worthwhile check is to verify that both polynomials in the
          basis have both seven and five as roots.)
        \partsitem Here there are three conditions on the cubics,
          that $a_0+7a_1+49a_2+343a_3=0$, that $a_0+5a_1+25a_2+125a_3=0$,
          and that $a_0+3a_1+9a_2+27a_3=0$.
          Gauss's Method 
          \begin{equation*}
            \begin{linsys}{4}
              a_0  &+  &7a_1  &+  &49a_2  &+  &343a_3  &=  &0  \\
              a_0  &+  &5a_1  &+  &25a_2  &+  &125a_3  &=  &0  \\    
              a_0  &+  &3a_1  &+  &9a_2   &+  &27a_3   &=  &0    
            \end{linsys}
            \grstep[-\rho_1+\rho_3]{-\rho_1+\rho_2}
            \repeatedgrstep{-2\rho_2+\rho_3}
            \begin{linsys}{4}
              a_0  &+  &7a_1  &+  &49a_2  &+  &343a_3  &=  &0  \\
                   &   &-2a_1 &-  &24a_2  &-  &218a_3  &=  &0  \\    
                   &   &      &   &8a_2   &+  &120a_3  &=  &0    
            \end{linsys}
          \end{equation*}
          yields the single free variable $a_3$, with 
          $a_2=-15a_3$, $a_1=71a_3$, and $a_0=-105a_3$.
          The parametrization is this.
          \begin{multline*} 
            \set{(-105a_3)+(71a_3)x+(-15a_3)x^2+(a_3)x^3\suchthat a_3\in\Re} \\
            =
            \set{a_3\cdot(-105+71x-15x^2+x^3)\suchthat a_3\in\Re}
          \end{multline*}
          Therefore, a natural candidate for the basis is 
          $\sequence{-105+71x-15x^2+x^3}$.
          It spans the space by the work above.
          It is clearly linearly independent because it is a one-element
          set (with that single element not the zero object of the space).
          Thus, any cubic through the three points $(7,0)$, $(5,0)$, and
          $(3,0)$ is a multiple of this one.
          (\textit{Comment.}
          As in the prior question, 
          a worthwhile check is to verify that plugging seven, five, and
          three into this polynomial yields zero each time.)
        \partsitem This is the trivial subspace of $\polyspace_3$.
          Thus, the  basis is empty $\sequence{}$.
      \end{exparts}
      \noindent\textit{Remark.}
      Alternatively, we could have derived the polynomial in the third item
      by multiplying out $(x-7)(x-5)(x-3)$.
     \end{answer}
  \item 
     We've seen that the result of reordering a basis can be another basis.
     Must it be? 
     \begin{answer}
       Yes.
       Linear independence and span are unchanged by reordering.
     \end{answer}
  \item  
    Can a basis contain a zero vector?
    \begin{answer}
      No linearly independent set contains a zero vector.  
    \end{answer}
  \recommended \item
    Let \( \sequence{\vec{\beta}_1,\vec{\beta}_2,\vec{\beta}_3} \)
    be a basis for a vector space.
    \begin{exparts}
      \partsitem Show that 
        \( \sequence{c_1\vec{\beta}_1,c_2\vec{\beta}_2,c_3\vec{\beta}_3} \)
        is a basis when \( c_1, c_2, c_3\neq 0 \).
        What happens when at least one \( c_i \) is $0$?
      \partsitem Prove that
        \( \sequence{\vec{\alpha}_1,\vec{\alpha}_2,\vec{\alpha}_3} \)
        is a basis where \( \vec{\alpha}_i=\vec{\beta}_1+\vec{\beta}_i \).
    \end{exparts}
    \begin{answer}
      \begin{exparts}
         \partsitem To show that it is linearly independent, note that if
           \( d_1(c_1\vec{\beta}_1)+d_2(c_2\vec{\beta}_2)+d_3(c_3\vec{\beta}_3)
               =\zero \)
           then
           \( (d_1c_1)\vec{\beta}_1+(d_2c_2)\vec{\beta}_2+(d_3c_3)\vec{\beta}_3
               =\zero \), which in turn 
           implies that each \( d_ic_i \) is zero.
           But with \( c_i\neq 0 \) that means that each \( d_i \) is zero.
           Showing that it spans the space is much the same; 
           because $\sequence{\vec{\beta}_1,\vec{\beta}_2,\vec{\beta}_3}$
           is a basis, and so spans the space, we can for any $\vec{v}$ write  
           \( \vec{v}=d_1\vec{\beta}_1+d_2\vec{\beta}_2+d_3\vec{\beta}_3 \),
           and then
           \( \vec{v}=(d_1/c_1)(c_1\vec{\beta}_1)
               +(d_2/c_2)(c_2\vec{\beta}_2)+(d_3/c_3)(c_3\vec{\beta}_3) \).

           If any of the scalars are zero then the result is not a basis,
           because it is not linearly independent.
         \partsitem Showing that 
           $\sequence{2\vec{\beta}_1,\vec{\beta}_1+\vec{\beta}_2,
             \vec{\beta}_1+\vec{\beta}_3}$ is linearly independent is easy. 
           To show that it spans the space, assume that
           \( \vec{v}=d_1\vec{\beta}_1+d_2\vec{\beta}_2+d_3\vec{\beta}_3 \).
           Then, we can represent the same \( \vec{v} \) with respect to
           \( \sequence{2\vec{\beta}_1,\vec{\beta}_1+\vec{\beta}_2,
                         \vec{\beta}_1+\vec{\beta}_3} \)
           in this way
           $\vec{v}=(1/2)(d_1-d_2-d_3)(2\vec{\beta}_1)
           +d_2(\vec{\beta}_1+\vec{\beta}_2)+d_3(\vec{\beta}_1+\vec{\beta}_3)$.
      \end{exparts}   
    \end{answer}
  \item 
    Find one vector $\vec{v}$ that will make each into a basis
    for the space.
    \begin{exparts*}
      \partsitem $\sequence{\colvec[r]{1 \\ 1},\vec{v}}$ in $\Re^2$
      \partsitem $\sequence{\colvec[r]{1 \\ 1 \\ 0},
                            \colvec[r]{0 \\ 1 \\ 0},\vec{v}}$ in $\Re^3$
      \partsitem $\sequence{x,1+x^2,\vec{v}}$ in $\polyspace_2$
    \end{exparts*} 
    \begin{answer}
      Each forms a linearly independent set if we omit $\vec{v}$.
      To preserve linear independence, we must expand the span of each.
      That is, we must determine the span of each (leaving $\vec{v}$ out),
      and then pick a $\vec{v}$ lying outside of that span.
      Then to finish, we must check that the result spans the entire given
      space.
      Those checks are routine.
      \begin{exparts}
        \partsitem Any vector that is not a multiple of the given one, 
          that is, any vector that is not on the line $y=x$ will do here.
          One is $\vec{v}=\vec{e}_1$.
        \partsitem By inspection, we notice that the vector $\vec{e}_3$ is
          not in the span of the set of the two given vectors.
          The check that the resulting set is a basis for $\Re^3$ is 
          routine.
        \partsitem For any member of the span 
          $\set{c_1\cdot(x)+c_2\cdot(1+x^2)\suchthat c_1,c_2\in\Re}$,
          the coefficient of $x^2$ equals the constant term.
          So we expand the span if we add a quadratic without this property,
          say, $\vec{v}=1-x^2$.
          The check that the result is a basis for $\polyspace_2$ is easy.  
      \end{exparts}
    \end{answer}
  \recommended \item Consider \( 2+4x^2, 1+3x^2, 1+5x^2 \in\polyspace_2\).
    \begin{exparts}
      \partsitem Find a linear relationship among the three.
      \partsitem Represent them with respect to 
        $B=\sequence{1-x,1+x,x^2}$.
      \partsitem Check that the same linear relationship holds among
        the representations, as in 
        \nearbylemma{lem:LinearRelRepresented}.
    \end{exparts}
    \begin{answer}
    \begin{exparts}
      \partsitem $1\cdot(2+4x^2)-3\cdot(1+3x^2)+1\cdot(1+5x^2)=0+0x+0x^2$
      \partsitem Simple calculation gives this.
        \begin{equation*}
          \rep{2+4x^2}{B}=\colvec{1 \\ 1 \\ 4}
          \quad
          \rep{1+3x^2}{B}=\colvec{1/2 \\ 1/2 \\ 3}
          \quad
          \rep{1+5x^2}{B}=\colvec{1/2 \\ 1/2 \\ 5}
        \end{equation*}
      \partsitem The check is straightforward.
        \begin{equation*}
          1\cdot\colvec{1 \\ 1 \\ 4}
          -3\cdot\colvec{1/2 \\ 1/2 \\ 3}
          +1\cdot\colvec{1/2 \\ 1/2 \\ 5}
          =\colvec{0 \\ 0 \\ 0}
        \end{equation*}
    \end{exparts}
    \end{answer}
  \recommended \item  
    Where
    \( \sequence{\vec{\beta}_1,\dots,\vec{\beta}_n } \)
    is a basis, show that in this equation
    \begin{equation*}
       c_1\vec{\beta}_1+\dots+c_k\vec{\beta}_k
       =
       c_{k+1}\vec{\beta}_{k+1}+\dots+c_n\vec{\beta}_n
    \end{equation*}
    each of the \( c_i \)'s is zero.
    Generalize.
    \begin{answer}
      To show that each scalar is zero, simply subtract
      \( c_1\vec{\beta}_1+\dots+c_k\vec{\beta}_k
          -c_{k+1}\vec{\beta}_{k+1}-\dots-c_n\vec{\beta}_n=\zero \).
      The obvious generalization is that in any equation involving only the
      \( \vec{\beta} \)'s, and in which each \( \vec{\beta} \) appears only
      once, each scalar is zero.
      For instance, an equation with a combination of 
      the even-indexed basis vectors
      (i.e., $\vec{\beta}_2$, $\vec{\beta}_4$, etc.) on the right and the
      odd-indexed basis vectors on the left also gives the conclusion that
      all of the coefficients are zero. 
    \end{answer}
  \item  
    A basis contains some of the vectors from a vector space; can it 
    contain them all?
    \begin{answer}
      No; no linearly independent set contains the zero vector.
    \end{answer}
  \item  
    \nearbytheorem{th:BasisIffUniqueRepWRT} shows that, with respect to a
    basis, every linear combination is unique.
    If a subset is not a basis, can linear combinations be not unique?
    If so, must they be?
    \begin{answer}
      Here is a subset of $\Re^2$ that is not a basis, and two different
      linear combinations of its elements that sum to the same vector.
      \begin{equation*}
        \set{\colvec[r]{1 \\ 2},\colvec[r]{2 \\ 4}}
        \qquad
        2\cdot\colvec[r]{1 \\ 2}+0\cdot\colvec[r]{2 \\ 4}
        =0\cdot\colvec[r]{1 \\ 2}+1\cdot\colvec[r]{2 \\ 4}
      \end{equation*}
      Thus, when a subset is not a basis, it can be the case that its
      linear combinations are not unique.

      But just because a subset is not a basis does not imply that its
      combinations must be not unique.
      For instance, this set 
      \begin{equation*}
        \set{\colvec[r]{1 \\ 2}}
      \end{equation*}
      does have the property that
      \begin{equation*}
        c_1\cdot\colvec[r]{1 \\ 2}
        =
        c_2\cdot\colvec[r]{1 \\ 2}
      \end{equation*}
      implies that $c_1=c_2$.
      The idea here is that this subset fails to be a basis because it fails
      to span the space; the proof of the theorem establishes that
      linear combinations are unique if and only if the subset is linearly
      independent.
     \end{answer}
  \item
    A square matrix is \definend{symmetric}\index{matrix!symmetric}%
    \index{symmetric matrix} 
    if for all indices \( i \) and
    \( j \), entry \( i,j \) equals entry \( j,i \).
    \begin{exparts}
      \partsitem Find a basis for the vector space of
        symmetric \( \nbyn{2} \) matrices.
      \partsitem Find a basis for the space of symmetric \( \nbyn{3} \)
        matrices.
      \partsitem Find a basis for the space of symmetric \( \nbyn{n} \)
        matrices.
    \end{exparts}
    \begin{answer}
      \begin{exparts}
        \partsitem Describing the vector space as
          \begin{equation*}
             \set{\begin{mat}
                     a  &b  \\
                     b  &c
                  \end{mat}  \suchthat a,b,c\in\Re}
          \end{equation*}
          suggests this for a basis.
          \begin{equation*}
            \sequence{
              \begin{mat}[r]
                1  &0  \\
                0  &0
              \end{mat},
              \begin{mat}[r]
                0  &0  \\
                0  &1
              \end{mat},
              \begin{mat}[r]
                0  &1  \\
                1  &0
              \end{mat}  }
          \end{equation*}
          Verification is easy.
        \partsitem This is one possible basis.
          \begin{equation*}
            \sequence{
              \begin{mat}[r]
                1  &0  &0  \\
                0  &0  &0  \\
                0  &0  &0
              \end{mat},
              \begin{mat}[r]
                0  &0  &0  \\
                0  &1  &0  \\
                0  &0  &0
              \end{mat},
              \begin{mat}[r]
                0  &0  &0  \\
                0  &0  &0  \\
                0  &0  &1
              \end{mat},
              \begin{mat}[r]
                0  &1  &0  \\
                1  &0  &0  \\
                0  &0  &0
              \end{mat},
              \begin{mat}[r]
                0  &0  &1  \\
                0  &0  &0  \\
                1  &0  &0
              \end{mat},
              \begin{mat}[r]
                0  &0  &0  \\
                0  &0  &1  \\
                0  &1  &0
              \end{mat}  }
          \end{equation*}
         \partsitem As in the prior two questions, we can form a basis from two
           kinds of  matrices.
           First are the matrices with a single one on the diagonal and all
           other entries zero (there are \( n \) of those matrices).
           Second are the matrices with two opposed off-diagonal entries
           are ones and all other entries are zeros.
           (That is, all entries in $M$ are zero except that 
           $m_{i,j}$ and $m_{j,i}$ are one.) 
      \end{exparts}  
    \end{answer}
  \item  
     We can show that every basis for $\Re^3$ contains the same
     number of vectors.
     \begin{exparts}
       \partsitem Show that no linearly independent subset of $\Re^3$
          contains more than three vectors.
       \partsitem Show that 
          no spanning subset of $\Re^3$ contains fewer than three vectors.
          \textit{Hint:} 
          recall how to calculate the span of a set and show that
          this method 
          cannot yield all of $\Re^3$ when we apply it to fewer than
          three vectors.
     \end{exparts}
     \begin{answer}
       \begin{exparts}
         \partsitem Any four vectors from $\Re^3$ are linearly related because
           the vector equation
           \begin{equation*}
             c_1\colvec{x_1 \\ y_1 \\ z_1}
             +c_2\colvec{x_2 \\ y_2 \\ z_2}
             +c_3\colvec{x_3 \\ y_3 \\ z_3}
             +c_4\colvec{x_4 \\ y_4 \\ z_4}
             =\colvec[r]{0 \\ 0 \\ 0}
           \end{equation*}
           gives rise to a linear system
           \begin{equation*}
             \begin{linsys}{4}
                x_1c_1  &+  &x_2c_2  &+  &x_3c_3  &+  &x_4c_4  &=  &0 \\
                y_1c_1  &+  &y_2c_2  &+  &y_3c_3  &+  &y_4c_4  &=  &0 \\
                z_1c_1  &+  &z_2c_2  &+  &z_3c_3  &+  &z_4c_4  &=  &0
             \end{linsys}
           \end{equation*}
           that is homogeneous (and so has a solution) and has 
           four unknowns but only three equations,
           and therefore has nontrivial solutions.
           (Of course, this argument applies to any subset of $\Re^3$ 
           with four or more vectors.) 
         \partsitem We shall do just the two-vector case.
           Given $x_1$, \ldots, $z_2$, 
           \begin{equation*}
             S=\set{\colvec{x_1 \\ y_1 \\ z_1},
             \colvec{x_2 \\ y_2 \\ z_2}} 
           \end{equation*}
           to decide which vectors
           \begin{equation*}
             \colvec{x \\ y \\ z}
           \end{equation*}
           are in the span of $S$, set up 
           \begin{equation*}
             c_1\colvec{x_1 \\ y_1 \\ z_1}
             +c_2\colvec{x_2 \\ y_2 \\ z_2}
             =\colvec{x \\ y \\ z}
           \end{equation*}
           and row reduce the resulting system.
           \begin{equation*}
             \begin{linsys}{2}
                x_1c_1  &+  &x_2c_2   &=  &x \\
                y_1c_1  &+  &y_2c_2   &=  &y \\
                z_1c_1  &+  &z_2c_2   &=  &z
             \end{linsys}
           \end{equation*}
           There are two
           variables $c_1$ and $c_2$ but three equations, so 
           when Gauss's Method finishes, on the bottom
           row there will be some relationship of the form
           $0=m_1x+m_2y+m_3z$. 
           Hence, vectors in the span of the two-element set $S$
           must satisfy some restriction.
           Hence the span is not all of $\Re^3$.
       \end{exparts}  
    \end{answer}
  \item 
    One of the exercises in the Subspaces subsection shows that the set 
    \begin{equation*}
      \set{\colvec{x \\ y \\ z}\suchthat x+y+z=1}
    \end{equation*}
    is a vector space under these operations.
    \begin{equation*}
      \colvec{x_1 \\ y_1 \\ z_1}+\colvec{x_2 \\ y_2 \\ z_2}
      =\colvec{x_1+x_2-1 \\ y_1+y_2 \\ z_1+z_2}
      \qquad
      r\colvec{x \\ y \\ z}=\colvec{rx-r+1 \\ ry \\ rz}
    \end{equation*}
    Find a basis.
    \begin{answer}
      We have (using these oddball operations with care)
      \begin{multline*}
        \set{\colvec{1-y-z \\ y \\ z}\suchthat y,z\in\Re}
         =\set{\colvec{-y+1 \\ y \\ 0}
               +\colvec{-z+1 \\ 0 \\ z}
              \suchthat y,z\in\Re}                           \\
         =\set{y\cdot\colvec{0 \\ 1 \\ 0}
               +z\cdot\colvec{0 \\ 0 \\ 1}
              \suchthat y,z\in\Re}
      \end{multline*}
      and so a natural candidate for a basis is this.
      \begin{equation*}
        \sequence{\colvec[r]{0 \\ 1 \\ 0},\colvec[r]{0 \\ 0 \\ 1}}
      \end{equation*}
      To check linear independence we set up
      \begin{equation*}
         c_1\colvec[r]{0 \\ 1 \\ 0}+c_2\colvec[r]{0 \\ 0 \\ 1}
         =\colvec[r]{1 \\ 0 \\ 0}
      \end{equation*}
      (the vector on the right is the zero object in this space).
      That yields the linear system
      \begin{equation*}
        \begin{linsys}{3}
          (-c_1+1)  &+  &(-c_2+1)  &-  &1  &=  &1  \\
             c_1    &   &          &   &   &=  &0  \\
                    &   &c_2       &   &   &=  &0
        \end{linsys}
      \end{equation*}
      with only the solution $c_1=0$ and $c_2=0$.
      Checking the span is similar. 
   \end{answer}
\end{exercises}














\subsection{Dimension}
The previous subsection defines a basis of a vector space and
shows that a space can have many different bases.
% For example, following the definition of a basis, we saw three
% different bases for $\Re^2$.
So we cannot talk about ``the'' basis for a vector space.
True, some vector spaces have bases that strike us as more natural
than others, for instance, $\Re^2$'s basis $\stdbasis_2$ 
% or $\Re^3$'s basis $\stdbasis_3$
or $\polyspace_2$'s basis $\sequence{1,x,x^2}$.
But for 
the vector space $\set{a_2x^2+a_1x+a_0\suchthat 2a_2-a_0=a_1}$, 
no particular basis leaps out at us as the natural one.
We cannot, in general, associate with a space any single basis that
best describes it.

We can however find something about the bases that
is uniquely associated with the space.
This subsection shows that 
any two bases for a space have the same number of elements.
So with each space we can associate a number, 
the number of vectors in any of its bases.

Before we start, we first
limit our attention to spaces where at least one basis has only finitely
many members.

\begin{definition} \label{df:FiniteDimensional}
%<*df:FiniteDimensional>
A vector space is \definend{finite-dimensional\/}%
\index{vector space!finite dimensional}\index{finite-dimensional vector space}
if it has a basis with only finitely many vectors.
%</df:FiniteDimensional>
\end{definition}

\noindent One space that is not finite-dimensional is
the set of polynomials with real coefficients,
\nearbyexample{ex:PolysOfAllFiniteDegrees}. 
This is not spanned by any finite subset since that would contain
a polynomial of largest degree but this space has polynomials of all degrees.
Such spaces are interesting and important but we will focus in a
different direction.
From now on we will study only finite-dimensional vector spaces.
In the rest of this book we shall take `vector space' to mean 
`finite-dimensional vector space'.

% \begin{remark}
% One reason for sticking to finite-dimensional spaces is so that 
% the representation of a vector with respect to a 
% basis is a finitely-tall vector and we can easily write it.
% Another reason is that 
% the statement `any infinite-dimensional vector space has a basis'
% is equivalent to a statement called the Axiom of Choice
% \cite{Blass84} and so covering this would 
% move us far past this book's scope.
% (A discussion of the Axiom of Choice is in the
% Frequently Asked Questions list for \texttt{sci.math}, and
% another accessible one is \cite{Rucker}.)  
% \end{remark}

% \begin{theorem} \label{th:AllBasesSameSize}
% %<*th:AllBasesSameSizeCoord>
% In any finite-dimensional vector space, all 
% bases have the same number of elements.\index{basis}
% %</th:AllBasesSameSizeCoord>
% \end{theorem}

% \begin{proof}
% Fix a vector space with at least one finite basis.
% From among all of this space's bases, choose
% one \( B=\sequence{\vec{\beta}_1,\dots,\vec{\beta}_n} \) of minimal size.
% We will show that any other basis
% \( D=\smash{\sequence{\vec{\delta}_1,\vec{\delta}_2,\ldots}} \)
% has the same number of members, $n$.
% Because \( B \) has minimal size, \( D \) has no fewer than \( n \) vectors.
% We will argue that it cannot have more than \( n \) vectors.

% So suppose that
% $\vec{\delta}_1$, \ldots{}, $\vec{\delta}_{n+1}$ are distinct.
% The assumption that $D$ is linearly independent gives that
% the only relationship
% $a_1\vec{\delta}+1+\dots+a_{n+1}\vec{\delta}_{n+1}=\vec{0}$
% is the one where each $a_i$ is zero.
% We shall get a contradiction by finding that 
% there are nontrivial relationships among these~$n+1$ vectors.

% Represent them with respect to~$B$.
% \begin{equation*}
%   \rep{\vec{\delta}_1}{B}=\colvec{c_{1,1} \\ \vdots \\ c_{n,1}}
%   \quad\cdots\quad
%   \rep{\vec{\delta}_{n+1}}{B}=\colvec{c_{1,n+1} \\ \vdots \\ c_{n,n+1}}  
% \end{equation*}
% \nearbylemma{lem:LinearRelRepresented} says that a relationship
% holds among the vectors
% $a_1\vec{\delta}_1+\dots+a_{n+1}\vec{\delta}_{n+1}=\vec{0}$ 
% if and only if the same relationship holds among the representations.
% \begin{equation*}
%   a_1\colvec{c_{1,1} \\ \vdots \\ c_{n,1}}
%   +\cdots+
%   a_{n+1}\colvec{c_{1,n+1} \\ \vdots \\ c_{n,n+1}}
%   =\colvec{0 \\ \vdots \\ 0}  
% \end{equation*}
% Here the $c_{i,j}$ are fixed, since they are the coefficients in the 
% representations of the given vectors, while we are looking for the~$a_i$.
% Thus the above is a homogeneous linear system
% \begin{equation*}
%   \begin{linsys}{3}
%     c_{1,1}a_1 &+ &\cdots &+ &c_{1,n+1}a_{n+1} &= &0 \\ 
%               &  &        &  &               &\vdots \\ 
%     c_{n,1}a_1 &+ &\cdots &+ &c_{n,n+1}a_{n+1} &= &0 
%   \end{linsys}
% \end{equation*}
% with more unknowns, $n+1$, than equations.
% Such a system does not have only one solution, it has infinitely many solutions.
% \end{proof}


To prove the main theorem we shall use a technical result, the Exchange Lemma.
We first illustrate it with an example.

\begin{example}
Here is a basis for $\Re^3$ and a vector given as a linear combination
of members of that basis.
\begin{equation*}
  B=\sequence{\colvec[r]{1 \\ 0  \\ 0},
              \colvec[r]{1 \\ 1 \\ 0},
              \colvec[r]{0 \\ 0 \\ 2}}
  \qquad
  \colvec[r]{1 \\ 2 \\ 0}
  =(-1)\cdot\colvec[r]{1 \\ 0  \\ 0}
   +2\colvec[r]{1 \\ 1 \\ 0}
   +0\cdot\colvec[r]{0 \\ 0 \\ 2}
\end{equation*}
Two of the basis vectors have non-zero coefficients.
Pick one, for instance the first.
Replace it with the vector that we've expressed as the combination
\begin{equation*}
  \hat{B}=\sequence{\colvec[r]{1 \\ 2  \\ 0},
              \colvec[r]{1 \\ 1 \\ 0},
              \colvec[r]{0 \\ 0 \\ 2}}
\end{equation*}
and the result is another basis for \( \Re^3 \).
\end{example}

\begin{lemma}[Exchange Lemma] \label{lm:ExchangeLemma}
%<*lm:ExchangeLemma>
Assume that 
\( B=\sequence{\vec{\beta}_1,\dots,\vec{\beta}_n} \) is a basis for a
vector space, and that for the vector \( \vec{v} \)
the relationship \( \vec{v}=\lincombo{c}{\vec{\beta}} \)
has \( c_i\neq 0 \).
Then exchanging \( \vec{\beta}_i \) for \( \vec{v} \) yields another
basis for the space.
%</lm:ExchangeLemma>
\end{lemma}

\begin{proof}
%<*pf:ExchangeLemma0>
Call the outcome of the exchange 
\( \hat{B}=\sequence{\vec{\beta}_1,\dots,\vec{\beta}_{i-1},\vec{v},
                       \vec{\beta}_{i+1},\dots,\vec{\beta}_n}  \).

We first show that $\hat{B}$ is linearly independent.
Any relationship 
\( d_1\vec{\beta}_1+\dots+d_i\vec{v}+\dots+d_n\vec{\beta}_n=\zero \)
among the members of $\hat{B}$, after substitution for $\vec{v}$,
\begin{equation*}
  d_1\vec{\beta}_1+\dots
    +d_i\cdot(c_1\vec{\beta}_1+\dots+c_i\vec{\beta}_i+\dots+c_n\vec{\beta}_n)
    +\dots+d_n\vec{\beta}_n
  =\zero
\tag*{($*$)}\end{equation*}
gives a linear relationship among the members of $B$.
The basis $B$ is linearly independent so the coefficient $d_ic_i$ of 
$\vec{\beta}_i$ is zero.
Because we assumed that $c_i$ is nonzero, $d_i=0$.
Using this in equation~$(*)$ gives that all of the other $d$'s are also
zero.
Therefore $\hat{B}$ is linearly independent.
%</pf:ExchangeLemma0>

%<*pf:ExchangeLemma1>
We finish by showing that $\hat{B}$ has the same span as $B$.
Half of this argument, that $\spanof{\smash{\hat{B}}}\subseteq\spanof{B}$, 
is easy; we can write any member 
$d_1\vec{\beta}_1+\dots+d_i\vec{v}+\dots+d_n\vec{\beta}_n$
of $\spanof{\smash{\hat{B}}}$ as
$
  d_1\vec{\beta}_1+\dots
    +d_i\cdot(c_1\vec{\beta}_1+\dots+c_n\vec{\beta}_n)
    +\dots+d_n\vec{\beta}_n
$,
which is a linear combination of linear combinations of members of $B$, and
hence is in $\spanof{B}$.
For the $\spanof{B}\subseteq\spanof{\smash{\hat{B}}}$ half of the argument,
recall that if
$\vec{v}=c_1\vec{\beta}_1+\dots+c_n\vec{\beta}_n$ with $c_i\neq 0$
then we can rearrange the equation to
$\vec{\beta}_i=(-c_1/c_i)\vec{\beta}_1+\dots+(1/c_i)\vec{v}+\dots
   +(-c_n/c_i)\vec{\beta}_n$.
Now, consider any member
$d_1\vec{\beta}_1+\dots+d_i\vec{\beta}_i+\dots+d_n\vec{\beta}_n$
of $\spanof{B}$, substitute for $\vec{\beta}_i$ its expression as a linear
combination of the members of $\hat{B}$, and recognize, 
as in the first half of this argument, that the result is a linear
combination of linear combinations of members of $\hat{B}$, and hence is in 
$\spanof{\smash{\hat{B}}}$.
%</pf:ExchangeLemma1>
\end{proof}

\begin{theorem} \label{th:AllBasesSameSize}
%<*th:AllBasesSameSize>
In any finite-dimensional vector space, all 
bases have the same number of elements.
%</th:AllBasesSameSize>
\end{theorem}\index{basis}

\begin{proof}
%<*pf:AllBasesSameSize0>
Fix a vector space with at least one finite basis.
Choose, from among all of this space's bases,
one \( B=\sequence{\vec{\beta}_1,\dots,\vec{\beta}_n} \) of minimal size.
We will show that any other basis
\( D=\smash{\sequence{\vec{\delta}_1,\vec{\delta}_2,\ldots}} \)
also has the same number of members, $n$.
Because \( B \) has minimal size, \( D \) has no fewer than \( n \) vectors.
We will argue that it cannot have more than \( n \) vectors.
%</pf:AllBasesSameSize0>

%<*pf:AllBasesSameSize1>
The basis \( B \) spans the space and \( \vec{\delta}_1 \) is in the space,
so \( \vec{\delta}_1 \) is a nontrivial linear combination of elements of
\( B \).
By the Exchange Lemma, we can swap \( \vec{\delta}_1 \) for a
vector from \( B \), resulting in a basis \( B_1 \), where one element is
\( \vec{\delta}_1 \) and all of the \( n-1 \) other elements 
are \( \vec{\beta} \)'s.
%</pf:AllBasesSameSize1>

%<*pf:AllBasesSameSize2>
The prior paragraph forms the basis step for an induction argument.
The inductive step starts with a basis \( B_k \) (for \( 1\leq k<n \))
containing \( k \) members of \( D \) and \( n-k \) members of \( B \).
We know that \( D \) has at least \( n \) members so there is a
\( \vec{\delta}_{k+1} \).
Represent it as a linear combination of elements of \( B_k \).
The key point:~in that representation, at least one of the nonzero scalars
must be associated with a \( \vec{\beta}_i \) or else that
representation would be a
nontrivial linear relationship among elements of the linearly independent
set \( D \).
Exchange \( \vec{\delta}_{k+1} \) for \( \vec{\beta}_i \) to get a new basis
\( B_{k+1} \) with one \( \vec{\delta} \) more and one \( \vec{\beta} \)
fewer than the previous basis \( B_k \).
%</pf:AllBasesSameSize2>

%<*pf:AllBasesSameSize3>
Repeat that until no \( \vec{\beta} \)'s remain, so
that \( B_n \) contains  
$\vec{\delta}_1,\dots,\vec{\delta}_n$.
Now, \( D \) cannot have more than these \( n \) vectors because
any \( \vec{\delta}_{n+1} \) that remains would be in the span of
\( B_n \) (since it is a basis) 
and hence would be a linear combination of the other $\vec{\delta}$'s, 
contradicting that $D$ is linearly independent.
%</pf:AllBasesSameSize3>
\end{proof}

\begin{definition}  \label{df:Dimension}
%<*df:Dimension>
The \definend{dimension}\index{dimension}\index{vector space!dimension}
of a vector space is the number of vectors in any of its bases.
%</df:Dimension>
\end{definition}

\begin{example}
Any basis for \( \Re^n \) has \( n \) vectors since the standard basis 
\( \stdbasis_n \) has \( n \) vectors.
Thus, this definition of `dimension' generalizes the most familiar use of 
term, that $\Re^n$ is $n$-dimensional.
\end{example}

\begin{example}
The space \( \polyspace_n \) of polynomials of degree at most $n$ 
has dimension \( n+1 \).
We can show this by exhibiting any basis\Dash $\sequence{1,x,\dots,x^n}$ 
comes to mind\Dash and counting its members.
\end{example}

\begin{example}
The space of functions 
$\set{a\cdot\cos\theta+b\cdot\sin\theta\suchthat a,b\in\Re}$  
of the real variable $\theta$ has dimension~$2$ since this space has the 
basis $\sequence{\cos\theta,\sin\theta}$.
\end{example}

\begin{example}
A trivial space is zero-dimensional since its basis is empty.
\end{example}

Again, although we sometimes say `finite-dimensional' for emphasis, from now on
we take all vector spaces to be finite-dimensional.
So in the next result the word `space' 
means `finite-dimensional vector space'.

\begin{corollary} \label{cor:NoLiSetGreatDim}
%<*co:NoLiSetGreatDim>
No linearly independent set can have a size greater than the dimension of the
enclosing space.
%</co:NoLiSetGreatDim>
\end{corollary}

\begin{proof}
%<*pf:NoLiSetGreatDim>
The proof of \nearbytheorem{th:AllBasesSameSize} 
never uses that \( D \) spans the space,
only that it is linearly independent.
%</pf:NoLiSetGreatDim>
\end{proof}


\begin{example} \label{ex:RefSubSpDiagram}
Recall the diagram from Example I.\ref{ex:SubspRThree} showing 
the subspaces of \( \Re^3 \).
Each subspace is described with a minimal spanning set, a basis.
The whole space has a basis with three members, 
the plane subspaces have bases with two members,
the line subspaces have bases with one member,
and the trivial subspace has a basis with zero members.

In that section we could not show that these are 
\( \Re^3 \)'s only subspaces.
We can show it now.
The prior corollary proves that 
There are no, say, five-dimensional subspaces of three-space.
Further, by \nearbydefinition{df:Dimension} the dimension of every space
is a whole number so there are no subspaces of $\Re^3$
that are somehow $1.5$-dimensional, between lines and planes.  
Thus the list of subspaces that we gave is exhaustive;
the only subspaces of \( \Re^3 \) are either three-\hbox{},
two-\hbox{}, one-\hbox{}, or zero-dimensional.
\end{example}


\begin{corollary} \label{cor:LIExpBas}
%<*co:LIExpBas>
Any linearly independent set can be expanded to make a basis.
%</co:LIExpBas>
\end{corollary}

\begin{proof}
%<*pf:LIExpBas>
If a linearly independent set 
is not already a basis then it must not span the space.
Adding to the set a vector that is not in the span 
will preserve linear independence by 
Lemma II.\ref{lm:AddVecLiSetIsLiIffVecNotInSpan}.
Keep adding until the resulting set does span the space, 
which the prior corollary
shows will happen after only a finite number of steps.
%</pf:LIExpBas>
\end{proof}

\begin{corollary}  \label{co:SpanningSetShrinksToABasis}
%<*co:SpanningSetShrinksToABasis>
Any spanning set can be shrunk to a basis.
%</co:SpanningSetShrinksToABasis>
\end{corollary}

\begin{proof}
%<*pf:SpanningSetShrinksToABasis>
Call the spanning set \( S \).
If \( S \) is empty then it is already a basis (the space must be a trivial
space).
If \( S=\set{\zero} \) then it can be shrunk to the empty basis,
thereby making it linearly independent, without changing its span.

Otherwise, $S$ contains a vector $\vec{s}_1$ with $\vec{s}_1\neq\zero$
and we can form a basis \( B_1=\sequence{\vec{s}_1} \).
If \( \spanof{B_1}=\spanof{S} \) then we are done.
If not then there is a \( \vec{s}_2\in\spanof{S} \) such that
\( \vec{s}_2\not\in\spanof{B_1} \).
Let \( B_2=\sequence{\vec{s}_1,\vec{s_2}} \);
by Lemma~II.\ref{lm:AddVecLiSetIsLiIffVecNotInSpan} this is linearly independent
so if \( \spanof{B_2}=\spanof{S} \) then we are done.

We can repeat this process until the spans are equal,
which must happen in at most finitely many steps.
%</pf:SpanningSetShrinksToABasis>
\end{proof}


\begin{corollary} \label{cor:NVecsRNSpanIffLI}
%<*co:NVecsRNSpanIffLI>
In an \( n \)-dimensional space, a set composed of \( n \) vectors is linearly 
independent if and
only if it spans the space.
%</co:NVecsRNSpanIffLI>
\end{corollary}

\begin{proof}
%<*pf:NVecsRNSpanIffLI>
First we will
show that a subset with \( n \) vectors is linearly independent if and only
if it is a basis.
The `if' is trivially true\Dash bases are linearly independent.
`Only if' holds
because a linearly independent set can be expanded to a basis, but a
basis has \( n \) elements, so this 
expansion is actually the set that we began with.

To finish, we will show that any subset with \( n \) vectors spans the space 
if and only if it is a basis.
Again, `if' is trivial.
`Only if'
holds because any spanning set can be shrunk to a basis, but a basis has
\( n \) elements and so this shrunken set is just the one we started with.
%</pf:NVecsRNSpanIffLI>
\end{proof}

The main result of this subsection, that all of 
the bases in a finite-dimensional
vector space have the same number of elements, is the single most important
result in this book. 
As \nearbyexample{ex:RefSubSpDiagram}
shows, it describes what vector spaces and subspaces there can be.

One immediate consequence brings us 
back to when we considered the two things that could be meant
by the term `minimal spanning set'.
At that point we defined `minimal' as linearly independent
but we noted that 
another reasonable interpretation of the term is  
that a spanning set is `minimal' when it has the fewest
number of elements of any set with the same span.
Now that 
we have shown that all bases have the same number of elements, we know
that the two senses of `minimal' are equivalent.


\begin{exercises}
  \item[{\em Assume that all spaces are finite-dimensional unless otherwise
    stated.}]
  \recommended \item  
    Find a basis for, and the dimension of, \(  \polyspace_2 \).
    \begin{answer}
      One basis is \( \sequence{1,x,x^2} \), and so
      the dimension is three.
    \end{answer}
  \item 
    Find a basis for, and the dimension of, the solution set of
    this system.
    \begin{equation*}
      \begin{linsys}{4}
         x_1  &-  &4x_2  &+  &3x_3  &-  &x_4  &=  &0  \\
        2x_1  &-  &8x_2  &+  &6x_3  &-  &2x_4 &=  &0  
      \end{linsys}
    \end{equation*}
    \begin{answer}
      The solution set is
      \begin{equation*}
        \set{\colvec{4x_2-3x_3+x_4 \\ x_2 \\ x_3 \\ x_4}
               \suchthat x_2,x_3,x_4\in\Re}
      \end{equation*}
      so a natural basis is this
      \begin{equation*}
       \sequence{\colvec[r]{4 \\ 1 \\ 0 \\ 0},
                   \colvec[r]{-3 \\ 0 \\ 1 \\ 0},
                   \colvec[r]{1 \\ 0 \\ 0 \\ 1}  }
      \end{equation*}  
      (checking linear independence is easy).
      Thus the dimension is three.
    \end{answer}
  \recommended \item Find a basis for, and the dimension of, each space.
    \begin{exparts}
      \partsitem
        $
          \set{\colvec{x \\ y \\ z \\ w}\in\Re^4
               \suchthat x-w+z=0}
        $
      \partsitem the set of $\nbym{5}{5}$ matrices whose only nonzero entries 
        are on the diagonal (e.g., in entry $1,1$ and $2,2$, etc.)
      \partsitem 
        $\set{a_0+a_1x+a_2x^2+a_3x^3\suchthat 
            \text{$a_0+a_1=0$ and $a_2-2a_3=0$}}\subseteq\polyspace_3$ 
    \end{exparts}
    \begin{answer}
      \begin{exparts}
        \partsitem
          Parametrize to get this description of the space.
          \begin{equation*}
            \set{\colvec{w-z \\ y \\ z \\ w}
                 =\colvec{0 \\ 1 \\ 0 \\ 0}y+
                 \colvec{-1 \\ 0 \\ 1 \\ 0}z+
                 \colvec{1 \\ 0 \\ 0 \\ 1}w
                 \suchthat y,z,w\in\Re}
          \end{equation*}
          That gives the space as the span of the three-vector set.
          To show the three vector set makes a basis we check that it is 
          linearly independent.
          \begin{equation*}
             \colvec{0 \\ 0 \\ 0 \\ 0}
                 =\colvec{0 \\ 1 \\ 0 \\ 0}c_1+
                 \colvec{-1 \\ 0 \\ 1 \\ 0}c_2+
                 \colvec{1 \\ 0 \\ 0 \\ 1}c_3
          \end{equation*}
          The second components give that $c_1=0$, and the third and fourth 
          components give that $c_2=0$ and~$c_3=0$.
          So one basis is this.
          \begin{equation*}
              \sequence{
                 \colvec{0 \\ 1 \\ 0 \\ 0},
                 \colvec{-1 \\ 0 \\ 1 \\ 0},
                 \colvec{1 \\ 0 \\ 0 \\ 1} 
              } 
          \end{equation*}
          The dimension is the number of vectors in a basis: $3$.
        \partsitem
          The natural parametrization is this.
          \begin{multline*}
            \set{
              \begin{mat}
                a &0 &0 &0 &0 \\
                0 &b &0 &0 &0 \\
                0 &0 &c &0 &0 \\
                0 &0 &0 &d &0 \\
                0 &0 &0 &0 &e 
              \end{mat}
              \suchthat a,\ldots,e\in\Re}   \\
            \set{
              \begin{mat}
                1 &0 &0 &0 &0 \\
                0 &0 &0 &0 &0 \\
                0 &0 &0 &0 &0 \\
                0 &0 &0 &0 &0 \\
                0 &0 &0 &0 &0 
              \end{mat}\cdot a
              +\begin{mat}
                0 &0 &0 &0 &0 \\
                0 &1 &0 &0 &0 \\
                0 &0 &0 &0 &0 \\
                0 &0 &0 &0 &0 \\
                0 &0 &0 &0 &0 
              \end{mat}\cdot b
              +\cdots
              \suchthat a,\ldots,e\in\Re}
          \end{multline*}
          Checking that the five-element set is linearly independent is trivial.
          So this is a basis; the dimension is $5$.
          \begin{equation*}
            \sequence{
              \begin{mat}
                1 &0 &0 &0 &0 \\
                0 &0 &0 &0 &0 \\
                0 &0 &0 &0 &0 \\
                0 &0 &0 &0 &0 \\
                0 &0 &0 &0 &0 
              \end{mat},
              \begin{mat}
                0 &0 &0 &0 &0 \\
                0 &1 &0 &0 &0 \\
                0 &0 &0 &0 &0 \\
                0 &0 &0 &0 &0 \\
                0 &0 &0 &0 &0 
              \end{mat},
              \,\ldots\,,
              \begin{mat}
                0 &0 &0 &0 &0 \\
                0 &0 &0 &0 &0 \\
                0 &0 &0 &0 &0 \\
                0 &0 &0 &0 &0 \\
                0 &0 &0 &0 &1 
              \end{mat}
          }
          \end{equation*}
        \partsitem
          The restrictions form a two-equations, four-unknowns linear system.
          Parametrizing that system to express the leading variables in 
          terms of those that are
          free gives $a_0=-a_1$, $a_1=a_1$, $a_2=2a_3$, and~$a_3=a_3$.
          \begin{multline*}
            \set{-a_1+a_1x+2a_3x^2+a_3x^3\suchthat a_1,a_3\in\Re}         \\
            =\set{(-1+x)\cdot a_1+(2x^2+x^3)\cdot a_3\suchthat a_1,a_3\in\Re}
          \end{multline*}
          That description shows that the space is the span of the two-element 
          set $\set{-1+x,x^2+x^3}$.  
          We will be done if we show the set is linearly independent.
          This relationship
          \begin{equation*}
            0+0x+0x^2+0x^3=(-1+x)\cdot c_1+(2x^2+x^3)\cdot c_2
          \end{equation*}
          gives that $c_1=0$ from the constant terms, and $c_2=0$ from the 
          cubic terms.
          One basis for the space is $\sequence{-1+x,2x^2+x^3}$.
          This is a two-dimensional space.

      \end{exparts}
    \end{answer}
  \item
    Find a basis for, and the dimension of, \( \matspace_{\nbyn{2}} \),
    the vector space of \( \nbyn{2} \) matrices.
    \begin{answer}
      For this space
      \begin{multline*}
        \set{\begin{mat}
               a  &b  \\
               c  &d
             \end{mat} \suchthat a,b,c,d\in\Re}            \\
        =\set{a\cdot\begin{mat}[r]
             1  &0  \\
             0  &0
           \end{mat}
           +\dots+
           d\cdot\begin{mat}[r]
             0  &0  \\
             0  &1
           \end{mat} \suchthat a,b,c,d\in\Re}
      \end{multline*}
      this is a natural basis.
      \begin{equation*}
        \sequence{
           \begin{mat}[r]
             1  &0  \\
             0  &0
           \end{mat},
           \begin{mat}[r]
             0  &1  \\
             0  &0
           \end{mat},
           \begin{mat}[r]
             0  &0  \\
             1  &0
           \end{mat},
           \begin{mat}[r]
             0  &0  \\
             0  &1
           \end{mat}  }
      \end{equation*}
      The dimension is four. 
    \end{answer}
  \item 
    Find the dimension of the vector space of matrices
    \begin{equation*}
      \begin{mat}
        a  &b  \\
        c  &d
      \end{mat}
    \end{equation*}
    subject to each condition.
    \begin{exparts}
      \item $a, b, c, d\in\Re$ 
      \item $a-b+2c=0$ and~$d\in\Re$
      \item $a+b+c=0$, $a+b-c=0$, and~$d\in\Re$
    \end{exparts}
    \begin{answer}
     \begin{exparts}
      \partsitem As in the prior exercise, the space $\matspace_{\nbyn{2}}$ 
        of matrices without restriction has this basis
        \begin{equation*}
         \sequence{
           \begin{mat}[r]
             1  &0  \\
             0  &0
           \end{mat},
           \begin{mat}[r]
             0  &1  \\
             0  &0
           \end{mat},
           \begin{mat}[r]
             0  &0  \\
             1  &0
           \end{mat},
           \begin{mat}[r]
             0  &0  \\
             0  &1
           \end{mat}  }
        \end{equation*}
        and so the dimension is four.
      \partsitem For this space
        \begin{multline*}
         \set{\begin{mat}
               a  &b  \\
               c  &d
             \end{mat} \suchthat \text{$a=b-2c$ and $d\in\Re$}}     \\
         =\set{b\cdot\begin{mat}[r]
             1  &1  \\
             0  &0
           \end{mat}
           +c\cdot\begin{mat}[r]
             -2  &0  \\
              1  &0
           \end{mat}
           +d\cdot\begin{mat}[r]
             0  &0  \\
             0  &1
           \end{mat} \suchthat b,c,d\in\Re}
        \end{multline*}
        this is a natural basis.
        \begin{equation*}
          \sequence{
            \begin{mat}[r]
              1  &1  \\
              0  &0
            \end{mat},
            \begin{mat}[r]
              -2  &0  \\
               1  &0
            \end{mat},
            \begin{mat}[r]
              0  &0  \\
              0  &1
            \end{mat}  }
        \end{equation*}
        The dimension is three. 
      \partsitem Gauss's Method applied to the two-equation linear system gives
        that $c=0$ and that $a=-b$.
        Thus, we have this description
        \begin{equation*}
         \set{\begin{mat}
               -b  &b  \\
                0  &d
             \end{mat} \suchthat b,d\in\Re}
         =\set{b\cdot\begin{mat}[r]
             -1  &1  \\
             0  &0
           \end{mat}
           +d\cdot\begin{mat}[r]
             0  &0  \\
             0  &1
           \end{mat} \suchthat b,d\in\Re}
        \end{equation*}
        and so this is a natural basis.
        \begin{equation*}
          \sequence{
            \begin{mat}[r]
              -1  &1  \\
               0  &0
            \end{mat},
            \begin{mat}[r]
              0  &0  \\
              0  &1
            \end{mat}  }
        \end{equation*}
        The dimension is two. 
     \end{exparts} 
    \end{answer}
  \recommended \item
    Find the dimension of this subspace of $\Re^2$.
    \begin{equation*}
      S=\set{\colvec{a+b \\ a+c}\suchthat a,b,c\in\Re}
    \end{equation*}
    \begin{answer}
      We cannot simply count the parameters.
      That is, the answer is not~$3$.
      Instead, observe that we can express every member
      $\binom{x}{y}\in\Re^2$ in the form
      \begin{equation*}
        \colvec{x \\ y}=\colvec{a+b \\ a+c}  
      \end{equation*}
      with the choice of $a=0$, $b=x$, and $c=y$ 
      (other choices are possible).
      So $S$ is the set $S=\Re^2$.
      It has dimension~$2$.
    \end{answer}
  \recommended \item 
    Find the dimension of each.
    \begin{exparts}
      \partsitem  The space of cubic polynomials $p(x)$ such
        that $p(7)=0$
      \partsitem  The space of cubic polynomials $p(x)$ such
        that $p(7)=0$ and~$p(5)=0$
      \partsitem  The space of cubic polynomials $p(x)$ such
        that $p(7)=0$, $p(5)=0$, and~$p(3)=0$
      \partsitem  The space of cubic polynomials $p(x)$ such
        that $p(7)=0$, $p(5)=0$, $p(3)=0$, and~$p(1)=0$
    \end{exparts}
    \begin{answer}
      The bases for these spaces are developed 
      in the answer set of the prior subsection.
      \begin{exparts}
        \partsitem One basis is \( \sequence{-7+x,-49+x^2,-343+x^3} \).
          The dimension is three.
        \partsitem One basis is $\sequence{35-12x+x^2,420-109x+x^3}$ so the
          dimension is two.
        \partsitem A basis is $\set{-105+71x-15x^2+x^3}$.
          The dimension is one.
        \partsitem This is the trivial subspace of $\polyspace_3$ and so the 
          basis is empty.
          The dimension is zero.
      \end{exparts}  
    \end{answer}
  \item 
     What is the dimension of the span of the set 
     $\set{\cos^2\theta,\sin^2\theta,\cos2\theta,\sin2\theta}$?
     This span is a subspace of the space of all real-valued functions of 
     one real variable.
     \begin{answer}
       First recall that $\cos2\theta=\cos^2\theta-\sin^2\theta$, and so 
       deletion of $\cos2\theta$ from this set leaves the span unchanged.
       What's left, the set  
       $\set{\cos^2\theta,\sin^2\theta,\sin2\theta}$, is linearly independent
       (consider the relationship
       $c_1\cos^2\theta+c_2\sin^2\theta+c_3\sin2\theta=Z(\theta)$
       where $Z$ is the zero function, and then take
       $\theta=0$, $\theta=\pi/4$, and $\theta=\pi/2$ to conclude that
       each $c$ is zero).
       It is therefore a basis for its span.
       That shows that the span is a dimension three vector space.
     \end{answer}
  \item  
    Find the dimension of \( \C^{47} \), the vector space
    of $47$-tuples of complex numbers.
    \begin{answer}
      Here is a basis
      \begin{equation*}
        \sequence{(1+0i,0+0i,\dots,0+0i),\,
                  (0+1i,0+0i,\dots,0+0i),(0+0i,1+0i,\dots,0+0i),\ldots }
      \end{equation*}   
      and so the dimension is \( 2\cdot 47=94 \).
    \end{answer}
  \item  
    What is the dimension of the vector space $\matspace_{\nbym{3}{5}}$ 
    of \( \nbym{3}{5} \) matrices?
    \begin{answer}
      A basis is
      \begin{equation*}
        \sequence{
          \begin{mat}[r]
            1  &0  &0  &0  &0  \\
            0  &0  &0  &0  &0  \\
            0  &0  &0  &0  &0
          \end{mat},
          \begin{mat}[r]
            0  &1  &0  &0  &0  \\
            0  &0  &0  &0  &0  \\
            0  &0  &0  &0  &0
          \end{mat},
          \ldots,
          \begin{mat}[r]
            0  &0  &0  &0  &0  \\
            0  &0  &0  &0  &0  \\
            0  &0  &0  &0  &1
          \end{mat}  }
      \end{equation*}
      and thus the dimension is \( 3\cdot 5=15 \).  
    \end{answer}
  \recommended \item 
    Show that this is a basis for $\Re^4$.
    \begin{equation*}
      \sequence{\colvec[r]{1 \\ 0 \\ 0 \\ 0},
        \colvec[r]{1 \\ 1 \\ 0 \\ 0},
        \colvec[r]{1 \\ 1 \\ 1 \\ 0},
        \colvec[r]{1 \\ 1 \\ 1 \\ 1} }
    \end{equation*}
    (We can use the results of this subsection to simplify this job.)
    \begin{answer}
       In a four-dimensional space a set of four vectors is linearly
       independent if and only if it spans the space.
       The form of these vectors makes linear independence easy to show
       (look at the equation of fourth components, then at the equation of 
       third components, etc.).  
    \end{answer}
  \item Decide if each is a basis for $\polyspace_{2}$.
    \begin{exparts*}
      \partsitem $\set{1,x^2,x^2-x}$
      \partsitem $\set{x^2+x,x^2-x}$
      \partsitem $\set{2x^2+x+1,2x+1,2}$
      \partsitem $\set{3x^2,-1,3x,x^2-x}$
    \end{exparts*}
    \begin{answer}
      By the results of this section, because $\polyspace_2$ has dimension~$3$,
      to show that a linearly independent set is a basis we need only observe
      that it has three members. 
      To show a set is not a basis we need
      only observe that it does not have three members (in this case
      we don't have to worry about linear independence). 
      \begin{exparts}
        \partsitem  This is a basis; it is linearly independent by inspection
           (the first element has no quadratic or linear term, the second has
           a quadratic but no linear term, and the third has a linear term)
           and it has three elements.
        \partsitem This is not a basis as it has only two elements.
        \partsitem This three-element set is a basis.
        \partsitem This is not a basis as it has four elements.
      \end{exparts}
    \end{answer}
  \item 
     Refer to \nearbyexample{ex:RefSubSpDiagram}.
     \begin{exparts}
       \partsitem Sketch a similar subspace diagram for $\polyspace_2$.
       \partsitem Sketch one for $\matspace_{\nbyn{2}}$.
     \end{exparts}
     \begin{answer}
      \begin{exparts}
       \partsitem The diagram for $\polyspace_2$ has four levels.
          The top level has the only three-dimensional subspace, 
          $\polyspace_2$ itself.
          The next level contains the two-dimensional subspaces
          (\emph{not} just the linear polynomials; any two-dimensional
          subspace, like those polynomials of the form $ax^2+b$).
          Below that are the one-dimensional subspaces.
          Finally, of course, is the only zero-dimensional subspace,
          the trivial subspace.
        \partsitem For $\matspace_{\nbyn{2}}$, the diagram has five levels,
          including subspaces of dimension four through zero.
     \end{exparts}
    \end{answer}
  \recommended \item \label{exer:DimDomToR}
    Where \( S \) is a set, the functions
    \( \map{f}{S}{\Re} \) form a vector space under the natural operations:
    the sum $f+g$ is the function given by 
    \( f+g\,(s)=f(s)+g(s) \) and the scalar product is
    \( r\cdot f \, (s)=r\cdot f(s) \).
    What is the dimension of the space resulting for each domain?
    \begin{exparts*}
      \partsitem \( S=\set{1} \)
      \partsitem \( S=\set{1,2} \)
      \partsitem \( S=\set{1,\ldots,n} \)
    \end{exparts*}
    \begin{answer}
      \begin{exparts*}
        \partsitem One
        \partsitem Two
        \partsitem \( n \)
      \end{exparts*}  
    \end{answer}
  \item  
    (See \nearbyexercise{exer:DimDomToR}.)
    Prove that this is an infinite-dimensional space:
    the set of all functions \( \map{f}{\Re}{\Re} \) under the natural
    operations.
    \begin{answer}
      We need only produce an infinite linearly independent set.
      One is such sequence is \( \sequence{f_1,f_2,\ldots} \) where
      \( \map{f_i}{\Re}{\Re} \) is
      \begin{equation*}
         f_i(x)=\begin{cases}
                   1  &\text{if \( x=i \)}  \\
                   0  &\text{otherwise}
                \end{cases}
      \end{equation*}  
      the function that has value $1$ only at $x=i$.
    \end{answer}
  \item  
    (See \nearbyexercise{exer:DimDomToR}.)
    What is the dimension of the vector space of functions
    $\map{f}{S}{\Re}$, under the natural operations, where the
    domain $S$ is the empty set?
    \begin{answer}
      A function is a set of ordered pairs
      $(x,f(x))$.
      So there is only one function with an empty domain, namely the empty set.
      A vector space with only one element a trivial vector space 
      and has dimension zero.
    \end{answer}
  \item  
    Show that 
    any set of four vectors in \( \Re^2 \) is linearly dependent.
    \begin{answer}
      Apply \nearbycorollary{cor:NoLiSetGreatDim}.
    \end{answer}
  \item  
    Show that
    \( \sequence{\vec{\alpha}_1,\vec{\alpha}_2,\vec{\alpha}_3}\subset\Re^3 \)
    is a basis if and only if there is no plane through the origin containing
    all three vectors.
    \begin{answer}
      A plane has the
      form $\set{\vec{p}+t_1\vec{v}_1+t_2\vec{v}_2\suchthat t_1,t_2\in\Re}$.
      (The first chapter also calls this a `$2$-flat', and contains
      a discussion of why this is equivalent to the
      description often taken in Calculus as the set of points $(x,y,z)$
      subject to a condition of the form $ax+by+cz=d$).
      When the plane passes through the origin we can take the particular
      vector $\vec{p}$ to be $\zero$.
      Thus, in the language we have developed in this chapter, a plane through
      the origin is the span of a set of two vectors.

      Now for the statement.
      Asserting that the three are not coplanar is the same as asserting that
      no vector lies in the span of the other two\Dash no vector is a linear
      combination of the other two.
      That's simply an assertion that the three-element set is linearly
      independent.
      By \nearbycorollary{cor:NVecsRNSpanIffLI}, that's equivalent to an
      assertion that the set is a basis for $\Re^3$ (more precisely, 
      any sequence made from the set's elements is a basis). 
    \end{answer}
  \item 
    % \begin{exparts}
      % \partsitem Prove that any subspace of a finite dimensional space
      %   has a basis. 
        Prove that any subspace of a finite dimensional space is 
        finite dimensional.
    % \end{exparts}
    \begin{answer}
      Let the space $V$ be finite dimensional and let
      $S$ be a subspace of $V$.

      If $S$ is not finite dimensional then it has a linearly independent 
      set that is infinite (start with the empty set and iterate adding 
      vectors that are not linearly dependent on the set; this process
      can continue for infinitely many steps or else $S$ would be 
      finite dimensional).
      But any linearly independent subset of~$S$ is a linearly independent
      subset of~$V$, contradicting \nearbycorollary{cor:NoLiSetGreatDim}
      % \begin{exparts}
      %   \partsitem The empty set is a linearly independent subset of $S$.
      %     By \nearbycorollary{cor:LIExpBas}, it can be expanded to a basis
      %     for the vector space $S$. 
      %   \partsitem Any basis for the subspace $S$ is a linearly independent 
      %     set in the superspace $V$.
      %     Hence it can be expanded to a basis for the superspace, which is
      %     finite dimensional.
      %     Therefore it has only finitely many members.  
      % \end{exparts}
    \end{answer}
  \item  
    Where is the finiteness of \( B \) used in
    \nearbytheorem{th:AllBasesSameSize}?
    \begin{answer}
      It ensures that we exhaust the \( \vec{\beta} \)'s.
      That is, it justifies the first sentence of the last paragraph. 
    \end{answer}
  \item
    Prove that if \( U \) and \( W \) are both three-dimensional
    subspaces of \( \Re^5 \) then \( U\intersection W \) is non-trivial.
    Generalize.
    \begin{answer}
       Let \( B_U \) be a basis for \( U \) and let \( B_W \)
       be a basis for \( W \).
       Consider the concatenation of the two basis sequences.
       If there is a repeated element then the intersection
       \( U\intersection W \) is nontrivial.
       Otherwise, 
       the set \( B_U\union B_W \) is linearly dependent as it is a
       six member subset of the five-dimensional space \( \Re^5 \).
       In either case some member of \( B_W \) is in the span of 
       \( B_U \), and
       thus \( U\intersection W \) is more than just the trivial space
       \( \set{\zero\,} \).

       Generalization:
       if \( U,W \) are subspaces of a vector space of dimension \( n \) and
       if \( \dim(U)+\dim(W)>n \) then they have a nontrivial
       intersection.  
    \end{answer}
  \item 
    A basis for a space consists of elements of that space.
    So we are naturally led to how the property `is a basis' 
    interacts with operations $\subseteq$ and $\cap$ and $\cup$.
    (Of course, a basis is actually a sequence that it is ordered, 
    but there is a natural extension of these operations.)
    \begin{exparts}
      \partsitem Consider first how bases might be related by $\subseteq$.
        Assume that \( U,W \) are subspaces of some vector space and
        that \( U\subseteq W \).
        Can there exist bases \( B_U \) for \( U \) and \( B_W \) for \( W \)
        such that \( B_U\subseteq B_W \)?
        Must such bases exist?

        For any basis \( B_U \) for \( U \), must there be a basis \( B_W \)
        for \( W \) such that \( B_U\subseteq B_W \)?

        For any basis \( B_W \) for \( W \), must there be a basis \( B_U \)
        for \( U \) such that \( B_U\subseteq B_W \)?

        For any bases \( B_U, B_W \) for \( U \) and \( W \), must  \( B_U \)
        be a subset of \( B_W \)?
      \partsitem Is the $\cap$ of bases a basis?
        For what space?
      \partsitem Is the $\cup$ of bases a basis?
        For what space?
      \partsitem What about the complement operation?
     \end{exparts}
     (\textit{Hint.} 
     Test any conjectures against some subspaces of \( \Re^3 \).)
     \begin{answer}
       First, note that a set is a basis for some space if and only
       if it is linearly independent, because in that case 
       it is a basis for its own span.
       \begin{exparts}
         \partsitem The answer to the question in the second paragraph
            is ``yes''
           (implying ``yes'' answers for both questions in the first
           paragraph).
           If \( B_U \) is a basis for \( U \) then
           \( B_U \) is a linearly independent subset of \( W \).
           Apply \nearbycorollary{cor:LIExpBas} to expand it to a basis for 
           \( W \).
           That is the desired \( B_W \).

           The answer to the question in the third paragraph is ``no'', which
           implies a ``no'' answer to the question of the fourth paragraph.
           Here is an example of a basis for a superspace with no sub-basis
           forming a basis for a subspace: in \( W=\Re^2 \), consider the
           standard basis \( \stdbasis_2 \).
           No sub-basis of $\stdbasis_2$ forms a basis for the 
           subspace \( U \) 
           of $\Re^2$ that is the line \( y=x \).
         \partsitem It is a basis (for its span) because the
           intersection of linearly
           independent sets is linearly independent (the intersection is a
           subset of each of the linearly independent sets).

           It is not, however, a basis for the intersection of the spaces.
           For instance, these are bases for \( \Re^2 \):
           \begin{equation*}
             B_1=\sequence{\colvec[r]{1 \\ 0},\colvec[r]{0 \\ 1}}
             \quad\text{and}\quad
             B_2=\sequence[r]{\colvec{2 \\ 0},\colvec[r]{0 \\ 2}}
           \end{equation*}
           and \( \Re^2\intersection\Re^2=\Re^2 \), but
           \( B_1\intersection B_2 \) is empty.
           All we can say is that the $\cap$ of the bases is a basis
           for a subset of the intersection of the spaces.
         \partsitem The $\cup$ of bases need not be a basis: in \( \Re^2 \)
           \begin{equation*}
             B_1=\sequence{\colvec[r]{1 \\ 0},\colvec[r]{1 \\ 1}}
             \quad\text{and}\quad
             B_2=\sequence{\colvec[r]{1 \\ 0},\colvec[r]{0 \\ 2}}
           \end{equation*}
           \( B_1\union B_2 \) is not linearly independent.
           A necessary and sufficient condition for a $\cup$ of two bases
           to be a basis 
           \begin{equation*}
             B_1\union B_2 \text{ is linearly independent }
             \quad\iff\quad
             \spanof{B_1\intersection B_2}=\spanof{B_1}\intersection
                                            \spanof{B_2}
           \end{equation*}
           it is easy enough to prove (but perhaps hard to apply).
         \partsitem The complement of a basis cannot be a basis
           because it contains the zero vector.
       \end{exparts}  
     \end{answer}
  \recommended \item 
    Consider how `dimension' interacts with `subset'.
    Assume \( U \) and \( W \) are both subspaces of some vector space, and
    that \( U\subseteq W \).
    \begin{exparts}
      \partsitem Prove that \( \dim (U)\leq\dim (W) \).
      \partsitem Prove that equality of dimension holds
        if and only if \( U=W \).
      \partsitem Show that the prior item does not hold if they are 
        infinite-dimensional.
    \end{exparts}
    \begin{answer}
       \begin{exparts}
          \partsitem A basis for \( U \) is a linearly independent set
            in \( W \)
            and so can be expanded via \nearbycorollary{cor:LIExpBas}
            to a basis for \( W \).
            The second basis has at least as many members as the first.
          \partsitem One direction is clear:~if \( V=W \) then they have the 
            same dimension.
            For the converse, let \( B_U \) be a basis for \( U \).
            It is a linearly independent subset of \( W \) and so can be
            expanded to a basis for \( W \).
            If \( \dim(U)=\dim(W) \) then this basis for \( W \) has no more
            members than does \( B_U \) and so equals \( B_U \).
            Since \( U \) and \( W \) have the same bases, they are equal.
          \partsitem Let \( W \) be the space of finite-degree polynomials and
            let \( U \) be the subspace of polynomials that have only
            even-powered terms.
            \begin{equation*}
               U=\set{a_0+a_1x^2+a_2x^4+\dots+a_nx^{2n}\suchthat
                          a_0,\ldots,a_n\in\Re} 
            \end{equation*}
            Both spaces have infinite dimension but \( U \) is a proper
            subspace.
       \end{exparts}  
     \end{answer}
    \item Here is an alternative proof of this section's main result,
      \nearbytheorem{th:AllBasesSameSize}.
      First is an example, then a lemma, then the theorem.
      \begin{exparts}
        \partsitem Express this vector from $\Re^3$ a 
          as a linear combination of members of the basis.
         \begin{equation*}
           B=\sequence{\colvec[r]{1 \\ 0  \\ 0},
             \colvec[r]{1 \\ 1 \\ 0},
             \colvec[r]{0 \\ 0 \\ 2}}
           \qquad
           \vec{v}=\colvec[r]{1 \\ 2 \\ 0}
         \end{equation*}
        \partsitem In that combination pick a basis vector with a 
          non-zero coefficient.
          Alter~$B$ by exchanging~$\vec{v}$ for that basis vector,
          to get a new sequence~$\hat{B}$.
          Check that $\hat{B}$ is also a basis for~$\Re^3$.
        \partsitem (Exchange Lemma)
          Assume that 
          \( B=\sequence{\vec{\beta}_1,\dots,\vec{\beta}_n} \) is a basis for a
          vector space, and that for the vector \( \vec{v} \)
          the relationship \( \vec{v}=\lincombo{c}{\vec{\beta}} \)
          has \( c_i\neq 0 \).
          Prove that exchanging  \( \vec{v} \) for \( \vec{\beta}_i \) 
          yields another basis for the space.
        \partsitem Use that, with induction, to prove
          \nearbytheorem{th:AllBasesSameSize}.
      \end{exparts}
      \begin{answer}
      \begin{exparts}
        \partsitem
          $\colvec[r]{1 \\ 2 \\ 0}
             =(-1)\cdot\colvec[r]{1 \\ 0  \\ 0}
              +2\colvec[r]{1 \\ 1 \\ 0}
              +0\cdot\colvec[r]{0 \\ 0 \\ 2}$
        \partsitem Two of the basis vectors are associated with 
           non-zero coefficients.
           We can for instance pick the first.
           \begin{equation*}
             \hat{B}=\sequence{\colvec[r]{1 \\ 2  \\ 0},
               \colvec[r]{1 \\ 1 \\ 0},
               \colvec[r]{0 \\ 0 \\ 2}}
           \end{equation*}
           Checking that it is another basis for \( \Re^3 \) is routine.
        \partsitem  
          Call the outcome of the exchange 
          \( \hat{B}=\sequence{\vec{\beta}_1,\dots,\vec{\beta}_{i-1},\vec{v},
            \vec{\beta}_{i+1},\dots,\vec{\beta}_n}  \).
          We first show that $\hat{B}$ is linearly independent.
          Any relationship 
          \( d_1\vec{\beta}_1+\dots+d_i\vec{v}+\dots+d_n\vec{\beta}_n=\zero \)
          among the members of $\hat{B}$, after substitution for $\vec{v}$,
          \begin{equation*}
          d_1\vec{\beta}_1+\dots
          +d_i\cdot(c_1\vec{\beta}_1+\dots+c_i\vec{\beta}_i+\dots+c_n\vec{\beta}_n)
          +\dots+d_n\vec{\beta}_n
          =\zero
          \tag*{($*$)}\end{equation*}
          gives a linear relationship among the members of $B$.
          The basis $B$ is linearly independent so the coefficient $d_ic_i$ of 
          $\vec{\beta}_i$ is zero.
          Because we assumed that $c_i$ is nonzero, $d_i=0$.
          Using this in equation~$(*)$ gives that all of the other $d$'s 
          are also zero.
          Therefore $\hat{B}$ is linearly independent.

          We finish by showing that $\hat{B}$ has the same span as $B$.
          Half of this argument, that 
          $\spanof{\smash{\hat{B}}}\subseteq\spanof{B}$, 
          is easy; we can write any member 
          $d_1\vec{\beta}_1+\dots+d_i\vec{v}+\dots+d_n\vec{\beta}_n$
          of $\spanof{\smash{\hat{B}}}$ as
          $
          d_1\vec{\beta}_1+\dots
          +d_i\cdot(c_1\vec{\beta}_1+\dots+c_n\vec{\beta}_n)
          +\dots+d_n\vec{\beta}_n
          $,
          which is a linear combination of linear combinations of members 
          of $B$, and hence is in $\spanof{B}$.
          For the $\spanof{B}\subseteq\spanof{\smash{\hat{B}}}$ half of the 
          argument, recall that if
          $\vec{v}=c_1\vec{\beta}_1+\dots+c_n\vec{\beta}_n$ with $c_i\neq 0$
          then we can rearrange the equation to
          $\vec{\beta}_i=(-c_1/c_i)\vec{\beta}_1+\dots+(1/c_i)\vec{v}+\dots
          +(-c_n/c_i)\vec{\beta}_n$.
          Now, consider any member
          $d_1\vec{\beta}_1+\dots+d_i\vec{\beta}_i+\dots+d_n\vec{\beta}_n$
          of $\spanof{B}$, substitute for $\vec{\beta}_i$ its expression 
          as a linear
          combination of the members of $\hat{B}$, and recognize, 
          as in the first half of this argument, that the result is a linear
          combination of linear combinations of members of $\hat{B}$, 
          and hence is in $\spanof{\smash{\hat{B}}}$.
        \partsitem Fix a vector space with at least one finite basis.
          Choose, from among all of this space's bases,
          one \( B=\sequence{\vec{\beta}_1,\dots,\vec{\beta}_n} \) 
          of minimal size.
          We will show that any other basis
          \( D=\smash{\sequence{\vec{\delta}_1,\vec{\delta}_2,\ldots}} \)
          also has the same number of members, $n$.
          Because \( B \) has minimal size, \( D \) has no fewer than 
          \( n \)~vectors.
          We will argue that it cannot have more than \( n \) vectors.

          The basis \( B \) spans the space and \( \vec{\delta}_1 \) 
          is in the space,
          so \( \vec{\delta}_1 \) is a nontrivial linear combination of 
          elements of \( B \).
          By the Exchange Lemma, we can swap \( \vec{\delta}_1 \) for a
          vector from \( B \), resulting in a basis \( B_1 \), 
          where one element is
          \( \vec{\delta}_1 \) and all of the \( n-1 \) other elements 
          are \( \vec{\beta} \)'s.

          The prior paragraph forms the basis step for an induction argument.
          The inductive step starts with a basis \( B_k \) (for \( 1\leq k<n \))
          containing \( k \) members of \( D \) and \( n-k \) members of 
          \( B \).
          We know that \( D \) has at least \( n \) members so there is a
          \( \vec{\delta}_{k+1} \).
          Represent it as a linear combination of elements of \( B_k \).
          The key point:~in that representation, 
          at least one of the nonzero scalars
          must be associated with a \( \vec{\beta}_i \) or else that
          representation would be a
          nontrivial linear relationship among elements of the linearly 
          independent set \( D \).
          Exchange \( \vec{\delta}_{k+1} \) for \( \vec{\beta}_i \) 
          to get a new basis
          \( B_{k+1} \) with one \( \vec{\delta} \) more and one 
          \( \vec{\beta} \) fewer than the previous basis \( B_k \).

          Repeat that until no \( \vec{\beta} \)'s remain, so
          that \( B_n \) contains  
          $\vec{\delta}_1,\dots,\vec{\delta}_n$.
          Now, \( D \) cannot have more than these \( n \) vectors because
          any \( \vec{\delta}_{n+1} \) that remains would be in the span of
          \( B_n \) (since it is a basis) 
          and hence would be a linear combination of the other 
          $\vec{\delta}$'s, contradicting that $D$ is linearly independent.
      \end{exparts}
      \end{answer}
  \puzzle \item 
    \cite{Sheffer}
    A library has $n$ books and $n+1$ subscribers. 
    Each subscriber read at least one book from the library. 
    Prove that there must exist two disjoint sets of subscribers who read 
    exactly the same books (that is, the union of the books 
    read by the subscribers in each set is the same).
    \begin{answer}
      (This answer is from a site comment by Yuzhou Gu.)
      For each person assign a vector $\vec{v}_i\in\R^n$, 
      where an entry is $1$ if the person reads the book, and $0$ otherwise. 
      Because the vectors are linear dependent, we will have a equation 
      of the form $\sum a_i v_i = 0$. 
      Then the set of $i$ such that $a_i>0$ and the set of $i$ such that 
      $a_i<0$ are two disjoint sets with same union of read books.
    \end{answer}
  \puzzle \item 
    \cite{Wohascum47}
    For any vector $\vec{v}$ in $\Re^n$ and any permutation $\sigma$
    of the numbers $1$, $2$, \ldots, $n$
    (that is, $\sigma$ is a rearrangement of those numbers into a new order), 
    define $\sigma(\vec{v})$ to be the vector whose components are 
    $v_{\sigma(1)}$, $v_{\sigma(2)}$, \ldots, and $v_{\sigma(n)}$
    (where $\sigma(1)$ is the first number in the rearrangement, etc.).
    Now fix $\vec{v}$ and let $V$ be the span of 
    $\set{\sigma(\vec{v})\suchthat \mbox{$\sigma$ permutes $1$, \ldots, $n$}}$.
    What are the possibilities for the dimension of $V$?
    \begin{answer}
      The possibilities for the dimension of $V$ are $0$, $1$, $n-1$,
      and $n$.

      To see this, first consider the case when all the coordinates of 
      $\vec{v}$ are equal.
      \begin{equation*}
         \vec{v}=\colvec{z \\ z \\ \vdots \\ z}
      \end{equation*}
      Then $\sigma(\vec{v})=\vec{v}$ for every permutation $\sigma$, so
      $V$ is just the span of $\vec{v}$, which has dimension $0$ or $1$
      according to whether $\vec{v}$ is $\zero$ or not.

      Now suppose not all the coordinates of $\vec{v}$ are equal; let $x$
      and $y$ with $x\neq y$ be among the coordinates of $\vec{v}$.
      Then we can find permutations $\sigma_1$ and $\sigma_2$ such that
      \begin{equation*}
        \sigma_1(\vec{v})=\colvec{x \\ y \\ a_3 \\ \vdots \\ a_n}
        \quad\text{and}\quad
        \sigma_2(\vec{v})=\colvec{y \\ x \\ a_3 \\ \vdots \\ a_n}
      \end{equation*}
      for some $a_3,\ldots,a_n\in\Re$.
      Therefore,
      \begin{equation*}
        \frac{1}{y-x}\bigl(\sigma_1(\vec{v})-\sigma_2(\vec{v})\bigr)=
        \colvec[r]{-1 \\ 1 \\ 0 \\ \vdotswithin{-1} \\ 0}
      \end{equation*}
      is in $V$.
      That is, $\vec{e}_2-\vec{e}_1\in V$, where $\vec{e}_1$, $\vec{e}_2$,
      \ldots, $\vec{e}_n$ is the standard basis for $\Re^n$.
      Similarly, $\vec{e}_3-\vec{e}_2$, \ldots, $\vec{e}_n-\vec{e}_1$
      are all in $V$.
      It is easy to see that the vectors
      $\vec{e}_2-\vec{e}_1$, $\vec{e}_3-\vec{e}_2$, \ldots, 
      $\vec{e}_n-\vec{e}_1$ are linearly independent (that is, form a linearly
      independent set), so $\dim V\geq n-1$.

      Finally, we can write
      \begin{align*}
        \vec{v} &=x_1\vec{e}_1+x_2\vec{e}_2+\dots+x_n\vec{e}_n  \\
                &=(x_1+x_2+\dots+x_n)\vec{e}_1+x_2(\vec{e}_2-\vec{e}_1)+\dots
                   +x_n(\vec{e}_n-\vec{e}_1)
      \end{align*}
      This shows that if $x_1+x_2+\dots+x_n=0$ then $\vec{v}$ is in the span
      of $\vec{e}_2-\vec{e}_1$, \ldots, $\vec{e_n}-\vec{e}_1$ (that is, is in
      the span of the set of those vectors); similarly, each $\sigma(\vec{v})$
      will be in this span, so $V$ will equal this span and $\dim V=n-1$.
      On the other hand, if $x_1+x_2+\cdots+x_n\neq 0$ then the above equation
      shows that $\vec{e}_1\in V$ and thus $\vec{e}_1,\dots,\vec{e}_n\in V$,
      so $V=\Re^n$ and $\dim V=n$.
    \end{answer}
\end{exercises}
\index{basis|)}



























\subsection{Vector Spaces and Linear Systems}
We will now reconsider linear systems and Gauss's Method, 
aided by the tools and terms of this chapter.
We will make three points.

For the first,
recall the insight from the Chapter One that
Gauss's Method works by taking linear
combinations of rows\Dash
if two matrices are
related by row operations $A\longrightarrow\cdots\longrightarrow B$ then
each row of $B$ is a linear combination of the rows of $A$. 
Therefore, the right setting in which to study row operations in general,
and Gauss's Method in particular, is the following vector space.

\begin{definition} \label{df:RowSpace}
%<*df:RowSpace>
The \definend{row space\/}\index{matrix!row space}\index{row space}
of a matrix is the span of the set of its rows.
The \definend{row rank\/}\index{row!rank}\index{matrix!row rank}\index{row rank}
is the dimension of this space, the number of linearly independent rows.
%</df:RowSpace>
\end{definition}

\begin{example}
If
\begin{equation*}
  A=\begin{mat}[r]
          2  &3  \\
          4  &6
       \end{mat}
\end{equation*}
then \( \rowspace{A} \) is this subspace of the space of
two-component row vectors.
\begin{equation*}
   \set{c_1\cdot\rowvec{2 &3}+c_2\cdot\rowvec{4  &6}
          \suchthat c_1,c_2\in\Re }
\end{equation*}
The second row vector is linearly dependent on the first and so we can
simplify the above description to 
$\set{c\cdot\rowvec{2 &3}\suchthat c\in\Re }$.
\end{example}

\begin{lemma}     \label{le:RowSpUnchByGR}
%<*lm:RowSpUnchByGR>
If two matrices \( A \) and \( B \) are related by a row operation
\begin{equation*}
  A\grstep{\rho_i\leftrightarrow\rho_j}B 
  \quad\text{or}\quad
  A\grstep{k\rho_i}B 
  \quad\text{or}\quad
  A\grstep{k\rho_i+\rho_j}B
\end{equation*}
(for $i\neq j$ and $k\neq 0$) then their row spaces are equal.
Hence, row-equivalent matrices have the same row space and therefore
the same row rank.
%</lm:RowSpUnchByGR>
\end{lemma}

\begin{proof}
%<*pf:RowSpUnchByGR0>
Corollary One.III.\ref{cor:RowsOfEqMatsLinCombos} 
shows that when \mbox{$A\longrightarrow B$} then each row of $B$
is a linear combination of the rows of \( A \).
That is, in the above terminology, each row of \( B \) 
is an element of the row space of $A$.
Then $\rowspace{B}\subseteq\rowspace{A}$ follows because a member of the
set $\rowspace{B}$ is a linear combination of the rows of $B$, so it
is a combination of combinations of the rows of $A$,
and by the Linear Combination Lemma is also a member of $\rowspace{A}$.
%</pf:RowSpUnchByGR0>

%<*pf:RowSpUnchByGR1>
For the other set containment, 
recall Lemma One.III.\ref{le:RowOpsRev}, that row operations are reversible
so
\mbox{$A\longrightarrow B$} if and only if
\mbox{$B\longrightarrow A$}.
Then 
$\rowspace{A}\subseteq\rowspace{B}$ follows as in the previous paragraph. 
%</pf:RowSpUnchByGR1>
% Thus the two sets are equal.
\end{proof}

Of course, Gauss's Method performs the row operations systematically,
with the goal of echelon form.

\begin{lemma}  \label{le:RowsEchMatLI}
%<*lm:RowsEchMatLI>
The nonzero rows of an echelon form matrix make up a linearly independent
set.
%</lm:RowsEchMatLI>
\end{lemma}

\begin{proof}
%<*pf:RowsEchMatLI>
Lemma~One.III.\ref{le:EchFormNoLinCombo}
says that no nonzero row of an echelon form matrix is
a linear combination of the other rows.
This result restates that using this chapter's terminology.
%</pf:RowsEchMatLI>
\end{proof}

Thus, in the language of this chapter,
Gaussian reduction works by eliminating linear dependences among rows, 
leaving the span unchanged, until no 
nontrivial linear relationships remain among the nonzero rows.
In short, Gauss's Method produces a basis for the row space.

\begin{example}
From any matrix, we can produce a basis for the row space by
performing Gauss's Method and taking the nonzero rows of the resulting
echelon form matrix.
For instance,
\begin{equation*}
  \begin{mat}[r]
    1  &3  &1  \\
    1  &4  &1  \\
    2  &0  &5
  \end{mat}
  \grstep[-2\rho_1+\rho_3]{-\rho_1+\rho_2}
  \repeatedgrstep{6\rho_2+\rho_3}
  \begin{mat}[r]
    1  &3  &1  \\
    0  &1  &0  \\
    0  &0  &3
  \end{mat}
\end{equation*}
produces the basis $\sequence{\rowvec{1 &3 &1},
            \rowvec{0 &1 &0},
            \rowvec{0 &0 &3} }$ for the row space.
This is a basis for the row space of both the starting and ending matrices, 
since the two row spaces are equal.
\end{example}

Using this technique, we can also find bases for spans 
not directly involving row vectors.

\begin{definition} \label{df:ColumnSpace}
%<*df:ColumnSpace>
The \definend{column space}\index{column!space}\index{matrix!column space}
of a matrix is the span of the set of its columns.
The \definend{column rank}\index{column!rank}\index{rank!column}
is the dimension of the column space, the number of linearly independent
columns.
%</df:ColumnSpace>
\end{definition}

Our interest in column spaces stems from our study of linear systems.
An example is that this system
\begin{equation*}
  \begin{linsys}{3}
    c_1  &+  &3c_2  &+  &7c_3  &=  &d_1  \\
   2c_1  &+  &3c_2  &+  &8c_3  &=  &d_2  \\
         &   &c_2   &+  &2c_3  &=  &d_3  \\
   4c_1  &   &      &+  &4c_3  &=  &d_4   
  \end{linsys}
\end{equation*}
has a solution if and only if the vector of \( d \)'s is a linear combination
of the other column vectors,
\begin{equation*}
  c_1\colvec[r]{1 \\ 2 \\ 0 \\ 4}
  +c_2\colvec[r]{3 \\ 3 \\ 1 \\ 0}
  +c_3\colvec[r]{7 \\ 8 \\ 2 \\ 4}
  =\colvec{d_1 \\ d_2 \\ d_3 \\ d_4}
\end{equation*}
meaning that the vector of \( d \)'s is in the column space of the
matrix of coefficients.

\begin{example}   \label{ex:BasisForColSpace}
Given this matrix,
\begin{equation*}
  \begin{mat}[r]
    1  &3  &7  \\
    2  &3  &8  \\
    0  &1  &2  \\
    4  &0  &4
  \end{mat}
\end{equation*}
to get a basis for the column space,
temporarily turn the columns into rows and reduce.
\begin{equation*}
    \begin{mat}[r]
       1  &2  &0  &4  \\
       3  &3  &1  &0  \\
       7  &8  &2  &4
    \end{mat}
  \grstep[-7\rho_1+\rho_3]{-3\rho_1+\rho_2}
  \repeatedgrstep{-2\rho_2+\rho_3}
  \begin{mat}[r]
     1  &2  &0  &4  \\
     0  &-3 &1  &-12\\
     0  &0  &0  &0
  \end{mat}
\end{equation*}
Now turn the rows back to columns.
\begin{equation*}
  \sequence{
     \colvec[r]{1 \\ 2 \\ 0 \\ 4},
     \colvec[r]{0 \\ -3 \\ 1 \\ -12} }
\end{equation*}
The result is a basis for the column space of the given matrix.
\end{example}

\begin{definition} \label{df:Transpose}
%<*df:Transpose>
The \definend{transpose}\index{transpose}\index{matrix!transpose}
of a matrix is the result of interchanging its rows and
columns, so that
column~\( j \) of the matrix \( A \) is row~\( j \)
of \( \trans{A} \) and vice versa.
%</df:Transpose>
\end{definition}
\noindent So we can summarize the prior example as ``transpose,
reduce, and transpose back.''

We can even, at the price of tolerating the as-yet-vague idea of vector
spaces being ``the same,''
use Gauss's Method to find bases for spans in other types of vector spaces.

\begin{example}
To get a basis for the span of
\( \set{x^2+x^4,2x^2+3x^4,-x^2-3x^4} \) in the space
\( \polyspace_4 \), think of these three polynomials 
as ``the same'' as the row vectors \( \rowvec{0 &0 &1 &0 &1} \),
\( \rowvec{0 &0 &2 &0 &3} \), and
\( \rowvec{0 &0 &-1 &0 &-3} \), apply Gauss's Method
\begin{equation*}
  \begin{mat}[r]
    0  &0  &1  &0  &1  \\
    0  &0  &2  &0  &3  \\
    0  &0  &-1 &0  &-3
  \end{mat}
  \grstep[\rho_1+\rho_3]{-2\rho_1+\rho_2}
  \repeatedgrstep{2\rho_2+\rho_3}
  \begin{mat}[r]
    0  &0  &1  &0  &1  \\
    0  &0  &0  &0  &1  \\
    0  &0  &0  &0  &0
  \end{mat}
\end{equation*}
and translate back to get the basis
\( \sequence{x^2+x^4,x^4} \).
(As mentioned earlier, 
we will make the phrase ``the same'' precise at the start of the next
chapter.)
\end{example}

Thus, the first point for this subsection is
that the tools of this chapter give us a more conceptual understanding of 
Gaussian reduction.

For the second point
observe that row operations on a matrix can change its column space.
\begin{equation*}
  \begin{mat}[r]
    1  &2  \\
    2  &4
  \end{mat}
  \grstep{-2\rho_1+\rho_2}
  \begin{mat}[r]
    1  &2  \\
    0  &0
  \end{mat}
\end{equation*}
The column space of the left-hand matrix contains vectors with
a second component that is nonzero
but the column space of the right-hand matrix 
contains only vectors whose second component is zero, so the two spaces 
are different.
This observation makes next result surprising.

\begin{lemma} \label{le:RowOpsNoChngColRnk}
%<*lm:RowOpsNoChngColRnk>
Row operations do not change the column rank.
%</lm:RowOpsNoChngColRnk>
\end{lemma}

\begin{proof}
%<*pf:RowOpsNoChngColRnk>
Restated, if $A$ reduces to $B$ 
then the column rank of $B$ equals the column rank of $A$.

This proof will be finished if we show that row operations do not affect 
linear relationships among columns, because the column rank is the size
of the largest set of unrelated columns.
That is, we will show that a relationship exists among columns
(such as that the fifth column is twice the
second plus the fourth) if and only if that relationship exists after
the row operation. 
But this is exactly the first theorem of this book,
Theorem One.I.\ref{th:GaussMethod}:~in a relationship
among columns,
\begin{equation*}
  c_1\cdot\colvec{a_{1,1} \\ a_{2,1} \\ \vdots \\ a_{m,1}}
   +\dots+
  c_n\cdot \colvec{a_{1,n} \\ a_{2,n} \\ \vdots \\ a_{m,n}}
   =\colvec[r]{0 \\ 0 \\ \vdotswithin{0} \\ 0}
\end{equation*}
row operations leave unchanged the set of solutions \( (c_1,\ldots,c_n) \).
%</pf:RowOpsNoChngColRnk>
\end{proof}

Another way to make the point that 
Gauss's Method has something to say about the column space as well as about 
the row space is with Gauss-Jordan reduction.
It ends with
the reduced echelon form of a matrix, as here.
\begin{equation*}
  \begin{mat}[r]
    1  &3  &1  &6  \\
    2  &6  &3  &16 \\
    1  &3  &1  &6
  \end{mat}
  \grstep{}\cdots\grstep{}
  \begin{mat}[r]
    1  &3  &0  &2  \\
    0  &0  &1  &4  \\
    0  &0  &0  &0
  \end{mat}
\end{equation*}
Consider the row space and the column space of this result.

The first point made earlier in this subsection says
that to get a basis for the row space we can just collect the rows with 
leading entries.
However, because this is in reduced echelon form, 
a basis for the column space is just as easy:~collect the
columns containing the leading entries,
\( \sequence{\vec{e}_1,\vec{e}_2} \).
% (Linear independence is obvious.
% As to span, the other columns are in the span 
% since they all have a third component of zero.)
Thus, for a reduced echelon form matrix
we can find bases for the row and column spaces
in essentially the same way, by taking the parts of the 
matrix, the rows or columns, containing the leading entries.

\begin{theorem} \label{th:RowRankEqualsColumnRank}
%<*th:RowRankEqualsColumnRank>
For any matrix, 
the row rank and column rank are equal.
%</th:RowRankEqualsColumnRank>
\end{theorem}

\begin{proof}
%<*pf:RowRankEqualsColumnRank0>
Bring the matrix to reduced echelon form.
Then the 
row rank equals the number of leading entries since that equals the
number of nonzero rows.
Then also, the number of leading entries
equals the column rank because the set of columns containing leading
entries consists of some of the \( \vec{e}_i \)'s from a standard
basis, and that set is linearly independent and spans the set of
columns.
Hence, in the reduced echelon form matrix, the row rank equals the column
rank, because each equals the number of leading entries.
%</pf:RowRankEqualsColumnRank0>

%<*pf:RowRankEqualsColumnRank1>
But \nearbylemma{le:RowSpUnchByGR} and \nearbylemma{le:RowOpsNoChngColRnk} 
show that the row rank and column rank are not changed
by using row operations to get to reduced echelon form.
Thus the row rank and the column rank of the original matrix are also equal.
%</pf:RowRankEqualsColumnRank1>
\end{proof}

\begin{definition} \label{df:Rank}
%<*df:Rank>
The \definend{rank\/}\index{rank} of a matrix is its row rank or column rank.
%</df:Rank>
\end{definition}

So the second point that we have made 
in this subsection is that the column space and row
space of a matrix have the same dimension.

Our final point is that the concepts that
we've seen arising naturally in the study of
vector spaces are exactly the ones that we have studied with linear systems.

\begin{theorem}  \label{th:RankVsSoltnSp}
%<*tm:RankVsSoltnSp>
For linear systems with \( n \) unknowns and with matrix of coefficients
\( A \), the statements
\begin{tfae}
   \item the rank of \( A \) is \( r \)
   \item the vector space of solutions of the associated homogeneous system has
     dimension \( n-r \)
\end{tfae}
are equivalent.
%</tm:RankVsSoltnSp>
\end{theorem}

\par\noindent So if the system has at least one particular solution 
then for the set of
solutions, the number of parameters equals $n-r$, the number of variables
minus the rank of the matrix of coefficients. 

\begin{proof}
%<*pf:RankVsSoltnSp>
The rank of \( A \) is \( r \) if and only if Gaussian reduction on \( A \)
ends with \( r \) nonzero rows.
That's true if and only if echelon form matrices row equivalent to \( A \)
have \( r \)-many leading variables.
That in turn holds if and only if there are \( n-r \) free variables.
%</pf:RankVsSoltnSp>
\end{proof}

\begin{corollary} \label{co:EquivToNonsingular}
%<*co:EquivToNonsingular>
Where the matrix $A$ is \( \nbyn{n} \), these statements 
\begin{tfae}
  \item the rank of \( A \) is \( n \)
  \item \( A \) is nonsingular
  \item the rows of \( A \) form a linearly independent set
  \item the columns of \( A \) form a linearly independent set
  \item any linear system whose matrix of coefficients is \( A \) has one and
    only one solution
\end{tfae}
are equivalent.
%</co:EquivToNonsingular>
\end{corollary}

\begin{proof}
%<*pf:EquivToNonsingular>
Clearly \( \text{(1)}\iff\text{(2)}\iff\text{(3)}\iff\text{(4)} \).
The last, \( \text{(4)}\iff\text{(5)} \), holds because a set of \( n \)
column vectors is linearly independent if and only if it is a basis for
\( \Re^n \), but the system
\begin{equation*}
  c_1\colvec{a_{1,1} \\ a_{2,1} \\ \vdots \\ a_{m,1}}
   +\dots+
  c_n\colvec{a_{1,n} \\ a_{2,n} \\ \vdots \\ a_{m,n}}
   =\colvec{d_1 \\ d_2 \\ \vdots \\ d_m}
\end{equation*}
has a unique solution for all choices of \( d_1,\dots,d_n\in\Re \) if and only
if the vectors of \( a \)'s on the left form a basis.
%</pf:EquivToNonsingular>
\end{proof}

\begin{remark}\cite{Munkres}  \label{rem:MunkresCommment}
Sometimes the results of this subsection are 
mistakenly remembered to say that the general solution
of an \( m \)~equations, \( n \)~unknowns system uses \( n-m \) parameters.
The number of equations is not the relevant number; rather, what matters
is the number of independent equations, the number of equations in a maximal
independent set.
Where there are \( r \) independent equations, the general solution involves
\( n-r \) parameters.
\end{remark}


\begin{exercises}
  \item 
    Transpose each.
    \begin{exparts*}
      \partsitem \( \begin{mat}[r]
                 2  &1  \\
                 3  &1
               \end{mat}  \)
      \partsitem \( \begin{mat}[r]
                 2  &1  \\
                 1  &3
               \end{mat}  \)
      \partsitem \( \begin{mat}[r]
                 1  &4  &3 \\
                 6  &7  &8
               \end{mat}  \)
      \partsitem \( \colvec[r]{0 \\ 0 \\ 0} \)
      \partsitem \( \rowvec{-1 &-2} \)
    \end{exparts*}
    \begin{answer}  
      \begin{exparts*}
        \partsitem \( \begin{mat}[r]
                   2  &3  \\
                   1  &1
                 \end{mat}  \)
        \partsitem \( \begin{mat}[r]
                   2  &1  \\
                   1  &3
                 \end{mat}  \)
        \partsitem \( \begin{mat}[r]
                   1  &6  \\
                   4  &7  \\
                   3  &8
                 \end{mat}  \)
        \partsitem \( \rowvec{0 &0 &0} \)
        \partsitem \( \colvec[r]{-1 \\ -2} \)
      \end{exparts*}  
    \end{answer}
  \recommended \item 
    Decide if the vector is in the row space of the matrix.
    \begin{exparts*}
       \partsitem \( \begin{mat}[r]
                  2  &1  \\
                  3  &1
                \end{mat}  \),
             \( \rowvec{1 &0} \)
       \partsitem \( \begin{mat}[r]
                  0  &1  &3  \\
                 -1  &0  &1  \\
                 -1  &2  &7
                \end{mat}  \),
             \( \rowvec{1 &1 &1} \)
    \end{exparts*}
    \begin{answer}  
       \begin{exparts}
         \partsitem Yes.
           To see if there are $c_1$ and $c_2$ such that
           \( c_1\cdot\rowvec{2 &1}+c_2\cdot\rowvec{3 &1}=\rowvec{1 &0} \),
           we solve
           \begin{equation*}
              \begin{linsys}{2}
                 2c_1  &+  &3c_2  &=  &1  \\
                  c_1  &+  &c_2   &=  &0  
              \end{linsys}
           \end{equation*}
           and get \( c_1=-1 \) and \( c_2=1 \). 
           Thus the vector is in the row space.
         \partsitem No.
           The equation
           \( c_1\rowvec{0 &1 &3}
              +c_2\rowvec{-1 &0 &1}
              +c_3\rowvec{-1 &2 &7}
              =\rowvec{1 &1 &1} \)
           has no solution.
           \begin{equation*}
             \begin{amat}[r]{3}
               0  &-1  &-1  &1  \\
               1  &0   &2   &1  \\
               3  &1   &7   &1
             \end{amat}
             \grstep{\rho_1\leftrightarrow\rho_2}
             \grstep{-3\rho_1+\rho_3}
             \repeatedgrstep{\rho_2+\rho_3}
             \begin{amat}[r]{3}
               1  &0   &2   &1  \\
               0  &-1  &-1  &1  \\
               0  &0   &0   &-1
             \end{amat}
           \end{equation*}
           Thus, the vector is not in the row space.
      \end{exparts}  
    \end{answer}
  \recommended \item 
    Decide if the vector is in the column space.
    \begin{exparts*}
      \partsitem \( \begin{mat}[r]
                 1  &1  \\
                 1  &1
               \end{mat}  \),
            \( \colvec[r]{1 \\ 3} \)
      \partsitem \( \begin{mat}[r]
                 1  &3  &1 \\
                 2  &0  &4 \\
                 1  &-3 &-3
               \end{mat}  \),
            \( \colvec[r]{1 \\ 0 \\ 0} \)
    \end{exparts*}
    \begin{answer}
       \begin{exparts}
          \partsitem No.
            To see if there are \( c_1,c_2\in\Re \) such that
            \begin{equation*}
              c_1\colvec[r]{1 \\ 1}
              +c_2\colvec[r]{1 \\ 1}
              =\colvec[r]{1 \\ -3}
            \end{equation*}
            we can use Gauss's Method on the resulting linear system.
            \begin{equation*}
              \begin{linsys}{2}
                 c_1  &+  &c_2  &=  &1  \\
                 c_1  &+  &c_2  &=  &3  
              \end{linsys}
              \grstep{-\rho_1+\rho_2}
              \begin{linsys}{2}
                 c_1  &+  &c_2  &=  &1  \\
                      &   &0    &=  &2  
              \end{linsys}
            \end{equation*}
            There is no solution and so the vector is not in the column space.
          \partsitem Yes.
            From this relationship
            \begin{equation*}
              c_1\colvec[r]{1 \\ 2 \\ 1}
              +c_2\colvec[r]{3 \\ 0 \\ -3}
              +c_3\colvec[r]{1 \\ 4 \\ -3}
              =\colvec[r]{1 \\ 0 \\ 0}
            \end{equation*}
            we get a linear system that, when we apply Gauss's Method,
            \begin{equation*}
              \begin{amat}[r]{3}
                1  &3  &1  &1  \\
                2  &0  &4  &0  \\
                1  &-3 &-3 &0
              \end{amat}
              \grstep[-\rho_1+\rho_3]{-2\rho_1+\rho_2}
              \repeatedgrstep{-\rho_2+\rho_3}
              \begin{amat}[r]{3}
                1  &3  &1  &1  \\
                0  &-6 &2  &-2 \\
                0  &0  &-6 &1
              \end{amat}
            \end{equation*}
            yields a solution.
            Thus, the vector is in the column space.
% sage: load('gauss_method.sage')
% sage: M = matrix(QQ, [[1,3, 1, 1], [2, 0, 4, 0], [1, -3, -3, 0]])
% sage: M
% [ 1  3  1  1]
% [ 2  0  4  0]
% [ 1 -3 -3  0]
% sage: gauss_method(M)
% [ 1  3  1  1]
% [ 2  0  4  0]
% [ 1 -3 -3  0]
% (' take', -2, 'times row', 1, 'plus row', 2)
% (' take', -1, 'times row', 1, 'plus row', 3)
% [ 1  3  1  1]
% [ 0 -6  2 -2]
% [ 0 -6 -4 -1]
% (' take', -1, 'times row', 2, 'plus row', 3)
% [ 1  3  1  1]
% [ 0 -6  2 -2]
% [ 0  0 -6  1]
      \end{exparts}  
    \end{answer}
  \recommended \item  
    Decide if the vector is in the column space of the matrix.
    \begin{exparts*}
      \partsitem \( \begin{mat}[r]
                   2  &1  \\
                   2  &5
                 \end{mat} \),~\( \colvec[r]{1 \\ -3}  \)
      \partsitem \( \begin{mat}[r]
                   4  &-8 \\
                   2  &-4
                 \end{mat} \),~\( \colvec[r]{0 \\ 1}  \)
      \partsitem \( \begin{mat}[r]
                   1  &-1  &1  \\
                   1  &1   &-1 \\
                  -1  &-1  &1
                \end{mat} \),~\( \colvec[r]{2 \\ 0 \\ 0}  \)
    \end{exparts*}
    \begin{answer}
      \begin{exparts}
        \partsitem Yes;
          we are asking if there are scalars \( c_1 \) and \( c_2 \) such that
          \begin{equation*}
            c_1\colvec[r]{2 \\ 2}+c_2\colvec[r]{1 \\ 5}=\colvec[r]{1 \\ -3}
          \end{equation*}
          which gives rise to a linear system
          \begin{equation*}
            \begin{linsys}{2}
              2c_1  &+  &c_2  &=  &1  \\
              2c_1  &+  &5c_2 &=  &-3
            \end{linsys}
            \grstep{-\rho_1+\rho_2}
            \begin{linsys}{2}
              2c_1  &+  &c_2  &=  &1  \\
                    &   &4c_2 &=  &-4
            \end{linsys}
          \end{equation*}
          and Gauss's Method produces \( c_2=-1 \) and \( c_1=1 \). 
          That is, there is indeed such a pair of scalars and so the vector
          is indeed in the column space of the matrix.
        \partsitem No;
          we are asking if there are scalars $c_1$ and $c_2$
          such that 
          \begin{equation*}
            c_1\colvec[r]{4 \\ 2}+c_2\colvec[r]{-8 \\ -4}=\colvec[r]{0 \\ 1}
          \end{equation*}
          and one way to proceed is to consider the resulting linear system
          \begin{equation*}
            \begin{linsys}{2}
              4c_1  &-  &8c_2  &=  &0  \\
              2c_1  &-  &4c_2  &=  &1
            \end{linsys}
          \end{equation*}
          that is easily seen to have no solution.
          Another way to proceed is to note
          that any linear combination of the columns on the left
          has a second component half as big as its first component, 
          but the  vector on the right does not meet that criterion.
        \partsitem Yes; we can simply observe that the vector
          is the first column minus the second.
          Or, failing that, setting up the relationship among the columns 
          \begin{equation*}
            c_1\colvec[r]{1 \\ 1 \\ -1}
             +c_2\colvec[r]{-1 \\ 1 \\ -1}
             +c_3\colvec[r]{1 \\ -1 \\ 1}
             =\colvec[r]{2 \\ 0 \\ 0}
          \end{equation*}
          and considering the resulting linear system
          \begin{equation*}
            \begin{linsys}{3}
              c_1  &-  &c_2  &+  &c_3  &=  &2  \\
              c_1  &+  &c_2  &-  &c_3  &=  &0  \\
             -c_1  &-  &c_2  &+  &c_3  &=  &0    
            \end{linsys}
            \grstep[\rho_1+\rho_3]{-\rho_1+\rho_2}
            \begin{linsys}{3}
              c_1  &-  &c_2  &+  &c_3  &=  &2  \\
                   &   &2c_2 &-  &2c_3 &=  &-2  \\
                   &   &-2c_2&+  &2c_3 &=  &2    
            \end{linsys}
            \grstep{\rho_2+\rho_3}
            \begin{linsys}{3}
              c_1  &-  &c_2  &+  &c_3  &=  &2  \\
                   &   &2c_2 &-  &2c_3 &=  &-2  \\
                   &   &     &   &0    &=  &0    
            \end{linsys}
          \end{equation*}
          gives the additional information (beyond that there is at least one
          solution) that there are infinitely many solutions.
          Parametrizing gives $c_2=-1+c_3$ and $c_1=1$, and so taking $c_3$ to 
          be zero gives a particular solution of $c_1=1$, $c_2=-1$, and
          $c_3=0$ (which is, of course, the observation made at the start).
      \end{exparts}  
    \end{answer}
  \recommended \item  
    Find a basis for the row space of this matrix.
    \begin{equation*}
      \begin{mat}[r]
         2  &0  &3  &4  \\
         0  &1  &1  &-1 \\
         3  &1  &0  &2  \\
         1  &0  &-4 &-1
       \end{mat}
    \end{equation*}
    \begin{answer}
     % Using Sage:
     % M = matrix(QQ, [[2,0,3,4], [0,1,1,-1], [3,1,0,2], [1,0,-4,-1]])
     % M1 = M.with_added_multiple_of_row(2,0,-3/2)
     % M2 = M1.with_added_multiple_of_row(3,0,-1/2)
     % M3 = M2.with_added_multiple_of_row(2,1,-1)
     % M4 = M3.with_added_multiple_of_row(3,2,-1)
     A routine Gaussian reduction
     \begin{equation*}
       \begin{mat}[r]
         2  &0  &3 &4   \\
         0  &1  &1  &-1 \\
         3  &1  &0  &2  \\
         1  &0  &-4  &-1
       \end{mat}
       \grstep[-(1/2)\rho_1+\rho_4]{-(3/2)\rho_1+\rho_3}
       \grstep{-\rho_2+\rho_3}
       \repeatedgrstep{-\rho_3+\rho_4}
       \begin{mat}[r]
         2  &0  &3      &4   \\
         0  &1  &1      &-1 \\
         0  &0  &-11/2  &-3  \\
         0  &0  &0      &0
       \end{mat}
     \end{equation*}
     suggests this basis
     $\sequence{\rowvec{2 &0 &3 &4},
         \rowvec{0 &1 &1 &-1},\rowvec{0 &0 &-11/2 &-3}}$.

      Another procedure, perhaps more convenient, is to swap rows first,
      % From Sage
      % M = matrix(QQ, [[2,0,3,4], [0,1,1,-1], [3,1,0,2], [1,0,-4,-1]])
      % M1 = M.with_swapped_rows(0,3)
      % M2 = M1.with_added_multiple_of_row(2,0,-3)
      % M3 = M2.with_added_multiple_of_row(3,0,-2)
      % M4 = M3.with_added_multiple_of_row(2,1,-1)
      % M5 = M4.with_added_multiple_of_row(3,2,-1)
      \begin{equation*}
        \grstep{\rho_1\swap\rho_4}
        \repeatedgrstep[-2\rho_1+\rho_4]{-3\rho_1+\rho_3}
        \repeatedgrstep{-\rho_2+\rho_3}
        \repeatedgrstep{-\rho_3+\rho_4}
        \begin{mat}[r]
           1  &0  &-4 &-1 \\
           0  &1  &1  &-1 \\
           0  &0  &11 &6  \\
           0  &0  &0  &0
         \end{mat}
      \end{equation*}
      leading to the basis
      \( \sequence{\rowvec{1 &0 &-4 &-1},\rowvec{0 &1 &1 &-1},
                   \rowvec{0 &0 &11 &6}} \).  
    \end{answer}
  \recommended \item 
    Find the rank of each matrix.
    \begin{exparts*}
      \partsitem \(
        \begin{mat}[r]
          2  &1  &3  \\
          1  &-1 &2  \\
          1  &0  &3
        \end{mat}  \)
      \partsitem \(
        \begin{mat}[r]
          1  &-1 &2  \\
          3  &-3 &6  \\
         -2  &2  &-4
        \end{mat}  \)
      \partsitem \(
        \begin{mat}[r]
          1  &3  &2  \\
          5  &1  &1  \\
          6  &4  &3
        \end{mat}  \)
      \partsitem \(
        \begin{mat}[r]
          0  &0  &0  \\
          0  &0  &0  \\
          0  &0  &0
        \end{mat}  \)
    \end{exparts*}
    \begin{answer}
      \begin{exparts}
         \partsitem This reduction
           \begin{equation*}
             \mbox{}
             \grstep[-(1/2)\rho_1+\rho_3]{-(1/2)\rho_1+\rho_2}
             \repeatedgrstep{-(1/3)\rho_2+\rho_3}
             \begin{mat}[r]
               2  &1     &3     \\
               0  &-3/2  &1/2   \\
               0  &0     &4/3
             \end{mat}
           \end{equation*}
           shows that the row rank, and hence the rank, is three.
         \partsitem Inspection of the columns shows that the others
           are multiples of the first (inspection of the rows shows the same
           thing).
           Thus the rank is one.
           
           Alternatively, the reduction
           \begin{equation*}
             \begin{mat}[r]
               1  &-1  &2  \\
               3  &-3  &6  \\   
               -2 &2   &-4
             \end{mat}
             \grstep[2\rho_1+\rho_3]{-3\rho_1+\rho_2}
             \begin{mat}[r]
               1  &-1  &2  \\
               0  &0   &0  \\   
               0  &0   &0 
             \end{mat}
           \end{equation*}
           shows the same thing.
         \partsitem This calculation
           \begin{equation*}
             \begin{mat}[r]
               1  &3  &2  \\
               5  &1  &1  \\
               6  &4  &3  
             \end{mat}
             \grstep[-6\rho_1+\rho_3]{-5\rho_1+\rho_2}
             \repeatedgrstep{-\rho_2+\rho_3}
             \begin{mat}[r]
               1  &3   &2  \\
               0  &-14 &-9 \\
               0  &0   &0  
             \end{mat}
           \end{equation*}
           shows that the rank is two.
         \partsitem The rank is zero.
       \end{exparts}  
     \end{answer}
  \item Give a basis for the column space of this matrix.
    Give the matrix's rank. 
    \begin{equation*}
      \begin{mat}
        1 &3 &-1 &2 \\
        2 &1 &1  &0 \\
        0 &1 &1  &4
      \end{mat}
    \end{equation*}
  \begin{answer}
    We want a basis for this span.
    \begin{equation*}
      \spanof{\colvec{1 \\ 2 \\ 0},
              \colvec{3 \\ 1 \\ 1},
              \colvec{-1 \\ 1 \\ 1},
              \colvec{2 \\ 0 \\ 4}}\subseteq\Re^3
    \end{equation*}
    The most straightforward approach 
    is to transpose those columns to rows, use Gauss's Method
    to find a basis for the span of the rows, and then transpose them back to 
    columns. 
    \begin{equation*}
      \begin{mat}
        1 &2 &0 \\
        3 &1 &1 \\
       -1 &1 &1 \\
        2 &0 &4
      \end{mat}
      \grstep[\rho_1+\rho_3 \\ -2\rho_1+\rho_4]{-3\rho_1+\rho_2}
      \grstep[-(4/5)\rho_2+\rho_4]{(3/5)\rho_2+\rho_3}
      \grstep{-2\rho_3+\rho_4}
      \begin{mat}
        1 &2  &0 \\
        0 &-5 &1 \\
        0 &0  &8/5 \\
        0 &0  &0
      \end{mat}
    \end{equation*}
    Discard the zero vector as showing that there was  a redundancy among the 
    starting vectors, to get this basis for the column space.
    \begin{equation*}
      \sequence{
        \colvec{1 \\ 2 \\ 0},
        \colvec{0 \\ -5 \\ 1},
        \colvec{0 \\ 0 \\ 8/5}
        }
    \end{equation*}
    The matrix's rank is the dimension of its column space, so it is three.
    (It is also equal to the dimension of its row space.)  
  \end{answer}
  \recommended \item 
    Find a basis for the span of each set.
    \begin{exparts}
      \partsitem \(
        \set{\rowvec{1 &3},
          \rowvec{-1 &3},
          \rowvec{1 &4},
          \rowvec{2 &1}  }\subseteq\matspace_{\nbym{1}{2}} \)
      \partsitem \(
         \set{\colvec[r]{1 \\2 \\1},
           \colvec[r]{3 \\ 1 \\ -1},
           \colvec[r]{1 \\ -3 \\ -3}  }\subseteq\Re^3  \)
      \partsitem \(  \set{1+x,1-x^2,3+2x-x^2}\subseteq\polyspace_3  \)
      \partsitem \( \set{
          \begin{mat}[r]
            1  &0  &1  \\
            3  &1  &-1
          \end{mat},
          \begin{mat}[r]
            1  &0  &3  \\
            2  &1  &4
          \end{mat},
          \begin{mat}[r]
           -1  &0  &-5 \\
           -1  &-1 &-9
          \end{mat}  }  \subseteq\matspace_{\nbym{2}{3}}  \)
    \end{exparts}
    \begin{answer}
      \begin{exparts}
        \partsitem This reduction
          \begin{equation*}
             \begin{mat}[r]
               1  &3  \\
              -1  &3  \\
               1  &4  \\
               2  &1
             \end{mat}
             \grstep[-\rho_1+\rho_3 \\ -2\rho_1+\rho_4]{\rho_1+\rho_2}
             \grstep[(5/6)\rho_2+\rho_4]{-(1/6)\rho_2+\rho_3}
             \begin{mat}[r]
               1  &3  \\
               0  &6  \\
               0  &0  \\
               0  &0
             \end{mat}
          \end{equation*}
          gives \( \sequence{\rowvec{1 &3},\rowvec{0 &6}} \).
        \partsitem Transposing and reducing
          \begin{equation*}
             \begin{mat}[r]
               1  &2  &1  \\
               3  &1  &-1 \\
               1  &-3 &-3
             \end{mat}
             \grstep[-\rho_1+\rho_3]{-3\rho_1+\rho_2}
             \begin{mat}[r]
               1  &2  &1  \\
               0  &-5 &-4 \\
               0  &-5 &-4
             \end{mat}
             \grstep{-\rho_2+\rho_3}
             \begin{mat}[r]
               1  &2  &1  \\
               0  &-5 &-4 \\
               0  &0  &0
             \end{mat}
          \end{equation*}
          and then transposing back gives this basis.
          \begin{equation*}
            \sequence{\colvec[r]{1 \\ 2 \\ 1},
              \colvec[r]{0 \\ -5 \\ -4}   }
          \end{equation*}
        \partsitem Notice first that the surrounding space is as 
          $\polyspace_3$, not $\polyspace_2$. 
          Then, taking the first polynomial $1+1\cdot x+0\cdot x^2+0\cdot x^3$ 
          to be ``the same'' as the row vector $\rowvec{1 &1 &0 &0}$, etc., 
          leads to
          \begin{equation*}
            \begin{mat}[r]
              1  &1  &0  &0 \\
              1  &0  &-1 &0 \\
              3  &2  &-1 &0
            \end{mat}
            \grstep[-3\rho_1+\rho_3]{-\rho_1+\rho_2}
            \repeatedgrstep{-\rho_2+\rho_3}
            \begin{mat}[r]
              1  &1  &0  &0 \\
              0  &-1 &-1 &0 \\
              0  &0  &0  &0
            \end{mat}
          \end{equation*}
          which yields the basis \( \sequence{1+x,-x-x^2} \).
        \partsitem Here ``the same'' gives
          \begin{equation*}
            \begin{mat}[r]
              1  &0  &1  &3  &1  &-1  \\
              1  &0  &3  &2  &1  &4   \\
             -1  &0  &-5  &-1 &-1 &-9
            \end{mat}
            \grstep[\rho_1+\rho_3]{-\rho_1+\rho_2}
            \repeatedgrstep{2\rho_2+\rho_3}
            \begin{mat}[r]
              1  &0  &1  &3  &1  &-1  \\
              0  &0  &2  &-1 &0  &5   \\
              0  &0  &0  &0  &0  &0
            \end{mat}
          \end{equation*}
          leading to this basis.
          \begin{equation*}
            \sequence{
            \begin{mat}[r]
              1  &0  &1  \\
              3  &1  &-1
            \end{mat},
            \begin{mat}[r]
              0  &0  &2  \\
             -1  &0  &5
            \end{mat}  }
          \end{equation*}
      \end{exparts}   
     \end{answer}
  \item Give a basis for the span of each set, in the natural vector space.
    \begin{exparts}
      \partsitem
         $\set{\colvec{1 \\ 1 \\ 3},
              \colvec{-1 \\ 2 \\ 0},
              \colvec{0  \\ 12 \\ 6}}$
      \partsitem $\set{x+x^2, 2-2x, 7, 4+3x+2x^2}$
    \end{exparts}
    \begin{answer}
      \begin{exparts}
        \partsitem
          Transpose the columns to rows, bring to echelon form 
          (and then lose any zero rows), and transpose back to columns.
          \begin{equation*}
            \begin{mat}
              1 &1  &3 \\
             -1 &2  &0 \\
              0 &12 &6
            \end{mat}
            \grstep{\rho_1+\rho_2}
            \grstep{-4\rho_2+\rho_3}
            \begin{mat}
              1 &1  &3 \\
              0 &3  &3 \\
              0 &0  &-6
            \end{mat}
          \end{equation*}
          One basis for the span is this.
          \begin{equation*}
            \sequence{
              \colvec{1 \\ 1 \\ 3},
              \colvec{0 \\ 3 \\ 3},
              \colvec{0 \\ 0 \\ -6}
          }
          \end{equation*}
       \partsitem
          As in the prior part we think of those as rows, to take advantage of 
          the work we've done with Gauss's Method.
          \begin{multline*}
            \begin{mat}
              0 &1  &1 \\
              2 &-2 &0 \\
              7 &0 &0 \\
              4 &3  &2
            \end{mat}
            \grstep{\rho_1\leftrightarrow\rho_2}
            \grstep[2\rho_1+\rho_4]{-(7/2)\rho_1+\rho_3}
            \grstep[-7\rho_2+\rho_4]{-7\rho_2+\rho_3}          \\
            \grstep{-(5/7)\rho_3+\rho_4}
            \begin{mat}
              2 &-2  &0 \\
              0 &1 &1 \\
              0 &0  &-7 \\
              0 &0 &0
            \end{mat}
          \end{multline*}
          One basis for the span of that set is 
          $\sequence{2-2x, x+x^2, -5x^2}$.
      \end{exparts}
    \end{answer}
  \item 
    Which matrices have rank zero?
    Rank one?
    \begin{answer}
      Only the zero matrices have rank of zero.
      The only matrices of rank one have the form
      \begin{equation*}
        \begin{mat}
           k_1\cdot \rho  \\
            \vdots  \\
           k_m\cdot \rho
        \end{mat}
      \end{equation*}
      where \( \rho \) is some nonzero row vector, and not all of the 
      \( k_i \)'s are zero.
      (\textit{Remark.}
      We can't simply say that all of the rows are multiples of the first
      because the first row might be the zero row.
      \textit{Another Remark.}
      The above also applies with `column' replacing `row'.)  
    \end{answer}
  \recommended \item
    Given \( a,b,c\in\Re \), what choice of \( d \) will cause this matrix 
    to have the rank of one?
    \begin{equation*}
      \begin{mat}
         a  &b  \\
         c  &d
      \end{mat}
    \end{equation*}
    \begin{answer}
      If \( a\neq 0 \) then a choice of \( d=(c/a)b \) will make the second
      row be a multiple of the first, specifically, \( c/a \) times the first.
      If \( a=0 \) and \( b=0 \) then any non-\( 0 \) choice for \( d \) 
      will ensure that
      the second row is nonzero.
      If \( a=0 \) and \( b\neq 0 \) and \( c=0 \) then any choice for \( d \)
      will do, since the matrix will automatically have rank one (even with
      the choice of $d=0$).
      Finally, if \( a=0 \) and \( b\neq 0 \) and \( c\neq 0 \) then
      no choice for \( d \) will suffice because the matrix is sure to have 
      rank two.   
    \end{answer}
  \item 
    Find the column rank of this matrix.
    \begin{equation*}
      \begin{mat}[r]
        1  &3  &-1  &5  &0  &4  \\
        2  &0  &1   &0  &4  &1
      \end{mat}
    \end{equation*}
    \begin{answer}
      The column rank is two.
      One way to see this is by inspection\Dash the column space consists of
      two-tall columns and so can have a dimension of at least two, and we
      can easily find two columns that together form a linearly independent
      set (the fourth and fifth columns, for instance).
      Another way to see this is to recall that 
      the column rank equals the row rank, and to perform Gauss's Method, 
      which leaves two nonzero rows.
    \end{answer}
  \item 
    Show that a linear system with at least one solution has at most 
    one solution if
    and only if the matrix of coefficients has rank equal to the number of
    its columns.
    \begin{answer}
      We apply \nearbytheorem{th:RankVsSoltnSp}.  
      The number of columns of a matrix of coefficients $A$ of a linear 
      system equals the number $n$ of unknowns.
      A linear system with at least one solution has at most one solution
      if and only if the space of solutions of the associated homogeneous
      system has dimension zero (recall:~in the 
      `$\text{General}=\text{Particular}+\text{Homogeneous}$' equation
       $\vec{v}=\vec{p}+\vec{h}$, provided that such a $\vec{p}$ exists,
       the solution $\vec{v}$ is unique if and only if the vector $\vec{h}$
       is unique, namely $\vec{h}=\zero$).
       But that means, by the theorem, that $n=r$.
    \end{answer}
  \recommended \item
    If a matrix is \( \nbym{5}{9} \), which set must be dependent, its set of
    rows or its set of columns?
    \begin{answer}
      The set of columns must be dependent because the rank of the matrix
      is at most five while there are nine columns.  
    \end{answer}
  \item 
    Give an example to show that, despite that they have the same 
    dimension, the row space and column space of a matrix need not be equal.
    Are they ever equal?
    \begin{answer}
      There is little danger of their being equal since the row space
      is a set of row vectors while the column space is a set of
      columns (unless the matrix is \( \nbyn{1} \), in which case 
      the two spaces must be equal).

      \textit{Remark.}
      Consider
      \begin{equation*}
        A=\begin{mat}[r]
            1  &3  \\
            2  &6
          \end{mat}
      \end{equation*}
      and note that the row space is the set of 
      all multiples of \( \rowvec{1 &3} \)
      while the column space consists of multiples of
      \begin{equation*}
        \colvec[r]{1 \\ 2}
      \end{equation*}
      so we also cannot argue that the two spaces must be simply 
      transposes of each other.  
    \end{answer}
  \item 
    Show that the set
    \( \set{(1,-1,2,-3),(1,1,2,0),(3,-1,6,-6)} \) does not have the
    same span as \( \set{(1,0,1,0),(0,2,0,3)} \).
    What, by the way, is the vector space?
    \begin{answer}
      First, the vector space is the set of four-tuples of real numbers, under
      the natural operations.
      Although this is not the set of four-wide row vectors, the difference
      is slight\Dash it is ``the same'' as that set.
      So we will treat the four-tuples like four-wide vectors.

      With that, one way to see that \( (1,0,1,0) \) is not in the span 
      of the first set is to note that this reduction
      \begin{equation*}
        \begin{mat}[r]
          1  &-1  &2  &-3  \\
          1  &1   &2  &0   \\
          3  &-1  &6  &-6
        \end{mat}
        \grstep[-3\rho_1+\rho_3]{-\rho_1+\rho_2}
        \repeatedgrstep{-\rho_2+\rho_3}
        \begin{mat}[r]
          1  &-1  &2  &-3  \\
          0  &2   &0  &3   \\
          0  &0   &0  &0
        \end{mat}
      \end{equation*}
      and this one
      \begin{equation*}
        \begin{mat}[r]
          1  &-1  &2  &-3  \\
          1  &1   &2  &0   \\
          3  &-1  &6  &-6  \\
          1  &0   &1  &0
        \end{mat}
        \grstep[-3\rho_1+\rho_3 \\ -\rho_1+\rho_4]{-\rho_1+\rho_2}
        \repeatedgrstep[-(1/2)\rho_2+\rho_4]{-\rho_2+\rho_3}
        \repeatedgrstep{\rho_3\leftrightarrow\rho_4}
        \begin{mat}[r]
          1  &-1  &2  &-3  \\
          0  &2   &0  &3   \\
          0  &0   &-1 &3/2 \\
          0  &0   &0  &0   
        \end{mat}
      \end{equation*}
      yield matrices differing in rank.
      This means that addition of $(1,0,1,0)$ to the set of the first three
      four-tuples increases the rank, and hence the span, of that set.
      Therefore $(1,0,1,0)$ is not already in the span.
    \end{answer}
  \recommended \item 
    Show that this set of column vectors
    \begin{equation*}
      \set{\colvec{d_1 \\ d_2 \\ d_3}
           \suchthat
           \text{there are $x$, $y$, and $z$ such that:\ }
           \begin{linsys}{3}
             3x  &+  &2y  &+  &4z  &=   &d_1   \\
              x  &   &    &-  &z   &=   &d_2   \\
             2x  &+  &2y  &+  &5z  &=   &d_3   
           \end{linsys}
           }
    \end{equation*}
    is a subspace of \( \Re^3 \).
    Find a basis.
    \begin{answer}
      It is a subspace because it is the column space of the matrix 
      \begin{equation*}
        \begin{mat}[r]
          3  &2  &4  \\
          1  &0  &-1 \\
          2  &2  &5
        \end{mat}
      \end{equation*}
      of coefficients.
      To find a basis for the column space, 
      \begin{equation*}
        \set{c_1\colvec[r]{3 \\ 1 \\ 2}
             +c_2\colvec[r]{2 \\ 0 \\ 2}
             +c_3\colvec[r]{4 \\ -1 \\ 5}
             \suchthat c_1,c_2,c_3\in\Re}
      \end{equation*}
      we eliminate linear relationships among the three column vectors 
      from the spanning set by transposing, reducing,
      \begin{equation*}
         \begin{mat}[r]
           3  &1  &2  \\
           2  &0  &2  \\
           4  &-1 &5
         \end{mat}
         \grstep[-(4/3)\rho_1+\rho_3]{-(2/3)\rho_1+\rho_2}
         \repeatedgrstep{-(7/2)\rho_2+\rho_3}
         \begin{mat}[r]
           3  &1     &2  \\
           0  &-2/3  &2/3  \\
           0  &0     &0
         \end{mat}
      \end{equation*}
      omitting the zero row, and transposing back.
      \begin{equation*}
        \sequence{ \colvec[r]{3 \\ 1 \\ 2},
                   \colvec[r]{0 \\ -2/3 \\ 2/3}  }
      \end{equation*}  
    \end{answer}
  \item 
    Show that the transpose operation is 
    linear:\index{transpose!is linear}
    \begin{equation*}
      \trans{(rA+sB)}  = r\trans{A}+s\trans{B}
    \end{equation*}
    for \( r,s\in\Re \) and \( A,B\in\matspace_{\nbym{m}{n}} \).
    \begin{answer}
      We can do this as a straightforward calculation.
      \begin{align*}
        \trans{(rA+sB)}
        &=\trans{\begin{mat}
             ra_{1,1}+sb_{1,1}  &\ldots &ra_{1,n}+sb_{1,n} \\
                          &\vdots            \\
             ra_{m,1}+sb_{m,1}  &\ldots &ra_{m,n}+sb_{m,n}
                \end{mat}  }  \\
        &=\begin{mat}
           ra_{1,1}+sb_{1,1}  &\ldots &ra_{m,1}+sb_{m,1}  \\
                            &\vdots                           \\
           ra_{1,n}+sb_{1,n}  &\ldots &ra_{m,n}+sb_{m,n}
         \end{mat}           \\
        &=\begin{mat}
           ra_{1,1}  &\ldots &ra_{m,1}  \\
                    &\vdots                       \\
           ra_{1,n}  &\ldots &ra_{m,n}
         \end{mat}           
        +\begin{mat}
           sb_{1,1}  &\ldots &sb_{m,1}  \\
                    &\vdots                       \\
           sb_{1,n}  &\ldots &sb_{m,n}
         \end{mat}           \\
        &=r\trans{A}+s\trans{B}
      \end{align*}  
    \end{answer}
  \recommended \item  
    In this subsection we have shown that 
    Gaussian reduction finds a basis for the row space.
    \begin{exparts}
      \partsitem Show that 
        this basis is not unique\Dash different reductions may yield
        different bases.
      \partsitem Produce matrices with equal row spaces but unequal numbers of
        rows.
      \partsitem Prove that two matrices have equal row spaces if and
        only if after
        Gauss-Jordan reduction they have the same nonzero rows.
    \end{exparts}
    \begin{answer}
      \begin{exparts}
        \partsitem These reductions give different bases.
          \begin{equation*}
            \begin{mat}[r]
              1  &2  &0  \\
              1  &2  &1
            \end{mat}
            \grstep{-\rho_1+\rho_2}
            \begin{mat}[r]
              1  &2  &0  \\
              0  &0  &1
            \end{mat}
            \qquad
            \begin{mat}[r]
              1  &2  &0  \\
              1  &2  &1
            \end{mat}
            \grstep{-\rho_1+\rho_2}
            \repeatedgrstep{2\rho_2}
            \begin{mat}[r]
              1  &2  &0  \\
              0  &0  &2
            \end{mat}
          \end{equation*}
        \partsitem An easy example is this.
          \begin{equation*}
            \begin{mat}[r]
              1  &2  &1  \\
              3  &1  &4
            \end{mat}
            \qquad
            \begin{mat}[r]
              1  &2  &1  \\
              3  &1  &4  \\
              0  &0  &0
            \end{mat}
          \end{equation*}
          This is a less simplistic example.
          \begin{equation*}
            \begin{mat}[r]
              1  &2  &1  \\
              3  &1  &4
            \end{mat}
            \qquad
            \begin{mat}[r]
              1  &2  &1  \\
              3  &1  &4  \\
              2  &4  &2  \\
              4  &3  &5
            \end{mat}
          \end{equation*}
        \partsitem 
          % We give two proofs.
          % The first is a bit lower-level but more constructive.
          % The second is higher level.

          % \textit{First proof.}
          % Assume that \( A \) and \( B \) are matrices with equal row spaces.
          % Construct a matrix \( C \) with the rows of \( A \) above the rows
          % of \( B \), and another matrix \( D \) with the rows of \( B \)
          % above the rows of \( A \).
          % \begin{equation*}
          %   C=\begin{mat}
          %        A \\ B
          %     \end{mat}
          %   \qquad
          %   D=\begin{mat}
          %        B \\ A
          %     \end{mat}
          % \end{equation*}
          % Observe that \( C \) and \( D \) are row-equivalent (via a sequence
          % of row-swaps) and so Gauss-Jordan reduce to the same reduced
          % echelon form matrix.

          % Because the row spaces are equal, the rows of \( B \)
          % are linear combinations of the rows of \( A \) so Gauss-Jordan
          % reduction on \( C \) simply turns the rows of \( B \) to zero rows
          % and thus the nonzero rows of $C$ are just the nonzero rows obtained
          % by Gauss-Jordan reducing \( A \).
          % The same can be said for the matrix \( D \)\Dash Gauss-Jordan
          % reduction on \( D \) gives the same non-zero rows as are produced
          % by reduction on \( B \) alone.
          % Therefore, 
          % \( A \) yields the same nonzero rows as \( C  \), which yields
          % the same nonzero rows as \( D \), which yields the same nonzero
          % rows as \( B \).

          % \textit{Second proof (N~Schenker).}
          Because the row spaces of \( A \) and \( B \) are equal, 
          the two are row equivalent. 
          Because each row equivalence class contains 
          a unique reduced echelon form 
          (Theorem One.III.\ref{th:ReducedEchelonFormIsUnique}), 
          the reduced echelon form of \( A \) 
          must equal the reduced echelon form of \( B \).
      \end{exparts}  
    \end{answer}
  \item 
    Why is there not a problem with \nearbyremark{rem:MunkresCommment} 
    in the case that \( r \) is bigger than \( n \)?
    \begin{answer}
      It cannot be bigger.  
    \end{answer}
  \item  
    Show that the row rank of an \( \nbym{m}{n} \) matrix is at most
    \( m \).
    Is there a better bound?
    \begin{answer}
      The number of rows in a maximal linearly independent set cannot
      exceed the number of rows.
      A better bound (the bound that is, in general, the best possible) is
      the minimum of \( m \) and \( n \), because the row rank equals the
      column rank.
     \end{answer}
  \item
    Show that the rank of a matrix equals the rank of its transpose.
    \begin{answer}
      Because the rows of a matrix \( A \) are the columns
      of \( \trans{A} \) the dimension of the row space of \( A \) equals
      the dimension of the column space of \( \trans{A} \).
      But the dimension of the row space of \( A \) is the rank of \( A \) and
      the dimension of the column space of \( \trans{A} \) is the rank of
      \( \trans{A} \).
      Thus the two ranks are equal.
    \end{answer}
  \item 
    True or false:~the column space of a matrix equals the row space of its
    transpose.
    \begin{answer}
      False.
      The first is a set of columns while the second is a set of rows.

      This example, however,
      \begin{equation*}
         A=\begin{mat}[r]
             1  &2  &3  \\
             4  &5  &6
           \end{mat},
         \qquad
         \trans{A}=\begin{mat}[r]
             1  &4  \\
             2  &5  \\
             3  &6
           \end{mat}
      \end{equation*}
      indicates that as soon as we have a formal meaning for ``the same'',
      we can apply it here:
      \begin{equation*}
        \colspace{A}=\spanof{\set{
                        \colvec[r]{1 \\ 4},
                        \colvec[r]{2 \\ 5},
                        \colvec[r]{3 \\ 6} }  }
      \end{equation*}
      while
      \begin{equation*}
         \rowspace{\trans{A}}=\spanof{\set{
                                 \rowvec{1 &4},
                                 \rowvec{2 &5},
                                 \rowvec{3 &6} }  }
      \end{equation*}  
      are ``the same'' as each other.
    \end{answer}
  \recommended \item  
    We have seen that a row operation may change the column space.
    Must it?
    \begin{answer}
      No.
      Here, Gauss's Method does not change the column space.
      \begin{equation*}
        \begin{mat}[r]
          1  &0  \\
          3  &1
        \end{mat}
        \grstep{-3\rho_1+\rho_2}
        \begin{mat}[r]
          1  &0  \\
          0  &1
        \end{mat}
      \end{equation*} 
    \end{answer}
  \item
    Prove that a linear system has a solution if and only if that system's
    matrix of coefficients has the same rank as its augmented matrix.
    \begin{answer}
      A linear system
      \begin{equation*}
        c_1\vec{a}_1+\dots+c_n\vec{a}_n=\vec{d}
      \end{equation*}
      has a solution if and only if \( \vec{d} \) is in the span of
      the set \( \set{\vec{a}_1,\dots,\vec{a}_n} \).
      That's true if and only if the column rank of the augmented matrix
      equals the column rank of the matrix of coefficients.
      Since rank equals the column rank, the system has a solution if and
      only if the rank of its augmented matrix equals the rank of its matrix
      of coefficients.  
    \end{answer}
  \item 
    An \( \nbym{m}{n} \) matrix has 
    \definend{full row rank}\index{full row rank}\index{row rank!full} 
    if its row rank
    is \( m \), and it has 
    \definend{full column rank}\index{full column rank}\index{column!rank!full}
    if its column rank is \( n \).
    \begin{exparts}
      \partsitem Show that 
        a matrix can have both full row rank and full column rank
        only if it is square.
      \partsitem Prove that the linear system with matrix of
        coefficients \( A \)
        has a solution for any \( d_1 \), \ldots, \( d_n \)'s on the 
        right side if and only if \( A \) has full row rank.
      \partsitem Prove that a homogeneous system
        has a unique solution if and only if its matrix of coefficients
        $A$ has full column rank.
      \partsitem Prove that the statement 
       ``if a system with matrix of coefficients \( A \) has
        any solution then it has a unique solution'' holds
        if and only if \( A \) has full column rank.
    \end{exparts}
    \begin{answer}
      \begin{exparts}
        \partsitem Row rank equals column rank so each is at most the 
          minimum of the number of rows and columns.
          Hence both can be full only if the number of rows equals the number
          of columns.
          (Of course, the converse does not hold:~a 
          square matrix need not have full row rank or full column rank.)
        \partsitem If \( A \) has full row rank then, no matter what 
          the right-hand
          side, Gauss's Method on the augmented matrix ends with a leading
          one in each row and none of those leading ones in the
          furthest right column (the ``augmenting'' column).
          Back substitution then gives a solution.

          On the other hand, if the linear system lacks a solution for some
          right-hand side it can only be because Gauss's Method leaves some row
          so that it has all zeroes on the left of the ``augmenting'' bar 
          and has a nonzero entry on the right.
          Thus, if $A$ does not have a solution for some right-hand sides,
          then \( A \) does not have full row rank because some of its rows
          have been eliminated.
        \partsitem The matrix \( A \) has full column rank if and only
          if its columns
          form a linearly independent set.
          That's equivalent to the existence of only the trivial linear
          relationship among the columns, so the only solution of the
          system is where each variable is $0$.
        \partsitem The matrix \( A \) has full column rank if and only
          if the set of
          its columns is linearly independent, and so forms a basis for
          its span.
          That's equivalent to the existence of a unique linear representation
          of all vectors in that span.
          That proves it, 
          since any linear representation of a vector in the span is a 
          solution of the linear system.
      \end{exparts}  
    \end{answer}
  \item 
    How would the conclusion of \nearbylemma{le:RowSpUnchByGR} change
    if Gauss's Method were changed to allow multiplying a row by zero?
    \begin{answer}
      Instead of the row spaces being the same, the row space of $B$
      would be a subspace (possibly equal to) the row space of $A$.
    \end{answer}
  \item
    What is the relationship between \( \rank(A) \) and \( \rank(-A) \)?
    Between \( \rank(A) \) and \( \rank(kA) \)?
    What, if any, is the relationship between \( \rank(A) \), \( \rank(B) \),
    and \( \rank(A+B) \)?
    \begin{answer}
      Clearly \( \rank(A)=\rank(-A) \) as Gauss's Method allows us to
      multiply all rows of a matrix by \( -1 \).
      In the same way, when \( k\neq 0 \) we have \( \rank(A)=\rank(kA) \).

      Addition is more interesting.
      The rank of a sum can be smaller than the rank of the summands.
      \begin{equation*}
        \begin{mat}[r]
          1  &2  \\
          3  &4
        \end{mat}
        +
        \begin{mat}[r]
         -1  &-2  \\
         -3  &-4
        \end{mat}
        =
        \begin{mat}[r]
         0   &0   \\
         0   &0
        \end{mat}
      \end{equation*}
      The rank of a sum can be bigger than the rank of the summands.
      \begin{equation*}
        \begin{mat}[r]
          1  &2  \\
          0  &0
        \end{mat}
        +
        \begin{mat}[r]
          0  &0   \\
          3  &4
        \end{mat}
        =
        \begin{mat}[r]
         1   &2   \\
         3   &4
        \end{mat}
      \end{equation*}
      But there is an upper bound (other than the size of the matrices).
      In general, \( \rank(A+B)\leq\rank(A)+\rank(B) \).

      To prove this, note that we can perform Gaussian elimination on
      \( A+B \) in either of two ways:
      we can first add \( A \) to \( B \) and then apply the appropriate
      sequence of reduction steps
      \begin{equation*}
        (A+B)
        \grstep{\text{step}_1}
        \cdots
        \grstep{\text{step}_k}
        \text{echelon form}
      \end{equation*}
      or we can get the same results by performing \( \text{step}_1 \) through
      \( \text{step}_k \) separately on \( A \) and \( B \), and then adding.
      The largest rank that we can end with in the second case is clearly 
      the sum of the ranks.
      (The matrices above give examples of both possibilities,  
       $\rank(A+B)<\rank(A)+\rank(B)$ and $\rank(A+B)=\rank(A)+\rank(B)$,
       happening.)
    \end{answer}
    % Relationship between rank(A), rank(B) and rank(cA+dB)?
\end{exercises}























\subsectionoptional{Combining Subspaces}
\index{direct sum|(}
\textit{This subsection is optional.
It is required only for the last sections of
Chapter Three and Chapter Five and for occasional exercises.
You can pass it over without loss of continuity.}

% This chapter opened with the definition of a vector space, and the middle
% consisted of a first analysis of the idea.
% This subsection closes the chapter by finishing the analysis, in the sense
% that `analysis' means ``method of determining the \ldots{} essential
% features of something by separating it into parts'' \cite{MacmillanDict}.

One way to understand something is to see how to build it from 
component parts.
For instance, we sometimes think of \( \Re^3 \) 
put together
from the \( x \)-axis, the \( y \)-axis, and \( z \)-axis.
In this subsection we will describe 
how to decompose a vector space into a combination of
some of its subspaces.
In developing this idea of subspace combination, we will keep the $\Re^3$
example in mind as a prototype.

Subspaces are subsets and sets combine via union.
But taking the combination operation for subspaces to be the simple set union
operation isn't what we want.
For instance, the union of the 
\( x \)-axis, the \( y \)-axis, and \( z \)-axis is not all of $\Re^3$.
In fact this union is not
a subspace because it is not closed under addition: this vector 
\begin{equation*}
  \colvec[r]{1 \\ 0 \\ 0}
  +
  \colvec[r]{0 \\ 1 \\ 0}
  +
  \colvec[r]{0 \\ 0 \\ 1}
  =
  \colvec[r]{1 \\ 1 \\ 1}
\end{equation*}
is in none of the three axes and hence is not in the union.
Therefore
to combine subspaces, 
in addition to the members of those subspaces, we must at least also
include all of their linear combinations.

\begin{definition}
Where \( W_1,\dots, W_k \) are subspaces of a vector space, their
\definend{sum}\index{sum!of subspaces}\index{subspace!sum}
is the span of their union
$W_1+W_2+\dots +W_k=\spanof{W_1\union W_2\union \cdots W_k}$.
\end{definition}

\noindent Writing `\( + \)'  
fits with the conventional practice of using this symbol for a
natural accumulation operation.

\begin{example} \label{ex:RThreeSumAxes}
Our $\Re^3$ prototype works with this.
Any vector $\vec{w}\in\Re^3$ is a linear combination
$c_1\vec{v}_1+c_2\vec{v}_2+c_3\vec{v}_3$ where $\vec{v}_1$ is a member of
the $x$-axis, etc., in this way
\begin{equation*}
    \colvec{w_1 \\ w_2 \\ w_3}
    =1\cdot\colvec{w_1 \\ 0 \\ 0}
    +1\cdot\colvec{0 \\ w_2 \\ 0}
    +1\cdot\colvec{0 \\ 0 \\ w_3}
\end{equation*}
and so $\text{$x$-axis}+\text{$y$-axis}+\text{$z$-axis}=\Re^3$.
\end{example}

\begin{example}
A sum of subspaces can be less than the entire space. 
Inside of  \( \polyspace_4 \), 
let \( L \) be the subspace of linear polynomials
$\set{a+bx\suchthat a,b\in\Re}$
and let \( C \) be the subspace of purely-cubic polynomials
\( \set{cx^3\suchthat c\in\Re}\).
Then \mbox{$L+C$} is not all of $\polyspace_4$.
Instead, 
\( \mbox{$L+C$}=\set{a+bx+cx^3\suchthat a,b,c\in\Re} \).
\end{example}

\begin{example} \label{exam:RThreeIsSumXYAndYZ}
A space can be described as a combination of subspaces in more than one way.
Besides the decomposition 
$\Re^3=\text{$x$-axis}+\text{$y$-axis}+\text{$z$-axis}$,
we can also write
$\Re^3=\text{$xy$-plane}+\text{$yz$-plane}$.
To check this, note that any $\vec{w}\in\Re^3$ can be written
as a linear combination of a member of the $xy$-plane and a member of the
$yz$-plane; here are two such combinations.  
\begin{equation*}
  \colvec{w_1 \\ w_2 \\ w_3}
  =1\cdot\colvec{w_1 \\ w_2 \\ 0}
   +1\cdot\colvec{0 \\ 0 \\ w_3}
  \qquad
  \colvec{w_1 \\ w_2 \\ w_3}
  =1\cdot\colvec{w_1 \\ w_2/2 \\ 0}
   +1\cdot\colvec{0 \\ w_2/2 \\ w_3}
\end{equation*}
\end{example}

The above definition gives one way in which we can think of a space 
as a combination of some of its parts.
However, the prior example shows that there is at least one interesting
property of our benchmark model that is not captured by
the definition of the sum of subspaces.
In the familiar decomposition of $\Re^3$,
we often speak of a vector's `$x$~part' or `$y$~part' or 
`$z$~part'.
That is, in our prototype each vector has a unique decomposition into
pieces from the parts making up the whole space.
But in the decomposition used in \nearbyexample{exam:RThreeIsSumXYAndYZ}, we
cannot refer to the ``$xy$~part'' of a vector\Dash these three sums
\begin{equation*}
  \colvec[r]{1 \\ 2 \\ 3}
  =\colvec[r]{1 \\ 2 \\ 0}
  +\colvec[r]{0 \\ 0 \\ 3}
  =\colvec[r]{1 \\ 0 \\ 0}
  +\colvec[r]{0 \\ 2 \\ 3} 
  =\colvec[r]{1 \\ 1 \\ 0}
  +\colvec[r]{0 \\ 1 \\ 3} 
\end{equation*}
all describe the vector as comprised of something from the first plane plus
something from the second plane, but the ``$xy$~part'' is different in each.

%What's happening here, what allows the vector's second entry 
%$2$ to be taken from the $xy$-plane or
%from the $yz$-plane or from a combination, is that the two planes overlap.
%Their intersection is the $y$-axis and so the middle component 
%can shift around.
%In terms of our `analysis', the problem with the sum of subspaces operation is
%that, while the expression
%$\Re^3=\text{$xy$-plane}+\text{$yz$-plane}$ describes the space as a
%combination of its parts, those parts are not separate.

%\begin{lemma}
%The intersection of subspaces is a subspace.\index{subspace!intersection}
%\end{lemma}

%\begin{proof}
%To show that it is a subspace, we need only show that it is nonempty and
%closed under linear combinations.
%It is nonempty because it contains the zero vector.
%For closure, consider a linear combination of vectors from the
%intersection. 
%We can view that as a linear combination of vectors from the first subspace,
%because each vector in the intersection is in the first subspace.
%As all of the vectors are in the first subspace, 
%their combination sums to a vector that is also in the first subspace.
%In the same way, the combination sums to a vector that is in each subspace, 
%and so is in the intersection.
%\end{proof}

%Therefore, an intersection of subspaces is never empty; the smallest that it
%can be is the trivial subspace.
%So the above reference to subspaces that ``overlap'' must be taken to mean
%that their intersection is larger than the trivial subspace, 
%not just that their intersection is nonempty.

That is, when we consider how $\Re^3$ is put together from the three axes
we might mean ``in such a way that every vector 
has at least one decomposition,'' which gives the definition above.
But if we take it to mean 
``in such a way that every vector has one and only one
decomposition'' then we need another condition on combinations.
To see what this condition is, recall that
vectors are uniquely represented in terms of a basis.
We can use this to break a space into a sum of subspaces
such that any vector in the space breaks uniquely into a sum of members of
those subspaces.

\begin{example} \label{exam:BenchDirSum}
Consider $\Re^3$ with its standard basis 
$\stdbasis_3=\sequence{\vec{e}_1,\vec{e}_2,\vec{e}_3}$.
The subspace with the basis $B_1=\sequence{\vec{e}_1}$ is the $x$-axis,
the subspace with the basis $B_2=\sequence{\vec{e}_2}$ is the $y$-axis, and
the subspace with the basis $B_3=\sequence{\vec{e}_3}$ is the $z$-axis.
The fact that any member of $\Re^3$ is expressible as a sum of vectors from
these subspaces
\begin{equation*}
  \colvec{x \\ y \\ z}
   =\colvec{x \\ 0 \\ 0}
    +\colvec{0 \\ y \\ 0}
    +\colvec{0 \\ 0 \\ z}
\end{equation*}
reflects the fact that $\stdbasis_3$ spans the space\Dash this 
equation
\begin{equation*}
  \colvec{x \\ y \\ z}
   =c_1\colvec[r]{1 \\ 0 \\ 0}
    +c_2\colvec[r]{0 \\ 1 \\ 0}
    +c_3\colvec[r]{0 \\ 0 \\ 1}
\end{equation*}
has a solution for any $x,y,z\in\Re$.
And the fact that each such expression is unique reflects that fact that
$\stdbasis_3$ is linearly independent, so any 
equation like the one above has a unique
solution. 
\end{example}

\begin{example}
We don't have to take the basis vectors one at a time, 
we can conglomerate them into larger sequences.
Consider again the space $\Re^3$ and the vectors from the standard basis $\stdbasis_3$.
The subspace with the basis $B_1=\sequence{\vec{e}_1,\vec{e}_3}$
is the $xz$-plane.
The subspace with the basis $B_2=\sequence{\vec{e}_2}$ is the $y$-axis.
As in the prior example, the fact that any member of the space is a sum of
members of the two subspaces in one and only one way
\begin{equation*}
  \colvec{x \\ y \\ z}
   =\colvec{x \\ 0 \\ z}
    +\colvec{0 \\ y \\ 0}
\end{equation*}
is a reflection of the fact that these vectors form a
basis\Dash this equation
\begin{equation*}
  \colvec{x \\ y \\ z}
   =(c_1\colvec[r]{1 \\ 0 \\ 0}
    +c_3\colvec[r]{0 \\ 0 \\ 1})
    +c_2\colvec[r]{0 \\ 1 \\ 0}
\end{equation*}
has one and only one solution for any $x,y,z\in\Re$.
\end{example}

% These examples illustrate 
% the natural way to decompose a space into a sum of subspaces so that
% each vector decomposes uniquely into a sum of vectors from the parts.

\begin{definition}
The \definend{concatenation}\index{concatenation of sequences}%
\index{sequence!concatenation} 
of the sequences
$
   B_1=\sequence{\vec{\beta}_{1,1},\dots,\vec{\beta}_{1,n_1}}$, 
  \ldots,
  $B_k=\sequence{\vec{\beta}_{k,1},\dots,\vec{\beta}_{k,n_k}}$ 
adjoins them into a single sequence.
\begin{equation*}
 \cat{\cat{\cat{B_1}{B_2}}{\cdots}}{B_k}=
   \sequence{\vec{\beta}_{1,1},\dots,\vec{\beta}_{1,n_1},
            \vec{\beta}_{2,1},\dots,\vec{\beta}_{k,n_k} } 
\end{equation*}
\end{definition}

\begin{lemma} \label{le:UniqDecIffBasisDec}
Let $V$ be a vector space that is the sum of some of its subspaces
$V=W_1+\dots+W_k$.
Let $B_1$, \ldots, $B_k$ be bases for these subspaces.
The following are equivalent.
\begin{tfae}
  \item The expression of any $\vec{v}\in V$ as a combination  
     $\vec{v}=\vec{w}_1+\dots+\vec{w}_k$ 
     with $\vec{w}_i\in W_i$ is unique.
  \item The concatenation $\cat{\cat{B_1}{\cdots}}{B_k}$ is a basis for $V$.
  \item Among 
    nonzero vectors from different $W_i$'s every linear relationship is 
    trivial.
% ; that is, 
%     any set of nonzero vectors $\set{\vec{w}_1,\dots,\vec{w}_k}$
%     with $\vec{w}_i\in W_i$ is linearly independent.
\end{tfae}
\end{lemma}

\begin{proof}
We will show that $\text{(1)}\implies\text{(2)}$, 
that $\text{(2)}\implies\text{(3)}$, 
and finally that $\text{(3)}\implies\text{(1)}$.
For these arguments, observe that we can pass from a combination of 
$\vec{w}$'s to a combination of $\vec{\beta}$'s
\begin{align*}
  d_1\vec{w}_1+\dots+d_k\vec{w}_k                                 
    &= d_1(c_{1,1}\vec{\beta}_{1,1}+\dots+c_{1,n_1}\vec{\beta}_{1,n_1})   \\
    &\qquad\hbox{}+\dots
        +d_k(c_{k,1}\vec{\beta}_{k,1}+\dots+c_{k,n_k}\vec{\beta}_{k,n_k}) \\
    &= d_1c_{1,1}\cdot\vec{\beta}_{1,1}
        +\dots
        +d_kc_{k,n_k}\cdot\vec{\beta}_{k,n_k}   
  \tag{$*$}
\end{align*}
and vice versa
(we can move from the bottom to the top by taking each $d_i$ to be $1$).

For $\text{(1)}\implies\text{(2)}$,
assume that all decompositions are unique.
We will show that $\cat{\cat{B_1}{\cdots}}{B_k}$ spans the
space and is linearly independent.
It spans the space because the assumption 
that $V=W_1+\dots+W_k$
means that every $\vec{v}$ can be expressed as 
$\vec{v}=\vec{w}_1+\dots+\vec{w}_k$, which translates by equation~($*$) to 
an expression of $\vec{v}$ as a linear combination of the $\vec{\beta}$'s
from the concatenation.
For linear independence, consider this linear relationship.
\begin{equation*}
  \zero=c_{1,1}\vec{\beta}_{1,1}+\dots+c_{k,n_k}\vec{\beta}_{k,n_k}
\end{equation*}
Regroup as in ($*$) (that is, move from bottom to top) to get
the decomposition $\zero=\vec{w}_1+\dots+\vec{w}_k$.
Because the zero vector obviously
has the decomposition $\zero=\zero+\dots+\zero$,
the assumption that decompositions are unique shows that each
$\vec{w}_i$ is the zero vector.
This means that
$c_{i,1}\vec{\beta}_{i,1}+\dots+c_{i,n_i}\vec{\beta}_{i,n_i}=\zero$, and
since each $B_i$ is a basis we have the desired conclusion that all of
the $c$'s are zero.

For $\text{(2)}\implies\text{(3)}$ assume that the concatenation of the
bases is a basis for the entire space.
Consider a linear relationship among nonzero vectors from different
$W_i$'s.
This might or might not involve a vector from $W_1$, or
one from $W_2$, etc., so we write it 
$\zero=\mbox{}\cdots+d_i\vec{w}_i+\cdots\mbox{}$.
As in equation~($*$) expand the vector.
\begin{align*}
  \zero   
   &= \mbox{}\dots
      +d_i(c_{i,1}\vec{\beta}_{i,1}+\dots+c_{i,n_i}\vec{\beta}_{i,n_i})
      +\cdots\mbox{}                                               \\        
   &= \mbox{}\dots
      +d_ic_{i,1}\cdot\vec{\beta}_{i,1}
      +\dots+
      d_ic_{i,n_i}\cdot\vec{\beta}_{i,n_i}
      +\cdots\mbox{}                                                      
\end{align*}
The linear independence of $\cat{\cat{B_1}{\cdots}}{B_k}$ gives that each
coefficient $d_ic_{i,j}$ is zero.
Since $\vec{w}_i$ is nonzero vector, at least one of the 
$c_{i,j}$'s is not zero, and thus $d_i$ is zero.
This holds for each $d_i$, and therefore the linear relationship is trivial.

Finally, for $\text{(3)}\implies\text{(1)}$,
assume that among nonzero vectors from different $W_i$'s any linear
relationship is trivial.
Consider two decompositions of a vector 
$\vec{v}=\mbox{}\cdots+\vec{w}_i+\cdots\mbox{}$ 
and $\vec{v}=\mbox{}\cdots+\vec{u}_j+\cdots\mbox{}$
where $\vec{w}_i\in W_i$ and $\vec{u}_j\in W_j$.
Subtract one from the other to get a linear relationship, something like
this (if there is no $\vec{u}_i$ or $\vec{w}_j$ then leave those out).
\begin{equation*}
  \zero
   =\mbox{}\cdots+(\vec{w}_i-\vec{u}_i)
     +\cdots
     +(\vec{w}_j-\vec{u}_j)+\cdots\mbox{}
\end{equation*}
The case assumption that statement~(3) holds implies that the terms each
equal the zero vector
 $\vec{w}_i-\vec{u}_i=\zero$.
Hence decompositions are unique.
\end{proof}

\begin{definition}
A collection of subspaces \( \set{W_1,\ldots, W_k} \) is
\definend{independent}\index{subspace!independence}
if no nonzero vector from any \( W_i \) is a linear combination of
vectors from the other subspaces \( W_1,\dots, W_{i-1},W_{i+1},\dots, W_k \).
\end{definition}

\begin{definition}
A vector space \( V \) is the
\definend{direct sum}\index{subspace!direct sum}\index{direct sum!definition}
(or \definend{internal direct sum}\index{internal direct sum}) 
of its subspaces \( W_1,\dots, W_k \) if
\( V=W_1+W_2+\dots +W_k \)
and the collection \( \set{W_1,\dots, W_k} \) is independent.
We write \( V=W_1\directsum W_2\directsum \cdots\directsum W_k \).
\end{definition}

\begin{example}
Our prototype works: 
$\Re^3=\text{$x$-axis}\directsum\text{$y$-axis}\directsum\text{$z$-axis}$. 
\end{example}

\begin{example}
The space of \( \nbyn{2} \) matrices is this direct sum.
\begin{equation*}
  \set{\begin{mat}
         a  &0  \\
         0  &d
       \end{mat}  \suchthat a,d\in\Re }
  \,\mbox{}\directsum\mbox{}\,
  \set{\begin{mat}
         0  &b  \\
         0  &0
       \end{mat}  \suchthat b\in\Re }
  \,\mbox{}\directsum\mbox{}\,
  \set{\begin{mat}
         0  &0  \\
         c  &0
       \end{mat}  \suchthat c\in\Re }
\end{equation*}
It is the direct sum of subspaces in many other ways as well; 
direct sum decompositions are not unique.
\end{example}

\begin{corollary} \label{cor:DirSumDimsAdd}
The dimension of a direct sum is the sum of the dimensions of its summands.
\end{corollary}

\begin{proof}
In \nearbylemma{le:UniqDecIffBasisDec},
the number of basis vectors in the concatenation equals the sum of
the number of vectors in the sub-bases.
\end{proof}

The special case of two subspaces is worth its own mention.

\begin{definition}
When a vector space is the direct sum of two of its subspaces then they are
\definend{complements}.\index{subspace!complementary}%
\index{complementary subspaces}
\end{definition}

\begin{lemma}
\label{le:DirectSumTwoSp}
A vector space \( V \) is the 
direct sum\index{direct sum!of two subspaces}
of two of its subspaces \( W_1 \) and \( W_2 \) if and only if it is the
sum of the two \( V=W_1+W_2 \) and their intersection is trivial
\( W_1\intersection W_2=\set{\zero\,} \).
\end{lemma}

\begin{proof}
Suppose first that $V=W_1\directsum W_2$.
By definition, $V$ is the sum of the two $V=W_1+ W_2$.
To show that their intersection is trivial 
let $\vec{v}$ be a vector from $W_1\intersection W_2$
and consider the equation $\vec{v}=\vec{v}$.
On that equation's left side is a member of $W_1$
and on the right is a member of $W_2$, which we can think of as 
a linear combination of members of $W_2$.
But the two spaces are independent 
so the only way that a member of $W_1$ can be a linear
combination of vectors from $W_2$ is if that member 
is the zero vector $\vec{v}=\zero$.

For the other direction, suppose that $V$ is the sum of two spaces
with a trivial intersection.
To show that $V$ is a direct sum of the two we need only show that
the spaces are independent\Dash that no nonzero member of the first is 
expressible as a linear combination of members of the second, and vice versa.
This holds because any relationship 
$\vec{w}_1=c_1\vec{w}_{2,1}+\dots+c_k\vec{w}_{2,k}$ (with $\vec{w}_1\in W_1$
and $\vec{w}_{2,j}\in W_2$ for all $j$) shows that the vector on the left is 
also in $W_2$, since the right side is a combination of members of $W_2$.
The intersection of these two spaces is trivial, so $\vec{w}_1=\zero$.
The same argument works for any $\vec{w}_2$.
\end{proof}

%\begin{corollary}
%\label{cor:CompsIffDimsAndIntTriv}
%Two subspaces \( W_1,W_2 \) are complements in \( V \) if and only if
%\( \dim(W_1)+\dim(W_2)=\dim(V) \) and
%\( W_1\intersection W_2=\set{\zero\,} \).
%\end{corollary}

%\begin{proof}
%This follows immediately from 
%\nearbycorollary{cor:DimSumEqDimSummandsMinDimInter}.
%\end{proof}

\begin{example}
In $\Re^2$ the \( x \)-axis and the \( y \)-axis are complements,
that is, $\Re^2=\text{$x$-axis}\directsum\text{$y$-axis}$.
This points out that subspace complement is slightly different 
than set complement;~the $x$ and~$y$ axes are not set complements 
because their intersection
is not the empty set. 

A space can have more than one pair of complementary subspaces;
another pair for~$\Re^2$ are the subspaces consisting of the 
lines \( y=x \) and \( y=2x \).
\end{example}

\begin{example}
In the space \( F=\set{a\cos\theta+b\sin\theta\suchthat a,b\in\Re} \),
the subspaces \( W_1=\set{a\cos\theta\suchthat a\in\Re} \) and
\( W_2=\set{b\sin\theta\suchthat b\in\Re} \) are complements.
The prior example noted that a space can be decomposed into more than
one pair of complements. 
In addition note that $F$ can has more than one pair of
complementary subspaces where the first in the pair is $W_1$\Dash another 
complement of \( W_1 \) is
\( W_3=\set{b\sin\theta+b\cos\theta \suchthat b\in\Re} \).
\end{example}

\begin{example}
In \( \Re^3 \), the \( xy \)-plane and the \( yz \)-planes are not
complements, which is the point of the discussion following
\nearbyexample{exam:RThreeIsSumXYAndYZ}.
One complement of the \( xy \)-plane is the \( z \)-axis.
% A complement of the \( yz \)-plane is the line through \( (1,1,1) \).
\end{example}

Here is a natural question that arises from \nearbylemma{le:DirectSumTwoSp}:~for
$k>2$ is
the simple sum $V=W_1+\dots+W_k$ also a direct sum if and only if
the intersection of the subspaces is trivial?

\begin{example}    \label{exam:DirSumThree}
If there are more than two subspaces then 
having a trivial intersection is not enough to 
guarantee unique decomposition (i.e., is not enough to ensure that the
spaces are independent). 
In \( \Re^3 \), let
\( W_1 \) be the \( x \)-axis, let \( W_2 \) be the \( y \)-axis, and let
$W_3$ be this.
\begin{equation*}
  W_3=\set{\colvec{q \\ q \\ r} \suchthat q,r\in\Re}
\end{equation*}
The check that \( \Re^3=W_1+W_2+W_3 \) is easy.
The intersection 
\( W_1\intersection W_2\intersection W_3 \) 
is trivial, but decompositions aren't unique.
\begin{equation*}
  \colvec{x \\ y \\ z}
  =\colvec{0 \\ 0 \\ 0}
  +\colvec{0 \\ y-x \\ 0}
  +\colvec{x \\ x \\ z}
  =\colvec{x-y \\ 0 \\ 0}
  +\colvec{0 \\ 0 \\ 0}
  +\colvec{y \\ y \\ z}
\end{equation*}
(This example also shows that this requirement is also not 
enough:~that all pairwise intersections of the subspaces be trivial.
See \nearbyexercise{exer:ThreeSubsPairwseNonTriv}.)
\end{example}

In this subsection we have seen two ways to regard a space as built up from
component parts.
Both are useful; in particular we will use the direct sum definition
at the end of the Chapter Five.

%In looking at how this overlap of subspaces interacts with the 
%`$+$' operation, we first consider the simplest case, the two-subspace case.
%To understand the relationship between the two subspaces $U$ and $W$, their
%sum $U+W$, and their intersection $U\intersection W$, we can ask about the
%bases.  
%Note that a basis for the intersection $U\intersection W$ is a linearly 
%independent subset of $U$ and also of $W$.
%Thus, we can think to expand it to make bases for the larger sets.

%\begin{lemma} \label{le:BasesSumTwoSubs}
%Let \( U \) and \( W \) be subspaces of a vector space.
%If the sequence \( \sequence{\vec{\beta}_1,\dots,\vec{\beta}_k} \) 
%is a basis for
%\( U\intersection W \), and
%the sequence \( \sequence{\vec{\mu}_1,\dots,\vec{\mu}_j,
%   \vec{\beta}_1,\dots,\vec{\beta}_k} \)
%is a basis for \( U \), and
%\( \sequence{\vec{\beta}_1,\dots,\vec{\beta}_k,
%   \vec{\omega}_1,\dots,\vec{\omega}_p} \)
%is a basis for \( W \), then
%\begin{equation*}
%   \sequence{\vec{\mu}_1,\dots,\vec{\mu}_j,
%             \vec{\beta}_1,\dots,\vec{\beta}_k,
%             \vec{\omega}_1,\dots,\vec{\omega}_p}
%\end{equation*}
%is a basis for \( U+W \).
%\end{lemma}

%\begin{proof}
%We write $B_{U\intersection W}$ for the basis for $U\intersection W$,
%we write $B_U$ for the basis for $U$, we write $B_W$ for the basis for $W$,
%and we write $B_{U+W}$ for the basis under consideration.

%To see that that $B_{U+W}$ spans $U+W$, observe that
%any vector $c\vec{u}+d\vec{w}$ from $U+W$ can be written as a linear
%combination of the vectors in $B_{U+W}$, simply by expressing $\vec{u}$ in 
%terms of $B_U$ and expressing $\vec{w}$ in terms of $B_W$.

%We finish by showing that $B_{U+W}$ is linearly independent. 
%Consider
%\begin{equation*}
%  c_1\vec{\mu}_1+\dots+c_{j+1}\vec{\beta}_1+\dots+c_{j+k+p}\vec{\omega}_p=
%    \zero
%\end{equation*}
%which can be rewritten in this way.
%\begin{equation*}
%  c_1\vec{\mu}_1+\dots+c_j\vec{\mu}_j=
%       -c_{j+1}\vec{\beta}_1-\dots-c_{j+k+p}\vec{\omega}_p
%\end{equation*}
%Note that the left side sums to a vector in \( U \) while 
%right side sums to a vector in \( W \), and thus both sides sum to
%a member of $U\intersection W$.
%Since the left side is a member of $U\intersection W$, 
%it is expressible in terms of the members of $B_{U\intersection W}$, 
%which gives the combination of $\vec{\mu}$'s from the left side above
%as equal to a combination of $\vec{\beta}$'s. 
%But, the fact that 
%the basis $B_U$ is linearly independent shows that any such combination
%is trivial, and in particular, the coefficients $c_1$, \ldots,
%$c_j$ fro the left side above are all zero.
%Similarly, the coefficients of the \( \vec{\omega} \)'s are all zero.
%This leaves the above equation as a linear relationship among the 
%\( \vec{\beta} \)'s, 
%but $B_{U\intersection W}$ is linearly independent, 
%and therefore all of the coefficients of the $\vec{\beta}$'s are also zero.
%\end{proof}

%\begin{example}
%Label the axes in \( \Re^4 \) with \( x \), \( y \), \( z \), and \( w \).
%If \( U \) is \( xyz \)-space
%(that is, the subspace spanned by the \( x \), \( y \), and
%\( z \) axes), and \( W \) is \( xyw \)-space, then this
%\begin{equation*}
%  B_{U\intersection W}
%     =\sequence{\colvec{1 \\ 0 \\ 0 \\ 0},
%                \colvec{0 \\ 2 \\ 0 \\ 0} }
%\end{equation*}
%is a basis for \( U\intersection W \).
%It can be extended to these bases
%\begin{equation*}
%  B_U=\sequence{\colvec{1 \\ 1 \\ 1 \\ 0},
%            \colvec{1 \\ 0 \\ 0 \\ 0},
%            \colvec{0 \\ 2 \\ 0 \\ 0} }
%  \quad\text{and}\quad
%  B_W=\sequence{\colvec{1 \\ 0 \\ 0 \\ 0},
%            \colvec{0 \\ 2 \\ 0 \\ 0},
%            \colvec{1 \\ 0 \\ 0 \\ -1} }
%\end{equation*}
%for \( U \) and \( W \).
%Then this 
%\begin{equation*}
%  B_{U+W}
%   =\sequence{\colvec{1 \\ 1 \\ 1 \\ 0},
%            \colvec{1 \\ 0 \\ 0 \\ 0},
%            \colvec{0 \\ 2 \\ 0 \\ 0},
%            \colvec{1 \\ 0 \\ 0 \\ -1} }
%\end{equation*}
%is a basis for \( U+W=\Re^4 \).
%\end{example}

%\begin{corollary}
%\label{cor:DimSumEqDimSummandsMinDimInter}
%For any subspaces \( U \) and \( W \) of a vector space,
%\begin{equation*}
%  \dim(U+W)=\dim(U)+\dim(W)-\dim(U\intersection W).
%\end{equation*}
%\end{corollary}

%\begin{proof}
%Using the notation from the lemma, $\dim (U\intersection W)=k$,
%and $\dim(U)=j+k$, and $\dim(W)=k+p$, and $\dim (U+W)=j+k+p$. 
%\end{proof}

%\begin{example}
%In the prior example, $\dim (U\intersection W)=2$,
%and $\dim(U)=3$, and $\dim (W)=3$, and $\dim (U+W)=\dim(\Re^4)=4$,
%and so the equation becomes $4=3+3-2$.     
%\end{example}

%\begin{example}
%In \nearbyexample{exam:RThreeIsSumXYAndYZ} the equation gives
%$\dim(\Re^3)=\dim(\text{$xy$-plane})+\dim(\text{$yz$-plane})
%    -\dim(\text{$y$-axis})$ which is, of course, simply $3=2+2-1$.
%\end{example}

%The lemma gives us a way to ensure that every vector $\vec{v}\in U+W$ has a 
%unique decomposition as the sum of a vector from $U$ and a vector from $W$.
%Recall that vectors are uniquely represented in terms of a basis.
%If $B_{U\intersection W}$ is the empty basis (that is, if the intersection of
%$U$ and $W$ is trivial) then in the representation
%$\vec{v}=c_1\vec{\mu}_1+\dots+c_{j+p}\vec{\omega}_{j+p}$ we can take
%$\vec{u}$ to be $c_1\vec{\mu}_1+\dots+c_{j}\vec{\mu}_{j}$
%and take $\vec{w}$ to be 
%$c_{j+1}\vec{\omega}_{j+1}+\dots+c_{j+p}\vec{\omega}_{j+p}$
%to have a unique $\vec{v}=\vec{u}+\vec{w}$.

%\begin{definition}
%\label{def:DirectSumTwoSp}
%A vector space \( V \) is the 
%\definend{direct sum}\index{subspace!direct sum}%
%\index{direct sum!of two subspaces}
%of two of its subspaces \( W_1 \) and \( W_2 \), written 
%$V=W_1\directsum W_2$, if
%\( W_1+W_2=V \) and~\( W_1\intersection W_2=\set{\zero\,} \).
%The two subspaces are
%\definend{complements} %
%\index{subspace!complementary}\index{complementary subspaces}
%in \( V \).
%\end{definition}

%\begin{corollary}
%\label{cor:CompsIffDimsAndIntTriv}
%Two subspaces \( W_1,W_2 \) are complements in \( V \) if and only if
%\( \dim(W_1)+\dim(W_2)=\dim(V) \) and
%\( W_1\intersection W_2=\set{\zero\,} \).
%\end{corollary}

%\begin{proof}
%This follows immediately from 
%\nearbycorollary{cor:DimSumEqDimSummandsMinDimInter}.
%\end{proof}

%\begin{example}
%The \( x \)-axis and the \( y \)-axis are complements in \( \Re^2 \),
%that is, $\Re^2=\text{$x$-axis}\directsum\text{$y$-axis}$.
%A space can have more than one pair of complementary subspaces;
%another pair here are the subspaces consisting of the 
%lines \( y=x \) and \( y=2x \).
%\end{example}

%\begin{example}
%In the space \( F=\set{a\cos\theta+b\sin\theta\suchthat a,b\in\Re} \),
%the subspaces \( W_1=\set{a\cos\theta\suchthat a\in\Re} \) and
%\( W_2=\set{b\sin\theta\suchthat b\in\Re} \) are complements.
%In addition to the fact that a space like $F$ can have more than one pair of
%complementary subspaces, inside of the space a single subspace like $W_1$ can
%have more than one subspace that is a complement to it---another 
%complement to \( W_1 \) is
%\( W_3=\set{b\sin\theta+c\cos\theta \suchthat b,c\in\Re} \).
%\end{example}

%\begin{example}
%In \( \Re^3 \), the \( xy \)-plane and the \( yz \)-planes are not
%complements, because they have a nontrivial intersection.
%One complement of the \( xy \)-plane is the \( z \)-axis.
%A complement of the \( yz \)-plane is the line through \( (1,1,1) \).
%\end{example}

%\begin{lemma}
%For a space $V$ and two subspaces $W_1$, $W_2$, every vector $\vec{v}$ in the
%space has a decomposition $\vec{v}=\vec{u}+\vec{w}$ that is unique (where
%$\vec{u}\in U$ and $\vec{w}\in W$) if and only if the space is the direct sum
%of the subspaces.
%\end{lemma}

%\begin{proof}
%The `if' direction is handled as discussed before the definition.
%If $V=W_1\directsum W_2$ then the intersection of the two spaces is trivial
%and so \nearbylemma{le:BasesSumTwoSubs} shows that there is a basis for
%the sum $W_1+W_2=V$ consisting of a basis for $W_1$ with a basis for $W_2$
%adjoined.
%The representation of any $\vec{v}\in V$ with respect to this basis is unique,
%and we can regroup inside of this representation to have a unique expression
%$\vec{v}=\vec{w}_1+\vec{w}_2$.

%For the `only if' direction, we assume that decompositions exist and are
%unique.
%The fact that for every $\vec{v}\in V$ a decomposition exists shows that
%the space is the sum of the subspaces.
%To show that the intersection of the spaces is trivial,
%consider a member of the intersection \ldots
  
%\end{proof}

%We finish this subsection by going to the case of more than two subspaces.
%This is a bit trickier that the two-subspace case; the right notion of
%``do not overlap'' is perhaps not evident.
%That is, we are looking for a definition of a space as a direct sum of its
%subspaces where the first condition is that the space be the sum of those 
%subspaces, and the second condition is strong enough to give that each vector
%in the space has a unique decomposition into parts, one from each subspace.

%A natural first guess at the second condition, based on the material above,
%is to require that the spaces have a trivial intersection.
%The next example shows that this guess is, however, not enough to guarantee
%unique decomposition. 

%\begin{example}    \label{exam:DirSumThree}
%In \( \Re^3 \), let
%\( W_1 \) be the \( x \)-axis, let \( W_2 \) be the \( y \)-axis, and let
%\begin{equation*}
%  W_3=\set{\colvec{q \\ q \\ r} \suchthat q,r\in\Re}.
%\end{equation*}
%The check that \( \Re^3=W_1+W_2+W_3 \) is easy.
%The intersection of all three is trivial
%\( W_1\intersection W_2\intersection W_3=\set{\zero} \) but
%decompositions aren't unique.
%\begin{equation*}
%  \colvec{x \\ y \\ z}
%  =\colvec{0 \\ 0 \\ 0}
%  +\colvec{0 \\ y-x \\ 0}
%  +\colvec{x \\ x \\ z}
%  =\colvec{x-y \\ 0 \\ 0}
%  +\colvec{0 \\ 0 \\ 0}
%  +\colvec{y \\ y \\ z}
%\end{equation*}
%(This example also shows that this requirement is also not 
%enough:~that all pairwise intersections of the subspaces be trivial.
%See \nearbyexercise{exer:ThreeSubsPairwseNonTriv}.)
%\end{example}

%To ensure that each vector in the sum has a unique decomposition,
%we need a still stronger condition.

%\begin{definition}
%A collection of subspaces \( \set{W_1,\ldots, W_k} \) is
%{\em independent\/}\index{subspace!independence}
%if no nonzero vector from any \( W_i \) is a linear combination of
%vectors from the other subspaces \( W_1,\dots, W_{i-1},W_{i+1},\dots, W_k \).
%\end{definition}

%\begin{definition}
%A vector space \( V \) is the
%{\em direct sum\/}\index{subspace!direct sum}\index{direct sum!definition}
%(or {\em internal direct sum\/}) of its subspaces \( W_1,\dots, W_k \) if
%\( W_1+W_2+\dots +W_k=V \)
%and the collection \( \set{W_1,\dots, W_k} \) is independent.
%We write \( V=W_1\directsum W_2\directsum \dots\directsum W_k \).
%\end{definition}

%\begin{example}
%The benchmark example fits under this definition:
%$\Re^3=\text{$x$-axis}\directsum\text{$y$-axis}\directsum\text{$z$-axis}$ 
%\end{example}

%\begin{example}
%The space of \( \nbyn{2} \) matrices is this direct sum
%\begin{equation*}
%  \set{\begin{mat}
%         a  &0  \\
%         0  &d
%       \end{mat}  \suchthat a,d\in\Re }
%  \,\mbox{}\directsum\mbox{}\,
%  \set{\begin{mat}
%         0  &b  \\
%         0  &0
%       \end{mat}  \suchthat b\in\Re }
%  \,\mbox{}\directsum\mbox{}\,
%  \set{\begin{mat}
%         0  &0  \\
%         c  &0
%       \end{mat}  \suchthat c\in\Re }
%\end{equation*}
%(it is the direct sum of subspaces in many other ways also; as with the case
%of only two subspaces, direct sum decompositions into three or more 
%subspaces need not be unique).
%\end{example}

%Of course, following the work in this subsection, that definition will only be
%justified if we support it with a result saying that under those conditions,
%each vector in the summation has a unique decomposition into a sum of vectors,
%one from each summand.
%For that, 
%we will relate the decomposion a vector space into a direct sum
%to a decomposition of some basis for that space.
%We already know that uniqueness of decomposition relates to bases.

%\begin{definition}
%The sequence {\em concatenation\/}\index{concatenation} of
%\( B_1=\sequence{\vec{\beta}_{1,1},\dots,\vec{\beta}_{1,n_1}} \),
%\dots, \( B_k=\sequence{\vec{\beta}_{k,1},\dots,\vec{\beta}_{k,n_k}} \) is
%\(
% \cat{\cat{B_1}{\cdots}}{B_k}=
%  \sequence{\vec{\beta}_{1,1},\dots,\vec{\beta}_{1,n_1},
%            \vec{\beta}_{2,1},\dots,\vec{\beta}_{k,n_k} }.
%\)
%\end{definition}

%\begin{corollary}
%\label{cor:DirSumIffBasesCat}
%If \( V \) has subspaces \( W_1,\dots, W_k \)
%with bases \( B_1,\dots, B_k \)
%then \( W_1\directsum\dots\directsum W_k=V \) if and only if
%\( \cat{\cat{B_1}{\cdots}}{B_k} \) is a basis for \( V \).
%\end{corollary}

%\begin{proof}
%We will show that 
%the sum \( W_1+\dots+W_k \) equals \( V \) 
%if and only if \( \cat{\cat{B_1}{\cdots}}{B_k} \) spans \( V \), and that
%the set \( \set{W_1,\dots, W_k} \) is independent if and only if
%\( \cat{\cat{B_1}{\cdots}}{B_k} \) is linearly independent.
%These two halves prove the `if and only if' equivalence.

%First~(1).
%Left to right:
%assume the subspaces sum to \( V \) so any vector in \( V \) can be represented
%as \( \vec{v}=\vec{w}_1+\dots+\vec{w}_k \),
%rewrite each \( \vec{w}_i \) as a sum of basis elements and
%conclude \( \cat{\cat{B_1}{\cdots}}{B_k} \) spans \( V \).
%Right to left: if \( \cat{\cat{B_1}{\cdots}}{B_k} \) spans \( V \)
%then any \( \vec{v}\in V \) is a linear combination of vectors from
%\( \cat{\cat{B_1}{\cdots}}{B_k} \).
%Inside that combination commute to put vectors from \( B_1 \) first,
%followed from vectors from \( B_2 \), etc.
%Sum to first part to some \( \vec{w}_1\in W_1\), sum the second part to
%a \( \vec{w}_2\in W_2 \), etc.
%The result has the desired form.

%Now we prove item~(2).
%`If' is easy.
%For `only if' suppose the set of subspaces is independent and consider a linear
%relationship among vectors in
%\( \cat{\cat{B_1}{\cdots}}{B_k} \).
%Rewrite that relationship so all the vectors from \( B_1 \) are on one side and
%the other vectors are on the other side.
%This expresses some member of \( W_1 \) as a linear combination of members of
%the other subspaces, and independence of the set of subspaces then implies
%this linear combination of members of \( B_1 \) is \( \zero \).
%Thus, in this relationship, all the scalars associated with members of
%\( B_1 \) are \( 0 \).
%Similarly every scalar in the combination is \( 0 \).
%\end{proof}

%Hence, as mentioned above, each direct sum decomposition of a vector space is
%associated with a decomposition of one of that space's bases.

%\begin{corollary}
%\label{cor:DirectSumIffUniqueSum}
%Where \( W_1,\dots, W_k \) are subspaces,
%\( V=W_1\directsum\dots\directsum W_k \) if and
%only if every vector \( \vec{v}\in V \) has one and only one representation
%\( \vec{v}=\vec{w}_1+\dots+\vec{w}_k \) where \( \vec{w}_i\in W_i \).
%\end{corollary}

%\begin{proof}
%A chain of double implications works:
%\( V=W_1\directsum\dots\directsum W_k \)
%if and only if
%\( \cat{\cat{B_1}{\cdots}}{B_k} \) is a basis for \( V \),
%which is true if and only if each vector in \( V \)
%has a unique representation as a sum of vectors from
%\( \cat{\cat{B_1}{\cdots}}{B_k} \).
%\end{proof}

%\begin{corollary}
%The dimension of a direct sum is the sum of the dimensions of its summands.
%\end{corollary}


\begin{exercises}
  \recommended \item
    Decide if \( \Re^2 \) is the direct sum of each \( W_1 \) and \( W_2 \).
    \begin{exparts}
       \partsitem \( W_1=\set{\colvec{x \\ 0}\suchthat x\in\Re} \),
             \( W_2=\set{\colvec{x \\ x}\suchthat x\in\Re} \)
       \partsitem \( W_1=\set{\colvec{s \\ s}\suchthat s\in\Re} \),
             \( W_2=\set{\colvec{s \\ 1.1s}\suchthat s\in\Re} \)
       \partsitem \( W_1=\Re^2 \), \( W_2=\set{\zero} \)
       \partsitem \( W_1=W_2=\set{\colvec{t \\ t}\suchthat t\in\Re} \)
       \partsitem \( W_1=\set{\colvec[r]{1 \\ 0}+\colvec{x \\ 0}
                                \suchthat x\in\Re} \),
             \( W_2=\set{\colvec[r]{-1 \\ 0}+\colvec[r]{0 \\ y}\suchthat y\in\Re} \)
    \end{exparts}
    \begin{answer}
       With each of these we can apply \nearbylemma{le:DirectSumTwoSp}.
       \begin{exparts}
         \partsitem Yes.
           The plane is the sum of this $W_1$ and $W_2$ because for any 
           scalars $a$ and $b$
           \begin{equation*}
             \colvec{a \\ b}
             =\colvec{a-b \\ 0}
             +\colvec{b \\ b}
           \end{equation*}
           shows that the general vector is a sum of vectors from the two
           parts.
           And, these two subspaces are (different) lines through the origin,
           and so have a trivial intersection.
         \partsitem Yes.
           To see that any vector in the plane is a combination of vectors
           from these parts, consider this relationship.
           \begin{equation*}
             \colvec{a \\ b}=c_1\colvec[r]{1 \\ 1}+c_2\colvec[r]{1 \\ 1.1}
           \end{equation*}
           We could now simply note that the set
           \begin{equation*}
             \set{\colvec[r]{1 \\ 1},\colvec[r]{1 \\ 1.1}}
           \end{equation*}
           is a basis for the space (because it is clearly linearly
           independent, and has size two in $\Re^2$), and thus there is one and
           only one solution to the above equation, implying that all
           decompositions are unique.
           Alternatively, we can solve
           \begin{equation*}
             \begin{linsys}{2}
               c_1  &+  &c_2    &=  &a  \\
               c_1  &+  &1.1c_2 &=  &b
             \end{linsys}
             \grstep{-\rho_1+\rho_2}
             \begin{linsys}{2}
               c_1  &+  &c_2    &=  &a  \\
                    &   &0.1c_2 &=  &-a+b
             \end{linsys}
           \end{equation*}
           to get that $c_2=10(-a+b)$ and $c_1=11a-10b$, and so we have
           \begin{equation*}
             \colvec{a \\ b}=\colvec{11a-10b \\ 11a-10b}
                             +\colvec{-10a+10b \\ 1.1\cdot(-10a+10b)}
           \end{equation*}
           as required.
           As with the prior answer, each of the two subspaces is a line
           through the origin, and their intersection is trivial.
         \partsitem Yes.
           Each vector in the plane is a sum in this way
           \begin{equation*}
             \colvec{x \\ y}=\colvec{x \\ y}+\colvec[r]{0 \\ 0}
           \end{equation*}
           and the intersection of the two subspaces is trivial.
         \partsitem No.
           The intersection is not trivial.
         \partsitem No. 
           These are not subspaces.
      \end{exparts}  
    \end{answer}
  \recommended \item
    Show that \( \Re^3 \) is the direct sum of the \( xy \)-plane with    
    each of these. 
    \begin{exparts}
      \partsitem the \( z \)-axis
      \partsitem the line
        \begin{equation*}
          \set{\colvec{z \\ z \\ z} \suchthat z\in\Re }
        \end{equation*}
    \end{exparts}
    \begin{answer}
      With each of these we can use \nearbylemma{le:DirectSumTwoSp}.      
      \begin{exparts}
        \partsitem Any vector in \( \Re^3 \) can be decomposed as this sum.
          \begin{equation*}
            \colvec{x \\ y \\ z}
            =\colvec{x \\ y \\ 0}
            +\colvec{0 \\ 0 \\ z}
          \end{equation*}
          And, the intersection of the $xy$-plane and the $z$-axis is the
          trivial subspace.
        \partsitem Any vector in \( \Re^3 \) can be decomposed as
          \begin{equation*}
            \colvec{x \\ y \\ z}
            =\colvec{x-z \\ y-z \\ 0}
            +\colvec{z \\ z \\ z}
          \end{equation*}
          and the intersection of the two spaces is trivial.
      \end{exparts}  
    \end{answer}
  \item 
    Is \( \polyspace_2 \) the direct sum of
    \( \set{a+bx^2\suchthat a,b\in\Re} \) and
    \( \set{cx\suchthat c\in\Re} \)?
    \begin{answer}
      It is.
      Showing that these two are subspaces is routine.
      To see that the space is the direct sum of these two, just note that
      each member of \( \polyspace_2 \) has the unique decomposition
      \( m+nx+px^2=(m+px^2)+(nx) \).  
     \end{answer}
  \recommended \item 
    In \( \polyspace_n \), the 
    \definend{even}\index{even functions} polynomials are the members
    of this set
    \begin{equation*}
      \mathcal{E}=
      \set{p\in\polyspace_n \suchthat \text{\( p(-x)=p(x) \) for all \( x \)}}
    \end{equation*}
    and the 
    \definend{odd}\index{odd function}
    polynomials are the members of this set.
    \begin{equation*}
      \mathcal{O}=
      \set{p\in\polyspace_n \suchthat \text{\( p(-x)=-p(x) \) 
                                              for all \( x \)}}
    \end{equation*}
    Show that these are complementary subspaces.
    \begin{answer}
      To show that they are subspaces is routine.
      We will argue they are complements with \nearbylemma{le:DirectSumTwoSp}.
      The intersection \( \mathcal{E}\intersection\mathcal{O} \) is trivial
      because the only polynomial satisfying both conditions $p(-x)=p(x)$ and
      $p(-x)=-p(x)$ is the zero polynomial.
      To see that the entire space is the sum of the subspaces 
      \( \mathcal{E}+\mathcal{O}=\polyspace_n \), 
      note that the polynomials
      \( p_0(x)=1 \), \( p_2(x)=x^2 \), \( p_4(x)=x^4 \), etc., are 
      in \( \mathcal{E} \) and also note that the polynomials
      \( p_1(x)=x \), \( p_3(x)=x^3 \), etc., are in \( \mathcal{O} \).
      Hence any member of \( \polyspace_n \) is a combination of members of
      \( \mathcal{E} \) and \( \mathcal{O} \).  
    \end{answer}
  \item 
    Which of these subspaces of \( \Re^3 \) 
    \begin{center}
     \begin{tabular}{l}
      $W_1$: the \( x \)-axis,\quad
      $W_2$:~the \( y \)-axis,\quad
      $W_3$:~the \( z \)-axis,             \\
      $W_4$:~the plane \( x+y+z=0 \),\quad
      $W_5$:~the \( yz \)-plane
     \end{tabular}
    \end{center}
    can be combined to
    \begin{exparts*}
      \partsitem sum to \( \Re^3 \)? 
      \partsitem direct sum to \( \Re^3 \)?
    \end{exparts*}
    \begin{answer}
     Each of these is $\Re^3$.
     \begin{exparts}
      \partsitem 
        These are broken into some separate lines for readability.
        \begin{center}
          \begin{tabular}{l} 
            $W_1+W_2+W_3$, 
              $W_1+W_2+W_3+W_4$, 
              $W_1+W_2+W_3+W_5$,  \\ \quad %line too long
              $W_1+W_2+W_3+W_4+W_5$, 
              $W_1+W_2+W_4$, 
              $W_1+W_2+W_4+W_5$,     \\ \quad %line too long
              $W_1+W_2+W_5$,         
              $W_1+W_3+W_4$, 
              $W_1+W_3+W_5$, 
              $W_1+W_3+W_4+W_5$,     \\ \quad 
              $W_1+W_4$, 
              $W_1+W_4+W_5$,         
              $W_1+W_5$,             \\ 
              $W_2+W_3+W_4$, 
              $W_2+W_3+W_4+W_5$,     
              $W_2+W_4$, 
              $W_2+W_4+W_5$,         \\ 
              $W_3+W_4$, 
              $W_3+W_4+W_5$,         \\ 
              $W_4+W_5$ 
          \end{tabular}
         \end{center}
     \partsitem
       $W_1\directsum W_2\directsum W_3$,    
       $W_1\directsum W_4$,    
       $W_1\directsum W_5$,    
       $W_2\directsum W_4$,    
       $W_3\directsum W_4$    
     \end{exparts}
    \end{answer}
  \recommended \item
    Show that \( \polyspace_n=\set{a_0 \suchthat a_0\in\Re}\directsum\dots
                              \directsum\set{a_nx^n\suchthat a_n\in\Re} \).
    \begin{answer}
      Clearly each is a subspace.
      The bases \( B_i=\sequence{x^i} \) for the subspaces, when concatenated,
      form a basis for the whole space.  
    \end{answer}
  \item 
    What is \( W_1+W_2 \) if \( W_1\subseteq W_2 \)?
    \begin{answer}
       It is \( W_2 \).  
    \end{answer}
  \item 
    Does \nearbyexample{exam:BenchDirSum} generalize?
    That is, is this true or false:~if a vector space \( V \) has a basis
    \( \sequence{\vec{\beta}_1,\dots ,\vec{\beta}_n} \) then
    it is the direct sum of the spans of the one-dimensional subspaces
    \( V=\spanof{\set{\vec{\beta}_1}}\directsum\dots
           \directsum\spanof{\set{\vec{\beta}_n}} \)?
    \begin{answer}
      True by \nearbylemma{le:UniqDecIffBasisDec}.  
    \end{answer}
  \item 
    Can \( \Re^4 \) be decomposed as a direct sum in two different ways?
    Can \( \Re^1 \)?
    \begin{answer}
      Two distinct direct sum decompositions of \( \Re^4 \) are easy
      to find.
      Two such are
      \( W_1=\spanof{\set{\vec{e}_1,\vec{e}_2}} \) and
      \( W_2=\spanof{\set{\vec{e}_3,\vec{e}_4}} \), and also
      \( U_1=\spanof{\set{\vec{e}_1}} \) and
      \( U_2=\spanof{\set{\vec{e}_2,\vec{e}_3,\vec{e}_4}} \).
      (Many more are possible, 
      for example \( \Re^4 \) and its trivial subspace.)

      In contrast, any partition of \( \Re^1 \)'s single-vector basis will
      give one basis with no elements and another with a single element.
      Thus any decomposition involves \( \Re^1 \) and its trivial
      subspace. 
    \end{answer}
  \item 
     This exercise makes the notation of writing `$+$' between sets more
     natural. 
     Prove that, where \( W_1,\dots, W_k \) are subspaces of a vector space,
     \begin{equation*}
       W_1+\dots+W_k
       =\set{\vec{w}_1+\vec{w}_2+\dots+\vec{w}_k
             \suchthat \vec{w}_1\in W_1,\dots,\vec{w}_k\in W_k},
     \end{equation*}
     and so the sum of subspaces is the subspace of all sums.
     \begin{answer}
       Set inclusion one way is easy:
       \( \set{\vec{w}_1+\dots+\vec{w}_k\suchthat \vec{w}_i\in W_i} \)
       is a subset of \( \spanof{W_1\union \dots \union W_k}  \)
       because each \( \vec{w}_1+\dots+\vec{w}_k \) is a sum of vectors from 
       the union.

       For the other inclusion, to any linear combination of vectors from
       the union apply commutativity of vector addition to
       put vectors from \( W_1 \) first, followed by vectors from \( W_2 \), 
       etc.
       Add the vectors from \( W_1 \) to get a \( \vec{w}_1\in W_1 \),
       add the vectors from \( W_2 \) to get a \( \vec{w}_2\in W_2 \), etc.
       The result has the desired form.
     \end{answer}
  \item \label{exer:ThreeSubsPairwseNonTriv}
    (Refer to \nearbyexample{exam:DirSumThree}.
    This exercise 
    shows that the requirement that pairwise intersections be trivial
    is genuinely stronger than the requirement only that the intersection of
    all of the subspaces be trivial.)
    Give a vector space and three subspaces $W_1$, $W_2$, and $W_3$
    such that the space is the sum of the subspaces, 
    the intersection of all three subspaces
    $W_1\intersection W_2\intersection W_3$ is trivial, 
    but the pairwise intersections 
    $W_1\intersection W_2$, $W_1\intersection W_3$, and $W_2\intersection W_3$
    are nontrivial.
    \begin{answer}
      One example is to take the space to be $\Re^3$, and to take the
      subspaces to be the $xy$-plane, the $xz$-plane, and the $yz$-plane.
    \end{answer}
  \item 
    Prove that if \( V=W_1\directsum\dots\directsum W_k \) then
    \( W_i\intersection W_j \) is trivial whenever \( i\neq j \).
    This shows that the first half of the proof of 
    \nearbylemma{le:DirectSumTwoSp} extends to the case of more than two
    subspaces.
    (\nearbyexample{exam:DirSumThree} shows that this implication does
    not reverse; the other half does not extend.)
    \begin{answer}
      Of course, the zero vector is in all of the subspaces, so the 
      intersection contains at least that one vector..  
      By the definition of direct sum the set
      \( \set{W_1,\dots,W_k} \) is independent and so no nonzero vector
      of \( W_i \) is a multiple of a member of \( W_j \), when
      \( i\neq j \).
      In particular, no nonzero vector from \( W_i \) equals a member of
      \( W_j \).
    \end{answer}
  \item 
    Recall that no linearly independent set contains the zero vector.
    Can an independent set of subspaces contain the trivial subspace?
    \begin{answer}
      It can contain a trivial subspace;
      this set of subspaces of \( \Re^3 \) is independent:
      \( \set{\set{\zero},\text{\( x \)-axis}} \).
      No nonzero vector from the trivial space \( \set{\zero} \) 
      is a multiple of a
      vector from the \( x \)-axis, simply because the trivial space has
      no nonzero vectors to be candidates for such a multiple (and
      also no nonzero vector from the \( x \)-axis is a multiple of
      the zero vector from the trivial subspace).  
     \end{answer}
  \recommended \item 
    Does every subspace have a complement?
    \begin{answer}
      Yes.
      For any subspace of a vector space we can take any
      basis \( \sequence{\vec{\omega}_1,\dots,\vec{\omega}_k} \) for that
      subspace and extend it to a basis
      \( \sequence{\vec{\omega}_1,\dots,\vec{\omega}_k,
                     \vec{\beta}_{k+1},\dots,\vec{\beta}_n} \) for the whole
      space.
      Then the complement of the original subspace has this basis
      \( \sequence{\vec{\beta}_{k+1},\dots,\vec{\beta}_n} \).  
     \end{answer}
  \recommended \item 
    Let \( W_1, W_2 \) be subspaces of a vector space.
    \begin{exparts}
      \partsitem Assume that 
        the set \( S_1 \) spans \( W_1 \), and that the set \( S_2 \) spans 
        \( W_2 \).
        Can \( S_1\union S_2 \) span \( W_1+W_2 \)?
        Must it?
      \partsitem Assume that \( S_1 \) is a linearly independent 
        subset of \( W_1 \)
        and that \( S_2 \) is a linearly independent subset of \( W_2 \).
        Can \( S_1\union S_2 \) be a linearly independent subset of
        \( W_1+W_2 \)?
        Must it?
    \end{exparts}
    \begin{answer}
      \begin{exparts}
        \partsitem It must.
          We can write any member of \( W_1+W_2 \) as \( \vec{w}_1+\vec{w}_2 \)
          where \( \vec{w}_1\in W_1 \) and \( \vec{w}_2\in W_2 \).
          As \( S_1 \) spans \( W_1 \), the vector \( \vec{w}_1 \) is a
          combination of members of \( S_1 \).
          Similarly \( \vec{w}_2 \) is a combination of members of \( S_2 \).
        \partsitem An easy way to see that it can be linearly independent 
          is to take each to be the empty set.
          On the other hand, in the space $\Re^1$, 
          if \( W_1=\Re^1 \) and \( W_2=\Re^1 \) and $S_1=\set{1}$ and
          $S_2=\set{2}$, then their union $S_1\union S_2$ is not independent.
      \end{exparts}  
     \end{answer}
  \item \label{exer:BasesSumTwoSubs}
    When we decompose a vector space as a direct sum, the dimensions
    of the subspaces add to the dimension of the space.
    The situation with a space that is given as the sum of its subspaces 
    is not as simple.
    This exercise considers the two-subspace special case.
    \begin{exparts}
      \partsitem For these subspaces of \( \matspace_{\nbyn{2}} \) find
        \( W_1\intersection W_2 \), \( \dim(W_1\intersection W_2) \),
        \( W_1+W_2 \), and \( \dim(W_1+W_2) \).
        \begin{equation*}
          W_1=\set{\begin{mat}
                    0  &0  \\
                    c  &d
                  \end{mat} \suchthat c,d\in\Re  }
          \qquad
         W_2=\set{\begin{mat}
                    0  &b  \\
                    c  &0
                  \end{mat} \suchthat b,c\in\Re  }
       \end{equation*}
     \partsitem Suppose that \( U \) and \( W \) are subspaces 
       of a vector space.
       Suppose that the sequence 
       \( \sequence{\vec{\beta}_1,\dots,\vec{\beta}_k} \) 
       is a basis for
       \( U\intersection W \).
       Finally, suppose that the prior sequence has been expanded to give
       a sequence \( \sequence{\vec{\mu}_1,\dots,\vec{\mu}_j,
       \vec{\beta}_1,\dots,\vec{\beta}_k} \)
       that is a basis for \( U \), and a sequence
       \( \sequence{\vec{\beta}_1,\dots,\vec{\beta}_k,
          \vec{\omega}_1,\dots,\vec{\omega}_p} \)
       that is a basis for \( W \). 
       Prove that this sequence
       \begin{equation*}
         \sequence{\vec{\mu}_1,\dots,\vec{\mu}_j,
              \vec{\beta}_1,\dots,\vec{\beta}_k,
              \vec{\omega}_1,\dots,\vec{\omega}_p}
       \end{equation*}
       is a basis for the sum \( U+W \).
      \partsitem Conclude that 
          $\dim (U+W)=\dim(U)+\dim(W)-\dim(U\intersection W)$.
      \partsitem Let \( W_1 \) and \( W_2 \) be eight-dimensional 
        subspaces of a ten-di\-men\-sion\-al space.
       List all values possible for \( \dim(W_1\intersection W_2) \).
    \end{exparts}
    \begin{answer}
      \begin{exparts}
        \partsitem The intersection and sum are
          \begin{equation*}
             \set{\begin{mat}
                    0  &0  \\
                    c  &0
                  \end{mat} \suchthat c\in\Re  }
             \qquad
             \set{\begin{mat}
                    0  &b  \\
                    c  &d
                   \end{mat} \suchthat b,c,d\in\Re  }
          \end{equation*}
        which have dimensions one and three.
      \partsitem We write $B_{U\intersection W}$ for the basis for 
        $U\intersection W$,
        we write $B_U$ for the basis for $U$, 
        we write $B_W$ for the basis for $W$,
        and we write $B_{U+W}$ for the basis under consideration.

        To see that $B_{U+W}$ spans $U+W$, observe that we can write
        any vector $c\vec{u}+d\vec{w}$ from $U+W$ as a linear
        combination of the vectors in $B_{U+W}$, 
        simply by expressing $\vec{u}$ in 
        terms of $B_U$ and expressing $\vec{w}$ in terms of $B_W$.

         We finish by showing that $B_{U+W}$ is linearly independent. 
         Consider
         \begin{equation*}
           c_1\vec{\mu}_1+\dots+c_{j+1}\vec{\beta}_1+\dots
              +c_{j+k+p}\vec{\omega}_p=\zero
         \end{equation*}
         which can be rewritten in this way.
         \begin{equation*}
           c_1\vec{\mu}_1+\dots+c_j\vec{\mu}_j=
             -c_{j+1}\vec{\beta}_1-\dots-c_{j+k+p}\vec{\omega}_p
         \end{equation*}
         Note that the left side sums to a vector in \( U \) while 
         right side sums to a vector in \( W \), and thus both sides sum to
         a member of $U\intersection W$.
         Since the left side is a member of $U\intersection W$, 
         it is expressible in terms of the members of $B_{U\intersection W}$, 
         which gives the combination of $\vec{\mu}$'s from the left side above
         as equal to a combination of $\vec{\beta}$'s. 
         But, the fact that 
         the basis $B_U$ is linearly independent shows 
         that any such combination
         is trivial, and in particular, the coefficients $c_1$, \ldots,
         $c_j$ from the left side above are all zero.
         Similarly, the coefficients of the \( \vec{\omega} \)'s are all zero.
         This leaves the above equation as a linear relationship among the 
         \( \vec{\beta} \)'s, 
         but $B_{U\intersection W}$ is linearly independent, 
         and therefore all of the coefficients of the $\vec{\beta}$'s are also
         zero.
        \partsitem Just count the basis vectors in the prior 
          item:~$\dim(U+W)=j+k+p$, and $\dim(U)=j+k$, and $\dim(W)=k+p$, 
          and $\dim(U\intersection W)=k$.
        \partsitem  We know that
          \( \dim(W_1+W_2)=\dim(W_1)+\dim(W_2)-\dim(W_1\intersection W_2) \).
          Because \( W_1\subseteq W_1+W_2 \), 
          we know that \( W_1+W_2 \) must have dimension greater than that of
          $W_1$, that is, must have dimension eight, nine, or ten.
          Substituting gives us three possibilities
          $8=8+8-\dim(W_1\intersection W_2)$ or
          $9=8+8-\dim(W_1\intersection W_2)$ or
          $10=8+8-\dim(W_1\intersection W_2)$.
          Thus \( \dim(W_1\intersection W_2) \) must be either 
          eight, seven, or six.
          (Giving examples to show that 
          each of these three cases is possible
          is easy, for instance in \( \Re^{10} \).)  
      \end{exparts}
    \end{answer}
  \item 
    Let \( V=W_1\directsum\cdots\directsum W_k \) and for each index
    \( i \) suppose that \( S_i \) is a linearly independent subset of
    \( W_i \).
    Prove that the union of the \( S_i \)'s is linearly independent.
    \begin{answer}
      Expand each \( S_i \) to a basis $B_i$ for \( W_i \).
      The concatenation of those bases $\cat{\cat{B_1}{\cdots}}{B_k}$
      is a basis for \( V \) and thus its members form a linearly independent
      set.
      But the union $S_1\union\cdots\union S_k$ is a subset 
      of that linearly independent set, and thus is
      itself linearly independent.  
     \end{answer}
  \item 
    A matrix is \definend{symmetric}\index{matrix!symmetric}%
    \index{symmetric matrix}
    if for each pair of indices \( i \) and
    \( j \), the \( i,j \) entry equals the \( j,i \) entry.
    A matrix is \definend{antisymmetric}\index{matrix!antisymmetric}%
    \index{antisymmetric matrix} 
    if each \( i,j \) entry is the
    negative of the \( j,i \) entry.
    \begin{exparts}
      \partsitem Give a symmetric $\nbyn{2}$ matrix and an antisymmetric
        $\nbyn{2}$ matrix.
        (\textit{Remark.}
        For the second one, be careful about the entries on the diagonal.)
      \partsitem What is the relationship between a square symmetric matrix and
        its transpose?
        Between a square antisymmetric matrix and its transpose?
      \partsitem Show that \( \matspace_{\nbyn{n}} \) is the direct sum of 
        the space of symmetric matrices and the space of antisymmetric 
        matrices.
    \end{exparts}
    \begin{answer}
      \begin{exparts}
        \partsitem Two such are these.
          \begin{equation*}
            \begin{mat}[r]
              1  &2  \\
              2  &3
            \end{mat}
            \qquad
            \begin{mat}[r]
              0  &1  \\
              -1 &0
            \end{mat}
          \end{equation*}
          For the antisymmetric one, entries on the diagonal must be zero.
        \partsitem A square symmetric matrix equals its transpose.
          A square antisymmetric matrix equals the negative of its transpose.
        \partsitem Showing that the two sets are subspaces is easy.
         Suppose that \( A\in\matspace_{\nbyn{n}} \).
         To express $A$ as a sum of a symmetric and an antisymmetric matrix,
         we observe that
         \begin{equation*}
            A=(1/2)(A+\trans{A}) + (1/2)(A-\trans{A})
         \end{equation*}
         and note the first summand is symmetric while the second is 
         antisymmetric.
         Thus \( \matspace_{\nbyn{n}} \) is the sum of the two subspaces.
         To show that the sum is direct, 
         assume a matrix \( A \) is both symmetric
         \( A=\trans{A} \) and antisymmetric \( A=-\trans{A} \).
         Then \( A=-A \) and so all of \( A \)'s entries are zeroes.  
      \end{exparts}
     \end{answer}
  \item 
    Let \( W_1,W_2,W_3 \) be subspaces of a vector space.
    Prove that \( (W_1\intersection W_2)+(W_1\intersection W_3)\subseteq
              W_1\intersection (W_2+W_3) \).
    Does the inclusion reverse?
    \begin{answer}
      Assume that 
      \( \vec{v}\in (W_1\intersection W_2)+(W_1\intersection W_3) \).
      Then \( \vec{v}=\vec{w}_2+\vec{w}_3 \) where
      \( \vec{w}_2\in W_1\intersection W_2 \) and
      \( \vec{w}_3\in W_1\intersection W_3 \).
      Note that \( \vec{w}_2,\vec{w}_3\in W_1 \) and, as a subspace is closed
      under addition, \( \vec{w}_2+\vec{w}_3\in W_1 \).
      Thus \( \vec{v}=\vec{w}_2+\vec{w}_3\in W_1\intersection(W_2+W_3) \).

      This example proves that the inclusion may be strict:
      in \( \Re^2 \) take \( W_1 \) to be the \( x \)-axis, take \( W_2 \)
      to be the \( y \)-axis, and take \( W_3 \) to be the line \( y=x \).
      Then \( W_1\intersection W_2 \) and \( W_1\intersection W_3 \) are
      trivial and so their sum is trivial.
      But \( W_2+W_3 \) is all of \( \Re^2 \) so
      \( W_1\intersection(W_2+W_3) \) is the \( x \)-axis.  
     \end{answer}
  \item 
    The example of the $x$-axis and the $y$-axis in
    \( \Re^2 \) shows that \( W_1\directsum W_2=V \) does not imply that
    \( W_1\union W_2=V \).
    Can \( W_1\directsum W_2=V \) and \( W_1\union W_2=V \) happen?
    \begin{answer}
      It happens when at least one of \( W_1,W_2 \) is trivial.
      But that is the only way it can happen.

      To prove this, assume that both are non-trivial,
      select nonzero vectors \( \vec{w}_1, \vec{w}_2 \) from each,
      and consider \( \vec{w}_1+\vec{w}_2 \).
      This sum is not in \( W_1 \) because
      \( \vec{w}_1+\vec{w}_2=\vec{v}\in W_1 \) would imply that
      \( \vec{w}_2=\vec{v}-\vec{w}_1 \) is in \( W_1 \),
      which violates the assumption of the independence of the subspaces.
      Similarly, \( \vec{w}_1+\vec{w}_2 \) is not in \( W_2 \).
      Thus there is an element of \( V \) that is not in \( W_1\union W_2 \).
     \end{answer}
%   \recommended \item
%     Our model for complementary subspaces, the $x$-axis and the
%     $y$-axis in \( \Re^2 \), has one property not used here.
%     Where \( U \) is a subspace of \( \Re^n \) we define the
%     \definend{orthocomplement}\index{subspace!orthocomplement} 
%     of \( U \) to be
%     \begin{equation*}
%       U^\perp=
%       \set{\vec{v}\in\Re^n \suchthat \text{\( \vec{v}\dotprod\vec{u}=0 \)
%                                             for all \( \vec{u}\in U \)}}
%     \end{equation*}
%     (read ``\( U \) perp'').
%     \begin{exparts}
%       \partsitem Find the orthocomplement of the \( x \)-axis in \( \Re^2 \).
%       \partsitem Find the orthocomplement of the \( x \)-axis in \( \Re^3 \).
%       \partsitem Find the orthocomplement of the \( xy \)-plane in \( \Re^3 \).
%       \partsitem Show that the orthocomplement of a subspace is a subspace.
%       \partsitem Show that if \( W \) is the orthocomplement of \( U \) then
%         \( U \) is the orthocomplement of \( W \).
%       \partsitem Prove that a subspace and its orthocomplement have a trivial
%         intersection.
%       \partsitem Conclude that for any \( n \) and subspace 
%         \( U\subseteq \Re^n \) we have that \( \Re^n=U\directsum U^\perp \).
%       \partsitem Show that 
%         \( \dim(U)+\dim(U^\perp) \) equals the dimension of the
%         enclosing space.
%     \end{exparts}
%     \begin{answer}
%       \begin{exparts}
%         \partsitem The set 
%           \begin{equation*}
%             \set{\colvec{v_1 \\ v_2} 
%                \suchthat \colvec{v_1 \\ v_2}\dotprod\colvec{x \\ 0}=0 
%                               \text{\ for all $x\in\Re$}}
%           \end{equation*}
%           is easily seen to be the $y$-axis.
%         \partsitem The \( yz \)-plane.
%         \partsitem The \( z \)-axis.
%         \partsitem 
%           Assume that \( U \) is a subspace of some \( \Re^n \).
%           Because $U^\perp$ contains the zero vector, since that
%           vector is perpendicular to everything, we need only show that the
%           orthocomplement is closed under linear combinations of two elements.
%           If \( \vec{w}_1, \vec{w}_2\in U^\perp \) then
%           \( \vec{w}_1\dotprod\vec{u}=0 \) and \( \vec{w}_2\dotprod\vec{u}=0 \)
%           for all \( \vec{u}\in U \).
%           Thus \( (c_1\vec{w}_1+c_2\vec{w}_2)\dotprod\vec{u}=
%               c_1(\vec{w}_1\dotprod\vec{u})+c_2(\vec{w}_2\dotprod\vec{u})=0 \)
%           for all \( \vec{u}\in U \) and so \( U^\perp \) is closed under
%           linear combinations.
%         \partsitem The only vector orthogonal to itself is the zero vector.
%         \partsitem This is immediate.
%         \partsitem To prove that the dimensions add, 
%           it suffices by \nearbycorollary{cor:DirSumDimsAdd} and 
%           \nearbylemma{le:DirectSumTwoSp} to show that 
%           \( U\intersection U^\perp \) is the trivial subspace
%           \( \set{\zero} \).
%           But this is one of the prior items in this problem.
%       \end{exparts}  
%      \end{answer}
  \item 
    Consider \nearbycorollary{cor:DirSumDimsAdd}.
    Does it work both ways\Dash that is,
    supposing that \( V=W_1+\dots+ W_k \),
    is \( V=W_1\directsum\cdots\directsum W_k \) if and only if
    \( \dim(V)=\dim(W_1)+\dots+\dim(W_k) \)?
    \begin{answer}
      Yes.
      The left-to-right implication is
      \nearbycorollary{cor:DirSumDimsAdd}.
      For the other direction,
      assume that \( \dim(V)=\dim(W_1)+\dots+\dim(W_k) \).
      Let \( B_1,\dots, B_k \) be bases for \( W_1,\dots, W_k \).
      As $V$ is the sum of the subspaces, we can write 
      any \( \vec{v}\in V \) as
      \( \vec{v}=\vec{w}_1+\cdots+\vec{w}_k \) and expressing each
      \( \vec{w}_i \) as a combination of vectors from the associated
      basis \( B_i \) shows that the concatenation
      \( \cat{B_1}{\cat{\cdots}{B_k}}  \) spans \( V \).
      Now, that concatenation has
      \( \dim(W_1)+\dots+\dim(W_k) \) members, and so it is a spanning set of 
      size \( \dim(V) \).
      The concatenation is therefore a basis for \( V \).
      Thus \( V \) is the direct sum.
     \end{answer}
  \item 
    We know that if \( V=W_1\directsum W_2 \) then there is a basis for
    \( V \) that splits into a basis for \( W_1 \) and a basis for
    \( W_2 \).
    Can we make the stronger statement that every basis for \( V \) splits into
    a basis for \( W_1 \) and a basis for \( W_2 \)?
    \begin{answer}
      No.
      The standard basis for \( \Re^2 \) does not split into bases for
      the complementary subspaces the line \( x=y \) and the line 
      \( x=-y \).  
    \end{answer}
  \item 
     We can ask about the algebra of the `$+$' operation.
     \begin{exparts}
       \partsitem Is it commutative; is \( W_1+W_2=W_2+W_1 \)?
       \partsitem Is it associative; is \( (W_1+W_2)+W_3=W_1+(W_2+W_3) \)?
       \partsitem Let \( W \) be a subspace of some vector space.
         Show that \( W+W=W \).
       \partsitem Must there be an identity element,
         a subspace \( I \) such that \( I+W=W+I=W \) 
         for all subspaces \( W \)? 
       \partsitem Does left-cancellation hold:~if 
          \( W_1+W_2=W_1+W_3 \) then \( W_2=W_3 \)?
          Right cancellation?
    \end{exparts}
    \begin{answer}
      \begin{exparts}
        \partsitem  Yes, \( W_1+W_2=W_2+W_1 \) for all subspaces \( W_1,W_2 \)
          because each side is the span of \( W_1\union W_2=W_2\union W_1 \).
        \partsitem This one is similar to the prior one\Dash each side 
          of that equation is the span of
          \( (W_1\union W_2)\union W_3=W_1\union (W_2\union W_3) \).
        \partsitem Because this is an equality between sets, we can show that
          it holds by mutual inclusion. 
          Clearly \( W\subseteq W+W \).
          For \( W+W\subseteq W \) just recall that every subset is closed 
          under addition so any sum of the form \( \vec{w}_1+\vec{w}_2 \) is in
          \( W \).
        \partsitem In each vector space, the identity element with respect 
          to subspace addition is the trivial subspace.
        \partsitem   Neither of left or right cancellation needs to hold.
          For an example, in \( \Re^3 \) take \( W_1 \) to be the 
          \( xy \)-plane, take \( W_2 \) to be the \( x \)-axis, 
          and take \( W_3 \) to be the \( y \)-axis.  
      \end{exparts}
     \end{answer}
  % 2014-Dec-19 I commented this out to make the pages fit better; it is perfectly fine
  % \item 
  %   Consider the algebraic properties of the direct sum operation.
  %   \begin{exparts}
  %      \partsitem Does direct sum commute: does \( V=W_1\directsum W_2 \) imply
  %        that \( V=W_2\directsum W_1 \)?
  %     \partsitem Prove that direct sum is associative:
  %       \( (W_1\directsum W_2)\directsum W_3=
  %           W_1\directsum(W_2\directsum W_3) \).
  %      \partsitem Show that \( \Re^3 \) is the direct sum of the three axes
  %        (the relevance here is that by the previous item,
  %         we needn't specify which two of the three axes are combined first).
  %      \partsitem Does the direct sum operation left-cancel:~does
  %        \( W_1\directsum W_2=W_1\directsum W_3 \) imply \( W_2=W_3 \)?
  %        Does it right-cancel?
  %      \partsitem There is an identity element with respect to this operation.
  %        Find it.
  %      \partsitem Do some, or all, subspaces have inverses with respect to this
  %        operation:~is there a subspace \( W \) of some vector space such
  %        that there is a subspace \( U \) with the property that
  %        \( U\directsum W \) equals the identity element from
  %        the prior item?
  %   \end{exparts}
  %   \begin{answer}
  %     \begin{exparts}
  %        \partsitem They are equal because for each, 
  %          \( V \) is the direct sum if
  %          and only if we can write each \( \vec{v}\in V \) in a unique
  %          way as a sum \( \vec{v}=\vec{w}_1+\vec{w}_2 \) and
  %          \( \vec{v}=\vec{w}_2+\vec{w}_1 \).
  %        \partsitem They are equal because for each, 
  %          \( V \) is the direct sum if
  %          and only if we can write each \( \vec{v}\in V \) in a unique
  %          way as a sum of a vector from each
  %          $\vec{v}=(\vec{w}_1+\vec{w}_2)+\vec{w}_3$
  %          and $\vec{v}=\vec{w}_1+(\vec{w}_2+\vec{w}_3)$.
  %        \partsitem We can decompose any vector in \( \Re^3 \) uniquely into
  %          the sum of a vector from each axis.
  %        \partsitem No.
  %          For an example, in \( \Re^2 \) take \( W_1 \) to be the
  %          \( x \)-axis, take \( W_2 \) to be the \( y \)-axis, and
  %          take \( W_3 \) to be the line \( y=x \).
  %        \partsitem In any vector space the trivial subspace acts as 
  %          the identity element with respect to direct sum.
  %        \partsitem In any vector space, only the trivial subspace has
  %          a direct-sum inverse (namely, itself).
  %          One way to see this is that dimensions add, and so increase.
  %     \end{exparts}  
  %    \end{answer}
\end{exercises}
\index{direct sum|)}

% END NATIVE SOURCE src/vs/vs3.tex

% BEGIN NATIVE SOURCE src/vs/fields.tex
% Chapter 2, Topic _Linear Algebra_ Jim Hefferon
%  http://joshua.smcvt.edu/linearalgebra
%  2001-Jun-11
\topic{Fields}
\index{field|(}
Computations involving only integers or only rational numbers are much
easier than those with real numbers.
Could other algebraic structures, such as the integers or the rationals,
work in the
place of \( \Re \) in the definition of a vector space?

If we take ``work'' to mean that the results of this chapter
remain true
then there is a natural list of conditions that a structure 
(that is, number system)
must have in order to work in the place of $\Re$.
A \definend{field}\index{field!definition} is a set 
\( \F \) with operations
`\( + \)'
and `\( \cdot \)' such that
\begin{tfae}
   \item for any \( a,b\in\F \) the result of \( a+b \) is in \( \F \), and
         \( a+b=b+a \), and
         if \( c\in\F \) then \( a+(b+c)=(a+b)+c \)
   \item for any \( a,b\in\F \) the result of \( a\cdot b \) is in \( \F \), and
         \( a\cdot b=b\cdot a \), and
         if \( c\in\F \) then \( a\cdot (b\cdot c)=(a\cdot b)\cdot c \)
   \item if \( a,b,c\in\F \) then \( a\cdot (b+c)=a\cdot b+a\cdot c \)
   \item there is an element \( 0\in\F \) such that
         if \( a\in\F \) then \( a+0=a \), and
         for each \( a\in\F \) there is an element \( -a\in\F \)
                such that \( (-a)+a=0 \)
   \item there is an element \( 1\in\F \) such that
         if \( a\in\F \) then \( a\cdot 1=a \), and
         for each element \( a\neq 0 \) of \( \F \)
                there is an element \( a^{-1}\in\F \)
                such that \( a^{-1}\cdot a=1 \).
\end{tfae}

For example, the algebraic structure 
consisting of the set of real numbers along with its usual
addition and multiplication operation is a field.
Another field is the set of rational numbers with its usual addition
and multiplication operations.
An example of an algebraic structure that is not a field is
the integers, because it fails the final condition.

Some examples are more surprising.
The set \( \mathbb{B}=\set{0,1} \) under these operations:
\begin{center}
  \begin{tabular}{c|cc}
    \( + \) &\( 0 \) &\( 1 \) \\
    \hline
    \( 0 \) &\( 0 \) &\( 1 \) \\
    \( 1 \) &\( 1 \) &\( 0 \)
  \end{tabular}
  \qquad
  \begin{tabular}{c|cc}
    \( \cdot \) &\( 0 \) &\( 1 \) \\
     \hline
       \( 0 \)  &\( 0 \) &\( 0 \)  \\
       \( 1 \)  &\( 0 \) &\( 1 \)
  \end{tabular}
\end{center}
is a field; see \nearbyexercise{exer:BinField}.

We could in this book develop Linear Algebra as the theory of
vector spaces with scalars from an arbitrary field.
In that case, 
almost all of the statements here would carry over by replacing
`\( \Re \)' with `\( \F \)', that is, by
taking coefficients, vector entries,
and matrix entries to be elements of \( \F \)
(the exceptions are statements involving distances or angles,
which would need additional development).
Here are some examples; each applies to a vector space \( V \)
over a field \( \F \).
\begin{itemize}
  %\makebox[\parindent]{\ \hfil\ }
  \item[$*$] For any \( \vec{v}\in V \) and \( a\in\F \),
  (i)~\( 0\cdot\vec{v}=\zero \),
  (ii)~\( -1\cdot\vec{v}+\vec{v}=\zero \),
  and (iii)~\( a\cdot\vec{0}=\vec{0} \).

  \item[$*$] The span, the set of linear combinations, of a subset of \( V \)
  is a subspace of \( V \).

  \item[$*$] Any subset of a linearly independent set is also 
     linearly independent.

  \item[$*$] In a finite-dimensional vector space, 
    any two bases have the same number of elements.
\end{itemize}
(Even statements that don't explicitly mention \( \F \) use
field properties in their proof.)

We will not develop vector spaces in this more general setting because
the additional abstraction can be a distraction.
The ideas we want to bring out already appear when we stick to the reals.

The exception is Chapter Five.
There we must factor polynomials,
so we will switch to considering vector spaces over the
field of complex numbers.

\begin{exercises}
  \item 
    Check that the real numbers form a field.
    \begin{answer}
      Going through the five conditions shows that they are all familiar from 
      elementary mathematics.
    \end{answer}
  \item 
    Prove that these are fields.
    \begin{exparts*}
       \partsitem The rational numbers $\Q$
       \partsitem The complex numbers  $\C$
    \end{exparts*}
    \begin{answer}
      As with the prior question, going through the five conditions 
      shows that
      for both of these structures, 
      the properties are familiar.
    \end{answer}
  \item 
     Give an example that shows that the integer number system
     is not a field.
     \begin{answer}
       The integers fail condition~(5).
       For instance, there is no multiplicative inverse for $2$\Dash while
       $2$ is an integer, $1/2$ is not.
     \end{answer}
  \item \label{exer:BinField} 
     Check that the set $\mathbb{B}=\set{0,1}$ is a field under the operations 
     listed above, 
     \begin{answer}
       We can do these checks by listing all of the possibilities.
       For instance, to verify the first half of condition~(2) we must check
       that the structure is closed under addition and that addition
       is commutative  $a+b=b+a$,
       we can check both of these 
       for all possible pairs $a$ and~$b$ because there
       are only four such pairs.  
       Similarly, for associativity, there are only eight triples $a$, $b$,
       $c$, and so the check is not too long.
       (There are other ways to do the checks; in particular, you may 
       recognize these operations as arithmetic modulo~$2$.
       But an exhaustive check is not onerous)
     \end{answer}
  \item 
     Give suitable operations to make the set $\set{0,1,2}$
     a field.     
     \begin{answer}
       These will do.
       \begin{center}
          \begin{tabular}{c|ccc}
             \( + \) &\( 0 \) &\( 1 \) &$2$ \\
             \hline
             \( 0 \) &\( 0 \) &\( 1 \) &$2$ \\
             \( 1 \) &\( 1 \) &\( 2 \) &$0$ \\
             $2$     &$2$     &$0$     &$1$
          \end{tabular}
          \qquad
          \begin{tabular}{c|ccc}
            \( \cdot \) &\( 0 \) &\( 1 \) &$2$ \\
            \hline
            \( 0 \)  &\( 0 \) &\( 0 \) &$0$ \\
            \( 1 \)  &\( 0 \) &\( 1 \) &$2$ \\
            $2$      &$0$     &$2$     &$1$
          \end{tabular}
       \end{center}
       As in the prior item, we could 
       verify that they satisfy the conditions  
       by listing all of the cases.
     \end{answer}
\end{exercises}
\index{field|)}
%\include{la}
%\end{document}

% END NATIVE SOURCE src/vs/fields.tex

\appendix
\chapter{Separate editorial source notes}
\section*{Vector Spaces: source notes}
\begin{enumerate}
\item In the proof of scalar identity for matrices, the final matrix contains sa, sb, sc, sd. Those entries must be a, b, c, d: the scalar is 1, not a free s.
\item In the closure-under-addition display for upper-triangular matrices, an equals sign is missing between the second addend and the resulting sum matrix.
\item For the operation r·(x,y)=(rx,0), the preceding distributivity counterexample does not work: both sides equal (0,0). The separately stated scalar-identity counterexample is valid and the answer “not a vector space” remains correct.
\item The requested arrowed-operation restatement is inconsistent in axioms (6)–(8): scalar multiplication still uses plain · in (6)–(7), and vector addition uses plain + in (8). These should use the newly defined arrowed scalar and vector operations; scalar addition r+s remains plain.
\item The span-membership row reduction omits a step. The two labelled row operations produce third row (0, 3/2, 7/2), not (0,0,3). Adding the resulting second row to the third then gives (0,0,3), so the “no solution” conclusion is unchanged.
\item In the three-dimensional superhero answer, the three device vectors span all of R³, not a plane. Their column matrix has determinant 18. The “no place to hide” conclusion and the displayed elimination are correct.
\item The span-extension proof should assume v belongs to span(S), not necessarily S. Its displayed linear-combination representation is exactly the weaker hypothesis needed by the proof.
\item The intersection proof changes its vector names: it assumes w and s belong to the intersection, then forms r v+s w. Use v and w consistently as vectors, with r and s as scalars.
\item The union-of-subspaces counterexample names ambient space R³ but displays two-entry coordinate vectors. Either use R² as the ambient space or append a zero third coordinate to each vector.
\item The condition given for span(S∪T)=span(S)∪span(T) is too strong. Equality holds exactly when one span contains the other, not only when one original set contains the other. For S=\{(1,0)\} and T=\{(2,0)\}, neither set contains the other but both spans are the same line. The following noncontainment argument confuses membership in a set with membership in its span.
\end{enumerate}
\section*{Linear Independence: source notes}
\begin{enumerate}
\item In the \{1+x,1-x\} example, R2 - R1 yields -2c2 = 0, not +2c2 = 0. The conclusion of independence is unchanged.
\item The general echelon-matrix statement needs the qualification nonzero rows. An echelon matrix with a zero row supplies an immediate counterexample. The displayed example itself has no zero row.
\item When isolating v from 0 = c1 s1 + ... + cn sn + c(n+1) v, each displayed coefficient of si needs a minus sign. The span-membership conclusion remains valid.
\item The first equation in the spanning-set reduction contains a plus sign in the empty c4 column, printing two successive plus signs. The intended first equation is c1 + c3 + 3c5 = 0.
\item In the polynomial exercise answer, the stated row operations give the last pivot -152/13, not -128/13. The independence conclusion remains correct.
\item In the three-row-vector exercise, the operation producing the zero third row must be R3 - (4/7)R2, not R3 + (4/7)R2. The final dependence vector is correct.
\item The union example defines S and T but its last sentence calls them S1 and S2. These are the same displayed sets; their union is dependent.
\item The question labels its proposed direction as if, but the proposed implication from an independent union to trivial intersection is the only-if direction. This note corrects the direction label only; it supplies no answer.
\item The supplied answer reverses the names if and only if: trivial intersection implies an independent union is the if direction; the reverse is only if. The corrected criterion and proof substance are unchanged.
\item The two-by-two determinant argument needs unique solution, not solution; every homogeneous system has a zero solution. Its combined fraction also needs denominator a, not d.
\item The three-by-three elimination step needs the negative multiplier -((ah-bg)/a) on row 2 to clear row 3. The displayed final pivot is the result of this subtraction.
\item In the a != 0, ae-bd = 0 case, the proof describes singularity while calling it nonsingularity. After the row swap the system is nonsingular exactly when both ah-bg and af-cd are nonzero, equivalently when the determinant is nonzero. Either factor zero, and the subsequent equal-zero calculation, establish singularity.
\item For the two vectors in R3, dependence requires ae-bd = ah-gb = dh-ge = 0. Equality of these three minors without the final zero condition is insufficient.
\end{enumerate}
\section*{Basis and Dimension: source notes}
\begin{enumerate}
\item The first polynomial-basis answer omits the contributions of c1 to the x and constant coefficients. Its proposed c2 and c3 therefore do not represent a general quadratic when a2 is nonzero. The stated set is still a basis. Correct coefficients (c1,c2,c3) are (a2; a0/2 + a1/4 - a2/4; -a0/2 + a1/4 + 3*a2/4).
\item The matrix-basis proof obtains c3=0 from the lower-right entries, not the lower-left entries as stated at line 918. The displayed third basis matrix has its sole nonzero entry at the lower right. The conclusion remains correct.
\item Choose s2 in S outside span(B1), and repeat that choice within S. Such an element exists whenever span(B1) is a proper subspace of span(S). B1 is initially a basis of its own span, not necessarily the full space.
\item Read the intermediate set as \{-1+x, 2x\textasciicircum{}2+x\textasciicircum{}3\}. The printed x\textasciicircum{}2+x\textasciicircum{}3 fails the exercise condition a2-2a3=0; the final displayed basis and dimension are correct.
\item The two displayed sets describe the same diagonal-matrix space; an equality sign is missing at the start of the second line. Preserve the original display and identify this separately.
\item Use a nontrivial relation among the concatenated bases. Independence within each basis implies that the partial combination from either basis is nonzero. These equal, opposite-signed partial combinations provide a nonzero vector in U intersect W; no individual basis vector need lie in the other space.
\item Read V=W in the first sentence of this part as U=W. The remainder of the argument already uses U and W.
\item The right-hand column in the first displayed equation should be (1,3), as in the question and subsequent system. The conclusion that it is outside the column space remains correct.
\item At the second arrow, read the row-four operation as -2 times row one plus row four. The displayed echelon matrix then follows. The later basis ending in -5x\textasciicircum{}2 is valid: scaling the nonzero third basis vector does not change its span.
\item Read “at least two” here as “at most two”. The fourth and fifth columns are independent, so the column rank is exactly two, as stated.
\item For unequal row counts, first append zero rows to the shorter matrix, or compare only the nonzero reduced rows. The exercise claims equality of nonzero rows, not equality of differently sized matrices. A proof can stack A over B and B over A: the two stacked matrices are row equivalent, and reducing either dependent extra block leaves the same nonzero rows.
\item In part (b), read the right-hand entries as d1 through dm. There is one right-hand entry for each of the m equations.
\item For the missing direction, if every right-hand vector in R\textasciicircum{}m is attainable, the columns span R\textasciicircum{}m. The column rank is then m, and equality of row and column ranks gives full row rank. The second printed paragraph is the contrapositive of the first direction.
\item In part (e), W1 is a subspace: shifting its parameter still describes the entire x-axis. W2 is not a subspace because it does not contain zero. The answer “No” remains correct.
\item Read “greater than” here as “at least”. Dimensions eight, nine and ten for the sum, and eight, seven and six for the intersection, are all possible. The printed final list is correct.
\item Read “every subset” as “every subspace”. Closure under addition is available because W is a subspace; the identity W+W=W is correct.
\end{enumerate}
\section*{Fields: source notes}
\begin{enumerate}
\item The displayed axiom list does not explicitly require the additive and multiplicative identities to be distinct. Its literal conditions also hold in the singleton algebra with 0=1; the usual field definition excludes this case. Record the missing explicit condition separately, without changing the pinned source.
\item The supplied binary-field answer calls addition closure and commutativity the first half of condition (2), but those appear in condition (1). The operations and the finite-field conclusion are correct.
\end{enumerate}

\end{document}
